% Created 2026-11-11 Wed 09:46 with Tuatara, template article
% Intended LaTeX compiler: xelatex
%
% The enhanced export's own look, for a note that sets none of its own
% (no #+latex_template:, #+setupfile: or #+latex_header:, and #+latex_class:
% none or one of LaTeX's own, article, report or book): A4
% with even margins, fonts found on this system (*latex-main-font* and the
% like choose others), a font for every character the others lack,
% finer spacing, paragraphs apart rather than indented, black links.
\documentclass[11pt]{article}
\usepackage[a4paper, top=25mm, bottom=25mm, left=28mm, right=28mm]{geometry}
\usepackage{graphicx}
\usepackage{longtable}
\usepackage{wrapfig}
\usepackage{rotating}
\usepackage[normalem]{ulem}
\usepackage{amsmath}
\usepackage{amssymb}
\usepackage{fontspec}
\usepackage{capt-of}
\usepackage{hyperref}
\usepackage{booktabs}
\usepackage{xltabular}
\usepackage{multicol}
\usepackage{float}
\usepackage{xcolor}
\usepackage{fvextra}
\usepackage{pict2e}
\makeatletter
\newcommand\tuatara@turn{\clearpage
  \setlength{\@tempdima}{\paperwidth}\setlength{\paperwidth}{\paperheight}\setlength{\paperheight}{\@tempdima}%
  \ifdefined\pdfpagewidth\pdfpagewidth=\paperwidth\pdfpageheight=\paperheight\fi
  \ifdefined\pagewidth\pagewidth=\paperwidth\pageheight=\paperheight\fi
  \setlength{\@tempdimb}{\paperwidth}\addtolength{\@tempdimb}{-\paperheight}%
  \addtolength{\textwidth}{\@tempdimb}\addtolength{\textheight}{-\@tempdimb}%
  \setlength{\hsize}{\textwidth}\setlength{\columnwidth}{\textwidth}\setlength{\linewidth}{\textwidth}%
  \setlength{\vsize}{\textheight}\setlength{\@colht}{\textheight}\setlength{\@colroom}{\textheight}%
  \ifdefined\headwidth\setlength{\headwidth}{\textwidth}\fi}
\newenvironment{tuataralandscape}{\tuatara@turn}{\tuatara@turn}
\makeatother
\DeclareRobustCommand{\redactedbox}[1]{{\sbox0{#1}\rule[-0.8ex]{\wd0}{2.5ex}}}
\DeclareRobustCommand{\breakmark}{\hbox{\setlength{\unitlength}{0.07em}\linethickness{0.035em}\begin{picture}(10,4)(0,-0.5)\qbezier(2,3)(0.5,3)(0.5,2)\qbezier(0.5,2)(0.5,1)(2,1)\put(2,1){\line(1,0){7.5}}\put(9.5,1){\line(-3,2){2}}\put(9.5,1){\line(-3,-2){2}}\end{picture}}}
\fvset{breaksymbolleft=\breakmark}
\DeclareRobustCommand{\codebreak}{\discretionary{}{\breakmark}{}}
\pdfstringdefDisableCommands{\let\codebreak\relax}
\setlength{\emergencystretch}{3em}
\tolerance=1000
\setmainfont{Liberation Serif}
\setsansfont{Liberation Sans}
\setmonofont[Scale=MatchLowercase,Contextuals=NoAlternate,Ligatures=NoCommon]{JetBrains Mono}
\newfontfamily\tuataraFallbackONE{Noto Sans Symbols 2}[BoldFont={Noto Sans Symbols 2},ItalicFont={Noto Sans Symbols 2},BoldItalicFont={Noto Sans Symbols 2}]
\newfontfamily\tuataraFallbackTWO{Noto Sans Math}[BoldFont={Noto Sans Math},ItalicFont={Noto Sans Math},BoldItalicFont={Noto Sans Math}]
\newfontfamily\tuataraFallbackTHREE{DejaVu Sans}[BoldFont={DejaVu Sans},ItalicFont={DejaVu Sans},BoldItalicFont={DejaVu Sans}]
\newfontfamily\tuataraFallbackFOUR{Cascadia Code NF}[BoldFont={Cascadia Code NF},ItalicFont={Cascadia Code NF},BoldItalicFont={Cascadia Code NF}]
\usepackage{microtype}
\usepackage{parskip}
\usepackage[most]{tcolorbox}
\newenvironment{try}{\begin{tcolorbox}[breakable,enhanced,colback=black!4,boxrule=0pt,arc=0pt]}{\end{tcolorbox}}
\usepackage{etoolbox}
\pretocmd{\section}{\clearpage}{}{}
\hypersetup{
 colorlinks=true,
 linkcolor=black,
 urlcolor=black,
 citecolor=black,
 pdfauthor={},
 pdftitle={Tuatara Manual},
 pdfkeywords={},
 pdfsubject={},
 pdfcreator={Tuatara},
 pdflang={English}}
\author{}
\date{}
\title{Tuatara Manual}
\begin{document}

\maketitle
\setcounter{tocdepth}{2}
\tableofcontents


This manual is written with modal editing on, first and foremost. Where a command also has a key that works with modal editing off, it is given as well, as the non-modal key.

\section{The screen}
\label{sec:41}
What Tuatara shows you, and the name of each part.

When Tuatara starts, its screen is divided into a few parts. Each part has a
name, and the rest of this manual uses those names. This chapter goes through
them from the large ones to the small ones. For each part it lists the keys
that act on it and the settings that change how it looks.

This manual is written with modal editing on, first and foremost. Tuatara has
it on until you turn it off, so the keys the manual gives first are those of
NORMAL state. Where a command also has a key that works with modal editing off,
the manual gives that too.

\subsection{Keys and names}
\label{sec:1}
Most commands have a key. A key written as \texttt{C-x} means hold down the Control
key and press x. A key written as \texttt{M-x} means hold down Alt (called Meta in
older manuals) and press x. On a terminal where Alt does not reach the
program, press Escape, let go, then press x. \texttt{S-} stands for Shift.

Several keys in a row, as in \texttt{C-x C-f}, are pressed one after the other. You
may keep Control held down through the whole sequence. A few keys have names
instead of letters: \texttt{RET} is Return (Enter), \texttt{TAB} is Tab, \texttt{SPC} is the space
bar, \texttt{DEL} is the Backspace key and \texttt{ESC} is Escape.

Every command also has a name, such as \texttt{find-\codebreak{}file}. You can run any command by
its name. Press \texttt{M-x}, type a few letters of the name, and press \texttt{RET}. The
list of choices that appears shows the keys of each command beside it. A
command that has no key can only be run this way, and the tables in this
manual say \texttt{M-x} and its name for it.

A setting is a name that starts and ends with a star, such as
\texttt{*mode-\codebreak{}line-\codebreak{}tuatara*}. Settings go in your init file, the file Tuatara reads
when it starts, as in \texttt{(setf *mode-\codebreak{}line-\codebreak{}tuatara* nil)}. Each table of settings
in this chapter gives the standard value, the one Tuatara has until you
change it.

Tuatara is a modal editor, which means a key can do different things
depending on a state the editor is in. There are two states you will meet
first. In NORMAL state, the letter keys are commands: \texttt{i} starts INSERT
state. In INSERT state, the letter keys write text. \texttt{C-g} or \texttt{ESC} takes you
back to NORMAL state. The mode line, described below, always tells you which
state you are in.

The tables of commands have two columns of keys. The Modal column holds the
keys of the modal layer, such as the letter keys of NORMAL state. They work only
with modal editing on. The Non-modal column holds the keys that work in both
states and with modal editing off: those with Control or Alt, the arrow keys
and the keys of a major mode. A side with no key says “none”. In the text, a
key that has a partner of the other kind is followed by it in brackets, as in
\texttt{j} (non-modal: \texttt{C-n}) and \texttt{C-n} (modal: \texttt{j}). To turn modal editing off, put
\texttt{(setf *modal-\codebreak{}editing* nil)} in your init file. Every key of the Non-modal
column then does what the table says, and typing inserts text everywhere.

Most letter keys of NORMAL state that have no Control or Alt key of their own
are also under \texttt{C-z}, in every state: \texttt{m w} selects a word, and \texttt{C-z m w} does
the same with modal editing off. The Non-modal column gives them, as \texttt{C-z m w}.
A command that has a key in both states, such as \texttt{w} (non-modal: \texttt{M-f}), is not
under \texttt{C-z} a second time.

\subsection{The parts of the screen}
\label{sec:4}
Figure 1.1 shows the whole editor with two windows and a question open, and
Figure 1.2 the same with no question open. Each part is named as this manual,
the help and the project’s support pages name it, so that a question about a
part can say which one it means.

\begin{figure}[H]
\caption{Figure 1.1. The parts of the frame while a question is open.}\label{org2}
\begin{Verbatim}[breaklines,breakautoindent=false,breakindent=1em,breakafter=/-.\,;:=,commandchars=\\\{\}]
 +- Frame ----------------------------------------------------------+
 |+-----------------------------+  +------------------------------+ |
 || eggs                        |  | Pancakes                     | | Window
 || bread                       |  | 2 eggs, 200 g flour, milk    | |
 || milk                        |  | Mix, rest, fry.              | |
 |+-----------------------------+  +------------------------------+ |
 || TUATARA[N] PD [-] list.txt  |  | TUATARA[N] PD [-] pancakes.. | | Mode line
 |+-----------------------------+  +------------------------------+ |
 || 3/12 Find file: ~/notes/sh                                    | | Minibuffer
 || shopping.txt                                                  | | Panel
 || sheep-notes.org                                               | |
 |+---------------------------------------------------------------+ |
 +------------------------------------------------------------------+
\end{Verbatim}
\end{figure}

\begin{figure}[H]
\caption{Figure 1.2. The frame with no question open.}\label{org3}
\begin{Verbatim}[breaklines,breakautoindent=false,breakindent=1em,breakafter=/-.\,;:=,commandchars=\\\{\}]
 +- Frame ----------------------------------------------------------+
 |+-----------------------------+  +------------------------------+ |
 || eggs                        |  | Pancakes                     | | Window
 || bread                       |  | 2 eggs, 200 g flour, milk    | |
 |+-----------------------------+  +------------------------------+ |
 || TUATARA[N] PD [-] list.txt  |  | TUATARA[N] PD [-] pancakes.. | | Mode line
 |+-----------------------------+  +------------------------------+ |
 || Wrote /home/you/list.txt                                      | | Echo area
 |+---------------------------------------------------------------+ |
 +------------------------------------------------------------------+
\end{Verbatim}
\end{figure}

The frame is everything Tuatara draws. In a window of its own it is that
window, and in a terminal it is the terminal’s screen.

A window is a part of the frame that shows one piece of text. Each window has
a mode line under it, a row that describes what the window shows.

The echo area is the last row of the frame. Tuatara says short things there,
such as \texttt{Wrote /home/\codebreak{}you/\codebreak{}list.\codebreak{}txt}.

The minibuffer is the row where Tuatara asks you a question. The question is
called the prompt, and what you type after it is your answer. While a question
is open, the minibuffer and the panel take the bottom of the frame in place of
the echo area: the minibuffer on top and the choices under it, or the
minibuffer on the last row when \texttt{*prompt-\codebreak{}above-\codebreak{}candidates*} is \texttt{nil}.

The panel is the rows at the bottom of the frame that list choices while a
question is open, and the keys you can press next after a key such as \texttt{C-x},
above the echo area. It is only there when it has something to say.

The cursor is the mark that shows where typed text goes. Other editors call
that place the point, and Tuatara’s help uses the word too.

The gutter is the narrow strip at the left edge of a window. It can hold line
numbers.

\subsection{Windows and buffers}
\label{sec:10}
A buffer is a piece of text Tuatara holds in memory. When you open a file,
Tuatara reads it into a buffer. Your changes go to the buffer, and the file on
disk changes only when you save with \texttt{C-x C-s}. Not every buffer has a file.
\texttt{*Messages*}, the log of what the echo area said, is a buffer. So is the list of
files you get from \texttt{find-\codebreak{}file}. Buffers that Tuatara makes itself have their
names between stars.

A window is a view of one buffer. Two windows may show the same buffer, each
at its own place, or different buffers. Closing a window never closes its
buffer, and killing a buffer closes only the windows that were made for it.

Tuatara puts new windows where you would expect. A help buffer opens in
another window. The \texttt{*Messages*} log opens across the bottom, 30\% of the
height, and takes the cursor so that you can read it; \texttt{q} closes it. A window
that already holds a buffer is never split into pieces smaller than
\texttt{*window-\codebreak{}min-\codebreak{}height*} rows or \texttt{*window-\codebreak{}min-\codebreak{}width*} columns.

\begin{try}
\textbf{Try it, with a shopping list and a recipe}
\begin{enumerate}
\item Press \texttt{C-x C-f}, type \texttt{\textasciitilde{}/shopping.\codebreak{}txt} and press \texttt{RET}. The file does not
   exist yet, so Tuatara makes an empty buffer with that name.
\item Press \texttt{i} to start INSERT state. Type \texttt{eggs}, \texttt{RET}, \texttt{bread}, \texttt{RET}, \texttt{milk}.
   Press \texttt{C-g} to go back to NORMAL state, then \texttt{C-x C-s} to save.
\item Press \texttt{C-x 3}. The window splits in two, side by side, and both show the
   shopping list. The left window keeps the cursor.
\item Press \texttt{C-x o} to move to the right window. Press \texttt{C-x C-f}, type
   \texttt{\textasciitilde{}/pancakes.\codebreak{}txt}, press \texttt{RET}, then press \texttt{i} and type a recipe.
\item Press \texttt{C-x o} to go back to the left window. Press \texttt{C-x 1}. The left window
   fills the frame again. The pancake buffer is still there: press \texttt{C-x b},
   type \texttt{pan} and press \texttt{RET} to bring it back.
\end{enumerate}
\end{try}

\subsubsection{Commands for windows}
\label{sec:5}
\begin{xltabular}{\textwidth}{@{}l l l >{\raggedright\arraybackslash\hsize=1.000\hsize}X@{}}
\toprule
\textbf{Modal} & \textbf{Non-modal} & \textbf{Command} & \textbf{What it does}\\
\midrule
\endhead
none & \texttt{C-x 2} & \texttt{split-\codebreak{}window-\codebreak{}below} & Two windows, one above the other. The upper one keeps the cursor. A number before it sets the upper window’s rows.\\
none & \texttt{C-x 3} & \texttt{split-\codebreak{}window-\codebreak{}right} & Two windows side by side. The left one keeps the cursor.\\
none & \texttt{C-x 0} & \texttt{delete-\codebreak{}window} & Closes this window. Its room goes to its neighbour.\\
\texttt{F} & \texttt{C-x 1} & \texttt{delete-\codebreak{}other-\codebreak{}windows} & Closes all other windows. \texttt{F} works in NORMAL state.\\
none & \texttt{C-x o} & \texttt{other-\codebreak{}window} & Moves to the next window, from the top left. A number moves that many windows on, and a negative one goes back.\\
none & \texttt{C-M-h} & \texttt{select-\codebreak{}window-\codebreak{}left} & Moves to the window on the left.\\
none & \texttt{C-M-l} & \texttt{select-\codebreak{}window-\codebreak{}right} & Moves to the window on the right.\\
none & \texttt{C-M-k} & \texttt{select-\codebreak{}window-\codebreak{}above} & Moves to the window above.\\
none & \texttt{C-M-j} & \texttt{select-\codebreak{}window-\codebreak{}below} & Moves to the window below.\\
none & \texttt{q}, \texttt{h} & \texttt{quit-\codebreak{}window} & Closes a window that was opened for a help buffer, a list or a log, or shows the buffer it had before. \texttt{h} works in lists.\\
none & \texttt{C-x 4 0} & \texttt{kill-\codebreak{}buffer-\codebreak{}and-\codebreak{}window} & Kills this buffer, asking first if it has unsaved changes, and closes its window.\\
none & \texttt{C-x 4 b} & \texttt{switch-\codebreak{}to-\codebreak{}buffer-\codebreak{}other-\codebreak{}window} & Asks for a buffer and shows it in another window.\\
none & \texttt{C-x 4 f}, \texttt{C-x 4 C-f} & \texttt{find-\codebreak{}file-\codebreak{}other-\codebreak{}window} & Opens a file in another window.\\
\bottomrule
\end{xltabular}

\subsubsection{Moving windows and changing their size}
\label{sec:6}
{\small
\begin{xltabular}{\textwidth}{@{}l >{\raggedright\arraybackslash\hsize=0.870\hsize}X l >{\raggedright\arraybackslash\hsize=1.130\hsize}X@{}}
\toprule
\textbf{Modal} & \textbf{Non-modal} & \textbf{Command} & \textbf{What it does}\\
\midrule
\endhead
none & \texttt{C-x w h} & \texttt{swap-\codebreak{}window-\codebreak{}left} & Swaps this window’s buffer with the one on the left. The cursor goes with the buffer.\\
none & \texttt{C-x w l} & \texttt{swap-\codebreak{}window-\codebreak{}right} & The same, with the window on the right.\\
none & \texttt{C-x w k} & \texttt{swap-\codebreak{}window-\codebreak{}above} & The same, with the window above.\\
none & \texttt{C-x w j} & \texttt{swap-\codebreak{}window-\codebreak{}below} & The same, with the window below.\\
none & \texttt{C-x w x} & \texttt{window-\codebreak{}swap-\codebreak{}states} & Swaps what this window shows with the next window.\\
none & \texttt{C-x w r} & \texttt{rotate-\codebreak{}windows} & Moves each buffer to the next window.\\
none & \texttt{C-x w R} & \texttt{rotate-\codebreak{}windows-\codebreak{}back} & Moves each buffer to the window before.\\
none & \texttt{C-x w .} & \texttt{rotate-\codebreak{}window-\codebreak{}layout-\codebreak{}clockwise} & Turns the arrangement of windows a quarter turn clockwise.\\
none & \texttt{C-x w ,} & \texttt{rotate-\codebreak{}window-\codebreak{}layout-\codebreak{}anticlockwise} & Turns the arrangement a quarter turn the other way.\\
none & \texttt{C-x w t} & \texttt{transpose-\codebreak{}window-\codebreak{}layout} & Windows that were stacked go side by side, and the other way round.\\
none & \texttt{M-x flip-\codebreak{}window-\codebreak{}layout-\codebreak{}horizontally} & \texttt{flip-\codebreak{}window-\codebreak{}layout-\codebreak{}horizontally} & Mirrors the arrangement from left to right.\\
none & \texttt{M-x flip-\codebreak{}window-\codebreak{}layout-\codebreak{}vertically} & \texttt{flip-\codebreak{}window-\codebreak{}layout-\codebreak{}vertically} & Mirrors the arrangement from top to bottom.\\
none & \texttt{C-x \textasciicircum{}} & \texttt{enlarge-\codebreak{}window} & Makes this window a row taller, or as many rows as the number you gave first.\\
none & \texttt{M-x shrink-\codebreak{}window} & \texttt{shrink-\codebreak{}window} & Makes this window a row shorter. Right after \texttt{C-x \textasciicircum{}}, \texttt{v} does it.\\
none & \texttt{C-x \}} & \texttt{enlarge-\codebreak{}window-\codebreak{}horizontally} & Makes this window a column wider.\\
none & \texttt{C-x \{} & \texttt{shrink-\codebreak{}window-\codebreak{}horizontally} & Makes this window a column narrower.\\
none & \texttt{C-x +} & \texttt{balance-\codebreak{}windows} & Gives all windows the same size.\\
none & \texttt{C-x -} & \texttt{shrink-\codebreak{}window-\codebreak{}if-\codebreak{}larger-\codebreak{}than-\codebreak{}buffer} & Shrinks the window to the height its text needs.\\
\bottomrule
\end{xltabular}
}

Each of the four keys that make a window larger or smaller says in the echo
area what share of the split every window has now, in per cent. Side by side
it reads \texttt{Widths, left to right: 52\%, 48\%}, and one above the other it reads
\texttt{Heights, top to bottom: 55\%, 45\%}. The shares may not add up to exactly 100,
as each is rounded. The echo area then lists the keys that repeat the change.
Right after one of the size commands, you can press a key alone to go on
resizing. \texttt{\}} and \texttt{>} make the window wider, \texttt{\{} and \texttt{<} narrower, \texttt{\textasciicircum{}} and \texttt{+}
taller, and \texttt{v} and \texttt{-} shorter. If the window has no neighbour in that
direction, the echo area says that it cannot be made larger or smaller. With
the mouse,
drag a mode line up or down, or the line between two windows left or right.

\subsubsection{Scrolling and finding out where buffers go}
\label{sec:7}
\begin{xltabular}{\textwidth}{@{}l >{\raggedright\arraybackslash\hsize=0.651\hsize}X l >{\raggedright\arraybackslash\hsize=1.349\hsize}X@{}}
\toprule
\textbf{Modal} & \textbf{Non-modal} & \textbf{Command} & \textbf{What it does}\\
\midrule
\endhead
none & \texttt{C-v}, \texttt{<next>} & \texttt{scroll-\codebreak{}up-\codebreak{}command} & Scrolls forward by nearly a screenful. A number scrolls that many rows.\\
none & \texttt{M-v}, \texttt{<prior>} & \texttt{scroll-\codebreak{}down-\codebreak{}command} & Scrolls back by nearly a screenful.\\
\texttt{z z} & \texttt{C-l} & \texttt{recenter} & Puts the cursor’s line in the middle of the window. \texttt{z z} works in NORMAL state.\\
none & \texttt{C-M-v} & \texttt{scroll-\codebreak{}other-\codebreak{}window} & Scrolls the next window forward.\\
none & \texttt{C-M-\codebreak{}S-v} & \texttt{scroll-\codebreak{}other-\codebreak{}window-\codebreak{}down} & Scrolls the next window back.\\
none & \texttt{<wheel-\codebreak{}up>}, \texttt{<wheel-\codebreak{}down>} & \texttt{scroll-\codebreak{}by-\codebreak{}wheel} & Scrolls the window under the mouse.\\
none & \texttt{M-x describe-\codebreak{}buffer-\codebreak{}places} & \texttt{describe-\codebreak{}buffer-\codebreak{}places} & Lists where each kind of buffer is shown.\\
\bottomrule
\end{xltabular}

\subsubsection{Commands for buffers}
\label{sec:8}
\begin{xltabular}{\textwidth}{@{}l l l >{\raggedright\arraybackslash\hsize=1.000\hsize}X@{}}
\toprule
\textbf{Modal} & \textbf{Non-modal} & \textbf{Command} & \textbf{What it does}\\
\midrule
\endhead
none & \texttt{C-x b} & \texttt{switch-\codebreak{}to-\codebreak{}buffer} & Asks for a buffer. A name that does not exist makes a new buffer. \texttt{RET} alone goes back to the last one.\\
none & \texttt{C-x x b} & \texttt{switch-\codebreak{}other-\codebreak{}buffer} & Goes to the buffer used before this one. Press it again to come back.\\
none & \texttt{C-x C-b} & \texttt{buffers} & Lists every buffer, grouped by project.\\
none & \texttt{C-x k} & \texttt{kill-\codebreak{}buffer} & Kills this buffer, asking first if it has unsaved changes.\\
none & \texttt{C-x C-s} & \texttt{save-\codebreak{}buffer} & Writes the buffer to its file.\\
\bottomrule
\end{xltabular}

\subsubsection{Settings}
\label{sec:9}
\begin{xltabular}{\textwidth}{@{}l l >{\raggedright\arraybackslash\hsize=1.000\hsize}X@{}}
\toprule
\textbf{Setting} & \textbf{Standard} & \textbf{What it changes}\\
\midrule
\endhead
\texttt{*window-\codebreak{}min-\codebreak{}height*} & \texttt{4} & The fewest rows, mode line included, a window may be left with when it is split or made smaller.\\
\texttt{*window-\codebreak{}min-\codebreak{}width*} & \texttt{10} & The fewest columns a window may be left with.\\
\texttt{*split-\codebreak{}height-\codebreak{}threshold*} & \texttt{80} & The rows a window needs before a buffer shown “in another window” splits it into one above the other. A lone window always splits when it has room.\\
\texttt{*split-\codebreak{}width-\codebreak{}threshold*} & \texttt{160} & The columns a window needs before it is split into two side by side for the same purpose.\\
\texttt{*help-\codebreak{}window-\codebreak{}select*} & \texttt{nil} & Set to \texttt{t} to move the cursor to a help window when it opens. \texttt{q} there closes it and goes back.\\
\texttt{*scroll-\codebreak{}margin*} & \texttt{0} & How many rows to keep between the cursor and the top or bottom of a window.\\
\texttt{*scroll-\codebreak{}recenter*} & \texttt{nil} & Set to \texttt{t} to put the cursor’s line in the middle of the window whenever it leaves the window or comes within the margin.\\
\bottomrule
\end{xltabular}

\subsection{Frames, windows of their own and terminals}
\label{sec:13}
Tuatara runs in two ways. Typed as \texttt{tuatara}, it runs in the terminal you
started it from. Typed as \texttt{tuatara --gui}, it opens a window of its own on the
desktop, drawn with GTK, the toolkit most Linux desktops use. Both have the
same keys, the same commands and the same parts. The window of its own has
things a terminal does not, such as pictures in every case, a font you can
change and ligatures where the font has them.

A frame is one such screen. You can have several frames. In a window of its
own each frame is a separate desktop window. A terminal has only one screen,
so its frames take turns: the one you work in is shown, and the others wait.
When frames take turns, the mode line starts the buffer’s name with the
frame’s name, F1, F2 and so on. Every frame has its own windows and its own
echo area. They all show the same buffers, so a change you make in one frame
is in the other at once.

Tuatara can also run in the background as a daemon, started with
\texttt{tuatara --daemon}, and be reached with \texttt{tuatara --client}, each client with
frames of its own. The sign D in the mode line, described below, says that the
editor you are in runs as a daemon.

\begin{try}
\textbf{Try it, with two frames}
\begin{enumerate}
\item In a terminal, start Tuatara with \texttt{tuatara \textasciitilde{}/shopping.\codebreak{}txt}.
\item Press \texttt{C-x 5 2}. A second frame opens showing the same file. In a terminal
   the new frame replaces the first on the screen, and the mode line now
   starts the buffer’s name with \texttt{F2}.
\item Press \texttt{C-x 5 o} to go to the other frame, and back again with the same
   keys. The mode line says \texttt{F1} in the first.
\item Press \texttt{C-x 5 0} in the second frame to close it. Closing the last frame is
   refused. Quit with \texttt{C-x C-c}.
\end{enumerate}
\end{try}

\subsubsection{Commands for frames and the window’s font}
\label{sec:11}
\begin{xltabular}{\textwidth}{@{}l >{\raggedright\arraybackslash\hsize=0.718\hsize}X l >{\raggedright\arraybackslash\hsize=1.282\hsize}X@{}}
\toprule
\textbf{Modal} & \textbf{Non-modal} & \textbf{Command} & \textbf{What it does}\\
\midrule
\endhead
none & \texttt{C-x 5 2} & \texttt{make-\codebreak{}frame} & Opens another frame showing this buffer.\\
none & \texttt{C-x 5 0} & \texttt{delete-\codebreak{}frame} & Closes this frame. The buffers stay. The last frame cannot be closed this way.\\
none & \texttt{C-x 5 1} & \texttt{delete-\codebreak{}other-\codebreak{}frames} & Closes all frames but this one.\\
none & \texttt{C-x 5 o} & \texttt{other-\codebreak{}frame} & Goes to the next frame. A number goes that many on, a negative one goes back.\\
none & \texttt{C-x 5 f} & \texttt{find-\codebreak{}file-\codebreak{}other-\codebreak{}frame} & Opens a file in another frame.\\
none & \texttt{C-x 5 b} & \texttt{switch-\codebreak{}to-\codebreak{}buffer-\codebreak{}other-\codebreak{}frame} & Shows a buffer in another frame.\\
none & \texttt{C-x 5 d} & \texttt{rimu-\codebreak{}other-\codebreak{}frame} & Lists a folder in another frame.\\
none & \texttt{C-x C-+}, \texttt{C-x C-=} & \texttt{text-\codebreak{}scale-\codebreak{}adjust} & Makes the font a step larger. After it, \texttt{+}, \texttt{=}, \texttt{-} and \texttt{0} alone go on.\\
none & \texttt{C-x C--} & \texttt{text-\codebreak{}scale-\codebreak{}adjust} & Makes the font a step smaller.\\
none & \texttt{C-x C-0} & \texttt{text-\codebreak{}scale-\codebreak{}adjust} & Goes back to the size the window started with.\\
none & \texttt{M-x text-\codebreak{}scale-\codebreak{}increase} & \texttt{text-\codebreak{}scale-\codebreak{}increase} & Makes the font a step larger.\\
none & \texttt{M-x text-\codebreak{}scale-\codebreak{}decrease} & \texttt{text-\codebreak{}scale-\codebreak{}decrease} & Makes the font a step smaller.\\
none & \texttt{C-<wheel-\codebreak{}up>}, \texttt{C-<wheel-\codebreak{}down>} & \texttt{mouse-\codebreak{}wheel-\codebreak{}text-\codebreak{}scale} & Larger or smaller as the wheel turns with Control held.\\
none & \texttt{M-x set-\codebreak{}frame-\codebreak{}font} & \texttt{set-\codebreak{}frame-\codebreak{}font} & Asks for a font family. It lasts until you quit.\\
none & \texttt{<f11>} & \texttt{toggle-\codebreak{}frame-\codebreak{}fullscreen} & Fills the screen, or goes back.\\
none & \texttt{M-<f10>} & \texttt{toggle-\codebreak{}frame-\codebreak{}maximized} & Makes the window as large as the screen allows, or goes back.\\
\bottomrule
\end{xltabular}

Font steps are 1.2 times each, from 0.4 to 15 times the starting size. In a
terminal the font belongs to the terminal, and these commands say so.

\subsubsection{Settings}
\label{sec:12}
\begin{xltabular}{\textwidth}{@{}l >{\raggedright\arraybackslash\hsize=0.711\hsize}X >{\raggedright\arraybackslash\hsize=1.289\hsize}X@{}}
\toprule
\textbf{Setting} & \textbf{Standard} & \textbf{What it changes}\\
\midrule
\endhead
\texttt{*startup-\codebreak{}buffer*} & \texttt{"*lisp-\codebreak{}scratch*"} & The buffer shown when Tuatara starts without files, either \texttt{"*lisp-\codebreak{}scratch*"} or \texttt{"*org-\codebreak{}scratch*"}.\\
\texttt{*auto-\codebreak{}scratch-\codebreak{}buffers*} & \texttt{("*lisp-\codebreak{}scratch*" "*org-\codebreak{}scratch*")} & The scratch buffers made without being asked for.\\
\texttt{*gui-\codebreak{}font*} & \texttt{nil} & The window’s font as Pango reads it, \texttt{"Adwaita Mono 12"} for one. \texttt{nil} takes the desktop’s monospace font.\\
\texttt{*gui-\codebreak{}padding*} & \texttt{8} & Pixels between the window’s edge and its text. One number, or a list of two, three or four, read as in CSS.\\
\texttt{*gui-\codebreak{}renderer*} & \texttt{:gl} & How the window is drawn: \texttt{:gl}, \texttt{:vulkan}, \texttt{:cairo}, or \texttt{nil} for GTK’s choice.\\
\texttt{*gui-\codebreak{}title*} & \texttt{"\textasciitilde{}a - Tuatara"} & The window’s title. \texttt{\textasciitilde{}a} is the name of the buffer shown.\\
\texttt{*window-\codebreak{}colours-\codebreak{}from-\codebreak{}desktop*} & \texttt{nil} & Set to \texttt{t} to take colours the theme leaves out from the desktop’s style.\\
\texttt{*tui-\codebreak{}padding*} & \texttt{0} & Cells between the terminal’s edge and the text, in the same forms as \texttt{*gui-\codebreak{}padding*}.\\
\texttt{*tui-\codebreak{}mouse*} & \texttt{t} & Whether the mouse works in a terminal. \texttt{nil} leaves it to the terminal.\\
\texttt{*tui-\codebreak{}colours*} & \texttt{:auto} & How many colours to use in a terminal: \texttt{:auto}, \texttt{:24-\codebreak{}bit}, \texttt{256}, \texttt{16}, \texttt{8} or \texttt{nil}.\\
\texttt{*tui-\codebreak{}pictures*} & \texttt{:auto} & Whether pictures show in a terminal: \texttt{:auto}, \texttt{t} or \texttt{nil}.\\
\bottomrule
\end{xltabular}

\subsection{The mode line}
\label{sec:24}
The mode line is the bar under every window. It says what the window shows
and in what state. The mode line of the window you are working in is drawn
brighter than the others. Here is a typical one.

\begin{figure}[H]
\caption{Figure 1.2. A mode line. The sign P is lit: a job is running.}
\begin{Verbatim}[breaklines,breakautoindent=false,breakindent=1em,breakafter=/-.\,;:=,commandchars=\\\{\}]
Dear Hana, thank you for the lovely weekend.
 TUATARA[N] PD [-] letter.txt  L1 C0  Text  ~/letters/
\end{Verbatim}
\end{figure}

Read left to right, the parts are these.

\begin{enumerate}
\item The modal state, in square brackets after the word TUATARA.
\item Up to four signs, P, U, S and D, always in the same place.
\item The marker that says whether the buffer is saved, [-], [*] or [\%].
\item Words about the buffer itself, such as Safe or Narrow, when they apply.
\item The buffer’s name.
\item The line and the column of the cursor.
\item The name of the buffer’s major mode, when it has one.
\item Anything more the buffer has to say, such as CRLF or Page 3 of 10.
\item The folder the buffer lives in, at the right end.
\end{enumerate}

Only the parts that apply are shown, and a mode line that is too narrow cuts
the buffer’s name in the middle, with a mark for what is left out.

\subsubsection{The modal state}
\label{sec:14}
The first part tells you what a key will do. \texttt{TUATARA[N]} is NORMAL state,
\texttt{TUATARA[I]} is INSERT state. \texttt{TUATARA[S]} is a support buffer, one that holds
a list, help or the \texttt{*Messages*} log. There you can read, move and select but
not type. \texttt{TUATARA[E]} is a buffer that holds no modal keys, such as the Lisp
prompt (IELM). With \texttt{*modal-\codebreak{}editing*} set to \texttt{nil} there is no
modal part at all.

With \texttt{*mode-\codebreak{}line-\codebreak{}tuatara*} set to \texttt{nil} the word TUATARA is left out and only
\texttt{[N]} remains. \texttt{C-x T t} switches it while Tuatara runs.

\subsubsection{The signs P, U, S and D}
\label{sec:15}
Four letters stand before the saved marker, side by side. Each is faint when
nothing is going on and bright, in the warning colour, while its work goes on.
They never move, so nothing after them jumps when one lights up. A sign stays
lit for half a second after its work ends, so that a short job can be seen.

\begin{xltabular}{\textwidth}{@{}l l >{\raggedright\arraybackslash\hsize=1.000\hsize}X@{}}
\toprule
\textbf{Sign} & \textbf{Setting} & \textbf{What it says}\\
\midrule
\endhead
P & \texttt{*processing-\codebreak{}indicator*} & Processing. A job is running for this buffer, or waits its turn. Jobs are things done in the background, such as a search through a large file, a sort or a save. \texttt{C-x j} lists them.\\
U & \texttt{*indexing-\codebreak{}indicator*} & Update. A folder you asked Tuatara to keep read in memory is being read.\\
S & \texttt{*searching-\codebreak{}indicator*} & Search. A search runs over such a folder.\\
D & \texttt{*daemon-\codebreak{}indicator*} & Daemon. This editor runs in the background as a daemon.\\
\bottomrule
\end{xltabular}

U and S only appear when \texttt{*indexed-\codebreak{}directories*} names some folder, and P and D
always appear. Set any of the four settings to \texttt{nil} to hide that sign, or to
another piece of text, such as \texttt{"indexing"} or an icon from a font that has
one, to show that instead.

\subsubsection{Saved or not}
\label{sec:16}
Next comes the marker for the buffer’s state.

\begin{xltabular}{\textwidth}{@{}l >{\raggedright\arraybackslash\hsize=1.000\hsize}X@{}}
\toprule
\textbf{Marker} & \textbf{Meaning}\\
\midrule
\endhead
{}[-] & The buffer is as it was saved or read.\\
{}[*] & You have changed it and not yet saved it.\\
{}[\%] & The buffer is read-only. Typing does not change it. \texttt{C-x C-q} switches that.\\
\bottomrule
\end{xltabular}

\subsubsection{Words about the buffer}
\label{sec:17}
Just before the buffer’s name, Tuatara may add a word that is lit up in a
colour of its own.

\begin{xltabular}{\textwidth}{@{}l >{\raggedright\arraybackslash\hsize=1.000\hsize}X@{}}
\toprule
\textbf{Word} & \textbf{Meaning}\\
\midrule
\endhead
Safe & The file was opened safely. Tuatara shows its text and leaves the file untouched.\\
Evidence & The file is in a folder listed in \texttt{*evidence-\codebreak{}folders*}. It is never written to.\\
Narrow & The buffer shows only a part of its text (\texttt{C-x x n}, \texttt{C-x x s}). \texttt{C-x x w} shows all of it again.\\
Look only & The buffer belongs to a git repository opened for looking at only. Commands that would change it are refused.\\
\bottomrule
\end{xltabular}

\subsubsection{The name, the line and the column}
\label{sec:18}
The buffer’s name follows. When it is long, it is cut in the middle and a
mark stands for what is left out, down to 12 columns. \texttt{*mode-\codebreak{}line-\codebreak{}name-\codebreak{}limit*}
sets a shorter limit.

After the name comes \texttt{L} and the line number, then \texttt{C} and the column, which
counts from 0. A number above 999 has commas, as in \texttt{L100,000,001}. When Tuatara
is still counting the lines of a very large file you see \texttt{L...} instead, and
the number appears when the count is done. The same holds for a column on a
line too long to measure at once.

After them, if the buffer has a major mode, comes its name, such as Org or
Terminal. A major mode is the set of habits Tuatara gives a buffer according to
what it holds: how it colours text, which keys it adds, and so on.

\subsubsection{Other things the mode line says}
\label{sec:19}
While keys are recorded, the word Rec stands before the signs (\texttt{C-x (}
starts recording). While recorded keys play over and over for more than a
second, \texttt{Played 1,200/10,000} says how many plays are done. The words below
follow the mode name, each when it applies.

\begin{xltabular}{\textwidth}{@{}>{\raggedright\arraybackslash\hsize=0.421\hsize}X >{\raggedright\arraybackslash\hsize=1.579\hsize}X@{}}
\toprule
\textbf{Shows} & \textbf{Meaning}\\
\midrule
\endhead
CRLF or CR & The file’s lines end the way Windows (CRLF) or old Macintosh files (CR) end them. Plain files with the usual ending show nothing.\\
Following & The buffer follows its file as it grows, like \texttt{tail -f} (\texttt{follow-\codebreak{}file}).\\
Counting lines 40\% & Tuatara is counting the lines of a large file.\\
Saving 50\% & A save is under way. Another buffer’s save shows with its name.\\
Changed on disk & Another program changed the file after you read it.\\
Deleted on disk & The file was removed.\\
Stale, read-only & The file was rewritten on disk while it was shown, so the buffer is read-only. \texttt{revert-\codebreak{}buffer} reads it again.\\
incomplete & A file from another machine stopped arriving before it was all there.\\
Wrap, Centred & The buffer wraps long lines (\texttt{C-x T l}) or shows its text in a column in the middle (\texttt{C-x T o}).\\
Rewind a1b2c3d & The buffer shows a file as an older git commit had it (\texttt{C-x G}). The commit’s date and its place in the history follow.\\
Conflicts 3 & The buffer holds git conflict markers, and says how many.\\
Hex 0x00000010 digits & The buffer shows the file’s bytes. \texttt{digits} or \texttt{text} says which side the cursor is on.\\
Page 3 of 10 & A PDF’s page, the number of pages and its zoom.\\
Sandboxed & What you see was made by a program shut in a sandbox.\\
Quick edit & A file opened from a folder list to be edited there.\\
Filtered by kiwi: 12 of 1,000 records & A search filter hides some lines of the file.\\
{\tuataraFallbackONE{}✘}2 {\tuataraFallbackONE{}⚠}1 & A language server found 2 errors and 1 warning.\\
Debug & Debug mode is on.\\
SV< & An input method is on, here the one for Swedish.\\
{}[0:22] (Task) & The Org clock is running on a task.\\
by time, by name & How a folder list is sorted.\\
\bottomrule
\end{xltabular}

\subsubsection{The folder}
\label{sec:20}
The last part is the folder the buffer lives in, written with \texttt{\textasciitilde{}} for your home
folder: \texttt{\textasciitilde{}/letters/}. A buffer with no file, such as a scratch buffer, shows the
folder named by \texttt{*default-\codebreak{}directory*}. Files on another machine show their
\texttt{/ssh:\codebreak{}pinetree:} path.

\begin{try}
\textbf{Try it, with a letter}
\begin{enumerate}
\item Press \texttt{C-x C-f}, type \texttt{\textasciitilde{}/letter.\codebreak{}txt} and press \texttt{RET}. Look at the mode line:
   \texttt{[-]}, the name, \texttt{L1 C0}.
\item Press \texttt{i}, type \texttt{Dear Hana,} and press \texttt{RET}. The mode line now says \texttt{[*]},
   and \texttt{TUATARA[I]}, and \texttt{L2 C0}.
\item Press \texttt{C-g} to leave INSERT state, then \texttt{C-x C-s}. The marker goes back to
   \texttt{[-]} and the echo area says that the file was written.
\item Press \texttt{C-x C-q}. The marker becomes \texttt{[\%]}. Now \texttt{C-x C-q} again.
\item Press \texttt{C-x T t}. The word TUATARA disappears from the mode line. Press the
   same keys to bring it back.
\item Press \texttt{M-x}, type \texttt{follow-\codebreak{}file} and press \texttt{RET}. The word \texttt{Following}
   appears. Run it again to stop it.
\end{enumerate}
\end{try}

\subsubsection{Commands}
\label{sec:21}
\begin{xltabular}{\textwidth}{@{}l l l >{\raggedright\arraybackslash\hsize=1.000\hsize}X@{}}
\toprule
\textbf{Modal} & \textbf{Non-modal} & \textbf{Command} & \textbf{What it does}\\
\midrule
\endhead
none & \texttt{C-x T t} & \texttt{toggle-\codebreak{}mode-\codebreak{}line-\codebreak{}tuatara} & Writes TUATARA before the modal state, or only the state.\\
none & \texttt{C-x C-q} & \texttt{read-\codebreak{}only-\codebreak{}mode} & Makes the buffer read-only, or editable again.\\
none & \texttt{C-x j} & \texttt{list-\codebreak{}jobs} & Lists the jobs running or waiting, which P stands for.\\
none & \texttt{C-x J} & \texttt{stop-\codebreak{}job} & Stops the job running in this buffer, after asking.\\
\bottomrule
\end{xltabular}

Clicking a window’s mode line with the mouse selects that window.

\subsubsection{Settings}
\label{sec:22}
\begin{xltabular}{\textwidth}{@{}l l >{\raggedright\arraybackslash\hsize=1.000\hsize}X@{}}
\toprule
\textbf{Setting} & \textbf{Standard} & \textbf{What it changes}\\
\midrule
\endhead
\texttt{*mode-\codebreak{}line-\codebreak{}tuatara*} & \texttt{t} & Writes TUATARA before the modal state.\\
\texttt{*modal-\codebreak{}editing*} & \texttt{t} & Modal editing on. With \texttt{nil} the mode line has no state and every key types.\\
\texttt{*mode-\codebreak{}line-\codebreak{}name-\codebreak{}limit*} & \texttt{nil} & The most columns the buffer’s name may take, or \texttt{nil} for as many as fit.\\
\texttt{*processing-\codebreak{}indicator*} & \texttt{"P"} & The sign for jobs. \texttt{nil} hides it.\\
\texttt{*indexing-\codebreak{}indicator*} & \texttt{"U"} & The sign for folders being read.\\
\texttt{*searching-\codebreak{}indicator*} & \texttt{"S"} & The sign for searches over those folders.\\
\texttt{*daemon-\codebreak{}indicator*} & \texttt{"D"} & The sign for a daemon.\\
\texttt{*default-\codebreak{}directory*} & \texttt{"\textasciitilde{}/"} & The folder of a buffer with no file. \texttt{nil} is the folder Tuatara started in.\\
\texttt{*terminal-\codebreak{}title-\codebreak{}limit*} & \texttt{40} & The most columns of the title a program in a terminal buffer gave itself, shown in the mode line.\\
\bottomrule
\end{xltabular}

\subsubsection{Faces}
\label{sec:23}
A face is a name for a way of drawing text: a colour, bold, and so on. The
theme in force, ao-dark as standard, gives each face its look. Change one in
your init file, for example \texttt{(set-\codebreak{}face :mode-\codebreak{}line-\codebreak{}sign :fg "\#e0a33a")}.

\begin{xltabular}{\textwidth}{@{}l l >{\raggedright\arraybackslash\hsize=1.000\hsize}X@{}}
\toprule
\textbf{Face} & \textbf{Standard} & \textbf{What it paints}\\
\midrule
\endhead
\texttt{:mode-\codebreak{}line} & theme & The mode line of the selected window.\\
\texttt{:mode-\codebreak{}line-\codebreak{}inactive} & theme & The mode lines of the other windows.\\
\texttt{:mode-\codebreak{}line-\codebreak{}directory} & theme & The folder at the right end.\\
\texttt{:mode-\codebreak{}line-\codebreak{}sign} & theme & A sign while its work goes on.\\
\texttt{:mode-\codebreak{}line-\codebreak{}sign-\codebreak{}idle} & theme & A sign while nothing happens. Worked out from the mode line’s colours unless you set it.\\
\texttt{:mode-\codebreak{}line-\codebreak{}narrowed} & theme & The word Narrow.\\
\texttt{:mode-\codebreak{}line-\codebreak{}evidence} & theme & The word Evidence.\\
\texttt{:mode-\codebreak{}line-\codebreak{}look-\codebreak{}only} & theme & The words Look only.\\
\texttt{:mode-\codebreak{}line-\codebreak{}rewind} & theme & Rewind and the commit’s short name.\\
\texttt{:mode-\codebreak{}line-\codebreak{}recording} & theme & Rec and Played.\\
\texttt{:org-\codebreak{}mode-\codebreak{}line-\codebreak{}clock} & theme & The running Org clock.\\
\bottomrule
\end{xltabular}

\subsection{The echo area and the messages log}
\label{sec:27}
The echo area is the last row of the frame. Tuatara writes a short message
there after most things you do, such as \texttt{Wrote /home/\codebreak{}you/\codebreak{}letter.\codebreak{}txt}, and the
message stays until the next key you press does something. It also shows the
keys of a long sequence while you are typing it, and asks questions that one
key answers, such as yes or no.

Every message is also kept in the \texttt{*Messages*} buffer. A message said twice in
a row is written once with a count, as in \texttt{Saved [3 times]}. The buffer is
read-only, and it holds the last 1000 lines. Messages said while a password or
other secret was at hand are shown but not kept.

\texttt{C-h e} opens the log across the bottom of the frame, 30\% of the height,
with the cursor in it. It is a read-only buffer, so \texttt{q} closes it. The log
is the place to look when a message went by too fast.

\begin{try}
\textbf{Try it, with a message gone by}
\begin{enumerate}
\item Press \texttt{C-x C-f}, open \texttt{\textasciitilde{}/shopping.\codebreak{}txt}, press \texttt{i}, type \texttt{tea} and press
   \texttt{C-g}. Then press \texttt{C-x C-s}. The echo area says that the file was written
   (on a file with no change it says that no changes need to be saved). The
   message disappears when you press the next key.
\item Press \texttt{C-h e}. The log opens at the bottom with the message in its last
   line. Press \texttt{q} to close it.
\item Press \texttt{C-h N} to read the notifications Tuatara or its packages have given
   this session, for example when an agenda entry’s time came.
\end{enumerate}
\end{try}

\subsubsection{Commands}
\label{sec:25}
\begin{xltabular}{\textwidth}{@{}l l l >{\raggedright\arraybackslash\hsize=1.000\hsize}X@{}}
\toprule
\textbf{Modal} & \textbf{Non-modal} & \textbf{Command} & \textbf{What it does}\\
\midrule
\endhead
none & \texttt{C-h e} & \texttt{view-\codebreak{}echo-\codebreak{}area-\codebreak{}messages} & Shows the \texttt{*Messages*} log.\\
none & \texttt{C-h N} & \texttt{view-\codebreak{}notifications} & Shows the \texttt{*Notifications*} log, each line dated.\\
none & \texttt{C-g} & \texttt{keyboard-\codebreak{}quit} & Stops what is going on, drops a selection and closes a question.\\
\bottomrule
\end{xltabular}

\subsubsection{Settings}
\label{sec:26}
\begin{xltabular}{\textwidth}{@{}l l >{\raggedright\arraybackslash\hsize=1.000\hsize}X@{}}
\toprule
\textbf{Setting} & \textbf{Standard} & \textbf{What it changes}\\
\midrule
\endhead
\texttt{*message-\codebreak{}log-\codebreak{}max*} & \texttt{1000} & Lines \texttt{*Messages*} keeps. The oldest go first. \texttt{nil} keeps none.\\
\texttt{*lsp-\codebreak{}hover-\codebreak{}echo*} & \texttt{t} & Shows in the echo area what the thing under the cursor is, as a language server knows it.\\
\texttt{*escape-\codebreak{}quits*} & \texttt{nil} & Set to \texttt{t} to have \texttt{ESC} do what \texttt{C-g} does everywhere.\\
\texttt{*help-\codebreak{}window-\codebreak{}select*} & \texttt{nil} & See the settings for windows.\\
\bottomrule
\end{xltabular}

\subsection{The minibuffer and its prompt}
\label{sec:31}
When a command needs something from you, a file name, a word to search for, a
yes or no, it asks in the minibuffer, a row at the bottom of the frame
(Figure 1.1). The words that ask are the prompt, such as \texttt{Find file:}, and
your answer follows it (Figure 1.3). The prompt is drawn in bold. A frame
that has a question open keeps its window text where it was, and its cursor
is drawn as an outline there while you type in the minibuffer.

\begin{figure}[H]
\caption{Figure 1.3. The parts of the minibuffer.}\label{org28}
\begin{Verbatim}[breaklines,breakautoindent=false,breakindent=1em,breakafter=/-.\,;:=,commandchars=\\\{\}]
 3/12 Find file: ~/notes/sh
 ---- ---------- ----------
 count prompt    answer
\end{Verbatim}
\end{figure}

Press \texttt{RET} to give the answer, or \texttt{C-g} or \texttt{ESC} to give up. Before the prompt
Tuatara may show a count, such as \texttt{3/12}, which means the third of twelve
choices. A search shows which match the cursor is at and how many there are.
\texttt{*/0} means nothing matches but what you typed can be taken, and \texttt{!/0} means it
cannot.

You edit the answer with keys much like the ones in a text buffer. These keys
work only while a question is open, so \texttt{M-x} does not list their commands.

\begin{xltabular}{\textwidth}{@{}l l l >{\raggedright\arraybackslash\hsize=1.000\hsize}X@{}}
\toprule
\textbf{Modal} & \textbf{Non-modal} & \textbf{Command} & \textbf{What it does}\\
\midrule
\endhead
none & \texttt{RET} & \texttt{minibuffer-\codebreak{}finish} & Gives the answer, taking the selected choice if there is one.\\
none & \texttt{C-j}, \texttt{M-RET} & \texttt{minibuffer-\codebreak{}submit-\codebreak{}input} & Gives the answer exactly as typed, whatever is selected.\\
none & \texttt{C-g}, \texttt{ESC} & \texttt{minibuffer-\codebreak{}cancel} & Gives up the question.\\
none & \texttt{DEL}, \texttt{C-h} & \texttt{minibuffer-\codebreak{}backward-\codebreak{}delete-\codebreak{}char} & Deletes the character before the cursor.\\
none & \texttt{C-b}, \texttt{<left>} & \texttt{minibuffer-\codebreak{}backward-\codebreak{}char} & Moves the cursor left.\\
none & \texttt{C-f}, \texttt{<right>} & \texttt{minibuffer-\codebreak{}forward-\codebreak{}char} & Moves the cursor right.\\
none & \texttt{C-a}, \texttt{<home>} & \texttt{minibuffer-\codebreak{}beginning} & Moves to the start of the answer.\\
none & \texttt{C-e}, \texttt{<end>} & \texttt{minibuffer-\codebreak{}end} & Moves to the end of the answer.\\
none & \texttt{C-k} & \texttt{minibuffer-\codebreak{}kill-\codebreak{}line} & Deletes the answer after the cursor.\\
none & \texttt{C-u} & \texttt{minibuffer-\codebreak{}clear} & Clears the answer. The prompt stays.\\
none & \texttt{C-<backspace>}, \texttt{M-DEL} & \texttt{minibuffer-\codebreak{}backward-\codebreak{}kill-\codebreak{}word} & Deletes the word before the cursor.\\
none & \texttt{C-y}, \texttt{C-V}, \texttt{S-<insert>} & \texttt{minibuffer-\codebreak{}paste} & Pastes the clipboard, or what you cut last.\\
none & \texttt{M-p} & \texttt{minibuffer-\codebreak{}previous-\codebreak{}history} & Brings back an earlier answer to this question.\\
none & \texttt{M-n} & \texttt{minibuffer-\codebreak{}next-\codebreak{}history} & Goes to a later answer.\\
\bottomrule
\end{xltabular}

A question the editor thinks is risky, such as deleting a file, needs the
whole word yes or no and then \texttt{RET}. A question whose answer is easy to undo
takes one key. A password is typed with stars shown, and is never kept.

A question that offers a text to start with selects it, so that typing
replaces it and \texttt{RET} takes it.

\subsubsection{Settings}
\label{sec:29}
\begin{xltabular}{\textwidth}{@{}l l >{\raggedright\arraybackslash\hsize=1.000\hsize}X@{}}
\toprule
\textbf{Setting} & \textbf{Standard} & \textbf{What it changes}\\
\midrule
\endhead
\texttt{*show-\codebreak{}completion-\codebreak{}count*} & \texttt{t} & Shows the place of the selected choice and how many there are before the prompt.\\
\texttt{*prompt-\codebreak{}above-\codebreak{}candidates*} & \texttt{t} & Shows the prompt above the list of choices. \texttt{nil} puts it on the last row, under them.\\
\texttt{*last-\codebreak{}used-\codebreak{}first*} & \texttt{t} & Questions for a file or folder start where you chose last and list that entry first.\\
\bottomrule
\end{xltabular}

\subsubsection{Faces}
\label{sec:30}
\begin{xltabular}{\textwidth}{@{}l l >{\raggedright\arraybackslash\hsize=1.000\hsize}X@{}}
\toprule
\textbf{Face} & \textbf{Standard} & \textbf{What it paints}\\
\midrule
\endhead
\texttt{:minibuffer-\codebreak{}prompt} & theme & The prompt.\\
\texttt{:cursor} & theme & The cursor in a window of Tuatara’s own.\\
\texttt{:cursor-\codebreak{}inactive} & theme & The cursor in a window that does not get the typing.\\
\bottomrule
\end{xltabular}

\subsection{The choices panel}
\label{sec:34}
When a question has a list of answers, Tuatara shows them in the panel above
the last row, and narrows them as you type. Completing means finishing your
answer from that list. Type any part of a name. Several words may be typed,
and a choice has to match all of them, in any order. The letters that matched
are coloured, one colour for each word.

\begin{figure}[H]
\caption{Figure 1.3. Choosing a file. The prompt is above the choices.}
\begin{Verbatim}[breaklines,breakautoindent=false,breakindent=1em,breakafter=/-.\,;:=,commandchars=\\\{\}]
Mix the batter and let it rest.
 TUATARA[N] PD [-] pancakes.txt  L2 C0  Text  ~/notes/
  2/3  Find file: ~/notes/sh
  shed-plan.txt
  shopping.txt
  sheep-notes.org
\end{Verbatim}
\end{figure}

The selected choice is drawn reversed. \texttt{C-n} or \texttt{<down>} select the next one and
\texttt{C-p} or \texttt{<up>} the one before. Above the first one is the line you typed, when
what you typed can be taken as it is. \texttt{TAB} puts the selected choice into the
answer. For a file name it replaces the part after the last slash, and for a
folder it goes into the folder, so that you can walk down a path with \texttt{TAB}.
\texttt{RET} gives the selected choice, and \texttt{C-j} gives what you typed.

At most 10 rows are used for the panel, the prompt row included. When there
are more choices, \texttt{M-t} makes the panel as tall as the frame allows, and
again as usual. A choice may
have a note beside it, such as the keys of a command, the mode of a buffer or
the permissions of a file. A title with a line beside it begins a group, such
as the folder that the files under it are in.

\begin{try}
\textbf{Try it, with a folder of notes}
\begin{enumerate}
\item Press \texttt{C-x C-f}. Tuatara starts where you chose last and lists what is there.
\item Type \texttt{\textasciitilde{}/} and look at the panel. Each choice is a file or a folder.
\item Type a few letters of a file’s name. The list shrinks.
\item Press \texttt{C-n} twice and \texttt{C-p} once. The marked line moves.
\item Press \texttt{TAB}. The name is put into the answer. If it was a folder, its
   contents are listed.
\item Press \texttt{M-t}. The panel grows. Press it again to shrink it.
\item Press \texttt{C-g} to give up.
\end{enumerate}
\end{try}

\subsubsection{Commands}
\label{sec:32}
\begin{xltabular}{\textwidth}{@{}l l l >{\raggedright\arraybackslash\hsize=1.000\hsize}X@{}}
\toprule
\textbf{Modal} & \textbf{Non-modal} & \textbf{Command} & \textbf{What it does}\\
\midrule
\endhead
none & \texttt{C-n}, \texttt{<down>} & \texttt{minibuffer-\codebreak{}next-\codebreak{}candidate} & Selects the next choice.\\
none & \texttt{C-p}, \texttt{<up>} & \texttt{minibuffer-\codebreak{}previous-\codebreak{}candidate} & Selects the previous choice.\\
none & \texttt{TAB}, \texttt{C-M-i}, \texttt{M-TAB} & \texttt{minibuffer-\codebreak{}complete} & Puts the selected choice into the answer. In a Lisp question it lists the names that fit the word at the cursor.\\
none & \texttt{M-t} & \texttt{minibuffer-\codebreak{}toggle-\codebreak{}tall} & Makes the panel as tall as it can be, or as usual.\\
\bottomrule
\end{xltabular}

\subsubsection{Settings}
\label{sec:33}
The settings for the prompt, \texttt{*prompt-\codebreak{}above-\codebreak{}candidates*} and
\texttt{*show-\codebreak{}completion-\codebreak{}count*}, are in the table in the section before this one.

\begin{xltabular}{\textwidth}{@{}l l >{\raggedright\arraybackslash\hsize=1.000\hsize}X@{}}
\toprule
\textbf{Face} & \textbf{Standard} & \textbf{What it paints}\\
\midrule
\endhead
\texttt{:completion-\codebreak{}selected} & theme & The selected choice.\\
\texttt{:completion-\codebreak{}annotation} & theme & Notes beside choices.\\
\texttt{:completion-\codebreak{}group-\codebreak{}title} & theme & The title of a group of choices.\\
\texttt{:completion-\codebreak{}group-\codebreak{}line} & theme & The line beside a title.\\
\texttt{:completion-\codebreak{}match-0} & theme & What the first word typed matches.\\
\texttt{:completion-\codebreak{}match-1} & theme & What the second word matches.\\
\texttt{:completion-\codebreak{}match-2} & theme & What the third word matches.\\
\texttt{:completion-\codebreak{}match-3} & theme & What the fourth and later words match.\\
\texttt{:completion-\codebreak{}key} & theme & The keys beside commands.\\
\texttt{:completion-\codebreak{}major-\codebreak{}mode} & theme & The major mode beside a buffer.\\
\texttt{:completion-\codebreak{}modified} & theme & The mark beside a buffer with unsaved changes.\\
\texttt{:completion-\codebreak{}safe} & theme & The mark beside a file that opens safely.\\
\bottomrule
\end{xltabular}

\subsection{Key hints}
\label{sec:37}
Press \texttt{C-x} and wait. The panel lists every key that can follow, with the
name of its command, so you do not have to remember them. Names that start
with a plus sign, such as \texttt{+rectangles and registers}, are further prefixes,
and pressing their key opens another list. A key whose command cannot run in
this buffer is dimmed, and pressing it says why. \texttt{C-h} right after a prefix
key lists every key under it in a buffer of its own. \texttt{C-g} closes the hints.

The hints take as many rows as the frame has room for, and the echo area
stays under them. If they do not all fit, the last row ends in \texttt{... C-h lists
all}.

\begin{try}
\textbf{Try it, with the keys under C-x}
\begin{enumerate}
\item Press \texttt{C-x} and wait a moment. The panel appears with keys such as \texttt{C-f},
   \texttt{C-s} and \texttt{5}.
\item Press \texttt{5}. The panel changes to the keys under \texttt{C-x 5}, the ones for frames.
\item Press \texttt{C-g}. The hints close and nothing was run.
\item Press \texttt{C-x} and then \texttt{C-h}. A buffer lists all the keys under \texttt{C-x}.
\end{enumerate}
\end{try}

\subsubsection{Commands}
\label{sec:35}
\begin{xltabular}{\textwidth}{@{}l l l >{\raggedright\arraybackslash\hsize=1.000\hsize}X@{}}
\toprule
\textbf{Modal} & \textbf{Non-modal} & \textbf{Command} & \textbf{What it does}\\
\midrule
\endhead
none & \texttt{C-h} after a prefix & none & Lists every key under the prefix in a buffer.\\
none & \texttt{M-x} & \texttt{execute-\codebreak{}extended-\codebreak{}command} & Asks for a command by name, with its keys beside each choice.\\
\bottomrule
\end{xltabular}

\subsubsection{Settings}
\label{sec:36}
\begin{xltabular}{\textwidth}{@{}l l >{\raggedright\arraybackslash\hsize=1.000\hsize}X@{}}
\toprule
\textbf{Setting} & \textbf{Standard} & \textbf{What it changes}\\
\midrule
\endhead
\texttt{*key-\codebreak{}hint-\codebreak{}rows*} & \texttt{nil} & The most rows for the hints, the echo row included. \texttt{nil} takes all the room.\\
\bottomrule
\end{xltabular}

\begin{longtable}{@{}l l l@{}}
\toprule
\textbf{Face} & \textbf{Standard} & \textbf{What it paints}\\
\midrule
\endhead
\texttt{:key-\codebreak{}hint-\codebreak{}key} & theme & The keys.\\
\texttt{:key-\codebreak{}hint-\codebreak{}command} & theme & The command names.\\
\texttt{:key-\codebreak{}hint-\codebreak{}prefix} & theme & The names of prefixes.\\
\texttt{:key-\codebreak{}hint-\codebreak{}unavailable} & theme & Keys whose command cannot run here.\\
\texttt{:help-\codebreak{}key-\codebreak{}binding} & theme & Keys named in help and messages.\\
\bottomrule
\end{longtable}

\subsection{The cursor, the gutter and header lines}
\label{sec:40}
The cursor is drawn as a block as standard. A bar or an underline can be
chosen, and each modal state may have its own shape and colour, so that you
see at once whether you are in NORMAL or INSERT state. In a window that does
not get the typing the cursor is drawn as an outline in a window of its own
and as an underline in a terminal.

The gutter is on the left edge of a window. It is empty as standard. Line
numbers appear in it when you ask: \texttt{M-x toggle-\codebreak{}line-\codebreak{}numbers} goes through none,
absolute and relative, where relative numbers show how many lines away each
line is. Where a long line wraps, the gutter shows a mark beside the extra
rows.

Tuatara has no tab bar. Buffers are reached by name (\texttt{C-x b}) or from the list
(\texttt{C-x C-b}). A window can have a header line at its top. A file of records,
such as a CSV, can keep its first row, the header, at the top of the window
while the rest scrolls under it.

\subsubsection{Commands}
\label{sec:38}
\begin{xltabular}{\textwidth}{@{}l >{\raggedright\arraybackslash\hsize=0.848\hsize}X l >{\raggedright\arraybackslash\hsize=1.152\hsize}X@{}}
\toprule
\textbf{Modal} & \textbf{Non-modal} & \textbf{Command} & \textbf{What it does}\\
\midrule
\endhead
none & \texttt{M-x toggle-\codebreak{}line-\codebreak{}numbers} & \texttt{toggle-\codebreak{}line-\codebreak{}numbers} & Goes through no numbers, absolute and relative in this buffer.\\
none & \texttt{M-x display-\codebreak{}line-\codebreak{}numbers-\codebreak{}mode} & \texttt{display-\codebreak{}line-\codebreak{}numbers-\codebreak{}mode} & Absolute line numbers on or off.\\
none & \texttt{M-x display-\codebreak{}relative-\codebreak{}line-\codebreak{}numbers-\codebreak{}mode} & \texttt{display-\codebreak{}relative-\codebreak{}line-\codebreak{}numbers-\codebreak{}mode} & Relative line numbers on or off.\\
none & \texttt{M-x global-\codebreak{}display-\codebreak{}line-\codebreak{}numbers-\codebreak{}mode} & \texttt{global-\codebreak{}display-\codebreak{}line-\codebreak{}numbers-\codebreak{}mode} & The same for every buffer.\\
none & \texttt{C-x / t} & \texttt{list-\codebreak{}toggle-\codebreak{}header} & Takes the first line of a file of records as its header, or as a record.\\
none & \texttt{C-x / L} & \texttt{list-\codebreak{}toggle-\codebreak{}header-\codebreak{}lock} & Keeps the header at the top of the window while the rest scrolls.\\
\bottomrule
\end{xltabular}

\subsubsection{Settings}
\label{sec:39}
\begin{xltabular}{\textwidth}{@{}l l >{\raggedright\arraybackslash\hsize=1.000\hsize}X@{}}
\toprule
\textbf{Setting} & \textbf{Standard} & \textbf{What it changes}\\
\midrule
\endhead
\texttt{*cursor-\codebreak{}shape*} & \texttt{:block} & The cursor’s shape: \texttt{:block}, \texttt{:bar} or \texttt{:underline}.\\
\texttt{*cursor-\codebreak{}colour*} & \texttt{nil} & The cursor’s colour. \texttt{nil} uses the text’s colour.\\
\texttt{*tuatara-\codebreak{}normal-\codebreak{}cursor-\codebreak{}shape*} & \texttt{nil} & The shape in NORMAL state. \texttt{nil} uses \texttt{*cursor-\codebreak{}shape*}.\\
\texttt{*tuatara-\codebreak{}normal-\codebreak{}cursor-\codebreak{}colour*} & \texttt{nil} & The colour in NORMAL state.\\
\texttt{*tuatara-\codebreak{}insert-\codebreak{}cursor-\codebreak{}shape*} & \texttt{nil} & The shape in INSERT state.\\
\texttt{*tuatara-\codebreak{}insert-\codebreak{}cursor-\codebreak{}colour*} & \texttt{nil} & The colour in INSERT state.\\
\texttt{*tuatara-\codebreak{}support-\codebreak{}cursor-\codebreak{}shape*} & \texttt{nil} & The shape in lists, help and the log.\\
\texttt{*tuatara-\codebreak{}support-\codebreak{}cursor-\codebreak{}colour*} & \texttt{nil} & The colour in those buffers.\\
\texttt{*line-\codebreak{}numbers-\codebreak{}default*} & \texttt{nil} & Line numbers for buffers that did not choose: \texttt{nil}, \texttt{:absolute} or \texttt{:relative}.\\
\texttt{*line-\codebreak{}numbers-\codebreak{}modes*} & \texttt{nil} & Line numbers by major mode, such as \texttt{((prog-\codebreak{}mode . :relative))}.\\
\texttt{*soft-\codebreak{}wrap-\codebreak{}indicator*} & \texttt{"{\tuataraFallbackTWO{}↪}"} & The gutter’s mark beside a wrapped row. \texttt{nil} for none.\\
\texttt{*list-\codebreak{}header-\codebreak{}lock*} & \texttt{nil} & Keeps a header at the top wherever there is one.\\
\texttt{*csv-\codebreak{}header*} & \texttt{t} & Takes a CSV’s first line as its header.\\
\bottomrule
\end{xltabular}

\begin{longtable}{@{}l l l@{}}
\toprule
\textbf{Face} & \textbf{Standard} & \textbf{What it paints}\\
\midrule
\endhead
\texttt{:line-\codebreak{}number} & theme & The line numbers.\\
\texttt{:wrap-\codebreak{}indicator} & theme & The mark beside a wrapped row.\\
\texttt{:header-\codebreak{}line} & theme & The header of a list or a table.\\
\texttt{:vertical-\codebreak{}border} & theme & The line between windows side by side.\\
\bottomrule
\end{longtable}

\section{Moving around}
\label{sec:67}
How to put the cursor where you want it.

Before you can change text you have to get to it. This chapter lists every way
Tuatara has of moving the cursor: by character, word, line, sentence and
paragraph, to the start or end of the text, to a line you name, and back to a
place you left. It also covers counts, which repeat a move, and the marks that
let you return. For each part it gives the keys and the settings that change
how the move behaves. Scrolling and recentring are in chapter 1.

Keys are written as chapter 1 explains: \texttt{C-x} is Control and x, \texttt{M-x} is Alt
and x, \texttt{RET} is Return. The tables say “none” for a command that has no key
and can only be run by name with \texttt{M-x}.

\subsection{Two sets of keys}
\label{sec:42}
Tuatara has two sets of keys for moving. They do the same jobs, and you may use
either one, or both.

The first set is the letter keys of NORMAL state: \texttt{h}, \texttt{j}, \texttt{k}, \texttt{l} and a
few more. They need no Control or Alt, so your hands stay on the letters. They
work only in NORMAL state, because in INSERT state a letter is text.

The second set uses Control and Alt, as Emacs does: \texttt{C-f}, \texttt{M-f}, \texttt{C-a} and
so on. These work in both states. They are the keys to reach for while you are
typing in INSERT state and do not want to leave it. The arrow keys, \texttt{<home>} and
\texttt{<end>} work in both states as well. With \texttt{*modal-\codebreak{}editing*} set to \texttt{nil} the
second set is all there is.

A move never changes the text. If a move cannot go further, such as \texttt{l}
(non-modal: \texttt{C-f}) at the very end of the text, the cursor stays where it is.

The mode line shows the line and the column of the cursor, \texttt{L3 C4}, and both
count the way chapter 1 describes: the first line is 1 and the first column
is 0.

\subsection{By character and by line}
\label{sec:44}
\texttt{h} and \texttt{l} (non-modal: \texttt{C-b}, \texttt{C-f}) move one character left and right. \texttt{j} and
\texttt{k} (non-modal: \texttt{C-n}, \texttt{C-p}) move one line down and up. At the end of a line,
\texttt{l} goes on to the start of the next one, and \texttt{h} at the start of a line goes
back to the end of the one above.

When you move down or up, Tuatara tries to keep the column you started from.
If the next line is shorter, the cursor goes to its end. When you go on
pressing \texttt{j} or \texttt{k} (non-modal: \texttt{C-n}, \texttt{C-p}), the cursor returns to the
original column as soon as a line is long enough, so a short line in the middle
does not pull you off your column. Any other key, even \texttt{h} (non-modal: \texttt{C-b}),
forgets the column.

If you have turned on soft wrap (\texttt{C-x T l}), a long line shows as several rows
on the screen. \texttt{j} and \texttt{k} (non-modal: \texttt{C-n}, \texttt{C-p}) then move one row, not one
whole line.

In the hex view, which shows the bytes of a file, \texttt{h} and \texttt{l} (non-modal: \texttt{C-b},
\texttt{C-f}) move one byte instead of one character.

\begin{try}
\textbf{Try it, with a recipe}
\begin{enumerate}
\item Press \texttt{C-x C-f}, type \texttt{\textasciitilde{}/scones.\codebreak{}txt} and press \texttt{RET}. Press \texttt{i} and type
   these five lines, ending each but the last with \texttt{RET}: \texttt{Scones}, \texttt{2 cups
   flour}, \texttt{4 tsp baking powder}, \texttt{1 cup milk}, \texttt{Mix well and bake for twelve
   minutes}. Press \texttt{C-g} to go back to NORMAL state.
\item Press \texttt{g g} (non-modal: \texttt{M-<}) to go to the top. The mode line says \texttt{L1 C0}.
\item Press \texttt{j j}. You are on line 3 (\texttt{L3 C0}). Press \texttt{l l l} and watch the column
   go to 3. Press \texttt{h} (non-modal: \texttt{C-b}) once to see it come back to 2.
\item Press \texttt{j j} to go to the last line, then \texttt{g l} (non-modal: \texttt{C-e}) for its
   end. The column is 36, the length of that line.
\item Press \texttt{k} (non-modal: \texttt{C-p}). The line above is shorter, so the cursor goes
   to its end, at column 10. Press \texttt{j} (non-modal: \texttt{C-n}). The cursor is back at
   column 36.
\item Press \texttt{C-u 3 k}. The echo area shows \texttt{C-u 3-} while Tuatara waits for the
   key, then the cursor moves up three lines at once.
\end{enumerate}
\end{try}

\subsubsection{Commands}
\label{sec:43}
\begin{xltabular}{\textwidth}{@{}l l l >{\raggedright\arraybackslash\hsize=1.000\hsize}X@{}}
\toprule
\textbf{Modal} & \textbf{Non-modal} & \textbf{Command} & \textbf{What it does}\\
\midrule
\endhead
\texttt{l} & \texttt{C-f}, \texttt{<right>} & \texttt{forward-\codebreak{}char} & Moves one character on. A count moves that many, or back when it is negative. In the hex view, one byte.\\
\texttt{h} & \texttt{C-b}, \texttt{<left>} & \texttt{backward-\codebreak{}char} & Moves one character back.\\
\texttt{j} & \texttt{C-n}, \texttt{<down>} & \texttt{next-\codebreak{}line} & Moves down one line, keeping the column when it can. Where long lines wrap, one row.\\
\texttt{k} & \texttt{C-p}, \texttt{<up>} & \texttt{previous-\codebreak{}line} & Moves up one line, in the same way.\\
none & \texttt{<down-\codebreak{}mouse-1>} & \texttt{mouse-\codebreak{}set-\codebreak{}point} & Puts the cursor where you click with the first mouse button, in the window you click.\\
\bottomrule
\end{xltabular}

In a rectangle (chapter 5), \texttt{l} and \texttt{h} (non-modal: \texttt{C-f}, \texttt{C-b}) move the
corner of the rectangle at the cursor along the line, and do not run on to the
next line.

\subsection{By word}
\label{sec:46}
\texttt{w} (non-modal: \texttt{M-f}) moves to the end of the next word. If the cursor is in
the gap between two words, it first passes the gap and then the next word. \texttt{b}
(non-modal: \texttt{M-b}) moves to the start of the word before the cursor. A count
does it that many times. A negative count reverses the direction.

\texttt{W} and \texttt{B} (non-modal: \texttt{C-M-f}, \texttt{C-M-b}) move over a whole bracketed group or a
quoted piece of text in one step. In a program, they move over one expression as
the language reads it. Elsewhere they move over a bracketed group or a symbol.
At a closing bracket they stop, because the expression around the cursor ends
there. Chapter 13 describes moving over the parts of a program.

Words are looked for no further than 64 MB away, so one stray key in a huge
file does not make Tuatara wait. If the limit is reached, Tuatara says so.

\subsubsection{Commands}
\label{sec:45}
\begin{xltabular}{\textwidth}{@{}l l l >{\raggedright\arraybackslash\hsize=1.000\hsize}X@{}}
\toprule
\textbf{Modal} & \textbf{Non-modal} & \textbf{Command} & \textbf{What it does}\\
\midrule
\endhead
\texttt{w} & \texttt{M-f} & \texttt{forward-\codebreak{}word} & Moves to the end of the next word. A negative count moves back.\\
\texttt{b} & \texttt{M-b} & \texttt{backward-\codebreak{}word} & Moves to the start of the word before. A negative count moves on.\\
\texttt{W} & \texttt{C-M-f} & \texttt{forward-\codebreak{}sexp} & Moves forward over one balanced expression, a bracketed group, a string or a word.\\
\texttt{B} & \texttt{C-M-b} & \texttt{backward-\codebreak{}sexp} & Moves back over one balanced expression.\\
\bottomrule
\end{xltabular}

\subsection{Within a line}
\label{sec:49}
Figure 2.1 shows where the keys for the line and the word land when the cursor
is on the f of \texttt{flour}.

\begin{figure}[H]
\caption{Figure 2.1. Where keys take the cursor from the f of flour.}\label{org47}
\begin{Verbatim}[breaklines,breakautoindent=false,breakindent=1em,breakafter=/-.\,;:=,commandchars=\\\{\}]
 Mix the flour and the milk.
 ^   ^   ^    ^             ^
 g h b   here w             g l
\end{Verbatim}
\end{figure}

\texttt{g h} goes to the start of the line, and \texttt{g l} goes to its end, just before the
line break. Where soft wrap shows a long line as several rows, they go to the
start and the end of the row the cursor is on, and \texttt{C-a} and \texttt{C-e} do the same.

\texttt{M-m} goes to the first character of the line that is not a space or a tab.
This is the place to start when the line is indented.

\subsubsection{Commands}
\label{sec:48}
\begin{xltabular}{\textwidth}{@{}l l l >{\raggedright\arraybackslash\hsize=1.000\hsize}X@{}}
\toprule
\textbf{Modal} & \textbf{Non-modal} & \textbf{Command} & \textbf{What it does}\\
\midrule
\endhead
\texttt{g h} & \texttt{C-a}, \texttt{<home>} & \texttt{beginning-\codebreak{}of-\codebreak{}line} & Moves to the start of the line, or of the row where long lines wrap. In the hex view, to the start of the row of bytes.\\
\texttt{g l} & \texttt{C-e}, \texttt{<end>} & \texttt{end-\codebreak{}of-\codebreak{}line} & Moves to the end of the line, before its line break.\\
none & \texttt{M-m} & \texttt{back-\codebreak{}to-\codebreak{}indentation} & Moves to the first character that is not a space or a tab.\\
\bottomrule
\end{xltabular}

\subsection{By sentence and paragraph}
\label{sec:52}
A sentence ends at a full stop, a question mark, an exclamation mark, an
ellipsis (…) or an interrobang ({\tuataraFallbackTHREE{}‽}) that is followed by a space or the end of a
line. Closing quote marks and brackets after the stop belong to the sentence.
A full stop after an abbreviation such as \texttt{Dr.} or \texttt{e.g.} does not end it.
Tuatara knows a long list of those. You can add your own, or take one away that
you would rather have end a sentence.

Some people type two spaces after the end of a sentence. If you do, set
\texttt{*sentence-\codebreak{}end-\codebreak{}double-\codebreak{}space*} to \texttt{t}, and a sentence then ends only where
two spaces, a tab or a line break follow the stop. A file can say which it uses
with a line \texttt{sentence-\codebreak{}end-\codebreak{}double-\codebreak{}space: t} in its Local Variables block, and
that wins for that file.

A paragraph ends at a blank line. \texttt{M-\}} moves to the start of the blank line
after the paragraph, or to the end of the text if there is none, and \texttt{M-\{}
moves to the start of the blank line before the paragraph, or to the start of
the text.

The sentence and paragraph moves have no letter keys in NORMAL state, but the
Alt keys work there. A count does the move that many times, and a negative
count goes the other way.

\begin{try}
\textbf{Try it, with a letter}
\begin{enumerate}
\item Press \texttt{C-x C-f}, type \texttt{\textasciitilde{}/letter-\codebreak{}to-\codebreak{}napier.\codebreak{}txt} and press \texttt{RET}. Press \texttt{i}
   and type \texttt{Dear Napier,} and \texttt{RET} \texttt{RET}. Then type \texttt{Thank you for the
   weekend. The cake was lovely. We got home at six.} and \texttt{RET} \texttt{RET}. Then type
   \texttt{Please send the photos when you can. Dr. Taupo says hello.} Press \texttt{C-g}.
\item Press \texttt{g g} (non-modal: \texttt{M-<}). Press \texttt{M-\}}. The cursor goes to the blank
   line after the first paragraph (\texttt{L2 C0}). Press \texttt{j} (non-modal: \texttt{C-n}) to go
   to the start of the second paragraph.
\item Press \texttt{M-e}. The cursor moves to just after \texttt{weekend.}. Press \texttt{M-e} again,
   and it goes after \texttt{lovely.}. Press \texttt{M-a} to go back to the start of that
   sentence.
\item Press \texttt{M-\}} to go to the blank line after this paragraph. Press \texttt{M-\{} and
   watch the cursor go back up.
\item Press \texttt{g e} (non-modal: \texttt{M->}), then \texttt{M-a}. The cursor goes to the start of
   \texttt{Dr. Taupo says hello.}. The full stop after \texttt{Dr} did not end a sentence.
\end{enumerate}
\end{try}

\subsubsection{Commands}
\label{sec:50}
\begin{xltabular}{\textwidth}{@{}l l l >{\raggedright\arraybackslash\hsize=1.000\hsize}X@{}}
\toprule
\textbf{Modal} & \textbf{Non-modal} & \textbf{Command} & \textbf{What it does}\\
\midrule
\endhead
none & \texttt{M-e} & \texttt{forward-\codebreak{}sentence} & Moves to the end of the sentence, or of the paragraph’s text.\\
none & \texttt{M-a} & \texttt{backward-\codebreak{}sentence} & Moves to the start of the sentence.\\
none & \texttt{M-\}} & \texttt{forward-\codebreak{}paragraph} & Moves to the start of the blank line after the paragraph.\\
none & \texttt{M-\{} & \texttt{backward-\codebreak{}paragraph} & Moves to the start of the blank line before the paragraph.\\
\bottomrule
\end{xltabular}

\subsubsection{Settings}
\label{sec:51}
\begin{xltabular}{\textwidth}{@{}l l >{\raggedright\arraybackslash\hsize=1.000\hsize}X@{}}
\toprule
\textbf{Setting} & \textbf{Standard} & \textbf{What it changes}\\
\midrule
\endhead
\texttt{*sentence-\codebreak{}end-\codebreak{}double-\codebreak{}space*} & \texttt{nil} & Set to \texttt{t} for two spaces after the end of a sentence. Filling keeps the two spaces too.\\
\texttt{*extra-\codebreak{}sentence-\codebreak{}abbreviations*} & \texttt{()} & Abbreviations of your own whose full stop ends no sentence, as in \texttt{(setf *extra-\codebreak{}sentence-\codebreak{}abbreviations* '("Ph.\codebreak{}D." "Inc."))}.\\
\texttt{*removed-\codebreak{}sentence-\codebreak{}abbreviations*} & \texttt{()} & Abbreviations from Tuatara’s list that should end a sentence after all, such as \texttt{"St."}.\\
\bottomrule
\end{xltabular}

\subsection{To the start or the end of the text}
\label{sec:54}
\texttt{g g} (non-modal: \texttt{M-<}) goes to the very start of the buffer and \texttt{g e}
(non-modal: \texttt{M->}) goes to its end. If the buffer is narrowed (\texttt{C-x x n}), they
go to the start and end of the part shown. In the hex view, they go to the first
and last byte.

These two moves leave the mark behind at the place you left, so that you can go
back (see below). They do not do it while a selection is showing, nor when
you give a count.

\subsubsection{Commands}
\label{sec:53}
\begin{longtable}{@{}l l l l@{}}
\toprule
\textbf{Modal} & \textbf{Non-modal} & \textbf{Command} & \textbf{What it does}\\
\midrule
\endhead
\texttt{g g} & \texttt{M-<} & \texttt{beginning-\codebreak{}of-\codebreak{}buffer} & Moves to the start of the buffer.\\
\texttt{g e} & \texttt{M->} & \texttt{end-\codebreak{}of-\codebreak{}buffer} & Moves to the end of the buffer.\\
\bottomrule
\end{longtable}

\subsection{To a line}
\label{sec:56}
\texttt{:} asks for a line number in the minibuffer, and \texttt{RET} takes you to the start
of that line. \texttt{M-g g} and \texttt{M-g M-g} do the same, and they work in both states.
You may group the digits as the mode line does, \texttt{100,000,000}. The mark is left
where you were, unless a selection is showing.

In a large file whose lines are still being counted, a line in the first 64 MB
is found at once. A line further in waits until the count has reached it. In
the hex view the same keys ask for a byte offset instead, and a number that
starts with \texttt{0x} is read as hexadecimal.

If the line you ask for is outside the part of a narrowed buffer, Tuatara
widens the buffer and says so.

Tuatara has no command that asks for a column. To get to a column, go to the
start of the line and press \texttt{l} (non-modal: \texttt{C-f}) with a count, as in
\texttt{C-u 12 l}, which works when the line is long enough.

\subsubsection{Commands}
\label{sec:55}
\begin{xltabular}{\textwidth}{@{}l l l >{\raggedright\arraybackslash\hsize=1.000\hsize}X@{}}
\toprule
\textbf{Modal} & \textbf{Non-modal} & \textbf{Command} & \textbf{What it does}\\
\midrule
\endhead
\texttt{:} & \texttt{M-g g}, \texttt{M-g M-g} & \texttt{goto-\codebreak{}line} & Asks for a line number and goes to the start of that line. In the hex view, asks for an offset.\\
\bottomrule
\end{xltabular}

\subsection{Scrolling and recentring}
\label{sec:57}
The keys that scroll the window and put the cursor’s line in the middle are
listed in chapter 1: \texttt{C-v} and \texttt{M-v} scroll by nearly a screenful, \texttt{C-l} and
\texttt{z z} recentre, and the mouse wheel scrolls the window it is over. The settings
\texttt{*scroll-\codebreak{}margin*} and \texttt{*scroll-\codebreak{}recenter*} are described there too. They decide
how the window follows the cursor when one of the moves in this chapter takes
the cursor out of sight.

\subsection{Counts}
\label{sec:59}
Most moves in this chapter take a count: you say how many times to do
the move. Press \texttt{C-u}, type the number, and then press the key of the move. In
NORMAL state, \texttt{C-u 5 j} moves five lines down. The echo area shows \texttt{C-u 5-}
while Tuatara waits for the key.

\texttt{C-u} with no number means 4, and \texttt{C-u C-u} means 16. A minus sign after \texttt{C-u}
makes the count negative, so \texttt{C-u - 2 w} moves back two words, as \texttt{b}
(non-modal: \texttt{M-b}) would. A number you type after \texttt{C-u C-u} starts a new count
of its own.

\texttt{.} (non-modal: \texttt{C-z .}) runs the last command again, with the count it had.
Press \texttt{.} after \texttt{C-u 3 j} and the cursor goes three more lines down.

Other chapters use the same count for other things, such as how many times to
insert a character. The count is always taken by the key that follows it.

\begin{try}
\textbf{Try it, with a long letter}
\begin{enumerate}
\item Open any file that has more than 30 lines. Press \texttt{g g} (non-modal: \texttt{M-<}) to
   go to the top.
\item Press \texttt{C-u 3 0 j}. The cursor is on line 31.
\item Press \texttt{C-u C-u k}. You have gone 16 lines up, to line 15.
\item Press \texttt{C-u - 2 w} to move back two words.
\item Press \texttt{j} (non-modal: \texttt{C-n}), then \texttt{.}. The last command, \texttt{j}, runs once
   more.
\end{enumerate}
\end{try}

\subsubsection{Commands}
\label{sec:58}
\begin{xltabular}{\textwidth}{@{}l l l >{\raggedright\arraybackslash\hsize=1.000\hsize}X@{}}
\toprule
\textbf{Modal} & \textbf{Non-modal} & \textbf{Command} & \textbf{What it does}\\
\midrule
\endhead
none & \texttt{C-u} & \texttt{universal-\codebreak{}argument} & Starts a count. Alone it is 4, twice it is 16, digits give a number and a minus sign negates.\\
\texttt{.} & \texttt{C-z .} & \texttt{repeat} & Runs the last command again with its count.\\
\bottomrule
\end{xltabular}

\subsection{Marks and going back}
\label{sec:63}
The mark is a place that Tuatara remembers, and that stays where it is while the
cursor moves on. The text between the mark and the cursor is the region, drawn
in the \texttt{:region} face. The jumps in this chapter leave the mark where the cursor
was, so a long jump can be taken back.

Each buffer keeps the marks it has had in a mark ring, which holds the last 16.
Tuatara also keeps a global mark ring, with the marks of the buffers you set
marks in.

Tuatara keeps a second list as well: the places the cursor was at before each
move, ten for each buffer. \texttt{S} (non-modal: \texttt{C-z S}) goes back to the place
before the cursor’s last move, and \texttt{S} again goes further back, round the list.
It is meant for taking back a key that jumped to the wrong place. The echo area
says \texttt{Position 2 of 7} so you know how far back you are. A count before \texttt{S}
skips that many places.

Registers and bookmarks keep a place for longer. A register (\texttt{C-x r p}, then a
character) keeps the place until Tuatara quits, and \texttt{C-x r j} with the same
character goes back to it, even after the text above it has changed. A
bookmark (\texttt{C-x m m}, then a name) is kept for good, and \texttt{C-x m b} goes to one by
its name.

\begin{try}
\textbf{Try it, with a long list}
\begin{enumerate}
\item Open a file of more than 30 lines. Press \texttt{g g} (non-modal: \texttt{M-<}), then \texttt{:}
   (non-modal: \texttt{M-g g}), type \texttt{30} and press \texttt{RET}. The echo area says
   \texttt{Mark set} and you are on line 30.
\item Press \texttt{g e} (non-modal: \texttt{M->}) to go to the end. Press \texttt{S}. You are on the
   place before the jump, line 30. Press \texttt{S} again, and you go further back.
\item Press \texttt{C-u C-SPC}. The cursor goes to the mark, line 30, where you pressed
   \texttt{g e}. Press \texttt{C-u C-SPC} once more and you are back at line 1.
\item Press \texttt{C-x r p}, then \texttt{a}. Press \texttt{g e}. Press \texttt{C-x r j}, then \texttt{a}. The cursor
   goes back to where you set the register.
\item Press \texttt{C-SPC} at line 1, then \texttt{j j j}, and the region is shown. Press
   \texttt{C-x C-x}. The cursor and the mark swap places. Press \texttt{C-g} to drop the
   region.
\end{enumerate}
\end{try}

\subsubsection{Commands}
\label{sec:60}
\begin{xltabular}{\textwidth}{@{}l >{\raggedright\arraybackslash\hsize=0.545\hsize}X l >{\raggedright\arraybackslash\hsize=1.455\hsize}X@{}}
\toprule
\textbf{Modal} & \textbf{Non-modal} & \textbf{Command} & \textbf{What it does}\\
\midrule
\endhead
\texttt{S} & \texttt{C-z S} & \texttt{jump-\codebreak{}back} & Goes to a place the cursor was before, and further back each time you press it.\\
none & \texttt{C-SPC} & \texttt{set-\codebreak{}mark-\codebreak{}command} & Sets the mark at the cursor and shows the region. Pressed again at once, it hides the region, and once more shows it. With \texttt{C-u}, goes back to the mark, as \texttt{pop-\codebreak{}to-\codebreak{}mark-\codebreak{}command} does.\\
none & \texttt{C-u C-SPC} & \texttt{pop-\codebreak{}to-\codebreak{}mark-\codebreak{}command} & Goes to the mark and takes the mark before it from the ring, so the next press goes further back.\\
none & \texttt{C-x C-SPC} & \texttt{pop-\codebreak{}global-\codebreak{}mark} & Goes to the newest mark of the global ring, in its buffer, and so back through the buffers where you set marks.\\
none & \texttt{C-x C-x} & \texttt{exchange-\codebreak{}point-\codebreak{}and-\codebreak{}mark} & Swaps the cursor and the mark.\\
none & \texttt{M-x push-\codebreak{}mark-\codebreak{}command} & \texttt{push-\codebreak{}mark-\codebreak{}command} & Sets the mark and shows the region. When the mark is at the cursor already, only shows the region.\\
none & \texttt{C-x r p} & \texttt{register-\codebreak{}point} & Asks for a character and keeps the cursor’s place under it.\\
none & \texttt{C-x r j} & \texttt{register-\codebreak{}jump} & Asks for a character and goes to the place kept under it.\\
none & \texttt{C-x m m} & \texttt{bookmark-\codebreak{}set} & Keeps the place under a name that you give, for good.\\
none & \texttt{C-x m b} & \texttt{bookmark-\codebreak{}jump} & Goes to a bookmark by its name.\\
\bottomrule
\end{xltabular}

\subsubsection{Settings}
\label{sec:61}
\begin{xltabular}{\textwidth}{@{}l l >{\raggedright\arraybackslash\hsize=1.000\hsize}X@{}}
\toprule
\textbf{Setting} & \textbf{Standard} & \textbf{What it changes}\\
\midrule
\endhead
\texttt{*mark-\codebreak{}ring-\codebreak{}max*} & \texttt{16} & How many marks each buffer keeps, and the global ring. The oldest goes first.\\
\texttt{*set-\codebreak{}mark-\codebreak{}command-\codebreak{}repeat-\codebreak{}pop*} & \texttt{nil} & Set to \texttt{t} for \texttt{C-SPC} pressed right after \texttt{C-u C-SPC} to go on back to the mark before, instead of setting a mark.\\
\bottomrule
\end{xltabular}

\subsubsection{Faces}
\label{sec:62}
\begin{longtable}{@{}l l l@{}}
\toprule
\textbf{Face} & \textbf{Standard} & \textbf{What it paints}\\
\midrule
\endhead
\texttt{:region} & theme & The text between the mark and the cursor.\\
\bottomrule
\end{longtable}

\subsection{Going to the next place in a list}
\label{sec:65}
Some commands make a list of places for you to visit: the errors of a program
that was built, the lines that a search through files found, and the problems
that a language server reports. Chapters 13 and 6 describe how those lists are
made. Two keys take you through the list one place at a time, from any buffer. A
place in a file opens that file at the line.

\subsubsection{Commands}
\label{sec:64}
\begin{xltabular}{\textwidth}{@{}l l l >{\raggedright\arraybackslash\hsize=1.000\hsize}X@{}}
\toprule
\textbf{Modal} & \textbf{Non-modal} & \textbf{Command} & \textbf{What it does}\\
\midrule
\endhead
none & \texttt{M-g n} & \texttt{next-\codebreak{}place} & Goes to the next place of the last list.\\
none & \texttt{M-g p} & \texttt{previous-\codebreak{}place} & Goes to the place before the one gone to last.\\
\bottomrule
\end{xltabular}

\subsection{In help, lists and other read-only buffers}
\label{sec:66}
The windows that show help, the \texttt{*Messages*} log and similar read-only text are
in what the mode line calls SUPPORT state. The reading keys of NORMAL state
work there too: \texttt{h}, \texttt{j}, \texttt{k}, \texttt{l}, \texttt{w}, \texttt{W}, \texttt{b}, \texttt{B}, \texttt{g g}, \texttt{g e}, \texttt{g h},
\texttt{g l}, \texttt{S}, \texttt{:} and \texttt{z z}. They move the cursor and change nothing.

A list that a command made, such as the list of buffers or of the files in a
folder, may use some of these letters for its own purposes. In a list \texttt{h} closes
the window, and \texttt{g} starts a move, so the list is made again with \texttt{r}. The
chapter for each list says what its keys do.

\section{Writing and changing text}
\label{sec:98}
How to type, correct, undo, repeat and tidy text.

This chapter is about the text itself. It shows how to type it, take it out,
change it in place, bring it back after a mistake, and do the same thing again
without typing it again. It ends with the commands that tidy a whole block of
lines: filling, sorting, spelling and recording keys. Chapter 2 shows how to
get to the place you want to change, and chapter 4 shows what happens to the
text you cut.

Keys are written as chapter 1 explains: \texttt{C-x} is Control and x, \texttt{M-x} is Alt
and x, \texttt{RET} is Return, \texttt{DEL} is Backspace. The tables say “none” for a command
that has no key and can only be run by name with \texttt{M-x}.

\subsection{The two states}
\label{sec:70}
In NORMAL state the letter keys are commands. \texttt{d} deletes, \texttt{u} (non-modal:
\texttt{C-/}) undoes, \texttt{i} starts typing. In INSERT state the letter keys write text.
Most of your time goes to NORMAL state, where you choose what to change, and you
go into INSERT state only to type. Figure 3.1 shows how the two connect.

\begin{figure}[H]
\caption{Figure 3.1. The two states and the keys that move between them.}\label{org68}
\begin{Verbatim}[breaklines,breakautoindent=false,breakindent=1em,breakafter=/-.\,;:=,commandchars=\\\{\}]
                  i  a  I  A  o  O  c
        +----------------------------------->+
        |                                    |
   +----+-----+                        +-----+------+
   |  NORMAL  |                        |   INSERT   |
   | TUATARA[N]                        | TUATARA[I] |
   |  letters |                        |  letters   |
   | are keys |                        |  are text  |
   +----+-----+                        +-----+------+
        ^                                    |
        +<-----------------------------------+
                   C-g  or  ESC
\end{Verbatim}
\end{figure}

The mode line says which state you are in. Each buffer keeps its own state,
so a buffer you left in INSERT state is still in it when you come back to it.

Everything you type between the key that starts INSERT state and the key that
ends it is one step for undo, line breaks included. If you type two lines and
press \texttt{u} (non-modal: \texttt{C-/}), both go.

\subsubsection{Commands for starting to type}
\label{sec:69}
\begin{xltabular}{\textwidth}{@{}l l l >{\raggedright\arraybackslash\hsize=1.000\hsize}X@{}}
\toprule
\textbf{Modal} & \textbf{Non-modal} & \textbf{Command} & \textbf{What it does}\\
\midrule
\endhead
\texttt{i} & none & \texttt{tuatara-\codebreak{}insert} & Starts typing at the cursor. A selection, if there is one, is dropped first.\\
\texttt{a} & none & \texttt{tuatara-\codebreak{}insert-\codebreak{}after} & Starts typing after the character under the cursor.\\
\texttt{I} & none & \texttt{tuatara-\codebreak{}insert-\codebreak{}line-\codebreak{}start} & Starts typing at the start of the line.\\
\texttt{A} & none & \texttt{tuatara-\codebreak{}insert-\codebreak{}line-\codebreak{}end} & Starts typing at the end of the line.\\
\texttt{o} & none & \texttt{tuatara-\codebreak{}open-\codebreak{}below} & Opens a new line below this one and starts typing on it. A number first opens that many lines, with the cursor on the first.\\
\texttt{O} & none & \texttt{tuatara-\codebreak{}open-\codebreak{}above} & Opens a new line above this one and starts typing on it. A number first opens that many lines.\\
\texttt{c} & none & \texttt{tuatara-\codebreak{}change} & Cuts the selection, or deletes the character under the cursor, and starts typing. See “Changing text” below.\\
\texttt{C-g}, \texttt{ESC} & none & \texttt{tuatara-\codebreak{}normal} & In INSERT state, goes back to NORMAL state. The selection is dropped, and a line you began and left blank loses the blanks that were put in to indent it.\\
none & \texttt{C-g}, \texttt{ESC} & \texttt{tuatara-\codebreak{}quit} & The key’s command in the global map. In INSERT state it does the same as \texttt{tuatara-\codebreak{}normal}, and in NORMAL state it drops the selection.\\
\bottomrule
\end{xltabular}

At the end of a line the cursor stands on the line break, so \texttt{a} types after
it, at the start of the next line. \texttt{A} types at the end of the line you are on.

\texttt{RET}, \texttt{DEL} and \texttt{<delete>} change no text in NORMAL state, so you cannot spoil
a file by leaning on them. \texttt{TAB} does nothing there either, except in the
modes named by \texttt{*tuatara-\codebreak{}normal-\codebreak{}tab-\codebreak{}modes*}, where it does what it does in
INSERT state. \texttt{tuatara-\codebreak{}normal-\codebreak{}return}, \texttt{tuatara-\codebreak{}normal-\codebreak{}inert} and
\texttt{tuatara-\codebreak{}normal-\codebreak{}tab} are the commands that stand behind those three keys. \texttt{RET}
follows a link in an Org buffer and does nothing in a plain text buffer.

Typing the two keys of \texttt{*insert-\codebreak{}escape-\codebreak{}keys*} quickly one after the other leaves
INSERT state too. The first goes into the text as soon as you press it, and is
taken out again when the second follows within a fifth of a second. Typed more
slowly, both letters stay. It works in ordinary editing buffers, but not in the
minibuffer, a terminal or at several cursors.

With \texttt{*modal-\codebreak{}editing*} set to \texttt{nil}, there are no states. Every buffer is
always typing, letters write text, and all the other keys and commands in this
chapter that start with Control or Meta still work.

\subsection{Typing}
\label{sec:73}
In INSERT state a key that stands for a character puts that character in at the
cursor. \texttt{self-\codebreak{}insert-\codebreak{}command} is the command behind that. A number before such
a key puts the character in that many times, so \texttt{C-u 20 =} types a row of
twenty equals signs.

Characters you type and delete one after another, in either state, join into one
step for undo. After typing “weekend” and deleting the “d” with \texttt{DEL}, one \texttt{u}
(non-modal: \texttt{C-/}) takes back all of it.

\begin{try}
\textbf{Try it, with a letter}
\begin{enumerate}
\item Press \texttt{C-x C-f}, type \texttt{\textasciitilde{}/letter.\codebreak{}txt} and press \texttt{RET}. The file is new, so
   the buffer is empty. Press \texttt{i}. The mode line shows \texttt{TUATARA[I]}.
\item Type \texttt{Dear Hana,}, press \texttt{RET} twice, then type \texttt{Thanks for the weekend in
   Taupo. The weather was kind.} Press \texttt{RET} and type \texttt{Love, Aroha}.
\item Press \texttt{DEL} three times. The last letters of the name go.
\item Press \texttt{C-g}. The mode line shows \texttt{TUATARA[N]} again. Press \texttt{DEL} once more.
   Nothing happens, because \texttt{DEL} changes no text in NORMAL state.
\item Press \texttt{u} (non-modal: \texttt{C-/}). All the typing goes in one step and the buffer
   is empty again. Press \texttt{U} (non-modal: \texttt{C-?}) and all of it comes back.
\item Press \texttt{A} to type at the end of the line. Type \texttt{oha.}, which finishes
   the name again. Press \texttt{C-g}.
\item Press \texttt{C-x C-s} to save.
\end{enumerate}
\end{try}

\subsubsection{Commands}
\label{sec:71}
\begin{xltabular}{\textwidth}{@{}l l l >{\raggedright\arraybackslash\hsize=1.000\hsize}X@{}}
\toprule
\textbf{Modal} & \textbf{Non-modal} & \textbf{Command} & \textbf{What it does}\\
\midrule
\endhead
none & any character key & \texttt{self-\codebreak{}insert-\codebreak{}command} & Puts the character in. With a number, that many times. With an input method on and no number, the character goes through its rules.\\
none & \texttt{RET}, \texttt{C-j} & \texttt{newline} & Starts a new line, indented as the buffer’s mode indents it. Spaces left at the end of the line before go. A number gives that many line breaks.\\
none & \texttt{C-o} & \texttt{open-\codebreak{}line} & Puts a line break in after the cursor, and the cursor stays before it. The new line is not indented.\\
none & \texttt{TAB} & \texttt{indent-\codebreak{}for-\codebreak{}tab-\codebreak{}command} & In plain text, Org and CSV, goes on to the next tab stop. In code, indents the line.\\
none & \texttt{C-z TAB} & \texttt{tab-\codebreak{}to-\codebreak{}tab-\codebreak{}stop} & Goes on to the next tab stop in any buffer. A number goes that many stops on.\\
none & \texttt{C-q} & \texttt{quoted-\codebreak{}insert} & Reads one more key and puts in the character it stands for. See “Special characters”.\\
none & \texttt{M-/} & \texttt{expand-\codebreak{}word} & Finishes the word before the cursor with the nearest word that starts the same way. Press it again for the next one.\\
none & \texttt{C-M-/} & \texttt{complete-\codebreak{}word} & Lists the words that start like the one before the cursor.\\
\bottomrule
\end{xltabular}

\subsubsection{Settings}
\label{sec:72}
\begin{xltabular}{\textwidth}{@{}l l >{\raggedright\arraybackslash\hsize=1.000\hsize}X@{}}
\toprule
\textbf{Setting} & \textbf{Standard} & \textbf{What it changes}\\
\midrule
\endhead
\texttt{*modal-\codebreak{}editing*} & \texttt{t} & \texttt{nil} turns NORMAL and INSERT state off.\\
\texttt{*insert-\codebreak{}escape-\codebreak{}keys*} & \texttt{nil} & Two characters, \texttt{"jk"} for example, that leave INSERT state when typed quickly together.\\
\texttt{*escape-\codebreak{}quits*} & \texttt{nil} & \texttt{t} makes \texttt{ESC} do everything \texttt{C-g} does: close a question, stop a wait, stop work in the background.\\
\texttt{*indent-\codebreak{}new-\codebreak{}lines*} & \texttt{t} & \texttt{nil} keeps new lines at column 0 instead of indenting them as the mode does.\\
\texttt{*word-\codebreak{}completion-\codebreak{}keep-\codebreak{}case*} & \texttt{t} & Words are found whatever their case and put in with the case you typed. \texttt{nil} finds only words that start exactly as typed.\\
\bottomrule
\end{xltabular}

\subsection{Deleting}
\label{sec:75}
In INSERT state, \texttt{DEL} deletes the character before the cursor and \texttt{C-d} or
\texttt{<delete>} the one after it. A letter with accents, or an emoji with a skin
tone, goes as one character. A number first deletes that many, and a negative
number deletes the other way. A Windows line end, a carriage return and a
newline, counts as one.

In NORMAL state \texttt{d} does the same job and also works on a selection. Without a
selection it deletes the character under the cursor. It keeps no copy, because
a character under the cursor is a slip being corrected and not a cut. With a
selection it cuts the text into the kill ring, which chapter 4 describes. A
number before \texttt{d} deletes that many characters.

There is no command that deletes a word in one key. In NORMAL state press \texttt{m w}
(non-modal: \texttt{C-z m w}) to select the word at the cursor, which chapter 5
explains, and then \texttt{d}. Press \texttt{m w} again before \texttt{d} to take one more word.
\texttt{M-z} cuts up to the next character you name.

\subsubsection{Commands}
\label{sec:74}
\begin{xltabular}{\textwidth}{@{}l l l >{\raggedright\arraybackslash\hsize=1.000\hsize}X@{}}
\toprule
\textbf{Modal} & \textbf{Non-modal} & \textbf{Command} & \textbf{What it does}\\
\midrule
\endhead
none & \texttt{DEL} & \texttt{delete-\codebreak{}backward-\codebreak{}char} & Deletes the character before the cursor. A number deletes that many.\\
none & \texttt{C-d}, \texttt{<delete>} & \texttt{delete-\codebreak{}char} & Deletes the character after the cursor. A negative number deletes backwards.\\
\texttt{d} & none & \texttt{delete-\codebreak{}selection-\codebreak{}or-\codebreak{}char} & In NORMAL state. Cuts the selection, or deletes the characters under the cursor. A line selection takes the whole lines.\\
\texttt{D} & \texttt{C-k} & \texttt{kill-\codebreak{}line} & Cuts the rest of the line. At the end of a line, cuts the line break, which joins the next line on. \texttt{D} is for NORMAL state.\\
none & \texttt{C-S-<backspace>} & \texttt{kill-\codebreak{}whole-\codebreak{}line} & Cuts the whole line and its line break. A number cuts that many lines.\\
none & \texttt{M-z} & \texttt{zap-\codebreak{}up-\codebreak{}to-\codebreak{}char} & Cuts up to, and not including, the next character you type.\\
none & \texttt{M-\textbackslash{}} & \texttt{delete-\codebreak{}horizontal-\codebreak{}space} & Deletes the spaces and tabs around the cursor. With a number first, only those before it.\\
none & \texttt{M-SPC} & \texttt{cycle-\codebreak{}spacing} & Leaves one space where there were several. Press it again for none. A third press brings the spaces back.\\
none & \texttt{M-x just-\codebreak{}one-\codebreak{}space} & \texttt{just-\codebreak{}one-\codebreak{}space} & Leaves one space, with the cursor after it. A number leaves that many.\\
none & \texttt{C-x C-o} & \texttt{delete-\codebreak{}blank-\codebreak{}lines} & On a run of blank lines, leaves one. On a lone blank line, deletes it. On a line with text, deletes the blank lines after it.\\
\bottomrule
\end{xltabular}

Where a fold hides text, a delete never removes what you cannot see. The fold
opens first, and the same key pressed again deletes.

\subsection{Changing text}
\label{sec:78}
\texttt{c} is delete and type in one key. With a selection, it cuts the selection and
starts INSERT state. Without one, it deletes the character under the cursor and
starts INSERT state. A line selection keeps its final line break, so what you
type takes the place of the line’s words and the line stays a line. The cut and
what you type after it are one step for undo, so \texttt{u} (non-modal: \texttt{C-/}) puts
back the old text whole.

To replace a word, select it with \texttt{m w} (non-modal: \texttt{C-z m w}), press \texttt{c}, type
the new word and press \texttt{C-g}. To replace the end of a line, press \texttt{D}
(non-modal: \texttt{C-k}) and then \texttt{A} and type. \texttt{p} with a selection pastes the kill
ring over it, and chapter 4 describes it.

\begin{try}
\textbf{Try it, with a recipe}
\begin{enumerate}
\item Press \texttt{C-x C-f}, type \texttt{\textasciitilde{}/pancakes.\codebreak{}txt}, \texttt{RET}. Press \texttt{i} and type the
   lines \texttt{Pancakes}, \texttt{eggs}, \texttt{milk}, \texttt{flour}, each ended with \texttt{RET} except the
   last. Press \texttt{C-g}.
\item The flour should come before the milk. The cursor is on the last line. Press
   \texttt{C-x C-t}. The flour line and the milk line trade places.
\item Press \texttt{g g} (non-modal: \texttt{M-<}) to go to the top, then \texttt{M-u}. The title
   becomes \texttt{PANCAKES}. Press \texttt{u} (non-modal: \texttt{C-/}) and it is \texttt{Pancakes} again.
\item Press \texttt{j} (non-modal: \texttt{C-n}) to go to the eggs line. Press \texttt{d}. The first
   letter goes. Press \texttt{.}. The next letter goes. Press \texttt{u} twice.
\item Press \texttt{g h} (non-modal: \texttt{C-a}) to go to the start of the line, then
   \texttt{C-u 3 d}. Three letters go. Press \texttt{u}.
\item On the eggs line press \texttt{g j}. The flour line joins the eggs line, with one
   space between. Press \texttt{u}.
\item Press \texttt{g g j A}, type \texttt{, large}, and press \texttt{C-g}. Press \texttt{g g} and \texttt{m w c},
   type \texttt{Thin pancakes} and press \texttt{C-g}. The title changes in one step.
\item Press \texttt{C-x C-s} to save.
\end{enumerate}
\end{try}

\subsubsection{Commands for case}
\label{sec:76}
\begin{xltabular}{\textwidth}{@{}l l l >{\raggedright\arraybackslash\hsize=1.000\hsize}X@{}}
\toprule
\textbf{Modal} & \textbf{Non-modal} & \textbf{Command} & \textbf{What it does}\\
\midrule
\endhead
none & \texttt{M-u} & \texttt{upcase-\codebreak{}word} & Puts the next word in capitals and moves over it. A number does that many words, and a negative number does the words before the cursor.\\
none & \texttt{M-l} & \texttt{downcase-\codebreak{}word} & Puts the next word in small letters.\\
none & \texttt{M-c} & \texttt{capitalize-\codebreak{}word} & Makes the first letter a capital and the rest small.\\
none & \texttt{C-x C-u} & \texttt{upcase-\codebreak{}region} & Puts the selected text in capitals.\\
none & \texttt{C-x C-l} & \texttt{downcase-\codebreak{}region} & Puts the selected text in small letters.\\
\bottomrule
\end{xltabular}

\subsubsection{Commands for swapping and joining}
\label{sec:77}
\begin{xltabular}{\textwidth}{@{}l l l >{\raggedright\arraybackslash\hsize=1.000\hsize}X@{}}
\toprule
\textbf{Modal} & \textbf{Non-modal} & \textbf{Command} & \textbf{What it does}\\
\midrule
\endhead
none & \texttt{C-t} & \texttt{transpose-\codebreak{}chars} & Swaps the two characters about the cursor and moves on past them. At the end of a line, swaps the last two. A number carries the character before the cursor that many characters on.\\
none & \texttt{M-t} & \texttt{transpose-\codebreak{}words} & Swaps the words about the cursor. A number carries a word that many words on.\\
none & \texttt{C-x C-t} & \texttt{transpose-\codebreak{}lines} & Swaps this line with the one before it. A number carries the line that many lines on, and a negative number carries it back.\\
none & \texttt{C-M-t} & \texttt{transpose-\codebreak{}sexps} & Swaps the balanced pieces about the cursor, such as two bracketed groups. Chapter 13 says more about it.\\
none & \texttt{M-\textasciicircum{}} & \texttt{join-\codebreak{}line} & Joins this line to the one before it. The line break goes, and the spaces around it become one space, or none before a closing bracket or after an opening one. With a number first, joins the next line to this one. With a selection, joins all the lines it touches.\\
\texttt{g j} & none & \texttt{join-\codebreak{}selected-\codebreak{}lines} & In NORMAL state. Pulls the next line up onto this one. A number pulls that many. With a selection, joins the lines it touches into one.\\
\bottomrule
\end{xltabular}

A number of 0 given to a transpose command swaps the thing at the cursor with
the thing at the mark.

\subsection{Undo, redo and repeat}
\label{sec:81}
Every change can be undone, back to the file as you opened it. Nothing is let
go unless you set \texttt{*undo-\codebreak{}limit*}. Undo also covers work Tuatara does in the
background. A replacement or a sort of a large file is one step, and it stands
where you gave it: while it runs you can keep typing, and \texttt{u} (non-modal: \texttt{C-/})
takes back your edits first and then the background job.

\subsubsection{Commands}
\label{sec:79}
\begin{xltabular}{\textwidth}{@{}l l l >{\raggedright\arraybackslash\hsize=1.000\hsize}X@{}}
\toprule
\textbf{Modal} & \textbf{Non-modal} & \textbf{Command} & \textbf{What it does}\\
\midrule
\endhead
\texttt{u} & \texttt{C-x u}, \texttt{C-/}, \texttt{C-\_} & \texttt{undo} & Undoes the last change. In a narrowed buffer a change outside the part shown is undone too, and the echo area says so.\\
\texttt{U} & \texttt{M-\_}, \texttt{C-?}, \texttt{C-M-\_} & \texttt{redo} & Redoes the last change you undid. The cursor goes after the text that comes back.\\
\texttt{.} & \texttt{C-z .} & \texttt{repeat} & Runs the last command again, with the number it had.\\
none & \texttt{C-u} & \texttt{universal-\codebreak{}argument} & Starts a number for the next command. See “Repeat counts”.\\
\bottomrule
\end{xltabular}

\texttt{u} and \texttt{U} (non-modal: \texttt{C-/}, \texttt{C-?}) work in NORMAL state. The other keys work
in both states.

\texttt{repeat} is the key for “do that again”. If the last command acted on a
selection, the selection is made again first, as large as before, from the
cursor. It is the same number of whole lines, or characters, or the same kind of
text piece, or the same size of rectangle. After \texttt{V J d} deletes two lines, \texttt{.}
deletes the next two. Pressing \texttt{.} again runs the command once more. It repeats
one command only, so for a longer run of keys, record them (see “Recording
keys”).

A number given to \texttt{.} replaces the one the command had. After \texttt{C-u 3 d}, a plain
\texttt{.} deletes three characters again, and \texttt{C-u 5 .} deletes five.

\subsubsection{Settings}
\label{sec:80}
\begin{xltabular}{\textwidth}{@{}l l >{\raggedright\arraybackslash\hsize=1.000\hsize}X@{}}
\toprule
\textbf{Setting} & \textbf{Standard} & \textbf{What it changes}\\
\midrule
\endhead
\texttt{*undo-\codebreak{}limit*} & \texttt{nil} & The most memory, in bytes, that one buffer’s undo history may hold. \texttt{nil} keeps every change. With a limit, the oldest changes go first, down to about three quarters of it, and never the newest.\\
\bottomrule
\end{xltabular}

\subsection{Repeat counts}
\label{sec:82}
A number before a command makes it act that many times, or on that many
things. Type \texttt{C-u}, then the digits, then the command: \texttt{C-u 3 d} deletes three
characters and \texttt{C-u 12 M-u} puts twelve words in capitals. \texttt{C-u} alone means 4,
\texttt{C-u C-u} means 16, and \texttt{C-u -} before a command gives it a negative number.
While you type the number, the echo area shows it.

Each command in this chapter says what its number does. A command whose entry
says nothing about a number does not use one.

\subsection{Special characters}
\label{sec:85}
\texttt{C-q} reads the next key and puts in the character it stands for, whatever that
key does otherwise. \texttt{C-q TAB} puts in a real tab character. \texttt{C-q C-l} puts in a
form feed, the character that ends a page, and \texttt{C-q RET} puts in a carriage
return. A number first puts the character in that many times.

For letters your keyboard lacks there are three ways. You can type a character
by its Unicode name with \texttt{C-x 8 RET}, by a digraph (two keys that stand for one
character) with \texttt{C-x 8 \&}, or by an input method that changes what keys type
while it is on.

\begin{try}
\textbf{Try it, with a card for a friend in Sweden}
\begin{enumerate}
\item In \texttt{\textasciitilde{}/letter.\codebreak{}txt}, press \texttt{g e} to go to the end, then \texttt{o} to start a new
   line. Type \texttt{Greetings from Akaroa} and press \texttt{C-g}.
\item Press \texttt{C-x 8 RET}, type \texttt{em dash} and press \texttt{RET}. The dash is put in at the
   cursor.
\item Press \texttt{C-x 8 \&}, then \texttt{e} and \texttt{'}. An é goes in.
\item Press \texttt{C-x RET C-\textbackslash{}}, type \texttt{swedish-\codebreak{}postfix} and press \texttt{RET}. Press \texttt{i} and
   type \texttt{smaaland}. The word reads \texttt{småland}. Press \texttt{C-g}.
\item Press \texttt{C-\textbackslash{}} to turn the method off and \texttt{C-\textbackslash{}} again to turn it on.
\item Press \texttt{C-x 8 -} to turn the method off.
\end{enumerate}
\end{try}

\subsubsection{Commands}
\label{sec:83}
\begin{xltabular}{\textwidth}{@{}l l l >{\raggedright\arraybackslash\hsize=1.000\hsize}X@{}}
\toprule
\textbf{Modal} & \textbf{Non-modal} & \textbf{Command} & \textbf{What it does}\\
\midrule
\endhead
none & \texttt{C-q} & \texttt{quoted-\codebreak{}insert} & Puts in the character the next key stands for.\\
none & \texttt{M-x insert-\codebreak{}tab} & \texttt{insert-\codebreak{}tab} & Puts in a tab character, whatever the tab settings say.\\
none & \texttt{C-x 8 RET} & \texttt{insert-\codebreak{}char} & Asks for a Unicode name, such as \texttt{em dash}, or a code in hex such as \texttt{e9}, \texttt{U+00E9} or \texttt{\#xe9}. The names that fit are listed as you type. A number first puts the character in that many times.\\
none & \texttt{C-x 8 \&} & \texttt{insert-\codebreak{}digraph} & Asks for the two keys of a digraph, \texttt{e'} for é, \texttt{->} for an arrow, \texttt{<<} for a guillemet. With a selection, the character wraps it in a pair, such as « ». It works from NORMAL state too.\\
none & \texttt{M-x list-\codebreak{}digraphs} & \texttt{list-\codebreak{}digraphs} & Lists every digraph, the common ones first.\\
none & \texttt{C-x 8 .} & \texttt{digraph-\codebreak{}input-\codebreak{}mode} & Turns digraphs on as you type: \texttt{\&} and the digraph’s keys give the character, so \texttt{\&e'} types é.\\
none & \texttt{C-x 8 -} & \texttt{deactivate-\codebreak{}input-\codebreak{}method} & Turns off whichever input method is on in this buffer.\\
none & \texttt{C-\textbackslash{}} & \texttt{toggle-\codebreak{}input-\codebreak{}method} & Turns this buffer’s input method off, or on again. The first time it uses \texttt{*default-\codebreak{}input-\codebreak{}method*}, or asks.\\
none & \texttt{C-x RET C-\textbackslash{}} & \texttt{set-\codebreak{}input-\codebreak{}method} & Asks for an input method, turns it on, and makes it the one \texttt{C-\textbackslash{}} switches.\\
none & \texttt{C-h I} & \texttt{describe-\codebreak{}input-\codebreak{}method} & Shows the rules of the input method that is on, or of one you name.\\
\bottomrule
\end{xltabular}

While an input method is on, the mode line carries its short name after the mode
name, \texttt{SV<} for Swedish.

\subsubsection{Settings}
\label{sec:84}
\begin{xltabular}{\textwidth}{@{}l l >{\raggedright\arraybackslash\hsize=1.000\hsize}X@{}}
\toprule
\textbf{Setting} & \textbf{Standard} & \textbf{What it changes}\\
\midrule
\endhead
\texttt{*default-\codebreak{}input-\codebreak{}method*} & \texttt{nil} & The method \texttt{C-\textbackslash{}} turns on in a buffer that has had none, such as \texttt{"swedish-\codebreak{}postfix"}. \texttt{nil} asks.\\
\texttt{*input-\codebreak{}method-\codebreak{}keys*} & none & Methods of your own, each a key, a label and a method’s name. \texttt{(setf *input-\codebreak{}method-\codebreak{}keys* '(("s" "swedish" "swedish-\codebreak{}postfix")))} makes \texttt{SPC i s} turn it on.\\
\bottomrule
\end{xltabular}

\subsection{Indentation and filling in plain text}
\label{sec:88}
\texttt{TAB} in plain text goes on to the next tab stop, which is every four columns
as standard. Spaces are put in unless you ask for tab characters. In a language
mode \texttt{TAB} indents the line as the language does, which chapter 13 describes.

Filling re-flows a paragraph so that no line is longer than the fill column,
which is 80 as standard. The words stay in order, and the text that starts each
line (an indent, a \texttt{>} for a quotation, a comment mark) is kept. Each item of a
list is filled as a paragraph of its own, and its lines go on under its words,
past the mark and any check box: \texttt{- [ ] a task} continues six columns in.

\subsubsection{Commands}
\label{sec:86}
\begin{xltabular}{\textwidth}{@{}l l l >{\raggedright\arraybackslash\hsize=1.000\hsize}X@{}}
\toprule
\textbf{Modal} & \textbf{Non-modal} & \textbf{Command} & \textbf{What it does}\\
\midrule
\endhead
\texttt{g Q} & \texttt{M-q} & \texttt{fill-\codebreak{}paragraph} & Fills the paragraph at the cursor. In code, fills the comment there. Asks for the column first, with the buffer’s offered. A number first answers the question.\\
\texttt{g q} & \texttt{C-z g q} & \texttt{fill-\codebreak{}region} & Fills the paragraphs of the lines the selection touches. The lines before and after it are left alone.\\
none & \texttt{C-x f} & \texttt{set-\codebreak{}fill-\codebreak{}column} & Sets the fill column of this buffer for as long as Tuatara runs. A number first sets it without asking, as in \texttt{C-u 72 C-x f}.\\
none & \texttt{C-z TAB} & \texttt{tab-\codebreak{}to-\codebreak{}tab-\codebreak{}stop} & Goes on to the next tab stop.\\
none & \texttt{C-M-\textbackslash{}} & \texttt{indent-\codebreak{}region} & Indents the selected lines as the mode indents them.\\
\texttt{>} & \texttt{C-z >} & \texttt{indent-\codebreak{}selection-\codebreak{}or-\codebreak{}line} & In NORMAL state. Indents the selected lines, or this line, as the mode does.\\
none & \texttt{M-m} & \texttt{back-\codebreak{}to-\codebreak{}indentation} & Moves to the first character of the line that is not a space or tab.\\
none & \texttt{M-x indent-\codebreak{}tabs-\codebreak{}mode} & \texttt{indent-\codebreak{}tabs-\codebreak{}mode} & Switches this buffer between spaces and tab characters. A positive number turns tabs on, and zero or less turns them off.\\
none & \texttt{M-=} & \texttt{count-\codebreak{}words} & Counts the lines, sentences, words and characters of the selection, or of the whole buffer.\\
\bottomrule
\end{xltabular}

Filling takes prose only. In an Org file, headlines, tables, property drawers,
comment lines and the inside of source blocks are left as they are. In Markdown,
tables and the lines between code fences stay. So \texttt{g q} (non-modal: \texttt{C-z g q})
over a whole note re-flows its paragraphs and leaves its tables and code alone.

A paragraph larger than \texttt{*fill-\codebreak{}paragraph-\codebreak{}limit*} is left and the echo area says
so, so filling never holds the editor up. A sentence ends at \texttt{.}, \texttt{?}, \texttt{!} or
\texttt{…} followed by a space or a line end, or at a blank line. Short forms such as
\texttt{Dr.} end no sentence.

In your init file two functions set the tab habits for one kind of file.
\texttt{(set-\codebreak{}tab-\codebreak{}width 'c-\codebreak{}mode 8)} makes tab stops 8 columns apart in C files, and
\texttt{(set-\codebreak{}indent-\codebreak{}tabs-\codebreak{}mode 'bash t)} puts tab characters in shell scripts. Their
first argument is a major mode, and \texttt{prog} stands for every programming
language.

\subsubsection{Settings}
\label{sec:87}
\begin{xltabular}{\textwidth}{@{}l l >{\raggedright\arraybackslash\hsize=1.000\hsize}X@{}}
\toprule
\textbf{Setting} & \textbf{Standard} & \textbf{What it changes}\\
\midrule
\endhead
\texttt{*tab-\codebreak{}width*} & \texttt{4} & The columns from one tab stop to the next, a whole number from 1 to 64.\\
\texttt{*indent-\codebreak{}tabs-\codebreak{}mode*} & \texttt{nil} & \texttt{t} puts tab characters in where \texttt{TAB} goes on to a tab stop. \texttt{nil} puts spaces in.\\
\texttt{*tab-\codebreak{}always-\codebreak{}indent*} & \texttt{:complete} & What \texttt{TAB} does in code. \texttt{:complete} indents the line and, when that changes nothing, completes the name. \texttt{t} only indents. \texttt{nil} indents when the cursor is in the leading blanks and otherwise goes to the next stop.\\
\texttt{*tuatara-\codebreak{}normal-\codebreak{}tab-\codebreak{}modes*} & \texttt{(org-\codebreak{}mode)} & The modes where \texttt{TAB} works in NORMAL state, and the modes made from them.\\
\texttt{*fill-\codebreak{}column*} & \texttt{80} & The column filled lines end by.\\
\texttt{*fill-\codebreak{}ask-\codebreak{}column*} & \texttt{t} & \texttt{nil} fills at once, without asking for the column.\\
\texttt{*fill-\codebreak{}paragraph-\codebreak{}limit*} & \texttt{1048576} & The most bytes a paragraph may span to be filled.\\
\texttt{*sentence-\codebreak{}end-\codebreak{}double-\codebreak{}space*} & \texttt{nil} & \texttt{t} for two spaces between sentences: a sentence ends at a full stop followed by two spaces, and filling keeps two.\\
\texttt{*extra-\codebreak{}sentence-\codebreak{}abbreviations*} & \texttt{()} & Abbreviations of your own whose full stop ends no sentence, such as \texttt{'("Ph.\codebreak{}D." "et al.")}.\\
\texttt{*removed-\codebreak{}sentence-\codebreak{}abbreviations*} & \texttt{()} & Abbreviations of Tuatara’s own list that should end a sentence after all, such as \texttt{'("St.")}.\\
\bottomrule
\end{xltabular}

\subsection{Spelling}
\label{sec:92}
\texttt{M-x spellcheck-\codebreak{}mode} underlines the words spelt wrong in the buffer you are in.
It looks only at the words you can see, so a large file costs no more than a
small one. In text and Org every word is checked except code, links, dates and
tags. In code only comments and strings are checked. It uses enchant, a spelling
library, and needs a dictionary for your language installed. It is off until you
turn it on, and it never runs in a file opened safely or shown as bytes.

\subsubsection{Commands}
\label{sec:89}
\begin{xltabular}{\textwidth}{@{}l >{\raggedright\arraybackslash\hsize=0.850\hsize}X l >{\raggedright\arraybackslash\hsize=1.150\hsize}X@{}}
\toprule
\textbf{Modal} & \textbf{Non-modal} & \textbf{Command} & \textbf{What it does}\\
\midrule
\endhead
none & \texttt{M-x spellcheck-\codebreak{}mode} & \texttt{spellcheck-\codebreak{}mode} & Marks the words spelt wrong in this buffer, or stops.\\
none & \texttt{M-x global-\codebreak{}spellcheck-\codebreak{}mode} & \texttt{global-\codebreak{}spellcheck-\codebreak{}mode} & The same for every buffer.\\
\texttt{z n} & \texttt{C-z z n} & \texttt{spellcheck-\codebreak{}next} & Goes to the next word spelt wrong. A number goes that many on.\\
\texttt{z p} & \texttt{C-z z p} & \texttt{spellcheck-\codebreak{}previous} & Goes to the word spelt wrong before the cursor.\\
\texttt{z =} & \texttt{C-z z =} & \texttt{spellcheck-\codebreak{}correct} & Offers suggestions for the word, lets you type a replacement, or keeps the word as right.\\
\texttt{z g} & \texttt{C-z z g} & \texttt{spellcheck-\codebreak{}add-\codebreak{}word} & Adds the word to your own word list for a language. Other programs that use enchant share that list.\\
\texttt{z f} & \texttt{C-z z f} & \texttt{spellcheck-\codebreak{}add-\codebreak{}word-\codebreak{}to-\codebreak{}file} & Keeps the word as right in this file only, in a \texttt{LocalWords:} line at the end of it, written as a comment.\\
none & \texttt{M-x spellcheck-\codebreak{}languages} & \texttt{spellcheck-\codebreak{}languages} & Sets the languages this buffer is checked in, such as \texttt{en\_NZ sv\_SE}.\\
\bottomrule
\end{xltabular}

\texttt{z n}, \texttt{z p}, \texttt{z =}, \texttt{z g} and \texttt{z f} work in NORMAL state. In the other states
they are \texttt{C-z z n}, \texttt{C-z z p}, \texttt{C-z z =}, \texttt{C-z z g} and \texttt{C-z z f}.

\subsubsection{Settings}
\label{sec:90}
\begin{xltabular}{\textwidth}{@{}l >{\raggedright\arraybackslash\hsize=0.681\hsize}X >{\raggedright\arraybackslash\hsize=1.319\hsize}X@{}}
\toprule
\textbf{Setting} & \textbf{Standard} & \textbf{What it changes}\\
\midrule
\endhead
\texttt{*global-\codebreak{}spellcheck-\codebreak{}mode*} & \texttt{nil} & \texttt{t} checks every buffer.\\
\texttt{*spellcheck-\codebreak{}modes*} & none & Per major mode, such as \texttt{'((org-\codebreak{}mode . t) (prog-\codebreak{}mode . :comments) (csv-\codebreak{}mode . nil))}. \texttt{:strings} checks strings only, and \texttt{:all} checks all of a code buffer.\\
\texttt{*spellcheck-\codebreak{}languages*} & the locale’s & A list such as \texttt{'("en\_\codebreak{}NZ" "sv\_\codebreak{}SE")}. A word is right when any of them takes it. \texttt{nil} takes the language of the system’s locale, else en\_US.\\
\texttt{*spellcheck-\codebreak{}excluded-\codebreak{}faces*} & code, links, dates, tags, keywords & In text, the faces whose words are not checked.\\
\texttt{*spellcheck-\codebreak{}comment-\codebreak{}faces*} & \texttt{(:comment)} & In code, the faces of comments that are checked.\\
\texttt{*spellcheck-\codebreak{}string-\codebreak{}faces*} & \texttt{(:string)} & In code, the faces of strings that are checked.\\
\bottomrule
\end{xltabular}

\subsubsection{Faces}
\label{sec:91}
\begin{xltabular}{\textwidth}{@{}l >{\raggedright\arraybackslash\hsize=0.444\hsize}X >{\raggedright\arraybackslash\hsize=1.556\hsize}X@{}}
\toprule
\textbf{Face} & \textbf{Standard} & \textbf{What it paints}\\
\midrule
\endhead
\texttt{:spellcheck} & wavy yellow underline & A word spelt wrong. \texttt{(set-\codebreak{}face :spellcheck :underline :dotted)} in your init file makes dots, and \texttt{:underline t} a plain line.\\
\bottomrule
\end{xltabular}

\subsection{Sorting and tidying lines}
\label{sec:95}
These commands work on the lines a selection touches. Without a selection they
do nothing and say so. A large selection is handled in the background, so
Tuatara never waits on it. With a number, a sort runs in descending order.

\begin{try}
\textbf{Try it, with a shopping list}
\begin{enumerate}
\item Press \texttt{C-x C-f}, type \texttt{\textasciitilde{}/shopping.\codebreak{}txt}, \texttt{RET}. Press \texttt{i} and type the lines
   \texttt{milk}, \texttt{eggs}, \texttt{bread}, \texttt{eggs}, \texttt{butter}, each ended with \texttt{RET} except the
   last. Press \texttt{C-g}.
\item Press \texttt{g g} (non-modal: \texttt{M-<}), then \texttt{V} and \texttt{J} four times to select all
   five lines.
\item Press \texttt{M-x}, type \texttt{sort-\codebreak{}lines} and press \texttt{RET}. The list reads bread,
   butter, eggs, eggs, milk.
\item Select the five lines again with \texttt{g g}, \texttt{V} and \texttt{J} four times, as the sort
   ended the selection. Press \texttt{M-x}, type \texttt{delete-\codebreak{}duplicate-\codebreak{}lines} and press
   \texttt{RET}. The second \texttt{eggs} goes, and the echo area says one line was deleted.
\item Press \texttt{u} (non-modal: \texttt{C-/}) twice to get the original list back.
\item Select the lines again, run \texttt{flush-\codebreak{}lines}, type \texttt{e} and press \texttt{RET}. Every
   line with an \texttt{e} in it goes, so only \texttt{milk} is left. Press \texttt{u}.
\end{enumerate}
\end{try}

\subsubsection{Commands}
\label{sec:93}
\begin{xltabular}{\textwidth}{@{}>{\raggedright\arraybackslash\hsize=0.447\hsize}X >{\raggedright\arraybackslash\hsize=1.553\hsize}X@{}}
\toprule
\textbf{Command} & \textbf{What it does}\\
\midrule
\endhead
\texttt{sort-\codebreak{}lines} & Sorts the lines by their text. Capital and small letters compare by their codes unless \texttt{*sort-\codebreak{}fold-\codebreak{}case*} is on.\\
\texttt{sort-\codebreak{}fields} & Sorts the lines by one field. Fields are apart by spaces. A number names the field, and a negative one counts from the end of the line. In a CSV file or a log it asks for a column and moves each record whole, the header staying first.\\
\texttt{sort-\codebreak{}numeric-\codebreak{}fields} & Sorts by the number at the start of a field. \texttt{0x} starts a hexadecimal number. A field with no number counts as 0.\\
\texttt{sort-\codebreak{}columns} & Sorts by the text between the columns where the selection starts and ends.\\
\texttt{sort-\codebreak{}paragraphs} & Sorts paragraphs, which are apart by blank lines.\\
\texttt{sort-\codebreak{}pages} & Sorts pages, which end with a form feed at the start of a line.\\
\texttt{sort-\codebreak{}regexp-\codebreak{}fields} & Sorts the pieces a pattern finds, by a key you name: \texttt{\textbackslash{}\&} for the whole piece or \texttt{\textbackslash{}1} to \texttt{\textbackslash{}9} for a group.\\
\texttt{sort-\codebreak{}lines-\codebreak{}reverse}, \texttt{sort-\codebreak{}fields-\codebreak{}reverse}, \texttt{sort-\codebreak{}numeric-\codebreak{}fields-\codebreak{}reverse}, \texttt{sort-\codebreak{}columns-\codebreak{}reverse}, \texttt{sort-\codebreak{}paragraphs-\codebreak{}reverse}, \texttt{sort-\codebreak{}pages-\codebreak{}reverse}, \texttt{sort-\codebreak{}regexp-\codebreak{}fields-\codebreak{}reverse} & The same sorts, largest first.\\
\texttt{reverse-\codebreak{}region} & Puts the lines in the opposite order.\\
\texttt{delete-\codebreak{}duplicate-\codebreak{}lines} & Deletes all but the first of each set of identical lines. \texttt{C-u} keeps the last. \texttt{C-u C-u} deletes only a line identical to the one before it. \texttt{C-u C-u C-u} keeps all blank lines. Says how many went.\\
\texttt{flush-\codebreak{}lines} & Asks for a pattern and deletes the lines that match it. \texttt{\textasciicircum{}\$} matches the empty lines.\\
\texttt{keep-\codebreak{}lines} & Asks for a pattern and deletes the lines that do not match it.\\
\bottomrule
\end{xltabular}

None of these has a key. Run them with \texttt{M-x}. \texttt{flush-\codebreak{}lines} and \texttt{keep-\codebreak{}lines} use
the same pattern rules as search, which chapter 6 describes. In a CSV file or a
log they first ask which column to try the pattern on, and a header line always
stays.

\subsubsection{Settings}
\label{sec:94}
\begin{xltabular}{\textwidth}{@{}l l >{\raggedright\arraybackslash\hsize=1.000\hsize}X@{}}
\toprule
\textbf{Setting} & \textbf{Standard} & \textbf{What it changes}\\
\midrule
\endhead
\texttt{*sort-\codebreak{}fold-\codebreak{}case*} & \texttt{nil} & \texttt{t} ignores the case of ASCII letters, so \texttt{b} sorts before \texttt{C}.\\
\bottomrule
\end{xltabular}

\subsection{Recording keys}
\label{sec:97}
When you want the same few keys on line after line, record them once and play
them back. \texttt{C-x (} starts recording, \texttt{C-x )} stops, and \texttt{C-x e} plays the keys
as if you had pressed them again. What you type in answer to a question is
recorded too, so a search for a word searches for the same word each time. The
mouse is not recorded, and a passphrase is never recorded. A play stops where a
passphrase was typed and leaves that question to you. While you record, the mode
line shows \texttt{Rec}.

A number before \texttt{C-x e} plays that many times, and 0 plays until a play fails. A
play fails when a command in it fails, a search finds nothing or a move cannot
go on. \texttt{e} pressed straight after a play plays once more, and again, until you
press another key. Each play is one step for undo. A long run shows its progress
in the mode line as \texttt{Played 120/300}, and \texttt{C-g} stops it between two plays.
\texttt{C-g} while recording gives the recording up, and the keys recorded before it
are kept. In INSERT state the first \texttt{C-g} only leaves INSERT state.

\begin{try}
\textbf{Try it, with a numbered list of birds}
\begin{enumerate}
\item Press \texttt{C-x C-f}, type \texttt{\textasciitilde{}/birds.\codebreak{}txt}, \texttt{RET}. Press \texttt{i}, type \texttt{kea}, \texttt{RET},
   \texttt{weka}, \texttt{RET}, \texttt{kiwi}, \texttt{RET}, \texttt{takahe}, \texttt{RET}, \texttt{kaka}. Press \texttt{C-g}, then
   \texttt{g g}.
\item Press \texttt{C-x (}. Press \texttt{M-x}, type \texttt{insert-\codebreak{}keys-\codebreak{}counter}, press \texttt{RET}. Then
   press \texttt{i}, type a full stop and a space, press \texttt{C-g}, then \texttt{C-n} and
   \texttt{C-a}, and press \texttt{C-x )}.
\item The first line now begins with \texttt{1. =. Press \textasciitilde{}C-u 4 C-x e\textasciitilde{}. The other four
   lines are numbered. The last play meets the end of the buffer when it
   moves down, so the echo area says =Stopped after 3 of 4 plays: End of
   buffer}, but the number is already in.
\item Press \texttt{u} (non-modal: \texttt{C-/}) four times. The four plays are taken back one at
   a time.
\end{enumerate}
\end{try}

\subsubsection{Commands}
\label{sec:96}
\begin{xltabular}{\textwidth}{@{}l >{\raggedright\arraybackslash\hsize=0.889\hsize}X l >{\raggedright\arraybackslash\hsize=1.111\hsize}X@{}}
\toprule
\textbf{Modal} & \textbf{Non-modal} & \textbf{Command} & \textbf{What it does}\\
\midrule
\endhead
none & \texttt{C-x (} & \texttt{start-\codebreak{}recording-\codebreak{}keys} & Starts recording keys.\\
none & \texttt{C-x )} & \texttt{stop-\codebreak{}recording-\codebreak{}keys} & Ends the recording. The keys are kept for \texttt{C-x e}.\\
none & \texttt{C-x e} & \texttt{play-\codebreak{}recorded-\codebreak{}keys} & Plays the recorded keys. A number plays them that many times.\\
none & \texttt{M-x insert-\codebreak{}keys-\codebreak{}counter} & \texttt{insert-\codebreak{}keys-\codebreak{}counter} & Types a number and adds 1 to it, so recorded keys number the lines as they go. A number first adds that number.\\
none & \texttt{M-x set-\codebreak{}keys-\codebreak{}counter} & \texttt{set-\codebreak{}keys-\codebreak{}counter} & Sets the number typed next.\\
none & \texttt{M-x set-\codebreak{}keys-\codebreak{}counter-\codebreak{}format} & \texttt{set-\codebreak{}keys-\codebreak{}counter-\codebreak{}format} & Sets how the counter is written: \texttt{\%d} for the number, \texttt{\%03d} for three digits with zeros, \texttt{\%3d} with spaces before, \texttt{\%-3d} with spaces after, and \texttt{Item \%d.} for \texttt{Item 7.}\\
none & \texttt{M-x name-\codebreak{}recorded-\codebreak{}keys} & \texttt{name-\codebreak{}recorded-\codebreak{}keys} & Gives the recorded keys a name, which then plays them as \texttt{M-x NAME}. The name lasts while Tuatara runs.\\
none & \texttt{M-x insert-\codebreak{}recorded-\codebreak{}keys-\codebreak{}definition} & \texttt{insert-\codebreak{}recorded-\codebreak{}keys-\codebreak{}definition} & Types the init file lines that make the keys a command again and bind it to a key. You put them in the init file yourself.\\
\bottomrule
\end{xltabular}

The lines \texttt{insert-\codebreak{}recorded-\codebreak{}keys-\codebreak{}definition} types look like this, with your own
name and keys.

\begin{Verbatim}[breaklines,breakautoindent=false,breakindent=1em,breakafter=/-.\,;:=,commandchars=\\\{\}]
(define-recorded-keys 'kill-after-date "C-s date RET C-k")
(bind-key "C-c d" 'kill-after-date)
\end{Verbatim}

\section{Cutting and pasting}
\label{sec:110}
How to move text from one place to another, and to other programs.

Cutting takes text out of the buffer and keeps it. Copying keeps text and
leaves it where it is. Pasting puts the kept text in at the cursor. Emacs
calls cutting “killing” and pasting “yanking”, and Tuatara’s command names
use those words, so you will meet them here. This chapter shows the keys that
cut, copy and paste in both states, where the cut text is kept, how it
reaches other programs, how to keep several pieces of text at once and how to
set lines aside so that one key takes them together.

Chapter 5 teaches selecting, rectangles and many cursors. Here a selection is
only a way to say which text a key acts on.

\subsection{The kill ring}
\label{sec:100}
Everything you cut or copy goes on the kill ring, a list of your latest cuts
with the newest first. A paste puts the newest one in. Figure 4.1 shows the
ring after three cuts.

\begin{figure}[H]
\caption{Figure 4.1. The kill ring after three cuts. A paste uses the first.}\label{org99}
\begin{Verbatim}[breaklines,breakautoindent=false,breakindent=1em,breakafter=/-.\,;:=,commandchars=\\\{\}]
 kill ring                          system clipboard
 +-----------------------+          +-----------------------+
 | 1  "milk"             | <------> | "milk"                |
 | 2  "200 g flour"      |          +-----------------------+
 | 3  "2 eggs"           |            other programs paste
 +-----------------------+            from here
   newest on top
\end{Verbatim}
\end{figure}

The ring keeps \texttt{*kill-\codebreak{}ring-\codebreak{}max*} cuts. When it is full, the oldest is
forgotten. Tuatara has no command that steps back through the older cuts, so
a paste always gives you the newest. To keep a piece of text while you go on
cutting and copying other things, put it in a register, described below.

Two cuts in a row, with no other command between them, join into one. Press
\texttt{C-k} (modal: \texttt{D}) twice and the second cut is added to the first, so that one
paste brings both back, in order.

A cut that is too large for the system clipboard (see \texttt{*clipboard-\codebreak{}size-\codebreak{}limit*}
below) stays where it is in the buffer instead of being copied, and the echo
area says so when you copy it. It stays good after you kill the buffer it came
from. Text from an encrypted file never goes on the clipboard, and it leaves
the ring when its buffer is closed.

A cut that would take part of the text a folded heading hides shows the fold
instead. Press the same key again to cut.

\subsection{Cutting, copying and pasting with the Emacs keys}
\label{sec:102}
These keys mean the same in NORMAL and INSERT state. They act on the region,
the text between the cursor and the mark. A selection made with \texttt{v} is a
region too. If you have not selected anything, \texttt{C-SPC} sets the mark at the
cursor, and the region grows as you move away from it.

\begin{try}
\textbf{Try it, with a shopping list}
\begin{enumerate}
\item Press \texttt{C-x C-f}, type \texttt{\textasciitilde{}/shopping.\codebreak{}txt} and press \texttt{RET}. Press \texttt{i}, then
   type \texttt{eggs}, \texttt{RET}, \texttt{bread}, \texttt{RET}, \texttt{milk}, \texttt{RET}, \texttt{butter}.
\item Press \texttt{M-<} (modal: \texttt{g g}) to go to the top. Press \texttt{C-k} (modal: \texttt{D}). The
   word \texttt{eggs} is cut and the line is empty. Press \texttt{C-k} again. The empty line
   goes, and \texttt{bread} is now on the first line.
\item Press \texttt{C-y}. The cuts come back as one, \texttt{eggs} and its line end, in front
   of \texttt{bread}. Press \texttt{C-/} (modal: \texttt{u}) to undo that.
\item Press \texttt{C-SPC}, then \texttt{C-n} (modal: \texttt{j}). The region covers the first line.
   Press \texttt{M-w}. The echo area says \texttt{Copied} and nothing is removed.
\item Press \texttt{M->} (modal: \texttt{g e}) to go to the end, then \texttt{C-y}. The copied line is
   pasted after \texttt{butter}, and as that line has no line end the two run together.
   Press \texttt{C-x C-s} to save.
\end{enumerate}
\end{try}

\subsubsection{Commands}
\label{sec:101}
\begin{xltabular}{\textwidth}{@{}l l l >{\raggedright\arraybackslash\hsize=1.000\hsize}X@{}}
\toprule
\textbf{Modal} & \textbf{Non-modal} & \textbf{Command} & \textbf{What it does}\\
\midrule
\endhead
none & \texttt{C-w} & \texttt{kill-\codebreak{}region} & Cuts the region. A rectangle is cut as a block that \texttt{C-y} pastes as a block.\\
none & \texttt{M-w} & \texttt{kill-\codebreak{}ring-\codebreak{}save} & Copies the region and leaves it in place. Lines that a search filter hides are left out of the copy.\\
none & \texttt{C-y} & \texttt{yank} & Pastes the newest cut at the cursor, or what another program copied to the system clipboard since. The mark is left at the start of the pasted text, so \texttt{C-x C-x} selects it.\\
\texttt{D} & \texttt{C-k} & \texttt{kill-\codebreak{}line} & Cuts from the cursor to the end of the line. At the end of a line it cuts the line end.\\
none & \texttt{C-S-<backspace>} & \texttt{kill-\codebreak{}whole-\codebreak{}line} & Cuts the whole line and its line end. A number cuts that many lines, a negative number that many lines back, and 0 cuts the line without its line end.\\
none & \texttt{M-z} & \texttt{zap-\codebreak{}up-\codebreak{}to-\codebreak{}char} & Asks for a character, then cuts up to the next one of that kind and leaves it. A number takes the Nth one, and a negative number goes back.\\
none & \texttt{M-x zap-\codebreak{}to-\codebreak{}char} & \texttt{zap-\codebreak{}to-\codebreak{}char} & The same, but the character is cut too.\\
none & \texttt{M-x kill-\codebreak{}sexp} & \texttt{kill-\codebreak{}sexp} & Cuts the bracketed group, or the word, after the cursor. A number cuts that many.\\
none & \texttt{C-M-\codebreak{}DEL} & \texttt{backward-\codebreak{}kill-\codebreak{}sexp} & Cuts the one before the cursor.\\
\bottomrule
\end{xltabular}

\texttt{C-d} and \texttt{DEL} delete a character without keeping it. They are in chapter 3.

\subsection{Cutting, copying and pasting in NORMAL state}
\label{sec:103}
In NORMAL state the letter keys do the same work on one key each. With a
selection they act on it. Without one they act on the character under the
cursor, or on as many characters as the number you gave first.

\begin{xltabular}{\textwidth}{@{}l l l >{\raggedright\arraybackslash\hsize=1.000\hsize}X@{}}
\toprule
\textbf{Modal} & \textbf{Non-modal} & \textbf{Command} & \textbf{What it does}\\
\midrule
\endhead
\texttt{d} & none & \texttt{delete-\codebreak{}selection-\codebreak{}or-\codebreak{}char} & Cuts the selection and keeps it. With no selection it deletes the character and keeps nothing, because a character under the cursor is a typo being fixed. A line selection is widened to whole lines.\\
\texttt{c} & none & \texttt{tuatara-\codebreak{}change} & Cuts the selection and starts INSERT state. With no selection it deletes the character, unkept, and starts INSERT state.\\
\texttt{y} & none & \texttt{copy-\codebreak{}selection-\codebreak{}or-\codebreak{}char} & Copies the selection, or the character. A line selection is copied with its line end, so that it pastes as a line of its own.\\
\texttt{p} & none & \texttt{yank-\codebreak{}over-\codebreak{}selection} & Pastes the newest cut, or what another program copied since. It replaces the selection if there is one. A number pastes that many copies. The replaced text is not kept.\\
\texttt{D} & \texttt{C-k} & \texttt{kill-\codebreak{}line} & Cuts to the end of the line, as \texttt{C-k} does.\\
\texttt{P} & \texttt{C-x r y} & \texttt{yank-\codebreak{}rectangle} & Pastes the last rectangle cut or copied from a rectangle. See chapter 5.\\
\bottomrule
\end{xltabular}

If the kill ring is empty, \texttt{p} says so and changes nothing.

\begin{try}
\textbf{Try it, with a letter}
\begin{enumerate}
\item 1. Press \texttt{C-x C-f}, type \texttt{\textasciitilde{}/letter.\codebreak{}txt} and press \texttt{RET}. Press \texttt{i} and type
   four lines: \texttt{Dear Hana,}, \texttt{See you at Easter.}, \texttt{Thanks for the weekend in
   Raglan.} and \texttt{Love, Aroha}, each followed by \texttt{RET} except the last. Press
   \texttt{C-g} to go back to NORMAL state. 2. The thanks should come before the other
   line. Press \texttt{g g j j} to go to the third line, then \texttt{V} (non-modal: \texttt{C-z V})
   to select it and \texttt{d} to cut it. The line is gone. 3. Press \texttt{k} (non-modal:
   \texttt{C-p}) to go up, \texttt{g h} (non-modal: \texttt{C-a}) to go to the start of the line,
   then \texttt{p}. The thanks line is pasted above \texttt{See you at Easter.} Press \texttt{u}
   (non-modal: \texttt{C-/}) to undo it. 4. Press \texttt{V}, then \texttt{y}. The echo area says
   \texttt{Copied}. Press \texttt{g g p}. The copy is pasted above \texttt{Dear Hana,}. Press \texttt{u} to
   undo it. 5. Press \texttt{g g D}. The text of the first line is cut and the line
   stays, empty. Press \texttt{u} to bring it back. The text that \texttt{D} (non-modal:
   \texttt{C-k}) cut is on the kill ring, and \texttt{p} would paste it.
\end{enumerate}
\end{try}

\subsection{The system clipboard and the primary selection}
\label{sec:105}
The system clipboard is the one every program on your desktop shares. What you
cut or copy in Tuatara goes there as well, so you can paste it in another
program. When another program has copied something since, \texttt{C-y} and \texttt{p} paste
that instead and put it first on the kill ring.

The primary selection is a second place that many Linux programs have. It
holds whatever text was selected last, with no copy command. In Tuatara, a
selection you make by dragging with the mouse becomes the primary selection
when you let go. Clicking the middle mouse button in another program pastes it.
In Tuatara, the middle button pastes the primary selection of any program at
the place you clicked, with the mark left at its start. If nothing is selected
anywhere, the echo area says \texttt{Nothing is selected anywhere}. A selection made
with the keys does not become the primary selection.

The same limits apply to both. Text larger than \texttt{*clipboard-\codebreak{}size-\codebreak{}limit*} is
not sent, and neither is any text from an encrypted file, and with
\texttt{*use-\codebreak{}system-\codebreak{}clipboard*} set to \texttt{nil} nothing is sent or taken.

\begin{xltabular}{\textwidth}{@{}l >{\raggedright\arraybackslash\hsize=0.636\hsize}X l >{\raggedright\arraybackslash\hsize=1.364\hsize}X@{}}
\toprule
\textbf{Modal} & \textbf{Non-modal} & \textbf{Command} & \textbf{What it does}\\
\midrule
\endhead
none & \texttt{<mouse-2>} & \texttt{mouse-\codebreak{}yank-\codebreak{}primary} & Pastes the primary selection where the middle button was pressed.\\
none & \texttt{M-x gpg-\codebreak{}copy-\codebreak{}to-\codebreak{}clipboard} & \texttt{gpg-\codebreak{}copy-\codebreak{}to-\codebreak{}clipboard} & Puts the text copied last on the system clipboard even when it came from an encrypted file. Other programs, a clipboard manager among them, can then read and keep it.\\
none & \texttt{C}, \texttt{C-c C} & \texttt{server-\codebreak{}toggle-\codebreak{}clipboard} & In the list of servers, shares the clipboard with the server on the line, or stops sharing it. See chapter 18.\\
\bottomrule
\end{xltabular}

\subsubsection{Settings}
\label{sec:104}
\begin{xltabular}{\textwidth}{@{}l l >{\raggedright\arraybackslash\hsize=1.000\hsize}X@{}}
\toprule
\textbf{Setting} & \textbf{Standard} & \textbf{What it changes}\\
\midrule
\endhead
\texttt{*use-\codebreak{}system-\codebreak{}clipboard*} & \texttt{t} & Whether cuts and copies go on the system clipboard, and whether a paste takes in what another program copied. \texttt{nil} keeps them in Tuatara.\\
\texttt{*clipboard-\codebreak{}size-\codebreak{}limit*} & \texttt{8388608} & The most characters a cut may have and still go on the clipboard and the primary selection. That is 8 MB. A larger cut stays in the buffer’s own text.\\
\texttt{*kill-\codebreak{}ring-\codebreak{}max*} & \texttt{120} & How many cuts the kill ring keeps. The oldest go first.\\
\bottomrule
\end{xltabular}

\subsection{Registers}
\label{sec:106}
A register keeps one thing under one character, for as long as Tuatara runs.
The character is a letter, a digit or a sign. A capital letter and a small
letter name different registers. A register can hold text, a block of columns,
a place in a file or a whole number. Keeping something in a register does not
touch the kill ring, so a register is where to put what you want to paste
later, however much you cut in between.

Every register command asks for its character with one key and no \texttt{RET}. \texttt{C-g}
gives up. Keeping again under the same character replaces what was there.

\begin{try}
\textbf{Try it, with a recipe}
\begin{enumerate}
\item Press \texttt{C-x C-f}, type \texttt{\textasciitilde{}/pancakes.\codebreak{}txt} and press \texttt{RET}. Press \texttt{i} and type
   \texttt{2 eggs, 200 g flour, milk}, \texttt{RET}, \texttt{Mix, rest, fry.}, \texttt{RET}.
\item Press \texttt{M-<} (modal: \texttt{g g}), \texttt{C-SPC}, \texttt{C-e} (modal: \texttt{g l}) to select the first
   line without its line end. Press \texttt{C-x r s}, then \texttt{e}. The echo area says
   \texttt{Register e keeps text}.
\item Press \texttt{C-n C-a C-k} to cut the text of the second line. Press \texttt{C-x r i},
   then \texttt{e}. The ingredients are written in at the cursor, although the kill
   ring holds the method.
\item Press \texttt{C-u 4 C-x r n}, then \texttt{n}. Register \texttt{n} keeps the number 4. Press
   \texttt{C-x r +}, then \texttt{n}. It now keeps 5. Press \texttt{C-x r i}, then \texttt{n} to write \texttt{5}
   into the text.
\item Press \texttt{C-x r l}. The list shows each register with what it keeps. Press \texttt{q}
   to close it.
\end{enumerate}
\end{try}

\begin{xltabular}{\textwidth}{@{}l l l >{\raggedright\arraybackslash\hsize=1.000\hsize}X@{}}
\toprule
\textbf{Modal} & \textbf{Non-modal} & \textbf{Command} & \textbf{What it does}\\
\midrule
\endhead
none & \texttt{C-x r s} & \texttt{register-\codebreak{}copy} & Keeps the selected text under a character, or the block of columns when a rectangle is selected.\\
none & \texttt{C-x r p} & \texttt{register-\codebreak{}point} & Keeps the place of the cursor. It follows the text through edits, and it is found again after the buffer is closed when the buffer had a file.\\
none & \texttt{C-x r n} & \texttt{register-\codebreak{}number} & Keeps a whole number: the number you gave first, or one it asks for.\\
none & \texttt{C-x r +} & \texttt{register-\codebreak{}add} & Adds to a number register: the number you gave first, or 1. A negative number takes away. For a text register it puts the selected text after what is kept.\\
none & \texttt{M-x register-\codebreak{}prepend} & \texttt{register-\codebreak{}prepend} & Puts the selected text before the text a register keeps.\\
none & \texttt{C-x r i} & \texttt{register-\codebreak{}insert} & Writes what a register keeps at the cursor. Text goes in as it was, with the mark before it. A block goes in as a block, each row as many times as the number you gave first. A number goes in as digits.\\
none & \texttt{C-x r j} & \texttt{register-\codebreak{}jump} & Goes to the place a register keeps.\\
none & \texttt{C-x r l} & \texttt{register-\codebreak{}list} & Lists the registers in a buffer, one per line. \texttt{RET} on a place goes there, and \texttt{q} or \texttt{h} closes the list.\\
\bottomrule
\end{xltabular}

A register that keeps something that is not the right kind says what it
keeps instead. A place worth keeping beyond this session is a bookmark rather
than a register.

\subsection{Banked lines}
\label{sec:108}
Banking sets whole lines aside, wherever they are in the text, and shows them
in a colour of their own. Then \texttt{y}, \texttt{d} or \texttt{p} takes all the banked lines at
once, together with the lines of any selection you have made. This is the
easy way to gather lines that are far apart. Banking works in NORMAL state.

You can bank a line, move away and bank another. Banked lines follow the text
through edits, and they stay banked until you take them out of the bank or
until \texttt{y}, \texttt{d} or \texttt{p} uses them. \texttt{C-g} leaves them as they are, and \texttt{c} does
not act on them.

\begin{try}
\textbf{Try it, with a shopping list}
\begin{enumerate}
\item 1. Open \texttt{\textasciitilde{}/shopping.\codebreak{}txt} from the first try, or make a new file with six
   lines: \texttt{eggs}, \texttt{bread}, \texttt{milk}, \texttt{apples}, \texttt{tea}, \texttt{butter}. Press \texttt{C-g} so
   that you are in NORMAL state. 2. Press \texttt{g g j} to go to \texttt{bread}, then \texttt{m l}
   (non-modal: \texttt{C-z m l}). The line changes colour and the echo area says
   \texttt{Banked this line (1 total)}, then what to do next. 3. Press \texttt{j j j} to go to
   \texttt{tea}, then \texttt{m l} again. Now two lines are banked. 4. Press \texttt{m u} on \texttt{tea}.
   It leaves the bank. Press \texttt{m l} to bank it again. 5. Press \texttt{d}. Both lines
   are cut together and leave the bank, and the echo area says \texttt{Deleted 2
   lines}. 6. Press \texttt{g g p}. \texttt{bread} and \texttt{tea} are pasted above \texttt{eggs}, as one
   cut. Press \texttt{C-x C-s} to save.
\end{enumerate}
\end{try}

\begin{xltabular}{\textwidth}{@{}l >{\raggedright\arraybackslash\hsize=0.743\hsize}X l >{\raggedright\arraybackslash\hsize=1.257\hsize}X@{}}
\toprule
\textbf{Modal} & \textbf{Non-modal} & \textbf{Command} & \textbf{What it does}\\
\midrule
\endhead
\texttt{m l} & \texttt{C-z m l} & \texttt{bank-\codebreak{}lines} & Banks the line, or the lines the selection touches, and drops the selection. On lines that are all banked already it takes them out again, and on lines partly banked it banks the rest. A number banks that many lines from the cursor. A line whose end a fold hides is banked with all that the fold hides. Lines a search filter hides are not banked.\\
\texttt{m u} & \texttt{C-z m u} & \texttt{bank-\codebreak{}release-\codebreak{}line} & Takes the line at the cursor out of the bank and keeps the others.\\
\texttt{m U} & \texttt{C-z m U} & \texttt{bank-\codebreak{}release-\codebreak{}section} & Takes out the whole run of banked lines the cursor is on, the lines beside each other at once. Banked lines elsewhere stay.\\
\texttt{m DEL}, \texttt{m <delete>} & \texttt{C-z m DEL}, \texttt{C-z m <delete>} & \texttt{bank-\codebreak{}clear} & Lets go of every banked line in the buffer. The text is not changed.\\
\bottomrule
\end{xltabular}

With lines banked, \texttt{y} copies them as one cut and \texttt{d} cuts them as one cut and
one change, so that one \texttt{u} (non-modal: \texttt{C-/}) puts them back. \texttt{p} puts the
newest cut in place of them, as one change. Each of the three takes the banked
lines out of the bank. A cut too large to paste over banked lines is refused,
and the echo area gives its size.

\subsubsection{Faces}
\label{sec:107}
\begin{longtable}{@{}l l l@{}}
\toprule
\textbf{Face} & \textbf{Standard} & \textbf{What it paints}\\
\midrule
\endhead
\texttt{:banked-\codebreak{}line} & theme & The banked lines, across the full width of the window.\\
\bottomrule
\end{longtable}

\subsection{Pasting into prompts}
\label{sec:109}
When Tuatara asks you a question in the minibuffer, you can paste the answer.

\begin{xltabular}{\textwidth}{@{}l l l >{\raggedright\arraybackslash\hsize=1.000\hsize}X@{}}
\toprule
\textbf{Modal} & \textbf{Non-modal} & \textbf{Command} & \textbf{What it does}\\
\midrule
\endhead
none & \texttt{C-y}, \texttt{C-V}, \texttt{S-<insert>} & \texttt{minibuffer-\codebreak{}paste} & Writes in the system clipboard’s text, without the line end after it. With nothing on the clipboard it writes the text cut last. A suggested answer that is highlighted is replaced.\\
\bottomrule
\end{xltabular}

\texttt{C-V} is Control, Shift and v together, the paste key of most terminals. Typing
in a prompt does not touch the kill ring: \texttt{C-k} (modal: \texttt{D}) deletes the rest of
the answer and \texttt{C-u} clears all of it, and neither keeps what it removed. A
password prompt takes only what is on the clipboard, puts it in the answer and
nowhere else, and never in the kill ring.

\begin{try}
\textbf{Try it, with a file name}
\begin{enumerate}
\item In another program, such as your file manager, copy a file name, for
   example \texttt{shopping.\codebreak{}txt}.
\item In Tuatara press \texttt{C-x C-f}. The prompt asks for a file.
\item Press \texttt{C-y}. The name is written after the prompt. Press \texttt{RET} to open it.
\end{enumerate}
\end{try}

In a terminal buffer, \texttt{terminal-\codebreak{}paste} pastes into the program that runs
there, and a middle click does the same with the primary selection
(\texttt{terminal-\codebreak{}paste-\codebreak{}primary}). Chapter 17 describes both, with the setting
\texttt{*terminal-\codebreak{}paste-\codebreak{}keys*}.

When you paste into the terminal that Tuatara itself runs in, with that
terminal’s own paste key, \texttt{insert-\codebreak{}pasted-\codebreak{}text} puts the text in at the cursor
as text. None of it is taken for keys or commands. Its line ends become the
buffer’s, and one \texttt{u} (non-modal: \texttt{C-/}) takes it all back. It has no key.

\section{Selecting, rectangles and many cursors}
\label{sec:130}
How to pick out text, a block of columns, or many places at once.

Most commands in an editor act on something: a word to delete, a line to
copy, a block to indent. A selection is the text the next command acts on. It
shows in the \texttt{:region} face, usually as reversed colours. This chapter teaches
how to make a selection by hand, by the shape of the text and by the shape of
the code. It then goes on to blocks of columns, called rectangles, and to
working in many places at once, with many cursors or with a pattern that holds
every match.

Chapter 4 showed what \texttt{d}, \texttt{y} and \texttt{p} do with a selection. Chapter 2 showed
how to move without selecting. Here the keys only make the selection, and a
few commands that work on it follow.

\subsection{Selections and the mark}
\label{sec:113}
Tuatara has two ways to select, and both end up as the same thing.

The Emacs way uses the mark. The mark is a second place in the buffer, kept
beside the cursor. The region is the text between the mark and the cursor.
\texttt{C-SPC} sets the mark where the cursor is and shows the region, and every
move after that stretches it. \texttt{C-g} lets go. \texttt{C-x C-x} swaps the cursor and
the mark, which is how you reach the other end of a region to change it.
Chapter 2 describes the mark and the mark ring, the list of places the mark
has been.

The NORMAL state way has keys for it. \texttt{v} starts a selection at the cursor and
any move stretches it. Pressing \texttt{v} again while it is live starts it afresh from
where the cursor is now. \texttt{V} (non-modal: \texttt{C-z V}) selects whole lines instead.
\texttt{\%} selects the whole buffer. You can use whichever you like in either state,
and a selection made by one is a region for the other.

Figure 5.1 shows a selection made with \texttt{v} and one made with \texttt{V}.

\begin{figure}[H]
\caption{Figure 5.1. The same cursor, with \texttt{v w w} and with \texttt{V J}. [ ] marks the selected text.}\label{org111}
\begin{Verbatim}[breaklines,breakautoindent=false,breakindent=1em,breakafter=/-.\,;:=,commandchars=\\\{\}]
 v w w, from the u of "tui":

   The t[ui sings at] dawn.
   A kereru eats berries.
   The fantail is small.

 V J, two whole lines, line ends included:

   [The tui sings at dawn.
   A kereru eats berries.]
   The fantail is small.
\end{Verbatim}
\end{figure}

With a selection live, most letter keys act on it: \texttt{d} cuts it, \texttt{c} cuts it
and starts INSERT state, \texttt{y} copies it. A line selection is widened to whole
lines before it is cut or copied, so a deleted line leaves no empty line
behind.

\begin{try}
\textbf{Try it, with a recipe}
\begin{enumerate}
\item 1. Press \texttt{C-x C-f}, type \texttt{\textasciitilde{}/pancakes.\codebreak{}txt} and press \texttt{RET}. Press \texttt{i} and type
   \texttt{Pancakes}, \texttt{RET}, \texttt{2 eggs, 200 g flour, 1 cup milk}, \texttt{RET}, \texttt{Mix and rest.},
   \texttt{RET}, \texttt{Fry in butter.} Press \texttt{C-g}. 2. Press \texttt{g g j} to go to the second
   line. Press \texttt{v}, then \texttt{w w}. The first words of the line are selected. Press
   \texttt{C-g} to let go. 3. Press \texttt{V}. The whole line is selected. Press \texttt{J}
   (non-modal: \texttt{C-z J}). The line below is selected too. Press \texttt{d}. Both lines
   are cut. Press \texttt{u} (non-modal: \texttt{C-/}) to undo. 4. Press \texttt{V} on the second
   line, then \texttt{K}: the line above joins the selection. Press \texttt{C-g}. (\texttt{K} alone,
   with nothing selected, only moves up.) 5. Press \texttt{\%}. The whole buffer is
   selected. Press \texttt{C-g}.
\end{enumerate}
\end{try}

\subsubsection{Keys that select}
\label{sec:112}
\begin{xltabular}{\textwidth}{@{}l l >{\raggedright\arraybackslash\hsize=0.509\hsize}X >{\raggedright\arraybackslash\hsize=1.491\hsize}X@{}}
\toprule
\textbf{Modal} & \textbf{Non-modal} & \textbf{Command} & \textbf{What it does}\\
\midrule
\endhead
\texttt{v} & none & \texttt{start-\codebreak{}selection} & Starts a selection at the cursor. Any move stretches it.\\
\texttt{V} & \texttt{C-z V} & \texttt{select-\codebreak{}lines} & Selects the cursor’s line. Pressed with a line selection live, it lets go.\\
\texttt{J}, \texttt{K} & \texttt{C-z J}, \texttt{C-z K} & \texttt{select-\codebreak{}next-\codebreak{}line} and \texttt{select-\codebreak{}previous-\codebreak{}line} & Move down or up a line, adding it to a line selection. Without one they only move.\\
\texttt{\%} & none & \texttt{select-\codebreak{}whole-\codebreak{}buffer} & Selects everything, the cursor at the start.\\
none & \texttt{C-x h} & \texttt{mark-\codebreak{}whole-\codebreak{}buffer} & The same with the Emacs keys, the cursor at the start and the mark at the end.\\
\bottomrule
\end{xltabular}

With a line selection, \texttt{j} and \texttt{k} (non-modal: \texttt{C-n}, \texttt{C-p}) move a
character-wise cursor, while \texttt{J} and \texttt{K} move by whole lines. Press \texttt{C-g} in
NORMAL state to let go of any selection.

\begin{xltabular}{\textwidth}{@{}l l l >{\raggedright\arraybackslash\hsize=1.000\hsize}X@{}}
\toprule
\textbf{Modal} & \textbf{Non-modal} & \textbf{Command} & \textbf{What it does}\\
\midrule
\endhead
\texttt{v} & none & \texttt{start-\codebreak{}selection} & Start a selection at point: any motion extends it.\\
\texttt{V} & \texttt{C-z V} & \texttt{select-\codebreak{}lines} & Start a line selection anchored on the current line, or cancel a live one.\\
\texttt{J} & \texttt{C-z J} & \texttt{select-\codebreak{}next-\codebreak{}line} & Move down, extending the line selection if one is active.\\
\texttt{K} & \texttt{C-z K} & \texttt{select-\codebreak{}previous-\codebreak{}line} & Move up, extending the line selection if one is active.\\
\texttt{\%} & none & \texttt{select-\codebreak{}whole-\codebreak{}buffer} & Select the whole buffer: point at the start, the mark at the end.\\
none & \texttt{C-SPC} & \texttt{set-\codebreak{}mark-\codebreak{}command} & Set the mark at point and activate the region (C-SPC).\\
none & \texttt{C-x C-x} & \texttt{exchange-\codebreak{}point-\codebreak{}and-\codebreak{}mark} & Swap point and mark.\\
none & \texttt{C-x h} & \texttt{mark-\codebreak{}whole-\codebreak{}buffer} & Select the whole buffer: mark at the end, point at the start; where point was goes on the mark ring, with the mark before.\\
none & \texttt{C-g} & \texttt{tuatara-\codebreak{}quit} & C-g: in INSERT state return to NORMAL; in NORMAL drop the selection.\\
\bottomrule
\end{xltabular}

\subsection{Text objects, the m keys}
\label{sec:116}
Typing \texttt{v w w w} works, but it is slow when you want exactly one word, one
sentence or the text inside a pair of brackets. Tuatara has keys that take
such a piece whole. They all begin with \texttt{m}, which you can remember as “mark”.
The piece is called a text object.

The small letter takes the piece as a word, a sentence or a paragraph.
Pressed again, it takes one more. The capital letter does the same but grows
the selection the other way, backward, for the objects that have a direction.
For brackets and quotes it takes the pair with its brackets instead of what
is inside.

\begin{xltabular}{\textwidth}{@{}l l >{\raggedright\arraybackslash\hsize=0.669\hsize}X >{\raggedright\arraybackslash\hsize=1.331\hsize}X@{}}
\toprule
\textbf{Modal} & \textbf{Non-modal} & \textbf{Command} & \textbf{What it selects}\\
\midrule
\endhead
\texttt{m w} & \texttt{C-z m w} & \texttt{select-\codebreak{}word} & The word at the cursor. In a gap between two words, the next one.\\
\texttt{m b} & \texttt{C-z m b} & \texttt{select-\codebreak{}word-\codebreak{}backward} & The same, growing backward.\\
\texttt{m W}, \texttt{m B} & \texttt{C-z m W}, \texttt{C-z m B} & \texttt{select-\codebreak{}symbol}, \texttt{select-\codebreak{}symbol-\codebreak{}backward} & A symbol: a name with its \texttt{-} and \texttt{\_}, so \texttt{shopping-\codebreak{}list} is one symbol.\\
\texttt{m s}, \texttt{m S} & \texttt{C-z m s}, \texttt{C-z m S} & \texttt{select-\codebreak{}sentence}, \texttt{select-\codebreak{}sentence-\codebreak{}backward} & The sentence, as far as its full stop.\\
\texttt{m p}, \texttt{m P} & \texttt{C-z m p}, \texttt{C-z m P} & \texttt{select-\codebreak{}paragraph}, \texttt{select-\codebreak{}paragraph-\codebreak{}backward} & The paragraph with one blank line next to it.\\
\texttt{m i} & \texttt{C-z m i} & \texttt{select-\codebreak{}inside-\codebreak{}pair} & What is inside a pair of brackets or quotes, without them.\\
\texttt{m a} & \texttt{C-z m a} & \texttt{select-\codebreak{}pair} & The pair of brackets or quotes and what is inside.\\
\texttt{m I}, \texttt{m A} & \texttt{C-z m I}, \texttt{C-z m A} & \texttt{select-\codebreak{}sexp-\codebreak{}inside}, \texttt{select-\codebreak{}sexp} & What a bracketed expression holds, and the expression with its brackets.\\
\bottomrule
\end{xltabular}

A paragraph selection takes one blank line with it, the one before it. When
you delete it, one blank line is left between its neighbours, which is what you
want most of the time.

The pairs that \texttt{m i} and \texttt{m a} know are in the setting \texttt{*pair-\codebreak{}delimiters*}.
Put the cursor on a bracket and press \texttt{m i}. The echo area tells you it
selected the content, and \texttt{m a} the content with the brackets. The pair of the
bracket under the cursor is taken at once. With the cursor anywhere else,
\texttt{m i} and \texttt{m a} ask for a bracket: press the one you mean, such as \texttt{(}.

\begin{try}
\textbf{Try it, with a letter}
\begin{enumerate}
\item Press \texttt{C-x C-f}, type \texttt{\textasciitilde{}/letter.\codebreak{}txt} and press \texttt{RET}. Press \texttt{i} and type
   \texttt{Dear Hana,}, \texttt{RET}, \texttt{The kereru ate (all 12) of the berries.}, \texttt{RET},
   \texttt{It was a good day. See you at Easter.} Press \texttt{C-g}.
\item Press \texttt{g g j w} to land on \texttt{kereru}. Press \texttt{m w} (non-modal: \texttt{C-z m w}). The
   word is selected. Press \texttt{m w} again and \texttt{kereru ate} is selected. Press
   \texttt{C-g}.
\item Move to the brackets: press \texttt{g g j}, then press \texttt{w} (non-modal: \texttt{M-f}) until
   the cursor is on the closing bracket after \texttt{12}. Press \texttt{m i}. The text
   inside, \texttt{all 12}, is selected without its brackets. Press \texttt{m a}, then \texttt{(},
   and the brackets come in. Press \texttt{c}, type \texttt{(many)} and press \texttt{C-g}. Press \texttt{u}
   to undo.
\item Press \texttt{j} (non-modal: \texttt{C-n}), then \texttt{m s}. The sentence \texttt{It was a good day.}
   is selected. Press \texttt{m s} again and the second sentence joins it. Press \texttt{C-g}.
\item Press \texttt{m p}. The paragraph is selected, with the blank line before it.
   Press \texttt{C-g}.
\end{enumerate}
\end{try}

\subsubsection{Wrapping a selection}
\label{sec:114}
With a selection live, a bracket or a quote wraps it. Press \texttt{m w} on a word
and then \texttt{"} to put quotes round the word. Press \texttt{"} once more while the word
is selected inside its quotes and they come off. This works with the other
brackets and quotes of \texttt{*pair-\codebreak{}delimiters*}, and on a rectangle it wraps each
row.

\begin{xltabular}{\textwidth}{@{}l l l >{\raggedright\arraybackslash\hsize=1.000\hsize}X@{}}
\toprule
\textbf{Modal} & \textbf{Non-modal} & \textbf{Command} & \textbf{What it does}\\
\midrule
\endhead
\texttt{m w} & \texttt{C-z m w} & \texttt{select-\codebreak{}word} & Select the word at point; from the gap between two words, the one ahead (the last one at the end of the text).\\
\texttt{m b} & \texttt{C-z m b} & \texttt{select-\codebreak{}word-\codebreak{}backward} & Select the word at point, or grow a word selection backward, as select-word grows it forward: from the gap between two words, the one behind.\\
\texttt{m W} & \texttt{C-z m W} & \texttt{select-\codebreak{}symbol} & Select the symbol at point (a name with its \_ and -, Lisp’s own characters in Lisp), punctuation at either end left out: from the gap between two, the one ahead, brackets and quotes crossed.\\
\texttt{m B} & \texttt{C-z m B} & \texttt{select-\codebreak{}symbol-\codebreak{}backward} & Select the symbol at point, or grow a symbol selection backward: from the gap between two, the one behind.\\
\texttt{m s} & \texttt{C-z m s} & \texttt{select-\codebreak{}sentence} & Select the sentence at point, as far as its full stop; from the gap between two, the one ahead.\\
\texttt{m S} & \texttt{C-z m S} & \texttt{select-\codebreak{}sentence-\codebreak{}backward} & Select the sentence at point, or grow a sentence selection backward: from the gap between two, the one behind.\\
\texttt{m p} & \texttt{C-z m p} & \texttt{select-\codebreak{}paragraph} & Select the paragraph at point with one blank line, the one before it (the one after, for the first paragraph), so that deleting it leaves one blank line between its neighbours: from a blank line between two, the one below.\\
\texttt{m P} & \texttt{C-z m P} & \texttt{select-\codebreak{}paragraph-\codebreak{}backward} & Select the paragraph at point, or grow a paragraph selection backward: from a blank line between two, the one above.\\
\texttt{m i} & \texttt{C-z m i} & \texttt{select-\codebreak{}inside-\codebreak{}pair} & Select what is inside a pair of delimiters, the delimiters left out: the pair of the delimiter at point, or the one whose half is typed next ((, ), [, “, ‘, =, * and the rest of \textbf{pair-delimiters}).\\
\texttt{m a} & \texttt{C-z m a} & \texttt{select-\codebreak{}pair} & Select a pair of delimiters and what it holds, as select-inside-pair finds it.\\
\texttt{m I} & \texttt{C-z m I} & \texttt{select-\codebreak{}sexp-\codebreak{}inside} & Select what is inside the bracketed expression around point, its brackets left out: point at the start, the mark at the end.\\
\texttt{m A} & \texttt{C-z m A} & \texttt{select-\codebreak{}sexp} & Select the bracketed expression around point, its brackets included: point at the start, the mark at the end.\\
\texttt{"}, \texttt{\$}, \texttt{'} & \texttt{C-z "}, \texttt{C-z \$}, \texttt{C-z '} & \texttt{wrap-\codebreak{}selection-\codebreak{}with-\codebreak{}key} & Wrap the selection in the delimiter pressed and its partner, or take them off when they stand just outside it: ” on a word quotes it, and ” on the word of a quoted word unquotes it.\\
\bottomrule
\end{xltabular}

\begin{xltabular}{\textwidth}{@{}l >{\raggedright\arraybackslash\hsize=0.990\hsize}X >{\raggedright\arraybackslash\hsize=1.010\hsize}X@{}}
\toprule
\textbf{Setting} & \textbf{Standard} & \textbf{What it changes}\\
\midrule
\endhead
\texttt{*pair-\codebreak{}delimiters*} & \texttt{((\#\textbackslash{}\{ . \#\textbackslash{}\}) (\#\textbackslash{}[ . \#\textbackslash{}]) (\#\textbackslash{}( . \#\textbackslash{})) (\#\textbackslash{}...} & The pairs of delimiters, each (OPEN .\\
\bottomrule
\end{xltabular}

\subsubsection{The Emacs keys for the same}
\label{sec:115}
These keys do much the same with Emacs’s names for them. They work in both
states.

\begin{xltabular}{\textwidth}{@{}l >{\raggedright\arraybackslash\hsize=0.553\hsize}X l >{\raggedright\arraybackslash\hsize=1.447\hsize}X@{}}
\toprule
\textbf{Modal} & \textbf{Non-modal} & \textbf{Command} & \textbf{What it does}\\
\midrule
\endhead
none & \texttt{M-@} & \texttt{mark-\codebreak{}word} & Sets the mark at the end of the next word. Again, one word more.\\
none & \texttt{M-h} & \texttt{mark-\codebreak{}paragraph} & Selects the paragraph. Again, the next one too.\\
none & \texttt{C-M-\codebreak{}SPC}, \texttt{C-M-@} & \texttt{mark-\codebreak{}sexp} & Sets the mark after the next bracketed expression.\\
none & \texttt{M-x mark-\codebreak{}end-\codebreak{}of-\codebreak{}sentence} & \texttt{mark-\codebreak{}end-\codebreak{}of-\codebreak{}sentence} & Sets the mark at the end of the sentence. A number sets it that many sentences on.\\
none & \texttt{M-x mark-\codebreak{}defun} & \texttt{mark-\codebreak{}defun} & Selects the definition the cursor is in, in a file of code.\\
\bottomrule
\end{xltabular}

\begin{xltabular}{\textwidth}{@{}l >{\raggedright\arraybackslash\hsize=0.553\hsize}X l >{\raggedright\arraybackslash\hsize=1.447\hsize}X@{}}
\toprule
\textbf{Modal} & \textbf{Non-modal} & \textbf{Command} & \textbf{What it does}\\
\midrule
\endhead
none & \texttt{M-@} & \texttt{mark-\codebreak{}word} & Set the mark at the end of the next word, point staying, and activate the region; with a count, that many words.\\
none & \texttt{M-h} & \texttt{mark-\codebreak{}paragraph} & Select the paragraph at point: point at the blank line before it (or the start of the text), the mark at the blank line after it.\\
none & \texttt{C-M-@}, \texttt{C-M-\codebreak{}SPC} & \texttt{mark-\codebreak{}sexp} & Set the mark after the next balanced expression, point staying, and activate it; with a count, that many expressions, or before point when negative.\\
none & \texttt{M-x mark-\codebreak{}end-\codebreak{}of-\codebreak{}sentence} & \texttt{mark-\codebreak{}end-\codebreak{}of-\codebreak{}sentence} & Set the mark at the end of the sentence, point staying, and activate the region; with a count, that many sentences.\\
none & \texttt{M-x mark-\codebreak{}defun} & \texttt{mark-\codebreak{}defun} & Select the definition point is in, or the one after: point at its start, with the comments right above it and the blank lines above those, and the mark after it.\\
none & \texttt{M-x push-\codebreak{}mark-\codebreak{}command} & \texttt{push-\codebreak{}mark-\codebreak{}command} & Set the mark at point and activate the region, the mark before going on the mark ring; when the mark is at point already, only activate it.\\
none & \texttt{M-x pop-\codebreak{}to-\codebreak{}mark-\codebreak{}command} & \texttt{pop-\codebreak{}to-\codebreak{}mark-\codebreak{}command} & Go to the mark, and take the mark before it off the mark ring as the mark.\\
none & \texttt{C-x C-SPC} & \texttt{pop-\codebreak{}global-\codebreak{}mark} & Go to the newest mark of the global mark ring, in its buffer, that mark going to the ring’s far end.\\
\bottomrule
\end{xltabular}

\subsection{The mark run}
\label{sec:117}
Sometimes one object is not enough. You want a word, then the next word, then
the symbol before it. A mark run lets you keep adding to a selection with single
keys. Press \texttt{M} (non-modal: \texttt{C-z M}), the \texttt{m} held down, and you are in a run.
The keys of a run are the ones you would have pressed after \texttt{m}, without the
\texttt{m}.

\begin{xltabular}{\textwidth}{@{}l l >{\raggedright\arraybackslash\hsize=1.000\hsize}X@{}}
\toprule
\textbf{Modal} & \textbf{Non-modal} & \textbf{What it does in a run}\\
\midrule
\endhead
\texttt{w}, \texttt{b} & none & Add the next word, or the one before.\\
\texttt{W}, \texttt{B} & none & Add a symbol forward or backward.\\
\texttt{s}, \texttt{S} & none & Add a sentence forward or backward.\\
\texttt{m p}, \texttt{m P} & none & Add a paragraph forward or backward.\\
\texttt{h}, \texttt{l}, \texttt{j}, \texttt{k} & none & Move the cursor end by one character or line, without ending the run.\\
\texttt{g h}, \texttt{g l} & none & Stretch the cursor end to the start or end of its line.\\
\texttt{g g}, \texttt{g e} & none & Stretch the cursor end to the start or end of the buffer.\\
\texttt{J}, \texttt{K} & none & Grow the run down or up by whole lines.\\
\texttt{*} & none & Swap the two ends, so that \texttt{h j k l} then move the other end.\\
\texttt{u}, \texttt{U} & none & Take the last step back, or put it forward again.\\
\texttt{M} & none & End the run and let go of the selection.\\
\bottomrule
\end{xltabular}

\texttt{v} and \texttt{V} are refused in a run, because a run is already making a
selection. To take a selection you made another way into a run, so that these
keys grow it from where it is, use \texttt{mark-\codebreak{}run-\codebreak{}adopt}.

\begin{try}
\textbf{Try it, with a shopping list}
\begin{enumerate}
\item Press \texttt{C-x C-f}, type \texttt{\textasciitilde{}/list.\codebreak{}txt} and press \texttt{RET}. Press \texttt{i} and type
   \texttt{fresh bread and butter from the shop}. Press \texttt{C-g} and \texttt{g g g h}.
\item Press \texttt{M}. The first word, \texttt{fresh}, is selected. Press \texttt{w} (non-modal: \texttt{M-f})
   twice. Now \texttt{fresh bread and} is selected.
\item Press \texttt{u} (non-modal: \texttt{C-/}). One step goes back. Press \texttt{U} (non-modal:
   \texttt{C-?}) and it is forward again.
\item Press \texttt{M} to end the run.
\end{enumerate}
\end{try}

The echo area names the keys of the run while it lasts. See also the section
on many cursors below, where the run goes on at every cursor.

\begin{xltabular}{\textwidth}{@{}l >{\raggedright\arraybackslash\hsize=0.578\hsize}X l >{\raggedright\arraybackslash\hsize=1.422\hsize}X@{}}
\toprule
\textbf{Modal} & \textbf{Non-modal} & \textbf{Command} & \textbf{What it does}\\
\midrule
\endhead
\texttt{M} & \texttt{C-z M} & \texttt{mark-\codebreak{}run-\codebreak{}toggle} & Start a mark run: the m prefix held down.\\
none & \texttt{M-x mark-\codebreak{}run-\codebreak{}adopt} & \texttt{mark-\codebreak{}run-\codebreak{}adopt} & Take the selection into a mark run: its keys then grow it from where it is.\\
none & \texttt{M-x mark-\codebreak{}run-\codebreak{}cancel} & \texttt{mark-\codebreak{}run-\codebreak{}cancel} & Drop the selection and end the mark run.\\
none & \texttt{M-x mark-\codebreak{}run-\codebreak{}exchange} & \texttt{mark-\codebreak{}run-\codebreak{}exchange} & Swap point and the mark, so that h j k l adjust the run’s other end.\\
none & \texttt{M-x mark-\codebreak{}run-\codebreak{}step-\codebreak{}back} & \texttt{mark-\codebreak{}run-\codebreak{}step-\codebreak{}back} & Put the mark run back where its last press found it; U puts it forward again.\\
none & \texttt{M-x mark-\codebreak{}run-\codebreak{}step-\codebreak{}forward} & \texttt{mark-\codebreak{}run-\codebreak{}step-\codebreak{}forward} & Put the mark run forward where u stepped it back from.\\
none & \texttt{M-x mark-\codebreak{}run-\codebreak{}left} & \texttt{mark-\codebreak{}run-\codebreak{}left} & Move point back a character, or the count’s, without ending the mark run.\\
none & \texttt{M-x mark-\codebreak{}run-\codebreak{}right} & \texttt{mark-\codebreak{}run-\codebreak{}right} & Move point forward a character, or the count’s, without ending the mark run.\\
none & \texttt{M-x mark-\codebreak{}run-up} & \texttt{mark-\codebreak{}run-up} & Move point up a line, or the count’s, without ending the mark run.\\
none & \texttt{M-x mark-\codebreak{}run-\codebreak{}down} & \texttt{mark-\codebreak{}run-\codebreak{}down} & Move point down a line, or the count’s, without ending the mark run.\\
none & \texttt{M-x mark-\codebreak{}run-\codebreak{}line-\codebreak{}forward} & \texttt{mark-\codebreak{}run-\codebreak{}line-\codebreak{}forward} & Select the line point is on, its line end included, or grow the run’s end down by whole lines: a run ending mid-line first takes the rest of that line.\\
none & \texttt{M-x mark-\codebreak{}run-\codebreak{}line-\codebreak{}backward} & \texttt{mark-\codebreak{}run-\codebreak{}line-\codebreak{}backward} & Select the line point is on, or grow the run’s start up by whole lines: a run starting mid-line first takes the start of that line.\\
none & \texttt{M-x mark-\codebreak{}run-\codebreak{}line-\codebreak{}start} & \texttt{mark-\codebreak{}run-\codebreak{}line-\codebreak{}start} & Stretch the run’s start back to the start of its line, or with no run go there.\\
none & \texttt{M-x mark-\codebreak{}run-\codebreak{}line-\codebreak{}end} & \texttt{mark-\codebreak{}run-\codebreak{}line-\codebreak{}end} & Stretch the run’s end on to the end of its line, or with no run go there.\\
none & \texttt{M-x mark-\codebreak{}run-\codebreak{}buffer-\codebreak{}start} & \texttt{mark-\codebreak{}run-\codebreak{}buffer-\codebreak{}start} & Stretch the run’s start back to the start of the text, or with no run go there.\\
none & \texttt{M-x mark-\codebreak{}run-\codebreak{}buffer-\codebreak{}end} & \texttt{mark-\codebreak{}run-\codebreak{}buffer-\codebreak{}end} & Stretch the run’s end on to the end of the text, or with no run go there.\\
\bottomrule
\end{xltabular}

\subsection{Growing a selection by the code}
\label{sec:118}
In a file whose language Tuatara reads, it knows the structure: which words make
an expression, which expressions make a statement, and so on up to the whole
file. The keys below select by that structure. \texttt{M-o}, \texttt{M-i}, \texttt{M-n} and \texttt{M-p}
work in every state; the \texttt{m} keys after them, in NORMAL state or under \texttt{C-z}.
They need a file that has a syntax tree. In a file with no tree, such as a plain
text file, the echo area says there is none.

\begin{xltabular}{\textwidth}{@{}l l l >{\raggedright\arraybackslash\hsize=1.000\hsize}X@{}}
\toprule
\textbf{Modal} & \textbf{Non-modal} & \textbf{Command} & \textbf{What it does}\\
\midrule
\endhead
\texttt{M-o} & \texttt{M-o} & \texttt{expand-\codebreak{}selection} & Selects the smallest piece of syntax around the selection, or at the cursor. Again, a larger one.\\
\texttt{M-i} & \texttt{M-i} & \texttt{shrink-\codebreak{}selection} & Goes back to the selection before the last \texttt{M-o}.\\
\texttt{M-n} & \texttt{M-n} & \texttt{select-\codebreak{}next-\codebreak{}sibling} & Selects the piece beside it, the next argument, statement or item.\\
\texttt{M-p} & \texttt{M-p} & \texttt{select-\codebreak{}previous-\codebreak{}sibling} & The one before it.\\
\bottomrule
\end{xltabular}

A number before \texttt{M-o} grows that many steps at once.

The \texttt{m} keys have a second set that picks out pieces of code. The small letter
takes what is inside, and the capital the whole piece. Pressed again at once,
they take the next one out.

\begin{xltabular}{\textwidth}{@{}l l >{\raggedright\arraybackslash\hsize=0.818\hsize}X >{\raggedright\arraybackslash\hsize=1.182\hsize}X@{}}
\toprule
\textbf{Modal} & \textbf{Non-modal} & \textbf{Command} & \textbf{What it selects}\\
\midrule
\endhead
\texttt{m f}, \texttt{m F} & \texttt{C-z m f}, \texttt{C-z m F} & \texttt{select-\codebreak{}function-\codebreak{}inside}, \texttt{select-\codebreak{}function} & A function’s body, or all of it.\\
\texttt{m t}, \texttt{m T} & \texttt{C-z m t}, \texttt{C-z m T} & \texttt{select-\codebreak{}type-\codebreak{}inside}, \texttt{select-\codebreak{}type} & The fields or items of a type, or all of it.\\
\texttt{m r}, \texttt{m R} & \texttt{C-z m r}, \texttt{C-z m R} & \texttt{select-\codebreak{}parameter-\codebreak{}inside}, \texttt{select-\codebreak{}parameter} & A parameter or argument, or the same with its comma.\\
\texttt{m c}, \texttt{m C} & \texttt{C-z m c}, \texttt{C-z m C} & \texttt{select-\codebreak{}comment-\codebreak{}inside}, \texttt{select-\codebreak{}comment} & A comment, or with the comment lines next to it.\\
\texttt{m y}, \texttt{m Y} & \texttt{C-z m y}, \texttt{C-z m Y} & \texttt{select-\codebreak{}test-\codebreak{}inside}, \texttt{select-\codebreak{}test} & A test’s body, or the whole test.\\
\texttt{m e}, \texttt{m E} & \texttt{C-z m e}, \texttt{C-z m E} & \texttt{select-\codebreak{}entry-\codebreak{}inside}, \texttt{select-\codebreak{}entry} & The key or value of an entry in JSON, YAML or TOML, or the whole entry.\\
\texttt{m x}, \texttt{m X} & \texttt{C-z m x}, \texttt{C-z m X} & \texttt{select-\codebreak{}element-\codebreak{}inside}, \texttt{select-\codebreak{}element} & What the tags of an HTML or XML element hold, or the element with its tags.\\
\bottomrule
\end{xltabular}

\begin{try}
\textbf{Try it, with a small Python file}
\begin{enumerate}
\item Press \texttt{C-x C-f}, type \texttt{\textasciitilde{}/count.py} and press \texttt{RET}. Press \texttt{i} and type
   \texttt{def count(birds, colour):}, \texttt{RET}, and then, without any spaces of your
   own, \texttt{found = [b for b in birds if b.colour =} colour]=, \texttt{RET},
   \texttt{return len(found)}. Press \texttt{C-g} and \texttt{C-x C-s}.
\item Press \texttt{g g} (non-modal: \texttt{M-<}) and \texttt{w w} to reach \texttt{count}, then \texttt{l}
   (non-modal: \texttt{C-f}) to land on \texttt{birds}.
\item Press \texttt{M-n}. The next parameter, \texttt{colour}, is selected. Press \texttt{M-p} to go
   back to \texttt{birds}.
\item Press \texttt{M-o}. The parameters are selected, with their brackets. Press \texttt{M-o}
   again. The whole function is selected, which in this small file is also
   the whole file. Press \texttt{M-o} once more and the echo area says that nothing
   larger can be selected.
\item Press \texttt{M-i} until you are on \texttt{birds} again. Press \texttt{C-g}.
\item Press \texttt{j} (non-modal: \texttt{C-n}) to go down into the body. Press \texttt{m f}. The body
   of the function is selected; on the \texttt{def} line itself it says there is no
   function around point. Press \texttt{C-g} and then \texttt{m F} for the whole function.
\end{enumerate}
\end{try}

\begin{xltabular}{\textwidth}{@{}l l l >{\raggedright\arraybackslash\hsize=1.000\hsize}X@{}}
\toprule
\textbf{Modal} & \textbf{Non-modal} & \textbf{Command} & \textbf{What it does}\\
\midrule
\endhead
none & \texttt{M-o} & \texttt{expand-\codebreak{}selection} & Select the smallest piece of syntax around the selection, or around the character at point: the word, then the expression holding it, then the statement, and so on out; in a buffer with no syntax tree, as its language server sees the code.\\
none & \texttt{M-i} & \texttt{shrink-\codebreak{}selection} & Go back to the selection before the last M-o, or else select the first piece of syntax inside the selection; without a selection, the smallest piece of syntax at point.\\
none & \texttt{M-n} & \texttt{select-\codebreak{}next-\codebreak{}sibling} & Select the next piece of syntax beside the selection, or beside the one at point: the next argument, statement or item, and after the last one, the next beside the piece around them.\\
none & \texttt{M-p} & \texttt{select-\codebreak{}previous-\codebreak{}sibling} & Select the piece of syntax before the selection, or before the one at point: the argument, statement or item before it, and before the first one, the one before the piece around them.\\
\texttt{m f} & \texttt{C-z m f} & \texttt{select-\codebreak{}function-\codebreak{}inside} & Select the body of the function around point: the block of a Rust fn or closure, the body of a shell function, the forms of a Lisp defun, defmacro, defgeneric, defmethod or lambda after its lambda list.\\
\texttt{m F} & \texttt{C-z m F} & \texttt{select-\codebreak{}function} & Select all of the function around point.\\
\texttt{m t} & \texttt{C-z m t} & \texttt{select-\codebreak{}type-\codebreak{}inside} & Select the body of the type around point: the fields, variants or items of a Rust struct, enum, union, trait or impl; in Lisp, what follows the name in a defclass, defstruct, define-condition or deftype.\\
\texttt{m T} & \texttt{C-z m T} & \texttt{select-\codebreak{}type} & Select all of the type around point.\\
\texttt{m r} & \texttt{C-z m r} & \texttt{select-\codebreak{}parameter-\codebreak{}inside} & Select the parameter or argument around point: of a Rust function’s parameters, a call’s arguments or a generic’s parameters; of a shell command’s arguments; of a Lisp lambda list or a call’s arguments.\\
\texttt{m R} & \texttt{C-z m R} & \texttt{select-\codebreak{}parameter} & Select the parameter or argument around point with the comma after it.\\
\texttt{m c} & \texttt{C-z m c} & \texttt{select-\codebreak{}comment-\codebreak{}inside} & Select the comment around point.\\
\texttt{m C} & \texttt{C-z m C} & \texttt{select-\codebreak{}comment} & Select the comment around point with the comment lines right before and after it.\\
\texttt{m y} & \texttt{C-z m y} & \texttt{select-\codebreak{}test-\codebreak{}inside} & Select the body of the test around point: a Rust function marked \#[test], a Python function whose name begins with test, a Go TestX or BenchmarkX, and the tests of Zig, Elixir, Dart, Julia and others, as the language’s rules know them.\\
\texttt{m Y} & \texttt{C-z m Y} & \texttt{select-\codebreak{}test} & Select all of the test around point.\\
\texttt{m e} & \texttt{C-z m e} & \texttt{select-\codebreak{}entry-\codebreak{}inside} & Select the key or the value at point of the entry around it: of a key and its value in JSON, YAML or TOML, a field of a struct, an item of a list, as the language’s rules know them.\\
\texttt{m E} & \texttt{C-z m E} & \texttt{select-\codebreak{}entry} & Select all of the entry around point, its key and value or the list’s item.\\
\texttt{m x} & \texttt{C-z m x} & \texttt{select-\codebreak{}element-\codebreak{}inside} & Select what is between the tags of the element around point: in HTML, XML, JSX, Svelte and Vue.\\
\texttt{m X} & \texttt{C-z m X} & \texttt{select-\codebreak{}element} & Select all of the element around point, its tags too.\\
\bottomrule
\end{xltabular}

\subsection{Rectangles}
\label{sec:122}
A rectangle is a block of columns over several lines. It is the right tool
for a column of numbers, a bullet at the start of each line, or text that
lines up. Figure 5.2 shows one.

\begin{figure}[H]
\caption{Figure 5.2. A rectangle made with \texttt{m v} and \texttt{j j j l l l}. The block is the first four columns of four lines.}\label{org119}
\begin{Verbatim}[breaklines,breakautoindent=false,breakindent=1em,breakafter=/-.\,;:=,commandchars=\\\{\}]
   column  0123456789
   line 1  [The ]tui sings
   line 2  [The ]kereru eats
   line 3  [The ]fantail is small
   line 4  [The ]morepork calls
\end{Verbatim}
\end{figure}

A rectangle covers the columns between the cursor and the mark on every line
from one to the other. A line that ends inside the block gives only the part
it has. The character under the cursor is part of the block, so three \texttt{l}
(non-modal: \texttt{C-f}) presses from the first column cover four columns.

\subsubsection{Rectangles in NORMAL state}
\label{sec:120}
\texttt{m v} (non-modal: \texttt{C-x r SPC}) starts a rectangle in NORMAL state, and any move
sizes it. \texttt{y} copies the block, \texttt{d} cuts it and \texttt{c} cuts it and starts INSERT
state with a cursor on every row, so that what you type goes on every line. \texttt{P}
(non-modal: \texttt{C-x r y}) pastes the block back as a block, one row to a line.
\texttt{m v} again, or \texttt{C-g}, ends the rectangle without doing anything.

\begin{try}
\textbf{Try it, with a price list}
\begin{enumerate}
\item Press \texttt{C-x C-f}, type \texttt{\textasciitilde{}/prices.\codebreak{}txt} and press \texttt{RET}. Press \texttt{i} and type
   four lines, each followed by \texttt{RET} but the last:
   \texttt{coffee  4.50}, \texttt{scone   5.50}, \texttt{pie     7.00} and \texttt{tea     3.80}. Each
   price starts in the ninth column. Press \texttt{C-g}.
\item Press \texttt{g g g h}, then \texttt{l} (non-modal: \texttt{C-f}) eight times. The cursor is on
   the \texttt{4}. Press \texttt{m v} (non-modal: \texttt{C-x r SPC}), then \texttt{j j j l l l}. The four
   prices are the block. The echo area says any move sizes the block.
\item Press \texttt{y}. The echo area says it copied a rectangle of four rows. Press
   \texttt{g e} and \texttt{P} (non-modal: \texttt{M->}, \texttt{C-x r y}). The prices are pasted again at
   the end of the last line, one row to a line, below it. Press \texttt{u} (non-modal:
   \texttt{C-/}) to undo.
\item Press \texttt{g g g h} and \texttt{l} eight times, then \texttt{m v j j j l l l} and \texttt{d}. The
   prices are cut. Press \texttt{g g g h} and \texttt{P}. The column of prices has moved
   to the front of the lines. Press \texttt{u} twice to bring it all back.
\item With the block made again, press \texttt{c}, type \texttt{n/a} and press \texttt{C-g} twice.
   Every price is now \texttt{n/a}. Press \texttt{u}.
\end{enumerate}
\end{try}

\subsubsection{Rectangles with the Emacs keys}
\label{sec:121}
The keys under \texttt{C-x r} work on the block between the cursor and the mark. You
can start a rectangle with \texttt{C-x r SPC}, which does what \texttt{m v} does. Pressing it
again, or \texttt{C-g}, cancels. Outside NORMAL state a new rectangle starts with no
width, and \texttt{C-f} and \texttt{C-n} (modal: \texttt{l}, \texttt{j}) size it.

\begin{xltabular}{\textwidth}{@{}l >{\raggedright\arraybackslash\hsize=0.865\hsize}X l >{\raggedright\arraybackslash\hsize=1.135\hsize}X@{}}
\toprule
\textbf{Modal} & \textbf{Non-modal} & \textbf{Command} & \textbf{What it does}\\
\midrule
\endhead
none & \texttt{C-x r k} & \texttt{kill-\codebreak{}rectangle} & Cuts the block to a store of its own, apart from the kill ring.\\
none & \texttt{C-x r M-w} & \texttt{copy-\codebreak{}rectangle-\codebreak{}as-\codebreak{}kill} & Copies the block to that store.\\
none & \texttt{C-x r d} & \texttt{delete-\codebreak{}rectangle} & Deletes the block and keeps nothing. The store keeps what it had.\\
\texttt{P} & \texttt{C-x r y} & \texttt{yank-\codebreak{}rectangle} & Pastes the store as a block.\\
none & \texttt{C-x r t} & \texttt{string-\codebreak{}rectangle} & Asks for a text and puts it on every row in place of the block.\\
none & \texttt{C-x r c} & \texttt{clear-\codebreak{}rectangle} & Turns the block into spaces.\\
none & \texttt{C-x r N} & \texttt{number-\codebreak{}rectangle-\codebreak{}lines} & Numbers the rows 1, 2, 3 at the block’s left column.\\
none & \texttt{M-x number-\codebreak{}rectangle-\codebreak{}lines-\codebreak{}replace} & \texttt{number-\codebreak{}rectangle-\codebreak{}lines-\codebreak{}replace} & Puts numbers in place of the block, with nothing after them.\\
\bottomrule
\end{xltabular}

\begin{try}
\textbf{Try it, with the same price list}
\begin{enumerate}
\item Open \texttt{\textasciitilde{}/prices.\codebreak{}txt} again. Press \texttt{g g g h} and \texttt{m v j j j}. The block is
   the first column of the four lines.
\item Press \texttt{C-x r N}. The lines are numbered 1 to 4 and the text moves over.
   Press \texttt{u} (non-modal: \texttt{C-/}).
\item Make the block again and press \texttt{C-x r t}. Type =- = and press \texttt{RET}. Every
   line now starts with a dash and a space, which take the place of its first
   letter. Press \texttt{u}.
\item Press \texttt{g g g h} and \texttt{l} (non-modal: \texttt{C-f}) eight times, then
   \texttt{m v j j j l l l} and \texttt{C-x r c}. The prices are blanked and the lines are
   shorter. Press \texttt{u}.
\item Make the block once more and press \texttt{C-x r k}. Press \texttt{g e} (non-modal: \texttt{M->}),
   then \texttt{C-x r y} (modal: \texttt{P}). The prices are pasted as a block at the end.
   Press \texttt{u u}.
\end{enumerate}
\end{try}

\begin{xltabular}{\textwidth}{@{}l >{\raggedright\arraybackslash\hsize=0.865\hsize}X l >{\raggedright\arraybackslash\hsize=1.135\hsize}X@{}}
\toprule
\textbf{Modal} & \textbf{Non-modal} & \textbf{Command} & \textbf{What it does}\\
\midrule
\endhead
\texttt{m v} & \texttt{C-x r SPC} & \texttt{rectangle-\codebreak{}mark-\codebreak{}mode} & Start a rectangle: the block of columns between point and the mark on each line from one to the other; any motion sizes it, and its corner can stand past the end of a short line.\\
none & \texttt{C-x r k} & \texttt{kill-\codebreak{}rectangle} & Kill the rectangle between the mark and point: the block of columns from one to the other on each line between, a short line’s part padded with spaces, goes to the rectangle store that C-x r y pastes from.\\
none & \texttt{C-x r M-w} & \texttt{copy-\codebreak{}rectangle-\codebreak{}as-\codebreak{}kill} & Copy the rectangle between the mark and point to the rectangle store that C-x r y pastes from, the text staying as it is.\\
none & \texttt{C-x r d} & \texttt{delete-\codebreak{}rectangle} & Delete the rectangle between the mark and point, keeping it nowhere: the rectangle store keeps what it had.\\
\texttt{P} & \texttt{C-x r y} & \texttt{yank-\codebreak{}rectangle} & Paste the rectangle store (what y, d or c last took from a rectangle) as a block of columns: from point, one row a line, or over a selected rectangle of as many rows in its place.\\
none & \texttt{C-x r t} & \texttt{string-\codebreak{}rectangle} & Replace the rectangle between the mark and point with a text read in the minibuffer, on each of its lines; with no width, the text goes in at its column on each line.\\
none & \texttt{C-x r c} & \texttt{clear-\codebreak{}rectangle} & Blank out the rectangle between the mark and point: its columns become spaces on each line, so that what comes after stays where it was; a line that ends inside the block is cut at its left column instead, leaving no spaces at its end, and a line shorter than the block is left alone.\\
none & \texttt{C-x r N} & \texttt{number-\codebreak{}rectangle-\codebreak{}lines} & Number the lines of the rectangle between the mark and point: 1, 2, 3 and on, each with a space after it, put in at the block’s left column, a line shorter than that padded out to it; the numbers are right-aligned to the width of the last, and the text of the block moves right to make room.\\
none & \texttt{M-x number-\codebreak{}rectangle-\codebreak{}lines-\codebreak{}replace} & \texttt{number-\codebreak{}rectangle-\codebreak{}lines-\codebreak{}replace} & Replace the rectangle between the mark and point with numbers, 1, 2, 3 and on down its lines, nothing after them: to renumber a list, or to turn its bullets into numbers.\\
\bottomrule
\end{xltabular}

\begin{xltabular}{\textwidth}{@{}l l >{\raggedright\arraybackslash\hsize=1.000\hsize}X@{}}
\toprule
\textbf{Setting} & \textbf{Standard} & \textbf{What it changes}\\
\midrule
\endhead
\texttt{*rectangle-\codebreak{}cursors-\codebreak{}limit*} & \texttt{10000} & The most rows M-x tuatara-change on a rectangle (M-x rectangle-mark-mode) puts a cursor on, or NIL for as many as the memory has room for: the block is emptied and what is typed appears on every row, as at the cursors of M-x add-cursor-below.\\
\bottomrule
\end{xltabular}

\subsection{Many cursors}
\label{sec:127}
\texttt{t} (non-modal: \texttt{C-z t}) puts another cursor on the line below, at the same
column. \texttt{T} (non-modal: \texttt{C-z T}) puts one above. From then on, nearly every key
you press works at every cursor. Move, select, type, delete: each key runs once
at each place. This is the fast way to do the same small edit on many lines.

Figure 5.3 shows four cursors made with three presses of \texttt{t}.

\begin{figure}[H]
\caption{Figure 5.3. Cursors after \texttt{t t t} from the start of the first line. A cursor is shown as a vertical bar.}\label{org123}
\begin{Verbatim}[breaklines,breakautoindent=false,breakindent=1em,breakafter=/-.\,;:=,commandchars=\\\{\}]
   |Hana Raglan
   |Aroha Kerikeri
   |Tane Napier
   |Mere Taupo
\end{Verbatim}
\end{figure}

A line shorter than the column gets its cursor at its end. If lines are
selected, the first \texttt{t} puts one cursor on every selected line instead. At the
start of each line when the selection was made with \texttt{V}, or in the cursor’s
column otherwise.

Press \texttt{i} and type, and the text appears at every cursor. Press \texttt{C-g} and you
are back in NORMAL state with the cursors still there. Press \texttt{C-g} again and
the cursors end. One \texttt{u} (non-modal: \texttt{C-/}) takes back what was done at all of
them.

\begin{try}
\textbf{Try it, with a list of names}
\begin{enumerate}
\item Press \texttt{C-x C-f}, type \texttt{\textasciitilde{}/names.\codebreak{}txt} and press \texttt{RET}. Press \texttt{i} and type four
   lines, each followed by \texttt{RET} but the last: \texttt{Hana Raglan}, \texttt{Aroha Kerikeri},
   \texttt{Tane Napier} and \texttt{Mere Taupo}. Press \texttt{C-g}, then \texttt{g g}.
\item Press \texttt{t t t}. There are four cursors, one on each line. Press \texttt{i}, type
   =- = and press \texttt{C-g}. Each name has a dash in front. Press \texttt{C-g} to end the
   cursors. Press \texttt{u} (non-modal: \texttt{C-/}) to take the dashes back.
\item Press \texttt{g g t t t} then \texttt{DEL}. The last cursor goes and three remain.
   Press \texttt{C-g} twice.
\item Press \texttt{g g t t t}, then \texttt{m w}. Each cursor selects the word it is on. Press
   \texttt{M-u} to put them in capitals. Press \texttt{u}, then \texttt{C-g} twice.
\item Press \texttt{g g t t t}, then \texttt{A}. Type \texttt{.} and press \texttt{C-g} twice. Each line now
   ends with a full stop.
\end{enumerate}
\end{try}

\subsubsection{What each key does at the cursors}
\label{sec:124}
While the cursors stand, a key in NORMAL state does what it always does, at
each cursor. The table lists the keys that need a word of explanation.

\begin{xltabular}{\textwidth}{@{}l l >{\raggedright\arraybackslash\hsize=1.000\hsize}X@{}}
\toprule
\textbf{Modal} & \textbf{Non-modal} & \textbf{What it does at every cursor}\\
\midrule
\endhead
\texttt{t}, \texttt{T} & none & Add one more cursor below or above.\\
\texttt{DEL} & none & Takes back the cursor added last.\\
moves, \texttt{m} keys & none & Move or select at each cursor on its own line.\\
\texttt{v}, \texttt{V} & none & Start a selection at each, or select each cursor’s line.\\
\texttt{i}, \texttt{a} & none & Type before or after the character, or the selection.\\
\texttt{I}, \texttt{A} & none & Type at the start or the end of the line.\\
\texttt{o}, \texttt{O} & none & Open a line below or above each and type on them all.\\
\texttt{c} & none & Empty each selection, or as many characters as a number says, then type.\\
\texttt{d}, \texttt{D} & none & Delete each selection or character, or kill to the end of the line.\\
\texttt{y}, \texttt{p} & none & Copy every selection as one kill, a line to a cursor. Paste: a kill of as many parts as there are cursors gives each its own.\\
\texttt{>}, \texttt{C} & none & Indent each line, or comment or uncomment it.\\
\texttt{M-u}, \texttt{M-l}, \texttt{M-c} & none & Put each selection in capitals, small letters, or capitalised.\\
\texttt{u}, \texttt{U} & none & Undo or redo, the cursors kept.\\
\texttt{.} & none & Run the last command again at every cursor.\\
\texttt{=} & none & Bring every cursor to the column of the first.\\
\texttt{M} & none & Start a mark run at every cursor, or end it.\\
\texttt{C-g} & none & Let go of the selections. Again, end the cursors.\\
\bottomrule
\end{xltabular}

A key the cursors do not know ends them and then does its own work.

The copy and paste rule is worth a moment. Say you select four lines, one at
each of four cursors, and press \texttt{y}. The four lines are one kill. \texttt{p} at the
same four cursors puts each line back at its own cursor. If you paste a
kill that has as many lines as there are cursors, each cursor gets one line.
In every other case each cursor gets the whole kill.

\texttt{M} at the cursors starts a mark run at each one. Each run begins as the
word at its cursor, or as the selection it had. \texttt{w} (non-modal: \texttt{M-f}) then adds
a word to every run at once, and \texttt{u} (non-modal: \texttt{C-/}) takes the last step back
at all of them.

\subsubsection{Limits}
\label{sec:125}
Every key but typing runs once at every cursor, so a great many cursors make
every key slow. Tuatara holds up to \texttt{*cursors-\codebreak{}limit*} cursors. A cursor in
sight beyond the limit still works as one. Those out of sight are written to in
the background as each thing is done, and they move only together: \texttt{h}, \texttt{l},
\texttt{g h} and \texttt{g l} (non-modal: \texttt{C-b}, \texttt{C-f}, \texttt{C-a}, \texttt{C-e}) move them all alike, and
\texttt{i}, \texttt{a}, \texttt{I}, \texttt{A}, \texttt{c} and \texttt{d} edit at them. Any other key says what works and
leaves them as they were.

\begin{xltabular}{\textwidth}{@{}l >{\raggedright\arraybackslash\hsize=0.744\hsize}X l >{\raggedright\arraybackslash\hsize=1.256\hsize}X@{}}
\toprule
\textbf{Modal} & \textbf{Non-modal} & \textbf{Command} & \textbf{What it does}\\
\midrule
\endhead
\texttt{t} & \texttt{C-z t} & \texttt{add-\codebreak{}cursor-\codebreak{}below} & Add a cursor on the line below, at the same column.\\
\texttt{T} & \texttt{C-z T} & \texttt{add-\codebreak{}cursor-\codebreak{}above} & Add a cursor on the line above, at the same column: the first press makes cursors of point and the line above, each one after adds one above the first, and a count adds that many (negative, below).\\
none & \texttt{M-x drop-\codebreak{}last-\codebreak{}cursor} & \texttt{drop-\codebreak{}last-\codebreak{}cursor} & Take back the cursor added last, keeping the others; the real cursor is never taken back, and C-g ends the cursors.\\
none & \texttt{M-x align-\codebreak{}cursors} & \texttt{align-\codebreak{}cursors} & Put every cursor at the real cursor’s column, or its line’s end where the line is shorter; with cursors selecting, wrap the selections in = as any delimiter does.\\
none & \texttt{M-x cursors-\codebreak{}select} & \texttt{cursors-\codebreak{}select} & Start a selection at every cursor, where each is: the motions then grow them.\\
none & \texttt{M-x cursors-\codebreak{}select-\codebreak{}lines} & \texttt{cursors-\codebreak{}select-\codebreak{}lines} & Select every cursor’s whole line, or let go where every cursor has.\\
none & \texttt{M-x cursors-\codebreak{}insert} & \texttt{cursors-\codebreak{}insert} & Type before every cursor, or its selection: what is typed at one appears at all.\\
none & \texttt{M-x cursors-\codebreak{}append} & \texttt{cursors-\codebreak{}append} & Type after the character under every cursor, or after its selection; a cursor at its line’s end types there.\\
none & \texttt{M-x cursors-\codebreak{}insert-\codebreak{}line} & \texttt{cursors-\codebreak{}insert-\codebreak{}line} & Type at the start of every cursor’s line.\\
none & \texttt{M-x cursors-\codebreak{}append-\codebreak{}line} & \texttt{cursors-\codebreak{}append-\codebreak{}line} & Type at the end of every cursor’s line.\\
none & \texttt{M-x cursors-\codebreak{}open-\codebreak{}below} & \texttt{cursors-\codebreak{}open-\codebreak{}below} & Open a line below every cursor’s line and type on all of them.\\
none & \texttt{M-x cursors-\codebreak{}open-\codebreak{}above} & \texttt{cursors-\codebreak{}open-\codebreak{}above} & Open a line above every cursor’s line and type on all of them.\\
none & \texttt{M-x cursors-\codebreak{}line-\codebreak{}end} & \texttt{cursors-\codebreak{}line-\codebreak{}end} & Put every cursor at the end of its line, keeping its selection.\\
none & \texttt{M-x cursors-\codebreak{}change} & \texttt{cursors-\codebreak{}change} & Empty every cursor’s selection, or the count’s characters at every cursor, then type there; a selection goes on the kill ring as one kill.\\
none & \texttt{M-x cursors-\codebreak{}delete} & \texttt{cursors-\codebreak{}delete} & Delete every cursor’s selection, a selected line with its newline, or the count’s characters at every cursor; one change, the cursors kept.\\
none & \texttt{M-x cursors-\codebreak{}kill-\codebreak{}line} & \texttt{cursors-\codebreak{}kill-\codebreak{}line} & Kill from every cursor to the end of its line, as one kill and one change; a cursor at its line’s end kills nothing.\\
none & \texttt{M-x cursors-\codebreak{}copy} & \texttt{cursors-\codebreak{}copy} & Copy every cursor’s selection, or the count’s characters at every cursor, as one kill, one line a cursor; the selections are let go.\\
none & \texttt{M-x cursors-\codebreak{}yank} & \texttt{cursors-\codebreak{}yank} & Paste at every cursor, over its selection.\\
none & \texttt{M-x cursors-\codebreak{}indent} & \texttt{cursors-\codebreak{}indent} & Indent every cursor’s line as the mode indents it.\\
none & \texttt{M-x cursors-\codebreak{}comment} & \texttt{cursors-\codebreak{}comment} & Comment or uncomment every cursor’s line.\\
none & \texttt{M-x cursors-\codebreak{}upcase} & \texttt{cursors-\codebreak{}upcase} & Put every cursor’s selection in upper case, or the count’s words from every cursor.\\
none & \texttt{M-x cursors-\codebreak{}downcase} & \texttt{cursors-\codebreak{}downcase} & Put every cursor’s selection in lower case, or the count’s words from every cursor.\\
none & \texttt{M-x cursors-\codebreak{}capitalise} & \texttt{cursors-\codebreak{}capitalise} & Capitalise every cursor’s selection, or the count’s words from every cursor.\\
none & \texttt{M-x cursors-\codebreak{}undo} & \texttt{cursors-\codebreak{}undo} & Undo, keeping the cursors: a change made at them all is taken back at all.\\
none & \texttt{M-x cursors-\codebreak{}redo} & \texttt{cursors-\codebreak{}redo} & Redo, keeping the cursors.\\
none & \texttt{M-x cursors-\codebreak{}repeat} & \texttt{cursors-\codebreak{}repeat} & Run the cursors’ last command again at every cursor, with the count it had: a motion, a selection, or deleting, pasting, killing and banking.\\
none & \texttt{M-x cursors-\codebreak{}replay} & \texttt{cursors-\codebreak{}replay} & Run the command of the keys just pressed at every cursor: the NORMAL state motion or selection they are on, or in a mark run, the run’s.\\
none & \texttt{M-x cursors-\codebreak{}run-\codebreak{}toggle} & \texttt{cursors-\codebreak{}run-\codebreak{}toggle} & Start a mark run at every cursor, or end it and let go of the selections: each cursor keeps what it selected, or selects its word, and the run’s keys w W b B s S grow and h l g h g l * adjust every selection at once; u and U step them back and forward.\\
none & \texttt{M-x cursors-\codebreak{}run-\codebreak{}step-\codebreak{}back} & \texttt{cursors-\codebreak{}run-\codebreak{}step-\codebreak{}back} & Put every cursor’s selection back where the run’s last press found it.\\
none & \texttt{M-x cursors-\codebreak{}run-\codebreak{}step-\codebreak{}forward} & \texttt{cursors-\codebreak{}run-\codebreak{}step-\codebreak{}forward} & Put every cursor’s selection forward again, a step u took back.\\
none & \texttt{M-x cursors-\codebreak{}move-\codebreak{}past-\codebreak{}limit} & \texttt{cursors-\codebreak{}move-\codebreak{}past-\codebreak{}limit} & Move cursors past their limit, those out of sight with those in sight, by the NORMAL state motion of the keys pressed: a character right or left (with a count, that many), or to the start or the end of the line.\\
none & \texttt{M-x cursors-\codebreak{}bank} & \texttt{cursors-\codebreak{}bank} & Bank every cursor’s line, or take them all out where every one is banked; the bank is there for y, d and f once the cursors end.\\
none & \texttt{M-x cursors-\codebreak{}unbank} & \texttt{cursors-\codebreak{}unbank} & Take every cursor’s line out of the bank, leaving every other bank alone.\\
\bottomrule
\end{xltabular}

\begin{xltabular}{\textwidth}{@{}l l >{\raggedright\arraybackslash\hsize=1.000\hsize}X@{}}
\toprule
\textbf{Setting} & \textbf{Standard} & \textbf{What it changes}\\
\midrule
\endhead
\texttt{*cursors-\codebreak{}limit*} & \texttt{10000} & The most cursors M-x add-cursor-below and M-x add-cursor-above make, or NIL for as many as the memory has room for: every key but typing runs once at every cursor, so a great many make every key slow; what is typed is written as it is typed only at the cursors in sight, the others when the typing ends.\\
\bottomrule
\end{xltabular}

\subsubsection{Showing the cursors}
\label{sec:126}
The cursors other than the real one are drawn in their own colours, one for
NORMAL state and one while you type. The places a split holds, described next,
have a face as well.

\begin{xltabular}{\textwidth}{@{}l >{\raggedright\arraybackslash\hsize=1.000\hsize}X@{}}
\toprule
\textbf{Face} & \textbf{What it paints}\\
\midrule
\endhead
\texttt{:region} & The active region (selection).\\
\texttt{:split-\codebreak{}cursor} & The cursors of a split other than the real one, in NORMAL state.\\
\texttt{:split-\codebreak{}cursor-\codebreak{}insert} & The cursors of a split other than the real one, in INSERT state.\\
\texttt{:split-\codebreak{}place} & The places a split holds: the matches of f, what the cursors select.\\
\bottomrule
\end{xltabular}

\subsection{Splitting a selection on what matches}
\label{sec:128}
A pattern can make the cursors for you. Select some text and press \texttt{f}
(non-modal: \texttt{C-z f}). Tuatara asks for a regular expression, a pattern as used
in a search (chapter 6). It holds every match inside the selection at once, one
place for each, and waits for you to say what to do with them.

\begin{xltabular}{\textwidth}{@{}l l l >{\raggedright\arraybackslash\hsize=1.000\hsize}X@{}}
\toprule
\textbf{Modal} & \textbf{Non-modal} & \textbf{Command} & \textbf{What it does with every match}\\
\midrule
\endhead
\texttt{i} & none & \texttt{split-\codebreak{}insert} & Types before every match.\\
\texttt{a} & none & \texttt{split-\codebreak{}append} & Types after every match.\\
\texttt{c} & none & \texttt{split-\codebreak{}change} & Empties every match and types there.\\
\texttt{d} & none & \texttt{split-\codebreak{}delete} & Deletes every match.\\
a bracket or quote & none & \texttt{split-\codebreak{}wrap-\codebreak{}key} & Wraps every match in the pair. Press it again to take it off.\\
\texttt{C-g} & none & \texttt{split-\codebreak{}quit} & Lets go and ends the split.\\
\bottomrule
\end{xltabular}

Without a selection, \texttt{f} searches the line the cursor is on. If you put \texttt{\%} in
front of the pattern, as in \texttt{\%kereru}, the whole buffer is searched whatever is
selected. What you type at one match appears at all of them, and the writing
is one change, so \texttt{u} (non-modal: \texttt{C-/}) takes it all back. Case is ignored when
the pattern has no capital letter. As you type the pattern, the number of
matches shows before the prompt.

If the selection is a rectangle, \texttt{f} searches its block. With a number first
it searches each row’s whole line.

\begin{try}
\textbf{Try it, with the recipe}
\begin{enumerate}
\item Open \texttt{\textasciitilde{}/pancakes.\codebreak{}txt}. Press \texttt{g g j V} to select the second line, the one
   with the quantities.
\item Press \texttt{f}, type \texttt{[0-9]+} and press \texttt{RET}. The three numbers are held. The
   echo area says how many places it holds and what you can do.
\item Press \texttt{(}. Each number is wrapped in brackets. Press \texttt{u}.
\item Press \texttt{V}, then \texttt{f}, type \texttt{,} and press \texttt{RET}. Press \texttt{c}, type \texttt{;} and press
   \texttt{C-g}. The commas are now semicolons. Press \texttt{u} (non-modal: \texttt{C-/}).
\item Press \texttt{\%} and \texttt{f}, type \texttt{pancake}, press \texttt{RET} and then \texttt{i}. Type \texttt{*} and
   press \texttt{C-g}. The word \texttt{Pancakes} has a star in front of it. Press \texttt{u}.
\end{enumerate}
\end{try}

\texttt{f} may miss a match of more than 4 KB that runs across several lines. For
text where that matters, \texttt{split-\codebreak{}selection-\codebreak{}thorough} finds every match of up
to a megabyte. It is slower, so it has no key of its own. A pattern with
many matches is limited by \texttt{*split-\codebreak{}limit*}, and \texttt{*split-\codebreak{}limit-\codebreak{}ask*} lets you
hold them all even so.

\begin{xltabular}{\textwidth}{@{}l >{\raggedright\arraybackslash\hsize=0.744\hsize}X l >{\raggedright\arraybackslash\hsize=1.256\hsize}X@{}}
\toprule
\textbf{Modal} & \textbf{Non-modal} & \textbf{Command} & \textbf{What it does}\\
\midrule
\endhead
\texttt{f} & \texttt{C-z f} & \texttt{split-\codebreak{}selection} & Hold every match of a regexp in the selection, then wait for a verb: i types before every match and a after, c empties them first, d deletes them, w or any delimiter wraps them (again, to take the pair off), and C-g lets go.\\
none & \texttt{M-x split-\codebreak{}selection-\codebreak{}thorough} & \texttt{split-\codebreak{}selection-\codebreak{}thorough} & Hold every match of a regexp in the selection, as f does, finding every match of up to a megabyte (about a million characters of ASCII) whole however many lines it spans, where f may miss one of more than 4 KB that runs across lines and needs text further on; the writing out of sight the same.\\
none & \texttt{M-x split-\codebreak{}wrap} & \texttt{split-\codebreak{}wrap} & Wrap every place of the split in the pair of the delimiter asked for, or take the pair off where it stands outside every place already.\\
\bottomrule
\end{xltabular}

\begin{xltabular}{\textwidth}{@{}l l >{\raggedright\arraybackslash\hsize=1.000\hsize}X@{}}
\toprule
\textbf{Setting} & \textbf{Standard} & \textbf{What it changes}\\
\midrule
\endhead
\texttt{*split-\codebreak{}limit*} & \texttt{100000} & The most places M-x split-selection holds, or NIL for as many as the memory has room for, 100000 as standard.\\
\texttt{*split-\codebreak{}limit-\codebreak{}ask*} & \texttt{nil} & True to be asked, when the places would pass their limit, whether to hold them all the same, the memory having room, rather than those in sight only.\\
\bottomrule
\end{xltabular}

\subsection{Summary of the NORMAL keys}
\label{sec:129}
For quick reference, in one place.

\begin{xltabular}{\textwidth}{@{}>{\raggedright\arraybackslash\hsize=0.396\hsize}X l >{\raggedright\arraybackslash\hsize=1.604\hsize}X@{}}
\toprule
\textbf{Modal} & \textbf{Non-modal} & \textbf{Use them to}\\
\midrule
\endhead
\texttt{v}, \texttt{V}, \texttt{J}, \texttt{K}, \texttt{\%} & none & Select by character, by line, or everything.\\
\texttt{m w}, \texttt{m s}, \texttt{m p}, \texttt{m i}, \texttt{m a} & none & Select by what the text is.\\
\texttt{M-o}, \texttt{M-i}, \texttt{M-n}, \texttt{M-p} & the same & Select by what the code is.\\
\texttt{M} & none & Keep adding to a selection with single keys.\\
\texttt{m v} & none & Start a rectangle.\\
\texttt{t}, \texttt{T} & none & Add a cursor below or above.\\
\texttt{f} & none & Hold every match of a pattern in the selection.\\
\texttt{C-g} & none & Let go. Press it again to end the cursors.\\
\bottomrule
\end{xltabular}

\section{Searching and replacing}
\label{sec:148}
Finding text, changing it, and narrowing what you look at to the lines that
matter.

A search answers the question “where is this?”. A replacement answers “change
this everywhere”. Both start from a pattern, a short description of the text
you want. This chapter begins with the search you will use all day, then
covers replacing, narrowing a list to its matches, searching a whole folder
of files, and going to a heading. A table at the end of each section lists
every command, setting and face that belongs to it.

\subsection{Searching a buffer}
\label{sec:139}
Press \texttt{C-s} and Tuatara asks \texttt{Regex search forward:} in the minibuffer, the
same row where it asks for a file name. Type a few letters. The cursor jumps
to the first match while you type, and the matches in sight are lit. \texttt{C-r}
does the same going backward, towards the start of the buffer.

You do not need to finish the word. Typing \texttt{kere} finds the first place that
has those four letters, and typing \texttt{kereru} moves to the first place with all
six. If you make a mistake, \texttt{DEL} takes a letter off the end and the cursor
goes back to where the shorter pattern matched.

Three keys end a search.

\begin{itemize}
\item \texttt{RET} ends it at the match. The cursor stays there, and the place where you
  began is saved as the mark, so \texttt{C-x C-x} takes you back.
\item \texttt{C-g} or \texttt{ESC} gives up. The cursor goes back to where you began.
\item \texttt{C-s} or \texttt{C-r} do not end it. They go to the next match, or the one before.
  At the last match, \texttt{C-s} goes round to the first.
\end{itemize}

Straight after \texttt{RET}, the keys \texttt{C-s} and \texttt{C-r} keep working: the echo area says
\texttt{Repeat with C-s, C-r}, and each press moves to the next or previous match
without asking again. Any other key ends that.

Figure 6.1 shows the screen while a search is open.

\begin{figure}[H]
\caption{Figure 6.1. A search for “ki” in a list of birds, at its first match.}\label{org131}
\begin{Verbatim}[breaklines,breakautoindent=false,breakindent=1em,breakafter=/-.\,;:=,commandchars=\\\{\}]
 +----------------------------------------------------------------+
 | kakapo                                                         |
 | kea                                                            |
 | [ki]wi        <- the current match, in :search-current         |
 | takahe                                                         |
 | [ki]ngfisher  <- other matches in sight, in :search-other      |
 |                                                                |
 +----------------------------------------------------------------+
 | TUATARA[N] PD [-] birds.txt L3 C0                              |
 +----------------------------------------------------------------+
 | 1/2    Regex search forward: ki                                |
 +----------------------------------------------------------------+
    ^                                                  ^
    place and number of matches                        what you typed
\end{Verbatim}
\end{figure}

The \texttt{1/2} before the prompt says the cursor is at the first of two matches.
It is exact however large the buffer is. In a very large buffer the count may
wait until you pause (see “The count before the prompt” below), and typing
never waits for it.

\subsubsection{Capital letters}
\label{sec:132}
A pattern that has no capital letters ignores case: \texttt{kiwi} finds Kiwi, KIWI and
kiwi. As soon as the pattern has a capital, case counts: \texttt{Kiwi} finds only
Kiwi. So you can search loosely most of the time, and type one capital when a
word is too common.

\subsubsection{Regular expressions and plain text}
\label{sec:133}
What you type after \texttt{C-s} is a regular expression, called a regex for short. It
is a pattern in which a few characters are special. A dot stands for any one
character, a star means “any number of the one before”, and so on. That lets
you search for “a digit in brackets” or “a line that starts with a date”.
Table 6.1 lists the forms you will use most. The regex handbook teaches the
whole language.

When the text you want is full of special characters, you do not want to
escape each one. Press \texttt{M-r} in the middle of a search to switch to plain
text. The prompt changes to \texttt{Literal search forward:} and every character
stands for itself, so \texttt{(4)} finds the four characters ( 4 ). Press \texttt{M-r} again
to switch back. \texttt{C-M-s} and \texttt{C-M-r} start a search that is plain text from the
first key.

A regex that is not finished, such as \texttt{kea(}, cannot be searched for. The
prompt says why in brackets, for example \texttt{[incomplete]}, and \texttt{RET} waits until
you mend it. Either fix it or press \texttt{M-r} to take it as plain text.

\begin{xltabular}{\textwidth}{@{}l >{\raggedright\arraybackslash\hsize=1.125\hsize}X >{\raggedright\arraybackslash\hsize=0.875\hsize}X l l@{}}
\caption{Table 6.1. The regex forms you will use most.}\\
\toprule
\textbf{Form} & \textbf{Finds} & \textbf{Example} &  & \\
\midrule
\endfirsthead
\caption[]{Table 6.1. The regex forms you will use most. (continued)}\\
\toprule
\textbf{Form} & \textbf{Finds} & \textbf{Example} &  & \\
\midrule
\endhead
\texttt{.} & Any one character & \texttt{k.a} finds kea, kia, k-a &  & \\
\texttt{\textbackslash{}.} & A real dot (a backslash makes a special character plain) & \texttt{3\textbackslash{}.5} finds 3.5 &  & \\
\texttt{\textbackslash{}d} & One digit & \texttt{\textbackslash{}d\textbackslash{}d} finds 07, 42 &  & \\
\texttt{\textbackslash{}w} & One letter, digit or underscore & \texttt{\textbackslash{}w+} finds a whole word &  & \\
\texttt{\textbackslash{}s} & One space or tab & \texttt{\textbackslash{}s+\$} finds spaces at a line’s end &  & \\
\texttt{[abc]} & One of the characters in brackets & \texttt{gr[ae]y} finds gray, grey &  & \\
\texttt{[\textasciicircum{}abc]} & One character not in the brackets & \texttt{[\textasciicircum{}0-9]} finds a non-digit &  & \\
\texttt{*} & The one before it, any number of times, even none & \texttt{ka*} finds k, ka, kaaa &  & \\
\texttt{+} & The one before it, at least once & \texttt{\textbackslash{}d+} finds 7, 2026 &  & \\
\texttt{?} & The one before it, once or not at all & \texttt{colou?r} finds color, colour &  & \\
\texttt{\textasciicircum{}}, \texttt{\$} & The start, the end of a line & \texttt{\textasciicircum{}Dear} finds lines starting with Dear &  & \\
\texttt{\textbackslash{}b} & The edge of a word & \texttt{\textbackslash{}bkea\textbackslash{}b} does not find kearns &  & \\
=a\textbackslash{} & b= & Either one & =tui\textbackslash{} & kea=\\
\texttt{(...)} & A group; later it can be called \texttt{\textbackslash{}1}, \texttt{\textbackslash{}2} & \texttt{(\textbackslash{}d+)-\codebreak{}(\textbackslash{}d+)} &  & \\
\bottomrule
\end{xltabular}

A pattern that is only an edge, such as \texttt{\textasciicircum{}\$} for an empty line, is a match too.
Searching for it goes from one empty line to the next.

\subsubsection{Starting from the word under the cursor}
\label{sec:134}
Put the cursor on a word and press \texttt{C-s}, then \texttt{M-.}. The word is added to the
pattern as a whole word, so that searching for \texttt{tui} never stops in the middle
of \texttt{tuition}. In plain text mode (\texttt{M-r}) the word goes in as it is.

\subsubsection{Search history}
\label{sec:135}
Every pattern you finish a search with is remembered while Tuatara runs. In
the prompt, \texttt{M-p} or the up arrow brings back the one before, and \texttt{M-n} or the
down arrow a later one. One history serves every question that asks for a
pattern: searching, replacing, and the filters below. Press \texttt{C-s} with nothing
typed and \texttt{RET}, and the last pattern is used again.

\subsubsection{Keeping the matches lit}
\label{sec:136}
After a search ends, its matches stay lit as you move on, which is handy when
you want to see how many there are and where. \texttt{C-g}, or \texttt{keyboard-\codebreak{}quit}, puts
them out. \texttt{C-x T s} lights them again, or takes them away if they are lit.

\subsubsection{The count before the prompt}
\label{sec:137}
Before the question, Tuatara shows a count such as \texttt{3/12}. In a search it says
the match the cursor is at and how many there are. In a question that offers
choices, such as the list of files for \texttt{C-x C-f}, it says which choice is
selected and how many there are. When nothing matches it says \texttt{0/0} and
the search says \texttt{No match}. It always stands in the same six columns, so what
you type does not move as you go from match to match.

Counting takes time in a large buffer, and Tuatara does it in the background so
that typing never lags. A text of up to a megabyte is counted at every key. A
larger one is counted at every key when Tuatara can go through it in about a
hundredth of a second. Otherwise the count waits until you stop typing for
\texttt{*typing-\codebreak{}pause*} seconds, half a second as standard. While a new pattern is
being counted, the old count stays in view. If you do not want the count, set
\texttt{*show-\codebreak{}completion-\codebreak{}count*} to nil.

\subsubsection{Matches that cross many lines}
\label{sec:138}
A regex search is fast because it looks only at the part of the text in view,
plus what it needs. It will find a match of more than 4 KB that runs across
lines only when the text after the match is already in view. For a block such
as \texttt{BEGIN[\textbackslash{}s\textbackslash{}S]*?END} use a thorough command instead. It finds every match of
up to a megabyte, at up to twice the time. The prompt says \texttt{thorough} when it
is in use. The thorough form of replacing is described in the next section.

\begin{try}
\textbf{Try it, with a list of names}
\begin{enumerate}
\item Press \texttt{C-x C-f}, type \texttt{\textasciitilde{}/birds.\codebreak{}txt} and press \texttt{RET}. Press \texttt{i} and type these
   six lines, \texttt{RET} after each: \texttt{Kereru}, \texttt{Tui at Akaroa}, \texttt{Kea (4)},
   \texttt{Fantail at Kaikoura}, \texttt{kiwi}, \texttt{Tui (2)}. Then \texttt{C-g} and \texttt{g g} to go to the
   top.
\item Press \texttt{C-s} and type \texttt{tui}. The cursor goes to the first Tui, because a
   pattern with no capital ignores case. The prompt says \texttt{1/2}. Press \texttt{C-s} to
   go to the second one and \texttt{RET} to stay.
\item Press \texttt{g g} (non-modal: \texttt{M-<}), then \texttt{C-s} and type \texttt{Kiwi} with the capital.
   Nothing matches \texttt{Kiwi} here, because the list has only \texttt{kiwi}, and the prompt
   says \texttt{0/0}. Delete the capital and type \texttt{k} instead: now it matches. Press
   \texttt{C-g}.
\item Press \texttt{C-s} and type \texttt{(2)}. In a regular expression the brackets only make
   a group, so the search stops at the digit 2 alone. Press \texttt{M-r}: the prompt
   now starts \texttt{Literal search forward:} and the match is the three characters
   \texttt{(2)} on the last line. Press \texttt{C-g}.
\item Press \texttt{g e} (non-modal: \texttt{M->}) to go to the end. Press \texttt{C-r} and type \texttt{k}.
   The search goes backward and stops at \texttt{kiwi}. Press \texttt{C-g}, then \texttt{C-s} and
   \texttt{M-p}: the pattern you used before comes back. Press \texttt{C-g}.
\end{enumerate}
\end{try}

{\small
\begin{xltabular}{\textwidth}{@{}l >{\raggedright\arraybackslash\hsize=0.604\hsize}X l >{\raggedright\arraybackslash\hsize=1.396\hsize}X@{}}
\toprule
\textbf{Modal} & \textbf{Non-modal} & \textbf{Command} & \textbf{What it does}\\
\midrule
\endhead
none & \texttt{C-s} & \texttt{regex-\codebreak{}search-\codebreak{}forward} & Search forward as you type for a regular expression (PCRE syntax); a pattern without upper-case letters ignores case, unless it names a case: :upper:, :lower:, \texttt{\textbackslash{}p\{Lu\}}, \texttt{\textbackslash{}p\{Ll\}} and \texttt{\textbackslash{}p\{Lt\}} heed it, as a capital letter does.\\
none & \texttt{C-r} & \texttt{regex-\codebreak{}search-\codebreak{}backward} & Search backward as you type for a regular expression; keys, replacing and the limit on matches of more than 4 KB as in regex-search-forward.\\
none & \texttt{C-M-s} & \texttt{literal-\codebreak{}search-\codebreak{}forward} & Search forward as you type for literal text: keys as in regex-search-forward, nothing typed special.\\
none & \texttt{C-M-r} & \texttt{literal-\codebreak{}search-\codebreak{}backward} & Search backward as you type for literal text; see literal-search-forward.\\
none & \texttt{M-x search-\codebreak{}repeat-\codebreak{}forward} & \texttt{search-\codebreak{}repeat-\codebreak{}forward} & During a search: the next match forward (with an empty string, the last search string).\\
none & \texttt{M-x search-\codebreak{}repeat-\codebreak{}backward} & \texttt{search-\codebreak{}repeat-\codebreak{}backward} & During a search: the next match backward.\\
none & \texttt{M-x search-\codebreak{}toggle-\codebreak{}regex} & \texttt{search-\codebreak{}toggle-\codebreak{}regex} & During a search, or a search of files (M-s g): switch between a regular expression and literal text.\\
none & \texttt{M-x search-\codebreak{}insert-\codebreak{}symbol} & \texttt{search-\codebreak{}insert-\codebreak{}symbol} & During a search: add the symbol at point where the search began to the search string, in a regular expression as a whole symbol, never found inside a longer one.\\
none & \texttt{M-x search-\codebreak{}exit} & \texttt{search-\codebreak{}exit} & End the search, leaving point at the match and the mark where the search began; when the string holds a replacement, carry it out; for a search filter, keep it on the list.\\
none & \texttt{M-x search-\codebreak{}abort} & \texttt{search-\codebreak{}abort} & Cancel the search and return to where it began; a search filter being typed goes, the list shown as the filters kept leave it.\\
none & \texttt{M-x search-\codebreak{}next-\codebreak{}match} & \texttt{search-\codebreak{}next-\codebreak{}match} & Go to the next match of the search just ended with RET, without its prompt: C-s right after the search does, again and again, the last match followed by the first.\\
none & \texttt{M-x search-\codebreak{}previous-\codebreak{}match} & \texttt{search-\codebreak{}previous-\codebreak{}match} & Go to the previous match of the search just ended with RET, without its prompt: C-r right after the search does, again and again, the first match followed by the last.\\
none & \texttt{C-x T s} & \texttt{search-\codebreak{}toggle-\codebreak{}highlight} & Light the matches of the last search that ended at a match again, in its buffer, or take them away when they are lit, as M-x keyboard-quit does.\\
none & \texttt{M-x regex-\codebreak{}search-\codebreak{}forward-\codebreak{}thorough} & \texttt{regex-\codebreak{}search-\codebreak{}forward-\codebreak{}thorough} & Search forward as regex-search-forward does, finding every match of up to a megabyte (about a million characters of ASCII) whole however many lines it spans, where regex-search-forward may miss one of more than 4 KB that runs across lines and needs text further on (a block of lines matched by \texttt{BEGIN[\textbackslash{}s\textbackslash{}S]*?END}, say).\\
none & \texttt{M-x regex-\codebreak{}search-\codebreak{}backward-\codebreak{}thorough} & \texttt{regex-\codebreak{}search-\codebreak{}backward-\codebreak{}thorough} & Search backward as regex-search-backward does, finding every match of up to a megabyte whole however many lines it spans, with the cost and the limit of regex-search-forward-thorough: a match of more than a megabyte is still cut short, or missed when it needs text further on.\\
\bottomrule
\end{xltabular}
}

\begin{xltabular}{\textwidth}{@{}l l >{\raggedright\arraybackslash\hsize=1.000\hsize}X@{}}
\toprule
\textbf{Setting} & \textbf{Standard} & \textbf{What it changes}\\
\midrule
\endhead
\texttt{*show-\codebreak{}completion-\codebreak{}count*} & \texttt{t} & True to show before the prompt of a completing read the place of the selected candidate and how many there are, as they change with what is typed: 3/120 (*/0 when nothing matches and the text typed can be taken, !/0 when it cannot); before a search’s, the match the cursor is at and how many there are.\\
\texttt{*typing-\codebreak{}pause*} & \texttt{1/2} & Seconds without a key typed before work over a text too large to go through at each key starts while a pattern is typed: the count of a search’s matches (M-x regex-search-forward) and a filter being typed (M-x search-filter-regex), its lines in the window shown at once.\\
\bottomrule
\end{xltabular}

\begin{xltabular}{\textwidth}{@{}l >{\raggedright\arraybackslash\hsize=1.000\hsize}X@{}}
\toprule
\textbf{Face} & \textbf{What it paints}\\
\midrule
\endhead
\texttt{:match} & A match found by search or replace.\\
\texttt{:search-\codebreak{}current} & The current match of an incremental search.\\
\texttt{:search-\codebreak{}other} & The other matches of the current search in view.\\
\texttt{:search-\codebreak{}outside-\codebreak{}scope} & Text outside the scope of a search or replacement limited to a line or the region, faded so that the scope stands out.\\
\bottomrule
\end{xltabular}

When a region is active, a search stays inside it, and the text outside is
faded (\texttt{:search-\codebreak{}outside-\codebreak{}scope}) so that you can see the limit. A pattern that
begins with \texttt{\%} searches the whole buffer instead.

\subsection{Replacing}
\label{sec:142}
Replacing uses the same prompt as searching. In NORMAL state press \texttt{r r}
(non-modal: \texttt{C-z r r}). The prompt reads \texttt{Regex replace on this line:}. Type the
pattern, then \texttt{RET}, which types the \texttt{/} for you, then the new text, then \texttt{RET}.
You can also type it all on one line as \texttt{pattern/\codebreak{}new text}.

What gets changed depends on what is selected. With nothing selected, only
the line the cursor is on, and only the first match on it. With a selection,
the selected lines. A \texttt{\%} in front of the pattern, as in \texttt{\%tui/\codebreak{}Tui}, means every
match in the whole buffer. A replacement is a single change, so one \texttt{u}
(non-modal: \texttt{C-/}) takes the whole thing back, however many places it changed.

While you type, Tuatara shows what will happen. Each match that is going away
is lit as \texttt{:replace-\codebreak{}going}, and the new text stands beside it in
\texttt{:replace-\codebreak{}coming}, before you press \texttt{RET}. Look before you press it.
Afterwards the echo area says how much was done, for instance \texttt{Replaced 3
occurrences on 3 lines}.

Typed at the C-s prompt, a pattern with a \texttt{/} in it also replaces, the way it
does in the Vim editor. A word after the last slash gives flags: \texttt{g} replaces
every match on a line, \texttt{c} asks at each match, \texttt{n} only counts what would be
replaced, and \texttt{i} or \texttt{I} ignores or keeps case.

\subsubsection{What goes in the new text}
\label{sec:140}
\begin{longtable}{@{}l l@{}}
\toprule
\textbf{Write} & \textbf{To get}\\
\midrule
\endhead
\texttt{\&}, \texttt{\textbackslash{}0} & The whole text that matched\\
\texttt{\textbackslash{}1} to \texttt{\textbackslash{}9} & What the matching bracket group found\\
\texttt{\textbackslash{}n}, \texttt{\textbackslash{}t} & A line break, a tab\\
\texttt{\textbackslash{}u}, \texttt{\textbackslash{}l} & The next letter in capitals, in small letters\\
\texttt{\textbackslash{}U}, \texttt{\textbackslash{}L} & Everything up to \texttt{\textbackslash{}E} in capitals, in small letters\\
\texttt{\textbackslash{}/} & A real slash (also in the pattern)\\
\bottomrule
\end{longtable}

If the pattern has no capital letters, the replacement follows the case of
each match: a match in capitals is replaced in capitals, and one that starts
with a capital gets a capital. So \texttt{\%tui/\codebreak{}kokako} turns Tui into Kokako.

\subsubsection{Replacing with questions}
\label{sec:141}
\texttt{r q} (non-modal: \texttt{C-z r q}) replaces a regular expression and asks about each
match, with \texttt{Query replacing ... with ...:} and the keys \texttt{y n ! . q}:

\begin{longtable}{@{}l l l@{}}
\toprule
\textbf{Modal} & \textbf{Non-modal} & \textbf{Does}\\
\midrule
\endhead
none & \texttt{y} & Replace this one and go to the next\\
none & \texttt{n} & Leave this one and go to the next\\
none & \texttt{!} & Replace this one and all the rest\\
none & \texttt{.} & Replace this one and stop\\
none & \texttt{q} & Stop\\
\bottomrule
\end{longtable}

\texttt{M-\%} is the same for plain text, with no special characters to think about.
\texttt{C-z r q} is the key for \texttt{r q} in every state.

When you want to be sure of each change, use the question form. When a
change is plain and mechanical, use \texttt{r r} and look at the preview.

If the text you replace can run over many lines, use the thorough forms of
the commands. They are described under searching, above.

\begin{try}
\textbf{Try it, with a letter}
\begin{enumerate}
\item Open \texttt{\textasciitilde{}/letter.\codebreak{}txt} and type this text: \texttt{Dear Hana,}, a line break,
   \texttt{We met the tui at Akaroa. The tui sings at dawn. Tui are rare.}, a line
   break, \texttt{Yours, Aroha}. Press \texttt{C-g} and \texttt{g g}.
\item Press \texttt{j} (non-modal: \texttt{C-n}) to move to the second line, then \texttt{r r} and type
   \texttt{tui/\codebreak{}kea}, then \texttt{RET}. Only the first match on this line changes. Press \texttt{u}
   (non-modal: \texttt{C-/}) to take it back.
\item Press \texttt{r r} and type \texttt{\%tui/\codebreak{}kea}, \texttt{RET}. Every tui of the whole letter
   changes, and the third one keeps its capital: \texttt{Kea are rare}. The echo area
   says \texttt{Replaced 3 occurrences on 1 line}. Press \texttt{u}.
\item Press \texttt{r q} (non-modal: \texttt{C-z r q}), type \texttt{\%tui/\codebreak{}kea}, \texttt{RET}. Press \texttt{y} at the
   first, \texttt{n} at the second and \texttt{!} at the third. Only the first and the last
   were changed. Press \texttt{u}.
\item Press \texttt{r r} and type \texttt{\%(\textbackslash{}w+) at (\textbackslash{}w+)/\textbackslash{}2 at \textbackslash{}1}, \texttt{RET}. Look at what
   \texttt{\textbackslash{}1} and \texttt{\textbackslash{}2} did: the two words either side of “at” swap places, as in
   \texttt{Akaroa at tui}. Press \texttt{u}.
\end{enumerate}
\end{try}

\begin{xltabular}{\textwidth}{@{}l >{\raggedright\arraybackslash\hsize=0.684\hsize}X l >{\raggedright\arraybackslash\hsize=1.316\hsize}X@{}}
\toprule
\textbf{Modal} & \textbf{Non-modal} & \textbf{Command} & \textbf{What it does}\\
\midrule
\endhead
\texttt{r r} & \texttt{C-z r r} & \texttt{replace-\codebreak{}regexp} & Replace a regular expression, as one edit: pattern/replacement, typed on one line or in two parts, as RET after the pattern types the / before the replacement.\\
none & \texttt{M-\%} & \texttt{query-\codebreak{}replace} & Replace literal text, asking at each match: y replaces, n skips, ! replaces all the rest, .\\
\texttt{r q} & \texttt{C-z r q} & \texttt{query-\codebreak{}replace-\codebreak{}regexp} & Replace a regular expression as C-z r r does, asking at each match: y replaces, n leaves it and goes to the next, ! replaces all the rest, .\\
none & \texttt{M-x replace-\codebreak{}regexp-\codebreak{}thorough} & \texttt{replace-\codebreak{}regexp-\codebreak{}thorough} & Replace a regular expression as replace-regexp does (the first match on each line, every match after a leading \%), finding every match of up to a megabyte whole however many lines it spans, with the cost and the limit of regex-search-forward-thorough: a match of more than a megabyte is still cut short, or missed when it needs text further on.\\
none & \texttt{M-x query-\codebreak{}replace-\codebreak{}regexp-\codebreak{}thorough} & \texttt{query-\codebreak{}replace-\codebreak{}regexp-\codebreak{}thorough} & Replace a regular expression asking at each match, as query-replace-regexp does, finding every match of up to a megabyte whole however many lines it spans, with the cost and the limit of regex-search-forward-thorough: a match of more than a megabyte is still cut short, or missed when it needs text further on.\\
\bottomrule
\end{xltabular}

\begin{xltabular}{\textwidth}{@{}l >{\raggedright\arraybackslash\hsize=1.000\hsize}X@{}}
\toprule
\textbf{Face} & \textbf{What it paints}\\
\midrule
\endhead
\texttt{:replace-\codebreak{}coming} & What a replacement being typed (r r, r q) puts in place of a match, shown beside the match before RET, as no text of the buffer: as a difference shows text added, on a ground of its own.\\
\texttt{:replace-\codebreak{}going} & A match a replacement being typed (r r, r q) takes away, shown before RET beside what takes its place: as a difference shows text taken away, on a ground of its own so that it reads the same over any mode’s colours.\\
\bottomrule
\end{xltabular}

\subsection{Narrowing a list to its matches}
\label{sec:143}
A list is any buffer made of lines you choose among: the buffer list from
\texttt{C-x C-b}, the list of jobs, the packages, the keys. Each such list can be cut
down to the lines that match a pattern, and each file of records, such as a
log or a table of values, can be cut down in the same way. This is called a
search filter. The filter hides lines. Nothing is deleted, and the file on
disk does not change.

Press \texttt{C-x / s} in a list and type a pattern. The list shrinks as you type, and
the count before the prompt says how many lines are left, for example
\texttt{1/5    Regex filter: dr}. \texttt{RET} keeps the filter. \texttt{C-g} drops the one you were
typing and keeps the earlier ones. Press \texttt{C-x / s} again and a second filter
narrows what the first left. Above the list a line (\texttt{:list-\codebreak{}filters}) says by
what it is narrowed and how many entries are left.

\texttt{C-x / l} lists your filters in a buffer of their own, each with the number of
entries it leaves. There, \texttt{d} takes the filter on the line away, and \texttt{D}
(non-modal: \texttt{C-k}) takes them all away. \texttt{C-x / r} takes them all away too, from
the narrowed list itself, and the whole list is shown again.

The same filter works on files whose lines are records, such as a log or a
CSV. The modes that count are named in \texttt{*line-\codebreak{}record-\codebreak{}modes*}. In these, the
first line of a CSV is a header and is always shown. \texttt{C-x / f} shows only the
records whose one column matches a pattern, and \texttt{C-x / h} hides them. These ask
which column first. The lists and the data chapters describe columns in
full.

Two keys delete instead of hide. \texttt{C-x / k} keeps the records that match and
deletes the others. \texttt{C-x / d} deletes the ones that match. Both can be undone
with \texttt{u} (non-modal: \texttt{C-/}). \texttt{C-x / t} makes the first line of a file of records
its header, or a record like the rest, and \texttt{C-x / L} keeps the header at the top
of the window as the rest scrolls.

You can still type in the lines you see. A change that would reach a hidden
line is refused with a message, so a filter never changes what you cannot
see. Press \texttt{C-x / r} to see everything again.

\begin{try}
\textbf{Try it, with the buffer list}
\begin{enumerate}
\item Press \texttt{C-x C-b} to list your buffers. If you did the earlier chapters it
   holds \texttt{shopping.\codebreak{}txt}, \texttt{pancakes.\codebreak{}txt}, \texttt{birds.\codebreak{}txt} and \texttt{letter.\codebreak{}txt}.
\item Press \texttt{C-x / s} and type \texttt{txt}. Lines that do not hold those letters go
   away, such as the starred buffers. The count in front of the prompt
   follows what you type. Press \texttt{RET}.
\item Press \texttt{C-x / s} again, type \texttt{pan} and press \texttt{RET}. This narrows the list
   to what the first filter left, so only \texttt{pancakes.\codebreak{}txt} remains. Look at
   the line above the list.
\item Press \texttt{C-x / l}. The two filters are listed with how many entries each
   leaves. Move to the first line, press \texttt{d} to take it away, then \texttt{q}.
\item Press \texttt{C-x / r}. Every buffer is back.
\end{enumerate}
\end{try}

\begin{xltabular}{\textwidth}{@{}l l l >{\raggedright\arraybackslash\hsize=1.000\hsize}X@{}}
\toprule
\textbf{Modal} & \textbf{Non-modal} & \textbf{Command} & \textbf{What it does}\\
\midrule
\endhead
none & \texttt{C-x / s} & \texttt{search-\codebreak{}filter-\codebreak{}regex} & Narrow the list here to the entries a regular expression matches, as it is typed: those whose line matches, or what more the list holds of them (a mailbox the headers and contents of its messages).\\
none & \texttt{C-x / l} & \texttt{search-\codebreak{}filter-\codebreak{}list} & List the search filters narrowing the list here (C-x / s) in a buffer of their own, each with how many entries it leaves; there d takes away the filter on the line and D all of them.\\
none & \texttt{d}, \texttt{C-c d} & \texttt{search-\codebreak{}filter-\codebreak{}remove} & Take away the search filter on the current line of the list of filters (C-x / l): the list it narrowed widens, the filters after it narrowing again what those before it leave, in the background where that takes time.\\
none & \texttt{D}, \texttt{C-c D}, \texttt{C-x / r} & \texttt{search-\codebreak{}filter-\codebreak{}remove-\codebreak{}all} & Take away every search filter narrowing the list here, or the list the list of filters here is for: the whole list is shown again.\\
none & \texttt{C-x / f} & \texttt{search-\codebreak{}filter-\codebreak{}show-\codebreak{}records} & Show only the records whose line a regular expression matches, hiding the others, as C-x / s does, deleting nothing.\\
none & \texttt{C-x / h} & \texttt{search-\codebreak{}filter-\codebreak{}hide-\codebreak{}records} & Hide the records whose line a regular expression matches, showing the others, deleting nothing: asked for as C-x / f asks, kept with the other filters.\\
none & \texttt{C-x / k} & \texttt{list-\codebreak{}keep-\codebreak{}records} & Keep the records whose line a regular expression matches and delete the rest, in a file of records, a line each: a CSV or TSV, a log, a JSON Lines file, todo.txt and the others \textbf{LINE-RECORD-MODES} names, the column asked for first as C-x / f asks.\\
none & \texttt{C-x / d} & \texttt{list-\codebreak{}delete-\codebreak{}records} & Delete the records whose line a regular expression matches, keeping the rest: asked for as C-x / k asks, in the same files.\\
none & \texttt{C-x / t} & \texttt{list-\codebreak{}toggle-\codebreak{}header} & Take the first line of a file of records as its header, or as a record like the rest; the same again switches back; with a count, a header when it is positive, a record when it is zero or less.\\
none & \texttt{C-x / L} & \texttt{list-\codebreak{}toggle-\codebreak{}header-\codebreak{}lock} & Keep the first line of a file of records, its header, at the top of the window as the rest scrolls under it, coloured as it is below; the same again lets it scroll away; with a count, kept when it is positive, let go when it is zero or less.\\
\bottomrule
\end{xltabular}

\begin{xltabular}{\textwidth}{@{}l >{\raggedright\arraybackslash\hsize=0.582\hsize}X >{\raggedright\arraybackslash\hsize=1.418\hsize}X@{}}
\toprule
\textbf{Setting} & \textbf{Standard} & \textbf{What it changes}\\
\midrule
\endhead
\texttt{*line-\codebreak{}record-\codebreak{}modes*} & \texttt{(tuatara.\codebreak{}defaults::\codebreak{}log-\codebreak{}mode tuatara.\codebreak{}defa...} & The major modes whose files hold a record a line, as a log does: search filters (M-x search-filter-regex) narrow them, and the modes made from them, by hiding the lines they leave out, as they narrow a CSV.\\
\texttt{*csv-\codebreak{}header*} & \texttt{t} & True when the first line of the files CSV mode shows is their header, naming the columns: always shown by search filters, never deleted by M-x list-keep-records, its names offered for a column.\\
\bottomrule
\end{xltabular}

\begin{xltabular}{\textwidth}{@{}l >{\raggedright\arraybackslash\hsize=1.000\hsize}X@{}}
\toprule
\textbf{Face} & \textbf{What it paints}\\
\midrule
\endhead
\texttt{:list-\codebreak{}filters} & The line above a list narrowed by search filters (C-x / s), saying by what and how many entries are left.\\
\bottomrule
\end{xltabular}

\subsection{Searching many files}
\label{sec:146}
\texttt{M-s g} finds text in many files at once. It looks in the project that holds
your file. A project is the folder above your file that holds a \texttt{.git} or
\texttt{.project} entry (see \texttt{*project-\codebreak{}root-\codebreak{}markers*}). If there is no such folder,
it looks in the folder of the file. In a directory listing it searches that
directory, and with a prefix argument (\texttt{C-u} first) it asks for a directory.

Type words and Tuatara lists every line that holds all of them, in any order,
each as \texttt{path:\codebreak{}line: text}. Nothing is listed before the third character
(\texttt{*search-\codebreak{}files-\codebreak{}min-\codebreak{}input*}). The files are read in the background while you
type, and Tuatara skips hidden files, backups, binary files and whatever the
project’s \texttt{.gitignore} leaves out. Each word is a regex, as for \texttt{C-s}, and \texttt{M-r}
switches to plain text, where the whole line you typed is found as it stands,
spaces and all. At most \texttt{*search-\codebreak{}files-\codebreak{}lines-\codebreak{}limit*} lines are listed. If the
search stops there, the first line of the list says so.

Move down the list with \texttt{C-n} (modal: \texttt{j}) and \texttt{RET} opens the file at that
line. With \texttt{*preview-\codebreak{}file-\codebreak{}search*} set to true, the file of the selected line is
shown meanwhile in your window, so you can read around a match without leaving
the search. After the search, \texttt{M-g n} and \texttt{M-g p} go to the next and previous
place in the list, the way they go through the errors of a compile.

A search of your files with formatted names, the ones that carry a date and
keywords in their name, is \texttt{C-x n g} (\texttt{search-\codebreak{}in-\codebreak{}formatted}), and works the
same way, with \texttt{*preview-\codebreak{}formatted-\codebreak{}search*} as its preview setting.

\subsubsection{Changing the lines a search found}
\label{sec:144}
While the question is open, \texttt{C-c C-l} turns the lines found so far into a
buffer of their own. \texttt{C-x C-q} there makes the text after \texttt{path:\codebreak{}line:}
changeable. You can edit the lines as in any text, but you cannot change a
path or a line number or split a line. \texttt{C-c C-c} writes the changed lines back
into their files. It runs in the background, and a line is written only when
the file still has the text that was found. A file that is open in a buffer is
changed there, as a step you can undo. Other files are saved and closed again,
unless \texttt{*search-\codebreak{}results-\codebreak{}save-\codebreak{}files*} is nil, in which case each stays changed
and unsaved for you to look at. \texttt{C-c C-k} throws the changes away and leaves
the files as they were.

Files that you opened safely, files you may not write and read-only buffers are
never written. The message log says why.

\subsubsection{Keeping folders read}
\label{sec:145}
A big folder takes a while to read. Name the folders you search often in
\texttt{*indexed-\codebreak{}directories*} in your init file, for example
\texttt{(setf *indexed-\codebreak{}directories* (list "\textasciitilde{}/\codebreak{}notes/"))}. Tuatara then keeps their file
names and text in memory. It reads them in the background when it starts, and
again when something changes on disk. After that \texttt{M-s g} in them answers at
once, and so does the search of your files with formatted names. While a folder
is being read the mode line shows \texttt{U} after the \texttt{P} sign, and while a kept
folder is searched it shows \texttt{S}. They stand always in the same place and are
faint when idle (\texttt{*indexing-\codebreak{}indicator*}, \texttt{*searching-\codebreak{}indicator*}).

By default what was read is also kept on disk between sessions
(\texttt{*indexes-\codebreak{}kept-\codebreak{}on-\codebreak{}disk*}), so that at the next start only the files that
changed are read again. Set it to nil for a folder whose text must not be
copied anywhere. \texttt{*search-\codebreak{}files-\codebreak{}limit*} is the most text Tuatara keeps for
one search or one kept folder, 128 MB as standard. Files past it are searched
from the disk, more slowly.

\begin{try}
\textbf{Try it, with a folder of notes}
\begin{enumerate}
\item Make a folder called \texttt{notes} in your home folder, with your file manager
   or with \texttt{mkdir \textasciitilde{}/notes} in a terminal. Open \texttt{\textasciitilde{}/notes/\codebreak{}akaroa.\codebreak{}txt} with
   \texttt{C-x C-f} and type \texttt{Tui and kereru near the harbour}. Save it with
   \texttt{C-x C-s}. Make a second file, \texttt{\textasciitilde{}/notes/\codebreak{}kaikoura.\codebreak{}txt}, with \texttt{Fantail and
   tui by the stream} and save it.
\item With \texttt{kaikoura.\codebreak{}txt} showing, press \texttt{M-s g} and type \texttt{tui}. Nothing is listed
   before the third letter. Then both lines are, each as \texttt{path:\codebreak{}line: text}.
\item Press \texttt{RET}. \texttt{akaroa.\codebreak{}txt} opens at its line. Press \texttt{M-g n} to go on to the
   next match, in \texttt{kaikoura.\codebreak{}txt}. (\texttt{C-n} before \texttt{RET} would take you to the
   second line of the list instead.)
\item Press \texttt{M-s g} and type \texttt{tui harbour}. Only the \texttt{akaroa.\codebreak{}txt} line is listed,
   because it holds both words. Press \texttt{C-c C-l} to put the lines in a buffer,
   \texttt{C-x C-q} to change them, replace \texttt{harbour} with \texttt{bay} on the line, then
   \texttt{C-c C-c} to write it back. Look at \texttt{akaroa.\codebreak{}txt}: it says \texttt{bay}.
\item To keep the folder read, add \texttt{(setf *indexed-\codebreak{}directories* (list
   "\textasciitilde{}/\codebreak{}notes/"))} to your init file and restart. The \texttt{U} in the mode
   line lights up while it is read, and is faint again when it is done.
\end{enumerate}
\end{try}

\begin{xltabular}{\textwidth}{@{}l >{\raggedright\arraybackslash\hsize=0.522\hsize}X l >{\raggedright\arraybackslash\hsize=1.478\hsize}X@{}}
\toprule
\textbf{Modal} & \textbf{Non-modal} & \textbf{Command} & \textbf{What it does}\\
\midrule
\endhead
none & \texttt{M-s g} & \texttt{search-\codebreak{}in-\codebreak{}files} & Find text in the files of the project the current buffer’s file is in (see \textbf{PROJECT-ROOT-MARKERS}), or under its directory when it is in none, or under the directory it lists (C-x d), or \textbf{DEFAULT-DIRECTORY} for another buffer with no file; with a prefix argument, in a directory asked for.\\
none & \texttt{C-x n g} & \texttt{search-\codebreak{}in-\codebreak{}formatted} & Find text in \textbf{FORMATTED-FILE-DIRECTORY}, where your files with formatted names are: every file under it, as search-in-files finds text in a directory; the lines with the words typed are listed as you type, the files read in the background meanwhile, M-r switching to literal text and back, and M-x minibuffer-finish visits one.\\
none & \texttt{M-x search-\codebreak{}files-\codebreak{}list} & \texttt{search-\codebreak{}files-\codebreak{}list} & While searching files (M-s g), list the lines found so far in a buffer of their own, one each as path:line: text, to go through with RET or to change at once, with C-x C-q, and write back into their files.\\
none & \texttt{C-x C-q} & \texttt{search-\codebreak{}results-\codebreak{}edit} & Make the text after the path and line number of the lines found changeable, with the keys that change text: the path and line number stay as they are, and a line is not split or joined.\\
none & \texttt{C-c C-c} & \texttt{search-\codebreak{}results-\codebreak{}finish} & Write the lines changed here into their files.\\
none & \texttt{C-c C-k} & \texttt{search-\codebreak{}results-\codebreak{}cancel} & Throw away the changes made to the lines found, and make them read-only again; the files are not touched.\\
none & \texttt{M-g n} & \texttt{next-\codebreak{}place} & Go to the next place of the list of places used last: the next line naming a file and line in a compile’s output (M-x compile), the next line a search of files found (M-s g), or the next place the language servers found something wrong in a buffer.\\
none & \texttt{M-g p} & \texttt{previous-\codebreak{}place} & Go to the place before the one gone to last in the list of places used last: a compile’s output (M-x compile), the lines a search of files found (M-s g), or the places the language servers found something wrong in a buffer, as M-g n goes to the next.\\
\bottomrule
\end{xltabular}

\begin{xltabular}{\textwidth}{@{}l l >{\raggedright\arraybackslash\hsize=1.000\hsize}X@{}}
\toprule
\textbf{Setting} & \textbf{Standard} & \textbf{What it changes}\\
\midrule
\endhead
\texttt{*project-\codebreak{}root-\codebreak{}markers*} & \texttt{(".\codebreak{}git" ".project")} & The names of the files or directories that mark the root of a project: M-x search-in-files searches all of the project a buffer’s file is in, from the nearest directory above the file holding one of them.\\
\texttt{*default-\codebreak{}directory*} & \texttt{"\textasciitilde{}/"} & The directory of a buffer with no file (\textbf{scratch}, the lists): a file name typed there is taken from it, and M-x search-in-files searches it.\\
\texttt{*search-\codebreak{}files-\codebreak{}min-\codebreak{}input*} & \texttt{3} & How many characters search-in-files and search-in-formatted wait for before searching.\\
\texttt{*search-\codebreak{}files-\codebreak{}lines-\codebreak{}limit*} & \texttt{2000} & The most matching lines search-in-files and search-in-formatted list, 2000 as standard.\\
\texttt{*search-\codebreak{}files-\codebreak{}limit*} & \texttt{134217728} & The most octets of text kept in memory for a search of many files (M-x search-in-files, M-x search-in-formatted) and for each directory of \textbf{INDEXED-DIRECTORIES} (and on disk with it, \textbf{INDEXES-KEPT-ON-DISK}): 128 MB as standard.\\
\texttt{*search-\codebreak{}results-\codebreak{}save-\codebreak{}files*} & \texttt{t} & True to save the files that lines found by M-x search-in-files are written back into (M-x search-results-finish), those that were not open too, closed again after; NIL to leave each changed and not saved in a buffer of its own, for you to look at and save (M-x save-buffer).\\
\texttt{*preview-\codebreak{}file-\codebreak{}search*} & \texttt{nil} & True to show, while you search the files of a project or a directory with M-x search-in-files, the file of the line selected in the window you are in, at that line, opened safely, as \textbf{PREVIEW-FILES} shows files.\\
\texttt{*preview-\codebreak{}formatted-\codebreak{}search*} & \texttt{nil} & True to show, while you search your files with formatted names with M-x search-in-formatted, the file of the line selected in the window you are in, at that line, opened safely, as \textbf{PREVIEW-FILES} shows files.\\
\texttt{*indexed-\codebreak{}directories*} & \texttt{nil} & Directories Tuatara keeps read in memory, so that M-x open-file-formatted, M-x search-in-formatted and M-x search-in-files answer at once there, and in the directories under them: the names of their files and the text of those a search reads (at most \textbf{SEARCH-FILES-LIMIT} octets each), read in the background when the editor starts and again whenever something in them changes on disk.\\
\texttt{*indexes-\codebreak{}kept-\codebreak{}on-\codebreak{}disk*} & \texttt{t} & True to keep what is read of \textbf{INDEXED-DIRECTORIES} on disk between sessions, in the cache directory (\$XDG\_CACHE\_HOME/tuatara/indexes/, readable by you alone): at the next start the directories are only listed again, and only the files changed meanwhile read, rather than all.\\
\texttt{*indexing-\codebreak{}indicator*} & \texttt{"U"} & What the mode line shows while directories are kept read (\textbf{INDEXED-DIRECTORIES}), always in the same place, before [-], after the sign of work for the buffer (\textbf{PROCESSING-INDICATOR}): in :mode-line-sign while one is being read in the background (at least \textbf{SIGN-LIT-SECONDS}), else faint, in :mode-line-sign-idle, so that nothing after it moves as reading starts and ends.\\
\texttt{*searching-\codebreak{}indicator*} & \texttt{"S"} & What the mode line shows after \textbf{INDEXING-INDICATOR} while directories are kept read, always in the same place: in :mode-line-sign while a search runs over a directory kept read (M-x search-in-files and the like; at least \textbf{SIGN-LIT-SECONDS}), else faint, in :mode-line-sign-idle.\\
\bottomrule
\end{xltabular}

\subsection{The outline of a buffer}
\label{sec:147}
\texttt{M-g o} lists the headings of the buffer you are in and goes to the one you
choose. In an Org file they are the headings. In Lisp they are the comments
that start with three semicolons and the top-level forms. In a language that
Tuatara has a grammar for, they are its definitions, such as functions and
types, and in Markdown its headings. Each is shown as its line, after the
line’s number, in the order of the text.

Type a few letters of a name and the list narrows to the headings that hold
them. The window behind follows the first one (\texttt{:outline-\codebreak{}preview}), and \texttt{RET}
goes there. \texttt{C-g} takes you back to where you were. Start with a digit and a
space to see only the outer headings: =2 = then the words shows the headings
two levels deep or less.

When a language server is running for the file and \texttt{*lsp-\codebreak{}outline*} is true,
the list uses the names the server knows, including fields and methods inside
types. Otherwise it uses the headings Tuatara finds itself.

\begin{try}
\textbf{Try it, with a long Org file}
\begin{enumerate}
\item Open any Org file with several headings, for instance one you made in the
   earlier chapters, or press \texttt{C-x C-f} and make \texttt{\textasciitilde{}/trip.\codebreak{}org} with
   \texttt{* Akaroa}, \texttt{** Harbour}, \texttt{* Kaikoura}, \texttt{** Whale watching} and
   \texttt{** Crayfish} on separate lines.
\item Press \texttt{M-g o}. Every heading is listed with its line number.
\item Type \texttt{whale}. Only the matching heading is left, and the window behind
   shows it. Press \texttt{RET} to go there.
\item Press \texttt{M-g o} again, type =1 = and a space, and see only the top headings.
   Press \texttt{C-g} to go back.
\end{enumerate}
\end{try}

\begin{xltabular}{\textwidth}{@{}l l l >{\raggedright\arraybackslash\hsize=1.000\hsize}X@{}}
\toprule
\textbf{Modal} & \textbf{Non-modal} & \textbf{Command} & \textbf{What it does}\\
\midrule
\endhead
none & \texttt{M-g o} & \texttt{buffer-\codebreak{}outline} & Go to a heading of the current buffer, chosen from a list of them: Org’s headings, Lisp’s comments of three semicolons and top-level forms, and in a language parsed with tree-sitter its definitions and headings (a function, a type, a Markdown heading), each shown as its line after its number, in the order of the text.\\
\bottomrule
\end{xltabular}

\begin{xltabular}{\textwidth}{@{}l l >{\raggedright\arraybackslash\hsize=1.000\hsize}X@{}}
\toprule
\textbf{Setting} & \textbf{Standard} & \textbf{What it changes}\\
\midrule
\endhead
\texttt{*lsp-\codebreak{}outline*} & \texttt{t} & True for M-x buffer-outline to list a file’s names as its language server knows them (functions, types, their fields and methods, nested), where it has one that does; else as the mode finds its headings.\\
\bottomrule
\end{xltabular}

\begin{longtable}{@{}l l@{}}
\toprule
\textbf{Face} & \textbf{What it paints}\\
\midrule
\endhead
\texttt{:outline-\codebreak{}preview} & The line of the heading shown while one is chosen in M-g o.\\
\bottomrule
\end{longtable}

\section{Buffers and sessions}
\label{sec:165}
Keeping track of what is open, and getting it back after you leave.

Chapter 1 explained that a buffer is a piece of text Tuatara holds in memory,
and that a window is a view of one. After an afternoon of work you may have a
dozen buffers open: a letter, a shopping list, some notes, a folder listing,
the help. This chapter is about managing them. It shows the list of all
buffers, how to get rid of the ones you are done with, how to find a file you
had open yesterday, and how to leave a place marked so that you can come back
to it. It also explains what Tuatara does so that a crash or a closed
terminal does not cost you your work: it can remember what was open, and it
keeps a copy of changes you have not saved. The last part shows how to look
at just one part of a buffer.

Before you start, one reminder about the state keys. \texttt{i} starts INSERT state,
and \texttt{C-g} or \texttt{ESC} returns to NORMAL state. Every key in this chapter is
pressed in NORMAL state unless it starts with \texttt{C-x}, which works in both.

\subsection{The buffer list}
\label{sec:156}
\texttt{C-x C-b} shows every buffer in a buffer of its own called \texttt{*Buffers*}. The
one you used last is at the top. You use the list when you have lost track
of what is open, when you want to close several buffers at once, or when you
want to save several files in one go.

Figure 7.1 shows a list for someone who has a folder of letters under version
control, a few loose notes, and the log of what the echo area said.

\begin{figure}[H]
\caption{Figure 7.1. The buffer list, a line for each buffer.}\label{org149}
\begin{Verbatim}[breaklines,breakautoindent=false,breakindent=1em,breakafter=/-.\,;:=,commandchars=\\\{\}]
 {\tuataraFallbackTWO{}▾} letters (3)
   >*  council.txt        412  Text    ~/letters/council.txt
       thanks.txt         190  Text    ~/letters/thanks.txt
   D   old-draft.txt        0  Text    ~/letters/old-draft.txt
 {\tuataraFallbackTWO{}▾} No project (3)
       shopping.txt        34  Text    ~/shopping.txt
       notes.org         2,061 Org     ~/notes.org
       *Messages*        1,210 Text
\end{Verbatim}
\end{figure}

Each buffer has one line. The line gives the buffer’s marks, its name, its
size, its major mode (the kind of text it is, such as Text or Org) and the
file it visits. The marks come first. A \texttt{>} means you marked the buffer for
a command to act on. A \texttt{D} means you marked it for deletion. A \texttt{*} means it
has changes that are not saved, and a \texttt{\%} means it is read-only. In the figure,
\texttt{council.\codebreak{}txt} is marked and has unsaved changes, and \texttt{old-\codebreak{}draft.\codebreak{}txt} is
marked for deletion.

The names of buffers that Tuatara made itself, such as \texttt{*Messages*}, are set
apart in their own colour, as are read-only buffers, help buffers and marked
names. The last line of the list adds up how many buffers there are and how
much text they hold.

\subsubsection{Sections}
\label{sec:150}
The list is divided into sections, one for each project. A project is a
folder that holds a \texttt{.git} or a \texttt{.project} file, with every file in it and
every folder below it. In the figure, \texttt{letters} is a project because
\texttt{\textasciitilde{}/letters/} is under version control. Buffers that belong to no project are
in a section called No project. The section of the project you came from is
first. Each heading says the project’s folder name and how many buffers the
section holds.

A heading starts with \texttt{{\tuataraFallbackTWO{}▾}} when its section is open and \texttt{{\tuataraFallbackTWO{}▸}} when it is folded.
The setting \texttt{*project-\codebreak{}root-\codebreak{}markers*} holds the file names that make a folder a
project. If you keep your own projects in folders marked some other way, add
the name of the marker file there.

\subsubsection{Folding}
\label{sec:151}
A long list is easier to read with the sections you do not need folded away.
A folded section shows only its heading. \texttt{TAB} folds the section the cursor is
in, or opens it again if it was folded. \texttt{S-TAB} folds every section, or opens
them all if every one is folded. On a heading, \texttt{l} and \texttt{RET} open the section,
and \texttt{h} folds it. On any other line \texttt{h} does nothing, so a stray press does
not close the list; \texttt{q} does that.

\subsubsection{Marks}
\label{sec:152}
The list works on one buffer at a time, or on several if you mark them.
Marking is how you say “these ones”. \texttt{m} puts a \texttt{>} on the buffer the cursor
is on and moves down a line, so you can mark a run of buffers by pressing \texttt{m}
over and over. A number before the key marks that many lines. \texttt{u} takes a mark
off the line and moves down, and \texttt{U} takes every mark off. \texttt{t} swaps the two
states: what was marked becomes unmarked and the other way round. That is a
quick way to say “everything except these”.

\texttt{d} puts a \texttt{D} on a buffer instead. It means “kill this one”. Nothing is killed
until you press \texttt{x}, which kills every buffer marked with \texttt{D}. It is the same
two-step as in a folder listing: first mark, then do it. Buffers with unsaved
changes are asked about first, so a mark does not lose your work.

There are two commands that act on the buffers marked with \texttt{>}. \texttt{s} saves the
marked buffers that have a file. \texttt{D}, with a capital, kills the marked buffers
after asking. If you have marked nothing, both act on the buffer the cursor is
on. On a section’s heading, \texttt{D} kills every buffer of that project.

\subsubsection{Opening, and keeping the list current}
\label{sec:153}
\texttt{RET} and \texttt{l} show the buffer on the cursor’s line in the window you came
from. \texttt{o} shows it in another window and leaves the list open beside it. \texttt{n}
and \texttt{p} move to the next and the previous buffer, skipping the headings, and
\texttt{j} and \texttt{k} (non-modal: \texttt{C-n}, \texttt{C-p}) move a line at a time. \texttt{r} makes the list
again. You rarely need that, because the list updates itself while it is on the
screen, and your marks stay.

When Tuatara runs in the background, serving several windows of its own (a
daemon, described in chapter 18), each of those windows is a client and the
list has a section for each client instead. \texttt{K} on a client’s section makes
that client the one that stays when its last buffer is closed. The setting
\texttt{*client-\codebreak{}kept-\codebreak{}without-\codebreak{}buffers*} says which one stays when you do not choose.
You can skip this until chapter 18.

\begin{try}
\textbf{Try it, with a letter and a shopping list}
\begin{enumerate}
\item Press \texttt{C-x C-f}, type \texttt{\textasciitilde{}/letter.\codebreak{}txt} and press \texttt{RET}. Press \texttt{i} and type
   \texttt{Dear Aroha,}. Press \texttt{C-g} and then \texttt{C-x C-s} to save.
\item Open \texttt{\textasciitilde{}/shopping.\codebreak{}txt} the same way and type \texttt{eggs}, \texttt{RET}, \texttt{bread}. Press
   \texttt{C-g} and save with \texttt{C-x C-s}.
\item Press \texttt{C-x C-b}. The two files are listed, the shopping list at the top
   because you used it last. The cursor is on the buffer you came from.
\item Press \texttt{j} (non-modal: \texttt{C-n}) to move to the letter. Press \texttt{m}. A \texttt{>} shows
   before its name and the cursor moves down a line. Press \texttt{k} (non-modal:
   \texttt{C-p}) and \texttt{u} to take the mark off, then \texttt{k} and \texttt{m} to put it back.
\item Move to the shopping list with \texttt{k k} and press \texttt{d}. A \texttt{D} shows.
   Press \texttt{x}. \texttt{shopping.\codebreak{}txt} is gone from the list.
\item Press \texttt{o} on the letter. It opens in another window and the list stays.
   Press \texttt{C-x o} to go back to the list, then \texttt{q} to close it.
\end{enumerate}
\end{try}

\subsubsection{Commands for the buffer list}
\label{sec:154}
\begin{xltabular}{\textwidth}{@{}l l l >{\raggedright\arraybackslash\hsize=1.000\hsize}X@{}}
\toprule
\textbf{Modal} & \textbf{Non-modal} & \textbf{Command} & \textbf{What it does}\\
\midrule
\endhead
none & \texttt{C-x C-b} & \texttt{buffers} & List every live buffer in \textbf{Buffers}, those used most recently first, in a section for each project (a folder holding .git or .project, \textbf{PROJECT-ROOT-MARKERS}: its files, the folders of it listed and the buffers standing in them), the project of the buffer used last first, then one for those in none.\\
none & \texttt{s} & \texttt{buffers-\codebreak{}save} & Save the marked buffers (>) that have a file, or with none marked the buffer on the current line.\\
none & \texttt{K}, \texttt{C-c K} & \texttt{buffers-\codebreak{}keep-\codebreak{}client} & Make the client whose folder the cursor is in, in the buffer list, the one that stays when its last buffer is closed, in place of the one \textbf{CLIENT-KEPT-WITHOUT-BUFFERS} names, until it ends.\\
none & \texttt{l}, \texttt{RET}, \texttt{C-c l} & \texttt{list-\codebreak{}visit} & Open the entry on the current line of the list: the buffer, the bookmark’s place, the file; on a section’s heading, fold or unfold the section.\\
none & \texttt{o}, \texttt{C-c o} & \texttt{list-\codebreak{}visit-\codebreak{}other} & Open the entry on the current line of the list in another window; on a section’s heading, fold or unfold the section.\\
none & \texttt{n}, \texttt{C-c n} & \texttt{list-\codebreak{}next-\codebreak{}entry} & Go to the next entry of the list, the start of its first line, passing over the headings of its sections.\\
none & \texttt{p}, \texttt{C-c p} & \texttt{list-\codebreak{}previous-\codebreak{}entry} & Go to the entry of the list before the one on the current line, the start of its first line, passing over the headings of its sections.\\
none & \texttt{m}, \texttt{C-c m} & \texttt{list-\codebreak{}mark-\codebreak{}entry} & Mark the entry on the current line of the list (>), for the list’s commands that act on those marked, and move down; with a count, as many entries, up the list for a negative one.\\
none & \texttt{d}, \texttt{C-c d} & \texttt{list-\codebreak{}mark-\codebreak{}for-\codebreak{}deletion} & Mark the entry on the current line of the list for deletion (D), for the list’s command that deletes them, and move down; with a count, as many entries, up the list for a negative one.\\
none & \texttt{u}, \texttt{C-c u} & \texttt{list-\codebreak{}unmark-\codebreak{}entry} & Take the mark off the entry on the current line of the list, of any kind, and move down; with a count, off as many entries, up the list for a negative one.\\
none & \texttt{U}, \texttt{C-c U} & \texttt{list-\codebreak{}unmark-\codebreak{}all} & Take every mark off the entries of the list, of any kind; in rimu, those of the files marked in other directories too.\\
none & \texttt{t} & \texttt{list-\codebreak{}toggle-\codebreak{}marks} & Mark the entries the list shows that are not marked and take the mark off those marked; those marked for deletion (D), or given a mark of another kind, stay so, as do the entries a filter or a folded section hides.\\
none & \texttt{x} & \texttt{list-\codebreak{}execute} & Delete the entries of the list marked for deletion (D), after asking, as the list deletes: the buffer list kills its buffers, rimu puts files in the trash, the list of jobs drops those waiting and stops those running.\\
none & \texttt{D}, \texttt{C-c D} & \texttt{list-\codebreak{}delete-\codebreak{}marked} & Delete the marked entries of the list (>), or with none marked the one on the current line, after asking, as the list deletes, as M-x list-execute does those marked for deletion; in rimu, with C-u, for good rather than to the trash.\\
none & \texttt{TAB} & \texttt{list-\codebreak{}toggle-\codebreak{}fold} & Fold the section of the list the current line is in, its heading left showing, or unfold it when it is folded.\\
none & \texttt{S-TAB} & \texttt{list-\codebreak{}toggle-\codebreak{}all-\codebreak{}folds} & Fold every section of the list, or unfold them all when every one is folded; a section not in the list now keeps its fold for when it comes back.\\
\bottomrule
\end{xltabular}

\subsubsection{Settings and faces}
\label{sec:155}
\begin{xltabular}{\textwidth}{@{}l l >{\raggedright\arraybackslash\hsize=1.000\hsize}X@{}}
\toprule
\textbf{Setting} & \textbf{Standard} & \textbf{What it changes}\\
\midrule
\endhead
\texttt{*project-\codebreak{}root-\codebreak{}markers*} & \texttt{(".\codebreak{}git" ".project")} & The names of the files or directories that mark the root of a project: M-x search-in-files searches all of the project a buffer’s file is in, from the nearest directory above the file holding one of them.\\
\texttt{*client-\codebreak{}kept-\codebreak{}without-\codebreak{}buffers*} & \texttt{:first} & Which clients of a daemon stay when their last buffer is closed: :FIRST, the one opened first of those running; :LAST, only the last one left; :ALL, every one; NIL, none.\\
\bottomrule
\end{xltabular}

\begin{xltabular}{\textwidth}{@{}l >{\raggedright\arraybackslash\hsize=1.000\hsize}X@{}}
\toprule
\textbf{Face} & \textbf{What it paints}\\
\midrule
\endhead
\texttt{:list-\codebreak{}marked} & An entry marked (>) in any list whose entries are marked, rimu’s too, for the commands that act on those marked: the mark and the name, in the colour of warnings.\\
\texttt{:list-\codebreak{}deletion} & An entry marked for deletion (D) in any list whose entries are marked, rimu’s too: the mark and the name, in the colour of errors.\\
\texttt{:list-\codebreak{}dim} & What is of less interest in a list, set back: help buffers and the headings of clients that ended in the buffer list (C-x C-b), packages not loaded, bookmarks whose file is gone, files ignored in rimu and in a git status, a value missing in the health report (M-x check-health), a line saying more are not shown.\\
\texttt{:list-\codebreak{}special} & Something Tuatara made rather than a file of yours, in a list: the names in stars in the buffer list (C-x C-b), \textbf{Messages} and \textbf{scratch}.\\
\bottomrule
\end{xltabular}

\subsection{Scratch buffers}
\label{sec:157}
Sometimes you want somewhere to type that is not a file: a number to add up,
a phone number to keep in front of you for a minute, a few lines you may
throw away. Tuatara has two buffers for this. They have no file, so they
are never saved, and Tuatara has no file to ask about when you quit.

\texttt{*lisp-\codebreak{}scratch*} is a buffer for Lisp, the language Tuatara is written in.
You can use it as plain text too. Type a Lisp expression, such as \texttt{(+ 2 3)},
and press \texttt{C-x C-e} with the cursor after it, and Tuatara runs it. \texttt{C-j} puts
the answer on the lines after the expression. That makes it a calculator that
shows its working. \texttt{*org-\codebreak{}scratch*} is the same idea for notes in Org, the
outline format described in chapter 14.

When Tuatara starts without a file, it shows one of them, chosen by
\texttt{*startup-\codebreak{}buffer*}. It also makes both scratch buffers for you, a list that
\texttt{*auto-\codebreak{}scratch-\codebreak{}buffers*} holds, and makes them again when you kill your last
buffer so that you are never left with an empty screen. If you have killed one
and want it back, run \texttt{lisp-\codebreak{}scratch} or \texttt{org-\codebreak{}scratch} with \texttt{M-x}; they make a
new one when none is there.

\begin{xltabular}{\textwidth}{@{}l l l >{\raggedright\arraybackslash\hsize=1.000\hsize}X@{}}
\toprule
\textbf{Modal} & \textbf{Non-modal} & \textbf{Command} & \textbf{What it does}\\
\midrule
\endhead
none & \texttt{M-x lisp-\codebreak{}scratch} & \texttt{lisp-\codebreak{}scratch} & Show \textbf{lisp-scratch}, the Lisp buffer that is not saved, where C-x C-e evaluates in the running editor and C-j puts what the expression before point printed and its value on the lines after it; when it was killed, a new one.\\
none & \texttt{M-x org-\codebreak{}scratch} & \texttt{org-\codebreak{}scratch} & Show \textbf{org-scratch}, the Org buffer for notes that are not saved; when it was killed, a new one.\\
\bottomrule
\end{xltabular}

\begin{xltabular}{\textwidth}{@{}l >{\raggedright\arraybackslash\hsize=0.615\hsize}X >{\raggedright\arraybackslash\hsize=1.385\hsize}X@{}}
\toprule
\textbf{Setting} & \textbf{Standard} & \textbf{What it changes}\\
\midrule
\endhead
\texttt{*startup-\codebreak{}buffer*} & \texttt{"*lisp-\codebreak{}scratch*"} & The buffer shown when the editor starts without files: “\textbf{lisp-scratch}” or “\textbf{org-scratch}” (or any buffer the init file makes).\\
\texttt{*auto-\codebreak{}scratch-\codebreak{}buffers*} & \texttt{("*lisp-\codebreak{}scratch*" "*org-\codebreak{}scratch*")} & The scratch buffers made without being asked for: at startup, and when the last buffer is killed.\\
\bottomrule
\end{xltabular}

A file name typed in a scratch buffer is taken from the folder in
\texttt{*default-\codebreak{}directory*}, which is your home folder unless you say otherwise.
Saving a scratch buffer with \texttt{C-x C-s} asks where to put it, as for any new
file.

\subsection{Saving and killing buffers}
\label{sec:158}
Killing a buffer removes it from memory. If it visited a file, the file stays
where it is. The usual way is \texttt{C-x k}, which kills the buffer you are in. If it
has changes you have not saved, Tuatara asks first. Chapter 1 mentioned \texttt{C-x 4
0}, which does the same and closes the window too.

When you have finished with a whole piece of work, \texttt{C-x p k} kills every
buffer of the project the current buffer is in: its files, its folder
listings and the buffers standing in them. It says how many have unsaved
changes and asks about each of those. It is the quick way to tidy up after
you have finished the letters and want to start on something else. The same
thing is done by pressing \texttt{D} on the heading of the project in the buffer
list.

\texttt{C-x s} goes the other way. It saves every file buffer with unsaved changes,
after asking once. It is handy after a rename that touched many files. Leaving
Tuatara with \texttt{C-x C-c} asks first when any buffer visiting a file has changes
that are not saved.

\begin{xltabular}{\textwidth}{@{}l l l >{\raggedright\arraybackslash\hsize=1.000\hsize}X@{}}
\toprule
\textbf{Modal} & \textbf{Non-modal} & \textbf{Command} & \textbf{What it does}\\
\midrule
\endhead
none & \texttt{C-x k} & \texttt{kill-\codebreak{}buffer} & Kill the current buffer, asking first when it has changes not saved.\\
none & \texttt{C-x p k} & \texttt{kill-\codebreak{}project-\codebreak{}buffers} & Kill every buffer of the project the current buffer is in (a folder above it holding one of \textbf{PROJECT-ROOT-MARKERS}, .git or .project), its files, its folders listed (C-x d), and the buffers standing in them; after asking, saying how many are modified, each modified one asked about as it goes.\\
none & \texttt{C-x s} & \texttt{save-\codebreak{}some-\codebreak{}buffers} & Save every buffer visiting a file with changes not saved, after asking once, the files named: the files a rename across a project changed, say.\\
none & \texttt{C-x C-c} & \texttt{save-\codebreak{}buffers-\codebreak{}kill-\codebreak{}editor} & Exit the editor, asking first if a buffer visiting a file, or one made by C-x n e, has unsaved changes.\\
\bottomrule
\end{xltabular}

\subsection{Recent files and places}
\label{sec:159}
Most days you work on the same few files. You do not have to type their names
each time. Tuatara keeps a list of the files you have opened, the latest
first, and \texttt{C-x C-r} asks you to pick from it. Each file is listed with its
permissions, its size and when it last changed, which helps when two files
have similar names. The list is kept between sessions, so on Monday morning
it still holds Friday’s files.

The list keeps \texttt{*recent-\codebreak{}files-\codebreak{}limit*} files. A file that no longer exists is
taken out when Tuatara starts. If you want to keep such files in the list,
shown as not found, set \texttt{*recent-\codebreak{}files-\codebreak{}auto-\codebreak{}cleanup*} to \texttt{nil}. To keep no
list at all, set \texttt{*remember-\codebreak{}recent-\codebreak{}files*} to \texttt{nil}.

Tuatara also remembers where the cursor was in each file you leave. When you
open the file again, the cursor is back at that line, even in a later session.
That is the setting \texttt{*save-\codebreak{}place-\codebreak{}mode*}, and \texttt{*save-\codebreak{}place-\codebreak{}limit*} sets how
many files it remembers. It is useful for a long document you read a little
at a time.

\begin{try}
\textbf{Try it, with a letter and some notes}
\begin{enumerate}
\item Press \texttt{C-x C-f} and open \texttt{\textasciitilde{}/letter.\codebreak{}txt}, then use \texttt{j} to go down a few
   lines. If the file is short, type a few lines first and save.
\item Open \texttt{\textasciitilde{}/notes.\codebreak{}txt} the same way. If it is a new file, type a word in it
   and save it, as a file that was never written is not remembered. Press
   \texttt{C-x k} to kill the notes buffer, then \texttt{C-x k} again for the letter.
\item Press \texttt{C-x C-r}. Both files are listed, \texttt{notes.\codebreak{}txt} first.
\item Pick the letter and press \texttt{RET}. It opens with the cursor at the line
   where you left it.
\end{enumerate}
\end{try}

\begin{xltabular}{\textwidth}{@{}l l l >{\raggedright\arraybackslash\hsize=1.000\hsize}X@{}}
\toprule
\textbf{Modal} & \textbf{Non-modal} & \textbf{Command} & \textbf{What it does}\\
\midrule
\endhead
none & \texttt{C-x C-r} & \texttt{find-\codebreak{}recent-\codebreak{}file} & Visit a file visited lately: the latest first, each with its permissions, size and when it last changed beside it.\\
\bottomrule
\end{xltabular}

\begin{xltabular}{\textwidth}{@{}l l >{\raggedright\arraybackslash\hsize=1.000\hsize}X@{}}
\toprule
\textbf{Setting} & \textbf{Standard} & \textbf{What it changes}\\
\midrule
\endhead
\texttt{*remember-\codebreak{}recent-\codebreak{}files*} & \texttt{t} & True to keep a list of the files you visit, the latest first, for find-recent-file (M-x find-recent-file, M-x find-recent-file) to offer; NIL adds nothing to it.\\
\texttt{*recent-\codebreak{}files-\codebreak{}limit*} & \texttt{25} & How many files the list of recent files keeps, 25 as standard.\\
\texttt{*recent-\codebreak{}files-\codebreak{}auto-\codebreak{}cleanup*} & \texttt{t} & True to take out of the list of recent files, when the editor starts, the files no longer there; NIL keeps them, shown as not found.\\
\texttt{*save-\codebreak{}place-\codebreak{}mode*} & \texttt{t} & True to remember where point was in each file you visit and put it back there when you visit the file again, in this session or a later one (on as standard); NIL to open files at their top.\\
\texttt{*save-\codebreak{}place-\codebreak{}limit*} & \texttt{400} & How many files’ places are kept, the latest first, 400 as standard.\\
\bottomrule
\end{xltabular}

\subsection{Bookmarks}
\label{sec:160}
A bookmark is a named place that Tuatara keeps for good. It is like a recent
file, but you choose it and give it a name. You would use one for the place in
a long letter where you stopped, or for the section of your notes that you
open every day. A bookmark remembers the file and the place in it, and it
finds the place again even when you have added text above it.

\texttt{C-x m m} sets a bookmark where the cursor is. It asks for a name, and
offers the file’s name. It then asks which categories the bookmark is in. A
category is a word of your own, such as \texttt{letters} or \texttt{house}; you may type
several, apart by commas, and the ones you used before are offered. \texttt{C-x m b}
asks for a name and goes to the place. \texttt{C-x m d} forgets a bookmark.

\texttt{C-x m l} opens the bookmark manager, a list that works like the buffer list.
It has a section for each category, with a heading saying how many bookmarks
the section holds, and one for the bookmarks in no category. A bookmark in two
categories appears in both. Folding, marking and the \texttt{RET}, \texttt{l} and \texttt{o} keys
work as described for the buffer list. In the manager, \texttt{r} gives the bookmark
under the cursor another name and \texttt{c} changes its categories. A bookmark
whose file has gone is shown faint.

A bookmark on a folder listing keeps the folder and the entry the cursor is
on, so you can bookmark a folder you visit often.

\begin{try}
\textbf{Try it, with some notes}
\begin{enumerate}
\item Open \texttt{\textasciitilde{}/notes.\codebreak{}txt} and type a few lines, such as \texttt{Plumber: Wanaka Plumbing},
   \texttt{Dentist: 9 May}, \texttt{Car service: Taupo Motors}, each ended with \texttt{RET}. Save.
\item Move to the \texttt{Dentist} line. Press \texttt{C-x m m}, type \texttt{dentist} and \texttt{RET}. At the
   question about categories, type \texttt{house} and \texttt{RET}. The question comes
   again, for one more category: press \texttt{RET} on the empty answer to finish.
\item Press \texttt{C-x C-r} and open some other file, so that the notes are out of
   sight. Press \texttt{C-x m b}, type \texttt{dent} and \texttt{RET}. You are back on the line.
\item Press \texttt{C-x m l}. The bookmark is in a section called \texttt{house}. Press \texttt{q} to
   close the list.
\end{enumerate}
\end{try}

\begin{xltabular}{\textwidth}{@{}l l l >{\raggedright\arraybackslash\hsize=1.000\hsize}X@{}}
\toprule
\textbf{Modal} & \textbf{Non-modal} & \textbf{Command} & \textbf{What it does}\\
\midrule
\endhead
none & \texttt{C-x m m} & \texttt{bookmark-\codebreak{}set} & Keep where point is in the current file (or the folder a directory list shows and the entry the cursor is on) under a name you give, the file’s name offered, then in the categories you choose, those used already offered, new ones typed (\textbf{BOOKMARK-SET-ASKS-CATEGORIES}).\\
none & \texttt{C-x m b} & \texttt{bookmark-\codebreak{}jump} & Go to a bookmark, chosen by its name (each listed with its file): its file visited, point where it was set, found again when the text above it changed.\\
none & \texttt{C-x m d} & \texttt{bookmark-\codebreak{}delete} & Forget a bookmark, chosen by its name.\\
none & \texttt{C-x m l} & \texttt{bookmark-\codebreak{}manager} & Manage the bookmarks in a list, as the buffer list (C-x C-b) manages buffers, in sections: one for each category (C-x m m asks for them, c here sets them), its heading saying how many it holds, then one for the bookmarks in none; a bookmark in two categories is in both.\\
none & \texttt{r}, \texttt{C-c r} & \texttt{bookmark-\codebreak{}rename} & Give the bookmark on the current line another name, in every section it is in; a name another bookmark has already is taken from it after asking.\\
none & \texttt{c}, \texttt{C-c c} & \texttt{bookmark-\codebreak{}set-\codebreak{}categories} & Say which categories the bookmark on the current line is in, or the marked ones (>) are, the sections holding them unfolded first: those used already offered, new ones typed, several apart by commas; the bookmark’s own written in first, to change.\\
\bottomrule
\end{xltabular}

\begin{xltabular}{\textwidth}{@{}l l >{\raggedright\arraybackslash\hsize=1.000\hsize}X@{}}
\toprule
\textbf{Setting} & \textbf{Standard} & \textbf{What it changes}\\
\midrule
\endhead
\texttt{*bookmark-\codebreak{}set-\codebreak{}asks-\codebreak{}categories*} & \texttt{t} & True for M-x bookmark-set to ask, after the name, which categories the bookmark is in (those used already offered, new ones typed); NIL never to ask.\\
\texttt{*bookmarks-\codebreak{}kept*} & \texttt{1000} & How many bookmarks you may have before setting another says you have more.\\
\bottomrule
\end{xltabular}

\subsection{Registers}
\label{sec:161}
A register is a place to keep text, a number or a position under a single
character, for as long as Tuatara runs. A register that keeps a position is a
short-term bookmark. Registers are described in chapter 4, with the kill ring
that they sit beside. Bookmarks are the choice when you want the place next
week.

\subsection{The session kept}
\label{sec:162}
When you close Tuatara with \texttt{C-x C-c}, everything that was open is gone
from the screen. The session is the name for what was open: the files and the
folder listings, and how the windows were split. If you turn on
\texttt{*keep-\codebreak{}session*}, Tuatara writes the session down as it stops, and
opens it again the next time it starts without a file named. You see your
windows as they were, each on its file, at the place where the cursor was.
It is off until you turn it on.

Put this line in your init file, the file described in chapter 8:

\begin{Verbatim}[breaklines,breakautoindent=false,breakindent=1em,breakafter=/-.\,;:=,commandchars=\\\{\}]
(setf *keep-session* t)
\end{Verbatim}

The session is also written now and then while Tuatara runs, as often as
\texttt{*keep-\codebreak{}session-\codebreak{}interval*} seconds when something has changed. If Tuatara
crashes, you lose at most that long. The file that holds it is in Tuatara’s
state folder and only you can read it. A file you opened with care, because
you do not trust it, is opened with the same care next time. A file that is
no longer there is left out.

If you start Tuatara with a file name, as in \texttt{tuatara letter.\codebreak{}txt}, you get that
file, and the session’s files are opened behind it. The session is not
replaced by that start.

\begin{try}
\textbf{Try it, with a letter and a shopping list}
\begin{enumerate}
\item Add \texttt{(setf *keep-\codebreak{}session* t)} to your init file and restart Tuatara.
\item Open \texttt{\textasciitilde{}/letter.\codebreak{}txt}. Press \texttt{C-x 3} to split the window, then \texttt{C-x o} to go
   to the new window and open \texttt{\textasciitilde{}/shopping.\codebreak{}txt} there.
\item Press \texttt{C-x C-c}. Tuatara leaves.
\item Start Tuatara again with no file name. Both windows are there, each on its
   own file. The echo area says the session was opened again.
\end{enumerate}
\end{try}

\begin{xltabular}{\textwidth}{@{}l l >{\raggedright\arraybackslash\hsize=1.000\hsize}X@{}}
\toprule
\textbf{Setting} & \textbf{Standard} & \textbf{What it changes}\\
\midrule
\endhead
\texttt{*keep-\codebreak{}session*} & \texttt{nil} & True for the editor to keep its session: the files and directories open and how the windows show them, written as it stops (and now and then while it runs), opened again as it was when it starts without files.\\
\texttt{*keep-\codebreak{}session-\codebreak{}interval*} & \texttt{30} & Seconds between the writings of the session while the editor runs, when it has changed (\textbf{KEEP-SESSION}), 30 as standard: what a crash can lose.\\
\bottomrule
\end{xltabular}

\subsection{Auto-save and getting changes back}
\label{sec:163}
Saving is your decision. Tuatara never writes your file until you press
\texttt{C-x C-s}. But it does not leave your unsaved typing exposed to a crash or a
slip. When a buffer that visits a file has changes you have not saved, Tuatara
writes them to a copy of their own, in a folder of its own, a few seconds
after you stop typing. If you keep typing without a pause, it still writes
one every minute. The file you are editing is never touched. When you save,
or kill the buffer with nothing to save, the copy is deleted.

If Tuatara stops without your saving, whether because of a crash, a power cut or
a quit without saving, the copy stays. When you open the file again, Tuatara
tells you it holds changes that are newer than the file. \texttt{M-x recover-\codebreak{}file}
takes them into the buffer, after asking. They arrive as one edit. \texttt{C-/} (modal:
\texttt{u}) takes it back if you change your mind, and \texttt{C-x C-s} writes it to the file
if you do not.

You can also see every copy that is waiting with \texttt{M-x list-\codebreak{}auto-\codebreak{}saves}. The
list shows when each copy was written, how big it is and which file it
belongs to, with the newest first. It marks and acts as the buffer list
does: \texttt{RET} recovers the copy on the line, after asking, and \texttt{x} and \texttt{D} delete
copies.

Auto-save is on as standard. A buffer larger than \texttt{*auto-\codebreak{}save-\codebreak{}size-\codebreak{}limit*}
octets is not copied, because each copy is written whole. Files opened with
care are never copied. \texttt{*auto-\codebreak{}save-\codebreak{}directory*} says where the copies go, and
you can point it at a folder on another drive. \texttt{*auto-\codebreak{}save-\codebreak{}idle*} and
\texttt{*auto-\codebreak{}save-\codebreak{}interval*} set the two waits.

\begin{try}
\textbf{Try it, with a letter}
\begin{enumerate}
\item Open \texttt{\textasciitilde{}/letter.\codebreak{}txt}. Press \texttt{i}, type \texttt{P.S. Thank you for the plums.} and
   press \texttt{C-g}. Do not save. Wait ten seconds.
\item Press \texttt{C-x C-c}. Tuatara asks about the unsaved file. Leave without
   saving it.
\item Start Tuatara again and open the letter. The echo area says it holds
   changes newer than the file.
\item Press \texttt{M-x}, type \texttt{recover-\codebreak{}file} and press \texttt{RET}. Answer yes. The postscript
   is back. Press \texttt{C-x C-s} to save it.
\end{enumerate}
\end{try}

\begin{xltabular}{\textwidth}{@{}l l l >{\raggedright\arraybackslash\hsize=1.000\hsize}X@{}}
\toprule
\textbf{Modal} & \textbf{Non-modal} & \textbf{Command} & \textbf{What it does}\\
\midrule
\endhead
none & \texttt{M-x list-\codebreak{}auto-\codebreak{}saves} & \texttt{list-\codebreak{}auto-\codebreak{}saves} & List the auto-saved copies of changes not saved (\textbf{AUTO-SAVE-DIRECTORY}), the newest first: after a crash or a quit without saving, what can be taken back.\\
none & \texttt{M-x recover-\codebreak{}file} & \texttt{recover-\codebreak{}file} & Take into the current buffer the changes auto-saved for its file (\textbf{AUTO-SAVE}) after asking, as one edit: C-/ takes it back, a save writes it to the file.\\
\bottomrule
\end{xltabular}

\begin{xltabular}{\textwidth}{@{}l l >{\raggedright\arraybackslash\hsize=1.000\hsize}X@{}}
\toprule
\textbf{Setting} & \textbf{Standard} & \textbf{What it changes}\\
\midrule
\endhead
\texttt{*auto-\codebreak{}save*} & \texttt{t} & True for buffers visiting files to have their changes not saved written now and then to a copy of their own in \textbf{AUTO-SAVE-DIRECTORY}, never to the file itself; M-x recover-file takes such a copy back.\\
\texttt{*auto-\codebreak{}save-\codebreak{}directory*} & \texttt{nil} & The folder auto-saved copies are kept in, made for you alone when it is not there: NIL for auto-save/ in Tuatara’s state folder (\$XDG\_STATE\_HOME/tuatara/, or \textasciitilde{}/.local/state/tuatara/), or a folder’s name, “\textasciitilde{}/backups/tuatara/“.\\
\texttt{*auto-\codebreak{}save-\codebreak{}idle*} & \texttt{5} & Seconds without a change after which a buffer’s changes not saved are auto-saved (\textbf{AUTO-SAVE}), 5 as standard.\\
\texttt{*auto-\codebreak{}save-\codebreak{}interval*} & \texttt{60} & The most seconds changes not saved go without an auto-save while they go on without a pause (\textbf{AUTO-SAVE}), 60 as standard.\\
\texttt{*auto-\codebreak{}save-\codebreak{}size-\codebreak{}limit*} & \texttt{268435456} & The largest buffer auto-saved, in octets, 256 MB as standard: a larger one is not copied (each copy would be written whole), so a crash loses its changes not saved.\\
\bottomrule
\end{xltabular}

\subsection{Narrowing to a part of a buffer}
\label{sec:164}
A long file can be a distraction. If you want to rewrite one paragraph of a
letter, you do not want a search to wander off into another part of the file,
or a replacement to change the other places by mistake. Narrowing shows only a
part of the buffer and keeps every command inside it. Moving, searching,
replacing, editing and selecting all stop at the edges of the part. The rest
of the buffer is out of sight and out of reach until you widen it again.

While a buffer is narrowed, the mode line says Narrow in a warning colour, so
you cannot forget it. Nothing is lost. Saving still writes the whole file, and
the line numbers stay those of the whole file.

\texttt{C-x x n} narrows to the text you have selected. \texttt{C-x x s} narrows to the
section the cursor is in. In an Org file that is the heading it is under, with
everything below it. In a program it is the function. \texttt{C-x x w} shows all of
the buffer again.

Going to a place outside the part, with a bookmark for example, widens the
buffer and says so. A command that would change text outside the part is
refused, and the echo area says so.

\begin{try}
\textbf{Try it, with a shopping list}
\begin{enumerate}
\item Open \texttt{\textasciitilde{}/shopping.\codebreak{}txt} and make it eight lines long: \texttt{eggs}, \texttt{bread},
   \texttt{milk}, \texttt{butter}, \texttt{apples}, \texttt{tea}, \texttt{rice} and \texttt{honey}, one per line. Save.
\item Press \texttt{M-<} (modal: \texttt{g g}) to go to the top, then \texttt{j j} to reach the line
   \texttt{milk}. Press \texttt{V} and \texttt{J J} to select \texttt{milk}, \texttt{butter} and \texttt{apples}.
\item Press \texttt{C-x x n}. Only those three lines are left on the screen, and Narrow
   shows in the mode line.
\item Press \texttt{C-s}, type \texttt{rice} and \texttt{RET}. It is not found, because it is out of
   reach.
\item Press \texttt{C-x x w}. All eight lines are back, and Narrow is gone.
\end{enumerate}
\end{try}

\begin{xltabular}{\textwidth}{@{}l l l >{\raggedright\arraybackslash\hsize=1.000\hsize}X@{}}
\toprule
\textbf{Modal} & \textbf{Non-modal} & \textbf{Command} & \textbf{What it does}\\
\midrule
\endhead
none & \texttt{C-x x n} & \texttt{narrow-\codebreak{}to-\codebreak{}region} & Show only the selected text, and keep every command to it: motion, search and replacement, editing and selecting stop at its edges, and the rest of the buffer is out of sight and out of reach until C-x x w.\\
none & \texttt{C-x x s} & \texttt{narrow-\codebreak{}to-\codebreak{}section} & Show only the section point is in, and keep every command to it, as C-x x n does for the selection: in Org, the subtree of the headline point is under, the headline with it; in code, the definition point is in (or the one after it), as C-M-a and C-M-e find it.\\
none & \texttt{C-x x w} & \texttt{widen} & Show all of the buffer again after C-x x n or C-x x s, every command reaching all of it, point staying where it is.\\
\bottomrule
\end{xltabular}

\begin{xltabular}{\textwidth}{@{}l >{\raggedright\arraybackslash\hsize=1.000\hsize}X@{}}
\toprule
\textbf{Face} & \textbf{What it paints}\\
\midrule
\endhead
\texttt{:mode-\codebreak{}line-\codebreak{}narrowed} & Narrow, in the mode line of a buffer narrowed to a part of it (C-x x n, C-x x s): lit up in the theme’s warning colour, so that it is seen at once that the rest of the buffer is out of sight and out of reach until C-x x w.\\
\bottomrule
\end{xltabular}

\section{Commands, keys and help}
\label{sec:185}
Running a command by its name, the keys that start longer keys, and the
help that Tuatara keeps for itself.

Everything Tuatara does is a command. Moving down a line is a command, so is
saving a file, and so is sorting a list. Most commands have a key, but many
do not, and nobody remembers all of them. This chapter is about the ways of
finding and running a command without knowing its key. It starts with \texttt{M-x},
which runs any command by name, and with the count that you can give to a
command. It then explains the prefix keys, the keys that start a longer key,
and the key hints that show you what can follow. The longest part is about
help: how to ask what a key does, what a command does, what a setting is and
where to look when you do not know the name. The end of the chapter has the
page that introduces Tuatara, the check of what your system has, and a short
look at running Lisp in the editor.

Before you start, one reminder about the state keys. \texttt{i} starts INSERT state,
and \texttt{C-g} or \texttt{ESC} returns to NORMAL state. The keys in this chapter that
start with \texttt{C-x}, \texttt{C-h} or \texttt{M-x} work in both states, and the others are
pressed in NORMAL state.

\subsection{Running a command by name}
\label{sec:168}
\texttt{M-x} asks for the name of a command in the minibuffer, the row at the bottom
of the screen where Tuatara asks its questions, and runs the command you
choose. It is the way to run anything that has no key, and the way to run
something you only need once in a while. You use it when you know what you
want done but not which key does it.

You do not type the whole name. As you type, Tuatara lists the commands that
match, and each one shows its keys and the first line of what it does. The
words you type may be in any order, so \texttt{lines sort} finds \texttt{sort-\codebreak{}lines}. The
first letters of the words also match, so \texttt{sbke} finds
\texttt{save-\codebreak{}buffers-\codebreak{}kill-\codebreak{}editor}. \texttt{C-n} and \texttt{C-p} move down and up the list, \texttt{TAB}
puts the chosen name into what you typed, and \texttt{RET} runs it. \texttt{C-g} gives up and
nothing runs. The commands you ran lately come first, so a command you use
often needs only a few letters.

Figure 8.1 shows the minibuffer after typing \texttt{sort lines} with some text
selected.

\begin{figure}[H]
\caption{Figure 8.1. M-x with its list of choices. The prompt shows the number of the chosen match and how many there are.}\label{org166}
\begin{Verbatim}[breaklines,breakautoindent=false,breakindent=1em,breakafter=/-.\,;:=,commandchars=\\\{\}]
 | 1/6  M-x sort lines                                             |
 | sort-lines          Sort the lines the selection touches by ...  |
 | sort-lines-reverse  Sort the lines the selection touches by ...  |
 | sort-fields         Sort the lines the selection touches by ...  |
 | sort-numeric-fields Sort the lines the selection touches by ...  |
 | sort-columns        Sort the lines from the one where the ...    |
\end{Verbatim}
\end{figure}

A command that needs something from you asks for it after you press \texttt{RET}, in
the same row. \texttt{sort-\codebreak{}lines}, for example, works on the lines you have
selected, and \texttt{flush-\codebreak{}lines} asks for the text to look for. A command made for
one kind of buffer says so when you run it in another, and tells you where
it works.

\begin{try}
\textbf{Try it, to find a command for sorting lines}
\begin{enumerate}
\item Press \texttt{C-x C-f}, type \texttt{\textasciitilde{}/shopping.\codebreak{}txt} and press \texttt{RET}. Press \texttt{i} and type
   \texttt{milk}, \texttt{RET}, \texttt{eggs}, \texttt{RET}, \texttt{bread}, \texttt{RET}, \texttt{apples}. Press \texttt{C-g}.
\item Press \texttt{M-<} (modal: \texttt{g g}) to go to the top. Press \texttt{V} to select the first
   line, then \texttt{J J J} to take in the other three.
\item Press \texttt{M-x} and type \texttt{sort lines}. The list shows \texttt{sort-\codebreak{}lines} first, with
   its first line of help beside it. Press \texttt{RET}.
\item The four lines are in alphabetical order: \texttt{apples}, \texttt{bread}, \texttt{eggs},
   \texttt{milk}. Press \texttt{C-/} (modal: \texttt{u}) to put them back.
\item Press \texttt{M-x} and type \texttt{sbke}. The list shows \texttt{save-\codebreak{}buffers-\codebreak{}kill-\codebreak{}editor}.
   Look at the key beside it, then press \texttt{C-g}. Nothing ran.
\end{enumerate}
\end{try}

\begin{xltabular}{\textwidth}{@{}l l l >{\raggedright\arraybackslash\hsize=1.000\hsize}X@{}}
\toprule
\textbf{Modal} & \textbf{Non-modal} & \textbf{Command} & \textbf{What it does}\\
\midrule
\endhead
none & \texttt{M-x} & \texttt{execute-\codebreak{}extended-\codebreak{}command} & Run a command chosen by name, with completion over every command: the candidates show each command’s keys and the first line of its documentation.\\
\bottomrule
\end{xltabular}

\subsubsection{The count before M-x}
\label{sec:167}
A count given before \texttt{M-x} goes to the command you choose. \texttt{C-u 3 M-x
upcase-\codebreak{}word} makes the next three words capital letters. The next section
says how a count is typed.

\subsection{Counts}
\label{sec:169}
A count tells a command how many times to do its work, or gives it a number to
work with. You type \texttt{C-u}, then the number, then the key or command. In NORMAL
state \texttt{C-u 5 j} moves five lines down. \texttt{C-u} alone means 4 and \texttt{C-u C-u} means
16, and a minus sign makes the count negative, which turns many commands
around. Chapter 2 shows counts with moves, and the other chapters name what a
count does for their commands, such as how many cursors to add or how many
lines to kill. The count is always taken by the command that follows it.

\texttt{.} (non-modal: \texttt{C-z .}) runs the last command again, with the count it had. It
is the quickest way to do the same small thing three times. After \texttt{C-u 4 j},
pressing \texttt{.} goes four more lines down. A command started with \texttt{M-x} is
repeated as well, with its count.

Inside a question, \texttt{C-u} does something else. In the minibuffer it clears what
you typed, and leaves the question as it is. That is useful when you started
a long name and want to begin again.

\begin{xltabular}{\textwidth}{@{}l l l >{\raggedright\arraybackslash\hsize=1.000\hsize}X@{}}
\toprule
\textbf{Modal} & \textbf{Non-modal} & \textbf{Command} & \textbf{What it does}\\
\midrule
\endhead
none & \texttt{C-u} & \texttt{universal-\codebreak{}argument} & Begin a numeric prefix argument: C-u alone means 4, C-u C-u 16, digits after it give their number, a leading minus negates.\\
\texttt{.} & \texttt{C-z .} & \texttt{repeat} & Run the last command again, other than repeat itself, with the prefix argument it had (one given to repeat is used instead) and as the key that ran it (C-x C-- shrinks again, it does not grow); repeat again, and it runs that command once more, as often as you like.\\
\bottomrule
\end{xltabular}

\subsection{Giving up}
\label{sec:170}
\texttt{C-g} is the key to press when you are unsure. It closes a question, closes a
list of choices, closes the key hints, drops a selection and leaves INSERT
state. Whatever you were in the middle of is abandoned, and nothing is run.

\texttt{ESC} is a key of its own as standard. It does what it is bound to, and where
it is bound to nothing it does what \texttt{C-g} does. It closes a question typed in
the minibuffer and leaves INSERT state, but it does not end a question that is
answered with one key. If you would rather have \texttt{ESC} do everything that \texttt{C-g}
does, set \texttt{*escape-\codebreak{}quits*} to \texttt{t} in your init file, the file described in
chapter 20.

\begin{xltabular}{\textwidth}{@{}l l l >{\raggedright\arraybackslash\hsize=1.000\hsize}X@{}}
\toprule
\textbf{Modal} & \textbf{Non-modal} & \textbf{Command} & \textbf{What it does}\\
\midrule
\endhead
none & \texttt{M-x keyboard-\codebreak{}quit} & \texttt{keyboard-\codebreak{}quit} & Abandon the current operation and deactivate the region; what stays shown until then goes, such as the matches of the last search.\\
\bottomrule
\end{xltabular}

\begin{longtable}{@{}l l l@{}}
\toprule
\textbf{Setting} & \textbf{Standard} & \textbf{What it changes}\\
\midrule
\endhead
\texttt{*escape-\codebreak{}quits*} & \texttt{nil} & True to have Escape do what C-g does everywhere.\\
\bottomrule
\end{longtable}

\subsection{Prefix keys}
\label{sec:171}
Most of the keys you press are one key. Longer keys are made by pressing a
prefix first and then another key, as in \texttt{C-x C-f}. A prefix key gathers a
family of commands under one key, so that a small number of keys at the top
leads to many commands. Tuatara uses four of them in particular ways.

\texttt{C-x} is the family of global commands. Its keys are the same in every buffer
and in both states: \texttt{C-x C-f} finds a file, \texttt{C-x C-s} saves it, \texttt{C-x b} goes to
another buffer, \texttt{C-x C-b} lists them. Under it are smaller families with a
name: \texttt{C-x w} for windows, \texttt{C-x r} for rectangles and registers, \texttt{C-x m} for
bookmarks, \texttt{C-x p} for projects, \texttt{C-x T} for things you switch on and off.
The chapter for each subject says which keys it uses.

\texttt{C-c} is the family of the buffer you are in. Its keys differ from mode to
mode. In a list of buffers, \texttt{C-c} shows what the list can do. In an Org note
it shows what Org can do, and in a PDF it shows how to turn a page. It is the
one prefix to try when you wonder what a particular buffer offers. In the
terminal that Tuatara can run, \texttt{C-c} is not a prefix, because it goes to the
program that is running there. The keys \texttt{C-c x}, \texttt{C-c s} and \texttt{C-c t} are never
given to a mode, and are left for your own commands.

\texttt{SPC} in NORMAL state is the leader, a prefix that belongs to you. As Tuatara
comes, nothing is under it. You put there the keys that you use most, in the
places that suit you, and nothing the editor updates will move them.
\texttt{C-x SPC} opens the same keys in every buffer and in both states, so that the
leader works in INSERT state too. If another key suits you better than \texttt{SPC},
set \texttt{*leader-\codebreak{}key*} to it, for example \texttt{","}. How to bind keys is the subject
of chapter 20.

\texttt{C-h} is the family of help, described below.

Two more prefixes you meet in NORMAL state are \texttt{g}, which starts moves such as
\texttt{g g} (non-modal: \texttt{M-<}) for the top of the buffer, and \texttt{m}, which starts the
selecting keys of chapter 5. They work the same way, and the key hints show
them.

Tuatara has a setting that turns the modal states off. With \texttt{*modal-\codebreak{}editing*}
set to \texttt{nil} every key writes text as it does in a plain editor, and the
Control and Meta keys do the commands. The prefixes in this section work in
the same way. \texttt{C-z} is a prefix of its own for this: under it are the commands
of NORMAL state’s selecting, cursor, spelling and replacing keys, on the same
letters (\texttt{C-z m f} for \texttt{m f}), and the key hints list them.

\subsection{Key hints}
\label{sec:173}
Press a prefix key and wait a moment. Under the text, a panel opens that lists
every key that can come next with the name of its command. You do not need to
remember the keys of a family; you need only remember the first key. Figure
8.2 shows the panel after \texttt{C-x} in a buffer that visits a letter.

\begin{figure}[H]
\caption{Figure 8.2. The key hints after C-x. A name that starts with a plus sign is another prefix.}\label{org172}
\begin{Verbatim}[breaklines,breakautoindent=false,breakindent=1em,breakafter=/-.\,;:=,commandchars=\\\{\}]
 +-------------------------------------------------------------+
 | C-x                                                         |
 | C-f  find-file          b  switch-to-buffer    +bookmarks   |
 | C-s  save-buffer        k  kill-buffer         +windows     |
 | C-b  buffers            u  undo                +rectangles  |
 | C-c  save-buffers-kill-editor                       ... C-h |
 +-------------------------------------------------------------+
 | C-x-                                                        |
\end{Verbatim}
\end{figure}

The names that start with a plus sign are further prefixes, and pressing their
key opens another panel one level down. A key whose command cannot run in the
buffer you are in is dimmed, and pressing it tells you why. \texttt{C-g} closes the
panel. If the keys do not all fit, the last row ends in \ldots{} and says that
\texttt{C-h} lists them all.

\texttt{C-h} pressed right after any prefix key lists every key under it in a buffer
of its own, the prefixes unfolded, so that you can read them at leisure:
\texttt{C-x C-h} for the global keys, \texttt{C-c C-h} for those of the buffer you are in.
\texttt{q} closes that buffer.

A panel of hints takes as many rows as the frame has room for. If you want it
shorter, set \texttt{*key-\codebreak{}hint-\codebreak{}rows*} to the most rows it may take. The panel’s
colours are the faces below.

\begin{try}
\textbf{Try it, with the keys of a buffer you do not know}
\begin{enumerate}
\item Press \texttt{C-x C-b} to open the list of buffers. The cursor is in it.
\item Press \texttt{C-c} and wait a moment. The panel lists what the list can do: mark
   a buffer, mark it for deletion, open one and so on, each with its key.
\item Press \texttt{C-g}. The panel closes and nothing was run.
\item Press \texttt{C-c C-h}. The same keys are listed in a buffer of its own, in a
   window beside the list, and the cursor stays in the list. Press \texttt{C-x o} to
   move to it and \texttt{q} to close it, then \texttt{q} again to close the list of
   buffers.
\item Press \texttt{C-x} and wait. Press \texttt{w} to see what is under \texttt{C-x w}, then \texttt{C-g}.
\end{enumerate}
\end{try}

\begin{xltabular}{\textwidth}{@{}l l >{\raggedright\arraybackslash\hsize=1.000\hsize}X@{}}
\toprule
\textbf{Setting} & \textbf{Standard} & \textbf{What it changes}\\
\midrule
\endhead
\texttt{*key-\codebreak{}hint-\codebreak{}rows*} & \texttt{nil} & Most rows the key hints take after a prefix key (C-x, SPC), the echo row under them included, or NIL for as many as the frame has room for, as a menu takes; where they do not all fit, the last row ends in \ldots{}\\
\bottomrule
\end{xltabular}

\begin{xltabular}{\textwidth}{@{}l >{\raggedright\arraybackslash\hsize=1.000\hsize}X@{}}
\toprule
\textbf{Face} & \textbf{What it paints}\\
\midrule
\endhead
\texttt{:key-\codebreak{}hint-\codebreak{}key} & Keys in the panel that lists the bindings under a prefix key.\\
\texttt{:key-\codebreak{}hint-\codebreak{}command} & Commands in the key-hint panel shown after a prefix key.\\
\texttt{:key-\codebreak{}hint-\codebreak{}prefix} & Prefix keys (further keymaps) in the key-hint panel.\\
\texttt{:key-\codebreak{}hint-\codebreak{}unavailable} & A key in the key hints or a menu whose command cannot run here, such as one that would change a git repository looked at only: dimmed, and pressing it says why.\\
\bottomrule
\end{xltabular}

\subsection{Menus}
\label{sec:174}
Some commands have more to say than one key can. Git’s commit command has
switches, options and several ways to finish. For these Tuatara opens a menu:
a panel of the same kind as the key hints, but with a title, and with keys that
are not all commands. There are three kinds of entry. An action runs a command
and closes the menu. A switch turns something on or off, as in “sign this
commit”, and keeps the menu open. An option asks you for a value, such as a
number of commits, and keeps the menu open too. A switch or option that is on
is drawn in its own colour, and the rest in a fainter one, so that you can see
at a glance what the action will be given.

\texttt{C-h} in a menu lists every entry in a buffer of its own, with the command or
argument behind each, including the entries that are left out of the panel for
now. \texttt{C-g} leaves a menu. A menu that holds changes you made asks first, so
that you do not lose them. The git chapter shows menus in use, and the
commands in them can be run from \texttt{M-x} without the menu.

\begin{xltabular}{\textwidth}{@{}l >{\raggedright\arraybackslash\hsize=0.531\hsize}X l >{\raggedright\arraybackslash\hsize=1.469\hsize}X@{}}
\toprule
\textbf{Modal} & \textbf{Non-modal} & \textbf{Command} & \textbf{What it does}\\
\midrule
\endhead
none & \texttt{M-x menu-\codebreak{}toggle-\codebreak{}switch} & \texttt{menu-\codebreak{}toggle-\codebreak{}switch} & Turn the switch of the open menu whose key was typed on, or off, or say it is left out of the menu now; the menu stays open.\\
none & \texttt{M-x menu-\codebreak{}set-\codebreak{}option} & \texttt{menu-\codebreak{}set-\codebreak{}option} & Give the option of the open menu whose key was typed a value, asked for, or take its value away when it has one, or ask it again with that value when the option says so (MENU-ENTRY-REPLACE), or say it is left out of the menu now; the menu stays open.\\
none & \texttt{M-x menu-\codebreak{}leave} & \texttt{menu-\codebreak{}leave} & Leave the open menu, asked first when anything in it was changed; it stays open unless the answer is yes.\\
none & \texttt{M-x menu-\codebreak{}describe} & \texttt{menu-\codebreak{}describe} & List in \textbf{Help} every entry of the open menu: its keys, what it says and the command or argument behind it, those not shown now too.\\
\bottomrule
\end{xltabular}

\begin{xltabular}{\textwidth}{@{}l >{\raggedright\arraybackslash\hsize=1.000\hsize}X@{}}
\toprule
\textbf{Face} & \textbf{What it paints}\\
\midrule
\endhead
\texttt{:menu-\codebreak{}argument-on} & An argument a menu has on (a switch, an option given a value): what its actions are given.\\
\texttt{:menu-\codebreak{}argument-\codebreak{}off} & An argument a menu has off: shown, not given.\\
\bottomrule
\end{xltabular}

\subsection{Asking what a key does}
\label{sec:175}
You press a key by accident and something happens, or you see a key in a
book and want to know what it is. \texttt{C-h k} answers. It waits for a key or a
longer key and then shows, in a window of its own, the command the key runs,
what the command does and where the command is defined. You do not need to
leave the key you were using; the question is answered the moment the key is
complete.

\texttt{C-h k} takes a prefix as well. Press \texttt{C-h k}, then \texttt{C-x}, and it waits for the
rest, then shows \texttt{C-x C-f} when you press \texttt{C-f}. A key with no command says it
is undefined.

\texttt{?} in NORMAL state, or \texttt{C-h b} anywhere, lists every key that is in effect in
the buffer you are in. The list goes keymap by keymap, the one for this buffer
first and then those it falls through to. A key that a mode has taken away is
shown as \texttt{(none here)}. It is a long list, but a search with \texttt{C-s} finds a key
or a command in it.

\begin{try}
\textbf{Try it, when you forgot what a key does}
\begin{enumerate}
\item Press \texttt{C-h k}. The echo area asks \texttt{Describe key:}.
\item Press \texttt{g e} (non-modal: \texttt{M->}). A second window opens, and the cursor stays
   in the one you were in. The new window says that \texttt{g e} runs the command
   \texttt{end-\codebreak{}of-\codebreak{}buffer}, lists the keys that run it, and says what it does.
\item Press \texttt{C-x 1} (modal: \texttt{F}) to leave only one window. Press \texttt{C-h k} again, and
   this time press \texttt{C-x} and then \texttt{C-s}. It names \texttt{save-\codebreak{}buffer}.
\item Press \texttt{?} (non-modal: \texttt{C-h b}) to see every key in force. Press \texttt{C-x o} to
   move into the list, then \texttt{C-s}, type \texttt{save} and \texttt{RET} to find the keys for
   saving. Close the list with \texttt{q}.
\end{enumerate}
\end{try}

\begin{xltabular}{\textwidth}{@{}l l l >{\raggedright\arraybackslash\hsize=1.000\hsize}X@{}}
\toprule
\textbf{Modal} & \textbf{Non-modal} & \textbf{Command} & \textbf{What it does}\\
\midrule
\endhead
none & \texttt{C-h k} & \texttt{describe-\codebreak{}key} & Read a key sequence and show the command it runs and that command’s documentation, the question taken away once the keys are read.\\
\texttt{?} & \texttt{C-h b} & \texttt{describe-\codebreak{}bindings} & List every key binding in effect, keymap by keymap, each followed by the keymaps it falls through to.\\
\bottomrule
\end{xltabular}

\subsection{Asking what a command does}
\label{sec:176}
\texttt{C-h x} asks for a command by name, with the same list as \texttt{M-x}, and shows the
command: what it does, which keys run it and in which file it is defined. You
use it when \texttt{M-x} offered a command whose first line was not enough, or when
you want to see the keys that a command has so that you can learn them. Every
command is offered, including those that belong to other kinds of buffers. The
file in the last line is a link. Press \texttt{RET} on it and the source opens beside
the help, with the cursor on the definition. A release carries its Lisp
sources, so this works in the AppImage and the distribution packages too.

\texttt{C-h f} is the same for any function of Lisp, not just for the commands. It
shows the parameters as well. It is for people who write their own commands.
In a Lisp buffer it offers the name at the cursor first.

\texttt{C-h v} does the same for a variable. A setting is a variable, so this is how
you ask what a setting is for and what it is now. Press \texttt{C-h v}, type
\texttt{*fill-\codebreak{}column*} and \texttt{RET}, and the help gives the value, 80 as Tuatara comes,
and what it is for. A very long value, such as a list of hundreds of items, is
cut after a few thousand characters, and says so.

The help is shown in a window of its own, and you stay where you were, in the
text. If you would rather be taken to the help, so that you can scroll it,
set \texttt{*help-\codebreak{}window-\codebreak{}select*} to \texttt{t}. \texttt{q} then closes the window and returns you
to where you were. Either way \texttt{q} closes a help window when you are in it, and
\texttt{TAB} goes from one link to the next, \texttt{RET} following a link.

\begin{xltabular}{\textwidth}{@{}l l l >{\raggedright\arraybackslash\hsize=1.000\hsize}X@{}}
\toprule
\textbf{Modal} & \textbf{Non-modal} & \textbf{Command} & \textbf{What it does}\\
\midrule
\endhead
none & \texttt{C-h x} & \texttt{describe-\codebreak{}command} & Show what a command does, which keys run it and where it is defined, the file a link that shows the definition.\\
none & \texttt{C-h f} & \texttt{describe-\codebreak{}function} & Show what a function, macro or command does: its parameters, its documentation, the keys that run it when it is a command, and where it is defined, the file a link that shows the definition.\\
none & \texttt{C-h v} & \texttt{describe-\codebreak{}variable} & Show a variable’s value and documentation, and where it is defined, the file a link that shows the definition.\\
\bottomrule
\end{xltabular}

\begin{xltabular}{\textwidth}{@{}l l >{\raggedright\arraybackslash\hsize=1.000\hsize}X@{}}
\toprule
\textbf{Setting} & \textbf{Standard} & \textbf{What it changes}\\
\midrule
\endhead
\texttt{*help-\codebreak{}window-\codebreak{}select*} & \texttt{nil} & True to move to the window a help buffer is shown in (M-x describe-function, M-x describe-key and the rest); M-x quit-window there then closes it and goes back; NIL (the standard) leaves you where you were.\\
\bottomrule
\end{xltabular}

\begin{xltabular}{\textwidth}{@{}l >{\raggedright\arraybackslash\hsize=1.000\hsize}X@{}}
\toprule
\textbf{Face} & \textbf{What it paints}\\
\midrule
\endhead
\texttt{:help-\codebreak{}key-\codebreak{}binding} & Keys named in help, documentation and messages: those that run a command now.\\
\texttt{:button} & A button: text that does something when RET is pressed on it or it is clicked: the source file a command’s help names, say.\\
\bottomrule
\end{xltabular}

\subsection{Finding a command when you do not know its name}
\label{sec:177}
You know what you want, such as turning a line of text into capital letters,
but not what it is called. Here are three ways to look.

The first is \texttt{M-x} itself. Type a word that is likely to be in the name and
read the list. Words may come in any order, so \texttt{upper} and \texttt{case} together
find the commands with both. The list shows the first line of what each one
does, which is often enough.

The second is \texttt{C-h a}, apropos. It asks for words and searches both the names
and the documentation of every command and setting, which finds things that
\texttt{M-x} cannot, because the word you chose is in the help and not in the name.
The matches appear in a buffer called \texttt{*Apropos*}, the commands first and then
the settings, with the first sentence of each one’s help. Each name is a link:
\texttt{RET} on it shows its full help, and \texttt{TAB} goes to the next name.

The third is the keys. \texttt{?} (non-modal: \texttt{C-h b}) lists everything that has a key,
and \texttt{C-c} or \texttt{C-x} with a wait shows what a family offers. A command that has a
key has usually been put there because it is common.

\begin{try}
\textbf{Try it, to find how to make words capital}
\begin{enumerate}
\item Press \texttt{C-h a}, type \texttt{capital} and press \texttt{RET}. The \texttt{*Apropos*} buffer lists
   the commands and settings that mention capital letters.
\item Press \texttt{TAB} until a name you like is selected, and \texttt{RET} on it. Its help
   opens.
\item Close the help with \texttt{q}, and the list with \texttt{q} again.
\item Press \texttt{M-x} and type \texttt{upcase}. The list shows the commands that change
   to upper case, and the keys of those that have one.
\item Press \texttt{C-g}.
\end{enumerate}
\end{try}

\begin{xltabular}{\textwidth}{@{}l l l >{\raggedright\arraybackslash\hsize=1.000\hsize}X@{}}
\toprule
\textbf{Modal} & \textbf{Non-modal} & \textbf{Command} & \textbf{What it does}\\
\midrule
\endhead
none & \texttt{C-h a} & \texttt{apropos-\codebreak{}command} & List the commands and variables whose names or documentation match the words typed, in \textbf{Apropos}.\\
\bottomrule
\end{xltabular}

\subsection{What a face looks like}
\label{sec:178}
A face is a name for a way of drawing text: its colour, whether it is bold,
and so on. Chapter 1 explained them. \texttt{M-x describe-\codebreak{}face} asks for a face by
name and shows all that is known of it: a sample drawn in it, what it
colours, what the theme and your init file set for it, the face it inherits
from, and each attribute as it is drawn now. \texttt{M-x list-\codebreak{}faces} shows every face,
grouped by the first word of its name, each with a sample and the first
sentence of what it colours. \texttt{RET} on a face in the list shows its full
description.

\texttt{M-x describe-\codebreak{}char} does the same for a character. It shows its number in
Unicode, its name, the bytes it takes in the file, the face it is drawn in
and, in a language Tuatara can parse, the part of the code it is in. It is
how you find out which of two letters that look alike you have.

\begin{xltabular}{\textwidth}{@{}l l l >{\raggedright\arraybackslash\hsize=1.000\hsize}X@{}}
\toprule
\textbf{Modal} & \textbf{Non-modal} & \textbf{Command} & \textbf{What it does}\\
\midrule
\endhead
none & \texttt{M-x describe-\codebreak{}face} & \texttt{describe-\codebreak{}face} & Show all that is known of a face: a sample drawn in it, what it colours, its definition, what the theme in force and your init file set for it, the face it inherits from and each of its attributes as it is drawn now.\\
none & \texttt{M-x list-\codebreak{}faces} & \texttt{list-\codebreak{}faces} & List every face a theme can colour in \textbf{Faces}, in sections by the first word of their names: each face’s name, a sample drawn in it as the theme in force draws it, and the first sentence of what it colours.\\
none & \texttt{M-x describe-\codebreak{}char} & \texttt{describe-\codebreak{}char} & Show in \textbf{Help} all that is known of the character at point: where it is, its code in decimal, octal and hexadecimal, its Unicode name and kind, its octets and the columns it takes, the face it shows in and the major mode.\\
\bottomrule
\end{xltabular}

\subsection{Messages, notifications and the log of what Tuatara says}
\label{sec:179}
Tuatara says short things in the echo area, one at a time, and they are gone
after a moment. \texttt{C-h e} shows them all in \texttt{*Messages*}, the latest at the end,
with repeats counted. It is the first place to look when something went by too
fast to read. The buffer keeps the last \texttt{*message-\codebreak{}log-\codebreak{}max*} lines.

\texttt{C-h N} shows \texttt{*Notifications*}: what was told to you during this session by
the agenda, when an entry marked for a reminder came due, and by packages.
Each one is a line with the date and time. Each notification is also sent to the
desktop by the program that \texttt{*notify-\codebreak{}program*} names, \texttt{notify-\codebreak{}send} as Tuatara
comes. Set it to \texttt{nil} to keep them in the buffer only.

\begin{xltabular}{\textwidth}{@{}l l l >{\raggedright\arraybackslash\hsize=1.000\hsize}X@{}}
\toprule
\textbf{Modal} & \textbf{Non-modal} & \textbf{Command} & \textbf{What it does}\\
\midrule
\endhead
none & \texttt{C-h e} & \texttt{view-\codebreak{}echo-\codebreak{}area-\codebreak{}messages} & Show the \textbf{Messages} buffer: the messages shown in the echo area so far, repeats counted, the latest at the end.\\
none & \texttt{C-h N} & \texttt{view-\codebreak{}notifications} & Show \textbf{Notifications}, what was told so far this session, the agenda’s entries tagged :notify: as their time came and what packages told, a line each with the date and time it was told and who told it, the latest at the end.\\
\bottomrule
\end{xltabular}

\begin{xltabular}{\textwidth}{@{}l l >{\raggedright\arraybackslash\hsize=1.000\hsize}X@{}}
\toprule
\textbf{Setting} & \textbf{Standard} & \textbf{What it changes}\\
\midrule
\endhead
\texttt{*message-\codebreak{}log-\codebreak{}max*} & \texttt{1000} & Lines the \textbf{Messages} buffer keeps of the messages shown, the oldest going first, 1000 as standard; NIL keeps none.\\
\texttt{*notify-\codebreak{}program*} & \texttt{"notify-\codebreak{}send"} & The program that shows a notification on the desktop: a name looked for on the PATH, or a path.\\
\bottomrule
\end{xltabular}

\subsection{The tutor and manual pages}
\label{sec:180}
\texttt{M-x tuatara-\codebreak{}tutor} opens the tutor, a buffer that teaches the modal keys by
making you use them. It is an ordinary buffer with its own instructions. Each
lesson tells you what to type and gives a line to practise on, marked with an
arrow, and the keys it names are the keys you really have, even if you moved
them. You edit the tutor as it says, and nothing you do in it is saved. Half
an hour with it is the quickest way to get used to the two states. The page
\texttt{*About Tuatara*} below has a button that opens it.

The system’s own manual pages, the ones that explain programs such as \texttt{ls} or
\texttt{grep} in a terminal, can be read in Tuatara. \texttt{M-x man} asks for a name, as in
\texttt{ls}, \texttt{ls(1)} or \texttt{1 ls}, and lists the names of the pages the system has as you
type. The page opens read-only in a window of its own, laid out for the width of
the window, with its headings, bold text and underlined text drawn in faces.
Every mention of another page, such as \texttt{chmod(1)}, is a link. \texttt{RET} or \texttt{l}
follows it, \texttt{TAB} goes to the next link, and \texttt{h} comes back to the page you
came from. \texttt{n} and \texttt{p} go to the next and the previous section, \texttt{r} lays the
page out again if you changed the size of the window, and \texttt{q} closes it.

If you open the source of a page, a file such as \texttt{tool.1}, \texttt{M-x man} offers
that file first, so that you can see what you wrote as it will look.

The page is made by the system’s \texttt{man} and \texttt{groff}, run with no network and
none of your files, and it is in English. If they are not installed,
\texttt{M-x check-\codebreak{}health} says so.

\begin{xltabular}{\textwidth}{@{}l l l >{\raggedright\arraybackslash\hsize=1.000\hsize}X@{}}
\toprule
\textbf{Modal} & \textbf{Non-modal} & \textbf{Command} & \textbf{What it does}\\
\midrule
\endhead
none & \texttt{M-x tuatara-\codebreak{}tutor} & \texttt{tuatara-\codebreak{}tutor} & Open the tutor: an ordinary buffer holding its own lessons, to learn the modal keys by editing it as it says, with every key the one you have.\\
none & \texttt{M-x man} & \texttt{man} & Show a manual page in a buffer of its own, read-only: ask for it by name (ls), with its section (ls(1)) or by section and name (1 ls), completing on the names of the pages the system has.\\
none & \texttt{n}, \texttt{C-c n} & \texttt{man-\codebreak{}next-\codebreak{}section} & Go to the start of the next section heading of the manual page.\\
none & \texttt{p}, \texttt{C-c p} & \texttt{man-\codebreak{}previous-\codebreak{}section} & Go to the start of the section heading before the line point is on.\\
none & \texttt{h}, \texttt{C-c h} & \texttt{man-\codebreak{}back} & Show the page a link was followed from; close the window when this page was not reached by a link.\\
\bottomrule
\end{xltabular}

\subsection{About Tuatara and what is there}
\label{sec:181}
\texttt{C-h C-a} opens =*About Tuatara*, the page that introduces the editor. It says
which version you have, who wrote it and under which licence, where its ideas
came from, where the sources are and where to report a fault, and it has
buttons for the ways in: the tutor, the help keys of this chapter, your init
file and the health check. \texttt{TAB} and \texttt{S-TAB} go from one button to the next,
\texttt{RET} presses one, and a click does too. Opening a web address from it asks
first and uses your desktop’s browser.

\texttt{C-h C-c} shows the licence, the GNU General Public License, version 3 or
later, and \texttt{C-h C-w} the part of it that says there is no warranty.
\texttt{M-x credits} shows the projects that Tuatara owes its ideas to and what it
took from each. \texttt{M-x tuatara-\codebreak{}version} says which Tuatara this is: the version,
the commit and date it was built from, the Steel Bank Common Lisp (SBCL) it
runs on and whether it shows in a window or a terminal. With \texttt{C-u} it writes
the line at the cursor instead of saying it, which is useful in a report of a
fault.

\begin{xltabular}{\textwidth}{@{}l l l >{\raggedright\arraybackslash\hsize=1.000\hsize}X@{}}
\toprule
\textbf{Modal} & \textbf{Non-modal} & \textbf{Command} & \textbf{What it does}\\
\midrule
\endhead
none & \texttt{C-h C-a} & \texttt{about-\codebreak{}tuatara} & Show \textbf{About Tuatara}, the page that introduces Tuatara.\\
none & \texttt{C-h C-c} & \texttt{describe-\codebreak{}copying} & Show the license Tuatara comes under.\\
none & \texttt{C-h C-w} & \texttt{describe-\codebreak{}no-\codebreak{}warranty} & Show the part of the license that says Tuatara comes with no warranty.\\
none & \texttt{M-x credits} & \texttt{credits} & Show where Tuatara’s ideas come from: the projects it owes, and what it took from each.\\
none & \texttt{M-x tuatara-\codebreak{}version} & \texttt{tuatara-\codebreak{}version} & Say which Tuatara this is: its version, the commit it was built from and when, the Steel Bank Common Lisp (SBCL) it runs on, and where it shows (a window, the terminal).\\
\bottomrule
\end{xltabular}

\subsection{Checking what your system has}
\label{sec:182}
Tuatara uses a few programs outside itself: for example, the sandbox that
protects it when it opens a file you do not trust, the programs that read PDFs
and Office files, and GnuPG for encryption. It also loads a grammar for each
programming language it colours and indents. \texttt{M-x check-\codebreak{}health} reports what
is there and what is not. It runs in the background, because some of it takes
a moment, and shows \texttt{*Health*} when it is ready. The top of the report says
where your init file is and what the sandbox and the programs need. A cross
beside a program the editor needs is a problem, and the message says how to
get it. A fainter cross beside a program that only adds a feature says what the
feature needs, and is nothing to worry about. Further down, a table has a row
for each language, with a tick when its colours, indentation and selections
work, and then the major modes and the files they take. In a window, the
report also measures how quickly each way of drawing is on your machine and
which is faster.

In the report \texttt{r} checks again and \texttt{q} closes it. Run it first when something
does not work, and send it with a report of a fault.

If a command fails and you want to see why, \texttt{M-x debug-\codebreak{}mode} keeps a record of
what the editor does: commands with their keys, background jobs, files parsed
and errors. \texttt{C-h E} shows it in \texttt{*Debug*}. The same lines go to a file that
survives a crash, so that it can be read afterwards. \texttt{Debug} in the mode line
says it is on, and \texttt{M-x debug-\codebreak{}mode} again turns it off. \texttt{*debug-\codebreak{}mode*} set to
\texttt{t} turns it on when Tuatara starts, and \texttt{*debug-\codebreak{}log-\codebreak{}limit*} is how many lines
\texttt{*Debug*} keeps.

\begin{xltabular}{\textwidth}{@{}l l l >{\raggedright\arraybackslash\hsize=1.000\hsize}X@{}}
\toprule
\textbf{Modal} & \textbf{Non-modal} & \textbf{Command} & \textbf{What it does}\\
\midrule
\endhead
none & \texttt{M-x check-\codebreak{}health} & \texttt{check-\codebreak{}health} & Show what each language and major mode has here.\\
none & \texttt{M-x debug-\codebreak{}mode} & \texttt{debug-\codebreak{}mode} & Switch debug mode on or off.\\
none & \texttt{C-h E} & \texttt{view-\codebreak{}debug-\codebreak{}buffer} & Show \textbf{Debug}, the reports of debug mode (M-x debug-mode), where \textbf{Messages} is shown, its end in sight.\\
\bottomrule
\end{xltabular}

\begin{xltabular}{\textwidth}{@{}l l >{\raggedright\arraybackslash\hsize=1.000\hsize}X@{}}
\toprule
\textbf{Setting} & \textbf{Standard} & \textbf{What it changes}\\
\midrule
\endhead
\texttt{*debug-\codebreak{}mode*} & \texttt{nil} & True to start the editor with debug mode on, as M-x debug-mode switches it: what the editor’s parts do is reported in \textbf{Debug}.\\
\texttt{*debug-\codebreak{}log-\codebreak{}limit*} & \texttt{10000} & The lines \textbf{Debug} keeps, the oldest going first, 10000 as standard; NIL keeps them all, and the buffer then grows for as long as debug mode is on.\\
\bottomrule
\end{xltabular}

\subsection{Lisp in the running editor}
\label{sec:183}
Tuatara is written in Common Lisp, and the Lisp you type runs inside the
editor. This is how an init file works. It is also a calculator, a way to look
at a value and a way to try something small. You do not need Lisp for any other
chapter, so you may skip this part.

\texttt{M-:} asks for one Lisp expression in the minibuffer and shows its value, in
any buffer. \texttt{M-: (+ 2 3) RET} shows \texttt{5}. In a Lisp buffer, such as
\texttt{*lisp-\codebreak{}scratch*} (chapter 7) or a file that ends in \texttt{.lisp}, \texttt{C-x C-e} runs the
expression before the cursor, \texttt{C-j} puts what it printed on the lines after it,
and \texttt{C-M-x} runs the whole top-level form the cursor is in. \texttt{M-x eval-\codebreak{}region}
runs the selection, and \texttt{M-x eval-\codebreak{}buffer} runs the whole buffer. They run only
in Lisp buffers, and run nothing in a file opened with care because you do
not trust it.

\texttt{M-x ielm} opens \texttt{*ielm*}, a buffer with a prompt where you type Lisp and
read the answer below. \texttt{RET} runs the input when it is whole expressions and
starts a new line when it is not. \texttt{M-p} and \texttt{M-n} bring back earlier inputs, and
\texttt{C-c C-u} clears what you typed. The modal keys are off in it, so you type at
once without \texttt{i}.

If Lisp meets an error, the echo area says what it was. \texttt{M-x
toggle-\codebreak{}debug-\codebreak{}on-\codebreak{}error} makes it show the calls under way too, in
\texttt{*Backtrace*}. That is for people who write Lisp, and is easy to leave off.

\begin{try}
\textbf{Try it, with a sum}
\begin{enumerate}
\item Press \texttt{M-:}, type \texttt{(* 12 7)} and press \texttt{RET}. The echo area shows
   \texttt{(* 12 7) => 84}, then the same number in other bases.
\item Press \texttt{M-x}, type \texttt{ielm} and press \texttt{RET}. After the prompt type
   \texttt{(reverse "kiwi")} and press \texttt{RET}. The answer is \texttt{"iwik"}.
\item Press \texttt{M-p}. The input comes back. Press \texttt{C-c C-u} to take it away.
\item Press \texttt{C-x k} and \texttt{RET} to close \texttt{*ielm*}. It has changed, so it asks;
   type \texttt{y} and press \texttt{RET}.
\end{enumerate}
\end{try}

\begin{xltabular}{\textwidth}{@{}l >{\raggedright\arraybackslash\hsize=0.585\hsize}X l >{\raggedright\arraybackslash\hsize=1.415\hsize}X@{}}
\toprule
\textbf{Modal} & \textbf{Non-modal} & \textbf{Command} & \textbf{What it does}\\
\midrule
\endhead
none & \texttt{M-:} & \texttt{eval-\codebreak{}expression} & Read a Lisp expression in the minibuffer and evaluate it in the running editor, in the current buffer’s package; show its value, or with a prefix argument (C-u) insert it at point.\\
none & \texttt{C-x C-e} & \texttt{eval-\codebreak{}last-\codebreak{}expression} & Evaluate the expression before point in the running editor and show its value in the echo area; with a prefix argument (C-u), insert the value after point.\\
none & \texttt{C-j} & \texttt{eval-\codebreak{}print-\codebreak{}last-\codebreak{}expression} & Evaluate the expression before point in the running editor and put what it printed and its value on the lines after it.\\
none & \texttt{C-M-x} & \texttt{eval-\codebreak{}top-\codebreak{}level-\codebreak{}form} & Evaluate the top-level form around point (the one starting with an open bracket at the start of a line) in the running editor and show its value.\\
none & \texttt{M-x eval-\codebreak{}region} & \texttt{eval-\codebreak{}region} & Evaluate the Lisp forms in the selection one after another, as a file’s are (IN-PACKAGE carries on to the next), and show the last one’s value.\\
none & \texttt{M-x eval-\codebreak{}buffer} & \texttt{eval-\codebreak{}buffer} & Evaluate every Lisp form in the current buffer one after another, as a file’s are, and show the last one’s value; not a buffer of more than \textbf{REWRITE-THRESHOLD} octets.\\
none & \texttt{M-x ielm} & \texttt{ielm} & Show \textbf{ielm}, where Lisp typed after the prompt is evaluated in the running editor: made with a prompt the first time, as it was after.\\
none & \texttt{RET} & \texttt{ielm-\codebreak{}return} & RET in \textbf{ielm}: after the last prompt, evaluate the input when it is whole expressions, else begin a new line of it; on an earlier input, evaluate that again.\\
none & \texttt{C-j} & \texttt{ielm-\codebreak{}send-\codebreak{}input} & Evaluate the input after \textbf{ielm}’s last prompt as it is, whole or not.\\
none & \texttt{M-p} & \texttt{ielm-\codebreak{}previous-\codebreak{}input} & Put the input before the one shown in \textbf{ielm}’s history in place of the input.\\
none & \texttt{M-n} & \texttt{ielm-\codebreak{}next-\codebreak{}input} & Put the input after the one shown in \textbf{ielm}’s history in place of the input.\\
none & \texttt{C-c C-u} & \texttt{ielm-\codebreak{}kill-\codebreak{}input} & Take away the input after \textbf{ielm}’s last prompt.\\
none & \texttt{C-c C-b} & \texttt{ielm-\codebreak{}change-\codebreak{}working-\codebreak{}buffer} & Choose the working buffer of \textbf{ielm}: the one the Lisp typed there sees as current, (insert “x”) putting text in it.\\
none & \texttt{TAB} & \texttt{ielm-\codebreak{}tab} & TAB in \textbf{ielm}: where only blanks come before point on its line of the input, indent the line as Lisp (never the prompt’s line); else complete the name at point from a list beside it, as C-M-i does.\\
none & \texttt{M-x toggle-\codebreak{}debug-\codebreak{}on-\codebreak{}error} & \texttt{toggle-\codebreak{}debug-\codebreak{}on-\codebreak{}error} & Switch debug-on-error on or off: while it is on, an error in a command, or in Lisp evaluated in the editor (C-x C-e, C-M-x, M-:, C-j), shows the calls under way when it was signalled in \textbf{Backtrace}, innermost first, besides the message.\\
\bottomrule
\end{xltabular}

\begin{xltabular}{\textwidth}{@{}l l >{\raggedright\arraybackslash\hsize=1.000\hsize}X@{}}
\toprule
\textbf{Setting} & \textbf{Standard} & \textbf{What it changes}\\
\midrule
\endhead
\texttt{*debug-\codebreak{}on-\codebreak{}error*} & \texttt{nil} & True to show the calls under way when a command or an evaluation meets an error, in \textbf{Backtrace}; M-x toggle-debug-on-error switches it.\\
\bottomrule
\end{xltabular}

\subsection{Binding keys yourself}
\label{sec:184}
The prefixes, the leader and the labels that the key hints show are yours to
set. You can bind a key to any command, and give it a label that the hints
show in place of the command’s name, in your init file. That is described in
chapter 20, together with the settings that move the leader.

\begin{xltabular}{\textwidth}{@{}l l >{\raggedright\arraybackslash\hsize=1.000\hsize}X@{}}
\toprule
\textbf{Setting} & \textbf{Standard} & \textbf{What it changes}\\
\midrule
\endhead
\texttt{*leader-\codebreak{}key*} & \texttt{"SPC"} & The key that opens the leader’s keys (bind-key’s :map :leader) in NORMAL and support states, in lists and in the calendar: SPC, or another written as bind-key writes keys (“,”, “C-c SPC”).\\
\texttt{*modal-\codebreak{}editing*} & \texttt{t} & True for Tuatara’s modal editing: NORMAL and INSERT states in editing buffers, NORMAL’s reading keys in support buffers, the leader on SPC.\\
\bottomrule
\end{xltabular}

\section{Files}
\label{sec:216}
Opening, saving and keeping track of the files on your disk, and on other
machines.

Chapter 1 explained that Tuatara shows a file in a buffer: a copy of the text
that you change freely, kept apart from the file until you save. This chapter
is about the file end of that. It shows how to open a file, how to make a new
one, and what happens when you save. It explains what to do when the file has
changed on disk behind your back, and how to give a file a new name. It also
covers the way Tuatara handles line endings and letters from other languages,
the file names that say what a file holds, and files that live on another
machine. It ends with the ways of opening a file from outside Tuatara: with
another program, and from the command line.

The folders that hold files are the subject of chapter 10, and very large
files of chapter 11. Files from a source you do not trust are in chapter 12.
The list of recent files, bookmarks and auto-save are in chapter 7.

Every key in this chapter is pressed in NORMAL state unless it starts with
\texttt{C-x}, which works in both states. \texttt{i} starts INSERT state, and \texttt{C-g} or \texttt{ESC}
brings you back.

\subsection{Opening a file}
\label{sec:189}
\texttt{C-x C-f} asks for a file name at the bottom of the screen. The question has a
folder filled in, and a list of what is in it. Type to narrow the list, and
press \texttt{RET} to open what you have chosen. If the name is a folder, you go into
it and the list shows what it holds. The question starts in the folder of the
last file or folder you chose, so when you work on one project you seldom type
more than the end of a name. \texttt{S-TAB} puts your home folder, \texttt{\textasciitilde{}/}, in place of
what is typed, and lists it. Typing \texttt{\textasciitilde{}/} after a folder does the same, starting
the path again from home.

Beside each file the list shows its permissions, its size and when it last
changed. Where a file is one you opened safely before, the list marks it, and
a buffer with changes you have not saved is marked as well.

Opening a file you already have open does not make a second buffer. It takes
you to the first. This holds for a file reached by two names, such as a link
and the file it leads to.

If the name you typed is not there, you get a new, empty file with that name.
Nothing is written yet. The next section says more.

Opening is quick even when the file is large, because Tuatara reads files in
the background. A very large file is opened in a special way, described in
chapter 11.

\subsubsection{Other places to open a file}
\label{sec:186}
\texttt{C-x 4 f} opens the file in another window and goes there, so you can read two
files side by side. \texttt{C-x 5 f} opens it in another frame, a window of its own,
which chapter 1 described. \texttt{C-x C-p} opens the file whose name is written in
the text, such as \texttt{\textasciitilde{}/letters/\codebreak{}council.\codebreak{}txt} in a note. A folder written there
opens as a list of its files.

When Tuatara runs in a window of its own, you can also drag files from your
file manager and drop them on it. Each one opens in the window it was dropped
on. If you want these opened safely, which is the careful way described in
chapter 12, set \texttt{*open-\codebreak{}dropped-\codebreak{}files-\codebreak{}safely*}. With \texttt{C-u} before \texttt{C-x C-f} a
file opens safely too. \texttt{M-x find-\codebreak{}file-\codebreak{}safely} does the same.

\texttt{M-x show-\codebreak{}file-\codebreak{}name} says in the echo area which file the buffer you are in
shows, with its folder. You use it when two files have the same name and you
want to be sure which one this is.

\texttt{C-x C-f} and the other ways above take the file name as you type it. A name
starting with \texttt{/ssh:} is a file on another machine, described later in this
chapter.

\subsubsection{The file’s own settings}
\label{sec:187}
A file can say how it should be shown. The first line of a program, for
instance, may hold \texttt{-*- mode: org -*-} to say that it is an Org file, or a
block of lines near its end may set the width of a tab. Tuatara reads those
settings when it opens the file, unless you turn off \texttt{*file-\codebreak{}local-\codebreak{}variables*}.
You will meet this in files written for Emacs. It is of no use to you until
you see it.

\begin{xltabular}{\textwidth}{@{}l >{\raggedright\arraybackslash\hsize=0.578\hsize}X l >{\raggedright\arraybackslash\hsize=1.422\hsize}X@{}}
\toprule
\textbf{Modal} & \textbf{Non-modal} & \textbf{Command} & \textbf{What it does}\\
\midrule
\endhead
none & \texttt{C-x C-f} & \texttt{find-\codebreak{}file} & Visit a file, completing its name against the folder being typed.\\
none & \texttt{C-x 4 f}, \texttt{C-x 4 C-f} & \texttt{find-\codebreak{}file-\codebreak{}other-\codebreak{}window} & Visit a file in another window and select it, completing its name as find-file does.\\
none & \texttt{C-x 5 f} & \texttt{find-\codebreak{}file-\codebreak{}other-\codebreak{}frame} & Visit a file in another frame and work in it, asking for its name as C-x C-f does; a folder is listed there.\\
none & \texttt{M-x find-\codebreak{}file-\codebreak{}safely} & \texttt{find-\codebreak{}file-\codebreak{}safely} & Open a file safely, for looking at a file from a source you do not trust: read-only, with no major mode, no hooks and no version control, so that nothing in the file is acted on; its file system left exactly as it was (no copy beside it), and its access time too where the system allows.\\
none & \texttt{C-x C-p} & \texttt{open-\codebreak{}path-\codebreak{}at-\codebreak{}point} & Open the file or folder whose path is at the cursor in the text, as C-x C-f opens it: a folder (\emph{etc}, \textasciitilde{}/notes/) in a list of it, a file \textbackslash{}(\textasciitilde{}/notes/todo.org, ./README.org, or a file: link’s) in its buffer.\\
none & \texttt{M-x open-\codebreak{}dropped-\codebreak{}files} & \texttt{open-\codebreak{}dropped-\codebreak{}files} & Visit the files dropped on the window, in the window they were dropped on, as C-x C-f does, the last one shown; safely, as C-u C-x C-f does, when \textbf{OPEN-DROPPED-FILES-SAFELY} is true.\\
none & \texttt{M-x show-\codebreak{}file-\codebreak{}name} & \texttt{show-\codebreak{}file-\codebreak{}name} & Say the full name of the file the current buffer shows, its directory and name, in the echo area, where \textbf{Messages} keeps it too: in a rimu list, that of the entry on the current line; in a buffer with no file, the directory it stands in.\\
\bottomrule
\end{xltabular}

\begin{xltabular}{\textwidth}{@{}l l >{\raggedright\arraybackslash\hsize=1.000\hsize}X@{}}
\toprule
\textbf{Setting} & \textbf{Standard} & \textbf{What it changes}\\
\midrule
\endhead
\texttt{*last-\codebreak{}used-\codebreak{}first*} & \texttt{t} & True for the questions that choose a file or a folder (M-x find-file, M-x rimu and the rest) to start in the folder of the path chosen last in any of them, and to list first, in each folder, what was chosen there last (or something in it), TAB then going the way last gone; the others follow in order of name.\\
\texttt{*open-\codebreak{}dropped-\codebreak{}files-\codebreak{}safely*} & \texttt{nil} & True to open files dropped on the window safely, as M-x universal-argument M-x find-file does: read-only, no mode, nothing written beside them.\\
\texttt{*file-\codebreak{}local-\codebreak{}variables*} & \texttt{t} & True for a file’s own settings (-\textbf{- mode: org -}- on its first line, a Local Variables block near its end) to be taken as it is opened, namely its major mode, tab width, tabs or spaces, fill column and spaces between sentences, and those a package takes; nothing else.\\
\texttt{*preview-\codebreak{}files*} & \texttt{nil} & True to show, while you choose a file with M-x find-file (and the other ways of finding a file) or M-x find-recent-file, the file selected in the window you are in, opened safely: read-only, no mode or hook run on it, its times untouched, a binary one in the hex view; what the window showed comes back when the question ends.\\
\bottomrule
\end{xltabular}

\begin{xltabular}{\textwidth}{@{}l >{\raggedright\arraybackslash\hsize=1.000\hsize}X@{}}
\toprule
\textbf{Face} & \textbf{What it paints}\\
\midrule
\endhead
\texttt{:completion-\codebreak{}modified} & The mark beside a buffer among completion candidates that has changes not saved.\\
\texttt{:completion-\codebreak{}safe} & The mark beside a file among completion candidates that was opened safely, and opens so again.\\
\texttt{:completion-\codebreak{}mode-\codebreak{}read} & The permission to read (r) beside a file among completion candidates.\\
\texttt{:completion-\codebreak{}mode-\codebreak{}write} & The permission to write (w) beside a file among completion candidates.\\
\texttt{:completion-\codebreak{}mode-\codebreak{}exec} & The permission to run or enter (x) beside a file among completion candidates.\\
\texttt{:completion-\codebreak{}mode-\codebreak{}none} & A permission a file does not give (-), beside it among completion candidates.\\
\texttt{:completion-\codebreak{}mode-\codebreak{}other} & Other letters of a file’s permissions (s, t, and the kinds of special files) beside it among completion candidates.\\
\texttt{:list-\codebreak{}date} & Dates and times in a list or beside completion candidates: when a file last changed, the date in a formatted name, a rimu list’s last commits (G), when a message was sent.\\
\texttt{:list-\codebreak{}size} & Sizes in a list or beside completion candidates: of files and buffers, and a rimu list’s sizes column (S).\\
\bottomrule
\end{xltabular}

\subsubsection{Seeing a file before you choose it}
\label{sec:188}
With \texttt{*preview-\codebreak{}files*} set to \texttt{t}, the question for a file shows the file under
the cursor in the window beside it while you move through the list. The file
is opened safely, so nothing in it is acted on, and the preview goes away when
you answer. It helps when you have several files with names like
\texttt{letter-\codebreak{}2.\codebreak{}txt} and \texttt{letter-\codebreak{}3.\codebreak{}txt}. It is off until you turn it on.

\begin{try}
\textbf{Try it, with a letter}
\begin{enumerate}
\item Press \texttt{C-x C-f}. The question shows a folder. Type \texttt{\textasciitilde{}/letters/} and the
   list shows what is in the folder, or nothing if it is not there yet.
\item Type \texttt{council.\codebreak{}txt} and press \texttt{RET}. The buffer is empty and the echo area
   says it is a new file. If \texttt{\textasciitilde{}/letters/} does not exist yet, it says that its
   folder is not there and that saving asks to make it.
\item Press \texttt{i}, type \texttt{Dear Mr Taupo,} and press \texttt{C-g}. The mode line shows \texttt{[*]},
   which means changes not saved.
\item Leave the buffer as it is. You will save it in the next section.
\end{enumerate}
\end{try}

\subsection{Creating a file}
\label{sec:190}
You do not make a file by a separate command. You open one with a name that is
not there. The buffer is empty and nothing exists on disk until you save. That
suits a letter or a note: open the name, type, save, or throw the buffer away.

If the folder you named does not exist either, the echo area says so when you
open the file, and saving asks \texttt{The folder ... is not there; make it?} with a
single key. \texttt{y} makes the folder and saves. \texttt{n} stops and writes nothing. You
can name a file three folders deep in one go this way. If Tuatara cannot make
the folder, because you may not write where it would go, the buffer is
read-only and the message says why.

A name that says a file is encrypted but is not encrypted by GnuPG, for
instance one ending \texttt{.age}, is refused when you save it, because the text
would be there to read. Names that end \texttt{.gpg} are encrypted every time they
are saved. 
The file gets the permissions your account normally gives a new file. A
folder listing, chapter 10, can show a file’s permissions.

\subsection{Saving}
\label{sec:193}
\texttt{C-x C-s} saves the buffer to its file. The echo area says \texttt{Wrote} and the
file’s name, and the mode line goes from \texttt{[*]} to \texttt{[-]}. If there is nothing to
save, it says \texttt{(No changes need to be saved)} and writes nothing. A buffer with
no file asks where to put it, in the same way as \texttt{C-x C-f} asks.

Chapter 7 described \texttt{C-x s}, which saves every file with changes after one
question, and \texttt{C-x C-c}, which asks before it leaves.

\subsubsection{What a save does with your file}
\label{sec:191}
Tuatara never edits the file in place. A save writes the whole text to a new
file beside the old one. It gives the new file the owner and the permissions of
the old one, and makes sure the text is on the disk, then renames the new file
over the old one. The file on disk is always either the old one or the new
one, even after a power cut, and a save that fails halfway leaves nothing
behind.

There is no backup copy, so you will find no \texttt{letter.\codebreak{}txt\textasciitilde{}} next to
\texttt{letter.\codebreak{}txt}. The safety net for a crash is the auto-saved copy of chapter 7,
and the safety net for a mistake is undo. A rewritten file is easier to see
than a backup would be, since Tuatara asks before it writes over a file that
changed behind its back.

A few files cannot be replaced by a new file without losing something. These
are files with more than one name (hard links), files whose owner or group your
account cannot give to a new file, and a few special files that the system
makes up. Tuatara writes these in place, after copying the old text to a spare
copy of its own. If the write fails, the spare copy is kept and its name is
said.

\subsubsection{When a save asks first}
\label{sec:192}
A save asks a question before it does something you may not want. It asks if:

\begin{itemize}
\item the file was changed on disk after you opened it, as \texttt{council.\codebreak{}txt changed on
  disk since it was read; overwrite it with this buffer?}. Say \texttt{n} and read the
  other version first, as described later in this chapter;
\item the file is write-protected: \texttt{File ... is write-\codebreak{}protected; try to save
  anyway?};
\item you save as a name that is already a file: \texttt{... exists; overwrite it?};
\item you save as a name that another buffer has open. Tuatara offers to write over
  it and kill that buffer.
\end{itemize}

A question that cannot be undone, such as writing over a file that exists, is
answered by typing \texttt{yes} or \texttt{no} and pressing \texttt{RET}. This is the same as for
leaving with changes unsaved, described in chapter 1.

A large buffer takes time to write. If the save takes longer than half a second,
Tuatara carries on in the background and the mode line shows \texttt{Saving} and a
percentage. You can keep typing. Leaving Tuatara waits for saves still running.
\texttt{*save-\codebreak{}in-\codebreak{}place*} is for very large files, which chapter 11 explains.

\begin{xltabular}{\textwidth}{@{}l l l >{\raggedright\arraybackslash\hsize=1.000\hsize}X@{}}
\toprule
\textbf{Modal} & \textbf{Non-modal} & \textbf{Command} & \textbf{What it does}\\
\midrule
\endhead
none & \texttt{C-x C-s} & \texttt{save-\codebreak{}buffer} & Save the current buffer to its file, asking for one if it has none, as C-x C-f asks.\\
\bottomrule
\end{xltabular}

\begin{xltabular}{\textwidth}{@{}l l >{\raggedright\arraybackslash\hsize=1.000\hsize}X@{}}
\toprule
\textbf{Setting} & \textbf{Standard} & \textbf{What it changes}\\
\midrule
\endhead
\texttt{*save-\codebreak{}in-\codebreak{}place*} & \texttt{nil} & True to save a large file by writing only what changed into the file itself, where that is quicker than writing it whole.\\
\bottomrule
\end{xltabular}

\subsection{Save as, and a new name}
\label{sec:194}
\texttt{C-x C-w} writes the buffer to another file. The question works like the one
for \texttt{C-x C-f}, and starts in the folder of the file you are in, with the name
to change. After the save, the buffer is the new file. The old file stays as
it was, which makes \texttt{C-x C-w} the way to copy a file: open it, press \texttt{C-x C-w},
and give the copy a name.

\texttt{M-x rename-\codebreak{}visited-\codebreak{}file} gives the file itself a new name, and the buffer
follows it. The question offers the old name to edit. A name alone keeps the
file in its folder. A path, such as \texttt{\textasciitilde{}/letters/\codebreak{}sent/\codebreak{}council.\codebreak{}txt}, moves it
there, provided that is on the same drive. Use this when you named a file in a
hurry. The file on another machine is the exception: this command leaves
that file where it is.

\begin{try}
\textbf{Try it, with a letter}
\begin{enumerate}
\item In the letter from the previous try, press \texttt{C-x C-s}. If \texttt{\textasciitilde{}/letters/} is
   not there yet, the echo area asks whether to make it: press \texttt{y}. It then
   says \texttt{Wrote \textasciitilde{}/letters/\codebreak{}council.\codebreak{}txt}, and the mode line shows \texttt{[-]}.
\item Press \texttt{C-x C-s} again. It says \texttt{(No changes need to be saved)}.
\item Press \texttt{C-x C-w}. The prompt offers the folder. Type
   \texttt{\textasciitilde{}/letters/\codebreak{}thanks.\codebreak{}txt}, which takes its place, and press \texttt{RET}. The buffer
   is \texttt{thanks.\codebreak{}txt} now, and \texttt{council.\codebreak{}txt} is as you saved it.
\item Run \texttt{M-x rename-\codebreak{}visited-\codebreak{}file}. The old name is offered. Change it to
   \texttt{thanks-\codebreak{}to-\codebreak{}aroha.\codebreak{}txt} and press \texttt{RET}. The echo area says it was renamed,
   and \texttt{thanks.\codebreak{}txt} is gone from the folder.
\item Press \texttt{C-x C-b}. The buffer list shows the new name and its path.
\end{enumerate}
\end{try}

\begin{xltabular}{\textwidth}{@{}l >{\raggedright\arraybackslash\hsize=0.542\hsize}X l >{\raggedright\arraybackslash\hsize=1.458\hsize}X@{}}
\toprule
\textbf{Modal} & \textbf{Non-modal} & \textbf{Command} & \textbf{What it does}\\
\midrule
\endhead
none & \texttt{C-x C-w} & \texttt{write-\codebreak{}file} & Write the current buffer to another file, asked for as C-x C-f asks (C-x C-w), from the folder of the path chosen last in a question for a path, else the buffer’s own folder.\\
none & \texttt{M-x rename-\codebreak{}visited-\codebreak{}file} & \texttt{rename-\codebreak{}visited-\codebreak{}file} & Give the file the current buffer visits a new name, asked for, its own offered to change: a name alone keeps it in its folder, a path (\textasciitilde{}/notes/a.org, ../a.org) moves it there on the same drive.\\
\bottomrule
\end{xltabular}

\subsection{Read-only files and buffers}
\label{sec:195}
A buffer is read-only when you cannot change its text. Tuatara makes a file’s
buffer read-only when your account may not write the file, and the mode line
shows \texttt{[\%]}. You can still look at it, search it and copy from it. A key that
would change the text is refused.

\texttt{C-x C-q} turns read-only on and off for the buffer you are in. Use it to guard
a file that you are only reading, so that a stray key does no harm, or to make
a read-only buffer editable when you know the file will let you save. A number
before the key sets it: a positive number makes the buffer read-only and zero
or a negative number makes it editable.

If you turn read-only off in a file you may not write, the save still asks
\texttt{File ... is write-\codebreak{}protected; try to save anyway?} and the system may refuse.

A file opened safely is read-only too, and so are files in evidence folders.
Those are the subject of chapter 12.

\begin{xltabular}{\textwidth}{@{}l l l >{\raggedright\arraybackslash\hsize=1.000\hsize}X@{}}
\toprule
\textbf{Modal} & \textbf{Non-modal} & \textbf{Command} & \textbf{What it does}\\
\midrule
\endhead
none & \texttt{C-x C-q} & \texttt{read-\codebreak{}only-\codebreak{}mode} & Make the current buffer read-only, or editable again if it is; with a count, read-only when it is positive, editable when it is zero or less.\\
none & \texttt{M-x end-\codebreak{}safe-\codebreak{}visit} & \texttt{end-\codebreak{}safe-\codebreak{}visit} & Open the file of this buffer, opened safely, as usual from now on, where it is.\\
\bottomrule
\end{xltabular}

\subsection{The file changed on disk}
\label{sec:198}
Another program may change a file you have open. A build may rewrite it, a
program may add to a log, or a copy may arrive from another machine. Tuatara
watches the files that you have open and tells you. The echo area and the
mode line say \texttt{Changed on disk}, \texttt{Deleted on disk} or \texttt{Stale, read-\codebreak{}only}. The
last word means that the file was rewritten while it was shown from the disk, so
the text on the screen is no longer what the file holds, and the buffer is
read-only until you read the file again.

In every case Tuatara asks before it does anything to your text. A file that
was changed on disk and a buffer with changes of your own are two versions
of the same file, and you choose which to keep.

\subsubsection{Reading it again}
\label{sec:196}
\texttt{M-x revert-\codebreak{}buffer} reads the file again and replaces the text in the
buffer. It asks first when the buffer has changes you have not saved. The
reading is one step in the undo history, so \texttt{C-/} (modal: \texttt{u}) takes it back if
you made a mistake. Point and the window stay on the same text, so you do not
lose your place. If the file only grew, only the new part is added.

Opening the file again, with \texttt{C-x C-f} or with \texttt{RET} in a list such as the git
status, also reads it again when its buffer has no changes of your own. With
changes of your own the buffer is left as it is, and the echo area says so. A
buffer that \texttt{M-x follow-\codebreak{}file} stopped following is not read again, and a file
deleted on disk is said.

In a buffer that Tuatara made for you, such as a list of buffers, \texttt{r} reads it
again from what it shows. \texttt{g r} does the same. This is how a list is brought up
to date. The buffer list does it by itself.

If you want to see how the two versions differ before you choose, use
\texttt{M-x diff-\codebreak{}files}, which compares the file with another one you name, or a copy
you saved under another name with \texttt{C-x C-w}.

\subsubsection{Keeping up by itself}
\label{sec:197}
Some files are meant to change under you: a log that a program is writing, or a
build’s output. For those you want the buffer to follow the file. \texttt{M-x
follow-\codebreak{}file} turns that on for the buffer you are in. What a program adds to the
end of the file is added to your buffer as it arrives, and if point is at the
end, it stays at the end, as \texttt{tail -f} does in a terminal. A file that is
rewritten is read again, if you have no changes of your own. The mode line says
\texttt{Following} while it is on. Run the command again to stop.

\texttt{*reload-\codebreak{}changed-\codebreak{}files*} does the same for every buffer. It is off as
standard, because most of the time you do not want the text to move under you.
With it on, a file that changed on disk while you have unsaved changes of your
own is not read again, since that would lose them. Tuatara says so, and you
decide with \texttt{M-x revert-\codebreak{}buffer}, which takes the file’s version, or \texttt{C-x C-s},
which writes yours over it.

\begin{try}
\textbf{Try it, with a recipe}
\begin{enumerate}
\item Open \texttt{\textasciitilde{}/recipe.\codebreak{}txt}, type \texttt{Pikelets: 1 cup flour, 1 egg, milk} and \texttt{RET}, so
   that the line has its end, and save with \texttt{C-x C-s}.
\item In a terminal, run \texttt{echo "Add a pinch of salt" >> \textasciitilde{}/recipe.\codebreak{}txt}.
\item Go back to Tuatara. The echo area, and the mode line, say that
   \texttt{recipe.\codebreak{}txt} changed on disk, and that \texttt{M-x revert-\codebreak{}buffer} reads it again.
\item Run \texttt{M-x revert-\codebreak{}buffer}. The new line is there. Press \texttt{C-/} (modal: \texttt{u}) and
   it is gone again, since a revert is one step of undo. Press \texttt{M-\_} (modal:
   \texttt{U}) to bring it back.
\item Run \texttt{M-x follow-\codebreak{}file}. The mode line says \texttt{Following}. Run
   \texttt{echo "Serve warm" >> \textasciitilde{}/recipe.\codebreak{}txt} in the terminal again. The line appears
   in Tuatara at once, with no command from you.
\item Run \texttt{M-x follow-\codebreak{}file} again to stop.
\end{enumerate}
\end{try}

\begin{xltabular}{\textwidth}{@{}l l l >{\raggedright\arraybackslash\hsize=1.000\hsize}X@{}}
\toprule
\textbf{Modal} & \textbf{Non-modal} & \textbf{Command} & \textbf{What it does}\\
\midrule
\endhead
none & \texttt{r}, \texttt{g r}, \texttt{C-c r} & \texttt{revert-\codebreak{}buffer} & Read the current buffer’s file again, replacing its text: one step that undo takes back, unless the text could no longer be trusted (then the history goes too).\\
none & \texttt{M-x follow-\codebreak{}file} & \texttt{follow-\codebreak{}file} & Keep the current buffer up with its file on disk, or stop.\\
\bottomrule
\end{xltabular}

\begin{xltabular}{\textwidth}{@{}l l >{\raggedright\arraybackslash\hsize=1.000\hsize}X@{}}
\toprule
\textbf{Setting} & \textbf{Standard} & \textbf{What it changes}\\
\midrule
\endhead
\texttt{*reload-\codebreak{}changed-\codebreak{}files*} & \texttt{nil} & True to have every buffer keep up with its file on disk by itself.\\
\bottomrule
\end{xltabular}

\subsection{Line endings and letters from other languages}
\label{sec:201}
\subsubsection{Line endings}
\label{sec:199}
A line in a text file ends with a code. Unix and Linux write one character, a
newline. Windows writes two, a carriage return and a newline. Very old
Macintosh files used a carriage return alone. You rarely see the difference
until you open a file from another system.

Tuatara looks at the first line end in a file and keeps to it. A file of
Windows lines shows \texttt{CRLF} in the mode line, and one of carriage return lines
shows \texttt{CR}. A file of plain newlines shows nothing, since that is the normal
case. When you press \texttt{RET}, or paste text, the new line ends the same way as
the rest of the file. A save writes every line end as it was read, so opening
and saving a Windows file does not turn it into a Unix file, and does not
change the lines you did not touch.

In a file whose lines end with a newline, a stray carriage return before the
newline shows as \texttt{\textasciicircum{}M}. That is how a file of mixed endings shows where it is
mixed. A carriage return that stands alone ends a line and shows no mark.

Line numbers, searching, sorting and every other command count lines the same
way, whatever the endings. A search for \texttt{\textasciicircum{}} or \texttt{\$} matches at every line end.

\subsubsection{Letters from other languages}
\label{sec:200}
Tuatara works in UTF-8, the way of writing letters that nearly all modern
files use. It holds accented letters, letters of other alphabets and
emoji, and it saves them back as they were. 

There is no setting for the encoding of a file, and no command to change it.
If a file is not UTF-8, Tuatara does not stop. Every part that is not valid
UTF-8 is kept as a character of its own, shown as \texttt{\textbackslash{}xE9} for the octet \texttt{E9}, and
is written back to the file exactly as it was. A letter with an accent in an
old file in the Latin-1 encoding is such a case. You see \texttt{\textbackslash{}xE9} where you would
expect an \texttt{e} with an accent, and you can edit around it, but the file is not
changed by opening and saving it. If you need to turn such a file into UTF-8,
use a tool such as \texttt{iconv} outside Tuatara.

A file that begins with a UTF-8 byte order mark opens normally. The mark is
kept when you save.

A setting described in chapter 12 shows characters that draw nothing as their
codes, in files opened safely. It helps you find what a file hides.

\subsection{Files with formatted names}
\label{sec:205}
A file name that says what the file holds is easy to find a year later. A
formatted name has three parts: the moment it was made, a title, and some
keywords. The scheme follows Denote, a package for Emacs, so such names read the
same in both editors.

\begin{figure}[H]
\caption{Figure 9.1. A formatted name, in three parts and an end.}\label{org202}
\begin{Verbatim}[breaklines,breakautoindent=false,breakindent=1em,breakafter=/-.\,;:=,commandchars=\\\{\}]
 20261109T091500--letter-to-the-council__letters_house.org
 └──────┬───────┘\discretionary{}{}{}  └──────────┬────\discretionary{}{}{}─────┘ └─────┬──────┘└┬\discretionary{}{}{}┘
   identifier              title             keywords    kind
   (date and time)    (lower case, with     (joined by _)
                         hyphens)
\end{Verbatim}
\end{figure}

The identifier is the year, month and day, a \texttt{T}, then the hour, minute and
second, so a folder of such names sorts by age by itself. Two dashes follow,
then the title in small letters with hyphens between the words. Two
underscores begin the keywords, which are joined by single underscores. The
file’s extension ends the name: \texttt{.org}, \texttt{.txt} or \texttt{.md}.

You keep all such files under one folder, set by \texttt{*formatted-\codebreak{}file-\codebreak{}directory*}.
It is your home folder until you set it. A good place is a folder of notes:

\begin{Verbatim}[breaklines,breakautoindent=false,breakindent=1em,breakafter=/-.\,;:=,commandchars=\\\{\}]
(setf *formatted-file-directory* "~/notes/")
\end{Verbatim}

\subsubsection{Making and finding}
\label{sec:203}
\texttt{C-x n n} makes a new file. It asks for a title, a subfolder, some keywords and
the kind of file, in that order. The keywords you used before are offered.
Press \texttt{RET} with nothing typed to skip a question. The new buffer is named from
what you answered, and begins with a header for its kind, holding the title,
the date, the keywords and the identifier. Nothing is written until you save
with \texttt{C-x C-s}. The date and identifier are those of the moment you asked, not
of the moment you saved.

\texttt{C-x n o} opens one of your files by any part of its name. The list is grouped
by kind, newest first, and shows each name as its date, its title and its
keywords, with the size and when it changed on the right. You type a few
letters of the title or a keyword, and \texttt{RET} opens the best match. The
preview of \texttt{*preview-\codebreak{}formatted-\codebreak{}files*} works here as it does for \texttt{C-x C-f}.

\texttt{C-x n f} is \texttt{C-x C-f} for such files. It opens a file by its path, and if the
name you type is not there, it makes a new file with a formatted name in the
folder you typed. It asks for the title and the keywords, and for the end of
the name, which you can have filled in by \texttt{*formatted-\codebreak{}file-\codebreak{}extension*}.

\subsubsection{Renaming}
\label{sec:204}
You can give an old file a formatted name. \texttt{C-x n r} asks for a title and
keywords, offering the ones the file already has, and renames the file. The
identifier is kept if the file has one, or made from the time the file last
changed. The extension is kept. A buffer that shows the file follows the new
name, and the list of recent files is updated.

\texttt{C-x n k} changes only the keywords. It keeps the titles and asks once for new
keywords, which it gives to the file, or to every marked file when you use it
in a folder listing. Use it to put a whole month of letters under \texttt{done}.

\begin{try}
\textbf{Try it, with some notes}
\begin{enumerate}
\item Add \texttt{(setf *formatted-\codebreak{}file-\codebreak{}directory* "\textasciitilde{}/\codebreak{}notes/")} to your init file and run
   it, or restart Tuatara. Make the folder in a terminal with \texttt{mkdir \textasciitilde{}/notes}.
\item Press \texttt{C-x n n}. Answer the title with \texttt{Pikelets for Sunday}, the folder
   with \texttt{RET}, the keywords with \texttt{recipe, house} and \texttt{RET} twice, and the kind
   with \texttt{RET} for Org.
\item The buffer is named like
   \texttt{20261109T091500-\codebreak{}-\codebreak{}pikelets-\codebreak{}for-\codebreak{}sunday\_\_\codebreak{}house\_\codebreak{}recipe.\codebreak{}org}. The keywords are
   in alphabetical order. Press \texttt{C-x C-s}.
\item Press \texttt{C-x n o}, type \texttt{pike} and \texttt{RET}. You are back in the file.
\item Press \texttt{C-x n k}, type \texttt{recipe, sunday} and press \texttt{RET} twice. The name now
   ends \texttt{\_\_recipe\_\codebreak{}sunday.\codebreak{}org}, the header is changed to match, and the buffer
   follows.
\end{enumerate}
\end{try}

\begin{xltabular}{\textwidth}{@{}l l l >{\raggedright\arraybackslash\hsize=1.000\hsize}X@{}}
\toprule
\textbf{Modal} & \textbf{Non-modal} & \textbf{Command} & \textbf{What it does}\\
\midrule
\endhead
none & \texttt{C-x n f} & \texttt{find-\codebreak{}file-\codebreak{}formatted} & Visit a file, or make a new one with a formatted name.\\
none & \texttt{C-x n o} & \texttt{open-\codebreak{}file-\codebreak{}formatted} & Open one of your files with formatted names, those in \textbf{FORMATTED-FILE-DIRECTORY} and its subdirectories, choosing among them by any part of the name: listed under their subdirectory or their kind, the ones changed last first, each as its date, its title and its keywords, with its size and when it changed at the right.\\
none & \texttt{C-x n n} & \texttt{new-\codebreak{}file-\codebreak{}formatted} & Make a new file with a formatted name in \textbf{FORMATTED-FILE-DIRECTORY}: asks what \textbf{FORMATTED-FILE-PROMPTS} says (as standard the title, the subdirectory, the keywords, the kind of file and the template), and names it as find-file-formatted does, with the header for its kind.\\
none & \texttt{C-x n r} & \texttt{rename-\codebreak{}formatted} & Give a file a formatted name: its identifier kept, or made from when it was last changed; a title and keywords asked for, the file’s own offered; its extension kept.\\
none & \texttt{C-x n k} & \texttt{rename-\codebreak{}formatted-\codebreak{}keywords} & Give files formatted names with other keywords, their titles kept, asked for once for them all: the current buffer’s file, or in a list of a directory the marked files or the one on the line.\\
\bottomrule
\end{xltabular}

\begin{xltabular}{\textwidth}{@{}l >{\raggedright\arraybackslash\hsize=0.909\hsize}X >{\raggedright\arraybackslash\hsize=1.091\hsize}X@{}}
\toprule
\textbf{Setting} & \textbf{Standard} & \textbf{What it changes}\\
\midrule
\endhead
\texttt{*formatted-\codebreak{}file-\codebreak{}directory*} & \texttt{"\textasciitilde{}/"} & The directory of your files with formatted names: open-file-formatted lists those in it and its subdirectories, new-file-formatted makes new ones there, and the keywords in their names are offered for every new file with a formatted name.\\
\texttt{*formatted-\codebreak{}file-\codebreak{}extension*} & \texttt{""} & The end of a new file’s name that find-file-formatted offers when it asks for one (New file EXTENSION): nothing typed there takes it.\\
\texttt{*formatted-\codebreak{}file-\codebreak{}keywords*} & \texttt{nil} & Keywords offered for a new file with a formatted name, besides the keywords in the names of the files in \textbf{FORMATTED-FILE-DIRECTORY}: none as standard.\\
\texttt{*formatted-\codebreak{}file-\codebreak{}sort-\codebreak{}keywords*} & \texttt{t} & True to put the keywords of a formatted name in alphabetical order.\\
\texttt{*formatted-\codebreak{}file-\codebreak{}prompts*} & \texttt{(:title :subdirectory :keywords :file-\codebreak{}ty...} & What new-file-formatted asks for a new file, in this order: :title, :subdirectory (of \textbf{FORMATTED-FILE-DIRECTORY}), :keywords, :file-type (one of \textbf{FORMATTED-FILE-TYPES}; else \textbf{FORMATTED-FILE-TYPE} is used) and :template (blank, or one of your note templates’ files, asked only when there are any).\\
\texttt{*formatted-\codebreak{}file-\codebreak{}type*} & \texttt{:org} & The kind of a new file when \textbf{FORMATTED-FILE-PROMPTS} does not ask, and the one offered first when it does: a name in \textbf{FORMATTED-FILE-TYPES}.\\
\texttt{*formatted-\codebreak{}file-\codebreak{}types*} & \texttt{((:\codebreak{}org ".org" :org "\#+title: \{title\} \#+d...} & The kinds of files with formatted names, each (NAME EXTENSION STYLE HEADER): the end of the file’s name, how its header writes the date and keywords, and the header a new file begins with.\\
\texttt{*preview-\codebreak{}formatted-\codebreak{}files*} & \texttt{nil} & True to show, while you choose one of your files with formatted names with M-x open-file-formatted, the file selected in the window you are in, opened safely, as \textbf{PREVIEW-FILES} shows files.\\
\bottomrule
\end{xltabular}

\begin{xltabular}{\textwidth}{@{}l >{\raggedright\arraybackslash\hsize=1.000\hsize}X@{}}
\toprule
\textbf{Face} & \textbf{What it paints}\\
\midrule
\endhead
\texttt{:completion-\codebreak{}formatted-\codebreak{}keywords} & The keywords of a file with a formatted name, from the name, in the list C-x n o shows.\\
\bottomrule
\end{xltabular}

The other commands that work on these files, such as searching their text and
writing links between them, are described with Org in chapter 14.

\subsection{Files on another machine}
\label{sec:212}
A file on a server can be opened as if it were here. Write the machine’s name
in front of the path, between \texttt{/ssh:} and a colon. The name is the one you give
to ssh: a plain name, \texttt{me@\codebreak{}pinetree}, or a name from your \texttt{\textasciitilde{}/.\codebreak{}ssh/\codebreak{}config}.

\begin{Verbatim}[breaklines,breakautoindent=false,breakindent=1em,breakafter=/-.\,;:=,commandchars=\\\{\}]
C-x C-f /ssh:pinetree:/etc/hosts RET
\end{Verbatim}

Nothing has to be installed on the server. Tuatara uses the file protocol that
every ssh server offers, through the ssh program already on your computer, with
your own keys. It sends nothing else to the machine: it runs no commands there.
Everything is done in the background, so Tuatara never waits for a slow
connection.

This is not the daemon of chapter 18, which runs the whole editor on the other
machine. Here Tuatara runs on your computer, and only the files are elsewhere.

\begin{figure}[H]
\caption{Figure 9.2. The file stays on the server; you work on a private copy.}\label{org206}
\begin{Verbatim}[breaklines,breakautoindent=false,breakindent=1em,breakafter=/-.\,;:=,commandchars=\\\{\}]
   Your computer                              The server, pinetree
  ┌───────────────\discretionary{}{}{}────────────────\discretionary{}{}{}───┐       ┌───────────────\discretionary{}{}{}─────┐
  │  Tuatara                         │       │                    │
  │   buffer  ◄── a copy only you ───┼─ ssh ─┤   /etc/hosts       │
  │               can read           │       │   stays where it   │
  │   opening reads it  ◄────────────┼──\discretionary{}{}{}─────┤   is              │
  │   C-x C-s writes it back ────────┼──────►\discretionary{}{}{}│                    │
  └───────────────\discretionary{}{}{}────────────────\discretionary{}{}{}───┘       └───────────────\discretionary{}{}{}─────┘
        both go through ssh, with your own keys
\end{Verbatim}
\end{figure}

Typing \texttt{/ssh:\codebreak{}pinetree:/\codebreak{}etc/} offers the names in that folder, as for a folder on
your own computer. \texttt{/ssh:\codebreak{}pinetree:} alone is your home folder there. A file
from a server opens as a local file does: a picture is drawn, and a PDF shows
its pages. Any other file that is not text opens as octets, in the hex view.

\subsubsection{Saving}
\label{sec:207}
\texttt{C-x C-s} saves as usual. The text goes to a new file beside the one on the
server, which takes the old file’s owner, group and permissions, and takes its
place only when the file there is still as you read it. If somebody changed it
meanwhile, the echo area says the file was not written, since it was changed
on the server since it was read. The buffer stays as it is, so nothing is
lost. \texttt{M-x revert-\codebreak{}buffer} reads the server’s version, and then you can copy
your changes over it. A new name is a new file on the server when you save.

A file with more than one name, or one whose owner a new file cannot be given,
is written in place, and the echo area says so. If such a write fails halfway,
the file there is partly written. The buffer keeps your text, and the next save
asks whether to write it over, whatever the file now holds. Answer \texttt{yes}
to write it whole.

Tuatara keeps the file you work on as a copy under \texttt{\textasciitilde{}/.\codebreak{}cache/\codebreak{}tuatara/\codebreak{}remote/},
which only you can read. The copy is deleted when you kill the buffer. A file on
another machine has no git marks, no language server, no auto-saved copy and no
build, since those would read what is around the file, and what is around it is
not here. A file’s own settings are taken only from a machine you trust, listed
by \texttt{M-x server-\codebreak{}list}, chapter 18. Each request has \texttt{*remote-\codebreak{}seconds*} to be
answered, 30 as standard. Past that, the work in hand fails and says so.

\subsubsection{The first time you meet a machine}
\label{sec:209}
When ssh meets a machine for the first time it asks if you trust it. Tuatara
asks the same question. While the question is open, a buffer below the windows
shows ssh’s own words in full, with the fingerprint of the machine’s key. At
the bottom is a short prompt for your answer. Completing a file name as you
type never asks anything. The question comes when you press \texttt{RET}. If a
connection takes more than a second, the echo area says \texttt{Connecting to
pinetree...} until it is made or has failed.

\begin{figure}[H]
\caption{Figure 9.3. The first connection: question above, prompt below.}\label{org208}
\begin{Verbatim}[breaklines,breakautoindent=false,breakindent=1em,breakafter=/-.\,;:=,commandchars=\\\{\}]
  ssh to pinetree

  The authenticity of host 'pinetree (10.0.0.5)' can't be established.
  ED25519 key fingerprint is:
  SHA256:LRMTH6kEUp1dDyMP\discretionary{}{}{}GJlAvQhgiX4cSO0J\discretionary{}{}{}X/aPMh350Sk
  This key is not known by any other names.

  Type yes or the fingerprint to accept, anything else refuses.
 TUATARA[N]  P  [-] *Question*  L1 C0  Help
 Accept this key (yes/no/fingerprint): █
\end{Verbatim}
\end{figure}

Compare the fingerprint with one you got from whoever runs the server. Then
type \texttt{yes} and press \texttt{RET}. You can type the fingerprint instead. Anything else
refuses, a plain \texttt{y} as well. A key that has changed since ssh last saw it is
never offered. ssh refuses the connection, and the echo area says the key
differs from the one ssh knows. That is the sign of a machine that has been
replaced, or of one pretending to be it. Ask the person who runs the server.

\subsubsection{Passwords}
\label{sec:210}
If your key is in your ssh agent, nothing is asked. A key with a passphrase is
asked for once, at a prompt saying it stays here, on your computer. If you log
in with a password, the prompt says the server asks, and that what you type
goes to the server. A jump host, if you use one, is named. What you type shows
as stars, and is passed to ssh and kept nowhere else: not in a history, not in
a recorded macro, and not in a log.

To paste a password from a password manager, copy it there and press \texttt{C-y} in
the prompt. In a terminal, the paste of the terminal works too. The password is
not put in the kill ring. \texttt{C-g} gives the question up, and the connection fails
with a message that the question was not answered. Nothing is left running.

\subsubsection{Folders and previews on a server}
\label{sec:211}
\texttt{C-x C-f /ssh:\codebreak{}pinetree:/\codebreak{}srv/\codebreak{}notes/ RET}, or \texttt{C-x d} with that name, lists the
folder as one is listed on your computer. Chapter 10 describes the list. What
would read many files, or keep watching them, is left out on a server: there is
no git column and no folder size, and the list is read again after each change
you make, not as files change.

A preview shows each file beside the list as you move. On a server each file
must come to your computer first, so the preview starts off. \texttt{P} in the list
turns it on, and moving to another file stops the transfer under way. A file
larger than \texttt{*remote-\codebreak{}preview-\codebreak{}size-\codebreak{}limit*} is not fetched, and the preview says
it is too large. The standard is 20 million octets. \texttt{RET} still opens it.
Files in evidence folders are not read for a preview, and a list opened safely
has no preview.

A server keeps no trash. A file deleted there is deleted for good, and the list
says so in its question.

Evidence folders may be on a server as well. An entry of \texttt{*evidence-\codebreak{}folders*}
may name a folder there, and files in it open safely and are never saved.
Chapter 12 describes it.

\begin{try}
\textbf{Try it, with notes on a server called pinetree}
These steps need a machine you can reach with ssh. Use your own name for
\texttt{pinetree}.

\begin{enumerate}
\item In a terminal, run \texttt{ssh pinetree} once, so you know the login works. Leave
   it again with \texttt{exit}.
\item In Tuatara, press \texttt{C-x C-f} and type \texttt{/ssh:\codebreak{}pinetree:}. Your home folder there
   is offered. Add \texttt{notes.\codebreak{}txt} and press \texttt{RET}. The buffer is empty.
\item Press \texttt{i}, type \texttt{Sheep: 41 counted on Friday} and press \texttt{C-g}. Press
   \texttt{C-x C-s}. The echo area says \texttt{Wrote /ssh:\codebreak{}pinetree:\codebreak{}notes.\codebreak{}txt}.
\item In another terminal, run \texttt{ssh pinetree 'echo "Gate fixed" >> notes.\codebreak{}txt'}.
   Back in Tuatara, press \texttt{i}, type \texttt{x} and press \texttt{C-g}, then \texttt{C-x C-s}. The
   echo area says the file changed on the server since it was read and was not
   written. Your buffer is unchanged.
\item Run \texttt{M-x revert-\codebreak{}buffer}. The line \texttt{Gate fixed} is there.
\item Press \texttt{C-x d} and type \texttt{/ssh:\codebreak{}pinetree:}, then \texttt{RET}. Your home folder opens
   as a list. Move to \texttt{notes.\codebreak{}txt} and press \texttt{P} to see it beside the list.
\item If you wish, press \texttt{C-x k} on the buffer. The private copy on your computer
   is deleted, and the file on pinetree is as you saved it.
\end{enumerate}
\end{try}

\begin{xltabular}{\textwidth}{@{}l l l >{\raggedright\arraybackslash\hsize=1.000\hsize}X@{}}
\toprule
\textbf{Modal} & \textbf{Non-modal} & \textbf{Command} & \textbf{What it does}\\
\midrule
\endhead
none & \texttt{P}, \texttt{C-c P} & \texttt{list-\codebreak{}toggle-\codebreak{}preview} & Turn the preview beside this list on or off: the entry on the current line shown in a window beside the list, on its right when there is room, else below it, following the cursor from line to line.\\
\bottomrule
\end{xltabular}

\begin{xltabular}{\textwidth}{@{}l l >{\raggedright\arraybackslash\hsize=1.000\hsize}X@{}}
\toprule
\textbf{Setting} & \textbf{Standard} & \textbf{What it changes}\\
\midrule
\endhead
\texttt{*remote-\codebreak{}seconds*} & \texttt{30} & Seconds a server reached for its files (/ssh:HOST:/path) has to answer each request, and ssh to connect: past them the connection is closed and the work in hand fails, saying so.\\
\texttt{*remote-\codebreak{}preview-\codebreak{}size-\codebreak{}limit*} & \texttt{20000000} & The most octets a file on another machine may hold to be shown in the preview of its list (M-x list-toggle-preview), 20 MB by default: all of it has to come from the server first.\\
\bottomrule
\end{xltabular}

\subsection{Opening a file with another program}
\label{sec:213}
Tuatara shows text, and also pictures, PDFs and some office documents. There is
work it does not do, such as printing, filling in a PDF form, or playing a
video. For that you can hand the file to the program your desktop uses for that
kind of file. \texttt{O} does it in a folder listing, in a picture and in a PDF’s view,
and \texttt{M-x open-\codebreak{}externally} does it for the file of the buffer you are in. Tuatara
runs \texttt{xdg-\codebreak{}open}, the command desktops provide for this, so the file opens in
what you would get by double-clicking it.

The other program is not shut in as Tuatara’s own readers are. It can read
everything you can. So Tuatara asks first, unless you turn off
\texttt{*open-\codebreak{}externally-\codebreak{}ask*}. It also looks at what the file is, not what its name
says. A program named \texttt{holiday.\codebreak{}jpg} is refused. Only the kinds of file in
\texttt{*open-\codebreak{}externally-\codebreak{}kinds*} are handed over: pictures, PDFs, sound and video,
and office documents.

To use another program for a kind, set \texttt{*open-\codebreak{}externally-\codebreak{}programs*} to a list
of a type and a program:

\begin{Verbatim}[breaklines,breakautoindent=false,breakindent=1em,breakafter=/-.\,;:=,commandchars=\\\{\}]
(setf *open-externally-programs*
      '(("application/pdf" . "zathura")
        ("video/*" "mpv" "--fs")))
\end{Verbatim}

The program given there is run as it is, with the file’s name after it.
\texttt{*open-\codebreak{}externally-\codebreak{}program*} names the command used for every other kind, which
is \texttt{xdg-\codebreak{}open} as standard.

\begin{xltabular}{\textwidth}{@{}l l l >{\raggedright\arraybackslash\hsize=1.000\hsize}X@{}}
\toprule
\textbf{Modal} & \textbf{Non-modal} & \textbf{Command} & \textbf{What it does}\\
\midrule
\endhead
none & \texttt{O}, \texttt{C-c O} & \texttt{open-\codebreak{}externally} & Open the file shown, in the desktop’s own program for its kind (xdg-open), or the one \textbf{OPEN-EXTERNALLY-PROGRAMS} gives its Multipurpose Internet Mail Extensions (MIME) type, for what Tuatara does not do: printing, forms, playing.\\
\bottomrule
\end{xltabular}

\begin{xltabular}{\textwidth}{@{}l >{\raggedright\arraybackslash\hsize=0.542\hsize}X >{\raggedright\arraybackslash\hsize=1.458\hsize}X@{}}
\toprule
\textbf{Setting} & \textbf{Standard} & \textbf{What it changes}\\
\midrule
\endhead
\texttt{*open-\codebreak{}externally-\codebreak{}ask*} & \texttt{t} & True for M-x open-externally to ask before opening a file in the desktop’s own program, which reads it outside the sandbox; NIL to open it at once.\\
\texttt{*open-\codebreak{}externally-\codebreak{}kinds*} & \texttt{(:png :jpeg :gif :webp :bmp :tiff :raw :...} & The kinds of file M-x open-externally opens in the desktop’s own program, as their first octets tell them (pictures, PDFs, sound and video, and Office documents, as standard): keywords as TUATARA.IO:OCTETS-KIND gives them, :png, :jpeg, :pdf, :docx and the like.\\
\texttt{*open-\codebreak{}externally-\codebreak{}program*} & \texttt{"xdg-\codebreak{}open"} & The program that opens a file in the desktop’s own program for its kind, given the file’s name: xdg-open, looked for on the PATH.\\
\texttt{*open-\codebreak{}externally-\codebreak{}programs*} & \texttt{nil} & Programs to open files and addresses in instead of the one the desktop would choose (xdg-open), by Multipurpose Internet Mail Extensions (MIME) type as xdg-mime names them: a list of (TYPE .\\
\bottomrule
\end{xltabular}

\subsection{Starting Tuatara with files}
\label{sec:215}
\texttt{tuatara letter.\codebreak{}txt} opens the file in the terminal you typed it in. Several
names open several buffers, the first one shown. \texttt{tuatara --gui letter.\codebreak{}txt}
opens a window of its own instead. Other files you had open, if you keep a
session as chapter 7 described, are opened behind the ones you named.

A number with a plus before a file name says which line to start on. \texttt{tuatara
+2 letter.\codebreak{}txt} opens the letter with the cursor on its second line. \texttt{+2:5} adds
a column. The numbers count from 1. A file with fewer lines opens at its last
line, and says so. You can name a line for each file: \texttt{tuatara +120 main.\codebreak{}lisp
+3 notes.\codebreak{}org}.

\texttt{tuatara --safe letter.\codebreak{}txt}, or \texttt{-R} for short, opens the files safely, as
described in chapter 12: read-only, with nothing in them acted on. It holds
for every file in that command. \texttt{--hex} opens them in the hex view instead of as
text. A single \texttt{-} reads what another program sends: \texttt{ls | tuatara -} shows the
output in a buffer called \texttt{*stdin*}. After \texttt{--}, every word is a file, so
\texttt{tuatara -- -odd.\codebreak{}txt} opens a file named \texttt{-odd.\codebreak{}txt}.

The table at the end of this section lists every option, as \texttt{tuatara --help}
does. The ones that start the daemon, make a client, compare files or merge
them are explained in chapters 15 and 18.

\begin{try}
\textbf{Try it, with a letter}
\begin{enumerate}
\item In a terminal, make a letter of five lines:
   \texttt{printf 'one\textbackslash{}ntwo\textbackslash{}nthree\codebreak{}\textbackslash{}nfour\textbackslash{}nfive\textbackslash{}n' > letter.\codebreak{}txt}.
\item Run \texttt{tuatara +4 letter.\codebreak{}txt}. The cursor is on \texttt{four}, and the mode line says
   \texttt{L4}.
\item Press \texttt{C-x C-c} to leave. Run \texttt{tuatara --safe letter.\codebreak{}txt}. The mode line
   says \texttt{Safe}. Press \texttt{i}. The buffer is read-only, so nothing can be typed.
\item Press \texttt{M-x end-\codebreak{}safe-\codebreak{}visit} and \texttt{RET}. The file opens as usual.
\item Leave with \texttt{C-x C-c}.
\end{enumerate}
\end{try}

\subsubsection{Options}
\label{sec:214}
\begin{xltabular}{\textwidth}{@{}>{\raggedright\arraybackslash\hsize=0.447\hsize}X >{\raggedright\arraybackslash\hsize=1.553\hsize}X@{}}
\toprule
\textbf{Option} & \textbf{What it does}\\
\midrule
\endhead
\texttt{--gui} & Open a window of its own instead (as when started with no terminal). Use it when the terminal is the limit: the window draws pictures and PDFs itself, shows 24-bit colour whatever TERM says, and takes the desktop’s clipboard through GTK. It needs a graphical session (WAYLAND\_DISPLAY or DISPLAY), which a console or an ssh login has not. How the window is drawn is the setting \textbf{gui-renderer}.\\
\texttt{--frames N} & Open N frames, from 1 to 16: with --gui each a window of its own, in the terminal screens shown one at a time (C-x 5 o goes to the next); the files named are shown one in each frame in turn, those left over open behind the first; --frames=N does as well. Use it to start with several views of the same buffers. With --gui they are windows you can put on different screens. In a terminal only one frame shows at a time, so more than one is for C-x 5 o; to see two views at once, split the window (C-x 2, C-x 3). The cost is small: measured once here, 16 frames on a small file added about 30 ms to the first screen and no memory worth naming.\\
\texttt{-c INIT-\codebreak{}FILE, --config INIT-\codebreak{}FILE} & Read INIT-FILE as the init file instead of \textasciitilde{}/.config/tuatara/init.lisp; --config=INIT-FILE does as well. Use it to keep several setups (one for writing, one for a work machine), to try a changed init file without touching the one you rely on, or to start from a file kept elsewhere. The file is read as the usual one is, so the packages it names are loaded. A file that is not there ends the start with status 2; the usual init file is not used in its place.\\
\texttt{+LINE[:\codebreak{}COLUMN]} & Before a file: open it with the cursor on that line (from 1; the last line when the file has fewer) and column (from 1; the C in the mode line counts from 0), as in tuatara +120:8 main.lisp +3 notes.org.\\
\texttt{--diff OLD NEW} & Open the differences between the two files, or two folders, as M-x diff-files shows them, and exit when you quit; for git, as git config difftool.tuatara.cmd ‘tuatara --diff “\$LOCAL” “\$REMOTE”‘.\\
\texttt{--merge LOCAL BASE REMOTE MERGED} & Merge the three files into MERGED, the conflicts marked there to resolve, then save it; exits 0 when MERGED is saved with no conflict left, else 1; for git, as git config mergetool.tuatara.cmd ‘tuatara --merge “\$LOCAL” “\$BASE” “\$REMOTE” “\$MERGED”’ and mergetool.tuatara.trustExitCode true.\\
\texttt{--safe, -R} & Open the files safely, whatever folder they are in: read only, no mode or hook acting on them, nothing written. Use it to look at a file you do not trust or must not alter: a download, the contents of a disk image, evidence. Nothing in the folder acts on the file and nothing is written, so opening it leaves it as it was. What you give up is what the folder’s own settings and tools would add, and saving changes.\\
\texttt{--hex} & Open the files in the hex view, never first as text. Use it for a file that is not text (a program, an image, a database), or for text whose bytes you must read as stored, where the text view could hide a stray byte or a different line ending. A byte takes more room on the screen than a character. Measured once here, a 1 GB file showed its first screen in the same time as hex or as text (about 0.35 s).\\
\texttt{-} & As a file: read standard input into a buffer, \textbf{stdin}, opened safely (tuatara - < data, or \ldots{} \(\vert{}\) tuatara -); the terminal stays the keyboard. Use it to look at what another program writes, in a pipe, before it goes on. The buffer is opened safely, as --safe does, so it is for reading and for taking parts of out.\\
\texttt{--tutor} & Start in the tutor.\\
\texttt{--daemon[=\codebreak{}NAME]} & Run one editor in the background, with no screen, for tuatara --client to show in terminals, and return once it is ready; it reads the init file and loads the packages as Tuatara does, and runs until M-x daemon-quit or tuatara --stop, or until it is sent SIGTERM, as at logout, which keeps the changes not saved as auto-saved copies and asks nothing; NAME, default when none is given, tells several apart; its log is \textasciitilde{}/.cache/tuatara/daemon-NAME.log, its session \textasciitilde{}/.local/state/tuatara/session-daemon-NAME. Use it for an editor that is already there when you want one: a client (tuatara --client) shows a frame in about 0.23 s, against about 0.32 s to start an editor of its own (measured once here, on a small file), and the buffers stay open when a terminal closes. Several names keep several separate sessions. The cost is a process that stays in memory until it is stopped, with whatever it has open: a large file stays held. It needs XDG\_RUNTIME\_DIR, where its socket is.\\
\texttt{--client[=\codebreak{}NAME]} & Show the daemon NAME in this terminal, or with --gui in a window of its own (C-x 5 2 opens another), a frame of its own, the files named opened in it, with +LINE, --safe, --hex, - and --frames as for Tuatara itself; --diff and --merge as well, with their exit codes; C-x C-c closes this client alone, its buffers kept in the daemon, and the daemon goes on; as git’s editor (GIT\_EDITOR=’tuatara --client’) the commit is made when the message is finished; exits 1 when no daemon of that name runs, or when it ends. Use it for the quicker start a daemon gives (about 0.23 s against about 0.32 s here) and one set of buffers in every terminal. In return, a change to the init file or a package reaches the daemon only when it is stopped and started again, and every client works in the one editor. As git’s editor, add --start so that it works when no daemon runs.\\
\texttt{--client=\codebreak{}HOST[:\codebreak{}NAME]} & Show the daemon NAME running on the machine HOST, reached with the system’s ssh, as --client shows one here: HOST as ssh takes it (build.example, me@build.example), a name from \textasciitilde{}/.ssh/config written with a colon after it (build: or build:work); its host key must be known already; the files named are that machine’s, a relative name from your home folder there; - is not taken, and pictures stay on each side, the clipboard, dropped files and your agents too unless HOST is trusted with them (M-x server-list); with --start, the daemon there is started first when none runs; --stop=HOST[:NAME] stops it. Use it to work on the machine that has the files, a build server or a remote box, with your own screen and keys. Tuatara must be installed there, and the system’s ssh is all it needs here. The editor runs there, so files are read at the speed of its disks, and each key crosses the network: a slow link shows as lag. The clipboard, dropped files and pictures stay on each side unless HOST is trusted with them, so a copy made there does not reach this machine’s clipboard by itself.\\
\texttt{--daemon-\codebreak{}socket[=\codebreak{}NAME]} & Say where the daemon NAME listens, which Tuatara this is and whether it runs, in lines a client on another machine reads over ssh, with where gpg here looks for its agent, and exit; with --start, start it first when none runs.\\
\texttt{--receive-\codebreak{}file PATH} & Write standard input, a file dropped on the window of a client on another machine, at PATH, from your home folder unless it begins with /, for you alone and under a name nothing has yet; say where, and exit.\\
\texttt{--start} & With --client: start the daemon first when none runs. Use it in a script or as git’s editor, where there may be no daemon yet: the first call starts one and later calls find it. Without it, a client with no daemon exits 1.\\
\texttt{--fallback} & With --client: when no daemon runs, start Tuatara as it starts without --client instead, the files opened in an editor of its own; the desktop’s Tuatara (client) starts it so. Use it where an editor must always open, as in a desktop entry. When no daemon runs it costs the ordinary start (about 0.32 s here) and gives an editor of its own, not shared with later clients.\\
\texttt{--stop[=\codebreak{}NAME]} & Stop the daemon NAME, asking in this terminal about the changes not saved in it, as M-x daemon-quit does; with no terminal, nothing is asked, and the changes are kept as auto-saved copies that M-x recover-file offers. Use it to free the memory of a daemon you no longer need, or to start it again after you change the init file or a package. With no terminal nothing is asked, so a script can stop it without waiting; what was not saved is kept as auto-saved copies.\\
\texttt{-Q, --no-\codebreak{}init} & Start without the init file and so without packages, with what Tuatara keeps between sessions (files visited lately, the session, bookmarks) kept apart in a folder that goes with the session, the real one neither read nor changed. Use it when something in the init file or a package gets in the way: a setting that slows the editor or a package that fails to load. You get the editor as it ships, so you can tell whether the fault is yours. What you do meanwhile is not remembered. It is no faster to start: with the starter init file the first screen took the same time with and without it (about 0.3 s, measured once).\\
\texttt{--generate-\codebreak{}config} & Write the init file Tuatara writes when there is none, every setting in it commented out, as \textasciitilde{}/.config/tuatara/generated-init.lisp, and exit; a file there already is left as it is. Use it for a fresh list of every setting after an update, to compare with your own init file. It never replaces a file that is there, and the list is all comments, so you copy out the few settings you want to change.\\
\texttt{--generate-\codebreak{}autostart} & Write the entry that has the desktop start the daemon as you log in, \textasciitilde{}/.config/autostart/tuatara-daemon.desktop, running tuatara --daemon, and exit; with --daemon=NAME, the daemon NAME’s, tuatara-daemon-NAME.desktop; a file there already is left as it is. Use it for a daemon that is ready from login, so that the first terminal or window you open starts quickly. It costs the daemon’s memory from the start of every session, whether or not you edit that day. Delete the file to undo it.\\
\texttt{--help, -h} & Show this and exit.\\
\texttt{--version, -V} & Show the version and exit.\\
\texttt{--} & The arguments after it are all files, even when one starts with a dash or a plus.\\
\texttt{--install-\codebreak{}desktop} & (the AppImage) copy it to \textasciitilde{}/.local/bin/tuatara and put Tuatara and Tuatara (client) in the desktop’s applications, the icon too. Use it to start the AppImage from the application menu. The copy at \textasciitilde{}/.local/bin/tuatara does not update itself: run this again from a new AppImage.\\
\texttt{--remove-\codebreak{}desktop} & (the AppImage) take them away again.\\
\texttt{--dark, --light} & (the AppImage) show the window’s title bar in the dark or the light style; with --install-desktop, for the desktop’s applications. Use it when the window’s title bar clashes with the desktop. It changes only the title bar’s style; the text and colours come from the theme.\\
\texttt{--dynamic-\codebreak{}space-\codebreak{}size SIZE} & The most memory Tuatara’s own data may take, such as 6GB; by default half the machine’s memory, from 2GB to 16GB, or \$TUATARA\_HEAP\_MB megabytes. Use it to set the heap for one run; \$TUATARA\_HEAP\_MB does for every run. That variable says what raising or lowering it does.\\
\bottomrule
\end{xltabular}

\section{Folders}
\label{sec:250}
Looking at a folder, and copying, moving, renaming and deleting what is in it.

A folder holds files, and sometimes other folders. When you want to find a
photo, tidy a pile of letters or clear out last year’s downloads, you need to
see what is there and then do something to it. Tuatara’s file manager is
called rimu. It shows a folder as a list, a line for each thing in it, and it
lets you pick some of those lines and act on them. The picking works as it
does in the buffer list of chapter 7, so most of what you learned there
applies here.

This chapter starts with opening a folder and moving about in it. Then it
shows how to change what the list shows, how to look at a file without
opening it, and how to choose entries. The rest is what you can do with the
choice: make, copy, move, rename, delete, pack into an archive and compare.
Chapter 9 explained how to open a single file. This chapter is about many
files at a time.

As in the other chapters, every key here is pressed in NORMAL state unless it
starts with \texttt{C-x}. \texttt{i} starts INSERT state, and \texttt{C-g} or \texttt{ESC} returns to
NORMAL. In a list you rarely need INSERT, because the keys are commands.

\subsection{Opening a folder}
\label{sec:219}
\texttt{C-x d} asks which folder you want and lists it. The question already holds
the folder of the buffer you are in, so you often only press \texttt{RET}. To go
elsewhere, type the rest of a path, or start again from your home folder with
\texttt{\textasciitilde{}/}. \texttt{TAB} completes a name as you type it.

If you are looking at a file and want to see what else is in its folder,
\texttt{C-x C-j} lists that folder with the cursor on the file. \texttt{C-x 5 d} does what
\texttt{C-x d} does in a new window of its own.

From a terminal you can start Tuatara on a folder, for example with
\texttt{tuatara \textasciitilde{}/Photos/\codebreak{}holiday/}, and the list is the first thing you see.

Figure 10.1 shows a list for a folder of holiday photos.

\begin{figure}[H]
\caption{Figure 10.1. A list of a folder, with two entries marked.}\label{org217}
\begin{Verbatim}[breaklines,breakautoindent=false,breakindent=1em,breakafter=/-.\,;:=,commandchars=\\\{\}]
  ~/Photos/holiday:
  Napier
  Rotorua
  Wanaka
> IMG_0412.jpg
  IMG_0413.jpg
> IMG_0414.jpg
  itinerary.txt
\end{Verbatim}
\end{figure}

The first line is the folder’s whole name. Below it come the entries, one to
a line, with folders first and then files. Names are in the order that \texttt{ls}
uses, so capital letters come before small ones. The two spaces before each
name are where marks go. In the figure, two photos have a \texttt{>}.

Names that start with a dot are hidden. Those are files that programs keep
for themselves, such as \texttt{.bashrc}, and you seldom want to see them. Folder
names have a colour of their own. Names of files that are only a by-product,
such as compiled files, are set back in a fainter colour.

The mode line at the bottom says \texttt{Rimu} and how the list is sorted. \texttt{q}
leaves the list and returns to the buffer you were in before.

Rimu keeps a list up to date. If another program adds, renames or removes a
file in the folder, the line appears or goes within a moment. \texttt{r} reads the
folder again at once, and your marks stay.

\subsubsection{Commands for opening a folder}
\label{sec:218}
\begin{xltabular}{\textwidth}{@{}l l l >{\raggedright\arraybackslash\hsize=1.000\hsize}X@{}}
\toprule
\textbf{Modal} & \textbf{Non-modal} & \textbf{Command} & \textbf{What it does}\\
\midrule
\endhead
none & \texttt{C-x d} & \texttt{rimu} & List a folder with rimu, asked for with completion from the folder of the path chosen last in a question for a path, else the current buffer’s: its name, then its entries a line each, folders first.\\
none & \texttt{C-x 5 d} & \texttt{rimu-\codebreak{}other-\codebreak{}frame} & List a folder with rimu in another frame and work in it, the folder asked for as C-x d asks.\\
none & \texttt{C-x C-j} & \texttt{rimu-\codebreak{}jump} & List the directory of the current buffer’s file with rimu, point on the file; a buffer with no file lists the directory it stands in.\\
\bottomrule
\end{xltabular}

\subsection{Moving about}
\label{sec:221}
Moving into a folder shows it in the same list. You do not collect a buffer
for every folder you pass through. The keys are the ones you know from
chapter 2, with a few more for the places a list has.

\begin{longtable}{@{}l l l@{}}
\toprule
\textbf{Modal} & \textbf{Non-modal} & \textbf{What it does}\\
\midrule
\endhead
\texttt{j}, \texttt{k} & \texttt{C-n}, \texttt{C-p} & Down and up a line\\
none & \texttt{l}, \texttt{RET} & Open the entry: a folder in this list, a file\\
none & \texttt{h} & Up to the folder that holds this one\\
\texttt{g g}, \texttt{g e} & \texttt{M-<}, \texttt{M->} & To the top of the list, to the end\\
none & \texttt{>}, \texttt{<} & To the next folder in the list, to the previous\\
none & \texttt{g f} & To an entry, by name, with completion\\
none & \texttt{g h} & To your home folder\\
none & \texttt{g l} & Back to the folder listed before\\
none & \texttt{q} & Leave the list\\
\bottomrule
\end{longtable}

When you go up with \texttt{h}, the cursor lands on the folder you came from. So you
can walk down into a folder and back out again without losing your place.
\texttt{g l} swaps between the last two folders, so pressing it twice returns to
where you started. This helps when you are working in two places.

\texttt{g f} is the quickest way to a name in a long list. Type a few letters of the
name, and \texttt{RET} puts the cursor on it.

\texttt{h j k l} only move in a list. The keys that move by word or select text do
nothing here, and the echo area says so. Each line is one whole entry, so
there is nothing to select inside it.

\texttt{RET} on a file opens it as \texttt{C-x C-f} does (chapter 9). A file that looks
like a program’s data rather than text opens in a safe view. \texttt{C-u RET} opens
any file that way. Chapter 12 explains what a safe view is.

\begin{try}
\textbf{Try it, with a folder of photos}
\begin{enumerate}
\item Press \texttt{C-x d}. The question holds the folder you are in. Type
   \texttt{\textasciitilde{}/Pictures/} after clearing it with \texttt{C-a C-k}, and press \texttt{RET}. Use any
   folder you like if you have no \texttt{Pictures}.
\item Press \texttt{j} (non-modal: \texttt{C-n}) a few times, then \texttt{l} on a folder to go into it.
\item Press \texttt{h}. You are back, with the cursor on the folder you left.
\item Press \texttt{g l} to go into it again, and \texttt{g l} once more to come back.
\item Press \texttt{g f}, type the first letters of a name and press \texttt{RET}.
\item Press \texttt{g h} to see your home folder, and \texttt{q} to leave the list.
\end{enumerate}
\end{try}

\subsubsection{Commands for moving}
\label{sec:220}
\begin{xltabular}{\textwidth}{@{}l l l >{\raggedright\arraybackslash\hsize=1.000\hsize}X@{}}
\toprule
\textbf{Modal} & \textbf{Non-modal} & \textbf{Command} & \textbf{What it does}\\
\midrule
\endhead
none & \texttt{l}, \texttt{RET}, \texttt{C-c l} & \texttt{rimu-\codebreak{}open} & Open the entry on the current line: a directory is listed in its place in this list; a file is visited as C-x C-f visits it, safely where that would be, and with a prefix argument (C-u) safely always.\\
none & \texttt{C-o} & \texttt{rimu-\codebreak{}display} & Show the entry on the current line in another window, staying in the list: a directory in a list of its own, a file as l opens it.\\
none & \texttt{h}, \texttt{C-c h} & \texttt{rimu-up} & List the directory holding this one, in this list, point on the one left; in a list of files found, list the directory they were found in.\\
none & \texttt{>}, \texttt{C-c >} & \texttt{rimu-\codebreak{}next-\codebreak{}directory} & Go to the next directory of this list.\\
none & \texttt{<}, \texttt{C-c <} & \texttt{rimu-\codebreak{}previous-\codebreak{}directory} & Go to the previous directory of this list.\\
none & \texttt{g f} & \texttt{rimu-\codebreak{}goto-\codebreak{}file} & Go to an entry of this list, its name asked for with completion; among those the search filters leave, while they narrow the list.\\
none & \texttt{g h} & \texttt{rimu-\codebreak{}home} & List your home directory in this list’s place.\\
none & \texttt{g l} & \texttt{rimu-\codebreak{}last-\codebreak{}directory} & List again the directory this list listed before, point on the entry it was on; again, back to this one, as cd - goes back and forth.\\
\bottomrule
\end{xltabular}

\subsection{What a list shows}
\label{sec:225}
At first a list shows only names. You can ask for more, one part at a time.
Each key switches its part on and off for the list you are in, and none of
them changes a file.

\texttt{I} shows details before each name: the same ones that \texttt{ls -l} gives, which
are permissions, the number of links, the owner, the group, the size and the
time of the last change. The inode number, which says where the file sits on
the disk, comes first. The line under the folder’s name then gives how much
free space the drive has. Here is a folder with details on.

\begin{figure}[H]
\caption{Figure 10.2. The same folder with details shown.}\label{org222}
\begin{Verbatim}[breaklines,breakautoindent=false,breakindent=1em,breakafter=/-.\,;:=,commandchars=\\\{\}]
  ~/Photos/holiday:
  total used in directory 1.9M available 59G
  1282547 drwxr-xr-x. 2 aroha aroha  100 Jan  4 14:11 Napier
  1282558 -rw-r--r--. 1 aroha aroha 640K Jan  4 14:20 IMG_0412.jpg
  1282559 -rw-r--r--. 1 aroha aroha 610K Jan  4 14:21 IMG_0413.jpg
  1282560 -rw-r--r--. 1 aroha aroha  31K Jan  5 09:02 itinerary.txt
\end{Verbatim}
\end{figure}

The permissions say who may read, write and run a file. The first letter is
\texttt{d} for a folder and \texttt{-} for a file. The next nine are three groups of three,
for the owner, for the owner’s group and for everyone else. In each group,
\texttt{r} means may read, \texttt{w} may write and \texttt{x} may run or, for a folder, may go
into it. So \texttt{-rw-\codebreak{}r--\codebreak{}r--} is a file its owner may read and change and everyone
else may only read. Rimu shows these but does not change them. If you would
rather read them as numbers, as the \texttt{chmod} command does, set
\texttt{*rimu-\codebreak{}permissions*} to \texttt{:octal}.

\texttt{H} shows the names that start with a dot. \texttt{s} swaps the order between by
name and by time, with the most recently changed first. The mode line says
which one is in force. \texttt{S} shows a size before each name. For a folder it is
the size of everything inside, and it is worked out in the background, so the
column fills in after a moment while you carry on. \texttt{G} shows, after each
name, when the file was last committed to git and the first line of what the
commit said. Chapter 15 describes git.

The standard for each of these is off, with folders first and the order by
name. The settings below change the standard for every list.

\begin{try}
\textbf{Try it, with a folder of letters}
\begin{enumerate}
\item Open a folder of letters with \texttt{C-x d}. If you have none, make one in a
   terminal first, or use any folder with several files.
\item Press \texttt{I}. Details appear. Press \texttt{I} again to hide them.
\item Press \texttt{s}. The list is now ordered by time, the newest first. Press \texttt{s}
   again to go back to names.
\item Press \texttt{S} and wait a moment for the sizes. Find the largest one.
\item Press \texttt{H} to see the hidden names, and \texttt{H} again to hide them.
\end{enumerate}
\end{try}

\subsubsection{Commands for what a list shows}
\label{sec:223}
\begin{xltabular}{\textwidth}{@{}l l l >{\raggedright\arraybackslash\hsize=1.000\hsize}X@{}}
\toprule
\textbf{Modal} & \textbf{Non-modal} & \textbf{Command} & \textbf{What it does}\\
\midrule
\endhead
none & \texttt{I}, \texttt{C-c I} & \texttt{rimu-\codebreak{}toggle-\codebreak{}details} & Show or hide the details before each name in this list: its inode number, as ls -i shows it, then its permissions, links, owner, group, size and when it last changed, as ls -l shows them, and where a symbolic link leads.\\
none & \texttt{H}, \texttt{C-c H} & \texttt{rimu-\codebreak{}toggle-\codebreak{}dotfiles} & Show or hide the entries of this list whose names start with a dot, .\\
none & \texttt{s} & \texttt{rimu-\codebreak{}toggle-\codebreak{}sort} & Order this list by name or by time, those changed last first; the mode line says which.\\
none & \texttt{S}, \texttt{C-c S} & \texttt{rimu-\codebreak{}toggle-\codebreak{}sizes} & Show or hide, before each name in this list, how much it takes: a file its size, a directory the size of everything in it and under it, hidden files too, folders of another drive mounted in it left out, worked out in the background while the list shows it.\\
none & \texttt{G}, \texttt{C-c G} & \texttt{rimu-\codebreak{}toggle-\codebreak{}git} & Show or hide, after each name in this list, the entry’s last commit in git’s history: when it was made and the first line of its message.\\
\bottomrule
\end{xltabular}

\subsubsection{Settings for what a list shows}
\label{sec:224}
\begin{xltabular}{\textwidth}{@{}l >{\raggedright\arraybackslash\hsize=0.571\hsize}X >{\raggedright\arraybackslash\hsize=1.429\hsize}X@{}}
\toprule
\textbf{Setting} & \textbf{Standard} & \textbf{What it changes}\\
\midrule
\endhead
\texttt{*rimu-\codebreak{}details*} & \texttt{nil} & True for rimu to show the details of each entry before its name: its inode number, as ls -i does, then as ls -l does its permissions, links, owner, group, size and when it last changed, and after the name where a symbolic link leads.\\
\texttt{*rimu-\codebreak{}permissions*} & \texttt{:symbolic} & How rimu’s details write each entry’s permissions: :SYMBOLIC, as ls -l does, drwxr-xr-x; :OCTAL, its kind’s letter and the number chmod takes, d755, -644, -4755 with the set-user, set-group or sticky bit.\\
\texttt{*rimu-\codebreak{}dotfiles*} & \texttt{nil} & True for rimu to show the entries whose names start with a dot, .\\
\texttt{*rimu-\codebreak{}sort*} & \texttt{:name} & The order rimu lists entries in: :NAME, by name, octet by octet as ls orders names in the C locale (capitals before small letters); :TIME, those changed last first.\\
\texttt{*rimu-\codebreak{}directories-\codebreak{}first*} & \texttt{t} & True for rimu to list directories, and the symbolic links that lead to directories, before everything else, as ls --group-directories-first does.\\
\texttt{*rimu-\codebreak{}directory-\codebreak{}sizes*} & \texttt{nil} & True for rimu’s lists to show each entry’s size before its name, a directory’s with everything under it; NIL, as standard, to show them only on M-x rimu-toggle-sizes in a list.\\
\texttt{*rimu-\codebreak{}git-\codebreak{}info*} & \texttt{nil} & True for rimu’s lists to show, after each name, the entry’s last commit in git’s history: when it was made and the first line of its message.\\
\texttt{*rimu-\codebreak{}hidden-\codebreak{}columns*} & \texttt{nil} & The columns rimu’s lists never show, however they are asked for: among the details :INODE, :PERMISSIONS, :LINKS, :OWNER, :GROUP, :SIZE and :TIME, and the columns of their own :SIZES (directories’ sizes) and :GIT (the last commits).\\
\texttt{*rimu-\codebreak{}ignored-\codebreak{}extensions*} & \texttt{(".o" ".so" ".a" ".fasl" ".elc" ".pyc" "...} & Endings of names rimu shows set back, in the :list-dim face: compiled files, backups (\textasciitilde{}) and the files some editors save as you type (\#).\\
\texttt{*rimu-\codebreak{}follow-\codebreak{}changes*} & \texttt{t} & True for rimu’s lists to be read again by themselves when their directory changes on disk: a file made, deleted, renamed or finished being written by another program, many at once as one; NIL to read one again only with M-x revert-buffer.\\
\texttt{*rimu-\codebreak{}reuse-\codebreak{}lists*} & \texttt{t} & True for a folder a bookmark goes to (M-x bookmark-jump), or a list goes to across a project’s edge, to be listed in a rimu list there is of its project, moved there, the one used last (for a folder in no project, one in none): a list of its own only when there is none.\\
\bottomrule
\end{xltabular}

\subsection{The preview beside the list}
\label{sec:229}
When you rest the cursor on an entry, Tuatara shows it in a window beside the
list: the start of a file, or for a folder the list of what is inside. It
follows the cursor as you move, so you can look through a folder of notes or
photos without opening any of them. On a wide screen the preview stands to
the right of the list, and on a narrow one below it. It closes when you leave
the list.

A text file shows its text. A picture shows the picture itself in the window
(\texttt{tuatara --gui}) and in terminals that can draw pictures, such as kitty. In
other terminals you see a line that describes it, such as its kind and
size. A PDF shows its first page. \texttt{T} switches between a file’s picture and
its text. For a PDF the text is the words on all of its pages, which is a
quick way to read one. Binary files show in the hex view, as bytes. The
preview shows only the first megabyte of a file, so a very large file does
not slow the list down.

\texttt{P} turns the preview off and on in the list you are in. If it gets in your
way, \texttt{P} is the quickest answer. To turn it off for every list, set
\texttt{*rimu-\codebreak{}preview*} to \texttt{nil} in your init file (chapter 20).

\texttt{C-o} shows the entry in another window and leaves the cursor in the list.
With the preview off it opens a window of its own for the purpose.

\subsubsection{Quick edit}
\label{sec:226}
Sometimes you only want to fix a word in a file. \texttt{o} turns the preview into
the file itself and puts the cursor in it. You can type there as in any
buffer. When you are done, press \texttt{q} in NORMAL state, or \texttt{C-c C-c} in any
state, and Tuatara asks whether to save what you changed. Then the preview
shows the file again and you are back in the list.

\begin{try}
\textbf{Try it, with a folder of letters}
\begin{enumerate}
\item In your folder of letters, move to a text file. A window opens beside the
   list and shows the start of it. If not, press \texttt{P}.
\item Press \texttt{P} to close the preview, and \texttt{P} again to open it.
\item Press \texttt{o}. The cursor is now in the file.
\item Press \texttt{g e} (non-modal: \texttt{M->}) to go to the end, \texttt{i} to type, and add a line
   such as \texttt{Thank you.} Press \texttt{C-g}.
\item Press \texttt{q}. Tuatara asks \texttt{Save letter.\codebreak{}txt? (y or n)}. Press \texttt{y}.
\item Press \texttt{k} and \texttt{j} (non-modal: \texttt{C-p}, \texttt{C-n}) to move over other entries and
   watch the preview change.
\end{enumerate}
\end{try}

Two more keys help when you are deciding what a file is. \texttt{g i} shows what the
system knows of it: its kind, size, owner, times and checksum. It leaves the
file’s access time as it was. \texttt{O} hands a copy of a picture, a PDF, a sound
or a video to your desktop’s own program, after asking, so you can play a
video that Tuatara cannot.

\subsubsection{Commands for the preview}
\label{sec:227}
\begin{xltabular}{\textwidth}{@{}l >{\raggedright\arraybackslash\hsize=0.444\hsize}X l >{\raggedright\arraybackslash\hsize=1.556\hsize}X@{}}
\toprule
\textbf{Modal} & \textbf{Non-modal} & \textbf{Command} & \textbf{What it does}\\
\midrule
\endhead
none & \texttt{P}, \texttt{C-c P} & \texttt{list-\codebreak{}toggle-\codebreak{}preview} & Turn the preview beside this list on or off: the entry on the current line shown in a window beside the list, on its right when there is room, else below it, following the cursor from line to line.\\
none & \texttt{T}, \texttt{C-c T} & \texttt{rimu-\codebreak{}toggle-\codebreak{}picture} & Switch the preview between the picture of a file and its text: the first page of a PDF or its text, an image or its octets in the hex view.\\
none & \texttt{o}, \texttt{C-c o} & \texttt{rimu-\codebreak{}quick-\codebreak{}edit} & Edit the file on the current line where the preview shows it, the quick edit: the preview becomes the file itself, opened as C-x C-f opens it (safely where that would be, and with C-u safely always), and the cursor goes there.\\
none & \texttt{M-x rimu-\codebreak{}quick-\codebreak{}edit-\codebreak{}finish} & \texttt{rimu-\codebreak{}quick-\codebreak{}edit-\codebreak{}finish} & Go back from the quick edit of this file to the preview beside the rimu list (as the cursor going back to the list does), asking first whether to save the changes made: y saves them; n lets them go (a buffer open before the quick edit keeps them); M-x keyboard-quit goes on with the quick edit.\\
none & \texttt{g i} & \texttt{show-\codebreak{}file-\codebreak{}metadata} & Show what the file shown is and what it says of itself, in a buffer of its own, which q closes: the file the buffer visits, the entry on the line in a list of a directory, the PDF or picture in its view; with C-u, or where there is none, asked for.\\
none & \texttt{O}, \texttt{C-c O} & \texttt{open-\codebreak{}externally} & Open the file shown, in the desktop’s own program for its kind (xdg-open), or the one \textbf{OPEN-EXTERNALLY-PROGRAMS} gives its Multipurpose Internet Mail Extensions (MIME) type, for what Tuatara does not do: printing, forms, playing.\\
\bottomrule
\end{xltabular}

\subsubsection{Settings for the preview}
\label{sec:228}
\begin{xltabular}{\textwidth}{@{}l l >{\raggedright\arraybackslash\hsize=1.000\hsize}X@{}}
\toprule
\textbf{Setting} & \textbf{Standard} & \textbf{What it changes}\\
\midrule
\endhead
\texttt{*rimu-\codebreak{}preview*} & \texttt{t} & True for rimu lists to show the entry on the current line in a window beside the list, once point has rested on it \textbf{RIMU-PREVIEW-DELAY} seconds: the start of a file, a binary one in the hex view; the list of a directory; what a named pipe, a socket or a device is, never opening it.\\
\texttt{*rimu-\codebreak{}preview-\codebreak{}delay*} & \texttt{0.1} & Seconds point rests on an entry of a rimu list before the preview shows it, so that moving over many lines quickly reads none of them, 0.1 as standard.\\
\texttt{*rimu-\codebreak{}preview-\codebreak{}size*} & \texttt{1048576} & The most of a file rimu’s preview shows, in octets from its start: 1 MB by default.\\
\texttt{*remote-\codebreak{}preview-\codebreak{}size-\codebreak{}limit*} & \texttt{20000000} & The most octets a file on another machine may hold to be shown in the preview of its list (M-x list-toggle-preview), 20 MB by default: all of it has to come from the server first.\\
\bottomrule
\end{xltabular}

\subsection{Marks}
\label{sec:232}
Marking works as it does in the buffer list. \texttt{m} puts a \texttt{>} on the entry and
moves down a line, so pressing it again and again marks a run. A number
before it marks that many. \texttt{u} takes the mark off, \texttt{U} takes every mark off,
and \texttt{t} turns the marks about, which is the quick way to say “everything but
these”. \texttt{d} puts a \texttt{D} on an entry instead, which means “delete this one”. It
does nothing until you press \texttt{x}.

With nothing marked, a command acts on the entry the cursor is on. With
entries marked, it acts on them.

A mark stays with its file when you go to another folder. You can mark two
photos in one folder, go to another and mark a third, and then copy all
three at once. This is useful and also something to watch. Under the folder’s
name the list says how many entries are marked in other folders, so you are
never surprised, and the question that a command asks counts all of them.
\texttt{U} takes away every mark in every folder.

\texttt{K} takes the lines of the marked entries out of the list. The files stay as
they are, and only the list shows fewer lines. It is a way to cross off what
you have dealt with and look at what is left. \texttt{r} brings the lines back.

\subsubsection{Filtering a list}
\label{sec:230}
A long list is easier to use when it shows only what you want. \texttt{C-x / s}
asks for a pattern and keeps the lines that contain it. If you type \texttt{jpg},
only the lines with \texttt{jpg} in them remain, and a line at the top says how
many of the entries are shown. A second \texttt{C-x / s} narrows it further.
\texttt{C-x / r} takes every filter away. \texttt{C-x / l} lists the filters in force.
Marks stay on entries that a filter hides, and a question that counts the
entries says how many are hidden, so nothing goes without your knowing.

With the details shown, \texttt{C-x / f} and \texttt{C-x / h} ask for a column first, such
as the owner or the size, and then keep or hide the entries whose value in
that column matches. Chapter 6 describes the pattern language.

\begin{try}
\textbf{Try it, with holiday photos}
\begin{enumerate}
\item Open your folder of holiday photos with \texttt{C-x d}.
\item Press \texttt{m} on three photos you like. Each gets a \texttt{>}.
\item Press \texttt{t}. Now every other entry is marked, and the three are not.
   Press \texttt{t} again to turn them back.
\item Press \texttt{U} to clear every mark.
\item Press \texttt{C-x / s}, type \texttt{jpg} and press \texttt{RET}. Only the pictures stay.
\item Press \texttt{C-x / r} to see the whole folder again.
\end{enumerate}
\end{try}

\subsubsection{Commands for marks and filters}
\label{sec:231}
\begin{xltabular}{\textwidth}{@{}l l l >{\raggedright\arraybackslash\hsize=1.000\hsize}X@{}}
\toprule
\textbf{Modal} & \textbf{Non-modal} & \textbf{Command} & \textbf{What it does}\\
\midrule
\endhead
none & \texttt{m}, \texttt{C-c m} & \texttt{list-\codebreak{}mark-\codebreak{}entry} & Mark the entry on the current line of the list (>), for the list’s commands that act on those marked, and move down; with a count, as many entries, up the list for a negative one.\\
none & \texttt{u}, \texttt{C-c u} & \texttt{list-\codebreak{}unmark-\codebreak{}entry} & Take the mark off the entry on the current line of the list, of any kind, and move down; with a count, off as many entries, up the list for a negative one.\\
none & \texttt{U}, \texttt{C-c U} & \texttt{list-\codebreak{}unmark-\codebreak{}all} & Take every mark off the entries of the list, of any kind; in rimu, those of the files marked in other directories too.\\
none & \texttt{t} & \texttt{list-\codebreak{}toggle-\codebreak{}marks} & Mark the entries the list shows that are not marked and take the mark off those marked; those marked for deletion (D), or given a mark of another kind, stay so, as do the entries a filter or a folded section hides.\\
none & \texttt{d}, \texttt{C-c d} & \texttt{list-\codebreak{}mark-\codebreak{}for-\codebreak{}deletion} & Mark the entry on the current line of the list for deletion (D), for the list’s command that deletes them, and move down; with a count, as many entries, up the list for a negative one.\\
none & \texttt{K}, \texttt{C-c K} & \texttt{rimu-\codebreak{}take-\codebreak{}out-\codebreak{}lines} & Take the lines of the entries marked in this directory out of this list, or with none marked here the line point is on, their marks taken away: the files stay as they are, only the list shows fewer, to look at what is left.\\
none & \texttt{C-c f m} & \texttt{rimu-\codebreak{}mark-\codebreak{}master} & Mark the file on the line of this rimu list as a master (M), a file that the files marked (m) are compared with by the commands that compare files, a photo of the person looked for among photos, and go to the next line; u takes the mark away.\\
none & \texttt{C-x / s} & \texttt{search-\codebreak{}filter-\codebreak{}regex} & Narrow the list here to the entries a regular expression matches, as it is typed: those whose line matches, or what more the list holds of them (a mailbox the headers and contents of its messages).\\
none & \texttt{C-x / l} & \texttt{search-\codebreak{}filter-\codebreak{}list} & List the search filters narrowing the list here (C-x / s) in a buffer of their own, each with how many entries it leaves; there d takes away the filter on the line and D all of them.\\
none & \texttt{D}, \texttt{C-c D}, \texttt{C-x / r} & \texttt{search-\codebreak{}filter-\codebreak{}remove-\codebreak{}all} & Take away every search filter narrowing the list here, or the list the list of filters here is for: the whole list is shown again.\\
none & \texttt{C-x / f} & \texttt{search-\codebreak{}filter-\codebreak{}show-\codebreak{}records} & Show only the records whose line a regular expression matches, hiding the others, as C-x / s does, deleting nothing.\\
none & \texttt{C-x / h} & \texttt{search-\codebreak{}filter-\codebreak{}hide-\codebreak{}records} & Hide the records whose line a regular expression matches, showing the others, deleting nothing: asked for as C-x / f asks, kept with the other filters.\\
\bottomrule
\end{xltabular}

\subsection{Making things}
\label{sec:234}
\texttt{+} makes a new folder in the folder the list shows. Type the name after the
folder that is offered. If you type \texttt{Napier/\codebreak{}Hawke/\codebreak{}2026}, Tuatara makes every
folder along the way that is not there yet.

\texttt{a} touches a name. The word comes from the \texttt{touch} command. A name that does
not exist becomes a new empty file, with the cursor on it so that \texttt{l} opens
it. A name that exists has its times set to now and nothing else changes. If
you press \texttt{RET} at the question without typing, the marked entries are
touched, or the one under the cursor.

\texttt{A} makes a file with a formatted name, such as one that starts with the date
and time and has your keywords at its end. It asks for a title and
keywords, and writes the header of the file for you.

\subsubsection{Commands for making things}
\label{sec:233}
\begin{xltabular}{\textwidth}{@{}l l l >{\raggedright\arraybackslash\hsize=1.000\hsize}X@{}}
\toprule
\textbf{Modal} & \textbf{Non-modal} & \textbf{Command} & \textbf{What it does}\\
\midrule
\endhead
none & \texttt{+}, \texttt{C-c +} & \texttt{rimu-\codebreak{}make-\codebreak{}directory} & Make a directory, asked for, in this list’s directory unless given elsewhere, with those above it that are not there; the list shows it.\\
none & \texttt{a}, \texttt{C-c a} & \texttt{rimu-\codebreak{}touch} & Touch a name asked for in this list’s directory, as touch does: a name not there is made a new, empty file, the list on it for l to open it; a file or directory there has its times set to now (a link’s own, not those of where it leads).\\
none & \texttt{A}, \texttt{C-c A} & \texttt{rimu-\codebreak{}make-\codebreak{}formatted-\codebreak{}file} & Make a new file with a formatted name in this list’s directory, asking what new-file-formatted asks but the subdirectory (the title, the keywords, the kind and the template), with the header of its kind, saved at once; the list shows it, point on it.\\
\bottomrule
\end{xltabular}

\subsection{Copying, moving and linking}
\label{sec:237}
\texttt{C} copies the marked entries, or the one under the cursor. \texttt{M} moves them.
The question asks where to, as \texttt{C-x C-f} asks for a file: you type a path,
and \texttt{TAB} completes it. If another rimu list is open in the next window, its
folder is offered first. That makes it easy to copy from one folder to
another that you can see.

A folder goes with everything in it, however many folders are inside. The
setting \texttt{*folder-\codebreak{}tree-\codebreak{}depth-\codebreak{}limit*} stops this at a very deep nesting, so that
a loop made by mistake cannot run for ever.

\texttt{M} moves only into a folder. To move a file and also give it a new name, use
\texttt{R} and a path, or copy it. When you copy a single entry you may type the
whole new path, name included.

A large copy runs in the background and the editor stays free. \texttt{P} lights at
the front of the mode line while it works. \texttt{C-x j} lists the jobs with how
far each has got, and \texttt{C-x J} stops one, taking away a file only half
copied.

Nothing is replaced without asking. If a file of that name is already there,
Tuatara asks whether to replace it, and the old one goes to the trash first,
so you can still get it back. Tuatara also never puts a file on top of
itself. Asking to copy a link onto the file it leads to is refused.

\texttt{Y} makes a symbolic link, which is a small file that says where another
file is. \texttt{L} makes a hard link, which is a second name for the very same
file. A change made through one name shows under the other, and the file
stays until both names are gone. A hard link can only be made on the same
drive, and never to a folder.

\texttt{=} compares two marked entries and shows what differs, line by line, as git
does. With two folders marked, it compares the files that have the same name.

\begin{try}
\textbf{Try it, with holiday photos}
\begin{enumerate}
\item In your folder of photos, press \texttt{m} on three photos.
\item Press \texttt{M}. At the question type the name of a folder, such as \texttt{Napier/},
   and press \texttt{RET}. If the folder is not there, Tuatara says so and moves
   nothing: make it with \texttt{+} and try again.
\item The three photos leave the list, and Tuatara says \texttt{Moved 3 files}.
\item Press \texttt{g f}, go to \texttt{Napier}, and press \texttt{l}. They are there.
\item Put the cursor on one photo and press \texttt{C}. Copy it to your home folder
   by typing \texttt{\textasciitilde{}/} after clearing the question with \texttt{C-a C-k}.
\item Press \texttt{C-x j} to see the job list. A small copy has finished already, so
   the list says there are no jobs.
\end{enumerate}
\end{try}

\subsubsection{Commands for copying, moving and linking}
\label{sec:235}
\begin{xltabular}{\textwidth}{@{}l l l >{\raggedright\arraybackslash\hsize=1.000\hsize}X@{}}
\toprule
\textbf{Modal} & \textbf{Non-modal} & \textbf{Command} & \textbf{What it does}\\
\midrule
\endhead
none & \texttt{C}, \texttt{C-c C} & \texttt{rimu-\codebreak{}copy} & Copy the marked entries, or the one on the line, to a place asked for: the directory of the other list beside this one offered first, or this one; one entry may be given a new name.\\
none & \texttt{M}, \texttt{C-c M} & \texttt{rimu-\codebreak{}move} & Move the marked entries, or the one on the line, into a directory asked for, navigating to it as C-x C-f does, only directories offered: the directory of the other list beside this one offered first.\\
none & \texttt{Y}, \texttt{C-c Y} & \texttt{rimu-\codebreak{}symlink} & Make symbolic links to the marked entries, or the one on the line, in a place asked for, as C asks: the directory of the other list beside this one offered first; one entry’s link may be given a new name.\\
none & \texttt{L}, \texttt{C-c L} & \texttt{rimu-\codebreak{}hardlink} & Give the marked files, or the one on the line, a second name in a place asked for, as C asks: one file, two names, on the same drive only, and no directory.\\
none & \texttt{=}, \texttt{C-c =} & \texttt{rimu-\codebreak{}diff} & Compare the two entries marked, from every directory they are marked in, the one shown first taken as the old, and show the differences in a buffer of their own, which q closes, as M-x diff-files does: unified, the old’s lines taken out after -, the other’s put in after +; two directories with all in them.\\
\bottomrule
\end{xltabular}

\subsubsection{Settings for copying and moving}
\label{sec:236}
\begin{xltabular}{\textwidth}{@{}l l >{\raggedright\arraybackslash\hsize=1.000\hsize}X@{}}
\toprule
\textbf{Setting} & \textbf{Standard} & \textbf{What it changes}\\
\midrule
\endhead
\texttt{*folder-\codebreak{}tree-\codebreak{}depth-\codebreak{}limit*} & \texttt{200} & The most folders, one inside another, that copying, moving or deleting a folder goes down through.\\
\bottomrule
\end{xltabular}

\subsection{Renaming}
\label{sec:239}
\texttt{R} renames the entry under the cursor in its place. The question offers the
old name, so you can edit it rather than type it all again. A name only: a
slash is refused, because moving is \texttt{M}’s job. A buffer that visits the file
follows it to its new name, and so do the recent files and other lists of the
same folder.

To rename many names at once, press \texttt{C-x C-q}. The names of the folder’s
entries appear as text in a buffer of their own, one to a line. You can now
change them with every key that edits text. You can add a word to the start
of each line with a many-cursor edit (chapter 5), or run a replacement over
all of them (chapter 6). When you press \texttt{C-c C-c}, every line that changed is
renamed. \texttt{C-c C-k} leaves the names as they were.

Each line stands for one file. If you delete a line, add one or join two,
nothing is renamed, and Tuatara tells you why. Before it renames anything it
checks all the names together, so a name used twice, or one that another file
already has, is reported and no file is renamed. You can swap two names.

\begin{try}
\textbf{Try it, with a folder of photos}
\begin{enumerate}
\item Open a folder that holds a few photos, such as \texttt{IMG\_\codebreak{}0412.\codebreak{}jpg}.
\item Press \texttt{C-x C-q}. The names appear as text, with the cursor on the first.
\item Press \texttt{I} to start INSERT at the beginning of the line. Type \texttt{Napier-}
   and press \texttt{C-g}.
\item Press \texttt{j} (non-modal: \texttt{C-n}) and repeat on the next line.
\item Press \texttt{C-c C-c}. Tuatara says how many were renamed and the list shows
   the new names.
\end{enumerate}
\end{try}

\subsubsection{Commands for renaming}
\label{sec:238}
\begin{xltabular}{\textwidth}{@{}l l l >{\raggedright\arraybackslash\hsize=1.000\hsize}X@{}}
\toprule
\textbf{Modal} & \textbf{Non-modal} & \textbf{Command} & \textbf{What it does}\\
\midrule
\endhead
none & \texttt{R}, \texttt{C-c R} & \texttt{rimu-\codebreak{}rename} & Rename the entry on the line, or the one marked, in its place, the new name asked for, the old offered: a name only, no directory; a name there already is refused.\\
none & \texttt{C-x C-q} & \texttt{rimu-\codebreak{}edit-\codebreak{}names} & Edit the names of this list’s entries as text, a line each, in its window, with every key that edits text: C-c C-c renames those whose line changed, C-c C-k leaves them.\\
none & \texttt{C-c C-c} & \texttt{rimu-\codebreak{}names-\codebreak{}finish} & Rename the entries whose line changed, and go back to the list: all or, when a line does not do, none, and why is said.\\
none & \texttt{C-c C-k} & \texttt{rimu-\codebreak{}names-\codebreak{}cancel} & Leave the names as they are, and go back to the list.\\
\bottomrule
\end{xltabular}

\subsection{Deleting, the trash and wiping}
\label{sec:242}
\texttt{D} puts the marked entries, or the one under the cursor, in the trash. This
is the same trash your desktop’s file manager uses, so a file you delete here
can be put back from there too. \texttt{x} does the same for the entries marked with
\texttt{D}. Deleting is a two-step in that case: first mark with \texttt{d}, then press \texttt{x}
when you are sure.

\texttt{C-u D} deletes for good, and asks twice. Nothing deleted this way can come
back. The setting \texttt{*rimu-\codebreak{}delete-\codebreak{}to-\codebreak{}trash*} decides what \texttt{D} does. If you set
it to \texttt{nil}, every deletion is for good, after asking twice.

\texttt{g t} lists the trash, as the buffer list lists buffers: a section for your
own trash and a section for each drive that has one. The thing that was put
in last is at the top, and each line says when it was deleted and where it
was. In the list, \texttt{R} puts the marked things back where they were, \texttt{D}
deletes them for good and \texttt{x} deletes those marked with \texttt{D}. \texttt{X} empties the
whole trash and asks twice. \texttt{RET} opens a thing safely so that you can look
at it before you decide. \texttt{h} or \texttt{q} returns to your folder.

A file that you trash on a USB stick goes to the stick’s own trash and comes
back to where it was. That is why the trash list has a section for each
drive. When you pull the stick out, its trash goes with it.

The trash only moves a file. The bytes stay on the drive until the space is
used again. If you really want a file’s contents gone, \texttt{g w} wipes the marked
files: it writes over them and then deletes them. This is slow, and it asks
once.

\begin{try}
\textbf{Try it, with some letters}
\begin{enumerate}
\item Make two new files in a folder with \texttt{a}: type \texttt{old-\codebreak{}1.\codebreak{}txt} and \texttt{RET}, then
   \texttt{old-\codebreak{}2.\codebreak{}txt} and \texttt{RET}.
\item Put the cursor on \texttt{old-\codebreak{}1.\codebreak{}txt}. Press \texttt{d}, which marks it and moves down,
   then \texttt{d} on the next, so both show \texttt{D}.
\item Press \texttt{x} and answer \texttt{y}. Both are in the trash.
\item Press \texttt{g t}. The two are at the top of the trash list.
\item Put the cursor on \texttt{old-\codebreak{}1.\codebreak{}txt} and press \texttt{R}. Answer \texttt{y}. It is back.
\item Press \texttt{q} to return. \texttt{old-\codebreak{}1.\codebreak{}txt} is in the list again.
\end{enumerate}
\end{try}

Do not press \texttt{X} while you practise. It empties the whole trash, not only
what you put there.

\subsubsection{Commands for deleting and the trash}
\label{sec:240}
\begin{xltabular}{\textwidth}{@{}l l l >{\raggedright\arraybackslash\hsize=1.000\hsize}X@{}}
\toprule
\textbf{Modal} & \textbf{Non-modal} & \textbf{Command} & \textbf{What it does}\\
\midrule
\endhead
none & \texttt{x} & \texttt{list-\codebreak{}execute} & Delete the entries of the list marked for deletion (D), after asking, as the list deletes: the buffer list kills its buffers, rimu puts files in the trash, the list of jobs drops those waiting and stops those running.\\
none & \texttt{D}, \texttt{C-c D} & \texttt{list-\codebreak{}delete-\codebreak{}marked} & Delete the marked entries of the list (>), or with none marked the one on the current line, after asking, as the list deletes, as M-x list-execute does those marked for deletion; in rimu, with C-u, for good rather than to the trash.\\
none & \texttt{g t} & \texttt{rimu-\codebreak{}trash} & List what is in the trash, as the buffer list (C-x C-b) lists buffers: a section for your own trash and one for each drive’s, the last to go there first, each thing’s name, when it went and where it was.\\
none & \texttt{R}, \texttt{C-c R} & \texttt{rimu-\codebreak{}trash-\codebreak{}restore} & Put the marked things (>), or with none marked the one on the current line, back where they were, after asking, where they go said; another name asked for when one’s is taken, the directories above it made again when they are gone.\\
none & \texttt{X}, \texttt{C-c X} & \texttt{rimu-\codebreak{}trash-\codebreak{}empty} & Empty the trash: everything the list of the trash shows deleted for good, after asking twice, in the background.\\
none & \texttt{g w} & \texttt{rimu-\codebreak{}wipe-\codebreak{}files} & Wipe the files marked (or the one on the line) as M-x wipe-file does, asked once: overwritten, then deleted.\\
\bottomrule
\end{xltabular}

\subsubsection{Settings for deleting}
\label{sec:241}
\begin{xltabular}{\textwidth}{@{}l l >{\raggedright\arraybackslash\hsize=1.000\hsize}X@{}}
\toprule
\textbf{Setting} & \textbf{Standard} & \textbf{What it changes}\\
\midrule
\endhead
\texttt{*rimu-\codebreak{}delete-\codebreak{}to-\codebreak{}trash*} & \texttt{t} & True for rimu’s deletions (M-x list-delete-marked, M-x list-execute) to move files and directories to the trash the desktop’s file managers use, from which they can be put back; NIL to delete them for good, after asking twice.\\
\bottomrule
\end{xltabular}

\subsection{Archives}
\label{sec:245}
An archive is one file that holds many files, often squeezed to take less
room. A zip of recipes that a friend sent you is an archive. Rimu works with
the common kinds: tar archives (plain, or ending \texttt{.gz}, \texttt{.bz2}, \texttt{.xz} or
\texttt{.zst}), zip, 7z and rar, ISO disc images and cpio. It tells the kind from
the first bytes of the file, so an archive with a wrong name still works.
Word, EPUB and Java files are zip archives and open as such.

\texttt{C-c e} shows what an archive holds, in a buffer of its own, without
unpacking it. \texttt{q} closes it. It is a way to look inside a zip before you
decide whether to unpack it.

\texttt{Z} extracts the archives marked, or the one under the cursor. It asks for a
folder, and offers one named after the archive, which it makes if it is not
there. Several marked archives go into the one folder you give. If an
archive is locked with a password, Tuatara asks for it and types nothing on
the screen. The password reaches the unpacking program on its standard
input, never on its command line, and a wrong one unpacks nothing.

\texttt{c} goes the other way. It packs the marked entries into one tar archive.
The question asks for a name, and its ending says how the archive is
squeezed: \texttt{.tar.gz}, \texttt{.tar.\codebreak{}bz2}, \texttt{.tar.xz}, \texttt{.tar.\codebreak{}zst}, or \texttt{.tar} for none.
Another ending is refused. Entries marked in several folders go into one
archive and keep their paths below the folder that holds them all.

Unpacking is careful because an archive can hold anything. It happens in the
background and shut in a sandbox. Nothing is written outside the folder you
chose, whatever names are inside, and a file already there is left as it is.
Everything is first unpacked into a new folder of its own and moved into
place after, so a half-finished unpack never damages what you had. Unpacking
stops before the drive has less than \texttt{*archive-\codebreak{}free-\codebreak{}reserve*} free, which is
a gigabyte unless you change it, so that a small archive that holds a great
deal cannot fill the drive. Archives inside an archive stay packed.

\begin{try}
\textbf{Try it, with a zip of recipes}
\begin{enumerate}
\item Put a zip file called \texttt{recipes.\codebreak{}zip} in a folder, or pack one: mark a few
   text files and press \texttt{c}, type \texttt{recipes.\codebreak{}tar.gz} after the folder offered
   and press \texttt{C-j}.
\item Press \texttt{U} to clear the marks, as \texttt{Z} would work on the marked files. Put
   the cursor on the archive and press \texttt{C-c e}. A list of what it holds opens.
\item Press \texttt{q} to close it.
\item Press \texttt{Z}. The question offers a folder named \texttt{recipes}. Press \texttt{C-j} to
   take it.
\item Tuatara says \texttt{Extracted 1 archive}. Press \texttt{l} on \texttt{recipes} and see the
   files.
\end{enumerate}
\end{try}

\subsubsection{Commands for archives}
\label{sec:243}
\begin{xltabular}{\textwidth}{@{}l l l >{\raggedright\arraybackslash\hsize=1.000\hsize}X@{}}
\toprule
\textbf{Modal} & \textbf{Non-modal} & \textbf{Command} & \textbf{What it does}\\
\midrule
\endhead
none & \texttt{c}, \texttt{C-c c} & \texttt{rimu-\codebreak{}archive-\codebreak{}create} & Pack the marked entries, from every directory they are marked in, or the one on the line, into one tar archive asked for, navigating to it as C-x C-f does: the directory of the other list beside this one offered first.\\
none & \texttt{Z}, \texttt{C-c Z} & \texttt{rimu-\codebreak{}archive-\codebreak{}extract} & Extract the archives marked, from every directory they are marked in, or the one on the line, into a directory asked for, navigating to it as C-x C-f does, made when it is not there: for one archive, a directory of its name beside it offered first.\\
none & \texttt{C-c e} & \texttt{rimu-\codebreak{}archive-\codebreak{}list} & Show what the archive shown holds, in a buffer of its own, which q closes: the entry on the line in a list of a directory (C-c e there); with C-u, or where there is none, asked for.\\
\bottomrule
\end{xltabular}

\subsubsection{Settings for archives}
\label{sec:244}
\begin{xltabular}{\textwidth}{@{}l l >{\raggedright\arraybackslash\hsize=1.000\hsize}X@{}}
\toprule
\textbf{Setting} & \textbf{Standard} & \textbf{What it changes}\\
\midrule
\endhead
\texttt{*archive-\codebreak{}free-\codebreak{}reserve*} & \texttt{1073741824} & The octets left free on the drive an archive is unpacked onto (1 GB): unpacking stops before the drive has less, so that an archive holding far more than it takes up cannot fill the drive.\\
\texttt{*archive-\codebreak{}listing-\codebreak{}limit*} & \texttt{67108864} & The most octets of an archive’s list of contents shown (64 MB); a longer one is not shown, and that is said.\\
\texttt{*archive-\codebreak{}seconds*} & \texttt{86400} & The seconds tar, 7-Zip or bsdtar may take to make, list or unpack one archive before it is ended: a day, for the largest; M-x stop-job ends it sooner.\\
\bottomrule
\end{xltabular}

\subsection{Drives and the folder tree}
\label{sec:246}
A drive is a disk, a USB stick or a network share that has been attached to
the computer, which is called mounting it. Rimu does not mount drives. Your
desktop does that when you plug one in, and then you list the drive’s folder
like any other.

Rimu knows about drives in a few places. A size worked out with \texttt{S} counts
what is inside a folder on the same drive only, so a drive attached inside it
is not added in. A hard link works only on one drive. The trash of a drive is
kept on that drive, as the last section explained.

A folder tree is a folder with its folders and theirs, as deep as they go.
\texttt{S} adds up a whole tree. Copying, moving or deleting a folder goes through
the entire tree. In a network or a mounted drive that is slow, but the
editor stays usable, because the work is done in the background.

\subsection{Finding, checking and encrypting}
\label{sec:248}
A list can also be made from a search. The commands below make a list of
results, and you work on that list as on any other: mark, copy, move, delete.

\texttt{M-x rimu-\codebreak{}find} lists the files whose names match a pattern, where \texttt{*} stands
for any characters and \texttt{?} for one. \texttt{*.pdf} finds every PDF below the folder
you give. \texttt{M-x rimu-\codebreak{}find-\codebreak{}text} lists the files whose text holds the words you
type. \texttt{M-x rimu-\codebreak{}duplicates} lists files that are the same as another, in
groups, so you can mark one of each pair for deletion.

The keys that start with \texttt{:} encrypt, decrypt and sign. \texttt{: e} and \texttt{: d}
encrypt a marked file into \texttt{FILE.\codebreak{}gpg} and decrypt it back. \texttt{: s} signs a
file and \texttt{: v} checks a signature.

\subsubsection{Commands for finding, checking and encrypting}
\label{sec:247}
\begin{xltabular}{\textwidth}{@{}l l l >{\raggedright\arraybackslash\hsize=1.000\hsize}X@{}}
\toprule
\textbf{Modal} & \textbf{Non-modal} & \textbf{Command} & \textbf{What it does}\\
\midrule
\endhead
none & \texttt{M-x rimu-\codebreak{}find} & \texttt{rimu-\codebreak{}find} & List in a rimu list the files and directories under a directory asked for whose names match a pattern: * stands for any characters, ? for one (\textbf{.pdf, report-202?.}); upper and lower case alike unless the pattern has an upper-case letter.\\
none & \texttt{M-x rimu-\codebreak{}find-\codebreak{}text} & \texttt{rimu-\codebreak{}find-\codebreak{}text} & List in a rimu list the files under a directory asked for whose text matches what you type, as M-s g matches a line: each word a regular expression or its own characters, all on one line.\\
none & \texttt{M-x rimu-\codebreak{}duplicates} & \texttt{rimu-\codebreak{}duplicates} & List in a rimu list the files that are the same as another, in groups: in the directories marked in this rimu list, or else under a directory asked for.\\
none & \texttt{: e} & \texttt{rimu-\codebreak{}encrypt-\codebreak{}files} & Encrypt the files marked (or the one on the line), each into FILE.gpg, to the keys or with a passphrase chosen in the minibuffer, asked once; and, when you say so, asked once too, wipe each file once it is encrypted (M-x wipe-file): with a passphrase, only after it, asked again, opens what was made.\\
none & \texttt{: d} & \texttt{rimu-\codebreak{}decrypt-\codebreak{}files} & Decrypt the files marked (or the one on the line), each into its name without .gpg, for you alone, after asking once: their text is then on disk.\\
none & \texttt{: s} & \texttt{rimu-\codebreak{}sign-\codebreak{}files} & Sign the files marked (or the one on the line) with your key, the kind of signature asked once.\\
none & \texttt{: v} & \texttt{rimu-\codebreak{}verify-\codebreak{}files} & Check the signatures of the files marked (or the one on the line), each said in turn and kept in \textbf{Messages}.\\
\bottomrule
\end{xltabular}

\subsection{Colours of a list}
\label{sec:249}
A list uses faces to tell its parts apart: marked entries, entries marked for
deletion, folders, sizes, dates and the rest. Links that lead nowhere have a
face of their own, so a broken one stands out. Chapter 20 shows how to change
a face.

\begin{xltabular}{\textwidth}{@{}l >{\raggedright\arraybackslash\hsize=1.000\hsize}X@{}}
\toprule
\textbf{Face} & \textbf{What it paints}\\
\midrule
\endhead
\texttt{:list-\codebreak{}marked} & An entry marked (>) in any list whose entries are marked, rimu’s too, for the commands that act on those marked: the mark and the name, in the colour of warnings.\\
\texttt{:list-\codebreak{}marked-\codebreak{}other} & An entry given a mark of another kind in a rimu list (a master’s M, RIMU-MARK-MASTER): the mark and the name.\\
\texttt{:list-\codebreak{}deletion} & An entry marked for deletion (D) in any list whose entries are marked, rimu’s too: the mark and the name, in the colour of errors.\\
\texttt{:list-\codebreak{}directory} & The names of folders in a list: in rimu (C-x d) and where a link to a folder leads, in the trash (M-x rimu-trash) and among bookmarks (C-x m l).\\
\texttt{:list-\codebreak{}size} & Sizes in a list or beside completion candidates: of files and buffers, and a rimu list’s sizes column (S).\\
\texttt{:list-\codebreak{}date} & Dates and times in a list or beside completion candidates: when a file last changed, the date in a formatted name, a rimu list’s last commits (G), when a message was sent.\\
\texttt{:list-\codebreak{}note} & What is said of an entry beside it in a list: the first line of a rimu list’s last commit messages, the first sentence of a face’s description in M-x list-faces, the details under a finding in a report.\\
\texttt{:list-\codebreak{}filters} & The line above a list narrowed by search filters (C-x / s), saying by what and how many entries are left.\\
\texttt{:list-\codebreak{}dim} & What is of less interest in a list, set back: help buffers and the headings of clients that ended in the buffer list (C-x C-b), packages not loaded, bookmarks whose file is gone, files ignored in rimu and in a git status, a value missing in the health report (M-x check-health), a line saying more are not shown.\\
\texttt{:rimu-\codebreak{}symlink} & The names of symbolic links in a rimu list.\\
\texttt{:rimu-\codebreak{}broken-\codebreak{}symlink} & Symbolic links in a rimu list that lead nowhere, and what they hold.\\
\texttt{:rimu-\codebreak{}special} & The names of named pipes, sockets and devices in a rimu list.\\
\texttt{:rimu-\codebreak{}perm-\codebreak{}write} & The permission to write (w) of the group and of others, among the details of an entry in a rimu list.\\
\texttt{:rimu-\codebreak{}set-id} & The set-user, set-group and sticky letters (s, S, t, T) among the permissions of an entry in a rimu list.\\
\bottomrule
\end{xltabular}

\section{Large files and work in the background}
\label{sec:278}
Opening a file of many gigabytes, and working on while Tuatara is busy.

Most files are small. A letter is a few thousand characters, and even a long
book is a few megabytes. Now and then you meet a file of another kind: a year
of a shop’s sales with a line for every till receipt, or the log a web server
has written for months. Such a file can hold a hundred million lines and take
more room than your computer has memory. Many editors refuse it, or make you
wait a minute before the first line shows.

Tuatara opens a file of any size at once and lets you work in it straight
away. This chapter explains how it does that and what you see on the screen
meanwhile. It then describes jobs, which is the name for work that takes a
while and is done in the background, such as replacing a word in six million
lines. You will learn how to see which jobs are running, how to stop one, and
what the letter P at the front of the mode line means. The last parts are about
saving a large file, about memory, and about a large file that lives on
another machine.

One rule holds through all of it. Tuatara never makes you wait. Whatever it is
doing, the next key you press is answered at once.

Every key in this chapter is pressed in NORMAL state unless it starts with
\texttt{C-x} or \texttt{M-}, which work in both states. \texttt{i} starts INSERT state, and \texttt{C-g}
or \texttt{ESC} returns to NORMAL state.

\subsection{Opening a very large file}
\label{sec:255}
You open a large file as you open any other, with \texttt{C-x C-f} and its name. The
first lines are on the screen at once, whether the file is 100 megabytes or
100 gigabytes.

The reason is that Tuatara does not read the file before showing it. A file
of up to 64 megabytes is read whole into memory, which takes no time you
would notice. A larger file stays on the disk, and Tuatara reads only the
part your window shows, with a little before and after it so that paging
down does not wait for the disk. The start and the end of the file are kept
ready too, so \texttt{g g} and \texttt{g e} (non-modal: \texttt{M-<}, \texttt{M->}) are instant. Figure 11.1
shows the idea.

\begin{figure}[H]
\caption{Figure 11.1. Only what you look at is read; the rest stays on disk.}\label{org251}
\begin{Verbatim}[breaklines,breakautoindent=false,breakindent=1em,breakafter=/-.\,;:=,commandchars=\\\{\}]
   sales.csv on the disk, 12 GB
  ┌────────┬──────\discretionary{}{}{}────────────────\discretionary{}{}{}─┬────────┬─────\discretionary{}{}{}────────────┬───\discretionary{}{}{}─────┐
  │ start  │      not read yet     │ window │   not read yet  │  end   │
  │ ready  │                       │ shown  │                 │ ready  │
  └────────┴──────\discretionary{}{}{}────────────────\discretionary{}{}{}─┴────────┴─────\discretionary{}{}{}────────────┴───\discretionary{}{}{}─────┘
   ────────────────\discretionary{}{}{}───►
   the lines are counted from the start, in the background
\end{Verbatim}
\end{figure}

You can type, move, search and save from the first moment. Nothing waits
until the whole file has been looked at.

\subsubsection{The line number while the lines are counted}
\label{sec:252}
To say which line the cursor is on, Tuatara has to know how many line ends
come before it. In a small file that is known as the file opens. In a large
file it would mean reading everything first, so Tuatara counts the lines in
the background while you work.

The mode line tells you about it. Where the line number normally stands, after
the buffer’s name, you may see \texttt{L…} with three dots instead of a number. To the
right, after the mode’s name, it says \texttt{Counting lines} and a percentage:

\begin{Verbatim}[breaklines,breakautoindent=false,breakindent=1em,breakafter=/-.\,;:=,commandchars=\\\{\}]
 TUATARA[N] P [-] sales.csv  L… C0  CSV  Counting lines 40%
\end{Verbatim}

\texttt{L…} means “the count has not reached this line yet”. It shows only for lines
the count has not passed. Near the start of the file you see the real number
straight away. If you jump to the end with \texttt{g e} (non-modal: \texttt{M->}) just after
opening, you see \texttt{L…} until the count gets there, and then the number takes its
place, such as \texttt{L100,000,001}. When the count is done, \texttt{Counting lines} goes
away.

The column can show \texttt{C…} in the same way. That happens on a single line so
long that measuring it would take time, such as a file of a gigabyte with no
line ends in it at all.

The count keeps up with you. If you type or delete lines while it runs, the
number it ends with is right for the text as you left it.

\subsubsection{Going to a line}
\label{sec:253}
\texttt{:} and \texttt{M-g g} ask for a line number and take you there, as chapter 2
explained. While the lines of a large file are still counted, a line the count
has passed is found at once. For a line further in, the echo area shows how
far the count has got until it reaches your line, and then you are there.
\texttt{C-g} gives up the wait and leaves the count running.

\subsubsection{Line numbers in the gutter}
\label{sec:254}
If you show line numbers beside the text (chapter 1), they follow the same
rule. The gutter shows a number for each line the count has reached.

\begin{try}
\textbf{Try it, with a shop’s sales}
A shop in Motueka keeps every sale as a line of a CSV file: the receipt’s
number, what was sold and how many. This makes a pretend file of six million
sales, about 100 MB.
\begin{enumerate}
\item In a terminal, type this on one line and press \texttt{RET}:
   \texttt{seq 1 6000000 | awk '\{ print \$1 ",kereru," \$1 \% 50 \}' > \textasciitilde{}/sales.\codebreak{}csv}
\item In Tuatara press \texttt{C-x C-f}, type \texttt{\textasciitilde{}/sales.\codebreak{}csv} and press \texttt{RET}. The first
   lines are there at once. The mode line says \texttt{Counting lines} and a
   percentage. A fast computer may finish the count before you can look.
\item Press \texttt{g e} (non-modal: \texttt{M->}) straight away. The last line,
   \texttt{6000000,\codebreak{}kereru,0}, shows at once. If the count is still running, the mode
   line says \texttt{L…}.
\item Without waiting, press \texttt{g g} (non-modal: \texttt{M-<}), then \texttt{i}, type \texttt{note,} and
   press \texttt{C-g}. The word is typed at the start of the first line at once. Press
   \texttt{u} to take it away.
\item Watch the mode line for a few seconds. \texttt{Counting lines} goes. Press
   \texttt{g e} again and the mode line shows the line number in full.
\end{enumerate}
\end{try}

\subsection{Jobs: work in the background}
\label{sec:262}
Some commands have a lot to do. Replacing a word in every line of
\texttt{sales.\codebreak{}csv} means changing six million lines. Sorting them means comparing
each with many others. If Tuatara did such work while you watched, the screen
would freeze until it was done.

Instead, work of that size becomes a job. A job runs in the background. You
give the command, and at once you can move, read, search and type again. When
the job is done, its result is taken into the buffer and the echo area tells
you, for example \texttt{Replaced 6,000,000 occurrences on 6,000,000 lines}.

You do not choose whether a command runs as a job. Tuatara decides by how much
there is to do. In a short letter a replacement is done before you lift your
finger. In a large file the same keys start a job, and the echo area says so:

\begin{Verbatim}[breaklines,breakautoindent=false,breakindent=1em,breakafter=/-.\,;:=,commandchars=\\\{\}]
Replacing in the buffer in the background; carry on (C-x J stops it)
\end{Verbatim}

\subsubsection{What runs as a job}
\label{sec:256}
These are the kinds of work that go to the background when they would take
time:

\begin{itemize}
\item counting the lines of a large file, described above
\item a replacement over a large text, and counting the matches of a search
\item sorting lines, and keeping or deleting the lines that match a pattern
\item a change at every row of a tall rectangle or at a great many cursors
\item a save that takes more than a moment
\item a count of words (\texttt{M-=}) and a checksum of a file (\texttt{M-x sha256-\codebreak{}sum})
\item copying, moving and deleting files from a folder listing (chapter 10)
\item reading a file from another machine
\end{itemize}

Other chapters describe these commands. What they share is told here: each
shows the sign P, each is listed with \texttt{C-x j}, and each can be stopped.

\subsubsection{The sign P}
\label{sec:257}
The mode line starts with a few signs that keep their place, described in
chapter 1. The first one is P, for processing. It belongs to the buffer shown
in that window. While a job runs in the buffer, or waits its turn there, P
is bright. The rest of the time it is faint, so faint that you hardly see it.
A glance at the front of the mode line tells you whether this buffer is still
being worked on.

P is lit too when another program is working for the buffer, such as a
language server reading a project (chapter 13).

The setting \texttt{*processing-\codebreak{}indicator*} holds what is shown. You can put another
character there, or \texttt{nil} to have no sign at all:

\begin{Verbatim}[breaklines,breakautoindent=false,breakindent=1em,breakafter=/-.\,;:=,commandchars=\\\{\}]
\color[HTML]{A85F00}(setf\normalcolor{} \color[HTML]{A85F00}*processing-indicator*\normalcolor{} \color[HTML]{0A7AB3}"*"\color[HTML]{A85F00})
\end{Verbatim}

\subsubsection{What a job will change is painted}
\label{sec:258}
While a job works through the buffer, the text it has still to change is
shown in a colour of its own, the face \texttt{:job-\codebreak{}queued}. As the job moves on, the
painted part shrinks. If a second job waits behind the first, the text it will
change is painted in another colour, \texttt{:job-\codebreak{}waiting}. You can see at a glance
which part of the file is finished and which is not.

For a replacement, the lines in your window are changed first, the moment you
press \texttt{RET}. The rest of the file follows in the background. What you look at
is never out of date.

\subsubsection{Jobs take turns}
\label{sec:259}
A buffer runs one job at a time. If you give a second command while the first
is still at work, the second waits, and the echo area says where it stands:

\begin{Verbatim}[breaklines,breakautoindent=false,breakindent=1em,breakafter=/-.\,;:=,commandchars=\\\{\}]
Queued, first in line: Sorting the buffer (C-x j lists the jobs)
\end{Verbatim}

It starts by itself when the one before it ends. A job that works line by line
also waits while the lines of a newly opened large file are counted. You do
not have to do anything, and you can keep typing while jobs wait.

Jobs in different buffers do not wait for each other. A sort in one file and
a save of another run side by side.

\subsubsection{Typing while a job runs}
\label{sec:260}
You may edit the buffer while a job works on it. What you type stays as you
typed it. A replacement that reaches a line you changed meanwhile leaves your
change alone and says so when it ends.

A sort and a filter are different, because they move whole lines about and
an edit in the middle would spoil the result. If you try to change text they
are working on, Tuatara asks whether to stop the job. Answer \texttt{y} to stop it
and make your change, or \texttt{n} to let it finish.

\subsubsection{Undo and jobs}
\label{sec:261}
Each job is one step to undo, however much it changed. One \texttt{u} takes back a
replacement of six million lines.

The step stands at the place in the history where you gave the command, even
if the job finished much later. Suppose you start a replacement, then start a
second one that has to wait, then type a word. The first \texttt{u} takes back the
word, the next the second replacement, the next the first. \texttt{U} brings them
back in the same order. When a job’s change comes back, the cursor stays
where it is. You are not thrown to the last line the job touched.

\subsection{The list of jobs}
\label{sec:265}
\texttt{C-x j} shows every job in a buffer called \texttt{*Jobs*}. Use it when P is lit and
you wonder what is going on, when a job takes longer than you expected, or
when you gave a command by mistake and want to drop it before its turn comes.

\begin{figure}[H]
\caption{Figure 11.2. The list of jobs: one runs, one waits.}\label{org263}
\begin{Verbatim}[breaklines,breakautoindent=false,breakindent=1em,breakafter=/-.\,;:=,commandchars=\\\{\}]
  Buffer        State                    Job
  ------        -----                    ---
  sales.csv     Replacing 12%            Replacing kereru with kea in the buffer
  sales.csv     waiting, first in line   Sorting the buffer
  access.log    Saving 60%               Saving access.log
\end{Verbatim}
\end{figure}

Each line is a job. The first column names the buffer it works in. The second
says what it is doing and how far it has got, or that it waits and where it
stands in the line. The third says in words what the job is. Jobs are listed
buffer by buffer, the running one first and those waiting below it in the
order they will run. The list keeps itself up to date while it is on the
screen.

The list works like the buffer list of chapter 7. \texttt{d} marks the job on the
cursor’s line with a \texttt{D}, and \texttt{x} carries it out after asking: a job that
waits is dropped, and one that runs is stopped. \texttt{m} marks a job and \texttt{u} takes
the mark off. \texttt{D} drops or stops the marked jobs, or the one on the line if
none is marked. \texttt{q} closes the list.

\begin{xltabular}{\textwidth}{@{}l l l >{\raggedright\arraybackslash\hsize=1.000\hsize}X@{}}
\toprule
\textbf{Modal} & \textbf{Non-modal} & \textbf{Command} & \textbf{What it does}\\
\midrule
\endhead
none & \texttt{C-x j} & \texttt{list-\codebreak{}jobs} & List the jobs in the background in \textbf{Jobs}: in each buffer those running, with how far they have got, and those waiting their turn, in the order they will run, their text painted in the buffer meanwhile.\\
none & \texttt{d}, \texttt{C-c d} & \texttt{list-\codebreak{}mark-\codebreak{}for-\codebreak{}deletion} & Mark the entry on the current line of the list for deletion (D), for the list’s command that deletes them, and move down; with a count, as many entries, up the list for a negative one.\\
none & \texttt{x} & \texttt{list-\codebreak{}execute} & Delete the entries of the list marked for deletion (D), after asking, as the list deletes: the buffer list kills its buffers, rimu puts files in the trash, the list of jobs drops those waiting and stops those running.\\
none & \texttt{m}, \texttt{C-c m} & \texttt{list-\codebreak{}mark-\codebreak{}entry} & Mark the entry on the current line of the list (>), for the list’s commands that act on those marked, and move down; with a count, as many entries, up the list for a negative one.\\
none & \texttt{u}, \texttt{C-c u} & \texttt{list-\codebreak{}unmark-\codebreak{}entry} & Take the mark off the entry on the current line of the list, of any kind, and move down; with a count, off as many entries, up the list for a negative one.\\
none & \texttt{D}, \texttt{C-c D} & \texttt{list-\codebreak{}delete-\codebreak{}marked} & Delete the marked entries of the list (>), or with none marked the one on the current line, after asking, as the list deletes, as M-x list-execute does those marked for deletion; in rimu, with C-u, for good rather than to the trash.\\
\bottomrule
\end{xltabular}

\subsubsection{Stopping a job}
\label{sec:264}
\texttt{C-x J}, with a capital J, stops the job running in the buffer you are in.
It asks first, so a slip of the finger costs nothing. The buffer is left as
it is, and a file that was being saved is left as it was before the save. If
other jobs wait behind the one you stopped, they start one after another. To
get rid of those as well, drop them in the list.

If nothing is running, \texttt{C-x J} says \texttt{Nothing is running in this buffer}.

One kind of save is always finished before it stops: a file being written in
place, described below. Stopping it half-way would leave the file spoiled.

\begin{xltabular}{\textwidth}{@{}l l l >{\raggedright\arraybackslash\hsize=1.000\hsize}X@{}}
\toprule
\textbf{Modal} & \textbf{Non-modal} & \textbf{Command} & \textbf{What it does}\\
\midrule
\endhead
none & \texttt{C-x J} & \texttt{stop-\codebreak{}job} & Stop the work running in the background in the current buffer, after asking.\\
\bottomrule
\end{xltabular}

\begin{try}
\textbf{Try it, with the shop’s sales again}
The shop finds that what it sold as kereru soft toys were kea all along.
\begin{enumerate}
\item Open \texttt{\textasciitilde{}/sales.\codebreak{}csv} from the first Try it and wait until \texttt{Counting lines}
   has gone.
\item Press \texttt{r r}, type \texttt{\%kereru/\codebreak{}kea} and press \texttt{RET}. The lines on the screen
   say \texttt{kea} at once. The echo area says the rest goes on in the background,
   and P lights up at the front of the mode line. A fast computer may end the
   job within a second. Then make the file ten times longer, with 60000000 in
   place of 6000000 in the command, and start again.
\item Press \texttt{C-x j} straight away. The list shows the replacement and how far it
   has got. Press \texttt{q}.
\item Press \texttt{g e} (non-modal: \texttt{M->}), then \texttt{g g} (non-modal: \texttt{M-<}). You move
   freely while the job works.
\item When the echo area says how many it replaced, press \texttt{u}. All six million
   lines say \texttt{kereru} again.
\item Press \texttt{r r}, type \texttt{\%kereru/\codebreak{}kea} and \texttt{RET} once more, and this time press
   \texttt{C-x J} and answer \texttt{y}. The job stops.
\item Press \texttt{C-x k} and answer that the changes may go. Nothing was saved, so
   the file on disk is as you made it.
\end{enumerate}
\end{try}

\begin{xltabular}{\textwidth}{@{}l l >{\raggedright\arraybackslash\hsize=1.000\hsize}X@{}}
\toprule
\textbf{Setting} & \textbf{Standard} & \textbf{What it changes}\\
\midrule
\endhead
\texttt{*processing-\codebreak{}indicator*} & \texttt{"P"} & What the mode line shows first, before [-] and the signs of directories kept read (\textbf{INDEXING-INDICATOR}), always in the same place: in :mode-line-sign while a job runs or waits its turn in the buffer, or another program works for it (a language server indexing), else faint, in :mode-line-sign-idle, so that nothing after it moves as work starts and ends.\\
\bottomrule
\end{xltabular}

\begin{xltabular}{\textwidth}{@{}l >{\raggedright\arraybackslash\hsize=1.000\hsize}X@{}}
\toprule
\textbf{Face} & \textbf{What it paints}\\
\midrule
\endhead
\texttt{:mode-\codebreak{}line-\codebreak{}sign} & The mode line’s signs before [-] while their work goes on: P while a job of the buffer runs or waits its turn (\textbf{processing-indicator}), U while a directory of \textbf{indexed-directories} is being read (\textbf{indexing-indicator}), S while one is searched (\textbf{searching-indicator}), D while the editor runs as a daemon (\textbf{daemon-indicator}); in the theme’s warning colour, to stand out.\\
\texttt{:mode-\codebreak{}line-\codebreak{}sign-\codebreak{}idle} & The mode line’s signs before [-] (P, U and S) while their work is not going on: faint, barely seen.\\
\texttt{:job-\codebreak{}queued} & Text a job in the background (a replacement, a sort, a rectangle’s change) has still to work through (for a replacement only its matches), all of what it changes again while the lines of its new text are counted, until it is taken in; and a job’s state in the list of jobs (C-x j).\\
\texttt{:job-\codebreak{}waiting} & Text a job waiting its turn behind another in the background will change.\\
\bottomrule
\end{xltabular}

\subsection{Searching a large file}
\label{sec:266}
\texttt{C-s} searches a large file the way it searches a small one (chapter 6). You
type, and the first match after the cursor is shown as you type. Matches in
the window are found at once, because that text is already read.

Two things differ, and both keep your typing quick.

The first is the count of matches beside the prompt. Counting every match in a
file of gigabytes takes a while, and doing it again for each letter you type
would be wasted work. So in a large text the count starts only when you stop
typing for a moment, half a second as standard, and it runs in the
background. The setting \texttt{*typing-\codebreak{}pause*} holds how long that moment is. Until
the count is done you can already press \texttt{RET} or go on to the next match.

The second is a match that lies far away, in a part of the file that has not
been read. The search then reads the file through, a piece at a time, in the
background. You can keep typing the pattern or give the search up with
\texttt{C-g} while it looks.

Commands that work on the lines a pattern finds run as jobs in a large file.
\texttt{M-x keep-\codebreak{}lines} keeps only the lines that match and \texttt{M-x flush-\codebreak{}lines} deletes
them, both within the selection, so select the whole buffer first with
\texttt{C-x h}.

\begin{try}
\textbf{Try it, with a huge log file}
A web server in Timaru writes a line to its log for each page it serves. Most
lines are fine. You want only the ones that say \texttt{error}.
\begin{enumerate}
\item In a terminal, make a pretend log of six million lines. Type this on one
   line and press \texttt{RET}:
   \texttt{seq 1 6000000 | awk '\{print \$1, (\$1 \% 50000 ? "ok" : "error")\}' > \textasciitilde{}/access.\codebreak{}log}
\item Open \texttt{\textasciitilde{}/access.\codebreak{}log} with \texttt{C-x C-f}. Press \texttt{C-s} and type \texttt{error}. The first
   match is far down the file. The search looks for it in the background and
   shows it when found. Stop typing, and the count of matches appears beside
   the prompt a moment later. Press \texttt{RET}.
\item Press \texttt{C-x h} to select everything, then \texttt{M-x}, type \texttt{keep-\codebreak{}lines} and
   \texttt{RET}. A log has columns, so the prompt asks \texttt{Filter by the column:}.
   Press \texttt{RET} to take the whole record, type \texttt{error} and \texttt{RET}. P lights up
   while the job works.
\item When it is done, the buffer holds 120 lines, each ending in \texttt{error}.
\item Press \texttt{u}. All six million lines are back.
\item Press \texttt{C-x k} and let the changes go.
\end{enumerate}
\end{try}

\subsection{Saving a large file}
\label{sec:268}
\texttt{C-x C-s} saves a large file as it saves any other. Tuatara waits half a
second for the save, so a quick one says \texttt{Wrote} at once. A save that takes
longer goes on in the background. The mode line shows \texttt{Saving} and a
percentage, P is lit, and you carry on working. A key you press during that
first half second is answered at once too, and the save carries on behind it.
If you leave Tuatara while a save is running, it waits for the save to end.

How long a save takes depends on the kind of drive, because of how Tuatara
keeps your file safe. A save does not write over the old file. It writes a new
file beside it and, when that is complete, puts it in the old one’s place. If
the power fails in the middle, you have either the old file or the new one,
whole. You never have a mix of both.

On some file systems, btrfs and XFS among them, the new file can share with
the old one every part that did not change. A save then takes only as long as
your change, even in a file of 10 GB. Tuatara does not go by the kind of file
system alone: the first time it writes on such a drive, it tries sharing with
two new empty files, which leaves nothing behind, and keeps the answer until
you quit. An older XFS made without reflink shares nothing, so there it saves
as on ext4.

On the others, ext4 among them, a new file means writing every part of it
again. Changing one line of a 10 GB file then costs about a minute of writing.
It happens in the background and you can go on working, but the disk is busy
all that time.

\subsubsection{Saving in place}
\label{sec:267}
If you often save very large files on such a drive, the setting
\texttt{*save-\codebreak{}in-\codebreak{}place*} shortens the wait. Put this in your init file:

\begin{Verbatim}[breaklines,breakautoindent=false,breakindent=1em,breakafter=/-.\,;:=,commandchars=\\\{\}]
\color[HTML]{A85F00}(setf\normalcolor{} \color[HTML]{A85F00}*save-in-place*\normalcolor{} \color[HTML]{A85F00}t)
\end{Verbatim}

Tuatara then writes only what changed, into the file itself, when that comes
to less than about two thirds of the file. That is the case for a change that
keeps the text the same length, and for any change in roughly the last third
of the file. One changed line near the end of a 10 GB file is saved in a
fraction of a second. A change near the start that makes the file longer or
shorter moves everything after it, so it is still saved as a new file, which
is the quicker way for it.

There is a price, and it is the reason the setting is off as standard. Before
it writes into the file, Tuatara notes the old contents of the parts it will
overwrite in a small file beside yours, named \texttt{.NAME.\codebreak{}tuatara-\codebreak{}journal}, and
removes it when the save is complete. If the computer loses power in the middle
of the writing, your file is half saved until Tuatara opens it again. The next
time you open, save or revert it, Tuatara puts the old contents back from the
journal and tells you that a save was cut short. Another program that reads
the file before then, a backup for one, may see it half saved.

Leave the setting off if other programs read your large files, or if you do
not mind the minute.

\begin{xltabular}{\textwidth}{@{}l l >{\raggedright\arraybackslash\hsize=1.000\hsize}X@{}}
\toprule
\textbf{Setting} & \textbf{Standard} & \textbf{What it changes}\\
\midrule
\endhead
\texttt{*save-\codebreak{}in-\codebreak{}place*} & \texttt{nil} & True to save a large file by writing only what changed into the file itself, where that is quicker than writing it whole.\\
\bottomrule
\end{xltabular}

\subsection{Memory, and pauses you might notice}
\label{sec:273}
You do not have to manage memory in Tuatara. This part tells you what to
expect, so that nothing surprises you.

\subsubsection{A file larger than your memory}
\label{sec:269}
A file of 11 GB opens on a laptop with 16 GB of memory, and so does one of
100 GB. Because a large file is used where it lies on the disk, opening it
takes almost no memory.

Work that reads the whole file, such as counting its lines, a search, a
replacement or a checksum, takes the file into memory a piece at a time and
lets each piece go once it is well past it. Your other programs are not pushed
aside while it runs. With a large file open and nothing to do, Tuatara uses no
processor time at all.

\subsubsection{Auto-save and large buffers}
\label{sec:270}
Tuatara keeps a copy of your unsaved changes so that a crash does not cost you
your work (chapter 7). Each such copy is the whole buffer, so copying a huge
one every few seconds would keep the disk busy for nothing. A buffer larger
than \texttt{*auto-\codebreak{}save-\codebreak{}size-\codebreak{}limit*}, 256 MB as standard, is not copied. In a file
that size, save with \texttt{C-x C-s} when you have made changes you want to keep.

\subsubsection{Undo}
\label{sec:271}
The undo history remembers every change so that \texttt{u} can take it back. A
change to all of a large file is remembered in a compact way, but a long day
of such changes adds up. As standard there is no limit. If memory runs
short, \texttt{*undo-\codebreak{}limit*} sets the most memory, in octets, that one buffer’s
history may hold.

\subsubsection{Collecting garbage}
\label{sec:272}
As it works, Tuatara makes short-lived data: the text of a message, the result
of a search. When that data is no longer needed it is garbage, and from time
to time Tuatara stops for a moment to clear it away. The moment is short,
usually too short to see.

Tuatara chooses that moment with care. When you have pressed no key for a
short while, 0.3 seconds as standard, it clears the garbage then, if enough
has gathered. The pause falls while you are thinking and does not come between
two keys. The setting \texttt{*idle-\codebreak{}collection-\codebreak{}seconds*} holds how long you must be
still. There is rarely a reason to change it.

\begin{xltabular}{\textwidth}{@{}l l >{\raggedright\arraybackslash\hsize=1.000\hsize}X@{}}
\toprule
\textbf{Setting} & \textbf{Standard} & \textbf{What it changes}\\
\midrule
\endhead
\texttt{*auto-\codebreak{}save-\codebreak{}size-\codebreak{}limit*} & \texttt{268435456} & The largest buffer auto-saved, in octets, 256 MB as standard: a larger one is not copied (each copy would be written whole), so a crash loses its changes not saved.\\
\texttt{*undo-\codebreak{}limit*} & \texttt{nil} & The most memory, in octets, the undo history of one buffer may hold; NIL, the default, for no limit.\\
\texttt{*idle-\codebreak{}collection-\codebreak{}seconds*} & \texttt{0.3} & Seconds you leave the editor alone, no key typed, after which the garbage made since the last collection is collected, when it is at least half of what makes one (\textbf{NURSERY-OCTETS}): the pause comes while you think, not between two keys.\\
\texttt{*typing-\codebreak{}pause*} & \texttt{1/2} & Seconds without a key typed before work over a text too large to go through at each key starts while a pattern is typed: the count of a search’s matches (M-x regex-search-forward) and a filter being typed (M-x search-filter-regex), its lines in the window shown at once.\\
\bottomrule
\end{xltabular}

\subsection{A large file on another machine}
\label{sec:276}
Chapter 9 showed how to open a file on a server by writing \texttt{/ssh:} and the
machine’s name before its path. A small file comes across whole and then
opens. A large one would keep you waiting for minutes on a slow line, so
Tuatara treats it differently.

A file on another machine that is larger than 64 MB opens as soon as its
first part has come. The rest goes on arriving in the background, and you can
read what is here while it does.

\begin{Verbatim}[breaklines,breakautoindent=false,breakindent=1em,breakafter=/-.\,;:=,commandchars=\\\{\}]
C-x C-f /ssh:akaroa:/var/log/access.log RET
\end{Verbatim}

\subsubsection{What you see while it arrives}
\label{sec:274}
A part of the file that has not come yet is shown as a single row in square
brackets, saying how much is on its way. When the part arrives, the row is
replaced by the text.

The transfer brings what your window looks at first. After that it brings the
start and the end of the file, and then works forward from the start. If you
jump to a place far ahead, that place comes next. You steer the transfer just
by looking.

While any of the file is missing:

\begin{itemize}
\item P is lit, and \texttt{C-x j} lists the transfer with how far it has got.
\item The buffer is read-only. Trying to change it says that the file has not all
  arrived.
\item The lines are not counted yet, so the mode line shows \texttt{L…}.
\item A search, a copy or a checksum that reaches a missing part stops and says
  that part of the file has not arrived. Tuatara never puts made-up text in
  the place of what is missing.
\end{itemize}

When the last part has come, the echo area says \texttt{All of access.\codebreak{}log has
arrived}. From then on the file behaves like any other large file: its lines
are counted, you can change it, and \texttt{C-x C-s} writes it back to the server.

\subsubsection{Stopping the transfer, or losing the connection}
\label{sec:275}
\texttt{C-x J} stops the transfer, after asking. What came stays on the screen and
can be read. The mode line then says \texttt{incomplete}, and the missing parts are
shown as not arrived. The same happens if the connection is lost. An
incomplete buffer is never written back to the server, so a half file cannot
replace a whole one. \texttt{M-x revert-\codebreak{}buffer} fetches the file again from the
start.

Killing the buffer ends the transfer and deletes the copy on your computer.

Each request to the server has \texttt{*remote-\codebreak{}seconds*} to be answered, 30 as
standard. A server that stays silent longer than that ends the transfer, and
the echo area says why.

\begin{try}
\textbf{Try it, with a log on a server}
You need a machine you can reach with ssh that holds a file larger than
64 MB. Here it is called \texttt{akaroa} and the file is \texttt{/var/\codebreak{}log/\codebreak{}access.\codebreak{}log}. Use
your own names.
\begin{enumerate}
\item Press \texttt{C-x C-f}, type \texttt{/ssh:\codebreak{}akaroa:/\codebreak{}var/\codebreak{}log/\codebreak{}access.\codebreak{}log} and press \texttt{RET}.
   The first lines show as soon as they have come. P is lit.
\item Press \texttt{g e} (non-modal: \texttt{M->}). The end of the file is there, since it was
   fetched early.
\item Press \texttt{C-x j}. The list shows the transfer and its percentage. Press \texttt{q}.
\item Press \texttt{i} and type a letter. Tuatara refuses, because the file has not all
   arrived. Press \texttt{C-g}.
\item Wait for the message that all of it has arrived, or press \texttt{C-x J} and
   answer \texttt{y} to stop. After stopping, the mode line says \texttt{incomplete}.
\item Press \texttt{C-x k}. The copy on your computer is deleted with the buffer.
\end{enumerate}
\end{try}

\begin{xltabular}{\textwidth}{@{}l l >{\raggedright\arraybackslash\hsize=1.000\hsize}X@{}}
\toprule
\textbf{Setting} & \textbf{Standard} & \textbf{What it changes}\\
\midrule
\endhead
\texttt{*remote-\codebreak{}seconds*} & \texttt{30} & Seconds a server reached for its files (/ssh:HOST:/path) has to answer each request, and ssh to connect: past them the connection is closed and the work in hand fails, saying so.\\
\bottomrule
\end{xltabular}

\subsection{In short}
\label{sec:277}
\begin{longtable}{@{}l l@{}}
\toprule
\textbf{You want to} & \textbf{Do this}\\
\midrule
\endhead
Open a file of any size & \texttt{C-x C-f}, as for any file\\
Know whether the buffer is busy & Look for a bright P\\
See what is running, and how far & \texttt{C-x j}\\
Stop the job in this buffer & \texttt{C-x J}\\
Drop a job that waits & \texttt{C-x j}, then \texttt{d} and \texttt{x}\\
Take back a whole job & \texttt{u} (non-modal: \texttt{C-/})\\
Save one line of a huge file quickly & \texttt{*save-\codebreak{}in-\codebreak{}place*}, on ext4\\
\bottomrule
\end{longtable}

\section{Files you do not trust}
\label{sec:303}
Looking at a file without letting anything in it happen, and keeping private
text private.

Most files you open come from you. Some do not: an attachment from an email you
were not expecting, a folder a client sent, a download, a disk someone handed
you. You want to read these, and you do not want to find out what they do. This
chapter is about reading a file without running or changing anything. It starts
with the safe way to open one file. It goes on to folders whose files are always
opened that way, to the programs that decode pictures and PDFs behind a wall,
and to git repositories you only look at. The last part is about the opposite
problem: a file that you trust but that nobody else should read, which you keep
encrypted.

Be modest about what this gives you. The safeguards stop Tuatara itself from
acting on a file and from changing it. They do not make a file harmless to the
rest of your computer. If you copy a program out of a folder and run it in a
terminal, no safeguard of an editor helps you. Each section below says what it
covers and what it leaves out.

\subsection{Opening one file safely}
\label{sec:283}
An ordinary open does some work for you. Tuatara guesses the kind of text
(letter, code, notes) and chooses a major mode, which colours the text and
changes some keys. It reads any settings the file gives for itself. It looks for
a git repository around the file. It keeps a copy of your unsaved changes. All
of that is help when the file is yours. For a file from a stranger it is also a
lot of ways for the file to take part.

Opening a file safely turns that work off. The file is shown as plain text and
it is read-only. No major mode and no hook runs on it. Version control is not
looked up. No auto-save copy is made (chapter 7), and the folder is left as it
was. The access time of a file of your own is left as it was too. The system
lets only a file’s owner read it without setting its access time, so reading
another person’s file still sets it. Mount such a drive read-only if that
matters. The mode line says \texttt{Safe} just before the buffer’s name.

Characters that draw nothing, such as a zero-width space, a soft hyphen or a
mark that turns text right to left, are shown by their codes. Control characters
and terminal escape sequences are shown as text, such as \texttt{\textasciicircum{}[}, and do nothing.
That does not close every way of hiding text. Variation selectors, the blank
Hangul and Braille characters and spaces such as the no-break space are drawn as
they are, and a letter of another alphabet that looks like a Latin one is not
marked.

There are four ways to open a file safely.

\begin{itemize}
\item \texttt{C-u C-x C-f}, a count before the ordinary command, asks for a file and opens
  it safely. \texttt{M-x find-\codebreak{}file-\codebreak{}safely} does the same.
\item \texttt{tuatara --safe FILE} from a terminal, or \texttt{-R} for short, opens every file
  named on the line that way. A file piped in with \texttt{ls | tuatara -} is a safe
  buffer too.
\item The setting \texttt{*safe-\codebreak{}directories*} lists folders whose files always open safely,
  however you open them. An example is \texttt{"\textasciitilde{}/\codebreak{}Downloads/"}.
\item A file that looks like a program’s data and not text opens safely by itself.
  The test is a null byte in its first 64 KB. Compiled programs, most archives
  and disk images have one, and so does text in UTF-16. A file with no null byte
  there, such as a script, or one whose first null byte comes later, opens the
  usual way. The setting \texttt{*open-\codebreak{}binary-\codebreak{}files-\codebreak{}safely*} turns this off. Files you
  drop on the window open the usual way unless \texttt{*open-\codebreak{}dropped-\codebreak{}files-\codebreak{}safely*} is
  true.
\end{itemize}

Files you opened safely are remembered. In the list of recent files (chapter 7)
they carry a mark, and \texttt{C-x C-r} opens them safely again.

Safe does not mean fixed. \texttt{C-x C-q} makes a safe buffer editable, for the case
where you really do want to write in it. It is still never auto-saved.
\texttt{M-x end-\codebreak{}safe-\codebreak{}visit} opens the file as usual: it gets its mode, its colours and
version control, and it is editable if you may write the file. Your cursor
position, your unsaved changes and your undo history stay as they are. A file
that is evidence (next section), or was when you opened it, refuses this, and
says that it stays safe. So does a buffer whose file name now leads to another
file than the one it read, as when a link was changed in the meantime.

Two more settings decide what a safe buffer shows.
\texttt{*invisible-\codebreak{}characters-\codebreak{}when-\codebreak{}safe*} is true as it comes, and shows the characters
that draw nothing. And \texttt{*text-\codebreak{}direction-\codebreak{}when-\codebreak{}safe*} is false, so that text
written right to left (Arabic, Hebrew) is shown in the order it is stored and
not the order it is read. The stored order is the one that tells you what is in
the file.

\subsubsection{What is never run}
\label{sec:279}
A safe buffer is for looking. A Lisp file opened safely has no Lisp mode, so
\texttt{C-x C-e} and \texttt{C-M-x} are not defined there, and \texttt{M-x} does not offer
\texttt{eval-\codebreak{}region} or \texttt{eval-\codebreak{}buffer}. If you turn on \texttt{M-x lisp-\codebreak{}mode} yourself, those
keys refuse, even after \texttt{C-x C-q}, and the echo area says that the buffer is
opened safely and that nothing was evaluated. A list of a folder in rimu
(chapter 10) opened safely refuses to pack, unpack, wipe, rename, move or delete
the file on the line.

\subsubsection{Why a file’s own settings are not taken}
\label{sec:280}
Some files carry instructions for the editor in the file itself. A first line
such as \texttt{-*- mode: org; fill-\codebreak{}column: 72 -*-}, or a block near the end that
begins \texttt{Local Variables:}, can choose the mode, the tab width, whether tabs or
spaces are used, the fill column and the spaces between sentences. These are
convenient in your own files.
Tuatara takes only those few settings. It never runs anything a file says: a
line that asks for an evaluation is left alone, and Tuatara says so. A package
can ask to take more settings, and it is given the plain text of each, never
anything to run.

A file opened safely is not read for these settings at all. For a file you do
not trust, the less it can say to the editor the better. The setting
\texttt{*file-\codebreak{}local-\codebreak{}variables*} is true as it comes. Set it to false and no file is
read for them, safe or not.

\begin{try}
\textbf{Try it, with an attachment you were not expecting}
\begin{enumerate}
\item In a terminal, make a note with a character in it that draws nothing:
   \texttt{printf 'From the council.\codebreak{}\textbackslash{}nMeet at the gate\textbackslash{}342\textbackslash{}200\textbackslash{}213 at noon.\codebreak{}\textbackslash{}n' > \textasciitilde{}/attachment.\codebreak{}txt}.
\item Start Tuatara and press \texttt{C-x C-f}. Type \texttt{\textasciitilde{}/attachment.\codebreak{}txt} and press \texttt{RET}.
   Look between \texttt{gate} and \texttt{at}. There is nothing to see. Press \texttt{C-x k} to close
   the buffer.
\item Press \texttt{C-u C-x C-f}, type \texttt{\textasciitilde{}/attachment.\codebreak{}txt} and press \texttt{RET}. The line now
   reads \texttt{Meet at the gate<200B> at noon.} and the mode line says \texttt{Safe}.
\item Press \texttt{i} and type a letter. The echo area says the buffer is read-only.
   Press \texttt{C-g}.
\item Press \texttt{M-x end-\codebreak{}safe-\codebreak{}visit} and \texttt{RET}. The file opens as usual and can be
   edited. Close it again with \texttt{C-x k}.
\item Press \texttt{C-x C-r}. The list marks \texttt{attachment.\codebreak{}txt} as one you opened safely.
   Press \texttt{C-g} to leave it.
\end{enumerate}
\end{try}

\subsubsection{Commands for opening safely}
\label{sec:281}
\begin{xltabular}{\textwidth}{@{}l >{\raggedright\arraybackslash\hsize=0.531\hsize}X l >{\raggedright\arraybackslash\hsize=1.469\hsize}X@{}}
\toprule
\textbf{Modal} & \textbf{Non-modal} & \textbf{Command} & \textbf{What it does}\\
\midrule
\endhead
none & \texttt{M-x find-\codebreak{}file-\codebreak{}safely} & \texttt{find-\codebreak{}file-\codebreak{}safely} & Open a file safely, for looking at a file from a source you do not trust: read-only, with no major mode, no hooks and no version control, so that nothing in the file is acted on; its file system left exactly as it was (no copy beside it), and its access time too where the system allows.\\
none & \texttt{M-x end-\codebreak{}safe-\codebreak{}visit} & \texttt{end-\codebreak{}safe-\codebreak{}visit} & Open the file of this buffer, opened safely, as usual from now on, where it is.\\
none & \texttt{C-x C-q} & \texttt{read-\codebreak{}only-\codebreak{}mode} & Make the current buffer read-only, or editable again if it is; with a count, read-only when it is positive, editable when it is zero or less.\\
none & \texttt{M-x open-\codebreak{}dropped-\codebreak{}files} & \texttt{open-\codebreak{}dropped-\codebreak{}files} & Visit the files dropped on the window, in the window they were dropped on, as C-x C-f does, the last one shown; safely, as C-u C-x C-f does, when \textbf{OPEN-DROPPED-FILES-SAFELY} is true.\\
\bottomrule
\end{xltabular}

\subsubsection{Settings and faces}
\label{sec:282}
\begin{xltabular}{\textwidth}{@{}l l >{\raggedright\arraybackslash\hsize=1.000\hsize}X@{}}
\toprule
\textbf{Setting} & \textbf{Standard} & \textbf{What it changes}\\
\midrule
\endhead
\texttt{*safe-\codebreak{}directories*} & \texttt{nil} & Directories whose files always open safely, as M-x find-file-safely opens a file: read-only, with no major mode, hooks or version control, and nothing written beside them; for example ’(“\textasciitilde{}/Evidence/”).\\
\texttt{*open-\codebreak{}binary-\codebreak{}files-\codebreak{}safely*} & \texttt{t} & True to open safely, as M-x find-file-safely does, a file that looks binary: one with a null (NUL) octet in its first 64 KB, as executables, archives and disk images have and text has not.\\
\texttt{*open-\codebreak{}dropped-\codebreak{}files-\codebreak{}safely*} & \texttt{nil} & True to open files dropped on the window safely, as M-x universal-argument M-x find-file does: read-only, no mode, nothing written beside them.\\
\texttt{*invisible-\codebreak{}characters-\codebreak{}when-\codebreak{}safe*} & \texttt{t} & True to show the characters that draw nothing as their codes (<FEFF>, <200B>) in every file opened safely (M-x find-file-safely) that has not chosen otherwise with M-x invisible-characters-mode; NIL to leave those buffers to \textbf{INVISIBLE-CHARACTERS-MODES} and \textbf{GLOBAL-INVISIBLE-CHARACTERS-MODE} as any other.\\
\texttt{*text-\codebreak{}direction-\codebreak{}when-\codebreak{}safe*} & \texttt{nil} & True to show right-to-left text in the order it is read in a file opened safely (M-x find-file-safely) too; NIL to show it there in the order it is stored, unless the buffer chose otherwise with M-x text-direction-mode.\\
\texttt{*file-\codebreak{}local-\codebreak{}variables*} & \texttt{t} & True for a file’s own settings (-\textbf{- mode: org -}- on its first line, a Local Variables block near its end) to be taken as it is opened, namely its major mode, tab width, tabs or spaces, fill column and spaces between sentences, and those a package takes; nothing else.\\
\bottomrule
\end{xltabular}

\begin{xltabular}{\textwidth}{@{}l >{\raggedright\arraybackslash\hsize=1.000\hsize}X@{}}
\toprule
\textbf{Face} & \textbf{What it paints}\\
\midrule
\endhead
\texttt{:completion-\codebreak{}safe} & The mark beside a file among completion candidates that was opened safely, and opens so again.\\
\bottomrule
\end{xltabular}

\subsection{Evidence folders}
\label{sec:290}
Sometimes one safe look is not enough. If you work with files that may be
evidence, such as a case folder from a client, nothing of yours must ever be
written among them. A stray backup file, an auto-save copy or a changed access
time can be a problem later. An evidence folder is a folder that you tell
Tuatara about, and from then on it does two things. Every file in it opens
safely, however it is opened. And nothing is written in it or below it by
anything Tuatara does.

The setting is \texttt{*evidence-\codebreak{}folders*} and it is empty as it comes, so no folder is
evidence until you say so. Put it in your init file (chapter 20). A good plan is
one folder for each case, with the evidence in one and your own work in another
beside it:

\begin{figure}[H]
\caption{Figure 12.1. A case folder. The evidence and your work sit side by side.}\label{org284}
\begin{Verbatim}[breaklines,breakautoindent=false,breakindent=1em,breakafter=/-.\,;:=,commandchars=\\\{\}]
~/cases/
├── 2026-0001/
│   ├── evidence/        everything here is evidence
│   │   ├── report.pdf
│   │   └── letters.zip
│   └── work/            your notes and copies, saved as usual
│       └── timeline.org
└── 2026-0002/
    ├── evidence/
    └── work/
\end{Verbatim}
\end{figure}

Never put your work folder inside the evidence folder. Everything below an
evidence folder is evidence, so a note kept there could not be saved.

\subsubsection{Naming the folders}
\label{sec:285}
The setting is a list. Each entry is written in one of three ways, and the list
may mix them.

\begin{xltabular}{\textwidth}{@{}>{\raggedright\arraybackslash\hsize=0.611\hsize}X >{\raggedright\arraybackslash\hsize=1.389\hsize}X@{}}
\toprule
\textbf{Entry} & \textbf{What it names}\\
\midrule
\endhead
\texttt{"\textasciitilde{}/\codebreak{}cases/*/\codebreak{}evidence/"} & A path. \texttt{*} stands for the name of any one folder, \texttt{**} for any folders at all, none too.\\
\texttt{(:name "evidence" :within "\textasciitilde{}/\codebreak{}cases/")} & Every folder called \texttt{evidence} inside \texttt{\textasciitilde{}/cases/}. Without \texttt{:within}, anywhere.\\
\texttt{(:regex "\textasciicircum{}[0-9]\{4\}-\codebreak{}[0-9]\{4\}\$" :within "\textasciitilde{}/\codebreak{}cases/" :then "evidence")} & In every folder of \texttt{\textasciitilde{}/cases/} whose whole name matches, the folder named by \texttt{:then}. Without \texttt{:then}, the matching folder itself.\\
\bottomrule
\end{xltabular}

For the layout in Figure 12.1, the first entry is enough:

\begin{Verbatim}[breaklines,breakautoindent=false,breakindent=1em,breakafter=/-.\,;:=,commandchars=\\\{\}]
\color[HTML]{A85F00}(setf\normalcolor{} \color[HTML]{A85F00}*evidence-folders*\normalcolor{} \color[HTML]{222222}'\color[HTML]{A85F00}(\color[HTML]{0A7AB3}"\color[HTML]{222222}~/cases/\color[HTML]{0A7AB3}*/evidence/"\color[HTML]{A85F00}))
\end{Verbatim}

Capital letters do not count in the comparison. A file is evidence by where it
really is, with every link followed. A link in your work folder that leads into
evidence does not make a loophole: the file behind the link is evidence. An
entry that cannot be read is left out and reported in \texttt{*Messages*} and by
\texttt{M-x check-\codebreak{}health}. The other entries still count. Give \texttt{:name} and \texttt{:regex}
entries a \texttt{:within}. Without it, Tuatara cannot tell from the setting alone
which repositories hold evidence and has to search for it.

\subsubsection{What happens to a file of evidence}
\label{sec:286}
A file in an evidence folder is opened safely whatever way you open it, whether
by \texttt{C-x C-f}, from a folder list, from a link, or on the command line. The mode
line adds the word \texttt{Evidence} before \texttt{Safe}, lit in the colour of errors, so
that you see at once where you are. \texttt{M-x end-\codebreak{}safe-\codebreak{}visit} refuses.

Nothing is written there. This covers saving, \texttt{C-x C-w}, an export to PDF or
Word, and Tuatara’s own backup, auto-save and lock files. It also covers making,
renaming, moving or deleting anything in the folder or below it. A refusal says
what could not be done, and ends with the folder and the entry of
\texttt{*evidence-\codebreak{}folders*} that made it evidence. A question that asks where to write,
such as \texttt{C-x C-w} or where to export to, does not start in an evidence folder.
It starts in the last folder you chose that is not one, or in your home folder.

A terminal that you open from a file of evidence starts in your home folder.
Then the shell’s prompt and helpers do not read what the folder holds, such as a
\texttt{.git} folder or an \texttt{.envrc} file.

\subsubsection{Asking, and copying out}
\label{sec:287}
\texttt{M-x describe-\codebreak{}evidence} tells you whether the file or folder here is evidence.
It looks at the entry on the line of a folder list, at a path written at the
cursor, or at the file of the buffer. It names the entry of the setting that
made it so. A link is evidence when what it leads to is. For a file with
another name that may be in an evidence folder, it says maybe. With the setting
empty, it says that no folder is named.

To work on a file, copy it out. \texttt{M-x copy-\codebreak{}from-\codebreak{}evidence} asks for a folder,
which must exist and may not be evidence, and copies the files marked in a
folder list, or the one on the line, or the file at the cursor, or the file of
the buffer. It works in the background. Whether a file is evidence is decided
again by the file it opens to copy, so a link changed in the meantime to lead
out of evidence copies nothing. Each copy is checked against its original by its
SHA-256 sum (a number worked out from the contents, so that a single changed
byte shows). A copy that differs is taken away. Nothing is written over a file
that is already there. Next to the copies is a log,
\texttt{copied-\codebreak{}from-\codebreak{}evidence-\codebreak{}DATE-\codebreak{}TIME.\codebreak{}log}, readable by you alone. It says for each
file what it was, where it came from, its size, both sums and the time. A folder
is not copied: mark the files in it.

Saving a file of evidence under a new name in your work folder with \texttt{C-x C-w} is
also a copy. The buffer then visits the copy. It is still opened safely, because
the text is the evidence’s, but the word \texttt{Evidence} is gone and you can edit it.
Use \texttt{M-x end-\codebreak{}safe-\codebreak{}visit} when you decide it is safe to open as usual.

\subsubsection{What the guard does not cover}
\label{sec:288}
The guard is Tuatara’s own. Be clear about its edges.

\begin{itemize}
\item A file with several names (a hard link) may have another name in an evidence
  folder. While evidence folders are named, Tuatara does not write such a file
  in place. It saves a new file and says so, and the other names keep what they
  held. Their link count and change time do change. If saving a new file is not
  possible, the save is refused. Such a file opens safely, \texttt{M-x end-\codebreak{}safe-\codebreak{}visit}
  refuses, and \texttt{M-x describe-\codebreak{}evidence} says maybe, but the mode line does not
  show \texttt{Evidence}.
\item The access time of a file that is not yours cannot be kept (see Opening one
  file safely). Git itself sets the access time of folders it reads, so
  \texttt{C-x g} in a repository of evidence can change those. Mount evidence read-only
  if access times must stay. To tell whether a path is evidence, Tuatara
  follows its links one at a time and stops at the first evidence folder, so it
  never reads a link inside one, which would change that link’s access time. A
  folder list of evidence does read the links it shows.
\item A folder list shows a preview of the file on the line. For a picture or a PDF
  in an evidence folder, that decodes a copy in the sandbox (next section).
\item Tuatara looks at where a folder really is just before it writes. A link or a
  drive changed in that moment is not seen until the next look.
\item Programs you run yourself write wherever they like. \texttt{M-x compile} and a shell
  in a terminal are not guarded.
\item A file opened in a background daemon (chapter 18) stays open there for as long
  as the daemon holds the buffer, even after the window that opened it is gone,
  and the drive it lives on cannot be unmounted. Do case work in plain
  \texttt{tuatara}.
\end{itemize}

A second guard is the system’s own. Make the evidence read-only for everyone
with \texttt{chmod -R a-w evidence/}. Tuatara’s guard does not replace it.

\begin{try}
\textbf{Try it, with a folder of case files}
\begin{enumerate}
\item In a terminal, make a case:
   \texttt{mkdir -p \textasciitilde{}/cases/\codebreak{}2026-0001/\codebreak{}evidence \textasciitilde{}/cases/\codebreak{}2026-0001/\codebreak{}work/\codebreak{}more} and
   \texttt{printf 'The report, as it was given.\codebreak{}\textbackslash{}n' > \textasciitilde{}/cases/\codebreak{}2026-0001/\codebreak{}evidence/\codebreak{}report.\codebreak{}txt}.
\item Open your init file with \texttt{C-x C-f \textasciitilde{}/.\codebreak{}config/\codebreak{}tuatara/\codebreak{}init.\codebreak{}lisp RET} and add
   \texttt{(setf *evidence-\codebreak{}folders* '("\textasciitilde{}/\codebreak{}cases/*/\codebreak{}evidence/"))}. Save, and start Tuatara
   again.
\item Open \texttt{\textasciitilde{}/cases/\codebreak{}2026-0001/\codebreak{}evidence/\codebreak{}report.\codebreak{}txt} with \texttt{C-x C-f}, not the safe
   way. The mode line shows \texttt{Evidence} and \texttt{Safe}.
\item Press \texttt{M-x describe-\codebreak{}evidence} and \texttt{RET}. It says yes, and names the entry.
\item Press \texttt{C-x C-q}, \texttt{i}, type a letter, \texttt{C-g}, then \texttt{C-x C-s}. Tuatara asks
   whether to save a write-protected file. Press \texttt{y} and \texttt{RET}. The save is
   refused and the message ends with \texttt{and nothing is written in it}.
\item Press \texttt{C-x C-w} and type \texttt{\textasciitilde{}/cases/\codebreak{}2026-0001/\codebreak{}work/\codebreak{}report.\codebreak{}txt} and \texttt{RET}. The
   copy is written. The buffer is now the copy: \texttt{Evidence} is gone.
\item Open the original again, and press \texttt{M-x copy-\codebreak{}from-\codebreak{}evidence}. Give
   \texttt{\textasciitilde{}/cases/\codebreak{}2026-0001/\codebreak{}work/\codebreak{}more/} as the folder. Look in it with \texttt{C-x d} for the
   copy and the log.
\end{enumerate}
\end{try}

\subsubsection{Commands, settings and faces}
\label{sec:289}
\begin{xltabular}{\textwidth}{@{}l >{\raggedright\arraybackslash\hsize=0.571\hsize}X l >{\raggedright\arraybackslash\hsize=1.429\hsize}X@{}}
\toprule
\textbf{Modal} & \textbf{Non-modal} & \textbf{Command} & \textbf{What it does}\\
\midrule
\endhead
none & \texttt{M-x describe-\codebreak{}evidence} & \texttt{describe-\codebreak{}evidence} & Say whether the file or folder here is evidence, and by which entry of \textbf{EVIDENCE-FOLDERS}: a rimu list’s entry on the line, a path written at point, or the buffer’s file; a link is evidence when what it leads to is.\\
none & \texttt{M-x copy-\codebreak{}from-\codebreak{}evidence} & \texttt{copy-\codebreak{}from-\codebreak{}evidence} & Copy files out of an evidence folder (\textbf{EVIDENCE-FOLDERS}) to a folder asked for, which may not be in one: in a rimu list the marked files, or the one on the line; elsewhere the file at point or the buffer’s file, each of them in an evidence folder.\\
\bottomrule
\end{xltabular}

\begin{xltabular}{\textwidth}{@{}l l >{\raggedright\arraybackslash\hsize=1.000\hsize}X@{}}
\toprule
\textbf{Setting} & \textbf{Standard} & \textbf{What it changes}\\
\midrule
\endhead
\texttt{*evidence-\codebreak{}folders*} & \texttt{nil} & The folders whose files are evidence: each opens safely however it is opened (read-only, no mode, hooks or version control run on it), and nothing is written in such a folder or below it, by saving, saving as, exporting, copying, moving, renaming, deleting or making folders, nor any of the editor’s own files (backups, auto-saves, locks).\\
\bottomrule
\end{xltabular}

\begin{xltabular}{\textwidth}{@{}l >{\raggedright\arraybackslash\hsize=1.000\hsize}X@{}}
\toprule
\textbf{Face} & \textbf{What it paints}\\
\midrule
\endhead
\texttt{:mode-\codebreak{}line-\codebreak{}evidence} & Evidence, in the mode line of a file in an evidence folder (\textbf{EVIDENCE-FOLDERS}): lit up in the theme’s error colour, so that it is seen at once that the file opened safely and nothing is written where it is.\\
\bottomrule
\end{xltabular}

\subsection{The sandbox around decoders}
\label{sec:293}
Pictures, PDFs and office files are not text. Showing one means decoding it, and
decoders have bugs: a crafted file has been known to take over the program that
reads it. Tuatara does not decode these files itself. It starts a separate
program of the system for each job, inside a sandbox, and takes back only the
result. A sandbox is a closed room: what happens inside does not reach the rest
of the machine. Tuatara uses a program called bubblewrap, whose name on the
command line is \texttt{bwrap}.

\begin{figure}[H]
\caption{Figure 12.2. What a decoder sees when it runs in the sandbox.}\label{org291}
\begin{Verbatim}[breaklines,breakautoindent=false,breakindent=1em,breakafter=/-.\,;:=,commandchars=\\\{\}]
  your files, your network, your terminal, other programs
  ────────────────\discretionary{}{}{}──────────  wall  ────────────────\discretionary{}{}{}───────
                      ┌───────────────\discretionary{}{}{}────────────────\discretionary{}{}{}─────┐
  copy of the file ──{\tuataraFallbackTWO{}▶}│  decoder (a system program)        │──{\tuataraFallbackTWO{}▶} pixels or text
  (in memory)         │  sees: /usr, read only             │    back to Tuatara
                      │  has: its own empty /tmp           │
                      │  no home, no network, no other     │
                      │  process, no way into the terminal │
                      │  ends after 10 seconds or 1 GB     │
                      └───────────────\discretionary{}{}{}────────────────\discretionary{}{}{}─────┘
\end{Verbatim}
\end{figure}

The decoder reads the file from a copy held in memory. It is given no name for
the file, so it cannot reach the other files beside it and cannot write to the
original. Beyond \texttt{/usr} it sees only a few files of \texttt{/etc}: the linker’s cache,
the alternatives and the font settings. Its time, its memory and the size of
what it may write are limited. On top of that the kernel refuses it system calls
that no decoder needs, such as making new rooms, loading kernel modules or
opening network sockets other than local ones. When the decoder finishes, when
you press \texttt{C-g} in its buffer, when you kill the buffer or when you leave
Tuatara, it is ended with everything it started. While a picture is drawn or the
text of a PDF is taken out, the work is a job: \texttt{C-x j} lists it and the mode
line lights \texttt{P}.

Where this helps: a PDF or picture that carries an attack against the decoder
finds a room with nothing in it to steal and no way out to send it. Where it
does not: it does not stop you from being misled by what the file says, and it
does not protect a program you start yourself. Some commands leave the sandbox
on purpose, such as \texttt{M-x open-\codebreak{}externally}, which hands the file to the desktop’s
own program. That command asks first, as long as \texttt{*open-\codebreak{}externally-\codebreak{}ask*} is
true.

Several programs use the same room. Code formatters (chapter 13) run inside it
with the project folder visible read-only, and no network. \LaTeX{} that makes a
PDF sees the document’s folder, the \TeX{} installation and your own fonts, all
read-only, and writes only in its own working folder.
Language grammars run in it too, where a sandbox exists. Without a sandbox, the
setting \texttt{*syntax-\codebreak{}helper-\codebreak{}unconfined*} decides whether the grammars of projects
you trust run on their own without one.

Press \texttt{M-x check-\codebreak{}health} to see whether \texttt{bwrap} was found and which version. If
it was not, or if you set \texttt{*sandbox-\codebreak{}program*} to nothing, Tuatara does not
decode: pictures and PDFs are shown as bytes and text. Your files are not at
risk for that reason. Exports to PDF still work, and say that \LaTeX{} ran without
a wall.

The limits are settings. A decoder gets 10 seconds (\texttt{*sandbox-\codebreak{}seconds*}) and 1
GB of memory (\texttt{*sandbox-\codebreak{}memory*}). All the decoders working on one file may use
2 GB together (\texttt{*sandbox-\codebreak{}total-\codebreak{}memory*}), where your system can hold them to it.
A file larger than 1 GB is not decoded (\texttt{*sandbox-\codebreak{}input-\codebreak{}limit*}). It opens as
bytes, and the echo area says why. A PDF that grew past the limit after it was
opened says so in its view when it is read again. \texttt{*sandbox-\codebreak{}seccomp*} is the
kernel’s refusal of unneeded calls, on x86-64 and AArch64. Keep it true unless a
decoder of yours fails because of it.

\begin{try}
\textbf{Try it, with a PDF from a stranger}
\begin{enumerate}
\item Open any PDF with \texttt{C-x C-f}. It is decoded in the sandbox and shown as pages
   (chapter 16).
\item While it loads, press \texttt{C-g}. In a shell, run
   \texttt{ps -eo pid,\codebreak{}args | grep [b]wrap | grep -v syntax-\codebreak{}helper}. Nothing of
   Tuatara’s is listed: \texttt{C-g} ended the decoder. The buffer says it was
   stopped, and \texttt{M-x revert-\codebreak{}buffer} reads it again. The syntax helpers you
   filtered out stay running, each in a sandbox of its own.
\item Press \texttt{M-x check-\codebreak{}health} and look for the line called Sandbox. It names
   \texttt{bwrap} and its version, and says whether the decoders are held to a total
   amount of memory.
\item Press \texttt{M-: *sandbox-\codebreak{}seconds* RET}. The echo area shows the number of seconds.
\end{enumerate}
\end{try}

\subsubsection{Commands and settings}
\label{sec:292}
\begin{xltabular}{\textwidth}{@{}l l l >{\raggedright\arraybackslash\hsize=1.000\hsize}X@{}}
\toprule
\textbf{Modal} & \textbf{Non-modal} & \textbf{Command} & \textbf{What it does}\\
\midrule
\endhead
none & \texttt{M-x check-\codebreak{}health} & \texttt{check-\codebreak{}health} & Show what each language and major mode has here.\\
\bottomrule
\end{xltabular}

\begin{xltabular}{\textwidth}{@{}l l >{\raggedright\arraybackslash\hsize=1.000\hsize}X@{}}
\toprule
\textbf{Setting} & \textbf{Standard} & \textbf{What it changes}\\
\midrule
\endhead
\texttt{*sandbox-\codebreak{}program*} & \texttt{"bwrap"} & The program that shuts in the programs decoding files for the editor, and \LaTeX{} making PDFs: bubblewrap, named bwrap, looked for on the PATH (or a full path to it).\\
\texttt{*sandbox-\codebreak{}seconds*} & \texttt{10} & The seconds a program decoding a file may run before it is ended, and the processor time it may take, 10 as standard.\\
\texttt{*sandbox-\codebreak{}memory*} & \texttt{1073741824} & The octets of memory a program decoding a file may take (1 GB); NIL for no limit.\\
\texttt{*sandbox-\codebreak{}total-\codebreak{}memory*} & \texttt{2147483648} & The octets of memory all the programs decoding one file may take together (2 GB), where the system can hold them to it (systemd-run, a scope of the user’s own); NIL for no such limit, each program still held to \textbf{SANDBOX-MEMORY}.\\
\texttt{*sandbox-\codebreak{}input-\codebreak{}limit*} & \texttt{1073741824} & The largest file a program decoding files is given (1 GB): it gets a copy in memory, which a larger file, or one claiming a huge size with nothing in it, would fill.\\
\texttt{*sandbox-\codebreak{}seccomp*} & \texttt{t} & True to have the kernel refuse a program decoding a file the system calls no decoder needs (see the file’s header), on x86-64 and AArch64; NIL to leave them to it, the rest of the sandbox kept.\\
\texttt{*format-\codebreak{}in-\codebreak{}sandbox*} & \texttt{t} & True for formatters to run shut in the sandbox: the project’s folder and the formatter’s own seen, read only, nothing written, no network.\\
\texttt{*syntax-\codebreak{}helper-\codebreak{}unconfined*} & \texttt{nil} & True to have the grammars run, where there is no sandbox, for the files of projects you trust, in a program of their own that is not shut in.\\
\bottomrule
\end{xltabular}

\subsection{Repositories you only look at}
\label{sec:295}
A git repository can run programs. Its hooks run when you commit. Its filters
and helpers run when you check out or compare files. A language server for its
code runs the project’s build scripts. A repository you downloaded may have set
these up for you. So Tuatara does not run any of it for a project you have not
said you trust.

The first time you open \texttt{C-x g} (the git status, chapter 15), start a language
server with \texttt{C-c L}, or compile in a project, Tuatara asks. The question says
everything a yes would allow. There are three answers.

\begin{xltabular}{\textwidth}{@{}l l >{\raggedright\arraybackslash\hsize=1.000\hsize}X@{}}
\toprule
\textbf{Modal} & \textbf{Non-modal} & \textbf{What it does}\\
\midrule
\endhead
none & \texttt{y} & Trust the project. Its programs may run. You are then asked once more, \texttt{y} and \texttt{RET}, because trusting cannot be taken back after programs ran.\\
none & \texttt{l} & Look only. Git reads the project and nothing it names is run. The answer is remembered.\\
none & \texttt{n} & Not now. Nothing runs, and you are asked again next time.\\
\bottomrule
\end{xltabular}

In a look-only repository the git status works and shows branches, commits and
changes. Git reads the working files and a private copy of the repository’s
settings and index. The repository’s own settings, hooks and filters never reach
git, and nothing in the repository is written, not even its index. Git does set
the access time of the folders it reads. The mode line of every git buffer says
\texttt{Look only} in the colour of warnings, and the first line of the status says
why. Commands that would change the repository are refused.

\texttt{M-x trust-\codebreak{}project} asks again, whatever you answered before. Use it to move a
project from look-only to trusted once you have read what is in it, or the other
way. A project can also be trusted in your init file. The setting
\texttt{*lsp-\codebreak{}trusted-\codebreak{}projects*} lists folders whose projects run their own programs
without asking. A link from a trusted folder to another does not give the other
any trust: a project is known by where it really is.

A repository that is, is inside, or holds an evidence folder is never trusted,
and it is never asked about. It is look-only for as long as the evidence is
there, even if it is in \texttt{*lsp-\codebreak{}trusted-\codebreak{}projects*}.

What this does not cover: it is a decision about Tuatara’s own work. If you open
a terminal inside the repository and run its build script, you have trusted it.

\begin{try}
\textbf{Try it, with a repository you were sent}
\begin{enumerate}
\item In a terminal, make a small repository to stand for a stranger’s:
   \texttt{mkdir \textasciitilde{}/sent \&\& cd \textasciitilde{}/sent \&\& git init \&\& printf 'hello\textbackslash{}n' > note.\codebreak{}txt \&\&
   git add note.\codebreak{}txt \&\& git commit -m first}.
\item In Tuatara, open \texttt{\textasciitilde{}/sent/\codebreak{}note.\codebreak{}txt} and press \texttt{C-x g}. The question about
   trusting the project appears. Press \texttt{l}.
\item The status opens. The mode line says \texttt{Look only}, and the first line of the
   status says why and that git only reads.
\item Press \texttt{M-x trust-\codebreak{}project} and \texttt{RET}. The same question appears. Press \texttt{n} to
   take the answer back. The status is look-only until you decide.
\end{enumerate}
\end{try}

\subsubsection{Command, setting and faces}
\label{sec:294}
\begin{xltabular}{\textwidth}{@{}l l l >{\raggedright\arraybackslash\hsize=1.000\hsize}X@{}}
\toprule
\textbf{Modal} & \textbf{Non-modal} & \textbf{Command} & \textbf{What it does}\\
\midrule
\endhead
none & \texttt{M-x trust-\codebreak{}project} & \texttt{trust-\codebreak{}project} & Ask again whether you trust the project of the current buffer (the git repository its file or folder is in, else its project, or the repository a git buffer shows), whatever was answered before.\\
\bottomrule
\end{xltabular}

\begin{xltabular}{\textwidth}{@{}l l >{\raggedright\arraybackslash\hsize=1.000\hsize}X@{}}
\toprule
\textbf{Setting} & \textbf{Standard} & \textbf{What it changes}\\
\midrule
\endhead
\texttt{*lsp-\codebreak{}trusted-\codebreak{}projects*} & \texttt{nil} & Folders whose projects you trust to run their own programs: their language servers start by themselves as their files are opened, git runs the programs they name for themselves (hooks, filters), and M-x compile runs their build commands without asking.\\
\bottomrule
\end{xltabular}

\begin{xltabular}{\textwidth}{@{}l >{\raggedright\arraybackslash\hsize=1.000\hsize}X@{}}
\toprule
\textbf{Face} & \textbf{What it paints}\\
\midrule
\endhead
\texttt{:mode-\codebreak{}line-\codebreak{}look-\codebreak{}only} & Look only, in the mode line of every git buffer of a repository looked at only (not trusted, or holding evidence): lit up in the theme’s warning colour, so that it is never taken for a repository you work in.\\
\texttt{:git-\codebreak{}look-\codebreak{}only} & The line at the top of a git status of a repository looked at only, saying why and that git only reads it.\\
\bottomrule
\end{xltabular}

\subsection{Encrypted files}
\label{sec:302}
Everything above is about protecting you from a file. This part is about
protecting a file from other people. Tuatara uses GnuPG (the program \texttt{gpg}) to
encrypt and sign. Encrypting turns text into something that only a key or a
passphrase can turn back. Signing adds a mark that shows who wrote a file and
that it has not been changed since.

\subsubsection{A file that ends in .gpg}
\label{sec:296}
A file whose name ends in \texttt{.gpg} opens decrypted. The text on disk stays
encrypted. The decrypted text lives in memory only. It is never auto-saved, and
it is not placed on the system clipboard unless you ask for that. Each save
encrypts it again. If the file is encrypted to a passphrase and not to a key,
gpg-agent’s prompt asks at once. When the minibuffer is what asks (next
section), the buffer waits until you press \texttt{RET} first.
\texttt{M-x gpg-\codebreak{}decrypt-\codebreak{}with-\codebreak{}passphrase} is the same.

Whom a file is encrypted to is a choice for each save.
\texttt{M-x gpg-\codebreak{}choose-\codebreak{}encryption-\codebreak{}keys} chooses again, between keys and a passphrase,
and the next save uses the new choice. \texttt{M-x gpg-\codebreak{}lock} makes gpg-agent forget
every passphrase it keeps, so the next file you open asks again. It does not
close buffers you already have open: kill those too, if you want their text gone
from memory.

\subsubsection{The passphrase prompt}
\label{sec:297}
Who asks for the passphrase is the setting \texttt{*gpg-\codebreak{}passphrase-\codebreak{}from*}. As it comes
(\texttt{:auto}), gpg-agent’s own window asks on a desktop, its prompt in the terminal
asks if Tuatara runs in one, and the minibuffer asks last. Either of the first
two keeps the passphrase out of Tuatara altogether. The minibuffer is the
fallback. It shows stars as you type, and the passphrase then passes through the
editor. Each character is wiped from memory as it is typed or deleted, and at
once if you give up with \texttt{C-g}. Still, the editor’s memory is not locked, a
look-alike program could ask in its place, and the garbage collector may have
copied text before it was wiped. So when the minibuffer asks, a window above it
says what that risks. \texttt{*gpg-\codebreak{}loopback-\codebreak{}warning*} is true as it comes. Turn it off
only when you know the risks by heart.

Which gpg runs is not left to chance. Tuatara runs the program found in the
system’s own folders (\texttt{/usr/\codebreak{}local/\codebreak{}bin}, \texttt{/usr/\codebreak{}bin}, \texttt{/bin}), never one that a
folder on your \texttt{PATH} might hold. \texttt{*gpg-\codebreak{}program*} names a different one, by a
full path from the root. Before each run Tuatara also checks that the folder gpg
takes its keys and options from (\texttt{\textasciitilde{}/.\codebreak{}gnupg}, or the folder \texttt{GNUPGHOME} names) is
yours, is not a link, and cannot be written by others. If the check fails, gpg
is not run, and the message says which folder is wrong and why.

\subsubsection{Encrypting and decrypting}
\label{sec:298}
\texttt{M-x gpg-\codebreak{}encrypt-\codebreak{}file} encrypts the file shown, or the one on the line of a
folder list, into \texttt{FILE.\codebreak{}gpg}. It asks whom to encrypt to in the minibuffer: your
keys, or a passphrase, with commas between several. It then offers to wipe the
original. Say no if you are not sure, since a wiped file cannot be recovered and
a solid state disk may keep old copies anyway. \texttt{M-x wipe-\codebreak{}file} does that wipe by
itself. \texttt{M-x gpg-\codebreak{}decrypt-\codebreak{}file} writes a decrypted copy to a file you name, after
warning you that its text is then on disk. To read an encrypted file, open the
\texttt{.gpg} file instead.

You can also encrypt part of a buffer. Select some text and use
\texttt{M-x gpg-\codebreak{}encrypt-\codebreak{}region}. The text becomes a block that starts with
\texttt{-----\codebreak{}BEGIN PGP MESSAGE-----}. This is useful for one paragraph of a note that
you still want to share as plain text. \texttt{M-x gpg-\codebreak{}decrypt-\codebreak{}region} turns the block
under the cursor, or the nearest one, back into text in place. The buffer is
then kept from the disk as an encrypted file is: no auto-save, no clipboard. A
save asks before writing the decrypted text unencrypted, so a habit of pressing
\texttt{C-x C-s} does not leave it on the disk.

\texttt{M-x gpg-\codebreak{}copy-\codebreak{}to-\codebreak{}clipboard} is the one way out. It puts the text you copied last
on the system clipboard, even if it came from an encrypted file. Other programs,
and clipboard managers, can read it and keep it.

\subsubsection{Signing and checking}
\label{sec:299}
\texttt{M-x gpg-\codebreak{}sign-\codebreak{}file} signs a file with your key. It asks which kind: a detached
signature, \texttt{FILE.\codebreak{}sig}, which leaves the file as it is; clear-signed text,
\texttt{FILE.\codebreak{}asc}; or a signed file, \texttt{FILE.\codebreak{}gpg}. \texttt{M-x gpg-\codebreak{}sign-\codebreak{}region} clear-signs the
selected text in place. \texttt{M-x gpg-\codebreak{}verify-\codebreak{}file} and \texttt{M-x gpg-\codebreak{}verify-\codebreak{}region} check
a signature. They say who signed, when, and how far that key is trusted.

Read the answer for what it says. A good signature tells you the text has not
changed since that key signed it. It does not tell you the text is safe, and it
tells you who signed only as far as you trust the key.

\subsubsection{Keys}
\label{sec:300}
\texttt{M-x gpg-\codebreak{}list-\codebreak{}keys} opens a list of your keys and those you have of others. The
list fills in as gpg answers. \texttt{M-x gpg-\codebreak{}keys-\codebreak{}details} shows one key’s
fingerprint, names and expiry. \texttt{M-x gpg-\codebreak{}import-\codebreak{}keys} and
\texttt{M-x gpg-\codebreak{}import-\codebreak{}keys-\codebreak{}region} add keys from a file or from selected text.
\texttt{M-x gpg-\codebreak{}keys-\codebreak{}export} writes the public part of marked keys to a file, and
\texttt{M-x gpg-\codebreak{}keys-\codebreak{}insert} puts it at the cursor in the buffer the list was opened
from. The folder list (chapter 10) has keys for the same jobs on files.

\begin{try}
\textbf{Try it, with an encrypted note}
\begin{enumerate}
\item In a terminal, make a note:
   \texttt{printf 'The alarm code is 4411.\codebreak{}\textbackslash{}n' > \textasciitilde{}/code.\codebreak{}txt}.
\item Open it with \texttt{C-x C-f}. Press \texttt{M-x gpg-\codebreak{}encrypt-\codebreak{}file} and \texttt{RET}. Choose a
   passphrase, press \texttt{RET} for no more, and give a passphrase twice. When it
   asks about wiping the original, answer \texttt{n}.
\item Open \texttt{\textasciitilde{}/code.\codebreak{}txt.\codebreak{}gpg} with \texttt{C-x C-f}. Give the passphrase if asked; gpg-agent
   may still hold it. The text appears. No auto-save copy of it is made.
\item Press \texttt{i}, add a line, press \texttt{C-g} and \texttt{C-x C-s}. Give the passphrase twice.
   The file on disk is encrypted again. In a terminal, \texttt{cat \textasciitilde{}/code.\codebreak{}txt.\codebreak{}gpg}
   shows nothing readable.
\item Press \texttt{M-x gpg-\codebreak{}lock}, kill the buffer with \texttt{C-x k} and open the file again.
   You are asked for the passphrase once more.
\item Open \texttt{\textasciitilde{}/code.\codebreak{}txt} and select its line. Press \texttt{M-x gpg-\codebreak{}encrypt-\codebreak{}region}, choose
   a passphrase and give it twice. The line becomes a block. Put the cursor in
   it and press \texttt{M-x gpg-\codebreak{}decrypt-\codebreak{}region}. Press \texttt{C-x C-s}: Tuatara asks whether
   to save the text decrypted. Answer \texttt{n} and press \texttt{RET}.
\end{enumerate}
\end{try}

\subsubsection{Commands and settings}
\label{sec:301}
\begin{xltabular}{\textwidth}{@{}l >{\raggedright\arraybackslash\hsize=0.850\hsize}X l >{\raggedright\arraybackslash\hsize=1.150\hsize}X@{}}
\toprule
\textbf{Modal} & \textbf{Non-modal} & \textbf{Command} & \textbf{What it does}\\
\midrule
\endhead
none & \texttt{M-x gpg-\codebreak{}encrypt-\codebreak{}file} & \texttt{gpg-\codebreak{}encrypt-\codebreak{}file} & Encrypt the file shown (in a list of a directory, the one on the line; with C-u, or where there is none, asked for) into FILE.gpg, to the keys or with a passphrase chosen in the minibuffer, commas between several.\\
none & \texttt{M-x gpg-\codebreak{}decrypt-\codebreak{}file} & \texttt{gpg-\codebreak{}decrypt-\codebreak{}file} & Decrypt the file shown (in a list of a directory, the one on the line) into a file asked for, the name without .gpg offered; after asking, since its text is then on disk where it can be read (made for you alone).\\
none & \texttt{M-x gpg-\codebreak{}encrypt-\codebreak{}region} & \texttt{gpg-\codebreak{}encrypt-\codebreak{}region} & Encrypt the text selected in place, as an armoured block (-----BEGIN PGP MESSAGE-----), to the keys or with a passphrase chosen in the minibuffer.\\
none & \texttt{M-x gpg-\codebreak{}decrypt-\codebreak{}region} & \texttt{gpg-\codebreak{}decrypt-\codebreak{}region} & Decrypt in place the armoured block (-----BEGIN PGP MESSAGE-----) in the text selected, or without a selection the one point is in, else the one nearest to point.\\
none & \texttt{M-x gpg-\codebreak{}decrypt-\codebreak{}with-\codebreak{}passphrase} & \texttt{gpg-\codebreak{}decrypt-\codebreak{}with-\codebreak{}passphrase} & Ask for the passphrase of the encrypted file this buffer waits for, shown as stars, and decrypt it with it.\\
none & \texttt{M-x gpg-\codebreak{}choose-\codebreak{}encryption-\codebreak{}keys} & \texttt{gpg-\codebreak{}choose-\codebreak{}encryption-\codebreak{}keys} & Choose again whom the encrypted file shown is encrypted to, or a passphrase, in the minibuffer: from its next save it is encrypted so.\\
none & \texttt{M-x gpg-\codebreak{}copy-\codebreak{}to-\codebreak{}clipboard} & \texttt{gpg-\codebreak{}copy-\codebreak{}to-\codebreak{}clipboard} & Put the text copied last on the system clipboard, even text of an encrypted file, which is otherwise kept within Tuatara: other programs (clipboard managers among them) can then read it, and keep it.\\
none & \texttt{M-x gpg-\codebreak{}lock} & \texttt{gpg-\codebreak{}lock} & Make gpg-agent forget every passphrase it keeps for a while (ten minutes, as it comes, its default-cache-ttl), of encrypted files and of your keys: the next to open one asks again.\\
none & \texttt{M-x wipe-\codebreak{}file} & \texttt{wipe-\codebreak{}file} & Wipe the file shown (in a list of a directory, the one on the line; with C-u, or where there is none, asked for) after asking: its contents overwritten with random octets, put on disk, then the file renamed and deleted, in the background, so that it is not found again where the disk writes in place.\\
none & \texttt{M-x gpg-\codebreak{}sign-\codebreak{}file} & \texttt{gpg-\codebreak{}sign-\codebreak{}file} & Sign the file shown (in a list of a directory, the one on the line) with your key: a detached signature (FILE.sig, the file left as it is), clear-signed text (FILE.asc) or a signed file (FILE.gpg), asked which.\\
none & \texttt{M-x gpg-\codebreak{}sign-\codebreak{}region} & \texttt{gpg-\codebreak{}sign-\codebreak{}region} & Clear-sign the text selected in place (-----BEGIN PGP SIGNED MESSAGE-----), with your key.\\
none & \texttt{M-x gpg-\codebreak{}verify-\codebreak{}file} & \texttt{gpg-\codebreak{}verify-\codebreak{}file} & Check the signature of the file shown (in a list of a directory, the one on the line): a detached signature (FILE.sig, FILE.asc) against its file, else the signature within the file; say who signed it, when, and how far that key is trusted.\\
none & \texttt{M-x gpg-\codebreak{}verify-\codebreak{}region} & \texttt{gpg-\codebreak{}verify-\codebreak{}region} & Check the signature of the clear-signed block (-----BEGIN PGP SIGNED MESSAGE-----) in the text selected, or without a selection the one point is in, else the one nearest to point: who signed it, when, and how far that key is trusted.\\
none & \texttt{M-x gpg-\codebreak{}import-\codebreak{}keys} & \texttt{gpg-\codebreak{}import-\codebreak{}keys} & Import the keys in a file asked for (armoured or not) into your keyring; say what was imported.\\
none & \texttt{M-x gpg-\codebreak{}import-\codebreak{}keys-\codebreak{}region} & \texttt{gpg-\codebreak{}import-\codebreak{}keys-\codebreak{}region} & Import the keys in the text selected (an armoured -----BEGIN PGP PUBLIC KEY BLOCK-----) into your keyring.\\
none & \texttt{M-x gpg-\codebreak{}list-\codebreak{}keys} & \texttt{gpg-\codebreak{}list-\codebreak{}keys} & List your GnuPG keys and those you have of others: there m marks one, M-x gpg-keys-details shows its details, M-x gpg-import-keys imports keys from a file, M-x gpg-keys-export exports, M-x gpg-keys-insert inserts and D deletes those marked.\\
none & \texttt{M-x gpg-\codebreak{}keys-\codebreak{}details} & \texttt{gpg-\codebreak{}keys-\codebreak{}details} & Show the details of the key at the cursor: its fingerprint, its names, its subkeys and what they can do, when they expire.\\
none & \texttt{M-x gpg-\codebreak{}keys-\codebreak{}export} & \texttt{gpg-\codebreak{}keys-\codebreak{}export} & Export the public part of the keys marked (or the one at the cursor), armoured, to a file asked for.\\
none & \texttt{M-x gpg-\codebreak{}keys-\codebreak{}insert} & \texttt{gpg-\codebreak{}keys-\codebreak{}insert} & Put the public part of the keys marked (or the one at the cursor), armoured, at point in the buffer the list was opened from.\\
\bottomrule
\end{xltabular}

\begin{xltabular}{\textwidth}{@{}l l >{\raggedright\arraybackslash\hsize=1.000\hsize}X@{}}
\toprule
\textbf{Setting} & \textbf{Standard} & \textbf{What it changes}\\
\midrule
\endhead
\texttt{*gpg-\codebreak{}program*} & \texttt{"gpg"} & The GnuPG program that decrypts and encrypts .gpg files, by its full path, from the root (/opt/gnupg/bin/gpg), or by name, then looked for only in the system’s own folders of programs (/usr/local/bin, /usr/bin, /bin), never in a folder of yours your PATH may name, nor in the folder the editor is in.\\
\texttt{*gpg-\codebreak{}passphrase-\codebreak{}from*} & \texttt{:auto} & Who asks for a passphrase gpg needs.\\
\texttt{*gpg-\codebreak{}loopback-\codebreak{}warning*} & \texttt{t} & When true, the minibuffer asking for a passphrase for gpg’s loopback shows above it, while it asks, a window saying why it asks, how the passphrase goes and what that risks (the garbage collector’s copies, memory not locked, no keyboard grab, a look-alike).\\
\bottomrule
\end{xltabular}

\section{Writing code}
\label{sec:342}
Colours, indentation, comments, brackets, finding your way and running the
program.

You can write code in Tuatara with nothing set up. A Python or shell script
opens coloured, a new line starts at the right depth, and a few keys move you
over whole functions. If you install a language server for your language, you
get more: errors marked as you type, a jump to where a name is defined, a
rename that changes only the real uses. This chapter goes through all of it in
the order you meet it when you start a small program: what Tuatara knows about
the language, how lines are laid out, how to comment code in and out, how to
move over brackets and definitions, how to find where something is defined, how
to run the program and read its errors, and how to settle a merge conflict in
a file of code.

The last part is about Lisp. Tuatara is written in Common Lisp, and the Lisp
you type can run inside the editor. If you do not write Lisp you can skip it,
though it is the quickest way to try a small idea.

Two reminders from earlier chapters. \texttt{i} starts INSERT state, and \texttt{C-g} or
\texttt{ESC} returns to NORMAL state. Every key in this chapter is pressed in NORMAL
state unless it starts with \texttt{C-x}, \texttt{C-c}, \texttt{C-M-} or \texttt{M-}, which work in both.
The \texttt{C-c} keys for code work in the buffers of programming languages. In a
text file or an Org file they are free for other things.

\subsection{What Tuatara knows about code}
\label{sec:308}
When you open a source file, Tuatara picks the language from the file’s name,
or from its first line when the name says nothing (a line such as
\texttt{\#!/\codebreak{}bin/sh}). The language gives the buffer its major mode, such as
\texttt{python-\codebreak{}mode} or \texttt{bash-\codebreak{}mode}, and the mode decides how the buffer behaves.

Two helpers then look at the code. Figure 13.1 shows which does what.

\begin{figure}[H]
\caption{Figure 13.1. The two helpers that read code.}\label{org304}
\begin{Verbatim}[breaklines,breakautoindent=false,breakindent=1em,breakafter=/-.\,;:=,commandchars=\\\{\}]
                 your file: sum.py
                        |
          +-------------+--------------+
          |                            |
  built in: the grammar        a program you install:
  (tree-sitter)                the language server
          |                            |
  reads this file only         reads the whole project
          |                            |
   colours                      errors and warnings
   indentation                  go to a definition
   pieces to select             every use of a name
   definitions, outline         rename everywhere
                                completion, hover, fixes
\end{Verbatim}
\end{figure}

The first helper is built in. A \emph{grammar} is a description of how a language is
put together. Tuatara has one for about three hundred languages and reads your
text with it as you type, building a \emph{syntax tree}: the function holds a
statement, the statement holds a call, the call holds its arguments. The tree
is what tells a keyword from a name and a string from a comment, so it gives
the colours. It also gives the indentation, the pieces that can be selected,
and where each definition starts and ends. All of this works with no setup and
with no other program installed.

The second helper is a \emph{language server}, a separate program for each
language, such as rust-analyzer for Rust or pylsp for Python. It knows the
whole project and not only the file you are in, so it can answer questions the
grammar cannot. It is optional. The sections on language servers and on finding
definitions say how to start one.

A grammar is code someone else wrote, and a file you were sent can be made to
confuse it. So Tuatara never runs a grammar inside the editor. A small program
of its own does it, shut in a sandbox with its memory held and no network. If a
grammar fails on a file, only that buffer loses its colours, and the echo area
says so. The grammars of the common languages are on from the start. For one
you use less, Tuatara asks the first time whether to turn it on. The setting
\texttt{*grammars-\codebreak{}on*} lists the languages that are on without asking.

Files you do not trust, such as evidence, are opened a different way, and
chapter 12 describes it. A file opened that way is not coloured by a grammar
and is never given to a language server.

\subsubsection{How much colour}
\label{sec:305}
Colour comes in three styles. \emph{Rich} colours everything the language’s rules
name, and it is the standard. \emph{Sparse} colours only some kinds of thing, which
suits a reader who finds a full palette tiring. \emph{None} shows plain text. \texttt{M-x
rotate-\codebreak{}colouring} moves the buffer’s language to the next style, in the order
rich, sparse, none, and rich again. Every buffer in that language changes with
it.

To start a language in another style, put a line in your init file (chapter 20
describes it): \texttt{(set-\codebreak{}colouring 'python :sparse)}. Sparse has four levels. Level
1 colours comments and strings. Level 2 adds keywords and functions. Level 3,
the standard, adds types and constants. Level 4 colours everything. Give the
level with \texttt{:level}, as in \texttt{(set-\codebreak{}colouring 'python :sparse :level 2)}. The
language can be named as \texttt{python}, as a mode such as \texttt{python-\codebreak{}mode}, as
\texttt{prog-\codebreak{}mode} for every programming language, or as \texttt{t} for every buffer.

\subsubsection{Seeing what a language has}
\label{sec:306}
\texttt{M-x check-\codebreak{}health} shows a table with a line for each language: whether its
grammar, colours, indentation, selections and definitions work, which formatter
and language server it knows, and a tick after each one that is installed on
your machine. Use it when something does not work, to see whether the language
is missing a piece or you are missing a program.

\subsubsection{Commands, settings and faces}
\label{sec:307}
\begin{xltabular}{\textwidth}{@{}l l l >{\raggedright\arraybackslash\hsize=1.000\hsize}X@{}}
\toprule
\textbf{Modal} & \textbf{Non-modal} & \textbf{Command} & \textbf{What it does}\\
\midrule
\endhead
none & \texttt{M-x check-\codebreak{}health} & \texttt{check-\codebreak{}health} & Show what each language and major mode has here.\\
none & \texttt{M-x rotate-\codebreak{}colouring} & \texttt{rotate-\codebreak{}colouring} & Colour the current buffer’s language in the next style: rich, then sparse, then none, then rich again.\\
\bottomrule
\end{xltabular}

\begin{xltabular}{\textwidth}{@{}l l >{\raggedright\arraybackslash\hsize=1.000\hsize}X@{}}
\toprule
\textbf{Setting} & \textbf{Standard} & \textbf{What it changes}\\
\midrule
\endhead
\texttt{*grammars-\codebreak{}on*} & \texttt{(:common)} & The languages whose grammars read text without asking: a file in such a language is coloured and parsed as it opens, and so is the text in it inside another file (a code block in a note or a Markdown file).\\
\texttt{*syntax-\codebreak{}helper-\codebreak{}unconfined*} & \texttt{nil} & True to have the grammars run, where there is no sandbox, for the files of projects you trust, in a program of their own that is not shut in.\\
\bottomrule
\end{xltabular}

The colours themselves come from the theme. These are the faces the syntax
colours use most.

\begin{longtable}{@{}l l@{}}
\toprule
\textbf{Face} & \textbf{What it paints}\\
\midrule
\endhead
\texttt{:variable} & Variables; plain by default.\\
\texttt{:punctuation} & Brackets, delimiters and separators; plain by default.\\
\texttt{:comment} & Comments.\\
\texttt{:string} & String and character literals.\\
\texttt{:keyword} & Keywords of the language.\\
\texttt{:function} & Names of functions where they are defined or called.\\
\texttt{:type} & Names of types and classes.\\
\texttt{:constant} & Constants and literals such as true, nil and numbers.\\
\texttt{:number} & Numeric literals.\\
\texttt{:operator} & Operators; plain by default.\\
\texttt{:variable.\codebreak{}parameter} & Parameters of functions and closures.\\
\texttt{:punctuation.\codebreak{}bracket} & Brackets, braces and parentheses.\\
\texttt{:builtin} & Built-in functions, types and variables.\\
\texttt{:preprocessor} & Preprocessor directives.\\
\bottomrule
\end{longtable}

\subsection{Indentation and tabs}
\label{sec:311}
\emph{Indenting} a line means putting blanks before it, so that the shape of the
code shows its structure. Tuatara indents a line by the language’s own rules,
read from the syntax tree. A Python line after \texttt{def f():} goes in one level, a
closing \texttt{\}} in C lines up with the line that opened it, and the lines of a long
Lisp call line up under its first argument.

You do not have to type the blanks yourself. A new line you make with \texttt{RET},
or in NORMAL state with \texttt{o} and \texttt{O}, begins at the indentation the language
wants. If you leave such a line empty, the blanks are taken away again when you
press \texttt{C-g}.

\texttt{TAB} indents the line the cursor is on. The cursor can be anywhere in the
line. When the line is already right, \texttt{TAB} does something else: it completes
the name before the cursor, as the section on language servers describes. The
setting \texttt{*tab-\codebreak{}always-\codebreak{}indent*} says what \texttt{TAB} does. Its standard value,
\texttt{:complete}, indents and then completes. With \texttt{t} it only indents. With \texttt{nil}
it indents only while the cursor is in the blanks at the start of the line, and
anywhere else it moves on to the next tab stop.

To lay out more than one line, select them and press \texttt{C-M-\textbackslash{}} (\texttt{indent-\codebreak{}region}).
In NORMAL state \texttt{>} (non-modal: \texttt{C-z >}, \texttt{indent-\codebreak{}selection-\codebreak{}or-\codebreak{}line}) does the
same for the selected lines, or for the current line when nothing is selected.
\texttt{M-m} (\texttt{back-\codebreak{}to-\codebreak{}indentation}) moves the cursor to the first character of the
line that is not a blank.

Lines inside a comment or a string that began on an earlier line are left as
they are, because their blanks are part of the text.

\subsubsection{Tabs and spaces}
\label{sec:309}
A \emph{tab stop} is a column that a tab character moves to. Tuatara puts them 4
columns apart. \texttt{*tab-\codebreak{}width*} changes that. Indentation is made of spaces unless
you say otherwise, which keeps a file looking the same in every program. Some
projects, a Makefile or a C project with its own style, use tab characters. For
those, \texttt{M-x indent-\codebreak{}tabs-\codebreak{}mode} switches the buffer between tabs and spaces, and
\texttt{*indent-\codebreak{}tabs-\codebreak{}mode*} sets the starting choice. \texttt{M-x insert-\codebreak{}tab} puts in one tab
character whatever the settings say, and \texttt{C-z TAB} (\texttt{tab-\codebreak{}to-\codebreak{}tab-\codebreak{}stop}) moves on
to the next tab stop with spaces or a tab as the buffer wants.

The init file has forms for one language at a time:

\begin{Verbatim}[breaklines,breakautoindent=false,breakindent=1em,breakafter=/-.\,;:=,commandchars=\\\{\}]
(set-indentation 'c 2)            ; two columns a level in C
(set-tab-width 'c-mode 8)         ; tab stops 8 apart in C files
(set-indent-tabs-mode 'bash t)    ; tabs in shell scripts
\end{Verbatim}

\texttt{set-\codebreak{}indentation} takes 1 to 16 columns for a language whose rules indent by
levels. \texttt{set-\codebreak{}tab-\codebreak{}width} takes 1 to 64, and \texttt{prog-\codebreak{}mode} in place of a language
means every programming language. A file can also say its own settings in its
first line or in a block of \texttt{Local Variables:} at its end, and Tuatara reads
them while \texttt{*file-\codebreak{}local-\codebreak{}variables*} is \texttt{t}.

Whole files are laid out by a \emph{formatter}, a program that sets the code out as
the language’s community expects, such as rustfmt for Rust, ruff for Python and
shfmt for shell scripts. The section on language servers covers it.

\begin{try}
\textbf{Try it, with a shell script}
\begin{enumerate}
\item Press \texttt{C-x C-f}, type \texttt{\textasciitilde{}/code-\codebreak{}practice/\codebreak{}count.sh} and press \texttt{RET}. Tuatara
   asks to make the folder when you save, if it is not there. Press \texttt{i} and
   type these seven lines, pressing \texttt{RET} at the end of each. Do not type any
   blanks in front. If Tuatara has already put some in, leave them.

\begin{Verbatim}[breaklines,breakautoindent=false,breakindent=1em,breakafter=/-.\,;:=,commandchars=\\\{\}]
   sum_up() \char123{}
   for n in "$@"; do
   total=$((total + n))
   done
   echo "$total"
   \char125{}
   sum_up 3 4 5
\end{Verbatim}

   Press \texttt{C-g}.
\item Press \texttt{\%} to select the whole buffer, then \texttt{C-M-\textbackslash{}}. The lines inside the
   function move in, the line inside the loop moves in again, \texttt{done} lines up
   with \texttt{for}, and the last line is at the left edge.
\item Press \texttt{C-g} and \texttt{C-x C-s} to save. Press \texttt{M-x}, type \texttt{compile} and press
   \texttt{RET}. Press \texttt{C-a} and \texttt{C-k} to clear the command offered, type \texttt{sh
   count.sh} and press \texttt{RET}. If Tuatara asks whether you trust the folder, it
   is your own, so type \texttt{y}, and at the second question type \texttt{y} and press
   \texttt{RET} (the section on running a build explains why it asks).
\item A window opens under the script with the name \texttt{*compile*}. The output is
   \texttt{12}, and the last line says the command finished with exit code 0.
\item Press \texttt{g g} and \texttt{j} (non-modal: \texttt{M-<}, \texttt{C-n}) to reach the \texttt{for} line. Press
   \texttt{I}, type four blanks and press \texttt{C-g}. The line sits too far in. Press \texttt{>}.
   It moves back to where the language wants it. Press \texttt{u} (non-modal: \texttt{C-/}) to
   undo that.
\item Try the tab setting. Press \texttt{M-x}, type \texttt{indent-\codebreak{}tabs-\codebreak{}mode} and press \texttt{RET}.
   Press \texttt{g g}, \texttt{j j} to reach the \texttt{total} line, \texttt{I}, type four blanks and
   press \texttt{C-g}. Press \texttt{>}, then \texttt{M-m}. The mode line says \texttt{C1}: the
   indentation is one tab character. With spaces it says \texttt{C4}. Run the command
   again to switch back. Lines already in the right column are left as they
   are, which is why you spoil one first.
\end{enumerate}
\end{try}

\subsubsection{Commands and settings}
\label{sec:310}
\begin{xltabular}{\textwidth}{@{}l l l >{\raggedright\arraybackslash\hsize=1.000\hsize}X@{}}
\toprule
\textbf{Modal} & \textbf{Non-modal} & \textbf{Command} & \textbf{What it does}\\
\midrule
\endhead
none & \texttt{TAB} & \texttt{indent-\codebreak{}for-\codebreak{}tab-\codebreak{}command} & Indent the current line as its major mode indents it: point in the blanks the line begins with goes to where its text begins, point in its text stays with the text.\\
none & \texttt{C-M-\textbackslash{}} & \texttt{indent-\codebreak{}region} & Indent each line the selection touches as its major mode indents it; the lines inside a comment or a string begun on a line before keep their blanks, which are text.\\
\texttt{>} & \texttt{C-z >} & \texttt{indent-\codebreak{}selection-\codebreak{}or-\codebreak{}line} & Indent the lines the selection touches, or the current line, as the language’s mode indents them.\\
none & \texttt{M-m} & \texttt{back-\codebreak{}to-\codebreak{}indentation} & Move point to the first character of its line that is not a space or a tab.\\
none & \texttt{C-z TAB} & \texttt{tab-\codebreak{}to-\codebreak{}tab-\codebreak{}stop} & Go on to the next tab stop, putting spaces in up to it, or a tab character where the buffer wants tabs (\textbf{indent-tabs-mode}, set-indent-tabs-mode, M-x indent-tabs-mode); with a count, that many stops.\\
none & \texttt{M-x insert-\codebreak{}tab} & \texttt{insert-\codebreak{}tab} & Insert a tab character at point, whatever the tab settings say; C-z TAB (tab-to-tab-stop) goes on to the next tab stop as they say.\\
none & \texttt{M-x indent-\codebreak{}tabs-\codebreak{}mode} & \texttt{indent-\codebreak{}tabs-\codebreak{}mode} & Switch the current buffer between tab characters and spaces where TAB and indentation go on to a tab stop; with a count, tabs when it is positive, spaces when it is zero or less.\\
\bottomrule
\end{xltabular}

\begin{xltabular}{\textwidth}{@{}l l >{\raggedright\arraybackslash\hsize=1.000\hsize}X@{}}
\toprule
\textbf{Setting} & \textbf{Standard} & \textbf{What it changes}\\
\midrule
\endhead
\texttt{*tab-\codebreak{}always-\codebreak{}indent*} & \texttt{:complete} & What M-x indent-for-tab-command does where the major mode indents lines (Lisp, Rust, C, shell scripts): :COMPLETE (the standard) indents the current line, and when that changes nothing completes the name at point, as M-x completion-at-point does in Lisp.\\
\texttt{*indent-\codebreak{}new-\codebreak{}lines*} & \texttt{t} & True to indent a line as it is begun: after RET, and on a line opened with NORMAL state’s M-x tuatara-open-below or M-x tuatara-open-above, the new line goes to its column as the language’s mode indents it (Lisp, Rust and shell scripts among them).\\
\texttt{*tab-\codebreak{}width*} & \texttt{4} & The columns from one tab stop to the next: a tab character shows its text on at the next stop, and M-x indent-for-tab-command goes on to it.\\
\texttt{*indent-\codebreak{}tabs-\codebreak{}mode*} & \texttt{nil} & True to put tab characters in where M-x indent-for-tab-command and indentation go on to a tab stop, NIL to put spaces in; for every buffer whose major mode says nothing of its own (set-indent-tabs-mode).\\
\texttt{*file-\codebreak{}local-\codebreak{}variables*} & \texttt{t} & True for a file’s own settings (-\textbf{- mode: org -}- on its first line, a Local Variables block near its end) to be taken as it is opened, namely its major mode, tab width, tabs or spaces, fill column and spaces between sentences, and those a package takes; nothing else.\\
\bottomrule
\end{xltabular}

\subsection{Comments}
\label{sec:313}
You comment out a line to stop the program running it without losing it, and
you write a comment to say something to the next reader, which may be you in
six months. Both are a single key.

\texttt{C-z C} (\texttt{toggle-\codebreak{}comment-\codebreak{}lines}) comments the current line. In NORMAL state
the single key \texttt{C} does the same. If you have selected lines, it comments every
line the selection touches. If the lines are all commented already, it takes
the comment marks away instead, so the same key switches the comment on and
off. The marker is the language’s own: \texttt{\#} in Python and shell scripts, \texttt{//} in
Rust and C, \texttt{;;} in Lisp. It goes at the indentation of the least indented
line, so the block stays tidy, and the cursor stays where it was.

If you like the cursor to move down to the next line afterwards, so that you
can press the key again and again to comment a run of lines, use
\texttt{toggle-\codebreak{}comment-\codebreak{}lines-\codebreak{}ret}. It has no key. Give it one in your init file:

\begin{Verbatim}[breaklines,breakautoindent=false,breakindent=1em,breakafter=/-.\,;:=,commandchars=\\\{\}]
(bind-key "C-z C" 'toggle-comment-lines-ret)
(bind-key "C" 'toggle-comment-lines-ret :map :normal)
\end{Verbatim}

\texttt{M-;} (\texttt{comment-\codebreak{}dwim}, “do what I mean”) does the sensible thing for where you
are. With lines selected it comments them as \texttt{C-z C} (modal: \texttt{C}) does. On a
line with no comment it starts one at the end of the code, with the marker and a
space, and leaves you in INSERT state ready to type the words. On a line that
has a comment already it moves you into it. In Python it writes
\texttt{\# adds the numbers} after the code, in Rust \texttt{// adds the numbers}.

The comment is also something you can select: \texttt{m c} selects the comment around
the cursor, and \texttt{m C} the comment with the comment lines right before and after
it.

\subsubsection{Commands}
\label{sec:312}
\begin{xltabular}{\textwidth}{@{}l >{\raggedright\arraybackslash\hsize=0.605\hsize}X l >{\raggedright\arraybackslash\hsize=1.395\hsize}X@{}}
\toprule
\textbf{Modal} & \textbf{Non-modal} & \textbf{Command} & \textbf{What it does}\\
\midrule
\endhead
\texttt{C} & \texttt{C-z C} & \texttt{toggle-\codebreak{}comment-\codebreak{}lines} & Comment the lines the selection touches, or the current line, or uncomment them when they are all commented: with the language’s marker (;; in Lisp, \emph{/ in Rust, \# in shell scripts and Org) at their smallest indentation, or each line between the two a comment has where it has an end (<!-- --> in HTML, /* *} in CSS).\\
none & \texttt{M-x toggle-\codebreak{}comment-\codebreak{}lines-\codebreak{}ret} & \texttt{toggle-\codebreak{}comment-\codebreak{}lines-\codebreak{}ret} & Comment the lines the region touches, or the line point is on, then going to the next line, or uncomment them when they are all commented; it has no key of its own.\\
none & \texttt{M-;} & \texttt{comment-\codebreak{}dwim} & Comment what suits where point is (M-;): with a region, its lines, or uncomment them when they all are (as M-x toggle-comment-lines-ret); on a line with a comment, go into it; on a blank line, begin a comment; after code, begin one at the line’s end, at \textbf{COMMENT-COLUMN}.\\
\texttt{m C} & \texttt{C-z m C} & \texttt{select-\codebreak{}comment} & Select the comment around point with the comment lines right before and after it.\\
\texttt{m c} & \texttt{C-z m c} & \texttt{select-\codebreak{}comment-\codebreak{}inside} & Select the comment around point.\\
\bottomrule
\end{xltabular}

\subsection{Brackets and pairs}
\label{sec:315}
Code is full of brackets, and a missing one is a common reason a program does
not run. Tuatara helps in three ways.

\emph{Showing the partner.} When the cursor is on a bracket, or just after a closing
one, the bracket that pairs with it is lit in the \texttt{:paren-\codebreak{}match} colour. A
bracket with no partner, or the wrong kind, is lit in \texttt{:paren-\codebreak{}mismatch}. It is
on as standard, and \texttt{M-x show-\codebreak{}paren-\codebreak{}mode} turns it off and on in every buffer.
When you type a closing bracket whose partner is out of sight, the echo area
tells you what it matches (\texttt{*blink-\codebreak{}matching-\codebreak{}paren*} controls that).

\emph{Closing them as you type.} \texttt{M-x electric-\codebreak{}pair-\codebreak{}mode} makes an opening bracket
write its closing half too, with the cursor between. Typing the closing half
when it is already there steps over it, and \texttt{DEL} between an empty pair takes
both away. It is off as standard. To turn it on for every session, put
\texttt{(setf *electric-\codebreak{}pair-\codebreak{}mode* t)} in your init file. It does not pair everything,
because some characters are usually plain text: the apostrophe in “don’t”, for
example. \texttt{*electric-\codebreak{}pair-\codebreak{}typing*} says which are paired. Its standard value,
\texttt{:safe}, pairs the delimiters of \texttt{*pair-\codebreak{}delimiters*} except those in
\texttt{*electric-\codebreak{}pair-\codebreak{}safe-\codebreak{}exclusions*}. \texttt{:all} pairs every one.

\emph{Putting a pair round something.} Select some text and press \texttt{M-(}
(\texttt{insert-\codebreak{}pair}) to put brackets round it. With nothing selected it writes an
empty pair with the cursor inside. On text that has the pair already, it takes
it off. \texttt{M-x delete-\codebreak{}pair} takes away the pair round the expression after the
cursor and keeps what was inside.

Selecting is just as direct. In NORMAL state \texttt{m a} selects a pair with what is
in it, and \texttt{m i} selects only what is inside. For bracketed expressions of code,
\texttt{m A} selects the whole expression and \texttt{m I} only its inside. These are the
same selection keys as in chapter 5, with the pair as the thing selected.

\begin{try}
\textbf{Try it, with a Python script}
\begin{enumerate}
\item Press \texttt{C-x C-f}, type \texttt{\textasciitilde{}/code-\codebreak{}practice/\codebreak{}sum.py} and press \texttt{RET}.
\item Press \texttt{i} and type this one line, then \texttt{C-g}:

\begin{Verbatim}[breaklines,breakautoindent=false,breakindent=1em,breakafter=/-.\,;:=,commandchars=\\\{\}]
   print(add_up([3, 4, 5]))
\end{Verbatim}
\end{enumerate}

Press \texttt{g g} (non-modal: \texttt{M-<}) then \texttt{O} to open a line above it. Type
\texttt{def add\_\codebreak{}up(numbers):}, press \texttt{RET} and type \texttt{return sum(numbers)}. Press \texttt{C-g}.
The second line is indented one level.
\begin{enumerate}
\item Put the cursor on the first \texttt{(} after \texttt{print}, with \texttt{g g}, \texttt{j j} and \texttt{l}
   if you need to move. The matching \texttt{)} at the end of the line is lit.
\item Press \texttt{m i}. The text inside the brackets is selected, and \texttt{M-o} makes the
   selection one piece larger. \texttt{C-g} lets go.
\item Press \texttt{M-x}, type \texttt{electric-\codebreak{}pair-\codebreak{}mode} and press \texttt{RET}. On a new line type
   \texttt{[}. You get \texttt{[]} with the cursor between. Type \texttt{1} and then \texttt{]}. The
   closing bracket is stepped over and you have \texttt{[1]}. Press \texttt{u} (non-modal:
   \texttt{C-/}) until the line is gone, and run the command again to switch pairing
   off.
\item Press \texttt{C-x C-s} to save.
\end{enumerate}
\end{try}

\subsubsection{Commands and settings}
\label{sec:314}
\begin{xltabular}{\textwidth}{@{}l >{\raggedright\arraybackslash\hsize=0.571\hsize}X l >{\raggedright\arraybackslash\hsize=1.429\hsize}X@{}}
\toprule
\textbf{Modal} & \textbf{Non-modal} & \textbf{Command} & \textbf{What it does}\\
\midrule
\endhead
none & \texttt{M-x show-\codebreak{}paren-\codebreak{}mode} & \texttt{show-\codebreak{}paren-\codebreak{}mode} & Switch on or off showing the bracket that pairs with the one at point, in every buffer: when on, the bracket before point when it closes, else the one after it when it opens, is shown with its partner in :paren-match, or in :paren-mismatch when the partner is of another kind or there is none.\\
none & \texttt{M-x electric-\codebreak{}pair-\codebreak{}mode} & \texttt{electric-\codebreak{}pair-\codebreak{}mode} & Turn closing delimiters as they are typed on or off.\\
none & \texttt{M-(} & \texttt{insert-\codebreak{}pair} & Put the pair of the delimiter of the key pressed around the region (M-( for parentheses), point on the first character inside; when the pair stands just outside the region already, take it off.\\
none & \texttt{M-x delete-\codebreak{}pair} & \texttt{delete-\codebreak{}pair} & Delete the brackets or quote marks around the balanced expression after point: the character there, after any spaces and tabs, and the one that closes the expression it begins.\\
\texttt{m a} & \texttt{C-z m a} & \texttt{select-\codebreak{}pair} & Select a pair of delimiters and what it holds, as select-inside-pair finds it.\\
\texttt{m i} & \texttt{C-z m i} & \texttt{select-\codebreak{}inside-\codebreak{}pair} & Select what is inside a pair of delimiters, the delimiters left out: the pair of the delimiter at point, or the one whose half is typed next ((, ), [, “, ‘, =, * and the rest of \textbf{pair-delimiters}).\\
\texttt{m A} & \texttt{C-z m A} & \texttt{select-\codebreak{}sexp} & Select the bracketed expression around point, its brackets included: point at the start, the mark at the end.\\
\texttt{m I} & \texttt{C-z m I} & \texttt{select-\codebreak{}sexp-\codebreak{}inside} & Select what is inside the bracketed expression around point, its brackets left out: point at the start, the mark at the end.\\
\bottomrule
\end{xltabular}

\begin{xltabular}{\textwidth}{@{}l >{\raggedright\arraybackslash\hsize=0.651\hsize}X >{\raggedright\arraybackslash\hsize=1.349\hsize}X@{}}
\toprule
\textbf{Setting} & \textbf{Standard} & \textbf{What it changes}\\
\midrule
\endhead
\texttt{*show-\codebreak{}paren-\codebreak{}mode*} & \texttt{t} & True to show the bracket that pairs with the one at point: the one before point when it closes, else the one after it when it opens.\\
\texttt{*blink-\codebreak{}matching-\codebreak{}paren*} & \texttt{t} & True to say, when a closing bracket is typed, what its partner is when it is out of sight (“Matches (defun foo ()”), or that it has none.\\
\texttt{*electric-\codebreak{}pair-\codebreak{}mode*} & \texttt{nil} & True to close delimiters as they are typed: an opening one writes its closing half after point, a closing half typed where it stands steps over it, and M-x delete-backward-char between an empty pair takes both.\\
\texttt{*electric-\codebreak{}pair-\codebreak{}typing*} & \texttt{:safe} & Which delimiters close themselves as they are typed: :safe, the pairs of \textbf{PAIR-DELIMITERS} less those \textbf{ELECTRIC-PAIR-SAFE-EXCLUSIONS} names; :all, every one of them; or a list of the opening characters, taken as it is (an empty list pairs nothing).\\
\texttt{*electric-\codebreak{}pair-\codebreak{}safe-\codebreak{}exclusions*} & \texttt{(\#\textbackslash{}< \#\textbackslash{}' \#\textbackslash{}` \#\textbackslash{}} \#\textbackslash{}* \#\textbackslash{}\textasciitilde{} \#\textbackslash{}\(\vert{}\) \#\textbackslash{}\textbackslash{} \#\textbackslash{}/ \#\textbackslash{}:\ldots{}= & The delimiters :safe leaves out: ordinary text far more often than a delimiter, such as the apostrophe in don’t, a colon, less-than.\\
\texttt{*pair-\codebreak{}delimiters*} & \texttt{((\#\textbackslash{}\{ . \#\textbackslash{}\}) (\#\textbackslash{}[ . \#\textbackslash{}]) (\#\textbackslash{}( . \#\textbackslash{})) (\#\textbackslash{}...} & The pairs of delimiters, each (OPEN .\\
\bottomrule
\end{xltabular}

\begin{xltabular}{\textwidth}{@{}l >{\raggedright\arraybackslash\hsize=1.000\hsize}X@{}}
\toprule
\textbf{Face} & \textbf{What it paints}\\
\midrule
\endhead
\texttt{:paren-\codebreak{}match} & The bracket at point and the one it pairs with (show-paren-mode).\\
\texttt{:paren-\codebreak{}mismatch} & A bracket at point that pairs with one of another kind, or with none (show-paren-mode).\\
\bottomrule
\end{xltabular}

\subsection{Moving by definitions and lists}
\label{sec:320}
Moving a character or a word at a time is slow in code. Code has a shape, and
Tuatara can move by the shape. A \emph{balanced expression} is a word, a string or a
bracketed group taken whole. A \emph{definition} is a function or a type: in Lisp a
bracket at the start of a line, and in Rust, C, Python and shell scripts
whatever the language’s grammar calls one. The keys use Control and Alt
together, written \texttt{C-M-}, and they come from Emacs. Figure 13.2 shows what some
of them do on a Lisp function, because Lisp, being nothing but brackets, shows
it most plainly.

\begin{figure}[H]
\caption{Figure 13.2. Where C-M-f, C-M-d and C-M-u take the cursor.}\label{org316}
\begin{Verbatim}[breaklines,breakautoindent=false,breakindent=1em,breakafter=/-.\,;:=,commandchars=\\\{\}]
  (defun average (numbers)
  ^       ^       ^
  |       |       C-M-f from here goes over (numbers)
  |       C-M-f from here goes over the name, then the list
  C-M-d here goes in, just after the bracket

  (defun average (numbers)
    (/ (sum numbers) (length numbers)))
        ^
        C-M-u here goes out, back to the bracket before (/
        C-M-f here goes over (sum numbers)
\end{Verbatim}
\end{figure}

\texttt{C-M-f} and \texttt{C-M-b} (\texttt{forward-\codebreak{}sexp}, \texttt{backward-\codebreak{}sexp}) move over the next or the
last balanced expression. \texttt{C-M-n} and \texttt{C-M-p} (\texttt{forward-\codebreak{}list}, \texttt{backward-\codebreak{}list})
do the same for bracketed groups only. \texttt{C-M-d} goes into the next group,
just after its opening bracket, and \texttt{C-M-u} goes back out of the group the
cursor is in, to its opening bracket. \texttt{M-x up-\codebreak{}list} goes out past the closing
bracket instead. In NORMAL state \texttt{W} and \texttt{B} do the work of \texttt{C-M-f} and \texttt{C-M-b}.

For whole definitions, \texttt{C-M-a} goes to the start of the one the cursor is in
(\texttt{beginning-\codebreak{}of-\codebreak{}defun}) and \texttt{C-M-e} to its end (\texttt{end-\codebreak{}of-\codebreak{}defun}). A count first
goes that many. \texttt{M-x mark-\codebreak{}defun} selects the whole definition with the comments
above it. Press it again to add the next one.

Some commands change the code by its pieces. \texttt{C-M-\codebreak{}SPC} or \texttt{C-M-@} (\texttt{mark-\codebreak{}sexp})
selects the next expression. \texttt{M-x kill-\codebreak{}sexp} cuts the expression after the
cursor and \texttt{C-M-\codebreak{}DEL} cuts the one before it, and \texttt{C-M-t} (\texttt{transpose-\codebreak{}sexps})
swaps the two either side of the cursor, which is a quick way to put two
arguments in the other order.

\subsubsection{Selecting by the code}
\label{sec:317}
In every state the syntax tree grows and shrinks a selection for you. \texttt{M-o}
(\texttt{expand-\codebreak{}selection}) selects the smallest piece around the cursor, and each
press after that selects the next larger one: the word, the call, the
statement, the function, the file. \texttt{M-i} (\texttt{shrink-\codebreak{}selection}) goes back a step.
\texttt{M-n} and \texttt{M-p} move the selection to the piece beside it, the next argument or
the next statement. A count before \texttt{M-o} grows several steps at once.

The \texttt{m} keys (non-modal: \texttt{C-z m} and the same letter) name a piece directly.
\texttt{m f} selects the body of the function around the cursor and \texttt{m F} the whole
function. \texttt{m t} and \texttt{m T} do the same for
a type, \texttt{m r} and \texttt{m R} for a parameter, and \texttt{m y} and \texttt{m Y} for a test. Which
pieces a language has depends on its grammar. Where the language has none of
the piece, the echo area says so.

\subsubsection{An outline of the file}
\label{sec:318}
\texttt{M-g o} (\texttt{buffer-\codebreak{}outline}) lists the definitions of the file, each with its
line number, in the order of the text. Type a few letters to narrow the list and
press \texttt{RET} to go there. In Lisp the list holds the top-level forms. Where a
language server runs, the list is its list of names, with functions, types and
fields nested.

\begin{try}
\textbf{Try it, with the shell script again}
\begin{enumerate}
\item Open \texttt{\textasciitilde{}/code-\codebreak{}practice/\codebreak{}count.sh} from the earlier box, and press \texttt{g g}.
\item Press \texttt{C-M-e}. The cursor goes to the end of the function, past the \texttt{\}}.
   Press \texttt{C-M-a}. It returns to the start of \texttt{sum\_up}.
\item Press \texttt{j} (non-modal: \texttt{C-n}) to go into the body, then \texttt{m f}. The body of the
   function is selected. Press \texttt{C-g}.
\item Put the cursor on the word \texttt{total} in the loop, with \texttt{j} and \texttt{l}. Press
   \texttt{M-o} a few times and watch the selection grow: first a small piece, then
   the loop, then the whole function. Press \texttt{M-i} to shrink it step by step.
\item Press \texttt{M-g o}. The list shows the definition. Press \texttt{C-g} to close it.
\end{enumerate}
\end{try}

\subsubsection{Commands}
\label{sec:319}
\begin{xltabular}{\textwidth}{@{}l l l >{\raggedright\arraybackslash\hsize=1.000\hsize}X@{}}
\toprule
\textbf{Modal} & \textbf{Non-modal} & \textbf{Command} & \textbf{What it does}\\
\midrule
\endhead
\texttt{W} & \texttt{C-M-f} & \texttt{forward-\codebreak{}sexp} & Move forward over a balanced expression.\\
\texttt{B} & \texttt{C-M-b} & \texttt{backward-\codebreak{}sexp} & Move backward over a balanced expression, the opposite of forward-sexp: a quote before the expression goes with it.\\
none & \texttt{C-M-n} & \texttt{forward-\codebreak{}list} & Move forward over the next bracketed group: the words, strings and comments before it passed over, the group taken whole.\\
none & \texttt{C-M-p} & \texttt{backward-\codebreak{}list} & Move backward over the bracketed group before point, to its opening bracket; the opposite of forward-list.\\
none & \texttt{M-x up-\codebreak{}list} & \texttt{up-\codebreak{}list} & Move forward out of the bracketed group around point, past its closing bracket; out of a string first to its end.\\
none & \texttt{C-M-u} & \texttt{backward-\codebreak{}up-\codebreak{}list} & Move back out of the bracketed group around point, to its opening bracket.\\
none & \texttt{C-M-d} & \texttt{down-\codebreak{}list} & Move forward into the next bracketed group, just past its opening bracket, the words, strings and comments before it passed over.\\
none & \texttt{C-M-a} & \texttt{beginning-\codebreak{}of-\codebreak{}defun} & Move back to the start of the definition point is in, or of the one before; with a count, that many definitions back, or forward to the start of the ones after when negative.\\
none & \texttt{C-M-e} & \texttt{end-\codebreak{}of-\codebreak{}defun} & Move forward to the end of the definition point is in, or of the one after, and on to the start of the next line when only blanks or a comment follow it there; with a count, that many definitions on, or back to the ends of those before when negative.\\
none & \texttt{M-x mark-\codebreak{}defun} & \texttt{mark-\codebreak{}defun} & Select the definition point is in, or the one after: point at its start, with the comments right above it and the blank lines above those, and the mark after it.\\
none & \texttt{C-M-@}, \texttt{C-M-\codebreak{}SPC} & \texttt{mark-\codebreak{}sexp} & Set the mark after the next balanced expression, point staying, and activate it; with a count, that many expressions, or before point when negative.\\
none & \texttt{M-x kill-\codebreak{}sexp} & \texttt{kill-\codebreak{}sexp} & Kill the balanced expression after point: to where forward-sexp goes.\\
none & \texttt{C-M-\codebreak{}DEL} & \texttt{backward-\codebreak{}kill-\codebreak{}sexp} & Kill the balanced expression before point: to where backward-sexp goes.\\
none & \texttt{C-M-t} & \texttt{transpose-\codebreak{}sexps} & Interchange the balanced expressions about point, point after them; with a count N, carry the expression before point past N expressions (back when N is negative); with 0, interchange the expressions at point and at the mark.\\
none & \texttt{M-o} & \texttt{expand-\codebreak{}selection} & Select the smallest piece of syntax around the selection, or around the character at point: the word, then the expression holding it, then the statement, and so on out; in a buffer with no syntax tree, as its language server sees the code.\\
none & \texttt{M-i} & \texttt{shrink-\codebreak{}selection} & Go back to the selection before the last M-o, or else select the first piece of syntax inside the selection; without a selection, the smallest piece of syntax at point.\\
none & \texttt{M-n} & \texttt{select-\codebreak{}next-\codebreak{}sibling} & Select the next piece of syntax beside the selection, or beside the one at point: the next argument, statement or item, and after the last one, the next beside the piece around them.\\
none & \texttt{M-p} & \texttt{select-\codebreak{}previous-\codebreak{}sibling} & Select the piece of syntax before the selection, or before the one at point: the argument, statement or item before it, and before the first one, the one before the piece around them.\\
\texttt{m F} & \texttt{C-z m F} & \texttt{select-\codebreak{}function} & Select all of the function around point.\\
\texttt{m f} & \texttt{C-z m f} & \texttt{select-\codebreak{}function-\codebreak{}inside} & Select the body of the function around point: the block of a Rust fn or closure, the body of a shell function, the forms of a Lisp defun, defmacro, defgeneric, defmethod or lambda after its lambda list.\\
\texttt{m T} & \texttt{C-z m T} & \texttt{select-\codebreak{}type} & Select all of the type around point.\\
\texttt{m t} & \texttt{C-z m t} & \texttt{select-\codebreak{}type-\codebreak{}inside} & Select the body of the type around point: the fields, variants or items of a Rust struct, enum, union, trait or impl; in Lisp, what follows the name in a defclass, defstruct, define-condition or deftype.\\
\texttt{m R} & \texttt{C-z m R} & \texttt{select-\codebreak{}parameter} & Select the parameter or argument around point with the comma after it.\\
\texttt{m r} & \texttt{C-z m r} & \texttt{select-\codebreak{}parameter-\codebreak{}inside} & Select the parameter or argument around point: of a Rust function’s parameters, a call’s arguments or a generic’s parameters; of a shell command’s arguments; of a Lisp lambda list or a call’s arguments.\\
\texttt{m Y} & \texttt{C-z m Y} & \texttt{select-\codebreak{}test} & Select all of the test around point.\\
\texttt{m y} & \texttt{C-z m y} & \texttt{select-\codebreak{}test-\codebreak{}inside} & Select the body of the test around point: a Rust function marked \#[test], a Python function whose name begins with test, a Go TestX or BenchmarkX, and the tests of Zig, Elixir, Dart, Julia and others, as the language’s rules know them.\\
none & \texttt{M-g o} & \texttt{buffer-\codebreak{}outline} & Go to a heading of the current buffer, chosen from a list of them: Org’s headings, Lisp’s comments of three semicolons and top-level forms, and in a language parsed with tree-sitter its definitions and headings (a function, a type, a Markdown heading), each shown as its line after its number, in the order of the text.\\
\bottomrule
\end{xltabular}

\subsection{Finding definitions}
\label{sec:322}
Reading code is mostly going from a name to where it is defined, and from a
definition to every place that uses it. Tuatara has a few ways, from the cheap
to the thorough.

\texttt{M-.} and \texttt{C-c d} (\texttt{xref-\codebreak{}find-\codebreak{}definitions}) go to where the name under the
cursor is defined. With a count, or where there is no name under the cursor, it
asks for one. When a name has several definitions they are listed in a buffer
called \texttt{*xref*}, and \texttt{RET} on one goes there. \texttt{M-,} or \texttt{C-c D}
(\texttt{xref-\codebreak{}go-\codebreak{}back}) takes you back to where you were, and pressing it again goes
further back. \texttt{C-M-,} (\texttt{xref-\codebreak{}go-\codebreak{}forward}) goes forward again.

In Lisp files this works with no setup, because the running Lisp knows its own
definitions. For other languages it needs a language server, as the next section
describes.

With a server, you can also ask more of the project:

\begin{itemize}
\item \texttt{C-c r} (\texttt{lsp-\codebreak{}find-\codebreak{}references}) lists every place the name is used, in
  \texttt{*xref*}.
\item \texttt{C-c g} (\texttt{lsp-\codebreak{}find-\codebreak{}symbol}) asks for a name and goes to a function or type of
  the project by it, with the name at the cursor offered.
\item \texttt{C-c T} goes to the definition of the type of the thing at the cursor, and
  \texttt{C-c I} to the implementations of a trait or an interface.
\item \texttt{C-c c} lists the places that call the function at the cursor and \texttt{C-c C} the
  functions it calls.
\end{itemize}

In an \texttt{*xref*} list, \texttt{n} and \texttt{p} show the next and the previous place in the
window you came from and leave the list where it is. \texttt{RET} or \texttt{l} goes to the
place. \texttt{o} shows it in the other window, and \texttt{h} or \texttt{q} closes the list.

For what no server knows, such as a word in a comment or a message in a string,
\texttt{M-s g} (\texttt{search-\codebreak{}in-\codebreak{}files}) searches the text of every file in the project as
you type. The lines it finds can be gone through with \texttt{M-g n} and \texttt{M-g p}, the
same keys as for a build’s errors, described below.

\subsubsection{Commands}
\label{sec:321}
\begin{xltabular}{\textwidth}{@{}l l l >{\raggedright\arraybackslash\hsize=1.000\hsize}X@{}}
\toprule
\textbf{Modal} & \textbf{Non-modal} & \textbf{Command} & \textbf{What it does}\\
\midrule
\endhead
none & \texttt{M-.}, \texttt{RET}, \texttt{C-c d} & \texttt{xref-\codebreak{}find-\codebreak{}definitions} & Go to the definition of the name at point.\\
none & \texttt{M-,}, \texttt{C-c D} & \texttt{xref-\codebreak{}go-\codebreak{}back} & Go back to where M-.\\
none & \texttt{C-M-,} & \texttt{xref-\codebreak{}go-\codebreak{}forward} & Go forward again to where M-, came back from.\\
none & \texttt{C-c r} & \texttt{lsp-\codebreak{}find-\codebreak{}references} & List where the thing at point is used, as its language server finds it, in \textbf{xref}: RET goes to one.\\
none & \texttt{C-c g} & \texttt{lsp-\codebreak{}find-\codebreak{}symbol} & Go to a function, type or other name of the current buffer’s project, as its language server finds them by what you type (the name at point offered): each is listed with what it is and where, and the one chosen is gone to.\\
none & \texttt{C-c T} & \texttt{lsp-\codebreak{}find-\codebreak{}type-\codebreak{}definition} & Go to the definition of the type of the thing at point (a variable’s type, a function’s result), as its language server finds it.\\
none & \texttt{C-c I} & \texttt{lsp-\codebreak{}find-\codebreak{}implementations} & Go to the implementations of the thing at point (of a trait, an interface, an abstract method), as its language server finds them: there at once when it is one, else listed in \textbf{xref}.\\
none & \texttt{C-c c} & \texttt{lsp-\codebreak{}incoming-\codebreak{}calls} & List the places that call the function at point, as its language server knows them, in \textbf{xref}: RET goes to one.\\
none & \texttt{C-c C} & \texttt{lsp-\codebreak{}outgoing-\codebreak{}calls} & List the functions the function at point calls, as its language server knows them, in \textbf{xref}: RET goes to one.\\
none & \texttt{M-s g} & \texttt{search-\codebreak{}in-\codebreak{}files} & Find text in the files of the project the current buffer’s file is in (see \textbf{PROJECT-ROOT-MARKERS}), or under its directory when it is in none, or under the directory it lists (C-x d), or \textbf{DEFAULT-DIRECTORY} for another buffer with no file; with a prefix argument, in a directory asked for.\\
\bottomrule
\end{xltabular}

\subsection{Language servers}
\label{sec:331}
A language server takes some time to start, and it runs in the background, so
you can go on typing while it starts. While it reads the project, a \texttt{P} shows
at the front of the mode line.

\subsubsection{Starting one, and trust}
\label{sec:323}
Tuatara knows the usual server for most languages and uses the ones it finds on
your \texttt{PATH}. Install one with your system’s packages or the language’s own
tool, such as rust-analyzer for Rust or python3-lsp-server for Python, and open
a file of that language again. \texttt{M-x check-\codebreak{}health} says which are installed.

A server runs the project’s own code: build scripts, macros, plugins. For that
reason nothing starts until you say you trust the project. \texttt{C-c L} (\texttt{lsp})
starts the servers for the buffer’s language and asks, the first time, whether
you trust the project. There are three answers.

\begin{itemize}
\item \texttt{y} trusts it. Its servers start by themselves whenever one of its files is
  opened, in later sessions too. Its build commands (\texttt{M-x compile}) run, and
  so do the programs git runs for it.
\item \texttt{l} looks only. Nothing of the project runs: no server, no build, and git
  only reads the repository. The answer is kept.
\item \texttt{n} leaves the question for now, and you are asked again next time.
\end{itemize}

\texttt{M-x trust-\codebreak{}project} asks again for the project of the current buffer, whatever
was said before. \texttt{M-x lsp-\codebreak{}forget-\codebreak{}project} takes the trust back and ends its
servers. To trust a whole folder of your own projects without being asked, put
it in \texttt{*lsp-\codebreak{}trusted-\codebreak{}projects*} in the init file, but never the folder where
files others send you are kept.

A \emph{project} is the folder that holds the language’s own mark, such as
\texttt{Cargo.\codebreak{}toml} for Rust, or else the nearest folder above holding \texttt{.git} or
\texttt{.project}. An empty \texttt{.project} file makes a folder a project.

If a language has no server installed, the commands that need one say so. For
example, \texttt{C-c d} says there is no way to find definitions in buffers of that
language, and nothing is broken. \texttt{M-x lsp-\codebreak{}restart} ends a stuck server and
starts it again, \texttt{M-x lsp-\codebreak{}shutdown} ends it, and \texttt{M-x lsp-\codebreak{}events} lists the
last messages sent to and from it, which is the place to look when a server
does not do what you expect.

\subsubsection{What the server tells you as you type}
\label{sec:324}
When the cursor rests in code, after \texttt{*lsp-\codebreak{}idle-\codebreak{}delay*} seconds, the echo area
shows the parameters of the call you are in, or else what the name under the
cursor is. A name and its other uses in the file are lit, and the mode line
counts the problems found after the mode’s name. A line with a problem has an
\texttt{E} (an error) or a \texttt{W} (a warning) beside it, and the faulty text is
underlined. A box under the cursor says what is wrong where it rests.

Some of these can be turned off. \texttt{*lsp-\codebreak{}signature-\codebreak{}help*} is the parameters,
\texttt{*lsp-\codebreak{}hover-\codebreak{}echo*} is what the name is, and \texttt{*lsp-\codebreak{}highlight-\codebreak{}symbol*} is the
lighting of a name’s other uses. \texttt{*lsp-\codebreak{}diagnostics-\codebreak{}display*} says how a problem
is said: \texttt{:box} under the cursor, \texttt{:echo} in the echo area, or \texttt{nil} not at all.
\texttt{*lsp-\codebreak{}diagnostics-\codebreak{}style*} is the underline: \texttt{:wave}, \texttt{:dotted}, \texttt{:line} or
\texttt{:background}.

The keys that go through the problems are \texttt{C-c n} and \texttt{C-c p}, which go to the
next and the previous one and say what it is, and \texttt{C-c l}, which lists all of
them in a window across the bottom. In that list \texttt{RET} opens the file at the
place.

\texttt{C-c h} (\texttt{lsp-\codebreak{}describe-\codebreak{}thing-\codebreak{}at-\codebreak{}point}) shows what the server knows of the thing
under the cursor, its type and its documentation, in the help window.

\subsubsection{Completion}
\label{sec:325}
\texttt{C-M-i} or \texttt{M-TAB} (\texttt{completion-\codebreak{}at-\codebreak{}point}) opens a list beside the cursor of
the names that fit what you have typed, each with what it is, and the chosen
one’s documentation beside the list. Typing narrows the list. \texttt{C-n} and \texttt{C-p}
choose, \texttt{TAB} or \texttt{RET} puts the chosen name in, and \texttt{C-g} closes the list. In
Lisp it offers the names of Lisp, and elsewhere what the server offers. Because
\texttt{*tab-\codebreak{}always-\codebreak{}indent*} is \texttt{:complete}, \texttt{TAB} also completes when the line is
already indented right.

Without a server you can still complete from the words you already have. \texttt{M-/}
(\texttt{expand-\codebreak{}word}) replaces the start of a word with the nearest word that begins
the same, from this buffer and then from others. Press it again for the next
nearest. \texttt{C-M-/} (\texttt{complete-\codebreak{}word}) shows those words as a list. Terminals send
that key as something else, and \texttt{C-M-i} lists the same words in a language with
no completion of its own.

\subsubsection{Rename, fixes and tidying}
\label{sec:326}
Renaming a function by searching for its name and replacing it is a classic way
to break a program. The same letters may be in a string or a comment, and you
miss the files you did not open. \texttt{C-c R} (\texttt{lsp-\codebreak{}rename}) asks the server instead.
It asks for the new name, then changes only the real uses, in every file of the
project. Each file’s changes are one step for \texttt{u} (non-modal: \texttt{C-/}) to take
back, and nothing is saved until you save (\texttt{C-x s} saves them all).

\texttt{C-c a} (\texttt{lsp-\codebreak{}code-\codebreak{}actions}) offers the server’s own fixes and changes for the
selection or the cursor, with those for what is wrong there first: fill in a
type, add an import, apply the suggestion for an error. \texttt{C-c o}
(\texttt{lsp-\codebreak{}organise-\codebreak{}imports}) puts the file’s imports in order and, where the server
does it, takes out the unused ones.

\texttt{C-c z} (\texttt{lsp-\codebreak{}fold}) hides the body of the block the cursor is in under its
first line, and shows it again when pressed on that line. \texttt{C-c Z}
(\texttt{lsp-\codebreak{}fold-\codebreak{}all}) hides every outermost block of the file, or shows them all
again when some are hidden.

Three more things can be shown within the code, each switched on and off by its
key. \texttt{C-c i} shows \emph{inlay hints}: the types of variables and the names of the
parameters a call fills in, written faint. \texttt{C-c S} colours each name by what the
server says it is, a parameter, a field, a macro, over the usual colours. \texttt{C-c
K} shows \emph{code lenses}, short notes such as \texttt{3 references} or \texttt{Run test} at the
end of a line, and \texttt{C-c k} runs the one on the current line. Their settings are
\texttt{*lsp-\codebreak{}inlay-\codebreak{}hints*}, \texttt{*lsp-\codebreak{}semantic-\codebreak{}tokens*} and \texttt{*lsp-\codebreak{}code-\codebreak{}lenses*}, which say
whether each is on as a file opens.

\subsubsection{Formatting}
\label{sec:327}
A \emph{formatter} is a program that sets code out as the language’s community has
agreed, so that every file of a project looks alike and a change shows only
what changed. \texttt{C-c f} (\texttt{format-\codebreak{}buffer}) runs the language’s formatter in the
background. It is shut in the sandbox, which lets it read the project and change
nothing else. It runs from the project’s folder, so a settings file of the
formatter’s own is found. If no formatter is installed the language server lays
the code out, and if neither is there Tuatara indents each line by its own
rules. The whole change is one step: \texttt{u} (non-modal: \texttt{C-/}) takes all of it
back. \texttt{C-c F} (\texttt{format-\codebreak{}region}) does the same for the selected lines.

To choose another formatter for a language, set \texttt{*language-\codebreak{}formatters*} in the
init file. \texttt{*format-\codebreak{}on-\codebreak{}save*} names the languages formatted before each save, or
\texttt{t} for every language with a formatter installed. A formatter that takes longer
than \texttt{*format-\codebreak{}seconds*} is stopped and nothing changes. \texttt{*format-\codebreak{}in-\codebreak{}sandbox*}
turns off the sandbox for a formatter that cannot work shut in.

\subsubsection{Settings of your own}
\label{sec:328}
\texttt{*lsp-\codebreak{}languages*} lists the languages whose servers are used.
\texttt{*language-\codebreak{}servers*} names a server of your own choice for a language in place
of the usual one. \texttt{*lsp-\codebreak{}wait*} is the seconds a command waits for an answer it
needs at once, and a key press stops the wait sooner. \texttt{*lsp-\codebreak{}watch-\codebreak{}files*} tells
servers about files other programs changed, and \texttt{*lsp-\codebreak{}watch-\codebreak{}skipped*} names the
folders it never watches, such as \texttt{target} and \texttt{node\_\codebreak{}modules}. \texttt{*lsp-\codebreak{}outline*}
makes \texttt{M-g o} list the server’s names.

\begin{try}
\textbf{Try it, with the Python script}
This needs a language server for Python, such as pylsp, ty or ruff. Without one,
the steps still show what Tuatara says.

\begin{enumerate}
\item Open \texttt{\textasciitilde{}/code-\codebreak{}practice/\codebreak{}sum.py}. Put the cursor on \texttt{add\_up} in the \texttt{print}
   line.
\item Press \texttt{M-x}, type \texttt{lsp} and press \texttt{RET}. If Tuatara asks whether you trust
   the folder, it is your own, so type \texttt{y} and press \texttt{RET}. The echo area says
   the server is starting and then that it is ready.
\item Press \texttt{C-c d}. The cursor jumps to the \texttt{def} line. Press \texttt{C-c D} to go back.
\item Press \texttt{C-c h}. The window shows what the server knows of the function. Press
   \texttt{C-x 1} (modal: \texttt{F}) when you have read it.
\item Press \texttt{C-c R}, type \texttt{total\_of} and press \texttt{RET}. The name changes in both
   places. Press \texttt{u} (non-modal: \texttt{C-/}) to undo both at once.
\item In the \texttt{return} line, press \texttt{A} and type \texttt{x}, then \texttt{C-g}, to make a
   mistake. After a moment an \texttt{E} shows beside the line and the mode line counts
   \texttt{1}. Press \texttt{C-c l} to list it, \texttt{RET} to go there, and \texttt{u} to take the mistake
   away.
\item Press \texttt{C-c f}. With a formatter installed, the file is laid out by it. The
   echo area names the formatter, or says why none could be used. Press \texttt{M-x},
   type \texttt{lsp-\codebreak{}shutdown} and press \texttt{RET} when you are done.
\end{enumerate}
\end{try}

\subsubsection{Commands}
\label{sec:329}
\begin{xltabular}{\textwidth}{@{}l >{\raggedright\arraybackslash\hsize=0.650\hsize}X l >{\raggedright\arraybackslash\hsize=1.350\hsize}X@{}}
\toprule
\textbf{Modal} & \textbf{Non-modal} & \textbf{Command} & \textbf{What it does}\\
\midrule
\endhead
none & \texttt{C-c L} & \texttt{lsp} & Start the language servers of the current buffer’s language for its project, after asking whether you trust the project: a server runs the project’s code (build scripts, macros, plugins).\\
none & \texttt{M-x trust-\codebreak{}project} & \texttt{trust-\codebreak{}project} & Ask again whether you trust the project of the current buffer (the git repository its file or folder is in, else its project, or the repository a git buffer shows), whatever was answered before.\\
none & \texttt{M-x lsp-\codebreak{}forget-\codebreak{}project} & \texttt{lsp-\codebreak{}forget-\codebreak{}project} & Trust the current buffer’s project no more: its language servers end, and no longer start by themselves.\\
none & \texttt{M-x lsp-\codebreak{}restart} & \texttt{lsp-\codebreak{}restart} & End the language servers the current buffer is open in and start them again; in a project not trusted, none is started (M-x lsp asks first).\\
none & \texttt{M-x lsp-\codebreak{}shutdown} & \texttt{lsp-\codebreak{}shutdown} & End the language servers the current buffer is open in; M-x lsp starts them again.\\
none & \texttt{M-x lsp-\codebreak{}events} & \texttt{lsp-\codebreak{}events} & List the last messages to and from the language servers the current buffer is open in (all those running when it is in none), the newest last, and what each said on its error output: to see what a server was asked and answered when it does not do what it should; the last 500 of each are kept.\\
none & \texttt{C-c l} & \texttt{lsp-\codebreak{}diagnostics} & List what the language servers found wrong in the open files, a line each, then in the files of their projects not open, in a window across the bottom; RET there goes to it, opening the file.\\
none & \texttt{C-c n} & \texttt{lsp-\codebreak{}next-\codebreak{}diagnostic} & Go to the next place the language servers found something wrong, and say what.\\
none & \texttt{C-c p} & \texttt{lsp-\codebreak{}previous-\codebreak{}diagnostic} & Go to the place before point where the language servers found something wrong, and say what.\\
none & \texttt{C-c h} & \texttt{lsp-\codebreak{}describe-\codebreak{}thing-\codebreak{}at-\codebreak{}point} & Show what the language server knows of the thing at point (its type, its documentation) in \textbf{Help}.\\
none & \texttt{C-M-i}, \texttt{M-TAB} & \texttt{completion-\codebreak{}at-\codebreak{}point} & Complete the word at point from a list beside it.\\
none & \texttt{M-/} & \texttt{expand-\codebreak{}word} & Put in place of the start of the word before point the nearest word that starts with it, looked for in this buffer, then in the others (those shown first).\\
none & \texttt{C-M-/} & \texttt{complete-\codebreak{}word} & List the words in the buffers that start with the text before point beside the cursor, the nearest first, each with the buffer it is in when that is not this one; M-/ puts them in one after another.\\
none & \texttt{C-c R} & \texttt{lsp-\codebreak{}rename} & Rename the thing at point everywhere the language server knows it is used, the new name asked for: each file’s changes one change undo takes back, nothing saved (C-x s saves them).\\
none & \texttt{C-c a} & \texttt{lsp-\codebreak{}code-\codebreak{}actions} & Offer the language server’s fixes and changes for the selection, or for point: those for what it found wrong there first; the one chosen is made.\\
none & \texttt{C-c o} & \texttt{lsp-\codebreak{}organise-\codebreak{}imports} & Have the language server put the file’s imports in order, as its own action for that does it (unused ones taken out, where it does that), once it answers, when the file is still as it was asked of.\\
none & \texttt{C-c z} & \texttt{lsp-\codebreak{}fold} & Hide the lines of the block the cursor is in (a function’s body, a struct, a run of imports), as its language server knows them, under the block’s first line; on a line whose block is hidden, show it again.\\
none & \texttt{C-c Z} & \texttt{lsp-\codebreak{}fold-\codebreak{}all} & Hide every outermost block of the file, as its language server knows them, under its first line; with any hidden, show all again.\\
none & \texttt{C-c i} & \texttt{lsp-\codebreak{}inlay-\codebreak{}hints-\codebreak{}mode} & Show the language servers’ inlay hints within the code, or no longer: the types of bindings, the names of the parameters arguments go to.\\
none & \texttt{C-c S} & \texttt{lsp-\codebreak{}semantic-\codebreak{}tokens-\codebreak{}mode} & Colour code as its language server sees it, or no longer: each name in the face of its kind (a parameter, a field, a macro), over the mode’s colours.\\
none & \texttt{C-c K} & \texttt{lsp-\codebreak{}code-\codebreak{}lenses-\codebreak{}mode} & Show what language servers say of lines (3 references, Run test) faint at each line’s end, or no longer; M-x lsp-code-lens runs one.\\
none & \texttt{C-c k} & \texttt{lsp-\codebreak{}code-\codebreak{}lens} & Run what a language server says of the current line (a code lens: Run test, 3 references), chosen when there are several, as the server’s own command.\\
none & \texttt{C-c f} & \texttt{format-\codebreak{}buffer} & Lay out the code of the current buffer as its language’s formatter does (\textbf{LANGUAGE-FORMATTERS}, else the language’s usual one: rustfmt, ruff, clang-format, shfmt\ldots{}), in the background, shut in the sandbox; the change is one that C-/ takes back whole.\\
none & \texttt{C-c F} & \texttt{format-\codebreak{}region} & Lay out the code selected as the language server does; with no server that lays out a part, each line selected indented as the mode indents it.\\
\bottomrule
\end{xltabular}

\subsubsection{Settings and faces}
\label{sec:330}
\begin{xltabular}{\textwidth}{@{}l >{\raggedright\arraybackslash\hsize=0.851\hsize}X >{\raggedright\arraybackslash\hsize=1.149\hsize}X@{}}
\toprule
\textbf{Setting} & \textbf{Standard} & \textbf{What it changes}\\
\midrule
\endhead
\texttt{*lsp-\codebreak{}trusted-\codebreak{}projects*} & \texttt{nil} & Folders whose projects you trust to run their own programs: their language servers start by themselves as their files are opened, git runs the programs they name for themselves (hooks, filters), and M-x compile runs their build commands without asking.\\
\texttt{*lsp-\codebreak{}languages*} & \texttt{t} & The languages whose language servers are used: T, every language with a server installed; a list of languages, those alone, ’(rust go python); or :NOT and languages, every one but those, ’(:not markdown latex).\\
\texttt{*language-\codebreak{}servers*} & \texttt{nil} & The language servers of languages you choose yourself, in place of the language’s usual ones: each (LANGUAGE SERVER\ldots{}), LANGUAGE a symbol (python, rust), each SERVER a known server’s name (“basedpyright”) or a list of a program and its arguments ((“my-server” ”--stdio”)); (LANGUAGE) alone for none.\\
\texttt{*lsp-\codebreak{}idle-\codebreak{}delay*} & \texttt{0.5} & The seconds the cursor rests before what is known of where it stands is shown: what a server found wrong there, the parameters of the call, what the thing is, 0.5 as standard.\\
\texttt{*lsp-\codebreak{}signature-\codebreak{}help*} & \texttt{t} & True for the echo area to show the parameters of the call the cursor is in, as the language server knows them, once the cursor rests there.\\
\texttt{*lsp-\codebreak{}hover-\codebreak{}echo*} & \texttt{t} & True for the echo area to show what the thing the cursor rests on is (its type, its signature), as the language server knows it, when the cursor is in no call whose parameters are shown.\\
\texttt{*lsp-\codebreak{}highlight-\codebreak{}symbol*} & \texttt{t} & True for the name the cursor rests on, and its other uses the language server finds in the file, to be shown in :symbol-highlight.\\
\texttt{*lsp-\codebreak{}diagnostics-\codebreak{}display*} & \texttt{:box} & How what a language server found where the cursor rests is said: :box, in a box under the cursor; :echo, in the echo area; NIL, not at all (the list of M-x lsp-diagnostics still has it).\\
\texttt{*lsp-\codebreak{}diagnostics-\codebreak{}style*} & \texttt{:wave} & How the text a language server finds something wrong in is marked: :wave, errors and warnings underlined wavy, information and hints dotted; :dotted, all dotted; :line, all with a straight line; :background, the text on a ground of its kind’s colour (the theme’s error, warning and information grounds), not underlined.\\
\texttt{*lsp-\codebreak{}inlay-\codebreak{}hints*} & \texttt{t} & True for language servers’ inlay hints to be shown within the code: the types of bindings, the names of the parameters arguments go to.\\
\texttt{*lsp-\codebreak{}semantic-\codebreak{}tokens*} & \texttt{nil} & True for code to be coloured as its language server sees it: each name in the face \textbf{LSP-SEMANTIC-FACES} gives its kind, over the mode’s colours.\\
\texttt{*lsp-\codebreak{}semantic-\codebreak{}faces*} & \texttt{(("namespace" . :namespace) ("type" . :t...} & The face each kind of name a language server tells apart is coloured in, as (KIND .\\
\texttt{*lsp-\codebreak{}code-\codebreak{}lenses*} & \texttt{nil} & True for what language servers say of lines (3 references, Run test) to be shown faint at each line’s end.\\
\texttt{*lsp-\codebreak{}outline*} & \texttt{t} & True for M-x buffer-outline to list a file’s names as its language server knows them (functions, types, their fields and methods, nested), where it has one that does; else as the mode finds its headings.\\
\texttt{*lsp-\codebreak{}wait*} & \texttt{2} & The seconds a command waits for a language server’s answer it needs at once (a completion, a definition), a key stopping the wait sooner, 2 as standard.\\
\texttt{*lsp-\codebreak{}watch-\codebreak{}files*} & \texttt{t} & True for language servers that ask to be told of their project’s files changed on disk, made or deleted, by programs other than the editor, to be told.\\
\texttt{*lsp-\codebreak{}watch-\codebreak{}skipped*} & \texttt{("target" "node\_\codebreak{}modules" "\_\_\codebreak{}pycache\_\_" "...} & The folders of a project whose files are never watched for its language servers, by name (hidden folders, whose names begin with a dot, never are): what a build makes.\\
\texttt{*language-\codebreak{}formatters*} & \texttt{nil} & The formatters of languages you choose yourself, in place of the language’s usual one: each (LANGUAGE PROGRAM ARGUMENT\ldots{}), the code given on its standard input and taken back laid out from its standard output; (LANGUAGE) alone for none, the language server or Tuatara’s own indentation laying it out instead.\\
\texttt{*format-\codebreak{}on-\codebreak{}save*} & \texttt{nil} & The languages whose buffers are formatted (M-x format-buffer) before each save, the save going on once the text is laid out: a list of language names, (rust python), or T for every language with a formatter installed.\\
\texttt{*format-\codebreak{}in-\codebreak{}sandbox*} & \texttt{t} & True for formatters to run shut in the sandbox: the project’s folder and the formatter’s own seen, read only, nothing written, no network.\\
\texttt{*format-\codebreak{}seconds*} & \texttt{30} & The seconds a formatter may take before it is stopped and nothing is changed, 30 as standard; a language server laying out code that has not answered by then is no longer waited for.\\
\texttt{*word-\codebreak{}completion-\codebreak{}keep-\codebreak{}case*} & \texttt{t} & True to find words that start with what is typed whatever the case, and to keep the case typed: with “foo” typed, a word Foobar is put in as foobar, and with “FOO” typed, it is put in as FOOBAR.\\
\bottomrule
\end{xltabular}

\begin{xltabular}{\textwidth}{@{}l >{\raggedright\arraybackslash\hsize=1.000\hsize}X@{}}
\toprule
\textbf{Face} & \textbf{What it paints}\\
\midrule
\endhead
\texttt{:diagnostic-\codebreak{}error} & The text a language server finds an error in: underlined, wavy.\\
\texttt{:diagnostic-\codebreak{}warning} & The text a language server warns of: underlined, wavy.\\
\texttt{:diagnostic-\codebreak{}info} & The text a language server says something of: underlined, dotted.\\
\texttt{:diagnostic-\codebreak{}hint} & The text a language server hints at: underlined, dotted.\\
\texttt{:diagnostic-\codebreak{}error-\codebreak{}ground} & The text a language server finds an error in, with \textbf{lsp-diagnostics-style} :background: on the theme’s error ground.\\
\texttt{:diagnostic-\codebreak{}warning-\codebreak{}ground} & The text a language server warns of, with \textbf{lsp-diagnostics-style} :background: on the theme’s warning ground.\\
\texttt{:diagnostic-\codebreak{}info-\codebreak{}ground} & The text a language server says something of, with \textbf{lsp-diagnostics-style} :background: on the theme’s information ground.\\
\texttt{:diagnostic-\codebreak{}hint-\codebreak{}ground} & The text a language server hints at, with \textbf{lsp-diagnostics-style} :background: faint.\\
\texttt{:symbol-\codebreak{}highlight} & The name the cursor rests on, and its other uses the language server finds in the window: bold.\\
\texttt{:inlay-\codebreak{}hint} & What a language server shows within the code, no text of it (a type, a parameter’s name): faint.\\
\texttt{:code-\codebreak{}lens} & What a language server says of a line (3 references, Run test), at its end: as an inlay hint.\\
\texttt{:fold-\codebreak{}ellipsis} & The ellipsis standing for folded text, as an Org subtree folded shows.\\
\bottomrule
\end{xltabular}

\subsection{Snippets}
\label{sec:333}
A \emph{snippet} is a piece of text with places to fill in, which you put in by
typing a short key and pressing \texttt{TAB}. It saves typing the same few lines every
time, such as the top of a script. Nothing in a snippet is ever run.

Your snippets are files in \texttt{\textasciitilde{}/.\codebreak{}config/\codebreak{}tuatara/\codebreak{}snippets/}, in a folder named
after the major mode: \texttt{python-\codebreak{}mode/}, \texttt{prog-\codebreak{}mode/} for every programming
language, or \texttt{fundamental-\codebreak{}mode/} for every buffer. A snippet file starts with
lines that name it and give its key, then a line of two dashes, then the text:

\begin{Verbatim}[breaklines,breakautoindent=false,breakindent=1em,breakafter=/-.\,;:=,commandchars=\\\{\}]
# name: Script header
# key: hdr
# --
# $\char123{}1:What this file does\char125{}
# Written on $CURRENT_YEAR-$CURRENT_MONTH-$CURRENT_DATE
$0
\end{Verbatim}

In the text, \texttt{\$1}, \texttt{\$2} and so on are the places to fill in, in order.
\texttt{\$\{1:\codebreak{}words\}} is a place that starts with some words, which your typing replaces.
\texttt{\$\{1|\codebreak{}a,\codebreak{}b|\}} offers a choice. \texttt{\$0} is where the cursor ends up. A name with a
dollar stands for a value: \texttt{\$CURRENT\_\codebreak{}YEAR}, \texttt{\$CURRENT\_\codebreak{}MONTH} and \texttt{\$CURRENT\_\codebreak{}DATE}
for today, \texttt{\$TM\_\codebreak{}FILENAME} for the file, \texttt{\$AUTHOR} for the name in
\texttt{*user-\codebreak{}full-\codebreak{}name*}.

Type the key and press \texttt{TAB}, and the text goes in with the cursor in the first
place. \texttt{TAB} moves on to the next place and \texttt{S-TAB} back to the one before.
After the last place the cursor goes to \texttt{\$0}. \texttt{C-c y} (\texttt{insert-\codebreak{}snippet}) chooses
a snippet from a list by name, each with its key, and with text selected the
selection goes where the snippet has \texttt{\$TM\_\codebreak{}SELECTED\_\codebreak{}TEXT}. \texttt{M-x copy-\codebreak{}snippet}
copies one of Tuatara’s own into your folder so that you can change it, and
yours takes the place of Tuatara’s.

\begin{try}
\textbf{Try it, with a Python script header}
\begin{enumerate}
\item In a terminal, make the snippet. Type this as one command.

\begin{Verbatim}[breaklines,breakautoindent=false,breakindent=1em,breakafter=/-.\,;:=,commandchars=\\\{\}]
   mkdir -p ~/.config/tuatara/snippets/python-mode && printf '# name: Script header\char92{}n# key: hdr\char92{}n# --\char92{}n# $\char123{}1:What this file does\char125{}\char92{}n# Written on $CURRENT_YEAR-$CURRENT_MONTH-$CURRENT_DATE\char92{}n$0\char92{}n' > ~/.config/tuatara/snippets/python-mode/header
\end{Verbatim}
\item Open \texttt{\textasciitilde{}/code-\codebreak{}practice/\codebreak{}sum.py}. Press \texttt{g g} and \texttt{O} to open a line above the
   first. Type \texttt{hdr} and press \texttt{TAB}.
\item Two comment lines appear. The first has \texttt{What this file does} ready to be
   typed over. Type \texttt{Adds up some numbers}, and press \texttt{TAB}. The cursor is on
   the line after, which is where \texttt{\$0} was.
\item Press \texttt{C-g} and \texttt{u} (non-modal: \texttt{C-/}) until the lines are gone. Delete the
   snippet file if you do not want to keep it.
\end{enumerate}
\end{try}

\subsubsection{Commands and settings}
\label{sec:332}
\begin{xltabular}{\textwidth}{@{}l >{\raggedright\arraybackslash\hsize=0.622\hsize}X l >{\raggedright\arraybackslash\hsize=1.378\hsize}X@{}}
\toprule
\textbf{Modal} & \textbf{Non-modal} & \textbf{Command} & \textbf{What it does}\\
\midrule
\endhead
none & \texttt{C-c y} & \texttt{insert-\codebreak{}snippet} & Put in a snippet chosen by its name among those of the current buffer’s mode (each listed with its key), at point; with the selection active, in its place, the selected text where it has \$TM\_SELECTED\_TEXT.\\
none & \texttt{M-x copy-\codebreak{}snippet} & \texttt{copy-\codebreak{}snippet} & Copy one of Tuatara’s own snippets, chosen by its name, into your snippets folder (the first of \textbf{snippet-directories}, in the folder of its mode) and visit it: your copy, changed as you like, takes the place of Tuatara’s.\\
none & \texttt{M-x snippet-\codebreak{}next-\codebreak{}field} & \texttt{snippet-\codebreak{}next-\codebreak{}field} & Go to the next field of the snippet being filled in; past the last, end it where it puts the cursor.\\
none & \texttt{M-x snippet-\codebreak{}previous-\codebreak{}field} & \texttt{snippet-\codebreak{}previous-\codebreak{}field} & Go back to the field of the snippet before the one the cursor is in.\\
\bottomrule
\end{xltabular}

\begin{xltabular}{\textwidth}{@{}l l >{\raggedright\arraybackslash\hsize=1.000\hsize}X@{}}
\toprule
\textbf{Setting} & \textbf{Standard} & \textbf{What it changes}\\
\midrule
\endhead
\texttt{*snippet-\codebreak{}directories*} & \texttt{nil} & The folders your snippets are in, each with a folder for each major mode (org-mode/, prog-mode/, fundamental-mode/), a snippet a file in it; NIL, as standard, for snippets/ in your configuration folder (\textasciitilde{}/.config/tuatara/snippets/, XDG\_CONFIG\_HOME’s when it is set).\\
\texttt{*snippet-\codebreak{}tab*} & \texttt{t} & True for a snippet’s key before the cursor and TAB to put the snippet in.\\
\texttt{*user-\codebreak{}full-\codebreak{}name*} & \texttt{nil} & Your full name, as “Ann Taranaki”: the author of an exported document with no \#+author:, \%author\% in a note template and \$AUTHOR in a snippet.\\
\bottomrule
\end{xltabular}

\subsection{Running a build or the tests}
\label{sec:335}
Writing code is a loop. You change something, run it, read what went wrong, go
to each place, fix it, and run it again. \texttt{M-x compile} runs the build or the
tests from inside the editor and turns every error in the output into a place
you can jump to, so the loop never leaves the keyboard.

\texttt{M-x compile} asks for a command. Type anything you would type in a shell: \texttt{make
-k}, \texttt{cargo build}, \texttt{python3 -m unittest} or \texttt{sh count.sh}. Tuatara runs it with
your shell from the root of the project the current file is in, or from the
file’s own folder when the file is in no project. The first time it offers \texttt{make
-k} (the setting \texttt{*compile-\codebreak{}command*}), and after that the command you used last.
\texttt{M-p} goes back through earlier ones. The command runs with no terminal, so a
program that asks for a password, such as \texttt{sudo} or \texttt{ssh}, cannot reach your
screen. Chapter 17 is for those.

A build runs the project’s own code, so in a project you have not trusted
Tuatara asks first, the same question the language servers ask. Answer \texttt{n} or
\texttt{l} and nothing runs.

The output appears in a buffer called \texttt{*compile*} as it comes, and you go on
working meanwhile. Read the first error and start fixing it while the rest is
still being built. When the command ends, a last line says how it ended and how
long it took. Only one compile runs at a time. If you start a second, Tuatara
asks and starts it when the first has finished. \texttt{C-c C-k} in \texttt{*compile*} stops a
command that is still running, \texttt{r} there runs it again, and \texttt{q} closes the
window. The output kept is capped at \texttt{*compile-\codebreak{}output-\codebreak{}limit*} octets.

Lines in the output that name a file and a line are \emph{places}. Tuatara knows the
form of the errors of GCC and Clang, of rustc and cargo, of Python, and the
\texttt{file:\codebreak{}line:} of \texttt{grep -n} and \texttt{make}. Errors are in the \texttt{:compile-\codebreak{}error} colour,
warnings in \texttt{:compile-\codebreak{}warning} and notes in \texttt{:compile-\codebreak{}note}. \texttt{M-g n}
(\texttt{next-\codebreak{}place}) goes to the next place and \texttt{M-g p} (\texttt{previous-\codebreak{}place}) to the one
before, from any buffer. At the end of the list the echo area says there are no
more. Tuatara counts columns from 0, while compilers usually count from 1, so a
compiler’s column 5 shows as column 4 in the mode line.

The same two keys go through any list of places made last: the lines found by a
search of files (\texttt{M-s g}) or the problems a language server found in a buffer.
In \texttt{*compile*}, \texttt{RET} on a line opens that place.

\begin{try}
\textbf{Try it, with a mistake in the Python script}
\begin{enumerate}
\item Open \texttt{\textasciitilde{}/code-\codebreak{}practice/\codebreak{}sum.py}. Press \texttt{g g}, then \texttt{j j} to reach the \texttt{print}
   line, then \texttt{A}, type \texttt{x} and press \texttt{C-g}. The line now reads
   \texttt{print(add\_\codebreak{}up([3, 4, 5]))x}, which is not valid Python.
\item Press \texttt{C-x C-s}. Press \texttt{M-x}, type \texttt{compile} and press \texttt{RET}. Press \texttt{C-a} and
   \texttt{C-k}, type \texttt{python3 sum.py} and press \texttt{RET}. If asked to trust the folder,
   type \texttt{y}, then \texttt{y} and \texttt{RET}.
\item \texttt{*compile*} opens under the file. It shows Python’s own words, a line
   beginning \texttt{File} with the file’s name and \texttt{line 3}, then the line itself and
   \texttt{SyntaxError}. The last line says the command failed with a non-zero exit
   code.
\item Press \texttt{M-g n}. The cursor goes to line 3 of \texttt{sum.py}, and the echo area shows
   the line from the output.
\item Press \texttt{u} (non-modal: \texttt{C-/}) to take the \texttt{x} away, \texttt{C-x C-s}, then \texttt{M-x},
   \texttt{compile} and \texttt{RET}, and \texttt{RET} again to run the command offered. The output
   is \texttt{12}.
\end{enumerate}
\end{try}

\subsubsection{Commands and settings}
\label{sec:334}
\begin{xltabular}{\textwidth}{@{}l l l >{\raggedright\arraybackslash\hsize=1.000\hsize}X@{}}
\toprule
\textbf{Modal} & \textbf{Non-modal} & \textbf{Command} & \textbf{What it does}\\
\midrule
\endhead
none & \texttt{M-x compile} & \texttt{compile} & Run a command asked for, a build or the tests (make -k, cargo build, pytest), with your shell from the root of the project the current buffer’s file is in (see \textbf{PROJECT-ROOT-MARKERS}), or from its directory when it is in none, and show what it prints in \textbf{compile} as it comes, the editor free meanwhile.\\
none & \texttt{C-c C-k} & \texttt{compile-\codebreak{}stop} & Stop the command running in this compile’s output, and every program it started: what it printed so far stays.\\
none & \texttt{M-g n} & \texttt{next-\codebreak{}place} & Go to the next place of the list of places used last: the next line naming a file and line in a compile’s output (M-x compile), the next line a search of files found (M-s g), or the next place the language servers found something wrong in a buffer.\\
none & \texttt{M-g p} & \texttt{previous-\codebreak{}place} & Go to the place before the one gone to last in the list of places used last: a compile’s output (M-x compile), the lines a search of files found (M-s g), or the places the language servers found something wrong in a buffer, as M-g n goes to the next.\\
\bottomrule
\end{xltabular}

\begin{xltabular}{\textwidth}{@{}l l >{\raggedright\arraybackslash\hsize=1.000\hsize}X@{}}
\toprule
\textbf{Setting} & \textbf{Standard} & \textbf{What it changes}\\
\midrule
\endhead
\texttt{*compile-\codebreak{}command*} & \texttt{"make -k "} & The command M-x compile offers the first time, before you have typed one; after that it offers the one typed last.\\
\texttt{*compile-\codebreak{}output-\codebreak{}limit*} & \texttt{268435456} & The most octets of a command’s output M-x compile keeps in \textbf{compile}: 256 MB as standard.\\
\bottomrule
\end{xltabular}

\begin{xltabular}{\textwidth}{@{}l >{\raggedright\arraybackslash\hsize=1.000\hsize}X@{}}
\toprule
\textbf{Face} & \textbf{What it paints}\\
\midrule
\endhead
\texttt{:compile-\codebreak{}error} & The lines of a compile’s output (M-x compile) that say an error, and the line saying the command failed: as errors are marked in text.\\
\texttt{:compile-\codebreak{}warning} & The lines of a compile’s output (M-x compile) that say a warning: as warnings are marked in text.\\
\texttt{:compile-\codebreak{}note} & The lines of a compile’s output (M-x compile) that say a note or give help: as information is marked in text.\\
\bottomrule
\end{xltabular}

\subsection{Merge conflicts}
\label{sec:337}
When two people change the same lines, or you change them on two branches, Git
cannot choose which to keep. It writes both versions into the file between
marker lines and stops, and you decide. The file then holds this:

\begin{Verbatim}[breaklines,breakautoindent=false,breakindent=1em,breakafter=/-.\,;:=,commandchars=\\\{\}]
<<<<<<< HEAD
name = "Aroha"
=======
name = "Hemi"
>>>>>>> other-branch
\end{Verbatim}

The first version, after \texttt{<<<<<<<}, is yours in a merge. In a rebase it is the
branch the commits are being made again onto. The second version, after \texttt{=====},
is the other one. When Git is set to show it, a third part after \texttt{|||||||} has
the lines as they were before either change, which is called the \emph{base}.

Putting this right by hand means deleting marker lines with care, and a marker
left behind breaks the build. Tuatara looks through each file you open for
conflicts. It colours the first side, the second side, the base and the marker
lines, and the mode line says how many there are. The keys under \texttt{C-c \textasciicircum{}} settle
each one in a single step, and each step goes on to the next conflict.

\begin{longtable}{@{}l l l@{}}
\toprule
\textbf{Modal} & \textbf{Non-modal} & \textbf{What it does}\\
\midrule
\endhead
none & \texttt{C-c \textasciicircum{} n} & Go to the next conflict.\\
none & \texttt{C-c \textasciicircum{} p} & Go to the one before.\\
none & \texttt{C-c \textasciicircum{} u} & Keep the first side, the upper one.\\
none & \texttt{C-c \textasciicircum{} l} & Keep the second side, the lower one.\\
none & \texttt{C-c \textasciicircum{} a} & Keep both sides, first then second.\\
none & \texttt{C-c \textasciicircum{} b} & Keep the base, the lines from before either change.\\
none & \texttt{C-c \textasciicircum{} =} & Show how the two sides differ, in a buffer of its own.\\
\bottomrule
\end{longtable}

Every choice is one step for \texttt{u} (non-modal: \texttt{C-/}), so a wrong choice is easy
to take back.

How the merge began, which branch is which, and how to finish it afterwards, is
in chapter 15. This chapter only covers the file.

\begin{try}
\textbf{Try it, with a small conflict}
\begin{enumerate}
\item In a terminal, make a file with a conflict in it. Type this as one command.

\begin{Verbatim}[breaklines,breakautoindent=false,breakindent=1em,breakafter=/-.\,;:=,commandchars=\\\{\}]
   printf 'greeting = "hello"\char92{}n<<<<<<<\discretionary{}{}{} HEAD\char92{}nname = "Aroha"\char92{}n=======\char92{}nname = "Hemi"\char92{}n>>>>>>> other-branch\char92{}nprint(gr\discretionary{}{}{}eeting, name)\char92{}n' > ~/code-practice/greet.py
\end{Verbatim}
\item Open \texttt{\textasciitilde{}/code-\codebreak{}practice/\codebreak{}greet.py}. The two sides and the marker lines are
   coloured, and the mode line says there is a conflict.
\item Press \texttt{g g} (non-modal: \texttt{M-<}) and then \texttt{C-c \textasciicircum{} n}. The cursor goes to the
   line with the \texttt{<<<<<<<} marker, and the echo area says which conflict it is
   and how many there are.
\item Press \texttt{C-c \textasciicircum{} =}. A buffer shows the lines only one side has. Press \texttt{q} to
   close it.
\item Press \texttt{C-c \textasciicircum{} l}. The marker lines and the first side are gone, and
   \texttt{name = "Hemi"} stays. Press \texttt{u} (non-modal: \texttt{C-/}) and the conflict is back.
   Press \texttt{C-c \textasciicircum{} a} to keep both lines instead, and \texttt{u} again.
\end{enumerate}
\end{try}

\subsubsection{Commands and faces}
\label{sec:336}
\begin{xltabular}{\textwidth}{@{}l >{\raggedright\arraybackslash\hsize=0.667\hsize}X l >{\raggedright\arraybackslash\hsize=1.333\hsize}X@{}}
\toprule
\textbf{Modal} & \textbf{Non-modal} & \textbf{Command} & \textbf{What it does}\\
\midrule
\endhead
none & \texttt{M-x conflict-\codebreak{}markers-\codebreak{}mode} & \texttt{conflict-\codebreak{}markers-\codebreak{}mode} & Show the conflicts a merge left in the current buffer, and give it the keys that put them right, or stop.\\
none & \texttt{M-x conflict-\codebreak{}next} & \texttt{conflict-\codebreak{}next} & Go to the start of the next conflict, its <<<<<<< line, and say which of how many it is.\\
none & \texttt{M-x conflict-\codebreak{}previous} & \texttt{conflict-\codebreak{}previous} & Go to the start of the conflict before point, its <<<<<<< line, and say which of how many it is.\\
none & \texttt{M-x conflict-\codebreak{}keep-\codebreak{}first-\codebreak{}side} & \texttt{conflict-\codebreak{}keep-\codebreak{}first-\codebreak{}side} & Put the conflict point is in right by keeping its first side, the lines after its <<<<<<< line, and taking away the rest and its marker lines; then go to the next conflict.\\
none & \texttt{M-x conflict-\codebreak{}keep-\codebreak{}second-\codebreak{}side} & \texttt{conflict-\codebreak{}keep-\codebreak{}second-\codebreak{}side} & Put the conflict point is in right by keeping its second side, the lines after its \texttt{=====} line, and taking away the rest and its marker lines; then go to the next conflict.\\
none & \texttt{M-x conflict-\codebreak{}keep-\codebreak{}both-\codebreak{}sides} & \texttt{conflict-\codebreak{}keep-\codebreak{}both-\codebreak{}sides} & Put the conflict point is in right by keeping both its sides, the first then the second, the base and the marker lines taken away; then go to the next conflict.\\
none & \texttt{M-x conflict-\codebreak{}keep-\codebreak{}base} & \texttt{conflict-\codebreak{}keep-\codebreak{}base} & Put the conflict point is in right by keeping its base, the lines after its \(\vert{}\)\(\vert{}\)\(\vert{}\)\(\vert{}\)\(\vert{}\)\(\vert{}\)\(\vert{}\) line, as they were before either side changed them, the sides and the marker lines taken away; then go to the next conflict.\\
none & \texttt{M-x conflict-\codebreak{}show-\codebreak{}sides-\codebreak{}diff} & \texttt{conflict-\codebreak{}show-\codebreak{}sides-\codebreak{}diff} & Show how the two sides of the conflict point is in differ, in a buffer of their own in another window, which q closes: the lines only the first side has after -, those only the second has after +, three lines around each change, the sides named by their marker lines.\\
\bottomrule
\end{xltabular}

\begin{xltabular}{\textwidth}{@{}l >{\raggedright\arraybackslash\hsize=1.000\hsize}X@{}}
\toprule
\textbf{Face} & \textbf{What it paints}\\
\midrule
\endhead
\texttt{:conflict-\codebreak{}first} & The first side of a conflict a merge left in a file, the lines after its <<<<<<< line: as the lines a difference takes away unless set.\\
\texttt{:conflict-\codebreak{}second} & The second side of a conflict a merge left in a file, the lines after its \texttt{=====} line: as the lines a difference adds unless set.\\
\texttt{:conflict-\codebreak{}base} & The base of a conflict a merge left in a file, the lines after its \(\vert{}\)\(\vert{}\)\(\vert{}\)\(\vert{}\)\(\vert{}\)\(\vert{}\)\(\vert{}\) line, as they were before either side changed them.\\
\texttt{:conflict-\codebreak{}marker} & The marker lines of a conflict a merge left in a file (M-x conflict-markers-mode): <<<<<<<, \(\vert{}\)\(\vert{}\)\(\vert{}\)\(\vert{}\)\(\vert{}\)\(\vert{}\)\(\vert{}\) and >>>>>>> with the names after them, and \texttt{=====}, out to the window’s edge.\\
\bottomrule
\end{xltabular}

\subsection{Lisp development}
\label{sec:340}
The Lisp you type runs in the editor itself. That is how the init file works,
and it is a quick way to try a small function or to look at a value. A file
ending in \texttt{.lisp} opens in Lisp mode, which has keys for running what you have
written. No language server is needed: the running Lisp knows its own
definitions, so \texttt{M-.} goes to them.

The keys that run text work only in Lisp buffers: a \texttt{.lisp} file, the scratch
buffer \texttt{*lisp-\codebreak{}scratch*} (\texttt{M-x lisp-\codebreak{}scratch}) and your init file.

\begin{itemize}
\item \texttt{C-x C-e} (\texttt{eval-\codebreak{}last-\codebreak{}expression}) runs the expression just before the cursor
  and shows its value in the echo area. A whole number is also shown in octal
  and hexadecimal.
\item \texttt{C-M-x} (\texttt{eval-\codebreak{}top-\codebreak{}level-\codebreak{}form}) runs the whole top-level form the cursor is
  in, that is, the one whose bracket is at the start of a line. A \texttt{defvar} in it
  takes its new value even when the variable already had one.
\item \texttt{M-x eval-\codebreak{}region} runs every form in the selection, and \texttt{M-x eval-\codebreak{}buffer}
  every form of the buffer, one after another, as a file’s forms are run when it
  is loaded.
\item \texttt{C-j} (\texttt{eval-\codebreak{}print-\codebreak{}last-\codebreak{}expression}) in \texttt{*lisp-\codebreak{}scratch*} runs the expression
  and writes its value into the buffer.
\item \texttt{M-:} (\texttt{eval-\codebreak{}expression}) works anywhere. It asks for an expression, shows its
  value, and runs only what you typed.
\item \texttt{C-c C-m} (\texttt{pp-\codebreak{}macroexpand-\codebreak{}last-\codebreak{}expression}) shows a macro call expanded once,
  laid out in a window of its own.
\end{itemize}

A file opened safely (chapter 12) runs nothing: the keys say so and do nothing.

\texttt{M-x ielm} opens \texttt{*ielm*}, a prompt where you type Lisp and read the answer at
once. The modal keys are off there, so you type without pressing \texttt{i}. \texttt{RET} runs
the input when it is a whole expression, and otherwise starts a new line of it.
\texttt{C-j} runs it as it is. \texttt{M-p} and \texttt{M-n} go back and forward through what you
typed before, \texttt{C-c C-u} clears the input, and \texttt{TAB} completes a name. \texttt{C-c C-b}
chooses the buffer that the Lisp you type sees as current, so that \texttt{(insert
"x")} puts text there.

When something fails, \texttt{M-x toggle-\codebreak{}debug-\codebreak{}on-\codebreak{}error} makes the next error show the
calls that were running when it happened, innermost first, in \texttt{*Backtrace*}.
The setting \texttt{*debug-\codebreak{}on-\codebreak{}error*} makes that the standard.

\begin{try}
\textbf{Try it, with a Lisp function that adds numbers}
\begin{enumerate}
\item Open \texttt{\textasciitilde{}/code-\codebreak{}practice/\codebreak{}sum.\codebreak{}lisp}, press \texttt{i} and type these lines. \texttt{RET}
   indents each line for you.

\begin{Verbatim}[breaklines,breakautoindent=false,breakindent=1em,breakafter=/-.\,;:=,commandchars=\\\{\}]
   (defun add-up (numbers)
     (reduce #'+ numbers))

   (add-up '(3 4 5))
\end{Verbatim}

   Press \texttt{C-g}.
\item Put the cursor inside the \texttt{defun} (\texttt{g g}) and press \texttt{C-M-x}. The echo area
   shows the form and then \texttt{ADD-UP}. The function now exists in the editor.
\item Go to the last line with \texttt{g e} (non-modal: \texttt{M->}), move to its end and press
   \texttt{C-x C-e}. The echo area shows the value, \texttt{12}, and then the number in octal
   and hexadecimal.
\item Put the cursor on \texttt{add-up} in the last line, with \texttt{C-a} (modal: \texttt{g h}) and
   \texttt{l l}, and press \texttt{M-.}. It goes to the \texttt{defun}. Press \texttt{M-,} to come back.
\item Press \texttt{M-x}, type \texttt{ielm} and press \texttt{RET}. After the prompt type
   \texttt{(add-up '(10 20))} and press \texttt{RET}. The answer is \texttt{30}. Press \texttt{M-p} and the
   same input comes back.
\item Press \texttt{M-x}, type \texttt{lisp-\codebreak{}scratch} and press \texttt{RET} for a buffer of your own to
   try things in. Press \texttt{g e} and \texttt{i}, type \texttt{(+ 1 2)} and press \texttt{C-j}. The
   value is written on the next line.
\end{enumerate}
\end{try}

\subsubsection{Commands}
\label{sec:338}
\begin{xltabular}{\textwidth}{@{}l >{\raggedright\arraybackslash\hsize=0.649\hsize}X l >{\raggedright\arraybackslash\hsize=1.351\hsize}X@{}}
\toprule
\textbf{Modal} & \textbf{Non-modal} & \textbf{Command} & \textbf{What it does}\\
\midrule
\endhead
none & \texttt{C-x C-e} & \texttt{eval-\codebreak{}last-\codebreak{}expression} & Evaluate the expression before point in the running editor and show its value in the echo area; with a prefix argument (C-u), insert the value after point.\\
none & \texttt{C-M-x} & \texttt{eval-\codebreak{}top-\codebreak{}level-\codebreak{}form} & Evaluate the top-level form around point (the one starting with an open bracket at the start of a line) in the running editor and show its value.\\
none & \texttt{M-x eval-\codebreak{}region} & \texttt{eval-\codebreak{}region} & Evaluate the Lisp forms in the selection one after another, as a file’s are (IN-PACKAGE carries on to the next), and show the last one’s value.\\
none & \texttt{M-x eval-\codebreak{}buffer} & \texttt{eval-\codebreak{}buffer} & Evaluate every Lisp form in the current buffer one after another, as a file’s are, and show the last one’s value; not a buffer of more than \textbf{REWRITE-THRESHOLD} octets.\\
none & \texttt{C-j} & \texttt{eval-\codebreak{}print-\codebreak{}last-\codebreak{}expression} & Evaluate the expression before point in the running editor and put what it printed and its value on the lines after it.\\
none & \texttt{M-:} & \texttt{eval-\codebreak{}expression} & Read a Lisp expression in the minibuffer and evaluate it in the running editor, in the current buffer’s package; show its value, or with a prefix argument (C-u) insert it at point.\\
none & \texttt{C-c C-m}, \texttt{C-c RET} & \texttt{pp-\codebreak{}macroexpand-\codebreak{}last-\codebreak{}expression} & Show what the macro call at point expands to, once (MACROEXPAND-1), pretty-printed in \textbf{Pp Macroexpand Output}: the call beginning at point when an opening bracket is there, else the one before point.\\
none & \texttt{M-x lisp-\codebreak{}scratch} & \texttt{lisp-\codebreak{}scratch} & Show \textbf{lisp-scratch}, the Lisp buffer that is not saved, where C-x C-e evaluates in the running editor and C-j puts what the expression before point printed and its value on the lines after it; when it was killed, a new one.\\
none & \texttt{M-x ielm} & \texttt{ielm} & Show \textbf{ielm}, where Lisp typed after the prompt is evaluated in the running editor: made with a prompt the first time, as it was after.\\
none & \texttt{C-j} & \texttt{ielm-\codebreak{}send-\codebreak{}input} & Evaluate the input after \textbf{ielm}’s last prompt as it is, whole or not.\\
none & \texttt{RET} & \texttt{ielm-\codebreak{}return} & RET in \textbf{ielm}: after the last prompt, evaluate the input when it is whole expressions, else begin a new line of it; on an earlier input, evaluate that again.\\
none & \texttt{M-p} & \texttt{ielm-\codebreak{}previous-\codebreak{}input} & Put the input before the one shown in \textbf{ielm}’s history in place of the input.\\
none & \texttt{M-n} & \texttt{ielm-\codebreak{}next-\codebreak{}input} & Put the input after the one shown in \textbf{ielm}’s history in place of the input.\\
none & \texttt{C-c C-u} & \texttt{ielm-\codebreak{}kill-\codebreak{}input} & Take away the input after \textbf{ielm}’s last prompt.\\
none & \texttt{C-c C-b} & \texttt{ielm-\codebreak{}change-\codebreak{}working-\codebreak{}buffer} & Choose the working buffer of \textbf{ielm}: the one the Lisp typed there sees as current, (insert “x”) putting text in it.\\
none & \texttt{TAB} & \texttt{ielm-\codebreak{}tab} & TAB in \textbf{ielm}: where only blanks come before point on its line of the input, indent the line as Lisp (never the prompt’s line); else complete the name at point from a list beside it, as C-M-i does.\\
none & \texttt{M-x toggle-\codebreak{}debug-\codebreak{}on-\codebreak{}error} & \texttt{toggle-\codebreak{}debug-\codebreak{}on-\codebreak{}error} & Switch debug-on-error on or off: while it is on, an error in a command, or in Lisp evaluated in the editor (C-x C-e, C-M-x, M-:, C-j), shows the calls under way when it was signalled in \textbf{Backtrace}, innermost first, besides the message.\\
\bottomrule
\end{xltabular}

\subsubsection{Settings}
\label{sec:339}
\begin{xltabular}{\textwidth}{@{}l l >{\raggedright\arraybackslash\hsize=1.000\hsize}X@{}}
\toprule
\textbf{Setting} & \textbf{Standard} & \textbf{What it changes}\\
\midrule
\endhead
\texttt{*debug-\codebreak{}on-\codebreak{}error*} & \texttt{nil} & True to show the calls under way when a command or an evaluation meets an error, in \textbf{Backtrace}; M-x toggle-debug-on-error switches it.\\
\bottomrule
\end{xltabular}

\subsection{The commands of each language}
\label{sec:341}
Each language has a command that switches the current buffer to its major mode,
for the rare time Tuatara guesses wrongly or you open a file with an unusual
name. Type \texttt{M-x} and the language’s name followed by \texttt{-mode}, for example
\texttt{M-x rust-\codebreak{}mode}. They are all listed here. Chapter 8 explains how to find a
command by its name.

\begin{xltabular}{\textwidth}{@{}l >{\raggedright\arraybackslash\hsize=0.875\hsize}X l >{\raggedright\arraybackslash\hsize=1.125\hsize}X@{}}
\toprule
\textbf{Modal} & \textbf{Non-modal} & \textbf{Command} & \textbf{What it does}\\
\midrule
\endhead
none & \texttt{M-x ada-\codebreak{}mode} & \texttt{ada-\codebreak{}mode} & Switch the current buffer to ada-mode.\\
none & \texttt{M-x adl-\codebreak{}mode} & \texttt{adl-\codebreak{}mode} & Switch the current buffer to adl-mode.\\
none & \texttt{M-x agda2-\codebreak{}mode} & \texttt{agda2-\codebreak{}mode} & Switch the current buffer to agda2-mode.\\
none & \texttt{M-x alloy-\codebreak{}mode} & \texttt{alloy-\codebreak{}mode} & Switch the current buffer to alloy-mode.\\
none & \texttt{M-x amber-\codebreak{}mode} & \texttt{amber-\codebreak{}mode} & Switch the current buffer to amber-mode.\\
none & \texttt{M-x asm-\codebreak{}mode} & \texttt{asm-\codebreak{}mode} & Switch the current buffer to asm-mode.\\
none & \texttt{M-x astro-\codebreak{}mode} & \texttt{astro-\codebreak{}mode} & Switch the current buffer to astro-mode.\\
none & \texttt{M-x awk-\codebreak{}mode} & \texttt{awk-\codebreak{}mode} & Switch the current buffer to awk-mode.\\
none & \texttt{M-x bash-\codebreak{}mode} & \texttt{bash-\codebreak{}mode} & Switch the current buffer to bash-mode.\\
none & \texttt{M-x basic-\codebreak{}mode} & \texttt{basic-\codebreak{}mode} & Switch the current buffer to basic-mode.\\
none & \texttt{M-x bass-\codebreak{}mode} & \texttt{bass-\codebreak{}mode} & Switch the current buffer to bass-mode.\\
none & \texttt{M-x bazel-\codebreak{}starlark-\codebreak{}mode} & \texttt{bazel-\codebreak{}starlark-\codebreak{}mode} & Switch the current buffer to bazel-starlark-mode.\\
none & \texttt{M-x bicep-\codebreak{}mode} & \texttt{bicep-\codebreak{}mode} & Switch the current buffer to bicep-mode.\\
none & \texttt{M-x bitbake-\codebreak{}mode} & \texttt{bitbake-\codebreak{}mode} & Switch the current buffer to bitbake-mode.\\
none & \texttt{M-x blade-\codebreak{}mode} & \texttt{blade-\codebreak{}mode} & Switch the current buffer to blade-mode.\\
none & \texttt{M-x blueprint-\codebreak{}mode} & \texttt{blueprint-\codebreak{}mode} & Switch the current buffer to blueprint-mode.\\
none & \texttt{M-x c++-\codebreak{}mode} & \texttt{c++-\codebreak{}mode} & Switch the current buffer to c++-mode.\\
none & \texttt{M-x c-mode} & \texttt{c-mode} & Switch the current buffer to c-mode.\\
none & \texttt{M-x c3-\codebreak{}mode} & \texttt{c3-\codebreak{}mode} & Switch the current buffer to c3-mode.\\
none & \texttt{M-x caddyfile-\codebreak{}mode} & \texttt{caddyfile-\codebreak{}mode} & Switch the current buffer to caddyfile-mode.\\
none & \texttt{M-x cairo-\codebreak{}mode} & \texttt{cairo-\codebreak{}mode} & Switch the current buffer to cairo-mode.\\
none & \texttt{M-x capnp-\codebreak{}mode} & \texttt{capnp-\codebreak{}mode} & Switch the current buffer to capnp-mode.\\
none & \texttt{M-x cel-\codebreak{}mode} & \texttt{cel-\codebreak{}mode} & Switch the current buffer to cel-mode.\\
none & \texttt{M-x chuck-\codebreak{}mode} & \texttt{chuck-\codebreak{}mode} & Switch the current buffer to chuck-mode.\\
none & \texttt{M-x circom-\codebreak{}mode} & \texttt{circom-\codebreak{}mode} & Switch the current buffer to circom-mode.\\
none & \texttt{M-x clarity-\codebreak{}mode} & \texttt{clarity-\codebreak{}mode} & Switch the current buffer to clarity-mode.\\
none & \texttt{M-x clojure-\codebreak{}mode} & \texttt{clojure-\codebreak{}mode} & Switch the current buffer to clojure-mode.\\
none & \texttt{M-x cmake-\codebreak{}mode} & \texttt{cmake-\codebreak{}mode} & Switch the current buffer to cmake-mode.\\
none & \texttt{M-x concerto-\codebreak{}mode} & \texttt{concerto-\codebreak{}mode} & Switch the current buffer to concerto-mode.\\
none & \texttt{M-x conf-\codebreak{}javaprop-\codebreak{}mode} & \texttt{conf-\codebreak{}javaprop-\codebreak{}mode} & Switch the current buffer to conf-javaprop-mode.\\
none & \texttt{M-x containerfile-\codebreak{}mode} & \texttt{containerfile-\codebreak{}mode} & Switch the current buffer to containerfile-mode.\\
none & \texttt{M-x cpon-\codebreak{}mode} & \texttt{cpon-\codebreak{}mode} & Switch the current buffer to cpon-mode.\\
none & \texttt{M-x crystal-\codebreak{}mode} & \texttt{crystal-\codebreak{}mode} & Switch the current buffer to crystal-mode.\\
none & \texttt{M-x csharp-\codebreak{}mode} & \texttt{csharp-\codebreak{}mode} & Switch the current buffer to csharp-mode.\\
none & \texttt{M-x css-\codebreak{}mode} & \texttt{css-\codebreak{}mode} & Switch the current buffer to css-mode.\\
none & \texttt{M-x cue-\codebreak{}mode} & \texttt{cue-\codebreak{}mode} & Switch the current buffer to cue-mode.\\
none & \texttt{M-x cylc-\codebreak{}mode} & \texttt{cylc-\codebreak{}mode} & Switch the current buffer to cylc-mode.\\
none & \texttt{M-x cython-\codebreak{}mode} & \texttt{cython-\codebreak{}mode} & Switch the current buffer to cython-mode.\\
none & \texttt{M-x d-mode} & \texttt{d-mode} & Switch the current buffer to d-mode.\\
none & \texttt{M-x dart-\codebreak{}mode} & \texttt{dart-\codebreak{}mode} & Switch the current buffer to dart-mode.\\
\bottomrule
\end{xltabular}

\begin{xltabular}{\textwidth}{@{}l >{\raggedright\arraybackslash\hsize=0.875\hsize}X l >{\raggedright\arraybackslash\hsize=1.125\hsize}X@{}}
\toprule
\textbf{Modal} & \textbf{Non-modal} & \textbf{Command} & \textbf{What it does}\\
\midrule
\endhead
none & \texttt{M-x dbml-\codebreak{}mode} & \texttt{dbml-\codebreak{}mode} & Switch the current buffer to dbml-mode.\\
none & \texttt{M-x debian-\codebreak{}control-\codebreak{}mode} & \texttt{debian-\codebreak{}control-\codebreak{}mode} & Switch the current buffer to debian-control-mode.\\
none & \texttt{M-x dhall-\codebreak{}mode} & \texttt{dhall-\codebreak{}mode} & Switch the current buffer to dhall-mode.\\
none & \texttt{M-x docker-\codebreak{}compose-\codebreak{}mode} & \texttt{docker-\codebreak{}compose-\codebreak{}mode} & Switch the current buffer to docker-compose-mode.\\
none & \texttt{M-x dotenv-\codebreak{}mode} & \texttt{dotenv-\codebreak{}mode} & Switch the current buffer to dotenv-mode.\\
none & \texttt{M-x doxyfile-\codebreak{}mode} & \texttt{doxyfile-\codebreak{}mode} & Switch the current buffer to doxyfile-mode.\\
none & \texttt{M-x dtd-\codebreak{}mode} & \texttt{dtd-\codebreak{}mode} & Switch the current buffer to dtd-mode.\\
none & \texttt{M-x dts-\codebreak{}mode} & \texttt{dts-\codebreak{}mode} & Switch the current buffer to dts-mode.\\
none & \texttt{M-x dune-\codebreak{}mode} & \texttt{dune-\codebreak{}mode} & Switch the current buffer to dune-mode.\\
none & \texttt{M-x dunstrc-\codebreak{}mode} & \texttt{dunstrc-\codebreak{}mode} & Switch the current buffer to dunstrc-mode.\\
none & \texttt{M-x earthfile-\codebreak{}mode} & \texttt{earthfile-\codebreak{}mode} & Switch the current buffer to earthfile-mode.\\
none & \texttt{M-x ebnf-\codebreak{}mode} & \texttt{ebnf-\codebreak{}mode} & Switch the current buffer to ebnf-mode.\\
none & \texttt{M-x eex-\codebreak{}mode} & \texttt{eex-\codebreak{}mode} & Switch the current buffer to eex-mode.\\
none & \texttt{M-x eiffel-\codebreak{}mode} & \texttt{eiffel-\codebreak{}mode} & Switch the current buffer to eiffel-mode.\\
none & \texttt{M-x ejs-\codebreak{}mode} & \texttt{ejs-\codebreak{}mode} & Switch the current buffer to ejs-mode.\\
none & \texttt{M-x elixir-\codebreak{}mode} & \texttt{elixir-\codebreak{}mode} & Switch the current buffer to elixir-mode.\\
none & \texttt{M-x elm-\codebreak{}mode} & \texttt{elm-\codebreak{}mode} & Switch the current buffer to elm-mode.\\
none & \texttt{M-x elvish-\codebreak{}mode} & \texttt{elvish-\codebreak{}mode} & Switch the current buffer to elvish-mode.\\
none & \texttt{M-x emacs-\codebreak{}lisp-\codebreak{}mode} & \texttt{emacs-\codebreak{}lisp-\codebreak{}mode} & Switch the current buffer to emacs-lisp-mode.\\
none & \texttt{M-x erb-\codebreak{}mode} & \texttt{erb-\codebreak{}mode} & Switch the current buffer to erb-mode.\\
none & \texttt{M-x erlang-\codebreak{}mode} & \texttt{erlang-\codebreak{}mode} & Switch the current buffer to erlang-mode.\\
none & \texttt{M-x esdl-\codebreak{}mode} & \texttt{esdl-\codebreak{}mode} & Switch the current buffer to esdl-mode.\\
none & \texttt{M-x feature-\codebreak{}mode} & \texttt{feature-\codebreak{}mode} & Switch the current buffer to feature-mode.\\
none & \texttt{M-x fennel-\codebreak{}mode} & \texttt{fennel-\codebreak{}mode} & Switch the current buffer to fennel-mode.\\
none & \texttt{M-x fga-\codebreak{}mode} & \texttt{fga-\codebreak{}mode} & Switch the current buffer to fga-mode.\\
none & \texttt{M-x fidl-\codebreak{}mode} & \texttt{fidl-\codebreak{}mode} & Switch the current buffer to fidl-mode.\\
none & \texttt{M-x fish-\codebreak{}mode} & \texttt{fish-\codebreak{}mode} & Switch the current buffer to fish-mode.\\
none & \texttt{M-x flatbuffers-\codebreak{}mode} & \texttt{flatbuffers-\codebreak{}mode} & Switch the current buffer to flatbuffers-mode.\\
none & \texttt{M-x forth-\codebreak{}mode} & \texttt{forth-\codebreak{}mode} & Switch the current buffer to forth-mode.\\
none & \texttt{M-x fortran-\codebreak{}mode} & \texttt{fortran-\codebreak{}mode} & Switch the current buffer to fortran-mode.\\
none & \texttt{M-x freebasic-\codebreak{}mode} & \texttt{freebasic-\codebreak{}mode} & Switch the current buffer to freebasic-mode.\\
none & \texttt{M-x fsharp-\codebreak{}mode} & \texttt{fsharp-\codebreak{}mode} & Switch the current buffer to fsharp-mode.\\
none & \texttt{M-x gdscript-\codebreak{}mode} & \texttt{gdscript-\codebreak{}mode} & Switch the current buffer to gdscript-mode.\\
none & \texttt{M-x gitattributes-\codebreak{}mode} & \texttt{gitattributes-\codebreak{}mode} & Switch the current buffer to gitattributes-mode.\\
none & \texttt{M-x gitconfig-\codebreak{}mode} & \texttt{gitconfig-\codebreak{}mode} & Switch the current buffer to gitconfig-mode.\\
none & \texttt{M-x github-\codebreak{}action-\codebreak{}mode} & \texttt{github-\codebreak{}action-\codebreak{}mode} & Switch the current buffer to github-action-mode.\\
none & \texttt{M-x gitignore-\codebreak{}mode} & \texttt{gitignore-\codebreak{}mode} & Switch the current buffer to gitignore-mode.\\
none & \texttt{M-x gitlab-\codebreak{}ci-\codebreak{}mode} & \texttt{gitlab-\codebreak{}ci-\codebreak{}mode} & Switch the current buffer to gitlab-ci-mode.\\
none & \texttt{M-x gleam-\codebreak{}mode} & \texttt{gleam-\codebreak{}mode} & Switch the current buffer to gleam-mode.\\
none & \texttt{M-x glimmer-\codebreak{}js-\codebreak{}mode} & \texttt{glimmer-\codebreak{}js-\codebreak{}mode} & Switch the current buffer to glimmer-js-mode.\\
\bottomrule
\end{xltabular}

\begin{xltabular}{\textwidth}{@{}l >{\raggedright\arraybackslash\hsize=0.583\hsize}X l >{\raggedright\arraybackslash\hsize=1.417\hsize}X@{}}
\toprule
\textbf{Modal} & \textbf{Non-modal} & \textbf{Command} & \textbf{What it does}\\
\midrule
\endhead
none & \texttt{M-x glimmer-\codebreak{}mode} & \texttt{glimmer-\codebreak{}mode} & Switch the current buffer to glimmer-mode.\\
none & \texttt{M-x glimmer-\codebreak{}ts-\codebreak{}mode} & \texttt{glimmer-\codebreak{}ts-\codebreak{}mode} & Switch the current buffer to glimmer-ts-mode.\\
none & \texttt{M-x glsl-\codebreak{}mode} & \texttt{glsl-\codebreak{}mode} & Switch the current buffer to glsl-mode.\\
none & \texttt{M-x gn-\codebreak{}mode} & \texttt{gn-\codebreak{}mode} & Switch the current buffer to gn-mode.\\
none & \texttt{M-x gnuplot-\codebreak{}mode} & \texttt{gnuplot-\codebreak{}mode} & Switch the current buffer to gnuplot-mode.\\
none & \texttt{M-x go-\codebreak{}mod-\codebreak{}mode} & \texttt{go-\codebreak{}mod-\codebreak{}mode} & Switch the current buffer to go-mod-mode.\\
none & \texttt{M-x go-\codebreak{}mode} & \texttt{go-\codebreak{}mode} & Switch the current buffer to go-mode.\\
none & \texttt{M-x go-\codebreak{}template-\codebreak{}mode} & \texttt{go-\codebreak{}template-\codebreak{}mode} & Switch the current buffer to go-template-mode.\\
none & \texttt{M-x go-\codebreak{}work-\codebreak{}mode} & \texttt{go-\codebreak{}work-\codebreak{}mode} & Switch the current buffer to go-work-mode.\\
none & \texttt{M-x godot-\codebreak{}resource-\codebreak{}mode} & \texttt{godot-\codebreak{}resource-\codebreak{}mode} & Switch the current buffer to godot-resource-mode.\\
none & \texttt{M-x gpr-\codebreak{}mode} & \texttt{gpr-\codebreak{}mode} & Switch the current buffer to gpr-mode.\\
none & \texttt{M-x graphql-\codebreak{}mode} & \texttt{graphql-\codebreak{}mode} & Switch the current buffer to graphql-mode.\\
none & \texttt{M-x graphviz-\codebreak{}dot-\codebreak{}mode} & \texttt{graphviz-\codebreak{}dot-\codebreak{}mode} & Switch the current buffer to graphviz-dot-mode.\\
none & \texttt{M-x gren-\codebreak{}mode} & \texttt{gren-\codebreak{}mode} & Switch the current buffer to gren-mode.\\
none & \texttt{M-x groovy-\codebreak{}mode} & \texttt{groovy-\codebreak{}mode} & Switch the current buffer to groovy-mode.\\
none & \texttt{M-x hare-\codebreak{}mode} & \texttt{hare-\codebreak{}mode} & Switch the current buffer to hare-mode.\\
none & \texttt{M-x haskell-\codebreak{}mode} & \texttt{haskell-\codebreak{}mode} & Switch the current buffer to haskell-mode.\\
none & \texttt{M-x haxe-\codebreak{}mode} & \texttt{haxe-\codebreak{}mode} & Switch the current buffer to haxe-mode.\\
none & \texttt{M-x hcl-\codebreak{}mode} & \texttt{hcl-\codebreak{}mode} & Switch the current buffer to hcl-mode.\\
none & \texttt{M-x hdl-\codebreak{}mode} & \texttt{hdl-\codebreak{}mode} & Switch the current buffer to hdl-mode.\\
none & \texttt{M-x heex-\codebreak{}mode} & \texttt{heex-\codebreak{}mode} & Switch the current buffer to heex-mode.\\
none & \texttt{M-x hocon-\codebreak{}mode} & \texttt{hocon-\codebreak{}mode} & Switch the current buffer to hocon-mode.\\
none & \texttt{M-x hoon-\codebreak{}mode} & \texttt{hoon-\codebreak{}mode} & Switch the current buffer to hoon-mode.\\
none & \texttt{M-x html-\codebreak{}mode} & \texttt{html-\codebreak{}mode} & Switch the current buffer to html-mode.\\
none & \texttt{M-x hurl-\codebreak{}mode} & \texttt{hurl-\codebreak{}mode} & Switch the current buffer to hurl-mode.\\
none & \texttt{M-x hy-\codebreak{}mode} & \texttt{hy-\codebreak{}mode} & Switch the current buffer to hy-mode.\\
none & \texttt{M-x hyprlang-\codebreak{}mode} & \texttt{hyprlang-\codebreak{}mode} & Switch the current buffer to hyprlang-mode.\\
none & \texttt{M-x ini-\codebreak{}mode} & \texttt{ini-\codebreak{}mode} & Switch the current buffer to ini-mode.\\
none & \texttt{M-x ink-\codebreak{}mode} & \texttt{ink-\codebreak{}mode} & Switch the current buffer to ink-mode.\\
none & \texttt{M-x inko-\codebreak{}mode} & \texttt{inko-\codebreak{}mode} & Switch the current buffer to inko-mode.\\
none & \texttt{M-x janet-\codebreak{}mode} & \texttt{janet-\codebreak{}mode} & Switch the current buffer to janet-mode.\\
none & \texttt{M-x java-\codebreak{}mode} & \texttt{java-\codebreak{}mode} & Switch the current buffer to java-mode.\\
none & \texttt{M-x jinja2-\codebreak{}mode} & \texttt{jinja2-\codebreak{}mode} & Switch the current buffer to jinja2-mode.\\
none & \texttt{M-x jq-\codebreak{}mode} & \texttt{jq-\codebreak{}mode} & Switch the current buffer to jq-mode.\\
none & \texttt{M-x js-\codebreak{}jsx-\codebreak{}mode} & \texttt{js-\codebreak{}jsx-\codebreak{}mode} & Switch the current buffer to js-jsx-mode.\\
none & \texttt{M-x js-\codebreak{}mode} & \texttt{js-\codebreak{}mode} & Switch the current buffer to js-mode.\\
none & \texttt{M-x json-\codebreak{}mode} & \texttt{json-\codebreak{}mode} & Switch the current buffer to json-mode.\\
none & \texttt{M-x json5-\codebreak{}mode} & \texttt{json5-\codebreak{}mode} & Switch the current buffer to json5-mode.\\
none & \texttt{M-x jsonc-\codebreak{}mode} & \texttt{jsonc-\codebreak{}mode} & Switch the current buffer to jsonc-mode.\\
none & \texttt{M-x jsonl-\codebreak{}mode} & \texttt{jsonl-\codebreak{}mode} & Switch the current buffer to jsonl-mode.\\
\bottomrule
\end{xltabular}

\begin{xltabular}{\textwidth}{@{}l >{\raggedright\arraybackslash\hsize=0.783\hsize}X l >{\raggedright\arraybackslash\hsize=1.217\hsize}X@{}}
\toprule
\textbf{Modal} & \textbf{Non-modal} & \textbf{Command} & \textbf{What it does}\\
\midrule
\endhead
none & \texttt{M-x jsonnet-\codebreak{}mode} & \texttt{jsonnet-\codebreak{}mode} & Switch the current buffer to jsonnet-mode.\\
none & \texttt{M-x julia-\codebreak{}mode} & \texttt{julia-\codebreak{}mode} & Switch the current buffer to julia-mode.\\
none & \texttt{M-x just-\codebreak{}mode} & \texttt{just-\codebreak{}mode} & Switch the current buffer to just-mode.\\
none & \texttt{M-x kcl-\codebreak{}mode} & \texttt{kcl-\codebreak{}mode} & Switch the current buffer to kcl-mode.\\
none & \texttt{M-x kconfig-\codebreak{}mode} & \texttt{kconfig-\codebreak{}mode} & Switch the current buffer to kconfig-mode.\\
none & \texttt{M-x kdl-\codebreak{}mode} & \texttt{kdl-\codebreak{}mode} & Switch the current buffer to kdl-mode.\\
none & \texttt{M-x koka-\codebreak{}mode} & \texttt{koka-\codebreak{}mode} & Switch the current buffer to koka-mode.\\
none & \texttt{M-x kotlin-\codebreak{}mode} & \texttt{kotlin-\codebreak{}mode} & Switch the current buffer to kotlin-mode.\\
none & \texttt{M-x koto-\codebreak{}mode} & \texttt{koto-\codebreak{}mode} & Switch the current buffer to koto-mode.\\
none & \texttt{M-x ld-\codebreak{}script-\codebreak{}mode} & \texttt{ld-\codebreak{}script-\codebreak{}mode} & Switch the current buffer to ld-script-mode.\\
none & \texttt{M-x lean4-\codebreak{}mode} & \texttt{lean4-\codebreak{}mode} & Switch the current buffer to lean4-mode.\\
none & \texttt{M-x less-\codebreak{}css-\codebreak{}mode} & \texttt{less-\codebreak{}css-\codebreak{}mode} & Switch the current buffer to less-css-mode.\\
none & \texttt{M-x lisp-\codebreak{}interaction-\codebreak{}mode} & \texttt{lisp-\codebreak{}interaction-\codebreak{}mode} & Switch the current buffer to lisp-interaction-mode.\\
none & \texttt{M-x lisp-\codebreak{}mode} & \texttt{lisp-\codebreak{}mode} & Switch the current buffer to lisp-mode.\\
none & \texttt{M-x literate-\codebreak{}haskell-\codebreak{}mode} & \texttt{literate-\codebreak{}haskell-\codebreak{}mode} & Switch the current buffer to literate-haskell-mode.\\
none & \texttt{M-x llvm-\codebreak{}mode} & \texttt{llvm-\codebreak{}mode} & Switch the current buffer to llvm-mode.\\
none & \texttt{M-x lpf-\codebreak{}mode} & \texttt{lpf-\codebreak{}mode} & Switch the current buffer to lpf-mode.\\
none & \texttt{M-x lua-\codebreak{}mode} & \texttt{lua-\codebreak{}mode} & Switch the current buffer to lua-mode.\\
none & \texttt{M-x luau-\codebreak{}mode} & \texttt{luau-\codebreak{}mode} & Switch the current buffer to luau-mode.\\
none & \texttt{M-x makefile-\codebreak{}mode} & \texttt{makefile-\codebreak{}mode} & Switch the current buffer to makefile-mode.\\
none & \texttt{M-x matlab-\codebreak{}mode} & \texttt{matlab-\codebreak{}mode} & Switch the current buffer to matlab-mode.\\
none & \texttt{M-x mermaid-\codebreak{}mode} & \texttt{mermaid-\codebreak{}mode} & Switch the current buffer to mermaid-mode.\\
none & \texttt{M-x meson-\codebreak{}mode} & \texttt{meson-\codebreak{}mode} & Switch the current buffer to meson-mode.\\
none & \texttt{M-x mojo-\codebreak{}mode} & \texttt{mojo-\codebreak{}mode} & Switch the current buffer to mojo-mode.\\
none & \texttt{M-x move-\codebreak{}mode} & \texttt{move-\codebreak{}mode} & Switch the current buffer to move-mode.\\
none & \texttt{M-x nasm-\codebreak{}mode} & \texttt{nasm-\codebreak{}mode} & Switch the current buffer to nasm-mode.\\
none & \texttt{M-x nearley-\codebreak{}mode} & \texttt{nearley-\codebreak{}mode} & Switch the current buffer to nearley-mode.\\
none & \texttt{M-x nestedtext-\codebreak{}mode} & \texttt{nestedtext-\codebreak{}mode} & Switch the current buffer to nestedtext-mode.\\
none & \texttt{M-x nginx-\codebreak{}mode} & \texttt{nginx-\codebreak{}mode} & Switch the current buffer to nginx-mode.\\
none & \texttt{M-x nickel-\codebreak{}mode} & \texttt{nickel-\codebreak{}mode} & Switch the current buffer to nickel-mode.\\
none & \texttt{M-x nim-\codebreak{}mode} & \texttt{nim-\codebreak{}mode} & Switch the current buffer to nim-mode.\\
none & \texttt{M-x nix-\codebreak{}mode} & \texttt{nix-\codebreak{}mode} & Switch the current buffer to nix-mode.\\
none & \texttt{M-x nushell-\codebreak{}mode} & \texttt{nushell-\codebreak{}mode} & Switch the current buffer to nushell-mode.\\
none & \texttt{M-x ocaml-\codebreak{}interface-\codebreak{}mode} & \texttt{ocaml-\codebreak{}interface-\codebreak{}mode} & Switch the current buffer to ocaml-interface-mode.\\
none & \texttt{M-x ocaml-\codebreak{}mode} & \texttt{ocaml-\codebreak{}mode} & Switch the current buffer to ocaml-mode.\\
none & \texttt{M-x odin-\codebreak{}mode} & \texttt{odin-\codebreak{}mode} & Switch the current buffer to odin-mode.\\
none & \texttt{M-x ohm-\codebreak{}mode} & \texttt{ohm-\codebreak{}mode} & Switch the current buffer to ohm-mode.\\
none & \texttt{M-x pascal-\codebreak{}mode} & \texttt{pascal-\codebreak{}mode} & Switch the current buffer to pascal-mode.\\
none & \texttt{M-x penrose-\codebreak{}mode} & \texttt{penrose-\codebreak{}mode} & Switch the current buffer to penrose-mode.\\
none & \texttt{M-x perl-\codebreak{}mode} & \texttt{perl-\codebreak{}mode} & Switch the current buffer to perl-mode.\\
\bottomrule
\end{xltabular}

\begin{xltabular}{\textwidth}{@{}l >{\raggedright\arraybackslash\hsize=0.870\hsize}X l >{\raggedright\arraybackslash\hsize=1.130\hsize}X@{}}
\toprule
\textbf{Modal} & \textbf{Non-modal} & \textbf{Command} & \textbf{What it does}\\
\midrule
\endhead
none & \texttt{M-x pest-\codebreak{}mode} & \texttt{pest-\codebreak{}mode} & Switch the current buffer to pest-mode.\\
none & \texttt{M-x php-\codebreak{}mode} & \texttt{php-\codebreak{}mode} & Switch the current buffer to php-mode.\\
none & \texttt{M-x picat-\codebreak{}mode} & \texttt{picat-\codebreak{}mode} & Switch the current buffer to picat-mode.\\
none & \texttt{M-x pip-\codebreak{}requirements-\codebreak{}mode} & \texttt{pip-\codebreak{}requirements-\codebreak{}mode} & Switch the current buffer to pip-requirements-mode.\\
none & \texttt{M-x pkl-\codebreak{}mode} & \texttt{pkl-\codebreak{}mode} & Switch the current buffer to pkl-mode.\\
none & \texttt{M-x ponylang-\codebreak{}mode} & \texttt{ponylang-\codebreak{}mode} & Switch the current buffer to ponylang-mode.\\
none & \texttt{M-x powershell-\codebreak{}mode} & \texttt{powershell-\codebreak{}mode} & Switch the current buffer to powershell-mode.\\
none & \texttt{M-x prisma-\codebreak{}mode} & \texttt{prisma-\codebreak{}mode} & Switch the current buffer to prisma-mode.\\
none & \texttt{M-x prolog-\codebreak{}mode} & \texttt{prolog-\codebreak{}mode} & Switch the current buffer to prolog-mode.\\
none & \texttt{M-x protobuf-\codebreak{}mode} & \texttt{protobuf-\codebreak{}mode} & Switch the current buffer to protobuf-mode.\\
none & \texttt{M-x proverif-\codebreak{}mode} & \texttt{proverif-\codebreak{}mode} & Switch the current buffer to proverif-mode.\\
none & \texttt{M-x prql-\codebreak{}mode} & \texttt{prql-\codebreak{}mode} & Switch the current buffer to prql-mode.\\
none & \texttt{M-x ptx-\codebreak{}mode} & \texttt{ptx-\codebreak{}mode} & Switch the current buffer to ptx-mode.\\
none & \texttt{M-x pug-\codebreak{}mode} & \texttt{pug-\codebreak{}mode} & Switch the current buffer to pug-mode.\\
none & \texttt{M-x purescript-\codebreak{}mode} & \texttt{purescript-\codebreak{}mode} & Switch the current buffer to purescript-mode.\\
none & \texttt{M-x python-\codebreak{}mode} & \texttt{python-\codebreak{}mode} & Switch the current buffer to python-mode.\\
none & \texttt{M-x ql-\codebreak{}mode} & \texttt{ql-\codebreak{}mode} & Switch the current buffer to ql-mode.\\
none & \texttt{M-x qml-\codebreak{}mode} & \texttt{qml-\codebreak{}mode} & Switch the current buffer to qml-mode.\\
none & \texttt{M-x quint-\codebreak{}mode} & \texttt{quint-\codebreak{}mode} & Switch the current buffer to quint-mode.\\
none & \texttt{M-x r-mode} & \texttt{r-mode} & Switch the current buffer to r-mode.\\
none & \texttt{M-x racket-\codebreak{}mode} & \texttt{racket-\codebreak{}mode} & Switch the current buffer to racket-mode.\\
none & \texttt{M-x rego-\codebreak{}mode} & \texttt{rego-\codebreak{}mode} & Switch the current buffer to rego-mode.\\
none & \texttt{M-x rescript-\codebreak{}mode} & \texttt{rescript-\codebreak{}mode} & Switch the current buffer to rescript-mode.\\
none & \texttt{M-x ripple-\codebreak{}mode} & \texttt{ripple-\codebreak{}mode} & Switch the current buffer to ripple-mode.\\
none & \texttt{M-x robot-\codebreak{}mode} & \texttt{robot-\codebreak{}mode} & Switch the current buffer to robot-mode.\\
none & \texttt{M-x ron-\codebreak{}mode} & \texttt{ron-\codebreak{}mode} & Switch the current buffer to ron-mode.\\
none & \texttt{M-x rpm-\codebreak{}spec-\codebreak{}mode} & \texttt{rpm-\codebreak{}spec-\codebreak{}mode} & Switch the current buffer to rpm-spec-mode.\\
none & \texttt{M-x ruby-\codebreak{}mode} & \texttt{ruby-\codebreak{}mode} & Switch the current buffer to ruby-mode.\\
none & \texttt{M-x rust-\codebreak{}mode} & \texttt{rust-\codebreak{}mode} & Switch the current buffer to rust-mode.\\
none & \texttt{M-x sage-\codebreak{}mode} & \texttt{sage-\codebreak{}mode} & Switch the current buffer to sage-mode.\\
none & \texttt{M-x scad-\codebreak{}mode} & \texttt{scad-\codebreak{}mode} & Switch the current buffer to scad-mode.\\
none & \texttt{M-x scala-\codebreak{}mode} & \texttt{scala-\codebreak{}mode} & Switch the current buffer to scala-mode.\\
none & \texttt{M-x scfg-\codebreak{}mode} & \texttt{scfg-\codebreak{}mode} & Switch the current buffer to scfg-mode.\\
none & \texttt{M-x scheme-\codebreak{}mode} & \texttt{scheme-\codebreak{}mode} & Switch the current buffer to scheme-mode.\\
none & \texttt{M-x sclang-\codebreak{}mode} & \texttt{sclang-\codebreak{}mode} & Switch the current buffer to sclang-mode.\\
none & \texttt{M-x scss-\codebreak{}mode} & \texttt{scss-\codebreak{}mode} & Switch the current buffer to scss-mode.\\
none & \texttt{M-x shellcheckrc-\codebreak{}mode} & \texttt{shellcheckrc-\codebreak{}mode} & Switch the current buffer to shellcheckrc-mode.\\
none & \texttt{M-x slang-\codebreak{}mode} & \texttt{slang-\codebreak{}mode} & Switch the current buffer to slang-mode.\\
none & \texttt{M-x slint-\codebreak{}mode} & \texttt{slint-\codebreak{}mode} & Switch the current buffer to slint-mode.\\
none & \texttt{M-x slisp-\codebreak{}mode} & \texttt{slisp-\codebreak{}mode} & Switch the current buffer to slisp-mode.\\
\bottomrule
\end{xltabular}

\begin{xltabular}{\textwidth}{@{}l >{\raggedright\arraybackslash\hsize=0.857\hsize}X l >{\raggedright\arraybackslash\hsize=1.143\hsize}X@{}}
\toprule
\textbf{Modal} & \textbf{Non-modal} & \textbf{Command} & \textbf{What it does}\\
\midrule
\endhead
none & \texttt{M-x smali-\codebreak{}mode} & \texttt{smali-\codebreak{}mode} & Switch the current buffer to smali-mode.\\
none & \texttt{M-x smithy-\codebreak{}mode} & \texttt{smithy-\codebreak{}mode} & Switch the current buffer to smithy-mode.\\
none & \texttt{M-x sml-\codebreak{}mode} & \texttt{sml-\codebreak{}mode} & Switch the current buffer to sml-mode.\\
none & \texttt{M-x snakemake-\codebreak{}mode} & \texttt{snakemake-\codebreak{}mode} & Switch the current buffer to snakemake-mode.\\
none & \texttt{M-x solidity-\codebreak{}mode} & \texttt{solidity-\codebreak{}mode} & Switch the current buffer to solidity-mode.\\
none & \texttt{M-x sourcepawn-\codebreak{}mode} & \texttt{sourcepawn-\codebreak{}mode} & Switch the current buffer to sourcepawn-mode.\\
none & \texttt{M-x spade-\codebreak{}mode} & \texttt{spade-\codebreak{}mode} & Switch the current buffer to spade-mode.\\
none & \texttt{M-x spicedb-\codebreak{}mode} & \texttt{spicedb-\codebreak{}mode} & Switch the current buffer to spicedb-mode.\\
none & \texttt{M-x sql-\codebreak{}mode} & \texttt{sql-\codebreak{}mode} & Switch the current buffer to sql-mode.\\
none & \texttt{M-x ssh-\codebreak{}config-\codebreak{}mode} & \texttt{ssh-\codebreak{}config-\codebreak{}mode} & Switch the current buffer to ssh-config-mode.\\
none & \texttt{M-x strictdoc-\codebreak{}mode} & \texttt{strictdoc-\codebreak{}mode} & Switch the current buffer to strictdoc-mode.\\
none & \texttt{M-x styx-\codebreak{}mode} & \texttt{styx-\codebreak{}mode} & Switch the current buffer to styx-mode.\\
none & \texttt{M-x svelte-\codebreak{}mode} & \texttt{svelte-\codebreak{}mode} & Switch the current buffer to svelte-mode.\\
none & \texttt{M-x sway-\codebreak{}mode} & \texttt{sway-\codebreak{}mode} & Switch the current buffer to sway-mode.\\
none & \texttt{M-x swift-\codebreak{}mode} & \texttt{swift-\codebreak{}mode} & Switch the current buffer to swift-mode.\\
none & \texttt{M-x systemd-\codebreak{}mode} & \texttt{systemd-\codebreak{}mode} & Switch the current buffer to systemd-mode.\\
none & \texttt{M-x systemverilog-\codebreak{}mode} & \texttt{systemverilog-\codebreak{}mode} & Switch the current buffer to systemverilog-mode.\\
none & \texttt{M-x t32-\codebreak{}mode} & \texttt{t32-\codebreak{}mode} & Switch the current buffer to t32-mode.\\
none & \texttt{M-x tablegen-\codebreak{}mode} & \texttt{tablegen-\codebreak{}mode} & Switch the current buffer to tablegen-mode.\\
none & \texttt{M-x tact-\codebreak{}mode} & \texttt{tact-\codebreak{}mode} & Switch the current buffer to tact-mode.\\
none & \texttt{M-x tcl-\codebreak{}mode} & \texttt{tcl-\codebreak{}mode} & Switch the current buffer to tcl-mode.\\
none & \texttt{M-x teal-\codebreak{}mode} & \texttt{teal-\codebreak{}mode} & Switch the current buffer to teal-mode.\\
none & \texttt{M-x templ-\codebreak{}mode} & \texttt{templ-\codebreak{}mode} & Switch the current buffer to templ-mode.\\
none & \texttt{M-x tera-\codebreak{}mode} & \texttt{tera-\codebreak{}mode} & Switch the current buffer to tera-mode.\\
none & \texttt{M-x textproto-\codebreak{}mode} & \texttt{textproto-\codebreak{}mode} & Switch the current buffer to textproto-mode.\\
none & \texttt{M-x thrift-\codebreak{}mode} & \texttt{thrift-\codebreak{}mode} & Switch the current buffer to thrift-mode.\\
none & \texttt{M-x tla-\codebreak{}mode} & \texttt{tla-\codebreak{}mode} & Switch the current buffer to tla-mode.\\
none & \texttt{M-x tolk-\codebreak{}mode} & \texttt{tolk-\codebreak{}mode} & Switch the current buffer to tolk-mode.\\
none & \texttt{M-x toml-\codebreak{}mode} & \texttt{toml-\codebreak{}mode} & Switch the current buffer to toml-mode.\\
none & \texttt{M-x tql-\codebreak{}mode} & \texttt{tql-\codebreak{}mode} & Switch the current buffer to tql-mode.\\
none & \texttt{M-x tsx-\codebreak{}mode} & \texttt{tsx-\codebreak{}mode} & Switch the current buffer to tsx-mode.\\
none & \texttt{M-x twig-\codebreak{}mode} & \texttt{twig-\codebreak{}mode} & Switch the current buffer to twig-mode.\\
none & \texttt{M-x typescript-\codebreak{}mode} & \texttt{typescript-\codebreak{}mode} & Switch the current buffer to typescript-mode.\\
none & \texttt{M-x typespec-\codebreak{}mode} & \texttt{typespec-\codebreak{}mode} & Switch the current buffer to typespec-mode.\\
none & \texttt{M-x typst-\codebreak{}mode} & \texttt{typst-\codebreak{}mode} & Switch the current buffer to typst-mode.\\
none & \texttt{M-x ungrammar-\codebreak{}mode} & \texttt{ungrammar-\codebreak{}mode} & Switch the current buffer to ungrammar-mode.\\
none & \texttt{M-x unison-\codebreak{}mode} & \texttt{unison-\codebreak{}mode} & Switch the current buffer to unison-mode.\\
none & \texttt{M-x uxntal-\codebreak{}mode} & \texttt{uxntal-\codebreak{}mode} & Switch the current buffer to uxntal-mode.\\
none & \texttt{M-x v-mode} & \texttt{v-mode} & Switch the current buffer to v-mode.\\
none & \texttt{M-x vala-\codebreak{}mode} & \texttt{vala-\codebreak{}mode} & Switch the current buffer to vala-mode.\\
\bottomrule
\end{xltabular}

\begin{xltabular}{\textwidth}{@{}l >{\raggedright\arraybackslash\hsize=0.694\hsize}X l >{\raggedright\arraybackslash\hsize=1.306\hsize}X@{}}
\toprule
\textbf{Modal} & \textbf{Non-modal} & \textbf{Command} & \textbf{What it does}\\
\midrule
\endhead
none & \texttt{M-x vento-\codebreak{}mode} & \texttt{vento-\codebreak{}mode} & Switch the current buffer to vento-mode.\\
none & \texttt{M-x verilog-\codebreak{}mode} & \texttt{verilog-\codebreak{}mode} & Switch the current buffer to verilog-mode.\\
none & \texttt{M-x vhdl-\codebreak{}mode} & \texttt{vhdl-\codebreak{}mode} & Switch the current buffer to vhdl-mode.\\
none & \texttt{M-x vhs-\codebreak{}mode} & \texttt{vhs-\codebreak{}mode} & Switch the current buffer to vhs-mode.\\
none & \texttt{M-x vimrc-\codebreak{}mode} & \texttt{vimrc-\codebreak{}mode} & Switch the current buffer to vimrc-mode.\\
none & \texttt{M-x vue-\codebreak{}mode} & \texttt{vue-\codebreak{}mode} & Switch the current buffer to vue-mode.\\
none & \texttt{M-x wast-\codebreak{}mode} & \texttt{wast-\codebreak{}mode} & Switch the current buffer to wast-mode.\\
none & \texttt{M-x wat-\codebreak{}mode} & \texttt{wat-\codebreak{}mode} & Switch the current buffer to wat-mode.\\
none & \texttt{M-x werk-\codebreak{}mode} & \texttt{werk-\codebreak{}mode} & Switch the current buffer to werk-mode.\\
none & \texttt{M-x wesl-\codebreak{}mode} & \texttt{wesl-\codebreak{}mode} & Switch the current buffer to wesl-mode.\\
none & \texttt{M-x wgsl-\codebreak{}mode} & \texttt{wgsl-\codebreak{}mode} & Switch the current buffer to wgsl-mode.\\
none & \texttt{M-x wit-\codebreak{}mode} & \texttt{wit-\codebreak{}mode} & Switch the current buffer to wit-mode.\\
none & \texttt{M-x woodpecker-\codebreak{}ci-\codebreak{}mode} & \texttt{woodpecker-\codebreak{}ci-\codebreak{}mode} & Switch the current buffer to woodpecker-ci-mode.\\
none & \texttt{M-x wren-\codebreak{}mode} & \texttt{wren-\codebreak{}mode} & Switch the current buffer to wren-mode.\\
none & \texttt{M-x xml-\codebreak{}mode} & \texttt{xml-\codebreak{}mode} & Switch the current buffer to xml-mode.\\
none & \texttt{M-x xtc-\codebreak{}mode} & \texttt{xtc-\codebreak{}mode} & Switch the current buffer to xtc-mode.\\
none & \texttt{M-x yaml-\codebreak{}mode} & \texttt{yaml-\codebreak{}mode} & Switch the current buffer to yaml-mode.\\
none & \texttt{M-x yara-\codebreak{}mode} & \texttt{yara-\codebreak{}mode} & Switch the current buffer to yara-mode.\\
none & \texttt{M-x yuck-\codebreak{}mode} & \texttt{yuck-\codebreak{}mode} & Switch the current buffer to yuck-mode.\\
none & \texttt{M-x zig-\codebreak{}mode} & \texttt{zig-\codebreak{}mode} & Switch the current buffer to zig-mode.\\
\bottomrule
\end{xltabular}

\section{Org}
\label{sec:393}
Notes, lists, tasks, dates, tables and documents, all in plain text.

Org is a way of writing plain text so that a program can help you with it. A
line that begins with a star is a heading. A line that begins with a hyphen is
an item of a list. A heading with the word \texttt{TODO} in front is a task. From a
handful of marks like these, Tuatara can fold a long note down to its
headings, tick off a shopping list, show you what is due this week, add up a
column of numbers and turn the whole note into a PDF or a Word document.

The file stays plain text throughout. You can read it in any program, send it
to anyone and keep it for decades. Org began in the editor Emacs, and files
written there open in Tuatara as they are. Tuatara’s Org is its own, written
for Tuatara, so nothing in this chapter needs Emacs.

This is a long chapter because Org does a lot. You do not need all of it. The
first three sections are enough to keep notes and lists. Read about tasks,
dates and the agenda when you want Tuatara to remind you what to do. Read
about tables when you have numbers, and about exporting when a note must
become a document for someone else.

One reminder about the state keys. \texttt{i} starts INSERT state, and \texttt{C-g} or \texttt{ESC}
returns to NORMAL state. Keys that start with \texttt{C-c}, \texttt{C-x} or \texttt{M-} work in
both. Where this chapter says “type”, do it in INSERT state.

\subsection{An Org file}
\label{sec:344}
A file whose name ends in \texttt{.org} opens in Org mode, and the mode line says
Org. Any other buffer becomes an Org buffer with \texttt{M-x org-\codebreak{}mode}. The scratch
buffer \texttt{*org-\codebreak{}scratch*} of chapter 7 is one already, and is a good place to try
things.

Here is a small note as you type it:

\begin{Verbatim}[breaklines,breakautoindent=false,breakindent=1em,breakafter=/-.\,;:=,commandchars=\\\{\}]
#+title: My garden

* Vegetables
The beds by the fence get the most sun.
** Tomatoes
Water *every* morning, /not/ at noon.
** Beans
* Fruit
\end{Verbatim}

The first line gives the note a title. Lines that begin with \texttt{\#+} tell Org
something about the note and are not part of its text. One star makes a
heading of the first level and two stars make a heading under it. There must
be a space after the stars. The text below a heading belongs to it, down to
the next heading.

Within the text a few marks change how words look, and how they come out in
an exported document.

\begin{longtable}{@{}l l@{}}
\toprule
\textbf{You type} & \textbf{It means}\\
\midrule
\endhead
\texttt{*bold*} & Bold\\
\texttt{/italic/} & Italic\\
\texttt{\_underlined\_} & Underlined\\
\texttt{\textasciitilde{}code\textasciitilde{}} & Code, such as a command’s name\\
\texttt{=verbatim=} & Text shown exactly as written\\
\texttt{[fn:1]} & A footnote, its text on a line beginning with the same mark\\
\bottomrule
\end{longtable}

\subsubsection{How a note is shown}
\label{sec:343}
Tuatara shows an Org file a little tidier than it is written. A heading shows
a round sign in place of its last star and the stars before it are hidden.
Each line is set in by the level of the heading it is under. None of this
changes the file, which still holds the stars at the left edge. This is the
indented view. \texttt{M-x org-\codebreak{}indent-\codebreak{}mode} turns it off for the buffer, so that you
see the file as it is, and on again.

If you would rather see the stars always, set \texttt{*org-\codebreak{}startup-\codebreak{}indented*} to
\texttt{nil} in your init file. A single note can say so for itself with the line
\texttt{\#+STARTUP: noindent}, and \texttt{\#+STARTUP: indent} asks for the indented view.
\texttt{*org-\codebreak{}bullets*} holds the signs that stand for the stars, one for each level,
and \texttt{nil} there shows the stars themselves. Lines like \texttt{\#+STARTUP:} are read
when the file is opened. After you change one, kill the buffer with \texttt{C-x k}
and open the file again.

\begin{xltabular}{\textwidth}{@{}l l l >{\raggedright\arraybackslash\hsize=1.000\hsize}X@{}}
\toprule
\textbf{Modal} & \textbf{Non-modal} & \textbf{Command} & \textbf{What it does}\\
\midrule
\endhead
none & \texttt{M-x org-\codebreak{}mode} & \texttt{org-\codebreak{}mode} & Switch the current buffer to org-mode.\\
none & \texttt{M-x org-\codebreak{}indent-\codebreak{}mode} & \texttt{org-\codebreak{}indent-\codebreak{}mode} & Show the current Org buffer indented, or as it is again: each line indented by the level of the headline it is under, a headline’s stars but the last hidden.\\
\bottomrule
\end{xltabular}

\begin{xltabular}{\textwidth}{@{}l l >{\raggedright\arraybackslash\hsize=1.000\hsize}X@{}}
\toprule
\textbf{Setting} & \textbf{Standard} & \textbf{What it changes}\\
\midrule
\endhead
\texttt{*org-\codebreak{}startup-\codebreak{}indented*} & \texttt{t} & True for Org files to open in the indented view (M-x org-indent-mode), as they do: each line shown indented by the level of the headline it is under, a headline’s stars but the last hidden, the file itself as it is; NIL for them to open as they are written.\\
\texttt{*org-\codebreak{}bullets*} & \texttt{("{\tuataraFallbackTWO{}◉}" "○" "{\tuataraFallbackONE{}✸}" "{\tuataraFallbackONE{}✿}")} & The signs Org headlines show in place of their last star, by level, the first for the first level and round again past the last; NIL to show the stars.\\
\texttt{*org-\codebreak{}hide-\codebreak{}leading-\codebreak{}stars*} & \texttt{t} & True to hide a headline’s stars but the last, with \textbf{ORG-BULLETS} or in the indented view; NIL to show them faint (the :org-leading-stars face).\\
\texttt{*org-\codebreak{}depth-\codebreak{}limit*} & \texttt{200} & The most elements and objects of Org text the parser reads one within another: what lies deeper is read as plain text, so that a file made of nothing but nesting cannot exhaust the parser.\\
\bottomrule
\end{xltabular}

\begin{xltabular}{\textwidth}{@{}l >{\raggedright\arraybackslash\hsize=1.000\hsize}X@{}}
\toprule
\textbf{Face} & \textbf{What it paints}\\
\midrule
\endhead
\texttt{:org-\codebreak{}document-\codebreak{}title} & The title an Org document gives itself (\#+title:).\\
\texttt{:org-\codebreak{}document-\codebreak{}info} & What an Org document says of itself: its author, date and the like.\\
\texttt{:org-\codebreak{}document-\codebreak{}info-\codebreak{}keyword} & The keywords before an Org document’s title, author and date (\#+title:).\\
\texttt{:org-\codebreak{}meta-\codebreak{}line} & Org’s keyword lines (\#+name:, \#+caption:\ldots{}).\\
\texttt{:org-\codebreak{}bold} & Bold text in Org (\textbf{bold}).\\
\texttt{:org-\codebreak{}italic} & Italic text in Org (\emph{italic}).\\
\texttt{:org-\codebreak{}underline} & Underlined text in Org (\_underline\_).\\
\texttt{:org-\codebreak{}code} & Code in Org text (\texttt{code}), and fixed-width lines (: \ldots{}).\\
\texttt{:org-\codebreak{}verbatim} & Verbatim text in Org (\texttt{verbatim}).\\
\texttt{:org-\codebreak{}footnote} & Org’s footnotes, \footnote{}.\\
\texttt{:org-\codebreak{}cite} & Org’s citations, [cite:@key].\\
\texttt{:org-\codebreak{}latex-\codebreak{}and-\codebreak{}related} & \LaTeX{} in Org text ($x^2$, \texttt{\textbackslash{}alpha}).\\
\bottomrule
\end{xltabular}

\subsection{Headings and folding}
\label{sec:349}
Headings are what make a long note manageable. You can hide everything under
a heading, move a heading with all that belongs to it, and jump to one by
name. A heading together with everything under it, its text and its own
headings, is called a subtree.

\subsubsection{Folding}
\label{sec:346}
Folding hides text without changing it. Press \texttt{TAB} on a heading and
everything under it is hidden. Three dots after the heading say there is
more. Press \texttt{TAB} again and you see the headings directly under it, each
still folded. A third press shows all of it, and a fourth folds it again.
The echo area names each step. Figure 14.1 shows the three.

\begin{figure}[H]
\caption{Figure 14.1. What TAB shows on the heading Vegetables, step by step.}\label{org345}
\begin{Verbatim}[breaklines,breakautoindent=false,breakindent=1em,breakafter=/-.\,;:=,commandchars=\\\{\}]
 Folded                 Children                    Subtree

 {\tuataraFallbackTWO{}◉} Vegetables...        {\tuataraFallbackTWO{}◉} Vegetables                {\tuataraFallbackTWO{}◉} Vegetables
 {\tuataraFallbackTWO{}◉} Fruit                  The beds by the fence       The beds by the fence
                          get the most sun.           get the most sun.
                          ○ Tomatoes...               ○ Tomatoes
                          ○ Beans                       Water every morning.
                        {\tuataraFallbackTWO{}◉} Fruit                       ○ Beans
                                                    {\tuataraFallbackTWO{}◉} Fruit
\end{Verbatim}
\end{figure}

\texttt{S-TAB} does the same for the whole note at once, wherever the cursor is. The
first press shows only the headings of the first level, which the echo area
calls Overview. The second shows every heading and no text, called Contents.
The third shows everything.

A folded part counts as one line. The cursor steps over it and never stands
inside it.

A note can say how it wants to open. \texttt{\#+STARTUP: overview} near the top opens
it folded to its first-level headings, and \texttt{\#+STARTUP: content} opens it with
every heading and no text. This is worth doing for a long file you open often.

Under a task Org keeps some records of its own, in blocks called drawers that
begin with a line such as \texttt{:LOGBOOK:} and end with \texttt{:END:}. They are folded
to one line when a file opens, because you seldom want to read them. \texttt{TAB} on
a drawer’s first line opens it.

\subsubsection{New headings}
\label{sec:347}
\texttt{M-RET} makes a new heading of the same level, after the subtree the cursor
is in. You can press it anywhere under a heading. It leaves you in NORMAL
state with the cursor after the new stars, so press \texttt{i} and type the heading.
\texttt{M-S-\codebreak{}RET} does the same and writes \texttt{TODO} for you, which starts a task.

\subsubsection{Moving and changing level}
\label{sec:348}
The keys that move the cursor, held with Alt, move headings.

\begin{xltabular}{\textwidth}{@{}l l >{\raggedright\arraybackslash\hsize=1.000\hsize}X@{}}
\toprule
\textbf{Modal} & \textbf{Non-modal} & \textbf{What it does to the heading the cursor is in}\\
\midrule
\endhead
none & \texttt{M-j} / \texttt{M-k} & Moves its subtree down or up, past the next heading of its level\\
none & \texttt{M-h} / \texttt{M-l} & Takes away or adds a star, on this heading alone\\
none & \texttt{M-H} / \texttt{M-L} & Takes away or adds a star on it and on every heading under it\\
\bottomrule
\end{xltabular}

Taking a star away is called promoting, and adding one is called demoting.
The arrow keys with Alt do the same as the letters. Each move is one step
for undo.

To jump to a heading by name, press \texttt{M-g o}, which chapter 13 describes. It
lists every heading of the note, and typing a few letters narrows the list.
To work on one heading without seeing the others, \texttt{C-x x s} narrows the
buffer to its subtree, as chapter 7 explains.

\begin{try}
\textbf{Try it, with notes about a garden}
\begin{enumerate}
\item Press \texttt{C-x C-f}, type \texttt{\textasciitilde{}/garden.\codebreak{}org} and press \texttt{RET}. Press \texttt{i} and type
   the note shown at the start of this chapter, eight lines from \texttt{\#+title:}
   to \texttt{* Fruit}. Press \texttt{C-g} and save with \texttt{C-x C-s}.
\item Move to the line \texttt{Vegetables} and press \texttt{TAB}. Its text and its two
   headings are hidden. Press \texttt{TAB} twice more to see the other two steps.
\item Press \texttt{S-TAB} three times and watch the echo area say Overview, Contents
   and Show all.
\item With the cursor on \texttt{Vegetables}, press \texttt{M-j}. The heading moves below
   \texttt{Fruit} and takes Tomatoes and Beans with it. Press \texttt{M-k} to move it back.
\item Move to \texttt{Beans} and press \texttt{M-l}. It gets a third star and is now under
   Tomatoes. Press \texttt{M-h} to put it back.
\item Press \texttt{M-RET}, then \texttt{i}, and type \texttt{Carrots}. Press \texttt{C-g}.
\end{enumerate}
\end{try}

\begin{xltabular}{\textwidth}{@{}l >{\raggedright\arraybackslash\hsize=0.773\hsize}X l >{\raggedright\arraybackslash\hsize=1.227\hsize}X@{}}
\toprule
\textbf{Modal} & \textbf{Non-modal} & \textbf{Command} & \textbf{What it does}\\
\midrule
\endhead
none & \texttt{TAB} & \texttt{org-\codebreak{}cycle} & On a headline, fold its subtree to its headline alone, then show its children’s headlines, then all of it, and fold it again, a step each time.\\
none & \texttt{S-TAB} & \texttt{org-\codebreak{}global-\codebreak{}cycle} & Cycle the whole buffer: its overview (the headlines of the first level), then its contents (every headline), then everything, and the overview again.\\
none & \texttt{M-RET} & \texttt{org-\codebreak{}meta-\codebreak{}return} & A new item or headline after what point is in and all under it: in a list, an item of the same kind after the one point is in and all under it (its bullet, the next number, [ ] when it has a checkbox).\\
none & \texttt{M-S-\codebreak{}RET} & \texttt{org-\codebreak{}insert-\codebreak{}todo-\codebreak{}heading} & A new TODO headline of this level after this subtree.\\
none & \texttt{M-j}, \texttt{M-<down>} & \texttt{org-\codebreak{}metadown} & In an Org table, move the row down; else move the subtree down (M-j).\\
none & \texttt{M-k}, \texttt{M-<up>} & \texttt{org-\codebreak{}metaup} & In an Org table, move the row up; else move the subtree up (M-k).\\
none & \texttt{M-h}, \texttt{M-<left>} & \texttt{org-\codebreak{}metaleft} & In an Org table, move the column left; else promote the headline (M-h).\\
none & \texttt{M-l}, \texttt{M-<right>} & \texttt{org-\codebreak{}metaright} & In an Org table, move the column right; else demote the headline (M-l).\\
none & \texttt{M-H}, \texttt{M-S-<left>} & \texttt{org-\codebreak{}shiftmetaleft} & In an Org table, delete the column; else promote the subtree (M-H).\\
none & \texttt{M-L}, \texttt{M-S-<right>} & \texttt{org-\codebreak{}shiftmetaright} & In an Org table, put an empty column before this one; else demote the subtree (M-L).\\
none & \texttt{M-x org-\codebreak{}move-\codebreak{}subtree-\codebreak{}down} & \texttt{org-\codebreak{}move-\codebreak{}subtree-\codebreak{}down} & Move the subtree point is in down, past its next sibling.\\
none & \texttt{M-x org-\codebreak{}move-\codebreak{}subtree-up} & \texttt{org-\codebreak{}move-\codebreak{}subtree-up} & Move the subtree point is in up, before its sibling above.\\
none & \texttt{M-x org-\codebreak{}do-\codebreak{}promote} & \texttt{org-\codebreak{}do-\codebreak{}promote} & Promote the headline point is in a level, its subtree’s other headlines as they are.\\
none & \texttt{M-x org-\codebreak{}do-\codebreak{}demote} & \texttt{org-\codebreak{}do-\codebreak{}demote} & Demote the headline point is in a level, its subtree’s other headlines as they are.\\
none & \texttt{M-x org-\codebreak{}promote-\codebreak{}subtree} & \texttt{org-\codebreak{}promote-\codebreak{}subtree} & Promote the headline point is in a level, and every headline under it.\\
none & \texttt{M-x org-\codebreak{}demote-\codebreak{}subtree} & \texttt{org-\codebreak{}demote-\codebreak{}subtree} & Demote the headline point is in a level, and every headline under it.\\
\bottomrule
\end{xltabular}

\begin{xltabular}{\textwidth}{@{}l l >{\raggedright\arraybackslash\hsize=1.000\hsize}X@{}}
\toprule
\textbf{Setting} & \textbf{Standard} & \textbf{What it changes}\\
\midrule
\endhead
\texttt{*org-\codebreak{}cycle-\codebreak{}hide-\codebreak{}drawers*} & \texttt{t} & True for drawers (:LOGBOOK:, :PROPERTIES:) and a task’s notes (ORG-NOTES-RANGE) to stay folded when TAB shows a headline’s or an item’s children or all of it; TAB on a drawer’s line, or on the first of the notes, opens it.\\
\texttt{*org-\codebreak{}fold-\codebreak{}metadata*} & \texttt{t} & True for an Org file to open with what is kept of each task folded: its drawers (:PROPERTIES:, :LOGBOOK:) and its notes (- State “DONE” \ldots{}, - Note taken on \ldots{}, CLOCK: lines) each to one line, those of headings and list items alike; its SCHEDULED: and DEADLINE: shown.\\
\texttt{*org-\codebreak{}cycle-\codebreak{}include-\codebreak{}plain-\codebreak{}lists*} & \texttt{t} & True for TAB on a list item with sub-items or text under its first line to fold it, as on a headline: folded, then its sub-items, then all of it; NIL for TAB there to indent.\\
\bottomrule
\end{xltabular}

\begin{xltabular}{\textwidth}{@{}l >{\raggedright\arraybackslash\hsize=1.000\hsize}X@{}}
\toprule
\textbf{Face} & \textbf{What it paints}\\
\midrule
\endhead
\texttt{:org-\codebreak{}level-1} & Org headlines of the first level, their stars and titles.\\
\texttt{:org-\codebreak{}level-2} & Org headlines of the second level.\\
\texttt{:org-\codebreak{}level-3} & Org headlines of the third level.\\
\texttt{:org-\codebreak{}level-4} & Org headlines of the fourth level.\\
\texttt{:org-\codebreak{}level-5} & Org headlines of the fifth level.\\
\texttt{:org-\codebreak{}level-6} & Org headlines of the sixth level.\\
\texttt{:org-\codebreak{}level-7} & Org headlines of the seventh level.\\
\texttt{:org-\codebreak{}level-8} & Org headlines of the eighth level and deeper.\\
\texttt{:org-\codebreak{}hide} & The stars of Org’s headlines hidden in the indented view: the ground’s colour (dim where it is not known).\\
\texttt{:org-\codebreak{}leading-\codebreak{}stars} & The stars of Org’s headlines before the last, shown faint (\textbf{org-hide-leading-stars} nil).\\
\texttt{:org-\codebreak{}archived} & Org headlines tagged ARCHIVE.\\
\texttt{:org-\codebreak{}drawer} & Org’s drawers’ first and last lines (:PROPERTIES:, :END:).\\
\texttt{:org-\codebreak{}property-\codebreak{}value} & The values in Org’s property drawers.\\
\bottomrule
\end{xltabular}

\subsection{Lists and checkboxes}
\label{sec:351}
A list is for things that do not each need a heading: what to buy, the steps
of a recipe, points to raise at a meeting.

A line that begins with a hyphen and a space is an item. A line that begins
with \texttt{1.} and a space is an item of a numbered list. An item set in further
than the one above it is a sub-item of that one. A line written as
\texttt{- term :: what it means} is an item of a description list, useful for a
short glossary.

\begin{Verbatim}[breaklines,breakautoindent=false,breakindent=1em,breakafter=/-.\,;:=,commandchars=\\\{\}]
- Weed the beds
- Mulch the feijoas
  - Bark, if there is any left
1. Soak the seeds overnight
2. Plant them a finger deep
\end{Verbatim}

\texttt{M-RET} on an item makes the next item of the same kind after it. In a
numbered list the items after it are numbered again. \texttt{C-c -} changes the
bullet of every item in the list to the next kind. The kinds are a hyphen, a
plus, a star, \texttt{1.} and \texttt{1)}, and a list at the left edge skips the star. That
is the quick way to turn a plain list into a numbered one.
\texttt{TAB} on an item that has sub-items or more lines under it folds it, as on a
heading.

\subsubsection{Checkboxes}
\label{sec:350}
An item that starts with \texttt{[ ]} has a checkbox, something to tick off. You do
not have to type all of it. Type the hyphen, a space and \texttt{[}, and Tuatara
completes the box and leaves the cursor after it.

\texttt{C-c C-c} on the item’s line ticks the box, which then reads \texttt{[X]}, or clears
it again. In NORMAL state \texttt{RET} on the line goes round three states: empty,
\texttt{[-]} for half done, and ticked.

Put \texttt{[/]} at the end of the heading above the list, or at the end of an item
that has sub-items, and Org counts for you. It shows how many boxes are
ticked out of how many there are, such as \texttt{[2/5]}, and keeps the count
current as you tick. \texttt{[\%]} shows a percentage instead.

When you tick a box, Tuatara notes the time in a small drawer under the item.
That lets the agenda show later what you finished on which day. If you do not
want these notes, set \texttt{*org-\codebreak{}log-\codebreak{}checkboxes*} to \texttt{nil}.

\begin{try}
\textbf{Try it, with a shopping list}
\begin{enumerate}
\item Press \texttt{C-x C-f}, type \texttt{\textasciitilde{}/shopping.\codebreak{}org} and press \texttt{RET}. Press \texttt{i} and type
   \texttt{* Shopping [/]}, then \texttt{RET}.
\item Type a hyphen, a space and \texttt{[}. The box completes itself. Type \texttt{eggs} and
   press \texttt{RET}. Do the same for \texttt{bread} and for \texttt{milk}.
\item Press \texttt{C-g}. Move to the line \texttt{bread} and press \texttt{C-c C-c}. The box reads
   \texttt{[X]} and the heading says \texttt{[1/3]}. A drawer under the item holds the
   time. It is shown open now and is folded when you open the file again.
\item Move to \texttt{eggs} and press \texttt{RET} twice. The box goes to \texttt{[-]} and then to
   \texttt{[X]}, and the heading says \texttt{[2/3]}.
\item Still on \texttt{eggs}, press \texttt{M-RET}, then \texttt{i}, and type \texttt{butter}. The new item
   has a box of its own. The count in the heading stays \texttt{[2/3]} until you next
   tick or clear a box. Press \texttt{C-g}.
\item Press \texttt{C-c -} a few times and watch the bullets change. Stop when they
   are hyphens again. Save with \texttt{C-x C-s}.
\end{enumerate}
\end{try}

The signs can be prettier than the plain text. \texttt{*org-\codebreak{}list-\codebreak{}bullets*} and
\texttt{*org-\codebreak{}checkbox-\codebreak{}signs*} name a character to show for each bullet and each
state of a box. The file itself keeps the plain marks.

\begin{xltabular}{\textwidth}{@{}l >{\raggedright\arraybackslash\hsize=0.435\hsize}X l >{\raggedright\arraybackslash\hsize=1.565\hsize}X@{}}
\toprule
\textbf{Modal} & \textbf{Non-modal} & \textbf{Command} & \textbf{What it does}\\
\midrule
\endhead
none & \texttt{C-c C-c} & \texttt{org-\codebreak{}ctrl-\codebreak{}c-\codebreak{}ctrl-c} & Do what Org’s C-c C-c does where point is: on a dynamic block’s \#+BEGIN: line, make the block again (M-x org-dblock-update).\\
none & \texttt{C-c -} & \texttt{org-\codebreak{}ctrl-\codebreak{}c-\codebreak{}minus} & Do what Org’s C-c - does where point is: in a table, a rule line below the row (org-table-insert-hline); in a list, the next kind of bullet (org-cycle-list-bullet).\\
none & \texttt{M-x org-\codebreak{}cycle-\codebreak{}list-\codebreak{}bullet} & \texttt{org-\codebreak{}cycle-\codebreak{}list-\codebreak{}bullet} & Give the list point is in its next kind of bullet (C-c -): - to + to * to 1.\\
none & \texttt{RET} & \texttt{org-\codebreak{}return} & RET in Org: on a link, follow it (\textbf{org-return-follows-link}); in a table, go to the row below, aligned; else a new line, as RET does in text.\\
\bottomrule
\end{xltabular}

\begin{xltabular}{\textwidth}{@{}l l >{\raggedright\arraybackslash\hsize=1.000\hsize}X@{}}
\toprule
\textbf{Setting} & \textbf{Standard} & \textbf{What it changes}\\
\midrule
\endhead
\texttt{*org-\codebreak{}complete-\codebreak{}checkboxes*} & \texttt{t} & True for [ typed right after a list item’s bullet and its space (- [) to make it a checkbox at once, - [ ] with point after the space, ready for the item’s text; NIL to type [ as anywhere else.\\
\texttt{*org-\codebreak{}return-\codebreak{}cycles-\codebreak{}checkboxes*} & \texttt{t} & True for M-x tuatara-normal-return in NORMAL state anywhere on a list item’s line with a checkbox to cycle it, [ ] to [-] to [X] and round, the counts above brought up to date (a link point is on is followed first); NIL to leave it to M-x org-ctrl-c-ctrl-c.\\
\texttt{*org-\codebreak{}log-\codebreak{}checkboxes*} & \texttt{t} & True to log when a list item’s checkbox is ticked or cleared, in a :LOGBOOK: drawer under the item: - State “DONE” from “TODO” and the time, as a TODO headline’s change of state is logged; NIL to log nothing.\\
\texttt{*org-\codebreak{}todo-\codebreak{}follows-\codebreak{}checkboxes*} & \texttt{t} & True for a TODO headline with a statistics cookie ([2/3], [66\%]) to become done by itself once all it counts are done, its TODO sub-headlines (or, with none, the checkboxes right under it), with CLOSED: and the time, and not done again when one is not yet done; up the tree, a headline above with a cookie following in turn.\\
\texttt{*org-\codebreak{}list-\codebreak{}bullets*} & \texttt{nil} & The signs list items show in place of their bullets, each (BULLET .\\
\texttt{*org-\codebreak{}checkbox-\codebreak{}signs*} & \texttt{nil} & The signs list items’ checkboxes show, each (CHECKBOX .\\
\bottomrule
\end{xltabular}

\begin{xltabular}{\textwidth}{@{}l >{\raggedright\arraybackslash\hsize=1.000\hsize}X@{}}
\toprule
\textbf{Face} & \textbf{What it paints}\\
\midrule
\endhead
\texttt{:org-\codebreak{}list-\codebreak{}bullet} & The signs list items show for their bullets (\textbf{org-list-bullets}).\\
\texttt{:org-\codebreak{}list-dt} & The terms of Org’s description lists (term :: description).\\
\texttt{:org-\codebreak{}checkbox} & The checkboxes of Org’s list items, [ ] and [X].\\
\bottomrule
\end{xltabular}

\subsection{Tasks}
\label{sec:356}
A checkbox is enough for a shopping list. A task is for work that takes
longer, that may have to wait for someone, or that has a date. A task is a
heading with a keyword after its stars:

\begin{Verbatim}[breaklines,breakautoindent=false,breakindent=1em,breakafter=/-.\,;:=,commandchars=\\\{\}]
* TODO Paint the fence
* WAIT Quote from the plumber
* DONE Ring the council
\end{Verbatim}

Four keywords come as standard. \texttt{TODO} is for something still to do. \texttt{WAIT}
is for something that is waiting on a person or an event. \texttt{DONE} is for
something finished, and \texttt{CANCELLED} for something you gave up. The first two
count as open and the last two as closed.

\texttt{C-c C-t} gives the heading the next keyword: none, \texttt{TODO}, \texttt{WAIT}, \texttt{DONE}
and none again. \texttt{S-<right>} and \texttt{S-<left>} go forward and back through the
same round. In NORMAL state \texttt{RET} on a task’s line also moves it to the next
keyword, and never leaves it without one. \texttt{CANCELLED} is left out of the
round so that you do not cancel something by pressing a key once too often.
\texttt{C-c C-x C} cancels the task the cursor is in.

When a task becomes \texttt{DONE}, Tuatara writes a line under it that says when:

\begin{Verbatim}[breaklines,breakautoindent=false,breakindent=1em,breakafter=/-.\,;:=,commandchars=\\\{\}]
* DONE Paint the fence
  CLOSED: [2026-10-13 Tue 10:00]
\end{Verbatim}

The agenda can list what you closed on each day from these lines. Taking
\texttt{DONE} away again removes the line. Set \texttt{*org-\codebreak{}log-\codebreak{}done*} to \texttt{nil} if you want
nothing written.

\subsubsection{Tasks with parts}
\label{sec:352}
A task can have smaller tasks under it, or a list of checkboxes. Put \texttt{[/]}
at the end of its heading and it counts them. When the last one is done, the
task becomes \texttt{DONE} by itself, and it becomes \texttt{TODO} again if you reopen one.
While a part is still open, \texttt{RET} passes over the closed keywords, so a task
is not marked done by accident with work left under it. \texttt{C-c C-t} still lets
you close it if you mean to.

\subsubsection{Notes on a task}
\label{sec:353}
\texttt{C-c C-z} asks for a note and files it under the task with the date and time.
Use it for what you would otherwise forget: “rang twice, no answer”, “the
paint is in the shed”. Notes go into the task’s \texttt{:LOGBOOK:} drawer, the
newest first, with the other records Tuatara keeps of the task. Open the
drawer with \texttt{TAB} on its first line.

\subsubsection{Keywords of your own}
\label{sec:354}
A note can name its own keywords on a line near the top. Those before the
bar are open and those after it are closed:

\begin{Verbatim}[breaklines,breakautoindent=false,breakindent=1em,breakafter=/-.\,;:=,commandchars=\\\{\}]
#+TODO: TODO NEXT WAIT | DONE CANCELLED
\end{Verbatim}

To change the keywords for every file that has no such line, set
\texttt{*org-\codebreak{}standard-\codebreak{}todo-\codebreak{}keywords*} in your init file. \texttt{*org-\codebreak{}todo-\codebreak{}keyword-\codebreak{}faces*}
gives a keyword a colour of its own.

\subsubsection{Priorities}
\label{sec:355}
A task can carry a priority, written by hand after its keyword as \texttt{[\#A]},
\texttt{[\#B]} or \texttt{[\#C]}. It is shown in its own colour, here and in the agenda, so
that the important tasks stand out.

\begin{xltabular}{\textwidth}{@{}l l l >{\raggedright\arraybackslash\hsize=1.000\hsize}X@{}}
\toprule
\textbf{Modal} & \textbf{Non-modal} & \textbf{Command} & \textbf{What it does}\\
\midrule
\endhead
none & \texttt{C-c C-t} & \texttt{org-\codebreak{}todo} & Change the TODO keyword of the headline point is in to the next of the document’s (\#+TODO:, else TODO, WAIT and DONE), none after the last, those of \textbf{org-todo-keywords-not-cycled} passed over (CANCELLED, which M-x org-cancel-task sets); a task done is noted with CLOSED: and when.\\
none & \texttt{S-<right>} & \texttt{org-\codebreak{}shiftright} & On a headline, its next TODO keyword, as M-x org-todo does.\\
none & \texttt{S-<left>} & \texttt{org-\codebreak{}shiftleft} & On a headline, its TODO keyword before, none before the first.\\
none & \texttt{C-c C-x C} & \texttt{org-\codebreak{}cancel-\codebreak{}task} & Cancel the task of the headline point is in: its keyword made CANCELLED, a done keyword that M-x org-todo and the other keys that go round the keywords pass over, noted with CLOSED: and when as a task done is.\\
none & \texttt{C-c C-z} & \texttt{org-\codebreak{}add-\codebreak{}note} & Add a note to the entry point is under (C-c C-z): the note is asked for, and logged where Org logs notes, at the top of its :LOGBOOK: drawer with \textbf{ORG-LOG-INTO-DRAWER}, else after its planning line and properties, in the form - Note taken on \textit{[2026-10-01 Thu 18:00]} \textbackslash{}\textbackslash{} and the note under it.\\
\bottomrule
\end{xltabular}

\begin{xltabular}{\textwidth}{@{}l >{\raggedright\arraybackslash\hsize=0.818\hsize}X >{\raggedright\arraybackslash\hsize=1.182\hsize}X@{}}
\toprule
\textbf{Setting} & \textbf{Standard} & \textbf{What it changes}\\
\midrule
\endhead
\texttt{*org-\codebreak{}standard-\codebreak{}todo-\codebreak{}keywords*} & \texttt{(("TODO" "WAIT") "DONE" "CANCELLED")} & The TODO keywords of an Org file with no \#+TODO: line of its own, as (NOT-DONE .\\
\texttt{*org-\codebreak{}todo-\codebreak{}keywords-\codebreak{}not-\codebreak{}cycled*} & \texttt{("CANCELLED" "CANCELED")} & The TODO keywords passed over as a task’s keyword goes round to the next: by M-x org-todo, M-x org-shiftright and M-x org-shiftleft, by RET on a task’s headline and by the agenda’s M-x org-agenda-todo.\\
\texttt{*org-\codebreak{}todo-\codebreak{}keyword-\codebreak{}faces*} & \texttt{(("WAIT" . :org-\codebreak{}wait) ("CANCELLED" . :or...} & The TODO keywords with a face of their own, each (KEYWORD .\\
\texttt{*org-\codebreak{}return-\codebreak{}cycles-\codebreak{}todo*} & \texttt{t} & True for M-x tuatara-normal-return in NORMAL state on a headline’s line with a TODO keyword to change it to the next, as M-x org-todo does, but never to none: TODO to WAIT to DONE and round, or round the document’s \#+TODO: keywords, those of \textbf{org-todo-keywords-not-cycled} passed over (a link point is on is followed first).\\
\texttt{*org-\codebreak{}log-\codebreak{}done*} & \texttt{:time} & What is noted when a task is done: :TIME (as standard) a CLOSED: line with the time, which the agenda’s log mode lists; NIL nothing.\\
\texttt{*org-\codebreak{}log-\codebreak{}into-\codebreak{}drawer*} & \texttt{t} & True for the notes logged of a task (a change of state, a repeat, a reschedule, a note taken) to go into the entry’s :LOGBOOK: drawer, where they fold with it (\textbf{ORG-FOLD-METADATA}); NIL for them to go right after its planning line and properties, outside any drawer.\\
\texttt{*org-\codebreak{}log-\codebreak{}repeat*} & \texttt{:time} & What Org notes when a task that repeats is done: :TIME the change of state and LAST\_REPEAT with the time; NIL nothing.\\
\texttt{*org-\codebreak{}log-\codebreak{}reschedule*} & \texttt{nil} & What Org notes when a task’s scheduled date is changed or taken away: NIL nothing, :TIME the date it had and when.\\
\texttt{*org-\codebreak{}log-\codebreak{}redeadline*} & \texttt{nil} & What Org notes when a task’s deadline is changed or taken away: NIL nothing, :TIME the deadline it had and when.\\
\bottomrule
\end{xltabular}

\begin{xltabular}{\textwidth}{@{}l >{\raggedright\arraybackslash\hsize=1.000\hsize}X@{}}
\toprule
\textbf{Face} & \textbf{What it paints}\\
\midrule
\endhead
\texttt{:org-\codebreak{}todo} & Org’s keywords of tasks not done (TODO).\\
\texttt{:org-\codebreak{}wait} & Org’s keyword of tasks open but waiting on something (WAIT), apart from TODO.\\
\texttt{:org-\codebreak{}done} & Org’s keywords of tasks done (DONE).\\
\texttt{:org-\codebreak{}cancelled} & Org’s keyword of tasks cancelled (CANCELLED), closed but apart from DONE.\\
\texttt{:org-\codebreak{}headline-\codebreak{}done} & The titles of Org headlines whose task is done.\\
\texttt{:org-\codebreak{}priority} & An Org headline’s priority, [\#A].\\
\texttt{:org-\codebreak{}log-\codebreak{}note} & The notes a task’s log keeps of its changes of state, wholly, the time included: - State “DONE” from “TODO” \textit{[2026-10-01 Thu 18:00]}.\\
\texttt{:org-\codebreak{}special-\codebreak{}keyword} & Org’s CLOSED: and CLOCK:.\\
\bottomrule
\end{xltabular}

\subsection{Tags}
\label{sec:357}
A tag is a word at the end of a heading, between colons, that says what the
heading is about:

\begin{Verbatim}[breaklines,breakautoindent=false,breakindent=1em,breakafter=/-.\,;:=,commandchars=\\\{\}]
* TODO Mow the lawn :outside:weekend:
\end{Verbatim}

Tags cut across the outline. A note about the house and a note about work can
both hold headings tagged \texttt{phone}, and the agenda can list every one of them
when you sit down to make calls. A heading also has the tags of the headings
above it, so tagging a whole project once is enough.

\texttt{C-c C-q} sets the tags of the heading the cursor is under. It offers the
tags the heading has. Type the words with spaces between them, or nothing to
remove them all. You can also type tags by hand.

As standard the tags sit one space after the heading’s text. If you like them
lined up at the right, set \texttt{*org-\codebreak{}tags-\codebreak{}column*} to a negative number, such as
\texttt{-77} for tags that end at column 77. \texttt{M-x org-\codebreak{}align-\codebreak{}tags} then puts the tags
of one heading in place, and \texttt{M-x org-\codebreak{}align-\codebreak{}all-\codebreak{}tags} those of every heading
in the buffer.

Two tags have a meaning to Tuatara. A heading tagged \texttt{notify} is announced on
your desktop when its time comes, as the next section explains. A heading
tagged \texttt{TOC} holds a table of contents, described under code blocks.

\begin{xltabular}{\textwidth}{@{}l >{\raggedright\arraybackslash\hsize=0.426\hsize}X l >{\raggedright\arraybackslash\hsize=1.574\hsize}X@{}}
\toprule
\textbf{Modal} & \textbf{Non-modal} & \textbf{Command} & \textbf{What it does}\\
\midrule
\endhead
none & \texttt{C-c C-q} & \texttt{org-\codebreak{}set-\codebreak{}tags-\codebreak{}command} & Set the tags of the headline point is under (C-c C-q): asked for with its tags offered, work urgent or :work:urgent:, nothing for none; they go at the end of its line, where \textbf{ORG-TAGS-COLUMN} says, by default one space after its text.\\
none & \texttt{M-x org-\codebreak{}align-\codebreak{}tags} & \texttt{org-\codebreak{}align-\codebreak{}tags} & Put the tags of the headline point is under where \textbf{ORG-TAGS-COLUMN} says: by default one space after its text.\\
none & \texttt{M-x org-\codebreak{}align-\codebreak{}all-\codebreak{}tags} & \texttt{org-\codebreak{}align-\codebreak{}all-\codebreak{}tags} & Put the tags of every headline of the buffer where \textbf{ORG-TAGS-COLUMN} says: by default one space after each one’s text, in one step to undo.\\
\bottomrule
\end{xltabular}

\begin{xltabular}{\textwidth}{@{}l l >{\raggedright\arraybackslash\hsize=1.000\hsize}X@{}}
\toprule
\textbf{Setting} & \textbf{Standard} & \textbf{What it changes}\\
\midrule
\endhead
\texttt{*org-\codebreak{}tags-\codebreak{}column*} & \texttt{0} & Where M-x org-set-tags-command and M-x org-align-tags put a headline’s tags: 0, one space after the headline’s text; negative, so that they end at that column (-77 to end at the 77th); positive, so that they start there.\\
\bottomrule
\end{xltabular}

\begin{longtable}{@{}l l@{}}
\toprule
\textbf{Face} & \textbf{What it paints}\\
\midrule
\endhead
\texttt{:org-\codebreak{}tag} & An Org headline’s tags, :work:urgent:.\\
\bottomrule
\end{longtable}

\subsection{Dates and scheduling}
\label{sec:361}
A date in an Org file is written in angle brackets, as \texttt{<2026-12-24 Thu>},
or with a time, as \texttt{<2026-12-24 Thu 14:00>}. You never have to type one out.
The commands of this section ask for a date and write it for you, weekday
and all.

There are three ways to put a date on something, and they mean different
things.

\begin{longtable}{@{}l l l l@{}}
\toprule
\textbf{What you want to say} & \textbf{How} & \textbf{Modal} & \textbf{Non-modal}\\
\midrule
\endhead
I plan to work on this that day & Schedule the task & none & \texttt{C-c C-s}\\
This must be finished by that day & Give it a deadline & none & \texttt{C-c C-d}\\
This happens then, whatever I do & A date in the text & none & \texttt{C-c .}\\
\bottomrule
\end{longtable}

Scheduling and deadlines are for tasks. They are written on a line of their
own right under the heading, beginning \texttt{SCHEDULED:} or \texttt{DEADLINE:}. A
scheduled task that you do not finish is carried forward in the agenda day
after day until you do. A deadline shows up in the agenda some days ahead,
14 unless you change \texttt{*org-\codebreak{}deadline-\codebreak{}warning-\codebreak{}days*}, so that it does not take
you by surprise.

A date in the text is for an event such as a flight, a birthday or a meeting.
\texttt{C-c .} puts one where the cursor is. The agenda shows it on its day and
nowhere else.

\texttt{C-c !} writes a date in square brackets, as \texttt{[2026-12-24 Thu]}. Such a date
is only a record. The agenda leaves it out. Use it to note when something
happened.

\subsubsection{Answering the date question}
\label{sec:358}
Every command that needs a date asks in the same way. A calendar of three
months opens below, with the cursor on today, and you can answer in either
of two ways.

You can point. Before you have typed anything, \texttt{l} and \texttt{h} move a day
forward and back and \texttt{j} and \texttt{k} move a week. \texttt{RET} writes the day at the
prompt, and a second \texttt{RET} accepts it. Between the two you can add a time,
such as \texttt{14:00}.

Or you can type, and Tuatara reads what you mean. The day it has understood
is shown as you type, so you can see that it is right before you press
\texttt{RET}.

\begin{xltabular}{\textwidth}{@{}>{\raggedright\arraybackslash\hsize=0.723\hsize}X >{\raggedright\arraybackslash\hsize=1.277\hsize}X@{}}
\toprule
\textbf{You type} & \textbf{It is read as}\\
\midrule
\endhead
nothing, only \texttt{RET} & Today\\
\texttt{fri} & The next Friday\\
\texttt{+3d}, \texttt{+2w}, \texttt{+1m} & Three days, two weeks or a month from today\\
\texttt{24} & The 24th of this month, or of the next once it is past\\
\texttt{dec 24}, \texttt{2026-12-24} & That date\\
\texttt{fri 9am}, \texttt{+2d 14:00} & A day with a time\\
\texttt{10:00-11:30} & Today, from then until then\\
\texttt{every week 10:00} & Each week at that time, starting today\\
\texttt{every second day} & Every other day\\
\texttt{first monday every month} & That day of each month\\
\bottomrule
\end{xltabular}

\subsubsection{Things that come round again}
\label{sec:359}
The last rows of the table make a date that repeats. A repeating date is
written with a mark after it, as in \texttt{<2026-10-13 Tue +1w>} for every week.
When you mark a repeating task \texttt{DONE}, it is not closed. Its date moves on to
the next time, it becomes \texttt{TODO} again, and a note of when you last did it is
kept under it. This is the way to keep track of the rubbish day, the rent
and watering the lemon tree.

If you cannot remember how to say a repeat in words, press \texttt{C-r} at the date
prompt. A menu opens where you choose the start, the time, how often and
when it ends, and it shows in plain words what your choices come to. \texttt{RET}
writes it.

\subsubsection{Being told on the desktop}
\label{sec:360}
Give a heading the tag \texttt{notify} and Tuatara shows a notification on your
desktop as its time comes, for as long as Tuatara is running. It tells you
ten minutes before and again at the time. An entry with a date and no time is
announced at eight in the morning. \texttt{*org-\codebreak{}notify-\codebreak{}before*} and
\texttt{*org-\codebreak{}notify-\codebreak{}day-\codebreak{}time*} change both. The notification is shown by the
program \texttt{notify-\codebreak{}send}, which chapter 8 mentions with the other notifications.

\begin{try}
\textbf{Try it, with a holiday plan}
\begin{enumerate}
\item Open \texttt{\textasciitilde{}/holiday.\codebreak{}org} with \texttt{C-x C-f}. Press \texttt{i} and type these four lines,
   then press \texttt{C-g}:
\begin{Verbatim}[breaklines,breakautoindent=false,breakindent=1em,breakafter=/-.\,;:=,commandchars=\\\{\}]
   * Holiday in Nelson
   ** TODO Book the ferry
   ** TODO Book the bach at Kaiteriteri
   ** Walk the Abel Tasman track
\end{Verbatim}
\item Move to \texttt{Book the ferry} and press \texttt{C-c C-d}. Type \texttt{+3d} and watch the
   calendar move three days on. Press \texttt{RET}. A line beginning \texttt{DEADLINE:}
   appears under the heading.
\item Move to \texttt{Book the bach} and press \texttt{C-c C-s}. Press \texttt{l} twice and \texttt{j} once,
   which is nine days from today. Press \texttt{RET} twice. The line under the
   heading begins \texttt{SCHEDULED:}.
\item Move to \texttt{Walk the Abel Tasman track} and press \texttt{o} to open a line under
   it. Press \texttt{C-c .}, type \texttt{dec 28 9am} and press \texttt{RET}. A date with a time
   is written into the text. Press \texttt{C-g}.
\item Move back to \texttt{Book the ferry} and press \texttt{C-c C-t} until it says \texttt{DONE}.
   A \texttt{CLOSED:} line joins the deadline. Save with \texttt{C-x C-s}.
\end{enumerate}
\end{try}

\begin{xltabular}{\textwidth}{@{}l l l >{\raggedright\arraybackslash\hsize=1.000\hsize}X@{}}
\toprule
\textbf{Modal} & \textbf{Non-modal} & \textbf{Command} & \textbf{What it does}\\
\midrule
\endhead
none & \texttt{C-c C-s} & \texttt{org-\codebreak{}schedule} & Schedule the task point is under (C-c C-s): the date is asked for (fri, +3d, 10-15, sep 30, 2026-12-24 14:00), what it reads shown as it is typed, and SCHEDULED: goes on its planning line, a repeater and warning it had kept.\\
none & \texttt{C-c C-d} & \texttt{org-\codebreak{}deadline} & Give the task point is under a deadline (C-c C-d): the date asked for as M-x org-schedule asks it, DEADLINE: on its planning line, a repeater and warning it had kept.\\
none & \texttt{C-c .} & \texttt{org-\codebreak{}time-\codebreak{}stamp} & Put in a timestamp at point (C-c .), the date asked for as M-x org-schedule asks it: \textit{<2026-10-13 Tue>}, or with a time \textit{<2026-10-13 Tue 14:00>}.\\
none & \texttt{C-c !} & \texttt{org-\codebreak{}time-\codebreak{}stamp-\codebreak{}inactive} & Put in an inactive timestamp at point (C-c !): as M-x org-time-stamp, written \textit{[2026-10-13 Tue]}, which the agenda leaves out.\\
\bottomrule
\end{xltabular}

These commands work only while a date is being asked for. You will seldom
call them by name, but they are what the keys above run.

\begin{xltabular}{\textwidth}{@{}l >{\raggedright\arraybackslash\hsize=0.700\hsize}X l >{\raggedright\arraybackslash\hsize=1.300\hsize}X@{}}
\toprule
\textbf{Modal} & \textbf{Non-modal} & \textbf{Command} & \textbf{What it does}\\
\midrule
\endhead
none & \texttt{M-x org-\codebreak{}calendar-\codebreak{}forward-\codebreak{}day} & \texttt{org-\codebreak{}calendar-\codebreak{}forward-\codebreak{}day} & In Org’s date prompt: the calendar’s day a day on, written at the prompt.\\
none & \texttt{M-x org-\codebreak{}calendar-\codebreak{}backward-\codebreak{}day} & \texttt{org-\codebreak{}calendar-\codebreak{}backward-\codebreak{}day} & In Org’s date prompt: the calendar’s day a day back, written at the prompt.\\
none & \texttt{M-x org-\codebreak{}calendar-\codebreak{}forward-\codebreak{}week} & \texttt{org-\codebreak{}calendar-\codebreak{}forward-\codebreak{}week} & In Org’s date prompt: the calendar’s day a week on, written at the prompt.\\
none & \texttt{M-x org-\codebreak{}calendar-\codebreak{}backward-\codebreak{}week} & \texttt{org-\codebreak{}calendar-\codebreak{}backward-\codebreak{}week} & In Org’s date prompt: the calendar’s day a week back, written at the prompt.\\
none & \texttt{M-x org-\codebreak{}calendar-\codebreak{}forward-\codebreak{}month} & \texttt{org-\codebreak{}calendar-\codebreak{}forward-\codebreak{}month} & In Org’s date prompt: the calendar’s day a month on, written at the prompt as the same day of the month, or the month’s last when it has fewer (31 January a month on is 28 February), as the calendar moves; +1m typed counts on past the month’s end instead (3 March).\\
none & \texttt{M-x org-\codebreak{}calendar-\codebreak{}backward-\codebreak{}month} & \texttt{org-\codebreak{}calendar-\codebreak{}backward-\codebreak{}month} & In Org’s date prompt: the calendar’s day a month back, written at the prompt as the same day of the month, or the month’s last when it has fewer, as the calendar moves; -1m typed counts on past the month’s end instead.\\
none & \texttt{M-x org-\codebreak{}calendar-\codebreak{}forward-\codebreak{}year} & \texttt{org-\codebreak{}calendar-\codebreak{}forward-\codebreak{}year} & In Org’s date prompt: the calendar’s day a year on, written at the prompt, so 29 February a year on is 28 February, as the calendar moves; +1y typed is 1 March.\\
none & \texttt{M-x org-\codebreak{}calendar-\codebreak{}backward-\codebreak{}year} & \texttt{org-\codebreak{}calendar-\codebreak{}backward-\codebreak{}year} & In Org’s date prompt: the calendar’s day a year back, written at the prompt, so 29 February a year back is 28 February, as the calendar moves; -1y typed is 1 March.\\
none & \texttt{M-x org-\codebreak{}calendar-\codebreak{}goto-\codebreak{}today} & \texttt{org-\codebreak{}calendar-\codebreak{}goto-\codebreak{}today} & In Org’s date prompt: the calendar on today, written at the prompt.\\
none & \texttt{M-x org-\codebreak{}date-\codebreak{}choose-\codebreak{}mark} & \texttt{org-\codebreak{}date-\codebreak{}choose-\codebreak{}mark} & In Org’s date question, answered in the calendar: mark the calendar’s day, beside those marked where the question takes several days, else in place of the one marked.\\
none & \texttt{M-x org-\codebreak{}date-\codebreak{}choose-\codebreak{}unmark} & \texttt{org-\codebreak{}date-\codebreak{}choose-\codebreak{}unmark} & In Org’s date question, answered in the calendar: take the mark off the calendar’s day.\\
none & \texttt{M-x org-\codebreak{}date-\codebreak{}choose-\codebreak{}write} & \texttt{org-\codebreak{}date-\codebreak{}choose-\codebreak{}write} & In Org’s date question, answered in the calendar: write the days chosen at the prompt (the calendar’s day when none is marked), to add a time or a repeat to; RET there takes it.\\
none & \texttt{M-x org-\codebreak{}date-\codebreak{}repeat-\codebreak{}menu} & \texttt{org-\codebreak{}date-\codebreak{}repeat-\codebreak{}menu} & At the date prompt of M-x org-schedule, M-x org-deadline and M-x org-time-stamp: leave it for the menu of a repeat (REPEAT-MENU), started from what was typed, or the calendar’s day.\\
none & \texttt{M-x repeat-\codebreak{}menu} & \texttt{repeat-\codebreak{}menu} & Open the menu of a repeat: its start, time, how it repeats and when it ends, what they come to shown under it; RET writes it.\\
none & \texttt{M-x repeat-\codebreak{}menu-\codebreak{}write} & \texttt{repeat-\codebreak{}menu-\codebreak{}write} & Write the repeat of the open repeat menu where the date prompt it came from would have written its date; say why when it cannot be written.\\
\bottomrule
\end{xltabular}

\begin{xltabular}{\textwidth}{@{}l l >{\raggedright\arraybackslash\hsize=1.000\hsize}X@{}}
\toprule
\textbf{Setting} & \textbf{Standard} & \textbf{What it changes}\\
\midrule
\endhead
\texttt{*org-\codebreak{}read-\codebreak{}date-\codebreak{}in-\codebreak{}calendar*} & \texttt{t} & True for Org’s date question (M-x org-schedule, M-x org-deadline, M-x org-time-stamp, a capture template’s \%\textasciicircum{}T) to begin in the calendar shown below it (\textbf{ORG-READ-DATE-POPUP-CALENDAR}).\\
\texttt{*org-\codebreak{}read-\codebreak{}date-\codebreak{}popup-\codebreak{}calendar*} & \texttt{t} & True for Org’s date prompt (M-x org-schedule, M-x org-deadline) to show the calendar below while it asks: the day read from what is typed marked there, S-<left> and S-<right> a day either way, S-<up> and S-<down> a week, M-S-<left> and M-S-<right> a month, M-S-<up> and M-S-<down> a year, C-.\\
\texttt{*org-\codebreak{}deadline-\codebreak{}warning-\codebreak{}days*} & \texttt{14} & How many days before its deadline a task is listed on today; a deadline’s own -3d says otherwise.\\
\texttt{*org-\codebreak{}notify-\codebreak{}tag*} & \texttt{"notify"} & The tag of the agenda’s entries told on the desktop as their time comes (a heading’s own or one it takes from those above it); NIL tells none.\\
\texttt{*org-\codebreak{}notify-\codebreak{}before*} & \texttt{10} & How many minutes before its time an entry tagged :notify: is told, as well as when the time comes; 0 or NIL for only then.\\
\texttt{*org-\codebreak{}notify-\codebreak{}day-\codebreak{}time*} & \texttt{"08:00"} & When an entry tagged :notify: with a date and no time of day is told on that day, or as Tuatara starts when that is later; NIL never tells them.\\
\bottomrule
\end{xltabular}

\begin{xltabular}{\textwidth}{@{}l >{\raggedright\arraybackslash\hsize=1.000\hsize}X@{}}
\toprule
\textbf{Face} & \textbf{What it paints}\\
\midrule
\endhead
\texttt{:org-\codebreak{}date} & Timestamps in Org (\textit{<2026-10-01 Thu>}).\\
\texttt{:org-\codebreak{}scheduled-\codebreak{}keyword} & Org’s SCHEDULED:, in the colour the agenda gives a task scheduled.\\
\texttt{:org-\codebreak{}deadline-\codebreak{}keyword} & Org’s DEADLINE:, in the colour the agenda gives a deadline.\\
\texttt{:calendar-\codebreak{}chosen} & Days chosen in the calendar for a date asked for, as Org’s date question has them chosen (M-x org-schedule, M-x org-time-stamp).\\
\bottomrule
\end{xltabular}

\subsection{The agenda}
\label{sec:368}
Dates scattered over a dozen files are of little use unless something
gathers them. The agenda does that. It reads your agenda files, finds every
task and every date in them, and shows them in one list in the order of the
calendar. You look at it in the morning to see what the day holds.

\subsubsection{Telling Tuatara which files to read}
\label{sec:362}
The agenda reads the files named in \texttt{*org-\codebreak{}agenda-\codebreak{}files*}. As it comes, that
is the one file \texttt{\textasciitilde{}/org/\codebreak{}agenda.\codebreak{}org}. Set it in your init file to the files or
folders where you keep tasks. A folder stands for every \texttt{.org} file in it.

\begin{Verbatim}[breaklines,breakautoindent=false,breakindent=1em,breakafter=/-.\,;:=,commandchars=\\\{\}]
(setf *org-agenda-files* '("~/org/" "~/holiday.org"))
\end{Verbatim}

\subsubsection{Opening it}
\label{sec:364}
\texttt{C-x a a} asks which view you want and waits for one more key.

\begin{xltabular}{\textwidth}{@{}l l >{\raggedright\arraybackslash\hsize=1.000\hsize}X@{}}
\toprule
\textbf{Modal} & \textbf{Non-modal} & \textbf{View}\\
\midrule
\endhead
none & \texttt{a} & The week, day by day\\
none & \texttt{t} & Every task that is still open, dated or not, with unticked checkboxes\\
none & \texttt{m} & The headings that match tags you type, such as \texttt{phone} or \texttt{work+urgent}\\
none & \texttt{u} & The open tasks that have no tag\\
none & \texttt{c} & The entries whose dates look wrong, each with the reason\\
\bottomrule
\end{xltabular}

\texttt{C-x a w}, \texttt{C-x a m} and \texttt{C-x a y} open the week, the month or the year
without asking. Figure 14.2 shows a week.

\begin{figure}[H]
\caption{Figure 14.2. The agenda of a week, seen on the Tuesday.}\label{org363}
\begin{Verbatim}[breaklines,breakautoindent=false,breakindent=1em,breakafter=/-.\,;:=,commandchars=\\\{\}]
Week-agenda (W42):
Monday     12 October 2026 W42
Tuesday    13 October 2026
  chores:     all day     Scheduled:  TODO Mow the lawn          :outside:
  chores:     all day     In   3 d.:  TODO Clean the gutters     :outside:
Wednesday  14 October 2026
Thursday   15 October 2026
  chores:      9:00 {\tuataraFallbackFOUR{}┄}{\tuataraFallbackFOUR{}┄}{\tuataraFallbackFOUR{}┄}{\tuataraFallbackFOUR{}┄}{\tuataraFallbackFOUR{}┄} Scheduled:  TODO Fix the gate
  holiday:    14:00 {\tuataraFallbackFOUR{}┄}{\tuataraFallbackFOUR{}┄}{\tuataraFallbackFOUR{}┄}{\tuataraFallbackFOUR{}┄}{\tuataraFallbackFOUR{}┄} Event:      Dentist in Motueka
Friday     16 October 2026
  chores:     all day     Deadline:   TODO Clean the gutters     :outside:
Saturday   17 October 2026
Sunday     18 October 2026
\end{Verbatim}
\end{figure}

Each line starts with the name of the file the entry is in, then its time or
\texttt{all day}, then what kind of line it is. \texttt{Scheduled:} is a task planned for
that day. \texttt{Deadline:} is a task due that day, and \texttt{In 3 d.:} on today is the
early warning of a deadline three days off. \texttt{Event:} is a date in the text.
The tags are at the right. In the figure, cleaning the gutters appears twice:
on Friday when it is due and on today as a warning.

On today the agenda also draws a time grid when something has a time, with a
line for the hour it is now, so that you see what is still ahead.

\subsubsection{Moving about}
\label{sec:365}
\texttt{j} and \texttt{k} move from line to line. \texttt{l} shows the next week and \texttt{h} the one
before. \texttt{g h} comes back to today. \texttt{D}, \texttt{w}, \texttt{m} and \texttt{Y} change the span to a
day, a week, a month or a year, and \texttt{g l} returns to the view you had before.
\texttt{RET} opens the entry on the line in its file, and \texttt{TAB} shows it in another
window while you stay in the agenda. \texttt{c} opens the calendar on the line’s
day, and \texttt{a} in the calendar shows the agenda of the week of the day the
cursor is on. \texttt{q} closes the agenda.

The agenda keeps itself current. It is drawn again when one of its files
changes, every ten minutes and just after midnight. \texttt{r} draws it again at
once.

\subsubsection{Changing tasks from the agenda}
\label{sec:366}
You do not have to open a file to deal with a task. The keys below act on
the entry of the line the cursor is on, and change it in its file.

\begin{longtable}{@{}l l l@{}}
\toprule
\textbf{Modal} & \textbf{Non-modal} & \textbf{What it does to the entry}\\
\midrule
\endhead
none & \texttt{T} & Gives it the next keyword, as \texttt{C-c C-t} does in the file\\
none & \texttt{C-c C-s} & Schedules it for another day\\
none & \texttt{d} & Gives it a deadline\\
none & \texttt{t} & Sets its tags\\
none & \texttt{n} & Adds a note to it\\
none & \texttt{I} / \texttt{O} & Starts and stops the clock on it\\
none & \texttt{u} & Undoes the last change made from this agenda\\
none & \texttt{s} & Saves every agenda file that has changes\\
\bottomrule
\end{longtable}

The usual keys from the file, such as \texttt{C-c C-t} and \texttt{C-c C-q}, work in the
agenda too. Changes made from the agenda are in the files’ buffers and are
not on disk until you press \texttt{s}. If you want them saved as they are made,
name the files in \texttt{*org-\codebreak{}agenda-\codebreak{}autosave-\codebreak{}files*}.

\texttt{L} turns on log mode. The agenda then also lists, on each day, the tasks you
closed and the time you clocked. It is a pleasant way to see on a Friday
what the week went on.

\texttt{i} takes down something new for the day the cursor is on, with capture,
which is described below.

\begin{try}
\textbf{Try it, with the holiday plan}
\begin{enumerate}
\item Tell the agenda about the holiday file for this session. Press \texttt{M-:},
   type \texttt{(push "\textasciitilde{}/\codebreak{}holiday.\codebreak{}org" *org-\codebreak{}agenda-\codebreak{}files*)} and press \texttt{RET}.
\item Press \texttt{C-x a a} and then \texttt{a}. The week is shown with the cursor on today.
   If booking the bach falls in the coming week, press \texttt{l} to see it there.
\item Press \texttt{g h} to come back. Press \texttt{C-x a a} and then \texttt{t}. Every open task is
   listed, \texttt{Book the bach at Kaiteriteri} among them.
\item Move to that line and press \texttt{T}. It becomes \texttt{WAIT}. Press \texttt{u} to take the
   change back.
\item Press \texttt{TAB}. The task is shown in its file in another window.
\item Press \texttt{s} to save, and \texttt{q} to close the agenda.
\end{enumerate}
\end{try}

\subsubsection{Views of your own}
\label{sec:367}
\texttt{*org-\codebreak{}agenda-\codebreak{}custom-\codebreak{}commands*} adds views under keys of your own to the
question \texttt{C-x a a} asks. Each has a key, a description and a match, such as
every open task tagged \texttt{phone}. The view under \texttt{u} is made this way, and its
entry in the setting is a pattern to copy.

\begin{xltabular}{\textwidth}{@{}l >{\raggedright\arraybackslash\hsize=0.558\hsize}X l >{\raggedright\arraybackslash\hsize=1.442\hsize}X@{}}
\toprule
\textbf{Modal} & \textbf{Non-modal} & \textbf{Command} & \textbf{What it does}\\
\midrule
\endhead
none & \texttt{a}, \texttt{C-c a}, \texttt{C-x a a} & \texttt{org-\codebreak{}agenda} & Show the agenda, asking which view: a the agenda of the week (or \textbf{ORG-AGENDA-SPAN}), t the TODO list, m the headlines a match selects (work+urgent-home, TODO=”NEXT”), or a view of your own (\textbf{ORG-AGENDA-CUSTOM-COMMANDS}, u the tasks with no tag), or c the entries whose dates look wrong (M-x org-agenda-check-dates), read from \textbf{ORG-AGENDA-FILES}.\\
none & \texttt{M-x org-\codebreak{}agenda-\codebreak{}list} & \texttt{org-\codebreak{}agenda-\codebreak{}list} & Show the agenda of this week (or \textbf{ORG-AGENDA-SPAN}): each day with its scheduled tasks, deadlines and timestamps.\\
none & \texttt{M-x org-\codebreak{}todo-\codebreak{}list} & \texttt{org-\codebreak{}todo-\codebreak{}list} & Show the TODO list: every task not done in \textbf{ORG-AGENDA-FILES}.\\
none & \texttt{M-x org-\codebreak{}tags-\codebreak{}view} & \texttt{org-\codebreak{}tags-\codebreak{}view} & Show the headlines of \textbf{ORG-AGENDA-FILES} an Org match selects, read in the minibuffer: work, +work-home, work\(\vert{}\)proj, TODO=”NEXT”, LEVEL>1, -\{.*\} for those with no tag.\\
none & \texttt{M-x org-\codebreak{}agenda-\codebreak{}check-\codebreak{}dates} & \texttt{org-\codebreak{}agenda-\codebreak{}check-\codebreak{}dates} & Show the entries of \textbf{ORG-AGENDA-FILES} whose dates look wrong, each with what is wrong (ORG-ENTRY-DATE-PROBLEM).\\
none & \texttt{D}, \texttt{C-c D} & \texttt{org-\codebreak{}agenda-\codebreak{}day-\codebreak{}view} & Show the agenda of one day: the day at point in the agenda, else today.\\
none & \texttt{w}, \texttt{C-c w}, \texttt{C-x a w} & \texttt{org-\codebreak{}agenda-\codebreak{}week-\codebreak{}view} & Show the agenda of a week: the week of the day at point in the agenda, else this week.\\
none & \texttt{m}, \texttt{C-c m}, \texttt{C-x a m} & \texttt{org-\codebreak{}agenda-\codebreak{}month-\codebreak{}view} & Show the agenda of a month: the month of the day at point in the agenda, else this month.\\
none & \texttt{Y}, \texttt{C-c Y}, \texttt{C-x a y} & \texttt{org-\codebreak{}agenda-\codebreak{}year-\codebreak{}view} & Show the agenda of a year: the year of the day at point in the agenda, else this year.\\
none & \texttt{l}, \texttt{C-c l} & \texttt{org-\codebreak{}agenda-\codebreak{}later} & Show the agenda’s next span: the next week in a week’s agenda, the next day in a day’s.\\
none & \texttt{h}, \texttt{C-c h} & \texttt{org-\codebreak{}agenda-\codebreak{}earlier} & Show the agenda’s previous span: the previous week in a week’s agenda, the previous day in a day’s.\\
none & \texttt{g h} & \texttt{org-\codebreak{}agenda-\codebreak{}goto-\codebreak{}today} & Show the agenda’s span holding today, point on today’s date.\\
none & \texttt{g l} & \texttt{org-\codebreak{}agenda-\codebreak{}previous-\codebreak{}view} & Go back to the view or span the agenda showed before this one, point on the line it was on, as cd - goes back to the directory before: pressed again, return.\\
none & \texttt{L}, \texttt{C-c L} & \texttt{org-\codebreak{}agenda-\codebreak{}log-\codebreak{}mode} & List what was closed and clocked on each day too, or no longer; \textbf{ORG-AGENDA-LOG-MODE} says how the agenda starts.\\
none & \texttt{c}, \texttt{C-c c} & \texttt{org-\codebreak{}agenda-\codebreak{}goto-\codebreak{}calendar} & Show the calendar on the day of the agenda line at point.\\
none & \texttt{a}, \texttt{C-c a} & \texttt{org-\codebreak{}calendar-\codebreak{}goto-\codebreak{}agenda} & Show the agenda of the week holding the calendar’s day at point, point on that day.\\
\bottomrule
\end{xltabular}

\begin{xltabular}{\textwidth}{@{}l >{\raggedright\arraybackslash\hsize=0.600\hsize}X l >{\raggedright\arraybackslash\hsize=1.400\hsize}X@{}}
\toprule
\textbf{Modal} & \textbf{Non-modal} & \textbf{Command} & \textbf{What it does}\\
\midrule
\endhead
none & \texttt{RET} & \texttt{org-\codebreak{}agenda-\codebreak{}switch-to} & Open the entry on the agenda’s line in this window, point on its headline.\\
none & \texttt{TAB} & \texttt{org-\codebreak{}agenda-\codebreak{}goto} & Show the entry on the agenda’s line in the other window and go there.\\
none & \texttt{T}, \texttt{C-c T}, \texttt{C-c C-t} & \texttt{org-\codebreak{}agenda-\codebreak{}todo} & Change the TODO keyword of the entry on the agenda’s line to its next, as C-c C-t does in its file: a task done is noted with CLOSED:, one that repeats moves on to its next date.\\
none & \texttt{C-c C-x C} & \texttt{org-\codebreak{}agenda-\codebreak{}cancel-\codebreak{}task} & Cancel the task on the agenda’s line, as C-c C-x C does in its file: its keyword made CANCELLED and noted with CLOSED:.\\
none & \texttt{C-c C-s} & \texttt{org-\codebreak{}agenda-\codebreak{}schedule} & Schedule the entry on the agenda’s line, the date asked for as C-c C-s asks it; with C-u, take its scheduling away.\\
none & \texttt{d}, \texttt{C-c d}, \texttt{C-c C-d} & \texttt{org-\codebreak{}agenda-\codebreak{}deadline} & Give the entry on the agenda’s line a deadline, the date asked for as C-c C-d asks it; with C-u, take it away.\\
none & \texttt{t}, \texttt{C-c C-q} & \texttt{org-\codebreak{}agenda-\codebreak{}set-\codebreak{}tags} & Set the tags of the entry on the agenda’s line, asked for as C-c C-q asks them.\\
none & \texttt{n}, \texttt{C-c n}, \texttt{C-c C-z} & \texttt{org-\codebreak{}agenda-\codebreak{}add-\codebreak{}note} & Add a note to the entry on the agenda’s line, as C-c C-z adds one.\\
none & \texttt{C-c C-w} & \texttt{org-\codebreak{}agenda-\codebreak{}refile} & Move the entry on the agenda’s line under another heading, as C-c C-w moves one.\\
none & \texttt{I}, \texttt{C-c I} & \texttt{org-\codebreak{}agenda-\codebreak{}clock-in} & Clock in on the entry on the agenda’s line, as M-x org-clock-in does: a clock running is clocked out first.\\
none & \texttt{O}, \texttt{C-c O} & \texttt{org-\codebreak{}agenda-\codebreak{}clock-\codebreak{}out} & Clock out, as C-c C-x C-o does, from the agenda: the clock’s file saved at once when it is one of \textbf{ORG-AGENDA-AUTOSAVE-FILES}.\\
none & \texttt{C}, \texttt{C-c C} & \texttt{org-\codebreak{}agenda-\codebreak{}clock-\codebreak{}cancel} & Cancel the clock, as C-c C-x C-q does, from the agenda: the clock’s file saved at once when it is one of \textbf{ORG-AGENDA-AUTOSAVE-FILES}.\\
none & \texttt{u}, \texttt{C-/}, \texttt{C-\_} & \texttt{org-\codebreak{}agenda-\codebreak{}undo} & Undo the last change made from this agenda, in the file it was made in, and draw the agenda again; when the file was changed there since, nothing is undone and you are told to undo in the file first.\\
none & \texttt{s} & \texttt{org-\codebreak{}agenda-\codebreak{}save-\codebreak{}all-\codebreak{}buffers} & Save the buffers of the agenda files (\textbf{org-agenda-files}, and those added to them) with changes not saved, those the agenda changed among them; other Org buffers are left as they are (M-x org-save-all-org-buffers saves them all).\\
none & \texttt{M-x org-\codebreak{}save-\codebreak{}all-\codebreak{}org-\codebreak{}buffers} & \texttt{org-\codebreak{}save-\codebreak{}all-\codebreak{}org-\codebreak{}buffers} & Save every Org buffer with changes not saved: those the agenda changed among them, and any other.\\
none & \texttt{i}, \texttt{C-c i} & \texttt{org-\codebreak{}capture-\codebreak{}on-\codebreak{}day} & Capture, as C-x a c does, for the day at point in the calendar or on the agenda (the day of its line, or of the date above it), so that the agenda lists the entry that day: the template’s \%t and \%T are that day’s date, and a template with neither gets that date on the line after its heading.\\
\bottomrule
\end{xltabular}

\begin{xltabular}{\textwidth}{@{}l >{\raggedright\arraybackslash\hsize=1.143\hsize}X >{\raggedright\arraybackslash\hsize=0.857\hsize}X@{}}
\toprule
\textbf{Setting} & \textbf{Standard} & \textbf{What it changes}\\
\midrule
\endhead
\texttt{*org-\codebreak{}agenda-\codebreak{}files*} & \texttt{("\textasciitilde{}/\codebreak{}org/\codebreak{}agenda.\codebreak{}org")} & The files the agenda reads: each a file, a folder (every .org file in it, those whose names begin with a dot left out) or a pattern of names in a folder, * for any characters and ? for one (“\textasciitilde{}/notes/*.org”), names beginning with a dot only for a pattern that does.\\
\texttt{*org-\codebreak{}agenda-\codebreak{}span*} & \texttt{:week} & How many days M-x org-agenda’s agenda shows: :day, :week, :month or :year, or a number of days.\\
\texttt{*org-\codebreak{}agenda-\codebreak{}start-\codebreak{}on-\codebreak{}weekday*} & \texttt{1} & The day a week’s agenda starts on, 1 for Monday to 6 for Saturday and 7 or 0 for Sunday; NIL to start on today.\\
\texttt{*org-\codebreak{}agenda-\codebreak{}log-\codebreak{}mode*} & \texttt{nil} & True for the agenda to list what was closed and clocked on each day too.\\
\texttt{*org-\codebreak{}agenda-\codebreak{}refresh-\codebreak{}minutes*} & \texttt{10} & Minutes after which an agenda shown is drawn again by itself, so that the time now in its grid and the days left to deadlines stay true; NIL for never.\\
\texttt{*org-\codebreak{}agenda-\codebreak{}autosave-\codebreak{}files*} & \texttt{nil} & The agenda files saved as soon as a change made from the agenda (a keyword, a date, tags, a note, the clock) changes them: NIL for none (M-x org-agenda-save-all-buffers saves them), T for every agenda file, or a list of files, folders and patterns as \textbf{ORG-AGENDA-FILES} takes them: ’(“\textasciitilde{}/notes/agenda.org”).\\
\texttt{*org-\codebreak{}agenda-\codebreak{}custom-\codebreak{}commands*} & \texttt{(("u" "Untagged tasks" :tags-\codebreak{}todo "-\{.*\}...} & Views of your own that M-x org-agenda offers beside its own: each (KEY DESCRIPTION TYPE MATCH :header HEADER).\\
\texttt{*org-\codebreak{}agenda-\codebreak{}time-\codebreak{}grid*} & \texttt{((:\codebreak{}daily :today :require-\codebreak{}timed) (800 100...} & Today’s time grid in the agenda: (WHEN TIMES AFTER-TIME LINE).\\
\texttt{*org-\codebreak{}agenda-\codebreak{}current-\codebreak{}time-\codebreak{}string*} & \texttt{"← now ────────────────\codebreak{}────────────────\codebreak{}─...} & The line that shows the time now in today’s time grid; NIL for none.\\
\texttt{*org-\codebreak{}agenda-\codebreak{}category-\codebreak{}width*} & \texttt{12} & How wide the agenda’s first column is, its category (the file’s name, or its \#+CATEGORY:) and a colon, in columns: a longer one is cut short with …, so that every line’s time and text stand in one column.\\
\texttt{*org-\codebreak{}agenda-\codebreak{}untimed-\codebreak{}text*} & \texttt{"all day"} & What the agenda’s time column says on a line with no time of day, faint: “all day”, so that every line’s kind and text stand in one column; ”” leaves the column empty, lined up all the same; NIL leaves it out, the line’s text where the time would have been.\\
\texttt{*org-\codebreak{}agenda-\codebreak{}event-\codebreak{}kind-\codebreak{}text*} & \texttt{"Event:"} & What the agenda’s column of kinds (Scheduled:, Deadline:) says on the line of an event, a date in the text of an entry with no TODO keyword (:org-agenda-event-kind), Event:.\\
\texttt{*org-\codebreak{}agenda-\codebreak{}task-\codebreak{}kind-\codebreak{}text*} & \texttt{"Task:"} & What the agenda’s column of kinds says on the line of a date in the text of a task, an entry with a TODO keyword, not on its SCHEDULED: or DEADLINE: line: Task:, told apart from an event’s (\textbf{ORG-AGENDA-EVENT-KIND-TEXT}, which says whether the column is there at all).\\
\texttt{*org-\codebreak{}agenda-\codebreak{}date-\codebreak{}in-\codebreak{}title*} & \texttt{nil} & True to keep the dates written in a heading or a checkbox’s line in its title on the agenda’s days, as Org shows them; NIL takes them out there, the day and the time column saying them already (the TODO list and matches keep them).\\
\bottomrule
\end{xltabular}

\begin{xltabular}{\textwidth}{@{}l >{\raggedright\arraybackslash\hsize=1.000\hsize}X@{}}
\toprule
\textbf{Face} & \textbf{What it paints}\\
\midrule
\endhead
\texttt{:org-\codebreak{}agenda-\codebreak{}structure} & The agenda’s headers: Week-agenda (W40):, Global list of TODO items.\\
\texttt{:org-\codebreak{}agenda-\codebreak{}date} & A day’s date in the agenda: Monday     28 September 2026.\\
\texttt{:org-\codebreak{}agenda-\codebreak{}date-\codebreak{}today} & Today’s date in the agenda.\\
\texttt{:org-\codebreak{}agenda-\codebreak{}date-\codebreak{}weekend} & The date of a Saturday or a Sunday in the agenda.\\
\texttt{:org-\codebreak{}agenda-\codebreak{}date-\codebreak{}weekend-\codebreak{}today} & Today’s date in the agenda, on a Saturday or a Sunday.\\
\texttt{:org-\codebreak{}agenda-\codebreak{}category} & The category an agenda line begins with, the file’s name or its CATEGORY (garden:).\\
\texttt{:org-\codebreak{}agenda-\codebreak{}time} & The time of an agenda line, 9:30-10:15.\\
\texttt{:org-\codebreak{}agenda-\codebreak{}untimed} & The time column of an agenda line with no time of day, all day (\textbf{org-agenda-untimed-text}).\\
\texttt{:org-\codebreak{}agenda-\codebreak{}event-\codebreak{}kind} & The kind of an agenda line of a date in an entry’s text, `Event:’ or `Task:’ for a task’s (\textbf{org-agenda-event-kind-text}, \textbf{org-agenda-task-kind-text}).\\
\texttt{:org-\codebreak{}agenda-\codebreak{}calendar-\codebreak{}event} & An entry at a time or on a range of days, and what was clocked, in the agenda.\\
\texttt{:org-\codebreak{}agenda-\codebreak{}tags} & The tags at the end of an agenda line, :work:meeting:.\\
\texttt{:org-\codebreak{}agenda-\codebreak{}done} & A task done, and what was closed, in the agenda.\\
\texttt{:org-\codebreak{}agenda-\codebreak{}current-\codebreak{}time} & The time now in the agenda’s time grid, ← now.\\
\texttt{:org-\codebreak{}time-\codebreak{}grid} & The lines of today’s time grid in the agenda, 8:00 {\tuataraFallbackFOUR{}┄}{\tuataraFallbackFOUR{}┄}{\tuataraFallbackFOUR{}┄}{\tuataraFallbackFOUR{}┄}{\tuataraFallbackFOUR{}┄}.\\
\texttt{:org-\codebreak{}scheduled} & A task scheduled for a day to come, in the agenda.\\
\texttt{:org-\codebreak{}scheduled-\codebreak{}today} & A task scheduled for today, in the agenda.\\
\texttt{:org-\codebreak{}scheduled-\codebreak{}previously} & A task scheduled for a day gone and not done, in the agenda: Sched.\\
\texttt{:org-\codebreak{}upcoming-\codebreak{}distant-\codebreak{}deadline} & A deadline early in its warning, in the agenda.\\
\texttt{:org-\codebreak{}upcoming-\codebreak{}deadline} & A deadline in the second half of its warning, in the agenda.\\
\texttt{:org-\codebreak{}imminent-\codebreak{}deadline} & A deadline today, or gone, in the agenda.\\
\texttt{:org-\codebreak{}checkbox-\codebreak{}todo} & TODO in the agenda for a list item’s checkbox kept as a task, apart from a headline’s TODO.\\
\bottomrule
\end{xltabular}

\subsection{Capture}
\label{sec:370}
A thought arrives while you are in the middle of something else: ring the
plumber, an idea for the letter, a book someone recommended. Capture takes it
down without making you leave what you are doing. You press a key, type the
thought in a small window, press another key, and you are back where you
were. The thought has been filed in the right place.

\texttt{C-x a c} starts it and shows the templates to choose from. A template says
what kind of entry to make and where to file it. Four come with Tuatara.

\begin{xltabular}{\textwidth}{@{}l l >{\raggedright\arraybackslash\hsize=0.951\hsize}X >{\raggedright\arraybackslash\hsize=1.049\hsize}X@{}}
\toprule
\textbf{Modal} & \textbf{Non-modal} & \textbf{Makes} & \textbf{Files it}\\
\midrule
\endhead
none & \texttt{t} & A task, after asking for its tags & Under Tasks in your agenda file\\
none & \texttt{n} & A quick note & Under Inbox in your agenda file\\
none & \texttt{m} & Notes of a meeting, with the clock running & Under Meetings in your agenda file\\
none & \texttt{j} & An entry in your journal & In today’s journal entry\\
\bottomrule
\end{xltabular}

“Your agenda file” is the first file named in \texttt{*org-\codebreak{}agenda-\codebreak{}files*}. After
you choose a template, a window called \texttt{*Capture*} opens below, in INSERT
state, with the template filled in. Its top line tells you the keys. Write
what you have to say and finish with one of them.

\begin{longtable}{@{}l l l@{}}
\toprule
\textbf{Modal} & \textbf{Non-modal} & \textbf{What it does}\\
\midrule
\endhead
none & \texttt{C-c C-c} & Files the entry where the template says\\
none & \texttt{C-c C-w} & Files it under a heading you choose instead\\
none & \texttt{C-c C-k} & Throws it away\\
\bottomrule
\end{longtable}

A captured task or note carries the time you wrote it and a link back to
where you were, so that “check this paragraph” still makes sense a week
later.

In the calendar and in the agenda, \texttt{i} captures for the day the cursor is on
and not for today. The entry gets that day’s date, so the agenda lists it
there. That is how you write down a dentist’s appointment three weeks away:
go to the day and press \texttt{i}.

\begin{try}
\textbf{Try it, with something to remember}
\begin{enumerate}
\item From any buffer, press \texttt{C-x a c} and then \texttt{t}.
\item The echo area asks for tags. Type \texttt{phone} and press \texttt{RET}.
\item In the window that opens, type \texttt{Ring Wanaka Plumbing about the tap}.
\item Press \texttt{C-c C-c}. The echo area says where it was filed and you are back
   where you started. If the folder \texttt{\textasciitilde{}/org/} does not exist yet, Tuatara
   asks whether to make it.
\item Press \texttt{C-x a a} and then \texttt{t}. The new task is in the list.
\end{enumerate}
\end{try}

\subsubsection{Templates of your own}
\label{sec:369}
A template is a small file. The part above the line \texttt{\# --} says what it is
called, which key chooses it and where its entries go. The part below is the
entry, with marks that are filled in when you capture.

\begin{Verbatim}[breaklines,breakautoindent=false,breakindent=1em,breakafter=/-.\,;:=,commandchars=\\\{\}]
# name: Book to read
# key: b
# file: ~/org/reading.org
# heading: To read
# --
* %?
Recommended by %^\char123{}Who recommended it\char125{}
%U
\end{Verbatim}

\begin{longtable}{@{}l l@{}}
\toprule
\textbf{Mark} & \textbf{Is filled with}\\
\midrule
\endhead
\texttt{\%?} & Nothing. The cursor is left here for you to write\\
\texttt{\%U}, \texttt{\%u} & Now, or today, as a record the agenda leaves out\\
\texttt{\%T}, \texttt{\%t} & Now, or today, as a date the agenda shows\\
\texttt{\%\textasciicircum{}T}, \texttt{\%\textasciicircum{}t} & A date that you are asked for\\
\texttt{\%\textasciicircum{}\{Words\}} & What you type when asked the words\\
\texttt{\%\textasciicircum{}g} & Tags that you are asked for\\
\texttt{\%a} & A link to where you were\\
\texttt{\%i} & The text that was selected\\
\bottomrule
\end{longtable}

Your templates live in the folder \texttt{capture-\codebreak{}templates} in your configuration
folder, one file each. The easy way to start is \texttt{M-x
org-\codebreak{}capture-\codebreak{}template-\codebreak{}copy}, which copies one of Tuatara’s there for you to
change. A template of yours with the same key as one of Tuatara’s replaces
it.

\begin{xltabular}{\textwidth}{@{}l >{\raggedright\arraybackslash\hsize=0.667\hsize}X l >{\raggedright\arraybackslash\hsize=1.333\hsize}X@{}}
\toprule
\textbf{Modal} & \textbf{Non-modal} & \textbf{Command} & \textbf{What it does}\\
\midrule
\endhead
none & \texttt{C-x a c} & \texttt{org-\codebreak{}capture} & Take a note, a task or a journal entry down from wherever you are: press the key of one of the capture templates (Tuatara’s and yours, in \textbf{CAPTURE-TEMPLATE-DIRECTORIES}), answer what it asks, write in the \textbf{Capture} buffer and file it with C-c C-c (or throw it away with C-c C-k); you are back where you were.\\
none & \texttt{C-c C-c} & \texttt{org-\codebreak{}capture-\codebreak{}finalize} & File the text of this \textbf{Capture} buffer where its template says, and go back to where you were; with its template’s clock, clock it.\\
none & \texttt{C-c C-w} & \texttt{org-\codebreak{}capture-\codebreak{}refile} & File the text of this \textbf{Capture} buffer under a heading chosen as C-c C-w chooses one, instead of where its template says.\\
none & \texttt{C-c C-k} & \texttt{org-\codebreak{}capture-\codebreak{}kill} & Throw away what this \textbf{Capture} buffer holds, and go back to where you were.\\
none & \texttt{M-x org-\codebreak{}capture-\codebreak{}template-\codebreak{}copy} & \texttt{org-\codebreak{}capture-\codebreak{}template-\codebreak{}copy} & Copy one of Tuatara’s own capture templates, chosen by its name, into your capture templates folder (the first of \textbf{CAPTURE-TEMPLATE-DIRECTORIES}) and visit it: your copy, changed as you like (where it files, its text), takes the place of Tuatara’s on the same key.\\
\bottomrule
\end{xltabular}

\begin{xltabular}{\textwidth}{@{}l l >{\raggedright\arraybackslash\hsize=1.000\hsize}X@{}}
\toprule
\textbf{Setting} & \textbf{Standard} & \textbf{What it changes}\\
\midrule
\endhead
\texttt{*capture-\codebreak{}template-\codebreak{}directories*} & \texttt{nil} & The folders your capture templates are in, a template a file, as M-x org-capture offers them with Tuatara’s own (t, n, j, m); NIL, as standard, for capture-templates/ in your configuration folder (\textasciitilde{}/.config/tuatara/capture-templates/, XDG\_CONFIG\_HOME’s when it is set).\\
\bottomrule
\end{xltabular}

\subsection{Moving an entry to another place}
\label{sec:371}
An inbox fills up. Refiling moves a heading, with everything under it, to
its proper home. \texttt{C-c C-w} on a heading asks where it should go and lists
the headings of your agenda files and of the file you are in, each written
as a path such as \texttt{house.\codebreak{}org/\codebreak{}Garden/\codebreak{}Fruit}. Type a few letters to narrow the
list and press \texttt{RET}. The heading is put at the end of what is under the one
you chose, one level below it.

With \texttt{C-u} first, the command goes to the heading you choose and moves
nothing. That makes it a quick way to get anywhere in your agenda files.
The same key works on a line in the agenda. \texttt{*org-\codebreak{}refile-\codebreak{}targets*} says
which files are offered and how deep their headings go.

\begin{xltabular}{\textwidth}{@{}l l l >{\raggedright\arraybackslash\hsize=1.000\hsize}X@{}}
\toprule
\textbf{Modal} & \textbf{Non-modal} & \textbf{Command} & \textbf{What it does}\\
\midrule
\endhead
none & \texttt{C-c C-w} & \texttt{org-\codebreak{}refile} & Move the heading point is under, with all under it, to another heading: chosen by its path (agenda.org/Commitments and Jobs) among the headings \textbf{ORG-REFILE-TARGETS} offers, put at the end of what is under it, a level below; a clock running in it goes on there.\\
\bottomrule
\end{xltabular}

\begin{xltabular}{\textwidth}{@{}l >{\raggedright\arraybackslash\hsize=0.736\hsize}X >{\raggedright\arraybackslash\hsize=1.264\hsize}X@{}}
\toprule
\textbf{Setting} & \textbf{Standard} & \textbf{What it changes}\\
\midrule
\endhead
\texttt{*org-\codebreak{}refile-\codebreak{}targets*} & \texttt{((:\codebreak{}agenda-\codebreak{}files . 3) (:this-\codebreak{}file . 3))} & Where M-x org-refile offers to move an entry: a list of (WHERE .\\
\bottomrule
\end{xltabular}

\subsection{Clocking}
\label{sec:372}
The clock records how long you spend on a task. It is for anyone who bills
by the hour, and for anyone curious where a day went.

\texttt{C-c C-x C-i} starts the clock on the task the cursor is under. A line
beginning \texttt{CLOCK:} is written in the task’s \texttt{:LOGBOOK:} drawer with the time
you started. While the clock runs, the mode line of every window shows the
time so far and the task’s name, so you cannot forget it. Starting the clock
on another task stops the first.

\texttt{C-x a o} stops the clock from any buffer. The line is finished with the
time you stopped and how long that was:

\begin{Verbatim}[breaklines,breakautoindent=false,breakindent=1em,breakafter=/-.\,;:=,commandchars=\\\{\}]
CLOCK: [2026-10-13 Tue 10:00]--[2026-10-13 Tue 10:45] =>  0:45
\end{Verbatim}

\texttt{C-x a j} takes you to the task that is being clocked, from wherever you
are. \texttt{C-c C-x C-q} stops the clock and removes its line, for when you
started it by mistake. In the agenda, \texttt{I} and \texttt{O} start and stop the clock
on the line’s entry.

If you leave Tuatara while the clock runs, it asks whether to stop it. If
you say no, the clock is found again at the next start and goes on counting.

To add the time up, run \texttt{M-x org-\codebreak{}clock-\codebreak{}report}. It writes a table of the
time spent on each heading, with a total, for the subtree the cursor is in
or for the whole file when the cursor is above the first heading. Run it
again later, or press \texttt{C-c C-c} on the table’s first line, to bring it up to
date.

\begin{Verbatim}[breaklines,breakautoindent=false,breakindent=1em,breakafter=/-.\,;:=,commandchars=\\\{\}]
#+BEGIN: clocktable :scope file :maxlevel 2
| Headline          |   Time |
|-------------------+--------|
| *Total time*      | *1:05* |
|-------------------+--------|
| Clean the gutters |   0:20 |
| Water the lemon   |   0:45 |
#+END:
\end{Verbatim}

\begin{xltabular}{\textwidth}{@{}l l l >{\raggedright\arraybackslash\hsize=1.000\hsize}X@{}}
\toprule
\textbf{Modal} & \textbf{Non-modal} & \textbf{Command} & \textbf{What it does}\\
\midrule
\endhead
none & \texttt{C-c C-x C-i} & \texttt{org-\codebreak{}clock-in} & Clock in on the task point is under (C-c C-x C-i): CLOCK: and the time go into its :LOGBOOK: drawer, made when it has none, after its planning line and properties; a clock running is clocked out first.\\
none & \texttt{C-x a o}, \texttt{C-c C-x C-o} & \texttt{org-\codebreak{}clock-\codebreak{}out} & Clock out (C-c C-x C-o, from any buffer C-x a o): the clock’s line finished with the time it ended and how long it ran; with none running, of the one left open in the agenda files (ORG-CLOCK-RUNNING).\\
none & \texttt{C-x a j}, \texttt{C-c C-x C-j} & \texttt{org-\codebreak{}clock-\codebreak{}goto} & Go to the task the clock runs on (C-c C-x C-j, from any buffer C-x a j): its file opened again when its buffer was killed; with none running, one left open in the agenda files, as after a restart, runs again (ORG-CLOCK-RUNNING).\\
none & \texttt{C-c C-x C-q} & \texttt{org-\codebreak{}clock-\codebreak{}cancel} & Stop the clock running and take its line away (C-c C-x C-q): its :LOGBOOK: drawer too, when nothing else is in it; saved when that is its buffer’s only change, as C-x a o does.\\
none & \texttt{M-x org-\codebreak{}clock-\codebreak{}report} & \texttt{org-\codebreak{}clock-\codebreak{}report} & Make the clock table point is in again, or put a new one in: its scope the subtree point is in, or the file before the first headline, its depth two levels.\\
\bottomrule
\end{xltabular}

\begin{xltabular}{\textwidth}{@{}l l >{\raggedright\arraybackslash\hsize=1.000\hsize}X@{}}
\toprule
\textbf{Setting} & \textbf{Standard} & \textbf{What it changes}\\
\midrule
\endhead
\texttt{*org-\codebreak{}clock-\codebreak{}into-\codebreak{}drawer*} & \texttt{t} & True for clock lines to go into a :LOGBOOK: drawer of their headline; NIL for them to go right under it.\\
\bottomrule
\end{xltabular}

\begin{xltabular}{\textwidth}{@{}l >{\raggedright\arraybackslash\hsize=1.000\hsize}X@{}}
\toprule
\textbf{Face} & \textbf{What it paints}\\
\midrule
\endhead
\texttt{:org-\codebreak{}mode-\codebreak{}line-\codebreak{}clock} & The clock running on an Org task, in the mode line: [0:22] (Task).\\
\bottomrule
\end{xltabular}

\subsection{Tables and formulas}
\label{sec:375}
A table in Org is made of lines that begin with a bar, with a bar between
each field. You type it roughly and Tuatara lines it up.

\begin{Verbatim}[breaklines,breakautoindent=false,breakindent=1em,breakafter=/-.\,;:=,commandchars=\\\{\}]
| Item  | Cost |
|-------+------|
| Rent  |  450 |
| Food  |  180 |
| Power |   95 |
\end{Verbatim}

\subsubsection{Making a table}
\label{sec:373}
Type a bar, the first field, a bar, the next field, and so on. Then press
\texttt{TAB}. The table is aligned, each column as wide as its widest field, and
the cursor goes to the next field. At the end of a row \texttt{TAB} starts a new
row. So filling in a table is typing a field, \texttt{TAB}, a field, \texttt{TAB}. Numbers
are set to the right of their column and words to the left.

\begin{xltabular}{\textwidth}{@{}l >{\raggedright\arraybackslash\hsize=0.499\hsize}X >{\raggedright\arraybackslash\hsize=1.501\hsize}X@{}}
\toprule
\textbf{Modal} & \textbf{Non-modal} & \textbf{What it does in a table}\\
\midrule
\endhead
none & \texttt{TAB} / \texttt{S-TAB} & Aligns, and goes to the next or the previous field\\
none & \texttt{RET} while typing & Goes to the same field of the row below\\
none & \texttt{C-c C-c} & Aligns\\
none & \texttt{C-c -} & Puts a rule below this row\\
none & \texttt{M-j} / \texttt{M-k} & Moves the row down or up\\
none & \texttt{M-h} / \texttt{M-l} & Moves the column left or right\\
none & \texttt{M-L} / \texttt{M-H} & Puts in an empty column, or deletes this one\\
none & \texttt{M-S-<down>} / \texttt{M-S-<up>} & Puts in an empty row above, or deletes this one\\
\bottomrule
\end{xltabular}

These are the keys that move headings, doing the matching thing to rows and
columns when the cursor is in a table. The rows above the first rule are the
table’s header.

\subsubsection{Formulas}
\label{sec:374}
A table can do sums. Its formulas are written on a line right under it that
begins \texttt{\#+TBLFM:}. A formula names a place in the table, an equals sign, and
how to work the value out.

\begin{Verbatim}[breaklines,breakautoindent=false,breakindent=1em,breakafter=/-.\,;:=,commandchars=\\\{\}]
| Crop     | Kg | Price | Value |
|----------+----+-------+-------|
| Tomatoes | 12 |  4.50 |   54. |
| Beans    |  3 |     6 |    18 |
| Total    |    |       |   72. |
#+TBLFM: $4=$2*$3::@>$4=vsum(@2..@-1)
\end{Verbatim}

There are two formulas here, with \texttt{::} between them. The first says that
column four is column two times column three, in every row below the rule.
The second says that the last row of column four is the sum of the rows
above it.

\begin{longtable}{@{}l l@{}}
\toprule
\textbf{In a formula} & \textbf{Means}\\
\midrule
\endhead
\texttt{\$2} & The second column\\
\texttt{@3} & The third row, with rules not counted\\
\texttt{@>}, \texttt{\$>} & The last row, the last column\\
\texttt{@-1} & The row above\\
\texttt{@2..@-1} & A range, here from the second row to the row above\\
\texttt{vsum}, \texttt{vmean} & The sum and the mean of a range\\
\texttt{vmax}, \texttt{vmin}, \texttt{vcount} & The largest, the smallest, and how many\\
\texttt{;\%.2f} after the formula & Write the answer with two decimals\\
\bottomrule
\end{longtable}

The numbers do not change by themselves when you change a field. Press
\texttt{C-c *} anywhere in the table to work the formulas out again, or \texttt{C-c C-c}
on the \texttt{\#+TBLFM:} line.

You do not have to write that line by hand. \texttt{C-c =} asks for the formula of
the column the cursor is in and adds it to the line for you. With \texttt{C-u}
first, it asks for a formula for the one field.

A field that cannot be worked out shows \texttt{\#ERROR}. Usually a field the
formula reads holds a word where a number should be, or the formula names a
column the table does not have.

\begin{try}
\textbf{Try it, with a small budget}
\begin{enumerate}
\item Open \texttt{\textasciitilde{}/budget.\codebreak{}org} with \texttt{C-x C-f}. Press \texttt{i} and type \texttt{|Item|\codebreak{}Cost}, then
   press \texttt{TAB}. The line becomes a tidy row and a second row starts.
\item Type \texttt{Rent}, \texttt{TAB}, \texttt{450}, \texttt{TAB}, \texttt{Food}, \texttt{TAB}, \texttt{180}, \texttt{TAB}, \texttt{Power},
   \texttt{TAB}, \texttt{95}, \texttt{TAB}, \texttt{Total}. Press \texttt{C-g}.
\item Move to the row \texttt{Item} and press \texttt{C-c -}. A rule is drawn under it.
\item Move to the row \texttt{Total} and press \texttt{o} to open a line under the table.
   Type \texttt{\#+TBLFM: @>\$2=\codebreak{}vsum(@\codebreak{}2..@-1)} and press \texttt{C-g}.
\item Press \texttt{C-c C-c} on that line. The field beside \texttt{Total} shows 725.
\item Move to the \texttt{180} of the Food row and change it to \texttt{200}. Press \texttt{C-c *}.
   The total becomes 745. Save with \texttt{C-x C-s}.
\end{enumerate}
\end{try}

\begin{xltabular}{\textwidth}{@{}l >{\raggedright\arraybackslash\hsize=0.711\hsize}X l >{\raggedright\arraybackslash\hsize=1.289\hsize}X@{}}
\toprule
\textbf{Modal} & \textbf{Non-modal} & \textbf{Command} & \textbf{What it does}\\
\midrule
\endhead
none & \texttt{M-x org-\codebreak{}table-\codebreak{}align} & \texttt{org-\codebreak{}table-\codebreak{}align} & Align the Org table point is in: each column as wide as its widest field, numbers to the right, rules as wide as the table.\\
none & \texttt{M-x org-\codebreak{}table-\codebreak{}insert-\codebreak{}hline} & \texttt{org-\codebreak{}table-\codebreak{}insert-\codebreak{}hline} & Put a rule below the row point is in, in an Org table (C-c -).\\
none & \texttt{M-S-<down>} & \texttt{org-\codebreak{}shiftmetadown} & In an Org table, put an empty row above this one (M-S-<down>).\\
none & \texttt{M-S-<up>} & \texttt{org-\codebreak{}shiftmetaup} & In an Org table, delete the row point is in (M-S-<up>).\\
none & \texttt{C-c *} & \texttt{org-\codebreak{}table-\codebreak{}recalculate} & Work out the formulas of the Org table point is in, its \#+TBLFM: line (C-c *): column formulas (\$3=\$1+\$2) below the first rule (in every row when it has none), field formulas (@4\$2=vsum(@2..@3)) and range formulas (@2\$2..@3\$3=\$1*10), worked out with the table arithmetic.\\
none & \texttt{C-c =} & \texttt{org-\codebreak{}table-\codebreak{}eval-\codebreak{}formula} & Ask for the formula of the column point is in, the one it has offered, and work it out (C-c =); with C-u, the formula of the field point is in.\\
\bottomrule
\end{xltabular}

\begin{xltabular}{\textwidth}{@{}l l >{\raggedright\arraybackslash\hsize=1.000\hsize}X@{}}
\toprule
\textbf{Setting} & \textbf{Standard} & \textbf{What it changes}\\
\midrule
\endhead
\texttt{*org-\codebreak{}formula-\codebreak{}depth-\codebreak{}limit*} & \texttt{200} & The most a formula’s parts may be nested, within brackets, signs and powers: a formula nested deeper is refused.\\
\bottomrule
\end{xltabular}

\begin{longtable}{@{}l l@{}}
\toprule
\textbf{Face} & \textbf{What it paints}\\
\midrule
\endhead
\texttt{:org-\codebreak{}table} & Org’s tables.\\
\texttt{:org-\codebreak{}formula} & Org’s table formulas (\#+TBLFM:).\\
\bottomrule
\end{longtable}

\subsection{Links}
\label{sec:378}
A link leads from one place to another: to a heading further down, to
another file, to a web page. It is written in double square brackets, with
the place it leads to first and the words to show second:

\begin{Verbatim}[breaklines,breakautoindent=false,breakindent=1em,breakafter=/-.\,;:=,commandchars=\\\{\}]
[[file:holiday.org][the holiday plan]]
\end{Verbatim}

Tuatara shows only the words, in the colour of a link. \texttt{C-c C-x L} shows
every link of the buffer as it is written, and the same key hides them again.

\begin{longtable}{@{}l l@{}}
\toprule
\textbf{Written} & \textbf{Leads to}\\
\midrule
\endhead
\texttt{[[*Vegetables]]} & The heading with that title, in this file\\
\texttt{[[shed key]]} & The place marked \texttt{<<shed key>>} in this file\\
\texttt{[[\#beds]]} & The heading whose \texttt{CUSTOM\_ID} property is \texttt{beds}\\
\texttt{[[file:\codebreak{}chores.\codebreak{}org]]} & Another file\\
\texttt{[[file:\codebreak{}chores.\codebreak{}org::*Mow the lawn]]} & A heading in another file\\
\texttt{[[https://\codebreak{}orgmode.\codebreak{}org][Org]]} & A web page\\
\bottomrule
\end{longtable}

\subsubsection{Making and following links}
\label{sec:377}
You seldom write a link by hand. The usual way has two steps. Go to the
place you want to link to and press \texttt{C-x a l}, which remembers a link to it.
Then go to where the link should be and press \texttt{C-c C-l}. It offers the links
you have remembered, the latest first, and then asks for the words to show.
If you have text selected, that text becomes the words. With \texttt{C-u} first,
\texttt{C-c C-l} asks for a file to link to. On a link that is already there, it
lets you change it.

\texttt{C-x a l} works in any file. Under an Org heading it remembers the heading.
In a file that is not Org it remembers the line.

To follow a link, press \texttt{RET} on it in NORMAL state, or \texttt{C-c C-o} in either
state. A link to a web page or a mail address opens in your desktop’s
program, and Tuatara asks before it does that. If you would rather \texttt{RET} did
not follow links, set \texttt{*org-\codebreak{}return-\codebreak{}follows-\codebreak{}link*} to \texttt{nil}.

Two more commands write a link that stands for a set of notes, the files
with formatted names of chapter 9 whose names match a pattern. Following
such a link lists the matching notes so that you can open one. \texttt{M-x
match-\codebreak{}formatted-\codebreak{}files-\codebreak{}link} writes Tuatara’s own kind, and \texttt{M-x
match-\codebreak{}denote-\codebreak{}files-\codebreak{}link} writes the kind that the Emacs package Denote
understands as well.

\begin{xltabular}{\textwidth}{@{}l >{\raggedright\arraybackslash\hsize=0.585\hsize}X l >{\raggedright\arraybackslash\hsize=1.415\hsize}X@{}}
\toprule
\textbf{Modal} & \textbf{Non-modal} & \textbf{Command} & \textbf{What it does}\\
\midrule
\endhead
none & \texttt{C-x a l} & \texttt{org-\codebreak{}store-\codebreak{}link} & Keep a link to where point is, for org-insert-link (M-x org-insert-link) to offer: under an Org headline, to it (by its CUSTOM\_ID, else its title); in one of your notes, to the note (denote:ID); in another file, to the line point is on.\\
none & \texttt{C-c C-l} & \texttt{org-\codebreak{}insert-\codebreak{}link} & Insert an Org link (C-c C-l): it asks for the link, offering those org-store-link (C-x a l) kept and the kinds of link (file asks for a file, denote for one of your notes), then for its label, one offered selected (the stored link’s, the note’s title) to take or type over; the text selected becomes the label instead.\\
none & \texttt{C-c C-o} & \texttt{org-\codebreak{}open-\codebreak{}at-\codebreak{}point} & Follow the link point is on (C-c C-o): to a headline, a custom ID or a target here, a file, an ID, a note by its identifier, or a web or mail address in the desktop’s program, after asking.\\
none & \texttt{C-c C-x L} & \texttt{org-\codebreak{}toggle-\codebreak{}link-\codebreak{}display} & Show this Org buffer’s links as written, \hyperref[org376]{label}, and its entities as written, \texttt{\textbackslash{}dollar\{\}}, or the links as their labels and the entities as their characters again (C-c C-x L).\\
none & \texttt{C-x n l} & \texttt{match-\codebreak{}formatted-\codebreak{}files-\codebreak{}link} & Insert a link to your files with formatted names whose names match a regular expression, asked for, the files it matches listed as it is typed: LABEL.\\
none & \texttt{M-x match-\codebreak{}denote-\codebreak{}files-\codebreak{}link} & \texttt{match-\codebreak{}denote-\codebreak{}files-\codebreak{}link} & Insert a link to your files with formatted names whose names match what is asked for, the files it matches listed as it is typed, as a denote: link, which other programs reading those links follow too, written LABEL; in Markdown [LABEL](denote:query-filenames:TEXT) and in text files <denote:query-filenames:TEXT>.\\
\bottomrule
\end{xltabular}

\begin{xltabular}{\textwidth}{@{}l l >{\raggedright\arraybackslash\hsize=1.000\hsize}X@{}}
\toprule
\textbf{Setting} & \textbf{Standard} & \textbf{What it changes}\\
\midrule
\endhead
\texttt{*org-\codebreak{}link-\codebreak{}descriptive*} & \texttt{t} & True for Org links to be shown as their labels, \hyperref[org376]{label} as label and \hyperref[org376]{target} as target, and entities as their characters, \texttt{\textbackslash{}dollar\{\}} as \$; NIL to show them as written.\\
\texttt{*org-\codebreak{}return-\codebreak{}follows-\codebreak{}link*} & \texttt{t} & True for RET on an Org link to follow it; M-x org-open-at-point follows one always.\\
\bottomrule
\end{xltabular}

\begin{longtable}{@{}l l@{}}
\toprule
\textbf{Face} & \textbf{What it paints}\\
\midrule
\endhead
\texttt{:org-\codebreak{}link} & Links in Org.\\
\texttt{:org-\codebreak{}target} & Org’s targets, \label{org376}.\\
\bottomrule
\end{longtable}

\subsection{Blocks and code}
\label{sec:382}
A block is a stretch of text set apart between a line that begins \texttt{\#+begin\_}
and one that begins \texttt{\#+end\_}. The word after the underscore says what kind of
block it is.

\begin{longtable}{@{}l l@{}}
\toprule
\textbf{Block} & \textbf{Is for}\\
\midrule
\endhead
\texttt{src} & Code, coloured as its language is\\
\texttt{example} & Text shown exactly as typed, such as what a program printed\\
\texttt{quote} & Words quoted from someone\\
\texttt{verse} & Poetry, where the line breaks matter\\
\texttt{center} & Lines set in the middle\\
\bottomrule
\end{longtable}

You do not type the two lines yourself. Type \texttt{<s} and press \texttt{TAB}, and a
source block is written around the cursor with a list of languages to pick
from. \texttt{<e}, \texttt{<q}, \texttt{<v} and \texttt{<c} do the same for the other four. These are
snippets, which chapter 13 describes.

\begin{Verbatim}[breaklines,breakautoindent=false,breakindent=1em,breakafter=/-.\,;:=,commandchars=\\\{\}]
#+begin_src shell
echo "Water the beds at dawn"
#+end_src
\end{Verbatim}

The code in a source block is coloured as it would be in a file of that
language. Org’s name for a language is sometimes different from Tuatara’s.
\texttt{*org-\codebreak{}src-\codebreak{}languages*} maps one to the other, and you can add your own.

\subsubsection{Writing the code out to files}
\label{sec:379}
A note can hold the code of a program or of your configuration, with the
explanation around it, and write the code out to the files that need it.
Give a source block the words \texttt{:tangle} and a file name on its first line:

\begin{Verbatim}[breaklines,breakautoindent=false,breakindent=1em,breakafter=/-.\,;:=,commandchars=\\\{\}]
#+begin_src shell :tangle water.sh
echo "Water the beds at dawn"
#+end_src
\end{Verbatim}

\texttt{C-c C-v t} then writes every such block to its file. Blocks that name the
same file go into it in order. Before anything is written, Tuatara lists the
files and says which are new and which would change, and asks. This suits
settings files well: the note says why each setting is there, and the file
the program reads is made from it.

\subsubsection{A table of contents}
\label{sec:380}
Tag a heading \texttt{TOC} and Tuatara keeps a table of contents under it, a list
of links to the note’s headings. It is made again each time you save. The
tag can say more: \texttt{TOC\_3} lists headings three levels deep. \texttt{M-x
toc-\codebreak{}org-\codebreak{}insert-\codebreak{}toc} makes it at any other time.

\subsubsection{Blocks that write themselves}
\label{sec:381}
A dynamic block is a block whose contents Tuatara writes. The clock table of
the section on clocking is one. \texttt{C-c C-x x} puts in a new one and asks which
kind. Apart from the clock table, the kinds list or include your files with
formatted names: links to those that match a pattern, links to the notes
that link to this one, or the text of the matching files themselves. \texttt{C-c
C-c} on a block’s first line, or \texttt{C-c C-x C-u} anywhere in it, writes it
again from how things stand now. \texttt{M-x org-\codebreak{}update-\codebreak{}all-\codebreak{}dblocks} does that for
every such block in the buffer.

{\small
\begin{xltabular}{\textwidth}{@{}l >{\raggedright\arraybackslash\hsize=0.880\hsize}X l >{\raggedright\arraybackslash\hsize=1.120\hsize}X@{}}
\toprule
\textbf{Modal} & \textbf{Non-modal} & \textbf{Command} & \textbf{What it does}\\
\midrule
\endhead
none & \texttt{C-c C-v t}, \texttt{C-c C-v C-t} & \texttt{tangle-\codebreak{}code-\codebreak{}files} & Write the code files of this Org file’s source blocks (C-c C-v t), each block with :tangle FILE (on its \#+begin\_src line, in a subtree’s property :header-args: or :header-args:LANGUAGE:, or in \#+PROPERTY: header-args for the file) goes to FILE.\\
none & \texttt{M-x toc-\codebreak{}org-\codebreak{}insert-\codebreak{}toc} & \texttt{toc-\codebreak{}org-\codebreak{}insert-\codebreak{}toc} & Make the table of contents of the current Org file again, under its first headline tagged :TOC: (as saving does, with \textbf{org-toc-on-save}), a list of links to its headlines, two levels deep unless the tag says how deep (:TOC\_3:), to GitHub’s anchors unless it says :TOC\_2\_org: for their titles; :TOC:QUOTE: puts it in a quote block.\\
none & \texttt{C-c C-x x} & \texttt{org-\codebreak{}dynamic-\codebreak{}block-\codebreak{}insert-\codebreak{}dblock} & Put a new dynamic block in (C-c C-x x): its kind asked for among those that can be (a clock table, links to files with formatted names\ldots{}), then what that kind asks.\\
none & \texttt{C-c C-x C-u} & \texttt{org-\codebreak{}dblock-\codebreak{}update} & Make the dynamic block point is in again (C-c C-x C-u): the text between \#+BEGIN: and \#+END: written anew by the writer \#+BEGIN: names, in the background, the editor free meanwhile.\\
none & \texttt{M-x org-\codebreak{}update-\codebreak{}all-\codebreak{}dblocks} & \texttt{org-\codebreak{}update-\codebreak{}all-\codebreak{}dblocks} & Make every dynamic block of the current Org buffer again: each written from the text as it is now, in the background, the editor free meanwhile, and put in together, in one step to undo.\\
none & \texttt{M-x insert-\codebreak{}formatted-\codebreak{}links-\codebreak{}block} & \texttt{insert-\codebreak{}formatted-\codebreak{}links-\codebreak{}block} & Put in a dynamic block of links to the files with formatted names whose names a regular expression matches (a block also written denote-links): asked for, the files it matches listed as it is typed; the links labelled with the files’ signatures and titles.\\
none & \texttt{M-x insert-\codebreak{}formatted-\codebreak{}missing-\codebreak{}links-\codebreak{}block} & \texttt{insert-\codebreak{}formatted-\codebreak{}missing-\codebreak{}links-\codebreak{}block} & Put in a dynamic block of links to the files with formatted names whose names a regular expression matches and that this file does not link to yet, (a block also written denote-missing-links): the regular expression asked for, the files it matches listed as it is typed.\\
none & \texttt{M-x insert-\codebreak{}formatted-\codebreak{}backlinks-\codebreak{}block} & \texttt{insert-\codebreak{}formatted-\codebreak{}backlinks-\codebreak{}block} & Put in a dynamic block of links to the files with formatted names that link to this one (a block also written denote-backlinks).\\
none & \texttt{M-x insert-\codebreak{}formatted-\codebreak{}files-\codebreak{}block} & \texttt{insert-\codebreak{}formatted-\codebreak{}files-\codebreak{}block} & Put in a dynamic block holding the text of the files with formatted names whose names a regular expression matches (a block also written denote-files): the regular expression and the order of the files asked for.\\
none & \texttt{M-x insert-\codebreak{}formatted-\codebreak{}files-\codebreak{}as-\codebreak{}headings-\codebreak{}block} & \texttt{insert-\codebreak{}formatted-\codebreak{}files-\codebreak{}as-\codebreak{}headings-\codebreak{}block} & Put in a dynamic block holding the text of the Org files with formatted names whose names a regular expression matches, each under a headline of its title and tags, its own headlines a level down (a block also written denote-files-as-headings): the regular expression and the order of the files asked for.\\
none & \texttt{M-x insert-\codebreak{}formatted-\codebreak{}sequence-\codebreak{}block} & \texttt{insert-\codebreak{}formatted-\codebreak{}sequence-\codebreak{}block} & Put in a dynamic block listing a file with a sequence for its signature and those under it (a block also written denote-sequence): the file asked for, and how many levels to list, its own the first.\\
\bottomrule
\end{xltabular}
}

\begin{xltabular}{\textwidth}{@{}l >{\raggedright\arraybackslash\hsize=0.571\hsize}X >{\raggedright\arraybackslash\hsize=1.429\hsize}X@{}}
\toprule
\textbf{Setting} & \textbf{Standard} & \textbf{What it changes}\\
\midrule
\endhead
\texttt{*org-\codebreak{}src-\codebreak{}languages*} & \texttt{(("emacs-\codebreak{}lisp" . "elisp") ("C++" . "cpp"...} & The language each name a source block is given (\#+begin\_src emacs-lisp) means, where Org’s name is not the language’s own: each (NAME .\\
\texttt{*code-\codebreak{}files-\codebreak{}make-\codebreak{}folders*} & \texttt{nil} & NIL for tangle-code-files to make the folders a file goes in only for blocks with :mkdirp yes, so that an Org file writes the same everywhere; true to make them for every block, as :mkdirp yes does (a block’s :mkdirp no still says not to).\\
\texttt{*org-\codebreak{}toc-\codebreak{}on-\codebreak{}save*} & \texttt{t} & True to make the table of contents of an Org file again each time it is saved, under its first headline tagged :TOC:; NIL to leave it to M-x toc-org-insert-toc.\\
\texttt{*org-\codebreak{}toc-\codebreak{}max-\codebreak{}depth*} & \texttt{2} & How deep the headlines of a table of contents go when its tag says nothing (:TOC:), 2 for those of one and two stars.\\
\texttt{*org-\codebreak{}toc-\codebreak{}style*} & \texttt{"gh"} & How a table of contents links to its headlines when its tag says nothing, “gh” for GitHub’s anchors (\#a-headline, which GitHub shows a README’s headlines by), “org” for their titles.\\
\bottomrule
\end{xltabular}

\begin{xltabular}{\textwidth}{@{}l >{\raggedright\arraybackslash\hsize=1.000\hsize}X@{}}
\toprule
\textbf{Face} & \textbf{What it paints}\\
\midrule
\endhead
\texttt{:org-\codebreak{}block} & The contents of Org’s blocks (\#+begin\_src\ldots{}\\
\texttt{:org-\codebreak{}block-\codebreak{}begin-\codebreak{}line} & The first line of an Org block, \#+begin\_src and its arguments.\\
\texttt{:org-\codebreak{}block-\codebreak{}end-\codebreak{}line} & The last line of an Org block, \#+end\_src.\\
\texttt{:org-\codebreak{}quote} & The contents of Org’s quote blocks.\\
\texttt{:org-\codebreak{}verse} & The contents of Org’s verse blocks.\\
\bottomrule
\end{xltabular}

\subsection{Exporting}
\label{sec:391}
A note is for you. A letter, a report or a set of minutes is for someone
else, who wants a PDF or a Word document. Exporting makes one from the note.
The headings become numbered sections, the lists become lists, the tables
become tables, and the marks for bold and italic become bold and italic.

Tuatara makes a PDF by way of \LaTeX{}, a typesetting system that produces
handsome pages. You do not need to know \LaTeX{}, but it must be installed: the
program \texttt{xelatex} or \texttt{lualatex}, from the collection called \TeX{} Live. \texttt{M-x
check-\codebreak{}health} tells you whether Tuatara found one. A Word document needs
nothing else installed.

The export commands have no keys. Run them with \texttt{M-x}. The command you ran
last is offered first there, so exporting again after a change is \texttt{M-x RET}.

\subsubsection{What goes at the top of the note}
\label{sec:383}
A few lines at the top say what the document is.

\begin{Verbatim}[breaklines,breakautoindent=false,breakindent=1em,breakafter=/-.\,;:=,commandchars=\\\{\}]
#+title: The fence between our sections
#+author: Aroha Ngata
#+date: 14 October 2026
#+options: toc:nil num:nil
\end{Verbatim}

The title, author and date are printed at the head of the document. Without
an \texttt{\#+author:} line, the name in \texttt{*user-\codebreak{}full-\codebreak{}name*} is used. The \texttt{\#+options:}
line here leaves out the table of contents and the numbers before headings,
which suits a letter. Leave the line out and a report gets both. A heading
tagged \texttt{noexport} is left out of the document with everything under it,
which is handy for notes to yourself.

\subsubsection{Making a PDF}
\label{sec:384}
\texttt{M-x org-\codebreak{}export-\codebreak{}latex-\codebreak{}pdf-\codebreak{}enhanced-\codebreak{}buffer} makes the PDF and shows it beside
your note. The work is done in the background, so you can go on writing. The
mode line shows that a job is running and the echo area says when the file
is written and how long it took. The PDF is put beside the note with the
same name, \texttt{letter.\codebreak{}pdf} for \texttt{letter.\codebreak{}org}. If you change the note and export
again, the PDF on the screen is read again at the page you were on.

Two more commands do the same and differ only in what happens afterwards.
\texttt{M-x org-\codebreak{}export-\codebreak{}latex-\codebreak{}pdf-\codebreak{}enhanced} makes the PDF and shows nothing. \texttt{M-x
org-\codebreak{}export-\codebreak{}latex-\codebreak{}pdf-\codebreak{}enhanced-\codebreak{}xdg} opens it in your desktop’s PDF reader.

The export takes the note as the buffer has it, saved or not. To export only
part of a note, select the lines first.

If \LaTeX{} fails, the echo area names the line of your note where it went
wrong, and a buffer called \texttt{*LaTeX errors*} lists each error. \texttt{RET} on one
shows that line of the note. An export is a background job like any other,
so chapter 11 applies to it. \texttt{C-x j} lists what is running and \texttt{C-x J}
stops it.

\subsubsection{Making a Word document}
\label{sec:385}
\texttt{M-x org-\codebreak{}export-\codebreak{}docx} writes the note as a Word document beside it, and
\texttt{M-x org-\codebreak{}export-\codebreak{}docx-\codebreak{}xdg} also opens it in your desktop’s word processor. The
document is laid out as the enhanced PDF is, with the same numbering and the
same tables, so you can send whichever the other person prefers.

\begin{try}
\textbf{Try it, with a letter}
\begin{enumerate}
\item Open \texttt{\textasciitilde{}/letter.\codebreak{}org} with \texttt{C-x C-f}. Press \texttt{i} and type the four lines
   shown under “What goes at the top of the note”, with your own name.
\item Leave an empty line and type the letter:
\begin{Verbatim}[breaklines,breakautoindent=false,breakindent=1em,breakafter=/-.\,;:=,commandchars=\\\{\}]
   Dear Mr Tane,

   Thank you for offering to share the cost of the new fence. The quote
   from Motueka Fencing is *$1,840* for the whole length.

   - Posts and rails: $1,100
   - Labour: $740

   Yours sincerely,

   Aroha Ngata
\end{Verbatim}
\item Press \texttt{C-g} and save with \texttt{C-x C-s}.
\item Press \texttt{M-x}, type \texttt{org-\codebreak{}export-\codebreak{}latex-\codebreak{}pdf-\codebreak{}enhanced-\codebreak{}buffer} and press \texttt{RET}.
   After a few seconds the letter is shown as a PDF at the right, with the
   title at the top, the sum in bold and the two costs as a list. On a narrow
   screen the PDF takes the whole window instead. Press \texttt{C-x k} to close it
   and return to the note before the next step.
\item Press \texttt{M-x}, type \texttt{org-\codebreak{}export-\codebreak{}docx} and press \texttt{RET}. The echo area says
   that \texttt{\textasciitilde{}/letter.\codebreak{}docx} was written.
\end{enumerate}
\end{try}

\subsubsection{Two kinds of export}
\label{sec:386}
The commands above all have “enhanced” in their name or, for Word, behave
that way. There is a second, plainer family without the word.

\begin{xltabular}{\textwidth}{@{}l >{\raggedright\arraybackslash\hsize=0.720\hsize}X >{\raggedright\arraybackslash\hsize=1.280\hsize}X@{}}
\toprule
\textbf{Family} & \textbf{Commands begin} & \textbf{What you get}\\
\midrule
\endhead
Plain & \texttt{org-\codebreak{}export-\codebreak{}latex-\codebreak{}pdf} & The PDF that Emacs’s Org would make from the same note\\
Enhanced & \texttt{org-\codebreak{}export-\codebreak{}latex-\codebreak{}pdf-\codebreak{}enhanced} & Tuatara’s own layout, with the additions below\\
\bottomrule
\end{xltabular}

Use the plain export for a note that must come out exactly as it did in
Emacs. Use the enhanced export for everything else. It adds four things:
tables that are fitted to the page, layout settings written in plain words,
templates that give a document your own look, and a way to black out names.

Each family also has a command that stops at the \LaTeX{} file and makes no
PDF, \texttt{M-x org-\codebreak{}export-\codebreak{}latex} and \texttt{M-x org-\codebreak{}export-\codebreak{}latex-\codebreak{}enhanced}, for people
who want to work on the \LaTeX{} themselves. If you have a \LaTeX{} file of your
own, \texttt{M-x export-\codebreak{}latex-\codebreak{}compile-\codebreak{}pdf} and its two companions make a PDF from
it in the same three ways.

\subsubsection{Layout in plain words}
\label{sec:387}
In the enhanced export a wide table is made to fit. Columns of sentences
wrap, numbers go to the right, and the header row is repeated when the table
runs on to another page. Most tables need nothing from you.

When you do want to change how something is set, you say so in a drawer
called \texttt{:EXPORT:}, in words. Put it right under a heading to affect the
heading and all under it, or just above a table, list, block or picture to
affect that.

\begin{Verbatim}[breaklines,breakautoindent=false,breakindent=1em,breakafter=/-.\,;:=,commandchars=\\\{\}]
* Costs
:EXPORT:
new page: before
spacing: 1.5
:END:
\end{Verbatim}

This starts the section on a new page and sets its lines one and a half
apart. The drawer itself is never printed. You do not need to remember the
settings. \texttt{M-x org-\codebreak{}insert-\codebreak{}export-\codebreak{}drawer} writes a drawer for whatever the
cursor is on, holding every setting that fits it, each set to \texttt{auto} for you
to change. \texttt{M-x describe-\codebreak{}export-\codebreak{}settings} explains them all, with examples.
A word in a drawer that is not a setting is shown in the colour of an error
as you type it.

\subsubsection{Your own look}
\label{sec:388}
A template gives a document its look: the fonts, the margins, a header with
a logo. A note chooses one with a line at the top, \texttt{\#+template:} and the
template’s name. Templates are kept in your configuration folder.

The simplest start is \texttt{M-x export-\codebreak{}latex-\codebreak{}template-\codebreak{}copy}. It copies one of the
built-in templates to your folder and opens it for you to change. If someone
gives you a Word or \LaTeX{} document and asks for reports that look like it,
\texttt{M-x export-\codebreak{}template-\codebreak{}from-\codebreak{}docx} and \texttt{M-x export-\codebreak{}template-\codebreak{}from-\codebreak{}latex} make a
template from that document.

A note template is the other half: the lines a new note of some kind starts
with, such as the \texttt{\#+template:} line and the headings every report has.
\texttt{C-x i} reads one into the buffer.

The fonts of the enhanced export are \texttt{*latex-\codebreak{}main-\codebreak{}font*} for the text,
\texttt{*latex-\codebreak{}sans-\codebreak{}font*} and \texttt{*latex-\codebreak{}mono-\codebreak{}font*}. Each is the name of a font
family installed on your machine. Left alone, Tuatara picks an installed
font for each.

\subsubsection{Hiding names}
\label{sec:389}
Sometimes a document must go out with a name blacked out. Select the name,
or put the cursor on it, press \texttt{C-c m} and choose \texttt{R}. The name is wrapped
in a mark and shown in its own colour. In the PDF and in the Word document
a black box stands where the name was, and the name is nowhere in the file.
\texttt{C-c M} takes the mark away again.

A box is of little use if the same name stands in plain text two pages
later. \texttt{M-x export-\codebreak{}check-\codebreak{}redactions} looks for every hidden text that also
appears unhidden and lists the places. The export runs the same check and
asks before it goes on. Only the enhanced export hides text. The plain
export refuses a note that asks for it.

\subsubsection{Notes from other people}
\label{sec:390}
An exported note can pull in other files: pictures it shows and setup files
it names. A note that someone sent you could name files of yours that way.
So the export takes files only from the note’s own folder and the folders
in it, and says what it left out. \texttt{*export-\codebreak{}trusted-\codebreak{}folders*} names further
folders you allow. Chapter 12 has more on files you do not trust.

{\small
\begin{xltabular}{\textwidth}{@{}l >{\raggedright\arraybackslash\hsize=0.694\hsize}X l >{\raggedright\arraybackslash\hsize=1.306\hsize}X@{}}
\toprule
\textbf{Modal} & \textbf{Non-modal} & \textbf{Command} & \textbf{What it does}\\
\midrule
\endhead
none & \texttt{M-x org-\codebreak{}export-\codebreak{}latex-\codebreak{}pdf-\codebreak{}enhanced} & \texttt{org-\codebreak{}export-\codebreak{}latex-\codebreak{}pdf-\codebreak{}enhanced} & Make the Org file of the current buffer into a PDF beside it, as M-x org-export-latex-pdf does, with Tuatara’s own improvements: its \#+template, tables that fit and redaction (M-x org-export-latex-enhanced).\\
none & \texttt{M-x org-\codebreak{}export-\codebreak{}latex-\codebreak{}pdf-\codebreak{}enhanced-\codebreak{}buffer} & \texttt{org-\codebreak{}export-\codebreak{}latex-\codebreak{}pdf-\codebreak{}enhanced-\codebreak{}buffer} & Make the Org file of the current buffer into a PDF as M-x org-export-latex-pdf-enhanced does, and show it as M-x org-export-latex-pdf-buffer shows its PDF.\\
none & \texttt{M-x org-\codebreak{}export-\codebreak{}latex-\codebreak{}pdf-\codebreak{}enhanced-\codebreak{}xdg} & \texttt{org-\codebreak{}export-\codebreak{}latex-\codebreak{}pdf-\codebreak{}enhanced-\codebreak{}xdg} & Make the Org file of the current buffer into a PDF as M-x org-export-latex-pdf-enhanced does, and open it in the desktop’s program as M-x org-export-latex-pdf-xdg opens its PDF.\\
none & \texttt{M-x org-\codebreak{}export-\codebreak{}latex-\codebreak{}enhanced} & \texttt{org-\codebreak{}export-\codebreak{}latex-\codebreak{}enhanced} & Write the Org file of the current buffer as \LaTeX{} as M-x org-export-latex does, with Tuatara’s own improvements (its \#+template, tables that fit, redaction): the \LaTeX{} file M-x org-export-latex-pdf-enhanced makes its PDF from.\\
none & \texttt{M-x org-\codebreak{}export-\codebreak{}docx} & \texttt{org-\codebreak{}export-\codebreak{}docx} & Write the Org note of this buffer as a Word document (.docx) beside it, or where \#+EXPORT\_FILE\_NAME says (its .docx): read as the enhanced PDF export reads it, its :EXPORT: drawers’ settings and all, what redact and protect hide a black bar, never its text, in the file; in the background.\\
none & \texttt{M-x org-\codebreak{}export-\codebreak{}docx-\codebreak{}xdg} & \texttt{org-\codebreak{}export-\codebreak{}docx-\codebreak{}xdg} & Write the Org note of this buffer as a Word document, as M-x org-export-docx does, and open it in the desktop’s program for Word documents.\\
none & \texttt{M-x org-\codebreak{}export-\codebreak{}html} & \texttt{org-\codebreak{}export-\codebreak{}html} & Write the Org note of this buffer as a web page (.html) beside it, or where \#+EXPORT\_FILE\_NAME says (its .html), one file with its stylesheet inside, in the Ao theme’s colours, light or dark as the reader’s system prefers: read as the enhanced PDF export reads it, what redact and protect hide a black bar, never its text.\\
\bottomrule
\end{xltabular}
}

\begin{xltabular}{\textwidth}{@{}l >{\raggedright\arraybackslash\hsize=0.722\hsize}X l >{\raggedright\arraybackslash\hsize=1.278\hsize}X@{}}
\toprule
\textbf{Modal} & \textbf{Non-modal} & \textbf{Command} & \textbf{What it does}\\
\midrule
\endhead
none & \texttt{M-x org-\codebreak{}export-\codebreak{}latex-\codebreak{}pdf} & \texttt{org-\codebreak{}export-\codebreak{}latex-\codebreak{}pdf} & Make the Org file of the current buffer into a PDF beside it, in the background: the plain export, the note set as standard \LaTeX{} sets it, with no template or layout of Tuatara’s own.\\
none & \texttt{M-x org-\codebreak{}export-\codebreak{}latex-\codebreak{}pdf-\codebreak{}buffer} & \texttt{org-\codebreak{}export-\codebreak{}latex-\codebreak{}pdf-\codebreak{}buffer} & Make the Org file of the current buffer into a PDF beside it, as M-x org-export-latex-pdf does, and show the PDF: down the right side of the screen when it is wide enough (\textbf{split-width-threshold} columns), else alone on it.\\
none & \texttt{M-x org-\codebreak{}export-\codebreak{}latex-\codebreak{}pdf-\codebreak{}xdg} & \texttt{org-\codebreak{}export-\codebreak{}latex-\codebreak{}pdf-\codebreak{}xdg} & Make the Org file of the current buffer into a PDF beside it, as M-x org-export-latex-pdf does, and open the PDF in the desktop’s program for PDFs (xdg-open), or the one \textbf{open-externally-programs} gives “application/pdf”, asking first unless \textbf{open-externally-ask} is NIL.\\
none & \texttt{M-x org-\codebreak{}export-\codebreak{}latex} & \texttt{org-\codebreak{}export-\codebreak{}latex} & Write the Org file of the current buffer as \LaTeX{} beside it (or where its \#+export\_file\_name says), the plain export with no template or layout of Tuatara’s own, in the background, without making a PDF: the \LaTeX{} file M-x org-export-latex-pdf makes its PDF from.\\
none & \texttt{M-x export-\codebreak{}latex-\codebreak{}compile-\codebreak{}pdf} & \texttt{export-\codebreak{}latex-\codebreak{}compile-\codebreak{}pdf} & Make the \LaTeX{} file of the current buffer into a PDF beside it, in the background, saving it first when it has changes, by xelatex, or by lualatex when xelatex fails (\textbf{latex-engines}), or first by the one its first lines name (“\% Intended \LaTeX{} compiler: lualatex”, “\%!TEX program = lualatex”), run as often as its cross-references need.\\
none & \texttt{M-x export-\codebreak{}latex-\codebreak{}compile-\codebreak{}pdf-\codebreak{}buffer} & \texttt{export-\codebreak{}latex-\codebreak{}compile-\codebreak{}pdf-\codebreak{}buffer} & Make the \LaTeX{} file of the current buffer into a PDF beside it, as M-x export-latex-compile-pdf does, and show the PDF: down the right side of the screen when it is wide enough (\textbf{split-width-threshold} columns), else alone on it.\\
none & \texttt{M-x export-\codebreak{}latex-\codebreak{}compile-\codebreak{}pdf-\codebreak{}xdg} & \texttt{export-\codebreak{}latex-\codebreak{}compile-\codebreak{}pdf-\codebreak{}xdg} & Make the \LaTeX{} file of the current buffer into a PDF beside it, as M-x export-latex-compile-pdf does, and open the PDF in the desktop’s program for PDFs (xdg-open), or the one \textbf{open-externally-programs} gives “application/pdf”, asking first unless \textbf{open-externally-ask} is NIL.\\
\bottomrule
\end{xltabular}

\begin{xltabular}{\textwidth}{@{}l >{\raggedright\arraybackslash\hsize=0.865\hsize}X l >{\raggedright\arraybackslash\hsize=1.135\hsize}X@{}}
\toprule
\textbf{Modal} & \textbf{Non-modal} & \textbf{Command} & \textbf{What it does}\\
\midrule
\endhead
none & \texttt{M-x org-\codebreak{}insert-\codebreak{}export-\codebreak{}drawer} & \texttt{org-\codebreak{}insert-\codebreak{}export-\codebreak{}drawer} & Write an :EXPORT: drawer for the table, list, block or picture point is on, or else for the heading it is under, holding the settings that fit it, each auto to change (M-x describe-export-settings says what each takes).\\
none & \texttt{M-x describe-\codebreak{}export-\codebreak{}settings} & \texttt{describe-\codebreak{}export-\codebreak{}settings} & Show the settings an :EXPORT: drawer may hold in an Org note, what each sets and what it takes: how the enhanced export sets a heading, a table, a list, a block or a picture, in words rather than \LaTeX{}.\\
none & \texttt{M-x export-\codebreak{}latex-\codebreak{}template-\codebreak{}copy} & \texttt{export-\codebreak{}latex-\codebreak{}template-\codebreak{}copy} & Copy a built-in template of the enhanced export, asked for with completion (article, its own look, or org, the plain export’s document), into your own templates, \textasciitilde{}/.config/tuatara/latex-templates/NAME.tex, and visit it there to change: from then on \#+template: NAME takes your copy.\\
none & \texttt{M-x export-\codebreak{}template-\codebreak{}from-\codebreak{}docx} & \texttt{export-\codebreak{}template-\codebreak{}from-\codebreak{}docx} & Make a Word document given to you a template to write reports in.\\
none & \texttt{M-x export-\codebreak{}template-\codebreak{}from-\codebreak{}latex} & \texttt{export-\codebreak{}template-\codebreak{}from-\codebreak{}latex} & Make a \LaTeX{} document given to you a template to write reports in.\\
none & \texttt{C-x i} & \texttt{insert-\codebreak{}formatted-\codebreak{}note-\codebreak{}template} & Read a note template, or any file, into the current buffer at point: one of your templates (blank, and each file in \textasciitilde{}/.config/tuatara/note-templates/), \%date\% and \%author\% filled in, or with “Another file\ldots{}” any file, as it is.\\
none & \texttt{C-c m} & \texttt{org-\codebreak{}macro-\codebreak{}wrap} & In an Org note, wrap the selection, or the word at point, in a macro chosen from a list: R, \texorpdfstring{\redactedbox{nnnnnnnnnn}}{[redacted]}, hidden in the enhanced export’s PDF as a black box; P, hidden there only with \#+options: protect:t; or one the note declares with a \#+macro: line, what it stands for shown beside it.\\
none & \texttt{C-c M} & \texttt{org-\codebreak{}macro-\codebreak{}unwrap} & In an Org note, take away the macro call point is in, its text left as it was: \texorpdfstring{\redactedbox{nnnnnnnnnn}}{[redacted]} becomes John Otago.\\
none & \texttt{M-x export-\codebreak{}check-\codebreak{}redactions} & \texttt{export-\codebreak{}check-\codebreak{}redactions} & Look for the texts this Org or \LaTeX{} file hides shown plainly elsewhere in it, and list each place in \textbf{Redactions}, where l shows it.\\
\bottomrule
\end{xltabular}

\begin{xltabular}{\textwidth}{@{}l >{\raggedright\arraybackslash\hsize=0.708\hsize}X >{\raggedright\arraybackslash\hsize=1.292\hsize}X@{}}
\toprule
\textbf{Setting} & \textbf{Standard} & \textbf{What it changes}\\
\midrule
\endhead
\texttt{*latex-\codebreak{}engines*} & \texttt{("xelatex" "lualatex")} & The programs that make a PDF of \LaTeX{}, tried in this order: when one fails, the next is tried.\\
\texttt{*latex-\codebreak{}main-\codebreak{}font*} & \texttt{nil} & The font the enhanced export sets a document’s text in, by its family’s name as fontconfig knows it (“Noto Serif”); NIL for the first installed of Libertinus Serif, Noto Serif, Source Serif 4, Liberation Serif, DejaVu Serif.\\
\texttt{*latex-\codebreak{}sans-\codebreak{}font*} & \texttt{nil} & The sans-serif font of the enhanced export, as \textbf{LATEX-MAIN-FONT}; NIL for the first installed of Libertinus Sans, Noto Sans, Source Sans 3, Liberation Sans, DejaVu Sans.\\
\texttt{*latex-\codebreak{}mono-\codebreak{}font*} & \texttt{nil} & The fixed-width font the enhanced export sets code in, as \textbf{LATEX-MAIN-FONT}; NIL for your window’s font (\textbf{gui-font}) when it is installed, else the first of JetBrains Mono, Cascadia Code, Noto Sans Mono, Source Code Pro, Liberation Mono, DejaVu Sans Mono.\\
\texttt{*logos-\codebreak{}directory*} & \texttt{nil} & The folder of the pictures your templates ask a note for, the logo its \#+logo: names (LOGOS-DIRECTORY); NIL, as standard, for logos/ in the configuration directory.\\
\texttt{*export-\codebreak{}trusted-\codebreak{}folders*} & \texttt{nil} & The folders, besides a note’s own folder and the folders in it and your note templates’ folder (NOTE-TEMPLATES-DIRECTORY), that exporting the note (to PDF or to a Word document) may take files from: the setup files its \#+SETUPFILE lines name and the pictures its \href{...}{\ldots{}} links show.\\
\texttt{*latex-\codebreak{}redaction-\codebreak{}commands*} & \texttt{("R" "D" "redact")} & The \LaTeX{} commands that hide their argument in your \LaTeX{} files, \texttt{\textbackslash{}\textbackslash{}R\{John} Otago\} for one, by their names without the backslash: what M-x export-check-redactions and M-x export-latex-compile-pdf look for in a \LaTeX{} file, as hidden text to be found nowhere else.\\
\bottomrule
\end{xltabular}

\begin{xltabular}{\textwidth}{@{}l >{\raggedright\arraybackslash\hsize=1.000\hsize}X@{}}
\toprule
\textbf{Face} & \textbf{What it paints}\\
\midrule
\endhead
\texttt{:org-\codebreak{}macro} & Macro calls in Org (\{\{\{name\}\}\}).\\
\texttt{:org-\codebreak{}macro-\codebreak{}text} & The text of an Org macro call shown while its braces are hidden (\{\{\{person(Ann Taranaki)\}\}\} shows Ann Taranaki, \{\{\{kbd(C-x)\}\}\} C-x), as a link’s label is in :org-link.\\
\bottomrule
\end{xltabular}

\subsection{Where to go from here}
\label{sec:392}
Org rewards starting small. Keep one file of notes and fold it. Add a list
with checkboxes when you next go shopping. When the tasks pile up, give them
dates and open the agenda each morning. The rest of this chapter will still
be here when you need a table or a PDF.

Many people keep their Org files under version control, so that every change
to their notes is kept and can be taken back. That is the subject of the
next chapter.

\section{Git}
\label{sec:437}
Keeping the history of a folder, and working on it with other people.

Git is a program that remembers every version of the files in a folder. You
tell it when a piece of work is done, and it keeps that state for good, with a
sentence from you saying what changed. Later you can look at any earlier
state, see who changed a line and why, try an idea without risk to the work
that is finished, and share the folder with someone else without mailing files
back and forth. It is used most for programs, but it works as well for a
folder of recipes, a set of letters or a small website.

Tuatara does not replace git. It runs the git on your computer for you and
shows what git says in buffers you can read and act on. Everything starts from
one buffer, the status, which \texttt{C-x g} opens. This chapter explains the status
first, then the daily round of staging and committing, then looking at
history, branches, sharing with a remote, and the tools for finding out when
and why something changed.

The examples use a folder of recipes, \texttt{\textasciitilde{}/recipes}, with a text file for each
recipe. You need git installed, and git needs to know your name once. If you
have never used it, type these two lines in a terminal, with your own name and
address:

\begin{Verbatim}[breaklines,breakautoindent=false,breakindent=1em,breakafter=/-.\,;:=,commandchars=\\\{\}]
git config --global user.name "Aroha Example"
git config --global user.email "aroha@example.org"
\end{Verbatim}

\subsection{What git keeps}
\label{sec:395}
A few words come up on every page of this chapter. They are git’s own, and
the same in every book about it.

A repository is a folder whose history git keeps. Git stores the history in a
hidden folder called \texttt{.git} inside it. Your files stay ordinary files that any
program can open.

A commit is one saved state of the whole folder, with a message, your name and
the time. Commits form a chain, each one pointing back to the one before. Every
commit has a name of its own made of letters and digits, called its hash, such
as \texttt{3f9c2ab}. You rarely type one; you pick the commit from a list.

Staging is how you choose what goes into the next commit. Suppose you fixed
the oven temperature in one recipe and began a new recipe in the same hour.
Those are two pieces of work. You stage the fix and commit it, then stage the
new recipe and commit that. Each commit then says one thing, which makes the
history easy to read a year later. What is staged is kept in a list that git
calls the index.

A branch is a line of commits with a name. A new repository has one, usually
called \texttt{main}. You make another when you want to try something, such as
rewriting every recipe in metric measures, and keep \texttt{main} as it is until the
new work is ready.

A remote is another copy of the repository, most often on a server. You push
your commits to it and fetch the commits other people pushed.

Figure 15.1 shows the way a change travels.

\begin{figure}[H]
\caption{Figure 15.1. From your files to a commit, and on to a remote.}\label{org394}
\begin{Verbatim}[breaklines,breakautoindent=false,breakindent=1em,breakafter=/-.\,;:=,commandchars=\\\{\}]
   your files          staged             commits            remote
 ┌────────────┐     ┌──────────┐     ┌──────────────┐\discretionary{}{}{}     ┌──────────┐
 │ scones.txt │  s  │ the next │ c c │ o──o──o──o   │ P p │ a copy   │
 │ (changed)  │ ──{\tuataraFallbackTWO{}▶} │ commit   │ ──{\tuataraFallbackTWO{}▶} │         main │ ──{\tuataraFallbackTWO{}▶} │ elsewhere│
 └────────────┘     └──────────┘     └──────────────┘\discretionary{}{}{}     └──────────┘
        ▲                u                                      │
        └───────────────\discretionary{}{}{}─┘ {\tuataraFallbackTWO{}◀}───────────── F u ────────────────\discretionary{}{}{}──┘
\end{Verbatim}
\end{figure}

The keys on the arrows are the ones you press in the status: \texttt{s} stages, \texttt{u}
unstages, \texttt{c c} commits, \texttt{P p} pushes and \texttt{F u} pulls. The rest of the
chapter takes them one at a time.

\subsection{The status}
\label{sec:399}
\texttt{C-x g} shows the state of the repository that the current file or folder is
in. It is the place you work with git from. It shows the branch you are on,
the files that have changed, what is staged, and the latest commits.

The first time you open the status of a project, Tuatara asks whether you
trust it, because git can run programs that a repository names. Chapter 12
explains the question and its three answers. For a folder you made yourself,
answer \texttt{y} and then type \texttt{yes}.

Figure 15.2 shows the status of the recipes after an afternoon’s work.

\begin{figure}[H]
\caption{Figure 15.2. A status: one new file, one changed, one staged.}\label{org396}
\begin{Verbatim}[breaklines,breakautoindent=false,breakindent=1em,breakafter=/-.\,;:=,commandchars=\\\{\}]
Head:     main  Add the scone recipe
Upstream: origin/main (1 ahead, 0 behind)

{\tuataraFallbackTWO{}▾} Untracked files (1)
  pavlova.txt
{\tuataraFallbackTWO{}▾} Unstaged changes (1)
  modified   scones.txt
{\tuataraFallbackTWO{}▾} Staged changes (1)
  modified   bread.txt
{\tuataraFallbackTWO{}▾} Recent commits (2)
  3f9c2ab Add the scone recipe
  81d04e7 Start the recipe folder
{\tuataraFallbackTWO{}▸} Ignored files
\end{Verbatim}
\end{figure}

The first line names the branch and the subject of its last commit. The
second appears when the branch follows a branch on a remote, its upstream, and
says how many commits you have that the remote lacks (ahead) and how many it
has that you lack (behind). When git is in the middle of something, such as a
merge that stopped, a third line beginning \texttt{Doing:} says what and which keys
go on with it or stop it.

Below are the sections. Each has a heading with a count, and is left out
when it would be empty.

\begin{xltabular}{\textwidth}{@{}l >{\raggedright\arraybackslash\hsize=1.000\hsize}X@{}}
\toprule
\textbf{Section} & \textbf{What it holds}\\
\midrule
\endhead
Unmerged files & Files with a conflict, while a merge or a rebase has stopped.\\
Untracked files & New files that git does not keep yet. A new folder is one line.\\
Unstaged changes & Changes to files git keeps, which you have not chosen yet.\\
Staged changes & What the next commit will hold.\\
Stashes & Work you put aside (see The stash below).\\
Unpushed to \ldots{} & Commits the remote does not have yet.\\
Recent commits & The latest commits of the branch.\\
Ignored files & Files git was told to leave alone. Folded until you open it.\\
\bottomrule
\end{xltabular}

The status is a list in sections, like the buffer list of chapter 7, and the
same keys fold it. These are the keys in the order you are likely to need
them.

\begin{xltabular}{\textwidth}{@{}l l >{\raggedright\arraybackslash\hsize=1.000\hsize}X@{}}
\toprule
\textbf{Modal} & \textbf{Non-modal} & \textbf{What it does}\\
\midrule
\endhead
\texttt{j}, \texttt{k} & \texttt{C-n}, \texttt{C-p} & Down and up a line.\\
none & \texttt{n}, \texttt{p} & To the next and the previous file, hunk or heading.\\
none & \texttt{TAB} & Shows a file’s changes under it, or hides them. On a heading, folds or opens the section.\\
none & \texttt{S-TAB} & Folds every section, or opens them all.\\
none & \texttt{RET}, \texttt{l} & Opens the file at the change, or shows the commit or stash on the line.\\
none & \texttt{s} & Stages the file, the hunk or the selected lines.\\
none & \texttt{u} & Unstages them.\\
none & \texttt{x} & Reverses a change, after asking. It cannot be undone.\\
none & \texttt{c} & Opens the commit menu.\\
none & \texttt{L}, \texttt{d} & Open the log menu and the diff menu.\\
none & \texttt{b}, \texttt{t}, \texttt{z} & Open the branch, tag and stash menus.\\
none & \texttt{m}, \texttt{R} & Open the merge and rebase menus.\\
none & \texttt{f}, \texttt{F}, \texttt{P} & Open the fetch, pull and push menus.\\
none & \texttt{r} & Reads the repository again.\\
none & \texttt{q}, \texttt{h} & Closes the status.\\
\bottomrule
\end{xltabular}

You seldom need \texttt{r}. The status reads itself again when you save a file of the
repository, and when git changes the repository from anywhere, a terminal
included. It waits at least \texttt{*git-\codebreak{}watch-\codebreak{}seconds*} between two readings. A file
made or changed by another program, without git, shows after \texttt{r}.

Staged files have a colour of their own so that you can tell them from the
ones not staged, a staged deletion has another, and the commits and stashes a
third. A long section shows its first \texttt{*git-\codebreak{}section-\codebreak{}lines*} files and then a
line saying how many more there are; \texttt{TAB} on that line shows more.
\texttt{*git-\codebreak{}recent-\codebreak{}commits*} sets how many commits the last section lists.

\subsubsection{Menus}
\label{sec:397}
Most keys in the status open a menu, a small panel that lists what you can
do next, each action with a key. \texttt{c} opens the commit menu, and \texttt{c} again
there starts a commit, so the whole thing is written \texttt{c c}. \texttt{C-g} closes a
menu without doing anything, and \texttt{C-h} in a menu lists every action with its
description.

Many menus also have switches. A switch changes how the next action runs, and
its key begins with \texttt{-}. In the push menu, \texttt{- n} turns on “Only try”, and
the push that follows only reports what it would send. An option is a switch
that takes a value, and its key begins with \texttt{=}. In the log menu, \texttt{= n} asks
how many commits to list.

To keep the menus short, they show the actions used most. \texttt{\textasciitilde{}} in the status
switches that status between the short menus and the full ones. When this
chapter names an action you cannot see in a menu, press \texttt{\textasciitilde{}} first. To have the
full menus always, put this in your init file:

\begin{Verbatim}[breaklines,breakautoindent=false,breakindent=1em,breakafter=/-.\,;:=,commandchars=\\\{\}]
(setf *git-menus* :full)
\end{Verbatim}

\texttt{*git-\codebreak{}slim-\codebreak{}actions*} is the list of actions the short menus show, if you
would sooner add one or two than see them all.

\begin{try}
\textbf{Try it, with a folder of recipes}
\begin{enumerate}
\item In a terminal, make the folder and a first recipe:
   \texttt{mkdir \textasciitilde{}/recipes \&\& cd \textasciitilde{}/recipes \&\& git init \&\& printf 'Scones\textbackslash{}n\textbackslash{}n2 cups
   flour\textbackslash{}n1 cup milk\textbackslash{}n' > scones.\codebreak{}txt}.
\item In Tuatara, open \texttt{\textasciitilde{}/recipes/\codebreak{}scones.\codebreak{}txt} and press \texttt{C-x g}. Answer the
   question about trust with \texttt{y}, then type \texttt{yes} and \texttt{RET}.
\item The status opens. The first line says the branch and \texttt{(no commit yet)}.
   \texttt{scones.\codebreak{}txt} is under Untracked files.
\item Press \texttt{n} twice to go to the file’s line, then \texttt{l}. The file opens. Press
   \texttt{C-x g} to come back.
\item Press \texttt{TAB} on the heading Untracked files. The section folds to one line.
   Press \texttt{TAB} again to open it.
\item Press \texttt{q}. The status closes. \texttt{C-x g} opens it again whenever you want it.
\end{enumerate}
\end{try}

\subsubsection{Commands, settings and faces of the status}
\label{sec:398}
\begin{xltabular}{\textwidth}{@{}l l l >{\raggedright\arraybackslash\hsize=1.000\hsize}X@{}}
\toprule
\textbf{Modal} & \textbf{Non-modal} & \textbf{Command} & \textbf{What it does}\\
\midrule
\endhead
none & \texttt{C-x g} & \texttt{git} & Show the state of the git repository the current file (or folder) is in: the branch, the files not tracked, changed and staged, the latest commits; C-h b there lists its keys.\\
none & \texttt{s} & \texttt{git-\codebreak{}stage} & Stage the file on the line of the status, or on a section’s title all of its files; on a hunk, the hunk; with lines selected in hunks, those lines.\\
none & \texttt{u}, \texttt{C-c u} & \texttt{git-\codebreak{}unstage} & Unstage the file on the line of the status, or on a section’s title all of its files, a rename by both its names; on a hunk, the hunk; with lines selected in hunks, those lines.\\
none & \texttt{x} & \texttt{git-\codebreak{}reverse} & Reverse the changes of the file on the line of the status (on a section’s title, of all its files), after asking, which says it cannot be undone: git kept no copy of them.\\
none & \texttt{TAB} & \texttt{git-\codebreak{}toggle-\codebreak{}diff} & Show the changes of the file on the line of the status under it, in hunks, the first \textbf{GIT-DIFF-LINES} of them, or no longer.\\
none & \texttt{l}, \texttt{RET}, \texttt{C-c l} & \texttt{git-\codebreak{}visit} & Visit the file on the line of the status, at the line of the hunk’s line point is on, a file in conflict at its first conflict, its conflicts shown (M-x conflict-markers-mode); or show the commit or the stash on it, or a submodule’s status; a file deleted is said to be.\\
none & \texttt{n}, \texttt{C-c n} & \texttt{git-\codebreak{}next-\codebreak{}section} & Go to the next file or hunk of the status, or the next title.\\
none & \texttt{p}, \texttt{C-c p} & \texttt{git-\codebreak{}previous-\codebreak{}section} & Go to the previous file or hunk of the status, or the previous title.\\
none & \textasciitilde{}, C-c \textasciitilde{} & \texttt{git-\codebreak{}toggle-\codebreak{}full-\codebreak{}menus} & Make the git menus of this status show every action, or only the ones most used again (\textbf{GIT-MENUS}, \textbf{GIT-SLIM-ACTIONS}).\\
\bottomrule
\end{xltabular}

\begin{xltabular}{\textwidth}{@{}l >{\raggedright\arraybackslash\hsize=0.519\hsize}X >{\raggedright\arraybackslash\hsize=1.481\hsize}X@{}}
\toprule
\textbf{Setting} & \textbf{Standard} & \textbf{What it changes}\\
\midrule
\endhead
\texttt{*git-\codebreak{}recent-\codebreak{}commits*} & \texttt{10} & How many of the latest commits M-x git shows.\\
\texttt{*git-\codebreak{}section-\codebreak{}lines*} & \texttt{1000} & How many files each section of a status (M-x git) lists at first, and how many more each M-x git-toggle-diff on the line after them lists.\\
\texttt{*git-\codebreak{}watch-\codebreak{}seconds*} & \texttt{2} & The least time, in seconds, between two readings of a status on the screen when git changes its repository, from a terminal too (GIT-WATCH-TICK), or NIL for a status never to follow it: then M-x revert-buffer reads it again.\\
\texttt{*git-\codebreak{}menus*} & \texttt{:slim} & How much the git menus (branch, worktree \ldots{}) show: :SLIM the actions \textbf{GIT-SLIM-ACTIONS} names, the ones most used, :FULL every one.\\
\texttt{*git-\codebreak{}slim-\codebreak{}actions*} & \texttt{(tuatara.\codebreak{}git::\codebreak{}git-\codebreak{}commit tuatara.\codebreak{}git::\codebreak{}gi...} & The actions the slim git menus show (\textbf{GIT-MENUS}), by command: push one to show it too, or take one away.\\
\bottomrule
\end{xltabular}

\begin{xltabular}{\textwidth}{@{}l >{\raggedright\arraybackslash\hsize=1.000\hsize}X@{}}
\toprule
\textbf{Face} & \textbf{What it paints}\\
\midrule
\endhead
\texttt{:git-\codebreak{}status-\codebreak{}branch} & The branch you are on, at the top of a git status: lit up, to be seen at once.\\
\texttt{:git-\codebreak{}status-\codebreak{}staged} & A file staged, under Staged changes in a git status, apart from those not staged yet: in the theme’s colour for added lines.\\
\texttt{:git-\codebreak{}status-\codebreak{}deleted} & A file deleted and staged, under Staged changes in a git status: in the theme’s colour for lines taken away.\\
\texttt{:git-\codebreak{}status-\codebreak{}commit} & The latest commits at the end of a git status, and the stashes: a colour of their own.\\
\bottomrule
\end{xltabular}

\subsection{Staging}
\label{sec:403}
Staging is choosing. \texttt{s} on a file’s line moves the file to Staged changes,
and \texttt{u} there moves it back. Neither changes the file itself. On a section’s
heading, \texttt{s} stages every file of the section at once, which is the quick way
when all of your changes belong together.

\subsubsection{Hunks and single lines}
\label{sec:400}
Often a file holds two changes that deserve a commit each. Press \texttt{TAB} on the
file and its changes appear under it. Lines that were added begin with \texttt{+},
and lines that were taken away begin with \texttt{-}. The changes come in groups
called hunks, each a run of changed lines with a few unchanged lines around
it, under a line that begins with \texttt{@@}.

With the cursor on a hunk, \texttt{s} stages that hunk alone. For something finer
still, select lines inside a hunk with \texttt{V} and then press \texttt{s}: only the
selected lines are staged. \texttt{u} works the same way under Staged changes.

A file with a great many changes shows the first \texttt{*git-\codebreak{}diff-\codebreak{}lines*} lines of
them, then a line saying how many more there are. \texttt{TAB} on that line shows
more. Staging the file still takes all of it.

\subsubsection{Reversing a change}
\label{sec:401}
\texttt{x} throws a change away. On a file under Unstaged changes it puts the file
back to how the last commit has it. On a hunk it removes that hunk’s change,
and with lines selected, those lines. Tuatara asks first and says that it
cannot be undone, because git never had a copy of the text you are discarding.

A file that is not tracked is treated more gently. \texttt{x} moves it to the trash,
and you can put it back from there (chapter 10).

\subsubsection{From the file itself}
\label{sec:402}
You do not have to open the status to stage. In a saved file, with the cursor
on a changed line, \texttt{M-x git-\codebreak{}stage-\codebreak{}hunk} stages the change at the cursor and
\texttt{M-x git-\codebreak{}unstage-\codebreak{}hunk} takes it back. \texttt{M-x git-\codebreak{}stage-\codebreak{}file} and
\texttt{M-x git-\codebreak{}unstage-\codebreak{}file} do the whole file. The same actions are in the file
menu, described under Blame and the file menu below.

A file in a repository also shows its changes in the gutter, the narrow strip
at the left of the window. There is a mark beside each line added since the
last commit, another beside each changed line, and one where lines were
deleted. The marks are worked out in the background when the file is first
shown and after each save. \texttt{M-x vc-\codebreak{}gutter-\codebreak{}mode} turns them off everywhere, or
on again. \texttt{M-x vc-\codebreak{}gutter-\codebreak{}refresh} works them out anew, which you may want
after a commit made in a terminal.

\begin{xltabular}{\textwidth}{@{}l >{\raggedright\arraybackslash\hsize=0.520\hsize}X l >{\raggedright\arraybackslash\hsize=1.480\hsize}X@{}}
\toprule
\textbf{Modal} & \textbf{Non-modal} & \textbf{Command} & \textbf{What it does}\\
\midrule
\endhead
none & \texttt{M-x git-\codebreak{}stage-\codebreak{}hunk} & \texttt{git-\codebreak{}stage-\codebreak{}hunk} & Stage the hunk at point: in the status (C-x g) the one on point’s line, as s does; in a file’s own buffer, saved, the change git sees at point’s line, as the file is on disk.\\
none & \texttt{M-x git-\codebreak{}unstage-\codebreak{}hunk} & \texttt{git-\codebreak{}unstage-\codebreak{}hunk} & Unstage the hunk at point: in the status (C-x g) the one on point’s line, as u does; in a file’s own buffer, saved, the change staged at point’s line.\\
none & \texttt{M-x git-\codebreak{}stage-\codebreak{}file} & \texttt{git-\codebreak{}stage-\codebreak{}file} & Stage the file you are in (as saved), or the one on point’s line of a status.\\
none & \texttt{M-x git-\codebreak{}unstage-\codebreak{}file} & \texttt{git-\codebreak{}unstage-\codebreak{}file} & Unstage the file you are in, or the one on point’s line of a status, in a repository with no commit yet too.\\
none & \texttt{M-x vc-\codebreak{}gutter-\codebreak{}mode} & \texttt{vc-\codebreak{}gutter-\codebreak{}mode} & Switch the version-control gutter on or off for all buffers; with a count, on when it is positive, off when it is zero or less.\\
none & \texttt{M-x vc-\codebreak{}gutter-\codebreak{}refresh} & \texttt{vc-\codebreak{}gutter-\codebreak{}refresh} & Recompute the version-control marks of the current buffer from git.\\
\bottomrule
\end{xltabular}

\begin{xltabular}{\textwidth}{@{}l l >{\raggedright\arraybackslash\hsize=1.000\hsize}X@{}}
\toprule
\textbf{Setting} & \textbf{Standard} & \textbf{What it changes}\\
\midrule
\endhead
\texttt{*git-\codebreak{}diff-\codebreak{}lines*} & \texttt{1000} & How many lines of a file’s changes M-x git-toggle-diff shows under it at first, and how many more each M-x git-toggle-diff on the line after them shows.\\
\bottomrule
\end{xltabular}

\begin{longtable}{@{}l l@{}}
\toprule
\textbf{Face} & \textbf{What it paints}\\
\midrule
\endhead
\texttt{:vc-\codebreak{}added} & Gutter mark for lines added since the last commit.\\
\texttt{:vc-\codebreak{}changed} & Gutter mark for lines changed since the last commit.\\
\texttt{:vc-\codebreak{}deleted} & Gutter mark after lines deleted since the last commit.\\
\bottomrule
\end{longtable}

\subsection{Committing}
\label{sec:407}
When what is staged is one finished piece of work, commit it. \texttt{c c} opens a
buffer for the message, with a list of what is staged under it on lines that
begin with \texttt{\#}. Those lines are for you to read and are left out of the
commit. The buffer starts in INSERT state, so you can type at once.

A good message has a short first line that says what changed, such as
\texttt{Lower the oven temperature for scones}. If there is more to say, leave an
empty line and write a paragraph on why. The lines are filled at
\texttt{*git-\codebreak{}commit-\codebreak{}fill-\codebreak{}column*} as you type.

\texttt{C-c C-c} makes the commit and closes the buffer. \texttt{C-c C-k} gives the message
up, and nothing is committed.

\begin{try}
\textbf{Try it, with the recipes}
\begin{enumerate}
\item Open the status of \texttt{\textasciitilde{}/recipes} with \texttt{C-x g}. Move to \texttt{scones.\codebreak{}txt} and press
   \texttt{s}. The file moves to Staged changes.
\item Press \texttt{c}, then \texttt{c}. The message buffer opens.
\item Type \texttt{Start the recipe folder}. Press \texttt{C-g} and then \texttt{C-c C-c}.
\item The status now has a section Recent commits with your commit in it, and the
   first line shows its subject.
\item Open \texttt{scones.\codebreak{}txt}, add the line \texttt{1 pinch salt} at the end and, a few lines
   up, change \texttt{2 cups flour} to \texttt{3 cups flour}. Save with \texttt{C-x C-s} and press
   \texttt{C-x g}.
\item Press \texttt{TAB} on \texttt{scones.\codebreak{}txt}. Move to the line \texttt{+1 pinch salt}, press \texttt{V} and
   then \texttt{s}. Only that line is staged.
\item Press \texttt{c c}, type \texttt{Add salt to the scones}, press \texttt{C-g} and \texttt{C-c C-c}. The
   change to the flour is still waiting under Unstaged changes.
\end{enumerate}
\end{try}

\subsubsection{Changing the last commit}
\label{sec:404}
Everyone commits and then sees a spelling mistake. The commit menu has three
ways to change the last commit, as long as you have not pushed it yet.

\begin{xltabular}{\textwidth}{@{}l l >{\raggedright\arraybackslash\hsize=1.000\hsize}X@{}}
\toprule
\textbf{Modal} & \textbf{Non-modal} & \textbf{What it does}\\
\midrule
\endhead
none & \texttt{c a} & Amend. What is staged goes into the last commit, and you edit its message.\\
none & \texttt{c w} & Reword. Only the message is edited.\\
none & \texttt{c e} & Extend. What is staged goes into the last commit, with the message as it is.\\
\bottomrule
\end{xltabular}

Use extend when you forgot a file, reword when the message is wrong, and amend
when both need fixing.

The full menu has two more, for a mistake further back. \texttt{c f} (fix up) and
\texttt{c s} (squash) make a small commit that is marked as belonging to an earlier
one you pick. A later rebase with the switch \texttt{- q} folds it into that commit.
Rebasing is described below.

The switches of the commit menu go to git with the commit. \texttt{- a} commits
every change to tracked files without staging first. \texttt{- s} adds a
\texttt{Signed-\codebreak{}off-by} line with your name, which some projects ask for.

\subsubsection{Signing commits}
\label{sec:405}
A signed commit carries proof that it was made with your key, so that others
can check it came from you. Tuatara commits the way git does with your own
settings, signing included. There is nothing to switch on in Tuatara. You tell
git once, in a terminal, which key to use and that every commit is to be
signed:

\begin{Verbatim}[breaklines,breakautoindent=false,breakindent=1em,breakafter=/-.\,;:=,commandchars=\\\{\}]
git config --global user.signingkey YOUR-KEY
git config --global commit.gpgsign true
\end{Verbatim}

From then on \texttt{C-c C-c} in a message buffer makes a signed commit. Git’s own
manual describes the kinds of key it accepts. A tag can be signed as well: in
the tag menu, turn on \texttt{- s} before you make the tag.

Signing and signing off are different things. The switch \texttt{- s} in the commit
menu only adds a line of text to the message.

To see which commits are signed, press \texttt{!} in the status and type
\texttt{log --format}”\%h \%G? \%s”\texttt{. A =G} before a subject means a good signature and
an \texttt{N} means none.

\subsubsection{Commands and settings for committing}
\label{sec:406}
\begin{xltabular}{\textwidth}{@{}l >{\raggedright\arraybackslash\hsize=0.480\hsize}X l >{\raggedright\arraybackslash\hsize=1.520\hsize}X@{}}
\toprule
\textbf{Modal} & \textbf{Non-modal} & \textbf{Command} & \textbf{What it does}\\
\midrule
\endhead
none & \texttt{c}, \texttt{C-c c} & \texttt{git-\codebreak{}commit-\codebreak{}menu} & Open the commit menu: write a commit, amend, reword or extend the last one, or make one fixing up or squashed into another; its switches go to git commit.\\
none & \texttt{M-x git-\codebreak{}commit} & \texttt{git-\codebreak{}commit} & Write a commit of what is staged: its message in a buffer of its own, what is staged listed under it as comments; C-c C-c commits it, C-c C-k gives it up.\\
none & \texttt{C-c C-c} & \texttt{git-\codebreak{}commit-\codebreak{}finish} & Commit with the message of this buffer (lines beginning with \# left out), as git commit does with your settings (a signature among them) and the arguments the commit was begun with (--amend for an amended one), and close it.\\
none & \texttt{C-c C-k} & \texttt{git-\codebreak{}commit-\codebreak{}cancel} & Give up this commit’s message: nothing is committed.\\
none & \texttt{M-x git-\codebreak{}commit-\codebreak{}amend} & \texttt{git-\codebreak{}commit-\codebreak{}amend} & Amend the last commit: what is staged goes into it, and its message, written in a buffer to begin with, is written again; C-c C-c amends it.\\
none & \texttt{M-x git-\codebreak{}commit-\codebreak{}reword} & \texttt{git-\codebreak{}commit-\codebreak{}reword} & Reword the last commit: its message, written in a buffer to begin with, is written again, nothing staged going into it; C-c C-c rewords it.\\
none & \texttt{M-x git-\codebreak{}commit-\codebreak{}extend} & \texttt{git-\codebreak{}commit-\codebreak{}extend} & Extend the last commit with what is staged, its message as it is.\\
none & \texttt{M-x git-\codebreak{}commit-\codebreak{}fixup} & \texttt{git-\codebreak{}commit-\codebreak{}fixup} & Make a commit of what is staged fixing up one asked for: a rebase with autosquash folds it into that one, its message left as it was.\\
none & \texttt{M-x git-\codebreak{}commit-\codebreak{}squash} & \texttt{git-\codebreak{}commit-\codebreak{}squash} & Make a commit of what is staged to be squashed into one asked for: a rebase with autosquash folds it in, its message added to that one’s.\\
none & \texttt{M-x git-\codebreak{}commit-\codebreak{}mode} & \texttt{git-\codebreak{}commit-\codebreak{}mode} & Switch the current buffer to git-commit-mode.\\
\bottomrule
\end{xltabular}

\begin{xltabular}{\textwidth}{@{}l l >{\raggedright\arraybackslash\hsize=1.000\hsize}X@{}}
\toprule
\textbf{Setting} & \textbf{Standard} & \textbf{What it changes}\\
\midrule
\endhead
\texttt{*git-\codebreak{}commit-\codebreak{}fill-\codebreak{}column*} & \texttt{72} & The column a git message’s lines are filled to: 72, as git recommends for the body; the summary on the first line is best kept to 50.\\
\bottomrule
\end{xltabular}

\subsection{The log and diffs}
\label{sec:410}
\subsubsection{The log}
\label{sec:408}
The log is the list of commits, the latest first. \texttt{L L} in the status lists
the commits of the branch you are on, a line each: the hash, the date, the
author, any branch or tag that sits at the commit, and the subject. \texttt{RET} or
\texttt{l} on a line shows the commit, with its message and every change it made.
\texttt{q} closes.

The log menu lists other things too.

\begin{xltabular}{\textwidth}{@{}l l >{\raggedright\arraybackslash\hsize=1.000\hsize}X@{}}
\toprule
\textbf{Modal} & \textbf{Non-modal} & \textbf{What it lists}\\
\midrule
\endhead
none & \texttt{L L} & The branch you are on.\\
none & \texttt{L o} & Another branch, or any revision you name.\\
none & \texttt{L f} & The commits that changed one file: the one on the line, or the current file.\\
none & \texttt{L b} & Every branch, the remotes’ included. (Full menu.)\\
none & \texttt{L a} & Every reference: branches, tags and stashes. (Full menu.)\\
none & \texttt{L r} & The reflog: where the branch has been, step by step. (Full menu.)\\
\bottomrule
\end{xltabular}

A revision is any way of naming a commit: a branch name, a tag, a hash, or
\texttt{HEAD}, which is git’s word for the commit you are on now.

The switches and options shape the list. \texttt{- g} draws the graph of branches
beside the commits, which helps after a merge. \texttt{= a} keeps one author’s
commits, \texttt{= s} those whose message contains a word, and \texttt{= d} and \texttt{= u} those
since or until a date; git understands dates such as \texttt{2 weeks ago}. The log
shows \texttt{*git-\codebreak{}log-\codebreak{}length*} commits unless \texttt{= n} says otherwise.

The reflog is worth knowing about on a bad day. Git records every move of the
branch, including the ones that seemed to throw commits away. A commit you
lost by a mistaken reset is still listed there for some weeks, and you can
check it out or make a branch at it.

\begin{xltabular}{\textwidth}{@{}l l l >{\raggedright\arraybackslash\hsize=1.000\hsize}X@{}}
\toprule
\textbf{Modal} & \textbf{Non-modal} & \textbf{Command} & \textbf{What it does}\\
\midrule
\endhead
none & \texttt{L}, \texttt{C-c L} & \texttt{git-\codebreak{}log-\codebreak{}menu} & Open the log menu: list the commits of the branch checked out, another, every branch or reference, or a file’s, or the reflog; its switches and options shape the list.\\
none & \texttt{M-x git-\codebreak{}log} & \texttt{git-\codebreak{}log} & List the commits of the branch checked out, the latest first, a line each (\textbf{GIT-LOG-LENGTH} of them); RET shows one.\\
none & \texttt{M-x git-\codebreak{}log-\codebreak{}other} & \texttt{git-\codebreak{}log-\codebreak{}other} & List the commits of a branch or revision asked for.\\
none & \texttt{M-x git-\codebreak{}log-\codebreak{}file} & \texttt{git-\codebreak{}log-\codebreak{}file} & List the commits that changed the file on point’s line of a status, or the current file, following it across renames.\\
none & \texttt{M-x git-\codebreak{}log-\codebreak{}branches} & \texttt{git-\codebreak{}log-\codebreak{}branches} & List the commits of every branch, the repository’s own and its remotes’.\\
none & \texttt{M-x git-\codebreak{}log-\codebreak{}all} & \texttt{git-\codebreak{}log-\codebreak{}all} & List the commits of every reference: branches, remotes’ branches, tags, stashes.\\
none & \texttt{M-x git-\codebreak{}reflog} & \texttt{git-\codebreak{}reflog} & List where HEAD has been, the latest first: each commit it was at and what moved it there (a commit, a checkout, a reset), as git’s reflog keeps it.\\
none & \texttt{l}, \texttt{RET}, \texttt{C-c l} & \texttt{git-\codebreak{}log-\codebreak{}show} & Show the commit on the current line of the log.\\
\bottomrule
\end{xltabular}

\begin{longtable}{@{}l l l@{}}
\toprule
\textbf{Setting} & \textbf{Standard} & \textbf{What it changes}\\
\midrule
\endhead
\texttt{*git-\codebreak{}log-\codebreak{}length*} & \texttt{200} & How many commits M-x git-log lists.\\
\bottomrule
\end{longtable}

\subsubsection{Diffs}
\label{sec:409}
A diff shows how two states differ, line by line, in the same form as the
changes under a file in the status. \texttt{d d} shows everything you have changed
since the last commit, in a buffer of its own. It is a good thing to read
through before you commit.

\begin{longtable}{@{}l l l@{}}
\toprule
\textbf{Modal} & \textbf{Non-modal} & \textbf{What it shows}\\
\midrule
\endhead
none & \texttt{d d} & Every change not committed, staged or not.\\
none & \texttt{d s} & What is staged: the changes the next commit will make.\\
none & \texttt{d u} & What is not staged. (Full menu.)\\
none & \texttt{d r} & How your files differ from a revision you name. (Full menu.)\\
none & \texttt{d R} & What changed between two revisions. (Full menu.)\\
none & \texttt{d c} & One commit.\\
\bottomrule
\end{longtable}

\texttt{M-x git-\codebreak{}diff-\codebreak{}file} shows the uncommitted changes of the file you are in.

Among the switches, \texttt{- w} ignores changes in spaces, which helps when a file
was only re-indented. \texttt{- s} shows the names of the files and how many lines
changed in each, without the lines. \texttt{= U} sets how many unchanged lines are
shown around each change.

Added and removed lines have their own colours. A diff larger than
\texttt{*differences-\codebreak{}limit*} octets is not shown, and Tuatara says so. Comparing two
files that are not in git is \texttt{M-x diff-\codebreak{}files}, described in chapter 9.

\begin{xltabular}{\textwidth}{@{}l >{\raggedright\arraybackslash\hsize=0.520\hsize}X l >{\raggedright\arraybackslash\hsize=1.480\hsize}X@{}}
\toprule
\textbf{Modal} & \textbf{Non-modal} & \textbf{Command} & \textbf{What it does}\\
\midrule
\endhead
none & \texttt{d}, \texttt{C-c d} & \texttt{git-\codebreak{}diff-\codebreak{}menu} & Open the diff menu: show the changes not committed, what is staged, what is not, the files against a revision, between two revisions, or a commit; its switches shape them.\\
none & \texttt{M-x git-\codebreak{}diff} & \texttt{git-\codebreak{}diff} & Show every change of the repository not yet committed, staged or not, in a buffer of its own; the diff menu’s switches shape it.\\
none & \texttt{M-x git-\codebreak{}diff-\codebreak{}staged} & \texttt{git-\codebreak{}diff-\codebreak{}staged} & Show what is staged: the changes the next commit makes.\\
none & \texttt{M-x git-\codebreak{}diff-\codebreak{}unstaged} & \texttt{git-\codebreak{}diff-\codebreak{}unstaged} & Show the changes not staged: those in the files beyond what is staged.\\
none & \texttt{M-x git-\codebreak{}diff-\codebreak{}revision} & \texttt{git-\codebreak{}diff-\codebreak{}revision} & Show how the files differ from a revision asked for (a branch, a tag, a commit).\\
none & \texttt{M-x git-\codebreak{}diff-\codebreak{}range} & \texttt{git-\codebreak{}diff-\codebreak{}range} & Show what changed between two revisions asked for: from the first to the second.\\
none & \texttt{M-x git-\codebreak{}diff-\codebreak{}commit} & \texttt{git-\codebreak{}diff-\codebreak{}commit} & Show a commit asked for (the one on the line first): what it says and its changes.\\
none & \texttt{M-x git-\codebreak{}diff-\codebreak{}file} & \texttt{git-\codebreak{}diff-\codebreak{}file} & Show the changes of the file you are in not yet committed, staged or not.\\
none & \texttt{M-x diff-\codebreak{}mode} & \texttt{diff-\codebreak{}mode} & Switch the current buffer to diff-mode.\\
\bottomrule
\end{xltabular}

\begin{xltabular}{\textwidth}{@{}l l >{\raggedright\arraybackslash\hsize=1.000\hsize}X@{}}
\toprule
\textbf{Setting} & \textbf{Standard} & \textbf{What it changes}\\
\midrule
\endhead
\texttt{*differences-\codebreak{}limit*} & \texttt{67108864} & The most octets of differences shown (64 MB); when there are more, none are shown, and that is said.\\
\bottomrule
\end{xltabular}

\begin{longtable}{@{}l l@{}}
\toprule
\textbf{Face} & \textbf{What it paints}\\
\midrule
\endhead
\texttt{:diff.\codebreak{}plus} & Lines a difference between files adds.\\
\texttt{:diff.\codebreak{}minus} & Lines a difference between files takes away.\\
\texttt{:diff.\codebreak{}delta} & Lines a difference between files changes.\\
\bottomrule
\end{longtable}

\subsection{Branches and tags}
\label{sec:415}
\subsubsection{Branches}
\label{sec:411}
A branch lets you work on something without disturbing what is finished.
Suppose you want to convert all the recipes to metric measures. It will take a
week, and in the meantime you still want to fix small things in the recipes
as they are. Make a branch called \texttt{metric}, do the conversion there one commit
at a time, and go back to \texttt{main} whenever you need to.

Going to a branch is called checking it out. Git then changes the files in
the folder to that branch’s state. If a buffer still shows a file’s old text
afterwards, \texttt{M-x revert-\codebreak{}buffer} reads the file again (chapter 9).

\begin{xltabular}{\textwidth}{@{}l l >{\raggedright\arraybackslash\hsize=1.000\hsize}X@{}}
\toprule
\textbf{Modal} & \textbf{Non-modal} & \textbf{What it does}\\
\midrule
\endhead
none & \texttt{b c} & Makes a new branch and checks it out. It asks for the name, then where to start; \texttt{RET} starts from where you are.\\
none & \texttt{b b} & Checks out a branch or a revision you name.\\
none & \texttt{b m} & Renames a branch.\\
none & \texttt{b d} & Deletes a branch, though never the one you are on.\\
\bottomrule
\end{xltabular}

Git refuses to check out another branch when that would overwrite changes you
have not committed, and Tuatara shows what it said. Commit them first, or put
them aside with the stash.

Deleting a branch whose commits are on no other branch would lose them.
Tuatara asks a second time in that case.

If you name a commit or a tag at \texttt{b b} where a branch was expected, git checks
out that state with no branch. The status then says \texttt{(detached)} on its first
line. You can look around, but commits made there belong to no branch. Check
out a branch to return.

The full menu has four more actions. \texttt{b n} makes a branch without checking
it out. \texttt{b x} moves a branch to another commit. \texttt{b u} sets a branch’s
upstream. \texttt{b s} is for the day you find you have been committing on \texttt{main}
when you meant to use a branch: it makes a new branch with the commits you
have not pushed, checks it out, and puts \texttt{main} back to where its upstream
is.

\texttt{y} in the status lists every branch, every branch of the remotes and every
tag, each with the commit it is at. \texttt{RET} on a line shows its log.

\begin{xltabular}{\textwidth}{@{}l >{\raggedright\arraybackslash\hsize=0.591\hsize}X l >{\raggedright\arraybackslash\hsize=1.409\hsize}X@{}}
\toprule
\textbf{Modal} & \textbf{Non-modal} & \textbf{Command} & \textbf{What it does}\\
\midrule
\endhead
none & \texttt{b}, \texttt{C-c b} & \texttt{git-\codebreak{}branch-\codebreak{}menu} & Open the branch menu: check out a branch or a revision, make one (checked out or not, or spun off with the work not pushed), rename, move, delete one, or set its upstream.\\
none & \texttt{M-x git-\codebreak{}branch-\codebreak{}and-\codebreak{}checkout} & \texttt{git-\codebreak{}branch-\codebreak{}and-\codebreak{}checkout} & Make a new branch, asked for, starting at a revision asked for (the branch checked out when nothing is typed), and check it out.\\
none & \texttt{M-x git-\codebreak{}checkout} & \texttt{git-\codebreak{}checkout} & Check out a branch or a revision, asked for: a branch by its name, else a revision (a commit, a tag, origin/main), which leaves HEAD detached.\\
none & \texttt{M-x git-\codebreak{}branch-\codebreak{}rename} & \texttt{git-\codebreak{}branch-\codebreak{}rename} & Rename a branch, asked for (the one checked out when nothing is typed), to a name asked for.\\
none & \texttt{M-x git-\codebreak{}branch-\codebreak{}delete} & \texttt{git-\codebreak{}branch-\codebreak{}delete} & Delete a branch, asked for, not the one checked out: one with commits no other branch has is deleted only after asking again, as they would be found only in git’s reflog then.\\
none & \texttt{M-x git-\codebreak{}branch-\codebreak{}create} & \texttt{git-\codebreak{}branch-\codebreak{}create} & Make a new branch, asked for, starting at a revision asked for, without checking it out.\\
none & \texttt{M-x git-\codebreak{}branch-\codebreak{}reset} & \texttt{git-\codebreak{}branch-\codebreak{}reset} & Move a branch, asked for, to a revision asked for, after asking: the branch checked out keeps the changes not committed (git refuses when they would be lost); its commits after the revision stay in git’s reflog.\\
none & \texttt{M-x git-\codebreak{}branch-\codebreak{}set-\codebreak{}upstream} & \texttt{git-\codebreak{}branch-\codebreak{}set-\codebreak{}upstream} & Set the upstream of a branch, asked for (the one checked out when nothing is typed), to a branch of a remote asked for: the one pull brings from and push sends to, and the status counts ahead and behind of.\\
none & \texttt{M-x git-\codebreak{}branch-\codebreak{}spin-\codebreak{}off} & \texttt{git-\codebreak{}branch-\codebreak{}spin-\codebreak{}off} & Make a new branch, asked for, of the work on the branch checked out and check it out: the commits the branch has that its upstream has not go with the new branch, and the branch goes back to where it and its upstream parted; changes not committed come along.\\
none & \texttt{y} & \texttt{git-\codebreak{}show-\codebreak{}refs} & List the repository’s branches (the one checked out marked *, each with its upstream), its remotes’ branches and its tags, a line each with the commit it is at and its subject, in sections folded as in every list in sections; RET on one lists its log.\\
none & \texttt{l}, \texttt{RET}, \texttt{C-c l} & \texttt{git-\codebreak{}refs-\codebreak{}log} & List the log of the branch or tag on the line of the list of references; on a section’s heading, fold or unfold the section.\\
\bottomrule
\end{xltabular}

\subsubsection{Merging}
\label{sec:413}
When the work on a branch is done, you merge it into \texttt{main}. Check out \texttt{main},
then press \texttt{m m} and name the branch. Git joins the two lines of work and
makes a commit that records the join. Figure 15.3 shows the shape.

\begin{figure}[H]
\caption{Figure 15.3. A branch made, worked on and merged back.}\label{org412}
\begin{Verbatim}[breaklines,breakautoindent=false,breakindent=1em,breakafter=/-.\,;:=,commandchars=\\\{\}]
                 o───o───o   metric
                {\tuataraFallbackTWO{}╱}         {\tuataraFallbackTWO{}╲}
   o───o───o───o───\discretionary{}{}{}────────o   main
               ▲           ▲
             b c         m m
\end{Verbatim}
\end{figure}

If \texttt{main} did not change while you worked on the branch, there is nothing to
join. Git moves \texttt{main} forward to the branch’s last commit, which is called a
fast-forward, and makes no extra commit.

The full menu has two more ways to merge. \texttt{m n} merges but stops before the
commit so that you can look first. \texttt{m s} brings in all of a branch’s changes
as one staged change, for you to commit as a single commit.
When both branches changed the same lines, the merge stops half way. That is
a conflict, and it has a section of its own below.

\begin{xltabular}{\textwidth}{@{}l >{\raggedright\arraybackslash\hsize=0.542\hsize}X l >{\raggedright\arraybackslash\hsize=1.458\hsize}X@{}}
\toprule
\textbf{Modal} & \textbf{Non-modal} & \textbf{Command} & \textbf{What it does}\\
\midrule
\endhead
none & \texttt{m}, \texttt{C-c m} & \texttt{git-\codebreak{}merge-\codebreak{}menu} & Open the merge menu: merge a branch, committed or not, squash its changes into one, or stop a merge stopped half way.\\
none & \texttt{M-x git-\codebreak{}merge} & \texttt{git-\codebreak{}merge} & Merge a branch, asked for, into the one checked out, committing it with git’s own message; files changed on both sides stop it half way, said so in the status (c c commits once they are put right, m a stops it).\\
none & \texttt{M-x git-\codebreak{}merge-\codebreak{}no-\codebreak{}commit} & \texttt{git-\codebreak{}merge-\codebreak{}no-\codebreak{}commit} & Merge a branch, asked for, without committing: the merge is left staged, to look at and commit (c c) or stop (m a).\\
none & \texttt{M-x git-\codebreak{}merge-\codebreak{}squash} & \texttt{git-\codebreak{}merge-\codebreak{}squash} & Bring the changes of a branch, asked for, into the files and stage them as one change, not a merge: c c commits them.\\
none & \texttt{M-x git-\codebreak{}merge-\codebreak{}continue} & \texttt{git-\codebreak{}merge-\codebreak{}continue} & Go on with the merge stopped half way, once the files git could not join are put right and staged.\\
none & \texttt{M-x git-\codebreak{}merge-\codebreak{}abort} & \texttt{git-\codebreak{}merge-\codebreak{}abort} & Stop the merge, after asking: the branch and the files go back to where they were before it began.\\
\bottomrule
\end{xltabular}

\subsubsection{Tags}
\label{sec:414}
A tag is a name for one commit that never moves. You would tag the state of
the recipes that you printed for the school fair as \texttt{fair-\codebreak{}2026}, so that you
can always get exactly that state back. \texttt{t t} asks for the revision, offering
the commit on the cursor’s line, and for the name. \texttt{t d} deletes a tag.

With the switch \texttt{- a} the tag carries a message of its own, and with \texttt{- s} it
is signed. Tags are not sent by an ordinary push. The full push menu sends
them: \texttt{P t} every tag, \texttt{P T} one.

\begin{xltabular}{\textwidth}{@{}l l l >{\raggedright\arraybackslash\hsize=1.000\hsize}X@{}}
\toprule
\textbf{Modal} & \textbf{Non-modal} & \textbf{Command} & \textbf{What it does}\\
\midrule
\endhead
none & \texttt{t} & \texttt{git-\codebreak{}tag-\codebreak{}menu} & Open the tag menu: tag a revision, or delete a tag; P t and P T in the push menu send tags.\\
none & \texttt{M-x git-\codebreak{}tag-\codebreak{}create} & \texttt{git-\codebreak{}tag-\codebreak{}create} & Tag a revision asked for (the commit on point’s line first, else HEAD) with a name asked for; the tag menu’s -a makes it annotated with a message asked for, -s signed (with a message too), -f moves a tag of that name.\\
none & \texttt{M-x git-\codebreak{}tag-\codebreak{}delete} & \texttt{git-\codebreak{}tag-\codebreak{}delete} & Delete a tag, asked for, here (not on any remote), after asking.\\
\bottomrule
\end{xltabular}

\begin{try}
\textbf{Try it, with the recipes}
\begin{enumerate}
\item In the status of \texttt{\textasciitilde{}/recipes}, press \texttt{b c}. Type \texttt{metric} and \texttt{RET}, then
   \texttt{RET} again to start from where you are. The first line now says \texttt{metric}.
\item Open \texttt{scones.\codebreak{}txt}, change \texttt{1 cup milk} to \texttt{250 ml milk} and save. Press
   \texttt{C-x g}, then \texttt{s} on the file, \texttt{c c}, type \texttt{Milk in millilitres}, \texttt{C-g} and
   \texttt{C-c C-c}.
\item Press \texttt{b b}, type \texttt{main} and \texttt{RET}. The first line says \texttt{main}, and the
   commit you just made is not among the recent commits. On this branch the
   milk is still in cups.
\item Press \texttt{m m}, type \texttt{metric} and \texttt{RET}. The commit is now on \texttt{main} too, and
   the file says \texttt{250 ml}.
\item Press \texttt{L}, then \texttt{- g} and \texttt{L}. The log shows the commits with the graph
   beside them. Press \texttt{q}.
\item Press \texttt{b d}, type \texttt{metric} and \texttt{RET}, then answer \texttt{y} and \texttt{RET} to the
   question. The branch is merged, so it can go.
\end{enumerate}
\end{try}

\subsection{The stash}
\label{sec:416}
Sometimes you are half way through a change and need the files clean for a
moment, for instance to check out \texttt{main} and fix a typing mistake someone
pointed out. The work is not ready for a commit. The stash is a shelf for
it. \texttt{z z} asks for a short message, puts every change aside and returns the
files to the last commit. The stash appears in the status under Stashes.

When you come back, \texttt{z p} (pop) puts the changes into the files again and
takes the stash off the shelf. \texttt{z a} (apply) does the same and keeps the
stash, for when you want the changes on two branches. \texttt{z d} drops a stash
without applying it, and that cannot be undone.

If the files changed in the meantime in the same places, popping can end in
a conflict. Tuatara then keeps the stash and says so, and nothing is lost.

New files that git does not track yet are left where they are unless you turn
on the switch \texttt{- u} first.

The full menu adds \texttt{z i}, which stashes only what is staged, \texttt{z x}, which
stashes everything and leaves the staged part in place, \texttt{z v}, which shows
what a stash holds, and \texttt{z b}, which makes a new branch out of a stash.

\begin{try}
\textbf{Try it, with the recipes}
\begin{enumerate}
\item Open \texttt{\textasciitilde{}/recipes/\codebreak{}scones.\codebreak{}txt} and add a line \texttt{Serve with jam} at the end.
   Save, and press \texttt{C-x g}. The file is under Unstaged changes.
\item Press \texttt{z z}, type \texttt{serving note} and \texttt{RET}. The change is gone from the
   file, and the status has a section Stashes.
\item Press \texttt{RET} on the stash to see what it holds, then \texttt{q}.
\item Press \texttt{z p} and pick the stash. The line is back in the file and the
   section Stashes is gone.
\end{enumerate}
\end{try}

\begin{xltabular}{\textwidth}{@{}l >{\raggedright\arraybackslash\hsize=0.426\hsize}X l >{\raggedright\arraybackslash\hsize=1.574\hsize}X@{}}
\toprule
\textbf{Modal} & \textbf{Non-modal} & \textbf{Command} & \textbf{What it does}\\
\midrule
\endhead
none & \texttt{z}, \texttt{C-c z} & \texttt{git-\codebreak{}stash-\codebreak{}menu} & Open the stash menu: put changes aside (all, only what is staged, or keeping it), apply, pop, drop or show a stash, or make a branch of one.\\
none & \texttt{M-x git-\codebreak{}stash} & \texttt{git-\codebreak{}stash} & Put every change aside as a stash, staged or not, with a message asked for; the files go back to the last commit.\\
none & \texttt{M-x git-\codebreak{}stash-\codebreak{}pop} & \texttt{git-\codebreak{}stash-\codebreak{}pop} & Bring the changes of a stash, asked for, back into the files and drop it; one whose changes meet others is kept, said so.\\
none & \texttt{M-x git-\codebreak{}stash-\codebreak{}apply} & \texttt{git-\codebreak{}stash-\codebreak{}apply} & Bring the changes of a stash, asked for, back into the files, the stash kept; git says when they meet changes of their own.\\
none & \texttt{M-x git-\codebreak{}stash-\codebreak{}drop} & \texttt{git-\codebreak{}stash-\codebreak{}drop} & Drop a stash, asked for, after asking: its changes are gone.\\
none & \texttt{M-x git-\codebreak{}stash-\codebreak{}index} & \texttt{git-\codebreak{}stash-\codebreak{}index} & Put only what is staged aside as a stash, with a message asked for; the changes not staged stay.\\
none & \texttt{M-x git-\codebreak{}stash-\codebreak{}keep-\codebreak{}index} & \texttt{git-\codebreak{}stash-\codebreak{}keep-\codebreak{}index} & Put every change aside as a stash, with a message asked for, keeping what is staged in place in the files and staged.\\
none & \texttt{M-x git-\codebreak{}stash-\codebreak{}show} & \texttt{git-\codebreak{}stash-\codebreak{}show} & Show a stash, asked for: the files it changes and its changes.\\
none & \texttt{M-x git-\codebreak{}stash-\codebreak{}branch} & \texttt{git-\codebreak{}stash-\codebreak{}branch} & Make a new branch, asked for, of a stash, asked for: checked out at the commit the stash was made on, its changes brought back, the stash dropped.\\
\bottomrule
\end{xltabular}

\subsection{Merge conflicts}
\label{sec:417}
A conflict happens when two lines of work changed the same lines of a file.
Git cannot know which is right, so it stops and asks you. It can happen in a
merge, a rebase, a cherry-pick, or when a stash comes back.

Nothing is broken when this happens. The status says what git is doing on its
\texttt{Doing:} line, with the keys that go on or stop, and it lists the files in
question first, in a section called Unmerged files. Each line says how the two
sides differ, for example \texttt{both changed}.

Git writes both versions into the file, between marker lines:

\begin{Verbatim}[breaklines,breakautoindent=false,breakindent=1em,breakafter=/-.\,;:=,commandchars=\\\{\}]
<<<<<<< HEAD
2 cups flour
=======
300 g flour
>>>>>>> metric
\end{Verbatim}

The lines after \texttt{<<<<<<<} are the first side. In a merge that is the branch
you are on. The lines after the row of equals signs are the second side, the
branch being merged in.

To settle a conflict, follow these steps.

\begin{enumerate}
\item Press \texttt{RET} on the file in the status. It opens with the cursor on its
   first conflict, with the two sides coloured.
\item For each conflict, decide. \texttt{C-c \textasciicircum{} u} keeps the first side, \texttt{C-c \textasciicircum{} l} the
   second and \texttt{C-c \textasciicircum{} a} both. \texttt{C-c \textasciicircum{} n} goes to the next conflict. You can
   also edit the lines by hand and delete the markers. Chapter 13 describes
   these keys in full.
\item Save the file with \texttt{C-x C-s}.
\item Press \texttt{C-x g} and then \texttt{s} on the file. Staging it tells git the file is
   settled, and it moves to Staged changes.
\item When no file is left under Unmerged files, finish: \texttt{m c} commits a merge,
   \texttt{R c} goes on with a rebase, and \texttt{A c} with a cherry-pick.
\end{enumerate}

If you would sooner not deal with it now, \texttt{m a} stops a merge and puts the
branch and the files back as they were before it began. \texttt{R a} and \texttt{A a} do
the same for a rebase and a cherry-pick. Each asks first.

\begin{try}
\textbf{Try it, with a conflict made on purpose}
\begin{enumerate}
\item In the status of \texttt{\textasciitilde{}/recipes}, press \texttt{b c}, type \texttt{grams} and \texttt{RET} \texttt{RET}. In
   \texttt{scones.\codebreak{}txt} change the flour line to \texttt{300 g flour}. Save, press \texttt{C-x g},
   stage with \texttt{s} and commit with \texttt{c c} and the message \texttt{Flour in grams}.
\item Press \texttt{b b}, type \texttt{main} and \texttt{RET}. In \texttt{scones.\codebreak{}txt}, run
   \texttt{M-x revert-\codebreak{}buffer} if the line still says grams, then change it to
   \texttt{2 cups plain flour}. Save, stage and commit as \texttt{Say which flour}.
\item Press \texttt{m m}, type \texttt{grams} and \texttt{RET}. The merge stops. The status shows a
   \texttt{Doing:} line and \texttt{scones.\codebreak{}txt} under Unmerged files.
\item Press \texttt{RET} on the file. If the line says \texttt{2 cups plain flour} and has no
   markers, the buffer is the old one: run \texttt{M-x revert-\codebreak{}buffer}. Both versions
   are there between the markers. Press \texttt{C-c \textasciicircum{} l} to keep \texttt{300 g flour}, then
   \texttt{C-x C-s}.
\item Press \texttt{C-x g}, then \texttt{s} on the file, then \texttt{m c}. If a message buffer
   opens, press \texttt{C-c C-c}. The merge is done.
\end{enumerate}
\end{try}

\subsection{Rebase, cherry-pick, revert and reset}
\label{sec:422}
These four change or reuse commits that already exist. You can work for a long
time without them. They are here for the day you need one.

\subsubsection{Rebase}
\label{sec:418}
Rebasing makes your commits again on top of another commit, as if you had
started your work later than you did. The usual reason is that \texttt{main} moved on
while you worked on a branch. A merge would join the two with an extra commit.
A rebase gives a single straight line instead, which many people find easier
to read. \texttt{R e} asks what to rebase on, and \texttt{R u} rebases on the branch’s
upstream.

If one of your commits clashes with the new base, the rebase stops with a
conflict. Settle it as described above and press \texttt{R c} to go on. The full
menu’s \texttt{R s} leaves that commit out, and \texttt{R a} stops the whole rebase and
puts everything back.

\texttt{R i} rebases as a list says, and is the tool for tidying your own commits
before you share them. It asks which commit to start after, then shows the
commits after it in a buffer, one per line, the oldest first:

\begin{Verbatim}[breaklines,breakautoindent=false,breakindent=1em,breakafter=/-.\,;:=,commandchars=\\\{\}]
pick 81d04e7 Add the pavlova
pick 5c2e9a1 Fix a typo in the pavlova
pick 9b7f310 Add the bread
\end{Verbatim}

Change the list to say what you want. Move a line to change the order. Change
the word \texttt{pick} to \texttt{reword} to give the commit a new subject, to \texttt{squash} or
\texttt{fixup} to fold it into the commit on the line above, or to \texttt{drop} to leave it
out. In the example, changing the second line to \texttt{fixup} makes the typo fix
part of the commit that added the pavlova. \texttt{C-c C-c} rebases as the list says,
asking for each new subject first. \texttt{C-c C-k} gives the list up and nothing
changes.

One rule keeps you out of trouble: rebase only commits you have not pushed. A
rebase makes new commits in place of the old ones, and anyone who already has
the old ones would be left with a history that no longer matches yours.

\begin{xltabular}{\textwidth}{@{}l >{\raggedright\arraybackslash\hsize=0.756\hsize}X l >{\raggedright\arraybackslash\hsize=1.244\hsize}X@{}}
\toprule
\textbf{Modal} & \textbf{Non-modal} & \textbf{Command} & \textbf{What it does}\\
\midrule
\endhead
none & \texttt{R}, \texttt{C-c R} & \texttt{git-\codebreak{}rebase-\codebreak{}menu} & Open the rebase menu: rebase the branch on its upstream or a revision, or go on with, skip a commit of, or stop a rebase stopped half way.\\
none & \texttt{M-x git-\codebreak{}rebase-\codebreak{}onto} & \texttt{git-\codebreak{}rebase-\codebreak{}onto} & Rebase the branch checked out on a revision asked for: its commits since they parted made again after it.\\
none & \texttt{M-x git-\codebreak{}rebase-\codebreak{}upstream} & \texttt{git-\codebreak{}rebase-\codebreak{}upstream} & Rebase the branch checked out on its upstream: its own commits made again after the upstream’s.\\
none & \texttt{M-x git-\codebreak{}rebase-\codebreak{}interactive} & \texttt{git-\codebreak{}rebase-\codebreak{}interactive} & Rebase the commits after a revision asked for (the commit on point’s line first) as a list says.\\
none & \texttt{C-c C-c} & \texttt{git-\codebreak{}rebase-\codebreak{}todo-\codebreak{}finish} & Rebase as this list says: each reword line asks the commit’s new subject first (the rest of its message kept), then git rebases in the background; the list’s folder, the new messages in it, is deleted once git is done.\\
none & \texttt{C-c C-k} & \texttt{git-\codebreak{}rebase-\codebreak{}todo-\codebreak{}cancel} & Give up this rebase’s list: nothing is rebased.\\
none & \texttt{M-x git-\codebreak{}rebase-\codebreak{}continue} & \texttt{git-\codebreak{}rebase-\codebreak{}continue} & Go on with the rebase stopped half way, once the files git could not join are put right and staged.\\
none & \texttt{M-x git-\codebreak{}rebase-\codebreak{}skip} & \texttt{git-\codebreak{}rebase-\codebreak{}skip} & Skip the commit the rebase stopped at and go on with the next.\\
none & \texttt{M-x git-\codebreak{}rebase-\codebreak{}abort} & \texttt{git-\codebreak{}rebase-\codebreak{}abort} & Stop the rebase, after asking: the branch and the files go back to where they were before it began.\\
none & \texttt{M-x git-\codebreak{}rebase-\codebreak{}mode} & \texttt{git-\codebreak{}rebase-\codebreak{}mode} & Switch the current buffer to git-rebase-mode.\\
\bottomrule
\end{xltabular}

\subsubsection{Cherry-pick}
\label{sec:419}
A cherry-pick copies one commit from another branch onto the branch you are
on. You would use it when a branch holds one fix you want now and the rest of
the branch is not ready. Press \texttt{A A} in the status and name the commit. A
branch’s name stands for its newest commit, and a commit on the cursor’s line
is offered first. The full menu’s \texttt{A n} brings the commit’s changes into your
files without committing them.

\begin{xltabular}{\textwidth}{@{}l >{\raggedright\arraybackslash\hsize=0.619\hsize}X l >{\raggedright\arraybackslash\hsize=1.381\hsize}X@{}}
\toprule
\textbf{Modal} & \textbf{Non-modal} & \textbf{Command} & \textbf{What it does}\\
\midrule
\endhead
none & \texttt{A}, \texttt{C-c A} & \texttt{git-\codebreak{}cherry-\codebreak{}pick-\codebreak{}menu} & Open the cherry-pick menu: make a commit again on this branch, or bring its changes in uncommitted; go on with, skip or stop one stopped half way.\\
none & \texttt{M-x git-\codebreak{}cherry-\codebreak{}pick} & \texttt{git-\codebreak{}cherry-\codebreak{}pick} & Make a commit, asked for (the one on point’s line of a status or a log first, or another branch’s newest by the branch’s name), again on the branch checked out, with its message.\\
none & \texttt{M-x git-\codebreak{}cherry-\codebreak{}pick-\codebreak{}no-\codebreak{}commit} & \texttt{git-\codebreak{}cherry-\codebreak{}pick-\codebreak{}no-\codebreak{}commit} & Bring the changes of a commit, asked for, into the files and stage them, not committed.\\
none & \texttt{M-x git-\codebreak{}cherry-\codebreak{}pick-\codebreak{}continue} & \texttt{git-\codebreak{}cherry-\codebreak{}pick-\codebreak{}continue} & Go on with the cherry-pick stopped half way, once the files git could not join are put right and staged.\\
none & \texttt{M-x git-\codebreak{}cherry-\codebreak{}pick-\codebreak{}skip} & \texttt{git-\codebreak{}cherry-\codebreak{}pick-\codebreak{}skip} & Skip the commit the cherry-pick stopped at and go on with the next.\\
none & \texttt{M-x git-\codebreak{}cherry-\codebreak{}pick-\codebreak{}abort} & \texttt{git-\codebreak{}cherry-\codebreak{}pick-\codebreak{}abort} & Stop the cherry-pick, after asking: the branch and the files go back to where they were before it began.\\
\bottomrule
\end{xltabular}

\subsubsection{Revert}
\label{sec:420}
A revert makes a new commit that undoes an earlier one. The earlier commit
stays in the history, and the new one says that it was taken back. This is the
safe way to undo something you have already pushed, because nothing that
others have is rewritten. \texttt{\_ \_} asks which commit, offering the one on the
cursor’s line.

\begin{xltabular}{\textwidth}{@{}l >{\raggedright\arraybackslash\hsize=0.553\hsize}X l >{\raggedright\arraybackslash\hsize=1.447\hsize}X@{}}
\toprule
\textbf{Modal} & \textbf{Non-modal} & \textbf{Command} & \textbf{What it does}\\
\midrule
\endhead
none & \texttt{\_}, \texttt{C-c \_} & \texttt{git-\codebreak{}revert-\codebreak{}menu} & Open the revert menu: make a commit undoing another, or undo it in the files uncommitted; go on with, skip or stop one stopped half way.\\
none & \texttt{M-x git-\codebreak{}revert} & \texttt{git-\codebreak{}revert} & Make a commit undoing a commit asked for (the one on point’s line of a status or a log first), with git’s message saying so.\\
none & \texttt{M-x git-\codebreak{}revert-\codebreak{}no-\codebreak{}commit} & \texttt{git-\codebreak{}revert-\codebreak{}no-\codebreak{}commit} & Undo the changes of a commit, asked for, in the files and stage that, not committed.\\
none & \texttt{M-x git-\codebreak{}revert-\codebreak{}continue} & \texttt{git-\codebreak{}revert-\codebreak{}continue} & Go on with the revert stopped half way, once the files git could not join are put right and staged.\\
none & \texttt{M-x git-\codebreak{}revert-\codebreak{}skip} & \texttt{git-\codebreak{}revert-\codebreak{}skip} & Skip the commit the revert stopped at and go on with the next.\\
none & \texttt{M-x git-\codebreak{}revert-\codebreak{}abort} & \texttt{git-\codebreak{}revert-\codebreak{}abort} & Stop the revert, after asking: the branch and the files go back to where they were before it began.\\
\bottomrule
\end{xltabular}

\subsubsection{Reset}
\label{sec:421}
A reset moves the branch itself to another commit. The commits after it are
no longer on the branch. It is the way to take back commits you have not
pushed. The actions differ in what happens to your files.

\begin{xltabular}{\textwidth}{@{}l l >{\raggedright\arraybackslash\hsize=1.000\hsize}X@{}}
\toprule
\textbf{Modal} & \textbf{Non-modal} & \textbf{What it keeps}\\
\midrule
\endhead
none & \texttt{X m} & Your files as they are. The undone commits’ changes show as not staged.\\
none & \texttt{X s} & Your files and what is staged. The changes show as staged. (Full menu.)\\
none & \texttt{X K} & Changes you have not committed; the files otherwise follow. (Full menu.)\\
none & \texttt{X H} & Nothing. Files and index go to the commit, and uncommitted changes are lost. It asks first.\\
none & \texttt{X f} & Puts one file back to how a revision has it, after asking.\\
\bottomrule
\end{xltabular}

\texttt{X m} is the one to reach for when you committed too early: the commit is
gone and its changes are in your files to stage again as you like. \texttt{X H} is
for when you want to throw work away, and it does. The reflog can still find
the commits afterwards, but changes that were never committed are gone.

\begin{xltabular}{\textwidth}{@{}l l l >{\raggedright\arraybackslash\hsize=1.000\hsize}X@{}}
\toprule
\textbf{Modal} & \textbf{Non-modal} & \textbf{Command} & \textbf{What it does}\\
\midrule
\endhead
none & \texttt{X}, \texttt{C-c X} & \texttt{git-\codebreak{}reset-\codebreak{}menu} & Open the reset menu: move the branch to a revision, keeping what is staged, the files, or changes not committed, or nothing; or put a file back to a revision.\\
none & \texttt{M-x git-\codebreak{}reset-\codebreak{}mixed} & \texttt{git-\codebreak{}reset-\codebreak{}mixed} & Move the branch checked out to a revision asked for, what is staged put back to it too, the files as they are (their changes shown as not staged).\\
none & \texttt{M-x git-\codebreak{}reset-\codebreak{}soft} & \texttt{git-\codebreak{}reset-\codebreak{}soft} & Move the branch checked out to a revision asked for, what is staged and the files as they are (the commits’ changes shown as staged).\\
none & \texttt{M-x git-\codebreak{}reset-\codebreak{}keep} & \texttt{git-\codebreak{}reset-\codebreak{}keep} & Move the branch checked out to a revision asked for, the files with it, keeping changes not committed; git refuses when they would be lost.\\
none & \texttt{M-x git-\codebreak{}reset-\codebreak{}hard} & \texttt{git-\codebreak{}reset-\codebreak{}hard} & Move the branch checked out to a revision asked for, what is staged and the files with it, after asking: changes not committed are lost.\\
none & \texttt{M-x git-\codebreak{}reset-\codebreak{}file} & \texttt{git-\codebreak{}reset-\codebreak{}file} & Put a file, the one on point’s line of a status or the current one, back to how a revision asked for has it, staged, after asking: its changes not committed are lost.\\
\bottomrule
\end{xltabular}

\subsection{Pulling and pushing}
\label{sec:423}
A remote is a second copy of the repository that you exchange commits with. It
may be on a hosting service, on a server at work, or in a folder on a shared
drive. A repository you copied from somewhere with \texttt{git clone} already has
one, called \texttt{origin}. To give a repository of your own a remote, press \texttt{M a}
in the status. It asks for a name, and \texttt{origin} is the custom for the first,
then for the address.

Three verbs cover the exchange.

Fetch downloads the commits the remote has and you lack, and changes nothing
in your files or your branch. It is always safe. \texttt{f u} fetches from the
upstream’s remote and \texttt{f a} from every remote. Afterwards the second line of
the status says how far behind you are.

Pull fetches and then brings the new commits into your branch. \texttt{F u} pulls
from the upstream. With no switch it only moves your branch forward when that
is all there is to do. If you and the remote both have new commits, it stops
and tells you, and you choose: \texttt{- r} before \texttt{u} rebases your commits on
theirs, and \texttt{- m} merges theirs into yours. Tuatara never makes a merge you
did not ask for.

Push sends your commits to the remote. \texttt{P p} pushes the branch to its
upstream. The first time, the branch has no upstream yet: press \texttt{P}, turn on
\texttt{- u}, then press \texttt{e} and name the remote. From then on \texttt{P p} is enough. If
the remote has commits you lack, the push is refused. Pull first, then push
again.

Fetching, pulling and pushing run in the background, so you can go on working.
The sign \texttt{P} at the front of the mode line is lit while they run (chapter 11).

A forced push, \texttt{- f} or \texttt{- F}, replaces what the remote has with what you
have. It always asks first, since it can remove other people’s commits from
the remote. Of the two, \texttt{- f} is the careful one: it refuses if the remote
has something you have not seen. \texttt{- n} only tries, and reports what a push
would send.

Tuatara does not ask for passwords itself. With an address that uses ssh,
your keys come from the ssh agent of your system. If the server refuses,
Tuatara shows git’s own words.

The remote menu looks after the remotes themselves: \texttt{M a} adds one and \texttt{M u}
gives one a new address. The full menu also renames, removes, and forgets
branches that the remote has deleted.

\begin{try}
\textbf{Try it, with a folder standing in for a server}
\begin{enumerate}
\item In a terminal, make an empty repository to play the server:
   \texttt{git init --bare /tmp/\codebreak{}recipes-\codebreak{}server.\codebreak{}git}.
\item In the status of \texttt{\textasciitilde{}/recipes}, press \texttt{M a}. Press \texttt{RET} at the name to take
   \texttt{origin}, then type \texttt{/tmp/\codebreak{}recipes-\codebreak{}server.\codebreak{}git} and \texttt{RET} for the address.
\item Press \texttt{P}, then \texttt{- u}, then \texttt{e}. Type \texttt{origin} and \texttt{RET}. The echo area
   reports the push when it is done, and the status gains a line beginning
   \texttt{Upstream:}.
\item Make a small change to a recipe, save, stage and commit it. The second
   line of the status says \texttt{1 ahead} and a section lists the commit as not
   pushed.
\item Press \texttt{P p}. The section goes, and you are level with the remote again.
\end{enumerate}
\end{try}

\begin{xltabular}{\textwidth}{@{}l >{\raggedright\arraybackslash\hsize=0.533\hsize}X l >{\raggedright\arraybackslash\hsize=1.467\hsize}X@{}}
\toprule
\textbf{Modal} & \textbf{Non-modal} & \textbf{Command} & \textbf{What it does}\\
\midrule
\endhead
none & \texttt{M}, \texttt{C-c M} & \texttt{git-\codebreak{}remote-\codebreak{}menu} & Open the remote menu: add a remote, rename one, give it a new address, remove it, or forget its branches gone.\\
none & \texttt{M-x git-\codebreak{}remote-\codebreak{}add} & \texttt{git-\codebreak{}remote-\codebreak{}add} & Add a remote: a name asked for (origin when nothing is typed and there is none) and its address; its branches are fetched with f e.\\
none & \texttt{M-x git-\codebreak{}remote-\codebreak{}set-\codebreak{}address} & \texttt{git-\codebreak{}remote-\codebreak{}set-\codebreak{}address} & Give a remote, asked for, a new address, its own shown to begin with.\\
none & \texttt{M-x git-\codebreak{}remote-\codebreak{}rename} & \texttt{git-\codebreak{}remote-\codebreak{}rename} & Rename a remote, asked for, to a name asked for.\\
none & \texttt{M-x git-\codebreak{}remote-\codebreak{}remove} & \texttt{git-\codebreak{}remote-\codebreak{}remove} & Remove a remote, asked for, after asking: its branches here go with it.\\
none & \texttt{M-x git-\codebreak{}remote-\codebreak{}prune} & \texttt{git-\codebreak{}remote-\codebreak{}prune} & Forget the branches of a remote, asked for, that it no longer has.\\
\bottomrule
\end{xltabular}

\begin{xltabular}{\textwidth}{@{}l l l >{\raggedright\arraybackslash\hsize=1.000\hsize}X@{}}
\toprule
\textbf{Modal} & \textbf{Non-modal} & \textbf{Command} & \textbf{What it does}\\
\midrule
\endhead
none & \texttt{f}, \texttt{C-c f} & \texttt{git-\codebreak{}fetch-\codebreak{}menu} & Open the fetch menu: fetch from the upstream remote, a remote asked for, or every remote.\\
none & \texttt{M-x git-\codebreak{}fetch} & \texttt{git-\codebreak{}fetch} & Fetch the commits of the branch’s upstream remote (or origin) that this repository has not, the branches checked out left as they are.\\
none & \texttt{M-x git-\codebreak{}fetch-\codebreak{}remote} & \texttt{git-\codebreak{}fetch-\codebreak{}remote} & Fetch from a remote asked for.\\
none & \texttt{M-x git-\codebreak{}fetch-\codebreak{}all} & \texttt{git-\codebreak{}fetch-\codebreak{}all} & Fetch from every remote.\\
\bottomrule
\end{xltabular}

\begin{xltabular}{\textwidth}{@{}l l l >{\raggedright\arraybackslash\hsize=1.000\hsize}X@{}}
\toprule
\textbf{Modal} & \textbf{Non-modal} & \textbf{Command} & \textbf{What it does}\\
\midrule
\endhead
none & \texttt{F}, \texttt{C-c F} & \texttt{git-\codebreak{}pull-\codebreak{}menu} & Open the pull menu: bring in the upstream’s commits, or a branch’s of a remote asked for.\\
none & \texttt{M-x git-\codebreak{}pull} & \texttt{git-\codebreak{}pull} & Bring in the upstream’s commits: only when the branch can simply move on to them (fast-forward), unless the pull menu says to rebase (-r) or merge (-m) when it cannot.\\
none & \texttt{M-x git-\codebreak{}pull-\codebreak{}from} & \texttt{git-\codebreak{}pull-\codebreak{}from} & Bring in the commits of a branch, asked for, of a remote asked for, as M-x git-pull does.\\
\bottomrule
\end{xltabular}

\begin{xltabular}{\textwidth}{@{}l >{\raggedright\arraybackslash\hsize=0.571\hsize}X l >{\raggedright\arraybackslash\hsize=1.429\hsize}X@{}}
\toprule
\textbf{Modal} & \textbf{Non-modal} & \textbf{Command} & \textbf{What it does}\\
\midrule
\endhead
none & \texttt{P}, \texttt{C-c P} & \texttt{git-\codebreak{}push-\codebreak{}menu} & Open the push menu: push the branch to its upstream or a remote asked for, every tag or one; its switches force, set the upstream, try only, skip the hooks.\\
none & \texttt{M-x git-\codebreak{}push} & \texttt{git-\codebreak{}push} & Push the branch to its upstream, in the background; the push menu’s switches go with it (a forced push asked first).\\
none & \texttt{M-x git-\codebreak{}push-\codebreak{}elsewhere} & \texttt{git-\codebreak{}push-\codebreak{}elsewhere} & Push the branch checked out to a remote asked for, under its own name; -u in the push menu makes that its upstream.\\
none & \texttt{M-x git-\codebreak{}push-\codebreak{}tags} & \texttt{git-\codebreak{}push-\codebreak{}tags} & Push every tag to the upstream remote (or origin).\\
none & \texttt{M-x git-\codebreak{}push-\codebreak{}tag} & \texttt{git-\codebreak{}push-\codebreak{}tag} & Push a tag, asked for, to a remote asked for.\\
\bottomrule
\end{xltabular}

\subsection{Blame and the file menu}
\label{sec:426}
\subsubsection{Who changed each line}
\label{sec:424}
Blame answers the question “where did this line come from?”. \texttt{M-x git-\codebreak{}blame}
lists the current file a line at a time, and before each line the commit that
last changed it, its author and its date. \texttt{RET} or \texttt{l} on a line shows that
commit in full, so you can read the message and see what else changed with it.
\texttt{q} closes. The name is git’s and sounds harsher than the tool is. Most often
the answer is yourself, two years ago, with a good reason in the message.

Blame reads the file as saved, so save first.

\subsubsection{The file menu}
\label{sec:425}
\texttt{M-x git-\codebreak{}file-\codebreak{}menu} opens a menu of git actions for the file you are in. It
works from the file’s own buffer, without a visit to the status, and from a
file’s line in the status.

\begin{xltabular}{\textwidth}{@{}l l >{\raggedright\arraybackslash\hsize=1.000\hsize}X@{}}
\toprule
\textbf{Modal} & \textbf{Non-modal} & \textbf{What it does}\\
\midrule
\endhead
none & \texttt{s}, \texttt{u} & Stages the file, and unstages it.\\
none & \texttt{S}, \texttt{U} & Stages the hunk at the cursor, and unstages it.\\
none & \texttt{d} & Shows the file’s uncommitted changes.\\
none & \texttt{L} & Lists the commits that changed the file.\\
none & \texttt{b} & Blame.\\
none & \texttt{r} & Goes back through the file’s history (see Rewind below).\\
none & \texttt{f} & Puts the file back to how a revision has it, after asking.\\
none & \texttt{R} & Renames the file, in git and on disk. (Full menu.)\\
none & \texttt{D} & Deletes the file from git and from disk, after asking. (Full menu.)\\
\bottomrule
\end{xltabular}

Renaming a file through git, with \texttt{R} here, keeps its history joined to the
new name, and the buffers that visit the file follow it.

If you use the menu often, give it a key in your init file. Chapter 20 shows
how.

\begin{xltabular}{\textwidth}{@{}l l l >{\raggedright\arraybackslash\hsize=1.000\hsize}X@{}}
\toprule
\textbf{Modal} & \textbf{Non-modal} & \textbf{Command} & \textbf{What it does}\\
\midrule
\endhead
none & \texttt{M-x git-\codebreak{}file-\codebreak{}menu} & \texttt{git-\codebreak{}file-\codebreak{}menu} & Open the menu of the file you are in (or the one on point’s line of a status): stage or unstage it or its hunk at point, show its changes, its log or who changed each line, go back through its history, put it back to a revision, rename or delete it.\\
none & \texttt{M-x git-\codebreak{}blame} & \texttt{git-\codebreak{}blame} & List the current file (as saved), or the one on point’s line of a status, a line each with the commit that last changed it, its author and date; RET on a line shows its commit.\\
none & \texttt{l}, \texttt{RET}, \texttt{C-c l} & \texttt{git-\codebreak{}blame-\codebreak{}show} & Show the commit that last changed the line of the blame point is on.\\
none & \texttt{M-x git-\codebreak{}rename-\codebreak{}file} & \texttt{git-\codebreak{}rename-\codebreak{}file} & Rename the file you are in, or the one on point’s line of a status, to a path asked for within the repository (into a folder there, when it names one); the rename is staged, and the buffers visiting the file follow it.\\
none & \texttt{M-x git-\codebreak{}delete-\codebreak{}file} & \texttt{git-\codebreak{}delete-\codebreak{}file} & Delete the file you are in, or the one on point’s line of a status, from the files and from git, after asking; the deletion is staged.\\
\bottomrule
\end{xltabular}

\subsection{Rewind: a file as it was}
\label{sec:427}
Rewind shows a file as it stood at each commit that changed it, one step at a
time. It is the quickest way to answer “what did this recipe say before I
rewrote it?”, and to copy a paragraph back out of the past.

Press \texttt{C-x G} in a file. A new buffer shows the file as its latest commit left
it, read-only, with the cursor on the line you were on. Then:

\begin{longtable}{@{}l l l@{}}
\toprule
\textbf{Modal} & \textbf{Non-modal} & \textbf{What it does}\\
\midrule
\endhead
none & \texttt{p} & Shows the file one commit older.\\
none & \texttt{n} & Shows it one commit newer.\\
none & \texttt{r} & Lets you choose a commit from the file’s log.\\
none & \texttt{c} & Shows the commit that left the file like this.\\
none & \texttt{C-c w} & Copies that commit’s short hash.\\
none & \texttt{C-x G} & Goes back to the file as it is now, on the same line.\\
none & \texttt{q} & Closes the buffer.\\
\bottomrule
\end{longtable}

While you rewind, the mode line says \texttt{Rewind} and the commit’s hash in the
warning colour, followed by the commit’s date and which of how many it is,
such as \texttt{3/5}. You cannot mistake the old text for the file as it is now.

The text is read-only, but you can select and copy from it as from any
buffer (chapter 4) and paste into the real file. The history follows the file
back through a rename. Changes you have not committed are not part of it,
since they are in no commit.

Rewind works on folders too. Press \texttt{C-x G} in a folder listing (chapter 10)
and the listing shows the folder as its latest commit left it, with files
that have since been deleted. \texttt{p} and \texttt{n} step through the commits as before.
\texttt{RET} on a file there goes through that file’s history from that commit, and
\texttt{q} comes back to the folder. Commands that would change files, such as
copying or renaming, are refused in such a listing, because the files in it
are from the past.

\begin{try}
\textbf{Try it, with the recipes}
\begin{enumerate}
\item Open \texttt{\textasciitilde{}/recipes/\codebreak{}scones.\codebreak{}txt}, which by now has several commits. Put the
   cursor on the flour line.
\item Press \texttt{C-x G}. The mode line says \texttt{Rewind} and a hash. The text is the
   file as last committed.
\item Press \texttt{p} a few times. The flour line changes back through grams and cups.
   At the oldest commit, the echo area says there is none before it.
\item Press \texttt{c} to read the commit that made the version you are looking at,
   then \texttt{q} to close it.
\item Press \texttt{n} to come forward, then \texttt{C-x G}. You are back in the real file.
\end{enumerate}
\end{try}

\begin{xltabular}{\textwidth}{@{}l >{\raggedright\arraybackslash\hsize=0.622\hsize}X l >{\raggedright\arraybackslash\hsize=1.378\hsize}X@{}}
\toprule
\textbf{Modal} & \textbf{Non-modal} & \textbf{Command} & \textbf{What it does}\\
\midrule
\endhead
none & \texttt{C-x G} & \texttt{git-\codebreak{}rewind} & Go back through git’s history a commit at a time, and come back.\\
none & \texttt{M-x git-\codebreak{}rewind-\codebreak{}older} & \texttt{git-\codebreak{}rewind-\codebreak{}older} & Show the file or folder of this history as the commit before left it.\\
none & \texttt{M-x git-\codebreak{}rewind-\codebreak{}newer} & \texttt{git-\codebreak{}rewind-\codebreak{}newer} & Show the file or folder of this history as the commit after left it.\\
none & \texttt{M-x git-\codebreak{}rewind-to} & \texttt{git-\codebreak{}rewind-to} & Show the file or folder of this history as a commit chosen from its log left it, the latest first, by its hash, date, author or subject.\\
none & \texttt{M-x git-\codebreak{}rewind-\codebreak{}show-\codebreak{}commit} & \texttt{git-\codebreak{}rewind-\codebreak{}show-\codebreak{}commit} & Show the commit that left the file or folder of this history as it is shown: what it says and its changes.\\
none & \texttt{M-x git-\codebreak{}rewind-\codebreak{}copy-\codebreak{}hash} & \texttt{git-\codebreak{}rewind-\codebreak{}copy-\codebreak{}hash} & Copy the short hash of the commit that left the file or folder of this history as it is shown.\\
\bottomrule
\end{xltabular}

\begin{xltabular}{\textwidth}{@{}l >{\raggedright\arraybackslash\hsize=1.000\hsize}X@{}}
\toprule
\textbf{Face} & \textbf{What it paints}\\
\midrule
\endhead
\texttt{:mode-\codebreak{}line-\codebreak{}rewind} & Rewind and the commit shown, in the mode line of a file or folder shown as a commit left it (C-x G): lit up in the theme’s warning colour, so that it is seen at once that this is not the file as it is now.\\
\bottomrule
\end{xltabular}

\subsection{Repositories you only look at}
\label{sec:428}
Git can be told by a repository to run programs: hooks, filters and other
helpers named in the repository’s own settings. That is fine in a folder of
your own and unwise in one a stranger sent you. For this reason the status
asks before it first opens a project, and one of the answers is \texttt{l}, look
only.

In a look-only repository you can read the status, the log, diffs, blame and
rewind as usual. The mode line of every git buffer says \texttt{Look only}. The keys
that would change the repository are dimmed in the menus, and pressing one
says why it is refused. A repository that holds evidence is always look-only.

Chapter 12 explains the question, how look-only reading works, and
\texttt{M-x trust-\codebreak{}project}, which asks again when you have changed your mind.

\subsection{Less common work}
\label{sec:436}
The status has menus for parts of git that most people need seldom or never.
They work like the menus above, and \texttt{C-h} in each lists its actions. This
section says what each is for, so that you know it is there when the day
comes.

\subsubsection{Worktrees}
\label{sec:429}
A worktree is a second folder that shows another branch of the same
repository. With one, you can have \texttt{main} open in one folder and \texttt{metric} in
another at the same time, without checking out back and forth. \texttt{Z} opens the
menu: \texttt{Z c} makes a worktree for a new branch, \texttt{Z g} shows the status of one,
and \texttt{Z d} deletes one.

\begin{xltabular}{\textwidth}{@{}l >{\raggedright\arraybackslash\hsize=0.609\hsize}X l >{\raggedright\arraybackslash\hsize=1.391\hsize}X@{}}
\toprule
\textbf{Modal} & \textbf{Non-modal} & \textbf{Command} & \textbf{What it does}\\
\midrule
\endhead
none & \texttt{Z}, \texttt{C-c Z} & \texttt{git-\codebreak{}worktree-\codebreak{}menu} & Open the worktree menu: make a worktree for a branch, or for a new one, show a worktree’s status, move or delete one.\\
none & \texttt{M-x git-\codebreak{}worktree-\codebreak{}branch} & \texttt{git-\codebreak{}worktree-\codebreak{}branch} & Make a new branch, asked for, starting at a revision asked for, and a worktree for it in a folder asked for.\\
none & \texttt{M-x git-\codebreak{}worktree-\codebreak{}checkout} & \texttt{git-\codebreak{}worktree-\codebreak{}checkout} & Make a worktree for a branch asked for, in a folder asked for: the branch checked out there, to work on beside this one.\\
none & \texttt{M-x git-\codebreak{}worktree-\codebreak{}status} & \texttt{git-\codebreak{}worktree-\codebreak{}status} & Show the status of a worktree, asked for, as C-x g shows a repository’s.\\
none & \texttt{M-x git-\codebreak{}worktree-\codebreak{}move} & \texttt{git-\codebreak{}worktree-\codebreak{}move} & Move a worktree, asked for, not the main one, to a folder asked for.\\
none & \texttt{M-x git-\codebreak{}worktree-\codebreak{}delete} & \texttt{git-\codebreak{}worktree-\codebreak{}delete} & Delete a worktree, asked for, not the main one, after asking: its folder goes, its branch stays.\\
\bottomrule
\end{xltabular}

\subsubsection{Bisect: finding the commit that broke something}
\label{sec:430}
Suppose the website looked right last month and looks wrong today, and fifty
commits lie between. Bisect finds the one at fault by halving. \texttt{B B} asks for
a commit that is bad and one that was good. Git checks out a commit in the
middle. You look, and say \texttt{B g} for good or \texttt{B b} for bad, and git halves
again. After about six answers it names the commit. \texttt{B r} ends the search
and returns you to where you were. If a command can tell good from bad, the
full menu’s \texttt{B x} lets it answer for you.

\begin{xltabular}{\textwidth}{@{}l l l >{\raggedright\arraybackslash\hsize=1.000\hsize}X@{}}
\toprule
\textbf{Modal} & \textbf{Non-modal} & \textbf{Command} & \textbf{What it does}\\
\midrule
\endhead
none & \texttt{B}, \texttt{C-c B} & \texttt{git-\codebreak{}bisect-\codebreak{}menu} & Open the bisect menu: look for the commit that brought a fault in, saying of each commit tried whether it is good or bad, or letting a command say.\\
none & \texttt{M-x git-\codebreak{}bisect-\codebreak{}start} & \texttt{git-\codebreak{}bisect-\codebreak{}start} & Begin looking for the commit that brought a fault in: a revision that has it (HEAD when nothing is typed) and one that has not are asked for, and git checks out one between to try; B g, B b and B s then say it is good, bad, or not to be tried.\\
none & \texttt{M-x git-\codebreak{}bisect-\codebreak{}good} & \texttt{git-\codebreak{}bisect-\codebreak{}good} & Say the commit checked out has not the fault: git tries one after it.\\
none & \texttt{M-x git-\codebreak{}bisect-\codebreak{}bad} & \texttt{git-\codebreak{}bisect-\codebreak{}bad} & Say the commit checked out has the fault: git tries one before it.\\
none & \texttt{M-x git-\codebreak{}bisect-\codebreak{}skip} & \texttt{git-\codebreak{}bisect-\codebreak{}skip} & Say the commit checked out cannot be tried: git tries one beside it.\\
none & \texttt{M-x git-\codebreak{}bisect-\codebreak{}run} & \texttt{git-\codebreak{}bisect-\codebreak{}run} & Let a command asked for, run by the shell in the repository, say of each commit tried whether it has the fault (it fails) or not (it succeeds, 125 for one not to be tried), until the first bad commit is found.\\
none & \texttt{M-x git-\codebreak{}bisect-\codebreak{}reset} & \texttt{git-\codebreak{}bisect-\codebreak{}reset} & Stop looking: HEAD goes back to where it was when the bisection began.\\
\bottomrule
\end{xltabular}

\subsubsection{Ignoring files}
\label{sec:431}
Some files do not belong in the history: backups, files a program makes as it
runs, anything private. The ignore menu, \texttt{i}, adds a file or a pattern such as
\texttt{*.bak} to the list of what git leaves alone. \texttt{i t} writes it in the
\texttt{.gitignore} file at the top of the repository, which is shared with everyone
who has the repository. \texttt{i p} keeps the rule on your computer only. The full
menu’s \texttt{i s} writes it in the \texttt{.gitignore} of the current folder. Ignored
files are listed in the last section of the status.

\begin{xltabular}{\textwidth}{@{}l >{\raggedright\arraybackslash\hsize=0.490\hsize}X l >{\raggedright\arraybackslash\hsize=1.510\hsize}X@{}}
\toprule
\textbf{Modal} & \textbf{Non-modal} & \textbf{Command} & \textbf{What it does}\\
\midrule
\endhead
none & \texttt{i}, \texttt{C-c i} & \texttt{git-\codebreak{}ignore-\codebreak{}menu} & Open the ignore menu: have git ignore a file or a pattern, for everyone (.gitignore at the top or in the file’s folder) or here alone.\\
none & \texttt{M-x git-\codebreak{}ignore-\codebreak{}shared} & \texttt{git-\codebreak{}ignore-\codebreak{}shared} & Have git ignore a pattern asked for (the file on point’s line, or the current one, offered), in the repository’s .gitignore, committed with it.\\
none & \texttt{M-x git-\codebreak{}ignore-\codebreak{}folder} & \texttt{git-\codebreak{}ignore-\codebreak{}folder} & Have git ignore a pattern asked for in the .gitignore of the current file’s folder.\\
none & \texttt{M-x git-\codebreak{}ignore-\codebreak{}private} & \texttt{git-\codebreak{}ignore-\codebreak{}private} & Have git ignore a pattern asked for here alone, in git’s own exclude file, which is never committed.\\
\bottomrule
\end{xltabular}

\subsubsection{Notes on commits}
\label{sec:432}
A note is a remark attached to a commit after the fact, without changing the
commit. \texttt{T T} writes one, for the commit on the cursor’s line or one you
name.

\begin{xltabular}{\textwidth}{@{}l l l >{\raggedright\arraybackslash\hsize=1.000\hsize}X@{}}
\toprule
\textbf{Modal} & \textbf{Non-modal} & \textbf{Command} & \textbf{What it does}\\
\midrule
\endhead
none & \texttt{T}, \texttt{C-c T} & \texttt{git-\codebreak{}notes-\codebreak{}menu} & Open the notes menu: write a note on a commit, take one away, or prune those on commits gone.\\
none & \texttt{M-x git-\codebreak{}notes-\codebreak{}edit} & \texttt{git-\codebreak{}notes-\codebreak{}edit} & Write a note on a commit asked for (the one on point’s line first, else HEAD), its note so far shown to begin with: a note says more of a commit without changing it.\\
none & \texttt{M-x git-\codebreak{}notes-\codebreak{}remove} & \texttt{git-\codebreak{}notes-\codebreak{}remove} & Take away the note on a commit asked for.\\
none & \texttt{M-x git-\codebreak{}notes-\codebreak{}prune} & \texttt{git-\codebreak{}notes-\codebreak{}prune} & Take away the notes on commits the repository no longer has.\\
\bottomrule
\end{xltabular}

\subsubsection{Patches and bundles}
\label{sec:433}
A patch is a commit written out as a file, to send to someone by mail or on a
stick. In the patch menu, \texttt{W c} writes commits as patch files and \texttt{W a}
applies a patch file to your files. The full menu’s \texttt{W s} saves your
uncommitted changes as a patch. The menu under \texttt{w} takes patch files in as
commits, with their authors and messages: \texttt{w w} starts, and \texttt{w c} and \texttt{w a}
go on or stop when one does not fit.

A bundle is a whole repository in a single file, for carrying to a computer
that has no network. \texttt{W b} writes one and \texttt{W v} checks one, both in the full
menu.

\begin{xltabular}{\textwidth}{@{}l >{\raggedright\arraybackslash\hsize=0.480\hsize}X l >{\raggedright\arraybackslash\hsize=1.520\hsize}X@{}}
\toprule
\textbf{Modal} & \textbf{Non-modal} & \textbf{Command} & \textbf{What it does}\\
\midrule
\endhead
none & \texttt{W}, \texttt{C-c W} & \texttt{git-\codebreak{}patch-\codebreak{}menu} & Open the patch menu: write commits as patch files, apply a patch file, save the changes not committed as one, or write or check a bundle.\\
none & \texttt{M-x git-\codebreak{}format-\codebreak{}patch} & \texttt{git-\codebreak{}format-\codebreak{}patch} & Write the commits after a revision asked for as patch files, a commit a file, in a folder asked for; the patch menu’s switches number them, add a cover letter, or sign them off (those for applying are left out).\\
none & \texttt{M-x git-\codebreak{}apply-\codebreak{}patch} & \texttt{git-\codebreak{}apply-\codebreak{}patch} & Apply a patch file asked for to the files, as git apply does; the patch menu’s switches stage it too (--index), only stage it (--cached), or join it three ways where it does not fit (those for writing patches are left out).\\
none & \texttt{M-x git-\codebreak{}save-\codebreak{}patch} & \texttt{git-\codebreak{}save-\codebreak{}patch} & Save every change not committed, staged or not, as a new patch file asked for, its octets as git writes them.\\
none & \texttt{M-x git-\codebreak{}bundle-\codebreak{}create} & \texttt{git-\codebreak{}bundle-\codebreak{}create} & Write the whole repository, every branch and tag, as a single bundle file asked for: one file to carry it, from which it can be cloned.\\
none & \texttt{M-x git-\codebreak{}bundle-\codebreak{}verify} & \texttt{git-\codebreak{}bundle-\codebreak{}verify} & Check a bundle file asked for: whether this repository has what it needs to take it in.\\
\bottomrule
\end{xltabular}

\begin{xltabular}{\textwidth}{@{}l l l >{\raggedright\arraybackslash\hsize=1.000\hsize}X@{}}
\toprule
\textbf{Modal} & \textbf{Non-modal} & \textbf{Command} & \textbf{What it does}\\
\midrule
\endhead
none & \texttt{w}, \texttt{C-c w} & \texttt{git-\codebreak{}am-\codebreak{}menu} & Open the menu of patches by mail: take patch files in as commits, or go on with, skip one of, or stop a taking in stopped half way.\\
none & \texttt{M-x git-am} & \texttt{git-am} & Take in patch files asked for, as sent by mail (git format-patch writes them), each made a commit with its author and message; one that does not fit stops it half way (w c goes on, w s skips it, w a stops it).\\
none & \texttt{M-x git-\codebreak{}am-\codebreak{}continue} & \texttt{git-\codebreak{}am-\codebreak{}continue} & Go on with the taking in of patches stopped half way, once the files git could not join are put right and staged.\\
none & \texttt{M-x git-\codebreak{}am-\codebreak{}skip} & \texttt{git-\codebreak{}am-\codebreak{}skip} & Skip the commit the taking in of patches stopped at and go on with the next.\\
none & \texttt{M-x git-\codebreak{}am-\codebreak{}abort} & \texttt{git-\codebreak{}am-\codebreak{}abort} & Stop the taking in of patches, after asking: the branch and the files go back to where they were before it began.\\
\bottomrule
\end{xltabular}

\subsubsection{Repositories inside repositories}
\label{sec:434}
A submodule is another repository kept in a folder of yours at one exact
commit, such as a theme that your website uses and someone else maintains.
\texttt{o} opens its menu: \texttt{o a} adds one, and \texttt{o u} checks them all out after you
have cloned a repository that has some.

A subtree also brings another repository into a folder, but copies its files
into your own history, so that people who clone yours need nothing extra. \texttt{O}
opens its menu.

A sparse checkout goes the other way. It shows only some folders of a large
repository in your files and leaves the rest in the history. \texttt{>} opens its
menu: \texttt{> s} chooses the folders and \texttt{> d} brings everything back.

\begin{xltabular}{\textwidth}{@{}l >{\raggedright\arraybackslash\hsize=0.681\hsize}X l >{\raggedright\arraybackslash\hsize=1.319\hsize}X@{}}
\toprule
\textbf{Modal} & \textbf{Non-modal} & \textbf{Command} & \textbf{What it does}\\
\midrule
\endhead
none & \texttt{o}, \texttt{C-c o} & \texttt{git-\codebreak{}submodule-\codebreak{}menu} & Open the submodule menu: add a submodule, check them out, take their addresses again, remove one or show its status.\\
none & \texttt{M-x git-\codebreak{}submodule-\codebreak{}add} & \texttt{git-\codebreak{}submodule-\codebreak{}add} & Add a repository, its address asked for, as a submodule at a path asked for (named after the address when nothing is typed), cloned there.\\
none & \texttt{M-x git-\codebreak{}submodule-\codebreak{}update} & \texttt{git-\codebreak{}submodule-\codebreak{}update} & Check every submodule out at the commit this repository records (cloning those not checked out yet), and theirs in turn; with -r in the submodule menu at the newest commit of the branch each follows instead.\\
none & \texttt{M-x git-\codebreak{}submodule-\codebreak{}sync} & \texttt{git-\codebreak{}submodule-\codebreak{}sync} & Take each submodule’s address again from .gitmodules, after it changed there.\\
none & \texttt{M-x git-\codebreak{}submodule-\codebreak{}status} & \texttt{git-\codebreak{}submodule-\codebreak{}status} & Show the status of a submodule, asked for (the one on point’s line first), as C-x g shows a repository’s.\\
none & \texttt{M-x git-\codebreak{}submodule-\codebreak{}remove} & \texttt{git-\codebreak{}submodule-\codebreak{}remove} & Remove a submodule, asked for, after asking: its folder and its entry in .gitmodules go, the removal staged.\\
\bottomrule
\end{xltabular}

\begin{xltabular}{\textwidth}{@{}l >{\raggedright\arraybackslash\hsize=0.520\hsize}X l >{\raggedright\arraybackslash\hsize=1.480\hsize}X@{}}
\toprule
\textbf{Modal} & \textbf{Non-modal} & \textbf{Command} & \textbf{What it does}\\
\midrule
\endhead
none & \texttt{O}, \texttt{C-c O} & \texttt{git-\codebreak{}subtree-\codebreak{}menu} & Open the subtree menu: bring another repository into a folder, pull its newer commits, push the folder’s, or split the folder’s history off.\\
none & \texttt{M-x git-\codebreak{}subtree-\codebreak{}add} & \texttt{git-\codebreak{}subtree-\codebreak{}add} & Bring another repository’s branch into a new folder of this one as a subtree, its history joined in (one commit with -s in the subtree menu).\\
none & \texttt{M-x git-\codebreak{}subtree-\codebreak{}pull} & \texttt{git-\codebreak{}subtree-\codebreak{}pull} & Bring the newer commits of another repository’s branch into the subtree folder.\\
none & \texttt{M-x git-\codebreak{}subtree-\codebreak{}push} & \texttt{git-\codebreak{}subtree-\codebreak{}push} & Send the commits that changed the subtree folder to another repository’s branch.\\
none & \texttt{M-x git-\codebreak{}subtree-\codebreak{}split} & \texttt{git-\codebreak{}subtree-\codebreak{}split} & Make a history of the commits that changed a folder asked for, as if it were a repository of its own, and say the commit it ends at (a branch of it with =b in the subtree menu).\\
\bottomrule
\end{xltabular}

\begin{xltabular}{\textwidth}{@{}l >{\raggedright\arraybackslash\hsize=0.490\hsize}X l >{\raggedright\arraybackslash\hsize=1.510\hsize}X@{}}
\toprule
\textbf{Modal} & \textbf{Non-modal} & \textbf{Command} & \textbf{What it does}\\
\midrule
\endhead
none & \texttt{>}, \texttt{C-c >} & \texttt{git-\codebreak{}sparse-\codebreak{}menu} & Open the sparse checkout menu: check out only some folders, more of them, everything again, or narrow it again.\\
none & \texttt{M-x git-\codebreak{}sparse-\codebreak{}set} & \texttt{git-\codebreak{}sparse-\codebreak{}set} & Check out only folders asked for (and the files at the top), the rest of the repository kept out of the files.\\
none & \texttt{M-x git-\codebreak{}sparse-\codebreak{}add} & \texttt{git-\codebreak{}sparse-\codebreak{}add} & Check out folders asked for too, in a sparse checkout.\\
none & \texttt{M-x git-\codebreak{}sparse-\codebreak{}reapply} & \texttt{git-\codebreak{}sparse-\codebreak{}reapply} & Keep out of the files again what the sparse checkout leaves out, after a merge or a rebase brought some in.\\
none & \texttt{M-x git-\codebreak{}sparse-\codebreak{}disable} & \texttt{git-\codebreak{}sparse-\codebreak{}disable} & Check the whole repository out again.\\
\bottomrule
\end{xltabular}

\subsubsection{Any git command}
\label{sec:435}
For anything the menus do not have, \texttt{!} in the status asks for the words you
would type after \texttt{git} in a terminal, such as \texttt{log --oneline -5}, runs git
with them in the repository and shows what it printed. No shell is involved,
so the words go to git exactly as you type them.

\begin{xltabular}{\textwidth}{@{}l l l >{\raggedright\arraybackslash\hsize=1.000\hsize}X@{}}
\toprule
\textbf{Modal} & \textbf{Non-modal} & \textbf{Command} & \textbf{What it does}\\
\midrule
\endhead
none & \texttt{!}, \texttt{C-c !} & \texttt{git-\codebreak{}run-\codebreak{}command} & Run git with the words asked for (log --oneline -5, say; quotes group words as in a shell) in the repository at hand, no shell between, and show what it printed; for what the menus have no action for.\\
\bottomrule
\end{xltabular}

Git is one of several things that work on files and text in Tuatara. The next
chapter, Data, pictures and PDFs, turns to files that are not plain text.

\section{Data, pictures and PDFs}
\label{sec:465}
Looking at files that are more than plain text.

Most of this manual is about text you write yourself. This chapter is about
files that come from somewhere else: a table of figures saved from a
spreadsheet, a photo from a camera, a form someone scanned and sent as a PDF, a
Word document. Tuatara opens all of them with \texttt{C-x C-f}, the same key as for
any file, and shows each in the way that suits it. A table gets a colour for
each column. A picture and a PDF are drawn. A file that is neither text nor
anything Tuatara can draw is shown as the raw numbers it is made of, in the hex
view.

The chapter starts with tables of data, because those you can also change:
filter, sort and tidy. Pictures, PDFs and Office documents are for looking at
only. Nothing in this chapter writes to them.

One reminder about the state keys. \texttt{i} starts INSERT state, and \texttt{C-g} or \texttt{ESC}
returns to NORMAL state. The single-letter keys in this chapter are pressed in
NORMAL state. Keys that start with \texttt{C-x} or \texttt{C-c} work in both.

\subsection{Delimited data: CSV and TSV}
\label{sec:444}
A spreadsheet program can save a sheet as plain text. Each row becomes a line,
and the cells of the row are written one after another with a comma between
them. That is a CSV file, for comma-separated values. A TSV file is the same
with a tab between the cells. Banks, shops, membership systems and nearly every
program that holds a table can write one. The character between the cells is
called the separator, a line is called a record, and a cell is called a field.

Such a file is easy for a program to read and hard for a person. The fields
are of different lengths, so the columns do not line up, and by the sixth comma
of a long line you no longer know which column you are in. Tuatara helps in two
ways. It gives each column its own colour, so that your eye can follow a column
down the screen. And the mode line says which column the cursor is in.

A file whose name ends in \texttt{.csv} opens in CSV mode. One that ends in \texttt{.tsv} or
\texttt{.tab} opens in TSV mode, which is CSV mode with the tab as the separator. For
a file with another ending, \texttt{M-x csv-\codebreak{}mode} or \texttt{M-x tsv-\codebreak{}mode} turns the mode on
by hand. Lines in these modes are never wrapped, because a record cut in two
would break the columns. A long record runs off the right edge of the window,
and the window scrolls sideways as the cursor moves along it.

\subsubsection{Columns are counted from 0}
\label{sec:439}
Figure 16.1 shows the membership list of a tramping club. The cursor is on
\texttt{Taupo}, in the second column.

\begin{figure}[H]
\caption{Figure 16.1. A CSV file, with the cursor’s column in the mode line.}\label{org438}
\begin{Verbatim}[breaklines,breakautoindent=false,breakindent=1em,breakafter=/-.\,;:=,commandchars=\\\{\}]
 name,town,joined
 Aroha Ngata,Wanaka,2019
 Tama Parata,Taupo,2021
 Mere Walker,Wanaka,2023
 Hemi Tane,Napier,2020

 column 0    column 1  column 2

 TUATARA[N] [-] members.csv  L3 C12  CSV  col 1 town
\end{Verbatim}
\end{figure}

The mode line says \texttt{CSV}, then \texttt{col} and the number of the column, then that
column’s name. The number starts at 0: the first column is col 0, the second
is col 1. The name comes from the first line of the file, which Tuatara takes
as the header. Wherever this chapter asks you for a column, you can answer with
the name or with the number, and the number is the one the mode line shows.

There are eight colours. In a table with more than eight columns, column 8 has
the colour of column 0 again, column 9 that of column 1, and so on. If you find
the colours too much, \texttt{M-x rotate-\codebreak{}colouring} goes on to a plainer style that
colours only the separators and the quoted fields, and then to no colour at
all. Chapter 13 describes these styles. To choose one for every CSV file, put
\texttt{(set-\codebreak{}colouring 'csv-\codebreak{}mode :sparse)} in your init file.

A field may hold the separator itself. An address such as \texttt{12 Lake Road,
Wanaka} has a comma in it. Programs write such a field between double quotes,
and Tuatara reads it as one field. A quoted field may even run over several
lines, and the mode line still names the right column on each of them.

\subsubsection{Moving in a table}
\label{sec:440}
There are no special keys for moving in a table. The keys of chapter 2 work as
in any text: \texttt{h} and \texttt{l} move along a record, \texttt{j} and \texttt{k} move from record to
record, and a search with \texttt{C-s} goes to a value. What changes is that you
always know where you are. Watch the mode line as you move along a line, and
the column number and name change each time you pass a separator.

To follow one column down a long table, turn on the highlight for the cursor’s
column with \texttt{C-x T c}, described in chapter 1. It only helps where the fields
above and below are about the same width.

\subsubsection{The separator}
\label{sec:441}
Tuatara works out the separator from the first line of the file. It counts the
commas, semicolons, tabs and pipes (\texttt{|}) there and takes the one there are most
of. Semicolons are common in files from countries that write a decimal comma.
You do not have to say anything for those files to work.

Now and then the guess is wrong, or the file uses something else, such as a
colon. \texttt{C-c C-s} asks for the separator of the buffer you are in. It offers
comma, semicolon, tab, pipe, colon and space by name, and \texttt{detect}, which makes
the guess again. Any other single character can be typed as itself. The echo
area then says, for example, \texttt{Fields separated by semicolon}, and the colours
change to match.

If all your files use one separator, set \texttt{*csv-\codebreak{}separator*} in your init file,
for example to \texttt{\#\textbackslash{};} for the semicolon. Tuatara then stops guessing.

\subsubsection{The header}
\label{sec:442}
The first line of most CSV files names the columns. Tuatara treats it as the
header: its names are shown in the mode line and offered when you are asked for
a column, a filter never hides it, and a sort leaves it where it is.

In a long table the header scrolls out of sight after the first screen.
\texttt{C-x / L} keeps it at the top of the window while the rest scrolls under it,
in the same colours as the columns below. Press it again to let the header
scroll away. To have every table open that way, set \texttt{*list-\codebreak{}header-\codebreak{}lock*} to
\texttt{t}.

Some files have no header. Their first line is a record like the others, and
you do not want it left out of a sort. \texttt{C-x / t} tells Tuatara that the first
line of this buffer is a record. The mode line then shows only the column
number, and columns are offered by number. Press it again to make the first
line the header once more. If none of your files has a header, set
\texttt{*csv-\codebreak{}header*} to \texttt{nil}.

\begin{try}
\textbf{Try it, with a club’s members}
\begin{enumerate}
\item Press \texttt{C-x C-f}, type \texttt{\textasciitilde{}/members.\codebreak{}csv} and press \texttt{RET}. Press \texttt{i} and type
   the five lines of figure 16.1, each ended with \texttt{RET}. Press \texttt{C-g} and save
   with \texttt{C-x C-s}.
\item The fields took their colours as soon as you typed the first comma. Press
   \texttt{M-<} (modal: \texttt{g g}) to go to the top. The mode line says \texttt{col 0 name}.
\item Press \texttt{j} twice, then \texttt{l} until the cursor is on \texttt{Taupo}. The mode line
   says \texttt{col 1 town}.
\item Press \texttt{C-c C-s}, type \texttt{semicolon} and press \texttt{RET}. There are no semicolons
   in the file, so each line is now one field and has one colour.
\item Press \texttt{C-c C-s} again, type \texttt{detect} and press \texttt{RET}. The echo area says
   \texttt{Fields separated by comma} and the columns are back.
\item Press \texttt{C-x / t}. The echo area says the first line is a record like the
   rest, and the mode line shows \texttt{col 1} with no name. Press \texttt{C-x / t} again
   to make it the header.
\end{enumerate}
\end{try}

\subsubsection{Commands, settings and faces for tables}
\label{sec:443}
\begin{xltabular}{\textwidth}{@{}l l l >{\raggedright\arraybackslash\hsize=1.000\hsize}X@{}}
\toprule
\textbf{Modal} & \textbf{Non-modal} & \textbf{Command} & \textbf{What it does}\\
\midrule
\endhead
none & \texttt{M-x csv-\codebreak{}mode} & \texttt{csv-\codebreak{}mode} & Switch the current buffer to csv-mode.\\
none & \texttt{M-x tsv-\codebreak{}mode} & \texttt{tsv-\codebreak{}mode} & Switch the current buffer to tsv-mode.\\
none & \texttt{C-c C-s} & \texttt{csv-\codebreak{}define-\codebreak{}separator} & Ask for the character that separates the fields of this buffer and show them separated by it.\\
none & \texttt{C-x / t} & \texttt{list-\codebreak{}toggle-\codebreak{}header} & Take the first line of a file of records as its header, or as a record like the rest; the same again switches back; with a count, a header when it is positive, a record when it is zero or less.\\
none & \texttt{C-x / L} & \texttt{list-\codebreak{}toggle-\codebreak{}header-\codebreak{}lock} & Keep the first line of a file of records, its header, at the top of the window as the rest scrolls under it, coloured as it is below; the same again lets it scroll away; with a count, kept when it is positive, let go when it is zero or less.\\
\bottomrule
\end{xltabular}

\begin{xltabular}{\textwidth}{@{}l l >{\raggedright\arraybackslash\hsize=1.000\hsize}X@{}}
\toprule
\textbf{Setting} & \textbf{Standard} & \textbf{What it changes}\\
\midrule
\endhead
\texttt{*csv-\codebreak{}separator*} & \texttt{:detect} & What separates the fields of the files CSV mode shows.\\
\texttt{*csv-\codebreak{}header*} & \texttt{t} & True when the first line of the files CSV mode shows is their header, naming the columns: always shown by search filters, never deleted by M-x list-keep-records, its names offered for a column.\\
\texttt{*list-\codebreak{}header-\codebreak{}lock*} & \texttt{nil} & True to keep the header of a file of records at the top of the window as the rest scrolls under it, wherever there is one: a CSV’s first line, a log’s once M-x list-toggle-header makes it a header; as M-x list-toggle-header-lock does for one buffer.\\
\bottomrule
\end{xltabular}

\begin{xltabular}{\textwidth}{@{}l >{\raggedright\arraybackslash\hsize=1.000\hsize}X@{}}
\toprule
\textbf{Face} & \textbf{What it paints}\\
\midrule
\endhead
\texttt{:csv-\codebreak{}column-0} & Delimited data: columns 0, 8, 16 \ldots{}\\
\texttt{:csv-\codebreak{}column-1} & Delimited data: columns 1, 9, 17 \ldots{}\\
\texttt{:csv-\codebreak{}column-2} & Delimited data: columns 2, 10, 18 \ldots{}\\
\texttt{:csv-\codebreak{}column-3} & Delimited data: columns 3, 11, 19 \ldots{}\\
\texttt{:csv-\codebreak{}column-4} & Delimited data: columns 4, 12, 20 \ldots{}\\
\texttt{:csv-\codebreak{}column-5} & Delimited data: columns 5, 13, 21 \ldots{}\\
\texttt{:csv-\codebreak{}column-6} & Delimited data: columns 6, 14, 22 \ldots{}\\
\texttt{:csv-\codebreak{}column-7} & Delimited data: columns 7, 15, 23 \ldots{}\\
\texttt{:csv-\codebreak{}separator} & Delimited data coloured sparsely (set-colouring): the separators between fields.\\
\bottomrule
\end{xltabular}

\subsection{Showing only some records}
\label{sec:446}
A table is most useful when you can ask it a question. Which members live in
Wanaka? What did the shop sell in March? Chapter 6 introduced search filters,
which narrow a list to the entries that match a pattern. In a table they do one
thing more: they can look at a single column. A search for \texttt{2021} in the whole
line would also find a phone number that happens to hold those digits. A filter
on the \texttt{joined} column finds only the people who joined that year.

\texttt{C-x / f} shows only the records that match. It asks two questions. The first
is the column, and each column is offered with the value it has in the first
record beside it, so you can see what it holds. The second is the pattern to
look for in that column. Records that do not match are hidden. They are still
in the file, and saving writes all of them.

You can name the column in several ways.

\begin{longtable}{@{}l l@{}}
\toprule
\textbf{You type} & \textbf{The pattern is tried on}\\
\midrule
\endhead
a name, such as \texttt{town} & The column the header gives that name.\\
\texttt{\$1} & Column 1, the second one. \texttt{\$0} is the first.\\
\texttt{\$*} or nothing at all & The whole record, as \texttt{C-x / s} does.\\
\bottomrule
\end{longtable}

The pattern follows the rules of search in chapter 6. Small letters match
capitals too, unless you type a capital yourself. Quotes around a field are
taken off before the pattern is tried, so you never have to type them.

\texttt{C-x / h} is the opposite. It hides the records that match and shows the rest.
Use it to get things out of the way, such as the cancelled orders.

While a filter is on, the mode line says by what the table is filtered and how
many records are left, as in \texttt{Filtered by kiwi: 12 of 1,000 records}. A second
filter narrows what the first left, so you can ask for Wanaka and then, of
those, for 2023. \texttt{C-x / l} lists the filters, and \texttt{d} on a line there takes
that one away. \texttt{C-x / r} takes them all away and shows the whole table. The
header is shown whatever the filters say.

You can edit the records you see. A change that would reach into a hidden
record is refused, so a filter never changes what you cannot see.

\subsubsection{Deleting records for good}
\label{sec:445}
Sometimes you want a smaller file, to send someone only their own town’s
members. \texttt{C-x / k} asks the same two questions and then deletes every record
that does not match, keeping the ones that do. \texttt{C-x / d} deletes the ones that
match. Both act on the whole table, or on the lines you have selected if there
is a selection. The header always stays. Each is one step for \texttt{C-/} (modal: \texttt{u})
to undo, and nothing reaches the file until you save.

When you are not sure of a pattern, try it with \texttt{C-x / f} first. You see what
would be kept, and nothing is lost if you got it wrong.

The same keys work in other files that hold one record on each line, such as a
log. The setting \texttt{*line-\codebreak{}record-\codebreak{}modes*} names the modes that count. A log has
no header until \texttt{C-x / t} makes its first line one, and its columns are its
words, apart by spaces.

\begin{try}
\textbf{Try it, with the club’s members again}
\begin{enumerate}
\item Open \texttt{\textasciitilde{}/members.\codebreak{}csv} from the last Try it.
\item Press \texttt{C-x / f}. At the question about the column, type \texttt{town} and press
   \texttt{RET}. At the next question type \texttt{wanaka} and press \texttt{RET}. Only Aroha and
   Mere are left under the header, and the mode line says how many of the
   four records are shown.
\item Press \texttt{C-x / f} again, answer \texttt{joined} and then \texttt{2023}. Only Mere is left.
\item Press \texttt{C-x / l}. Both filters are listed. Press \texttt{q} to close the list.
\item Press \texttt{C-x / r}. All four members are back.
\item Press \texttt{C-x / h}, answer \texttt{\$1} and then \texttt{napier}. Hemi is hidden and the
   other three are shown. Press \texttt{C-x / r} again.
\end{enumerate}
\end{try}

\begin{xltabular}{\textwidth}{@{}l l l >{\raggedright\arraybackslash\hsize=1.000\hsize}X@{}}
\toprule
\textbf{Modal} & \textbf{Non-modal} & \textbf{Command} & \textbf{What it does}\\
\midrule
\endhead
none & \texttt{C-x / f} & \texttt{search-\codebreak{}filter-\codebreak{}show-\codebreak{}records} & Show only the records whose line a regular expression matches, hiding the others, as C-x / s does, deleting nothing.\\
none & \texttt{C-x / h} & \texttt{search-\codebreak{}filter-\codebreak{}hide-\codebreak{}records} & Hide the records whose line a regular expression matches, showing the others, deleting nothing: asked for as C-x / f asks, kept with the other filters.\\
none & \texttt{C-x / s} & \texttt{search-\codebreak{}filter-\codebreak{}regex} & Narrow the list here to the entries a regular expression matches, as it is typed: those whose line matches, or what more the list holds of them (a mailbox the headers and contents of its messages).\\
none & \texttt{C-x / l} & \texttt{search-\codebreak{}filter-\codebreak{}list} & List the search filters narrowing the list here (C-x / s) in a buffer of their own, each with how many entries it leaves; there d takes away the filter on the line and D all of them.\\
none & \texttt{d}, \texttt{C-c d} & \texttt{search-\codebreak{}filter-\codebreak{}remove} & Take away the search filter on the current line of the list of filters (C-x / l): the list it narrowed widens, the filters after it narrowing again what those before it leave, in the background where that takes time.\\
none & \texttt{D}, \texttt{C-c D}, \texttt{C-x / r} & \texttt{search-\codebreak{}filter-\codebreak{}remove-\codebreak{}all} & Take away every search filter narrowing the list here, or the list the list of filters here is for: the whole list is shown again.\\
none & \texttt{C-x / k} & \texttt{list-\codebreak{}keep-\codebreak{}records} & Keep the records whose line a regular expression matches and delete the rest, in a file of records, a line each: a CSV or TSV, a log, a JSON Lines file, todo.txt and the others \textbf{LINE-RECORD-MODES} names, the column asked for first as C-x / f asks.\\
none & \texttt{C-x / d} & \texttt{list-\codebreak{}delete-\codebreak{}records} & Delete the records whose line a regular expression matches, keeping the rest: asked for as C-x / k asks, in the same files.\\
\bottomrule
\end{xltabular}

\begin{xltabular}{\textwidth}{@{}l >{\raggedright\arraybackslash\hsize=0.582\hsize}X >{\raggedright\arraybackslash\hsize=1.418\hsize}X@{}}
\toprule
\textbf{Setting} & \textbf{Standard} & \textbf{What it changes}\\
\midrule
\endhead
\texttt{*line-\codebreak{}record-\codebreak{}modes*} & \texttt{(tuatara.\codebreak{}defaults::\codebreak{}log-\codebreak{}mode tuatara.\codebreak{}defa...} & The major modes whose files hold a record a line, as a log does: search filters (M-x search-filter-regex) narrow them, and the modes made from them, by hiding the lines they leave out, as they narrow a CSV.\\
\bottomrule
\end{xltabular}

\begin{xltabular}{\textwidth}{@{}l >{\raggedright\arraybackslash\hsize=1.000\hsize}X@{}}
\toprule
\textbf{Face} & \textbf{What it paints}\\
\midrule
\endhead
\texttt{:list-\codebreak{}filters} & The line above a list narrowed by search filters (C-x / s), saying by what and how many entries are left.\\
\bottomrule
\end{xltabular}

\subsection{Sorting a table}
\label{sec:447}
The sorting commands of chapter 3 know about tables. In a CSV or TSV file,
\texttt{M-x sort-\codebreak{}fields} asks which column to sort by, in the same way a filter asks,
and then puts the records in the order of that column. Each record moves
whole, and the header stays first.

\texttt{sort-\codebreak{}fields} compares the values as text. That is right for names and towns.
It is wrong for numbers, because as text \texttt{100} comes before \texttt{25}: the
comparison looks at the first character and \texttt{1} is before \texttt{2}. For a column of
prices, counts or years, use \texttt{M-x sort-\codebreak{}numeric-\codebreak{}fields}, which reads the number
at the start of each value. A value with no number counts as 0.

Both sort from small to large. \texttt{sort-\codebreak{}fields-\codebreak{}reverse} and
\texttt{sort-\codebreak{}numeric-\codebreak{}fields-\codebreak{}reverse} sort from large to small, which is what you want
for “the biggest sales first”.

Like every sorting command, these act on the lines you have selected and do
nothing without a selection. To sort the whole table, select it all first with
\texttt{C-x h}. A sort is refused while a filter hides some of the lines, so press
\texttt{C-x / r} before you sort. A sort is one step to undo with \texttt{C-/} (modal: \texttt{u}).

\begin{try}
\textbf{Try it, with a shop’s sales}
\begin{enumerate}
\item Press \texttt{C-x C-f}, type \texttt{\textasciitilde{}/sales.\codebreak{}csv} and press \texttt{RET}. Press \texttt{i} and type
   these lines, each ended with \texttt{RET}:
\begin{Verbatim}[breaklines,breakautoindent=false,breakindent=1em,breakafter=/-.\,;:=,commandchars=\\\{\}]
   date,item,town,price
   2026-03-02,gumboots,Motueka,89
   2026-03-02,rain jacket,Nelson,240
   2026-03-05,wool socks,Motueka,24
   2026-03-09,thermos,Takaka,55
   2026-03-11,tent pegs,Nelson,9
\end{Verbatim}
   Press \texttt{C-g} and save with \texttt{C-x C-s}.
\item Press \texttt{C-x h} to select everything. Press \texttt{M-x}, type \texttt{sort-\codebreak{}fields} and
   press \texttt{RET}. At the question about the column, type \texttt{price} and \texttt{RET}.
   The prices are now in the order 24, 240, 55, 89, 9. That is the order of
   text.
\item Press \texttt{C-x h} again and run \texttt{M-x sort-\codebreak{}numeric-\codebreak{}fields} on \texttt{price}. Now
   the order is 9, 24, 55, 89, 240, and the header is still the first line.
\item Press \texttt{C-x h} and run \texttt{M-x sort-\codebreak{}numeric-\codebreak{}fields-\codebreak{}reverse} on \texttt{price}. The
   rain jacket is at the top.
\item Press \texttt{C-x / f}, answer \texttt{town} and then \texttt{nelson}. The two sales in Nelson
   are shown, the dearest first. Press \texttt{C-x / r} to see them all again.
\end{enumerate}
\end{try}

\begin{xltabular}{\textwidth}{@{}l >{\raggedright\arraybackslash\hsize=0.650\hsize}X l >{\raggedright\arraybackslash\hsize=1.350\hsize}X@{}}
\toprule
\textbf{Modal} & \textbf{Non-modal} & \textbf{Command} & \textbf{What it does}\\
\midrule
\endhead
none & \texttt{M-x sort-\codebreak{}fields} & \texttt{sort-\codebreak{}fields} & Sort the lines the selection touches by one field of each, as text: fields are separated by spaces and tabs, and the prefix argument names the field (the first by default; a negative one counts from the end of the line).\\
none & \texttt{M-x sort-\codebreak{}fields-\codebreak{}reverse} & \texttt{sort-\codebreak{}fields-\codebreak{}reverse} & Sort the lines the selection touches by one field of each, as text, in descending order: as M-x sort-fields, the prefix argument naming the field (a negative one counting from the end of the line).\\
none & \texttt{M-x sort-\codebreak{}numeric-\codebreak{}fields} & \texttt{sort-\codebreak{}numeric-\codebreak{}fields} & Sort the lines the selection touches by the number at the start of one field of each (the prefix argument’s, as for sort-fields): 0x and hexadecimal digits, or a number in \textbf{SORT-NUMERIC-BASE}, decimal by default and then with a fraction and exponent; a field without a number counts as 0.\\
none & \texttt{M-x sort-\codebreak{}numeric-\codebreak{}fields-\codebreak{}reverse} & \texttt{sort-\codebreak{}numeric-\codebreak{}fields-\codebreak{}reverse} & Sort the lines the selection touches by the number at the start of one field of each, largest first: as M-x sort-numeric-fields, the prefix argument naming the field.\\
\bottomrule
\end{xltabular}

Chapter 11 explains what happens when the table is very large. A filter or a
sort of millions of records runs in the background, and you can go on working
while it does.

\subsection{The hex view}
\label{sec:452}
Every file is a row of numbers, each from 0 to 255. Such a number is called an
octet, or a byte. In a text file the numbers stand for letters: 97 is \texttt{a}. In
a photo or a program they stand for something else, and shown as letters they
are nonsense. The hex view shows the numbers themselves. Hex is short for
hexadecimal, a way of writing numbers with sixteen digits, 0 to 9 and then a
to f, so that every octet is exactly two digits: 97 is written \texttt{61}.

Few people need the hex view every day. It is there for the times you want to
know what is really in a file. Is this file that will not open really a
picture, or something else with a picture’s name? Why do two lines that look
the same compare as different? Is there something invisible at the start of
the file?

\texttt{M-x hex-\codebreak{}mode} shows the buffer you are in as hex, and the same command shows
it as text again. Figure 16.2 shows the start of the club’s membership list.

\begin{figure}[H]
\caption{Figure 16.2. The hex view: offset, sixteen octets, and their text.}\label{org448}
\begin{Verbatim}[breaklines,breakautoindent=false,breakindent=1em,breakafter=/-.\,;:=,commandchars=\\\{\}]
 00000000  6e 61 6d 65 2c 74 6f 77  6e 2c 6a 6f 69 6e 65 64  name,town,joined
 00000010  0a 41 72 6f 68 61 20 4e  67 61 74 61 2c 57 61 6e  .Aroha Ngata,Wan
 00000020  61 6b 61 2c 32 30 31 39  0a 54 61 6d 61 20 50 61  aka,2019.Tama Pa

 offset    the octets in hex                                 the same as text
\end{Verbatim}
\end{figure}

Each row shows sixteen octets. On the left is the offset: how far into the
file the row starts, also in hex. In the middle are the sixteen octets. On the
right are the characters they make. An octet that is not a letter, digit or
sign you could print is shown as a dot. In the figure, the first octet of the
second row is \texttt{0a}, the end of a line, and it shows as a dot on the right.

The motion keys move by octets and rows: \texttt{h} and \texttt{l} one octet, \texttt{j} and \texttt{k} one
row. The octet under the cursor is marked in both places at once, in the hex
digits and in the text beside them, so you can see which character a number
stands for. The mode line shows \texttt{Hex}, the offset of the cursor, and \texttt{digits}
or \texttt{text} for the side the cursor is on. \texttt{TAB} moves the cursor from one side
to the other. \texttt{:} (non-modal: \texttt{M-g g}) asks for an offset and goes there. Type
it in decimal, or with \texttt{0x} in front for hex.

Only the rows on the screen are worked out, so the hex view of a very large
file opens at once.

\subsubsection{Searching for octets}
\label{sec:449}
A text search cannot find what cannot be typed. \texttt{M-x hex-\codebreak{}search-\codebreak{}forward} asks
for octets in hex, with or without spaces between them, and goes to the next
place after the cursor that holds them. \texttt{ff d8 ff} finds the start of a JPEG
photo, for example, and \texttt{25 50 44 46} finds \texttt{\%PDF}, the first four octets of a
PDF. The echo area says where it was found, in hex and in decimal.
\texttt{M-x hex-\codebreak{}search-\codebreak{}backward} looks before the cursor.

\subsubsection{Changing octets}
\label{sec:450}
The hex view can change a file, one octet over another. In INSERT state, with
the cursor on the hex side, typing hex digits writes over the octets there.
With the cursor on the text side, typing writes characters over them.
\texttt{M-x hex-\codebreak{}set-\codebreak{}byte} asks for one octet in hex and writes it at the cursor, then
moves on.

The size of the file never changes in the hex view. The keys that would add or
remove text say so and do nothing, because one octet too many would shift
everything after it. As always, nothing reaches the file until you save.

A buffer that is narrowed (chapter 7) is not shown in hex. Widen it first with
\texttt{C-x x w}.

\subsubsection{Files that open in hex by themselves}
\label{sec:451}
Tuatara looks at the start of every file it opens. If it finds an octet with
the value 0 in the first 64 KB, the file is not text. Programs, archives and
disk images have such octets and text does not. A file like that opens in the
hex view straight away, where showing it as text would fill the screen with
rubbish. If you would sooner see such files as text, set
\texttt{*hex-\codebreak{}view-\codebreak{}binary-\codebreak{}files*} to \texttt{nil}.

\begin{try}
\textbf{Try it, with the club’s members once more}
\begin{enumerate}
\item Open \texttt{\textasciitilde{}/members.\codebreak{}csv} and press \texttt{M-<} to go to the top. Press \texttt{M-x}, type
   \texttt{hex-\codebreak{}mode} and press \texttt{RET}. The echo area says \texttt{Hex view on}, and the
   screen looks like figure 16.2.
\item Press \texttt{l} four times. The cursor is on \texttt{2c}, and the comma on the right
   is marked with it. The mode line shows the offset \texttt{0x00000004}.
\item Press \texttt{TAB}. The cursor is now on the comma in the text. Press \texttt{TAB}
   again to go back to the digits.
\item Press \texttt{M-x}, type \texttt{hex-\codebreak{}search-\codebreak{}forward} and press \texttt{RET}. Type
   \texttt{57 61 6e 61 6b 61}, which spells \texttt{Wanaka}, and press \texttt{RET}. The echo
   area says \texttt{Found at 0x1d (29)}.
\item Press \texttt{:} (non-modal: \texttt{M-g g}), type \texttt{0} and press \texttt{RET} to go back to the
   start.
\item Run \texttt{M-x hex-\codebreak{}mode} again. The echo area says \texttt{Hex view off} and the table
   is back in its colours.
\end{enumerate}
\end{try}

\begin{xltabular}{\textwidth}{@{}l >{\raggedright\arraybackslash\hsize=0.542\hsize}X l >{\raggedright\arraybackslash\hsize=1.458\hsize}X@{}}
\toprule
\textbf{Modal} & \textbf{Non-modal} & \textbf{Command} & \textbf{What it does}\\
\midrule
\endhead
none & \texttt{M-x hex-\codebreak{}mode} & \texttt{hex-\codebreak{}mode} & Show the current buffer as a hex dump, or as text again: each row the offset, sixteen octets in hex and the text they make, a dot for an octet that is no printable ASCII character.\\
none & \texttt{M-x hex-\codebreak{}switch-\codebreak{}column} & \texttt{hex-\codebreak{}switch-\codebreak{}column} & In the hex view, go from the hex digits to the text beside them, or back: typing then writes characters there, or hex digits here, over the octets.\\
none & \texttt{M-x hex-\codebreak{}set-\codebreak{}byte} & \texttt{hex-\codebreak{}set-\codebreak{}byte} & In the hex view, write one octet over the one at the cursor, asking for it in hex (41 is A), and move on to the next.\\
none & \texttt{M-x hex-\codebreak{}search-\codebreak{}forward} & \texttt{hex-\codebreak{}search-\codebreak{}forward} & Find the next place after the cursor that holds the octets asked for in hex (“de ad be ef”), and show it in the hex view.\\
none & \texttt{M-x hex-\codebreak{}search-\codebreak{}backward} & \texttt{hex-\codebreak{}search-\codebreak{}backward} & Find the last place before the cursor that holds the octets asked for in hex, and show it in the hex view.\\
\bottomrule
\end{xltabular}

\begin{xltabular}{\textwidth}{@{}l l >{\raggedright\arraybackslash\hsize=1.000\hsize}X@{}}
\toprule
\textbf{Setting} & \textbf{Standard} & \textbf{What it changes}\\
\midrule
\endhead
\texttt{*hex-\codebreak{}view-\codebreak{}binary-\codebreak{}files*} & \texttt{t} & True to show a file that looks binary in the hex view (M-x hex-mode) when it is opened: one with a null (NUL) octet in its first 64 KB, as executables, archives and disk images have and text has not.\\
\bottomrule
\end{xltabular}

\begin{xltabular}{\textwidth}{@{}l >{\raggedright\arraybackslash\hsize=1.000\hsize}X@{}}
\toprule
\textbf{Face} & \textbf{What it paints}\\
\midrule
\endhead
\texttt{:hex-\codebreak{}offset} & The offsets at the start of the hex view’s rows.\\
\texttt{:hex-\codebreak{}mirror} & The octet at the hex view’s cursor, marked in both its hex digits and its character, the cursor on the one typing goes to.\\
\texttt{:hex-\codebreak{}unprintable} & The dots the hex view’s text column shows for octets that are not printable.\\
\bottomrule
\end{xltabular}

\subsection{Pictures}
\label{sec:455}
Open a picture with \texttt{C-x C-f} and Tuatara draws it. It knows the common kinds,
among them PNG, JPEG, GIF, WebP, BMP, TIFF, HEIF and AVIF, and the newer JPEG
2000 and JPEG XL. It tells the kind from what is in the file. A photo that
someone named \texttt{beach.\codebreak{}txt} is still shown as a photo.

The picture is drawn as large as the window allows and never larger than it
really is. A line above it says what it is and how many pixels wide and high.
The buffer is read-only and has nothing to save. \texttt{q} closes it.

This is for a quick look: is this the photo I meant, what is on this scan,
which of these twelve files is the one with the lake. For more than a look,
\texttt{O} hands the picture to the program your desktop uses for pictures, as
chapter 9 describes. Chapter 10 shows how a folder listing can preview each
picture as you move down the list, which is the fast way through a folder of
photos.

Drawing happens in the background. A large photo takes a moment, and in that
moment you can go on to something else. The mode line shows that work is going
on, and \texttt{C-x j} lists it as \texttt{Drawing} with the file’s name. Killing the buffer
stops the drawing.

The program that turns the file into a picture runs shut in a sandbox, where
it can reach neither the network nor your other files. A picture from a
stranger can therefore be looked at without worry. Chapter 12 explains the
sandbox.

\subsubsection{Pictures in a terminal}
\label{sec:453}
In its own window Tuatara always shows pictures. In a terminal it depends on
the terminal. Some terminals, such as kitty and Ghostty, can draw pictures, and
Tuatara asks the terminal when it starts. The setting \texttt{*tui-\codebreak{}pictures*} is
\texttt{:auto} as standard, which means “show them where the terminal says it can”.
Set it to \texttt{t} if your terminal can draw but does not say so, or to \texttt{nil} to
have no pictures. Inside tmux or screen, \texttt{:auto} shows none.

Where pictures cannot be shown, a picture opens in the hex view, as the last
section of this chapter describes.

\subsubsection{Very large pictures}
\label{sec:454}
Before it draws anything, Tuatara reads the size the file claims to have and
works out how much memory the drawing needs. A picture that needs more than
\texttt{*picture-\codebreak{}memory*} is not drawn, and the line above says how large it is and
what it would take. The standard is 4 GB, enough for a photo of about 85
million pixels. This guards against a small file that claims to be an enormous
picture.

\begin{try}
\textbf{Try it, with a holiday photo}
\begin{enumerate}
\item Find a photo on your computer, such as one from your phone. This example
   calls it \texttt{\textasciitilde{}/Pictures/\codebreak{}beach.\codebreak{}jpg}.
\item Press \texttt{C-x C-f}, type the photo’s name and press \texttt{RET}. The photo is
   drawn. The line above it says it is a JPEG image and gives its size.
\item Press \texttt{C-x 3} to split the window. The photo is drawn in both windows, in
   each as large as fits. Press \texttt{C-x 1} (modal: \texttt{F}) to go back to one window.
\item Press \texttt{M-x}, type \texttt{show-\codebreak{}file-\codebreak{}metadata} and press \texttt{RET}. A new buffer says
   what the photo keeps of itself: the camera, the time it was taken and,
   if the camera stored it, the place. Press \texttt{q} to close it.
\item Press \texttt{q} to close the photo.
\end{enumerate}
\end{try}

\begin{xltabular}{\textwidth}{@{}l l l >{\raggedright\arraybackslash\hsize=1.000\hsize}X@{}}
\toprule
\textbf{Modal} & \textbf{Non-modal} & \textbf{Command} & \textbf{What it does}\\
\midrule
\endhead
none & \texttt{O}, \texttt{C-c O} & \texttt{open-\codebreak{}externally} & Open the file shown, in the desktop’s own program for its kind (xdg-open), or the one \textbf{OPEN-EXTERNALLY-PROGRAMS} gives its Multipurpose Internet Mail Extensions (MIME) type, for what Tuatara does not do: printing, forms, playing.\\
\bottomrule
\end{xltabular}

\begin{xltabular}{\textwidth}{@{}l l >{\raggedright\arraybackslash\hsize=1.000\hsize}X@{}}
\toprule
\textbf{Setting} & \textbf{Standard} & \textbf{What it changes}\\
\midrule
\endhead
\texttt{*tui-\codebreak{}pictures*} & \texttt{:auto} & Whether pictures (a preview of an image, a PDF’s pages) are shown in the terminal.\\
\texttt{*picture-\codebreak{}memory*} & \texttt{4294967296} & The most memory the program drawing a large picture is given (4 GB): a picture’s own size, read from its file, tells how much it needs (about 32 octets a pixel, a hundred-megapixel PNG some 3.9 GB), never less than \textbf{SANDBOX-MEMORY}; one needing more than this is not drawn, and says so.\\
\bottomrule
\end{xltabular}

\subsection{PDFs}
\label{sec:461}
A PDF is a document fixed on its pages: a bank statement, a manual, a form
from the council. Open one with \texttt{C-x C-f} and Tuatara shows its pages as they
would print. You can turn the pages, make them larger, search the text, and
copy from it. You cannot change a PDF or fill in a form in it. For that, \texttt{O}
opens it in your desktop’s PDF program.

The mode line says \texttt{PDF}, the page you are on and how many there are, and how
the page is fitted, as in \texttt{Page 1 of 2, fit page}. Like pictures, the pages
are drawn in the background by a program shut in the sandbox. The pages near
the one you are reading are drawn ahead, so turning a page is quick. \texttt{q} closes
the PDF. When you open the same PDF another day, it opens at the page and zoom
where you left it.

\subsubsection{Turning pages}
\label{sec:456}
\texttt{l} goes to the next page and \texttt{h} to the one before, and so do \texttt{PageDown} and
\texttt{PageUp} (written \texttt{<next>} and \texttt{<prior>} in the tables). \texttt{g h} goes to the
first page and \texttt{g l} to the last. \texttt{g p} asks for a page number and goes
there, which is the quick way to page 212 of a manual.

\texttt{j} and \texttt{k} scroll down and up within a page. At the bottom of a page, \texttt{j}
goes on to the top of the next, so you can read a whole document with one key.
When the whole page fits in the window there is nothing to scroll, and \texttt{j} and
\texttt{k} turn a page at a time. The mouse wheel does the same as \texttt{j} and \texttt{k}.
\texttt{g g} shows the top of the page and \texttt{g e} its end.

\subsubsection{Zoom}
\label{sec:457}
A PDF opens with each page fitted whole into the window. That is good for
seeing the layout and often too small to read. \texttt{w} fits the width of the page
to the window. The text is then as large as it can be without scrolling
sideways, and you read down the page with \texttt{j}. \texttt{p} fits the whole page again.

\texttt{=} or \texttt{+} makes the page larger, by a quarter each time, and \texttt{-} makes it
smaller. The echo area says the zoom, such as \texttt{Zoom: 125\%}. \texttt{0} goes back to
the fit you chose with \texttt{p} or \texttt{w}. When a page is wider than the window, \texttt{H}
and \texttt{L} scroll it left and right.

\texttt{d} shows the pages dark: the white of the paper becomes the background colour
of your theme and the black of the text becomes its text colour. It is easier
on the eyes at night. It holds for every PDF until you press \texttt{d} again or
leave Tuatara. To have it always, set \texttt{*pdf-\codebreak{}view-\codebreak{}dark-\codebreak{}pages*} to \texttt{t}.

\subsubsection{Searching, and the text of a PDF}
\label{sec:458}
\texttt{C-s} asks for text and searches the PDF from the page you are on. It turns to
the page where the text is found, and the echo area shows the page number and
the line. \texttt{C-s} and then \texttt{RET} with nothing typed finds the next place. As in
every search, small letters match capitals too.

\texttt{T} switches from the pages to the text of the PDF, as plain text in the same
buffer, with a line marking where each page starts. There the keys move as in
any text. You can search it in every way chapter 6 describes, select a
paragraph and copy it with the keys of chapter 4. \texttt{T} goes back to the pages,
at the page the cursor was on. \texttt{PageDown}, \texttt{PageUp} and \texttt{g p} go by pages in
the text too.

A scanned PDF is a special case. A scanner makes a photo of each page, and a
photo of words holds no text. Such a PDF shows its pages like any other, but a
search finds nothing and \texttt{T} shows a single line saying that there is no text
in it.

\subsubsection{What the PDF says of itself, and files inside it}
\label{sec:459}
\texttt{i} shows what the PDF says of itself, in a help buffer: its title and author,
the programs that made it, when it was made and changed, how many pages it has
and their size, whether it is locked with a password, and whether it holds
forms or scripts. This is useful when you wonder where a document came from.
\texttt{q} closes the help.

A PDF can carry other files inside it, as an email carries attachments. An
invoice may hold its own figures as a spreadsheet, for example. \texttt{a} lists the
files embedded in the PDF in a buffer beside it. There, \texttt{n} and \texttt{p} move from
one to the next, and \texttt{RET} or \texttt{s} saves the one on the line. Tuatara asks where
to put it and writes it there as it is. It never opens or runs an embedded
file. If the PDF has none, the echo area says so.

To see the raw octets of a PDF, run \texttt{M-x hex-\codebreak{}mode} in its pages.

\begin{try}
\textbf{Try it, with a scanned form}
\begin{enumerate}
\item Find a PDF of a few pages. A form you scanned is a good one, such as an
   enrolment form for a school in Timaru. This example calls it
   \texttt{\textasciitilde{}/Documents/\codebreak{}form.\codebreak{}pdf}.
\item Press \texttt{C-x C-f}, type its name and press \texttt{RET}. The first page is shown
   whole, and the mode line ends with \texttt{Page 1 of} and the number of pages.
\item Press \texttt{l} to go to page 2, then \texttt{g p}, type \texttt{1} and press \texttt{RET} to go
   back.
\item Press \texttt{w}. The page is now as wide as the window. Press \texttt{j} a few times
   to read down it.
\item Press \texttt{=} twice. The echo area says \texttt{Zoom: 125\%} and then \texttt{Zoom: 156\%}.
   Press \texttt{L} and \texttt{H} to move across the page. Press \texttt{0} and then \texttt{p} to see
   the whole page again.
\item Press \texttt{T}. If the form was scanned, a line says there is no text in it.
   If it was made on a computer, you see its text. Press \texttt{T} to go back to
   the pages.
\item Press \texttt{i} to read what the PDF says of itself. A scan often names the
   scanner as the program that made it. Press \texttt{q} to close the help.
\item Press \texttt{a}. Most forms have no embedded files, and the echo area says so.
\item Press \texttt{q} to close the PDF.
\end{enumerate}
\end{try}

\subsubsection{Commands and settings for PDFs}
\label{sec:460}
\begin{xltabular}{\textwidth}{@{}l >{\raggedright\arraybackslash\hsize=0.667\hsize}X l >{\raggedright\arraybackslash\hsize=1.333\hsize}X@{}}
\toprule
\textbf{Modal} & \textbf{Non-modal} & \textbf{Command} & \textbf{What it does}\\
\midrule
\endhead
none & \texttt{l}, \texttt{C-c l} & \texttt{pdf-\codebreak{}view-\codebreak{}next-\codebreak{}page} & Turn to the next page of the PDF.\\
none & \texttt{h}, \texttt{C-c h} & \texttt{pdf-\codebreak{}view-\codebreak{}previous-\codebreak{}page} & Turn to the page before in the PDF.\\
none & \texttt{<next>} & \texttt{pdf-\codebreak{}view-\codebreak{}next-\codebreak{}page-\codebreak{}anywhere} & Turn to the next page of the PDF; in its text, go to where it begins.\\
none & \texttt{<prior>} & \texttt{pdf-\codebreak{}view-\codebreak{}previous-\codebreak{}page-\codebreak{}anywhere} & Turn to the page before in the PDF; in its text, go to where it begins.\\
none & \texttt{g h} & \texttt{pdf-\codebreak{}view-\codebreak{}first-\codebreak{}page} & Turn to the first page of the PDF.\\
none & \texttt{g l} & \texttt{pdf-\codebreak{}view-\codebreak{}last-\codebreak{}page} & Turn to the last page of the PDF.\\
none & \texttt{g p} & \texttt{pdf-\codebreak{}view-\codebreak{}goto-\codebreak{}page} & Turn to a page of the PDF by its number, asked for; in its text, go to where it begins.\\
none & \texttt{g g} & \texttt{pdf-\codebreak{}view-\codebreak{}page-\codebreak{}top} & Show the top of the page.\\
none & \texttt{g e} & \texttt{pdf-\codebreak{}view-\codebreak{}page-\codebreak{}end} & Show the bottom of the page.\\
none & \texttt{j}, \texttt{C-c j} & \texttt{pdf-\codebreak{}view-\codebreak{}scroll-\codebreak{}down} & Scroll the page down, and at its end on to the top of the next.\\
none & \texttt{k}, \texttt{C-c k} & \texttt{pdf-\codebreak{}view-\codebreak{}scroll-up} & Scroll the page up, and at its top on to the end of the page before.\\
none & \texttt{H}, \texttt{C-c H} & \texttt{pdf-\codebreak{}view-\codebreak{}scroll-\codebreak{}left} & Scroll a page wider than the window to the left.\\
none & \texttt{L}, \texttt{C-c L} & \texttt{pdf-\codebreak{}view-\codebreak{}scroll-\codebreak{}right} & Scroll a page wider than the window to the right.\\
none & \texttt{<wheel-\codebreak{}up>}, \texttt{<wheel-\codebreak{}down>} & \texttt{pdf-\codebreak{}view-\codebreak{}scroll-\codebreak{}by-\codebreak{}wheel} & Scroll the page as the mouse wheel turned, as M-x pdf-view-scroll-down and M-x pdf-view-scroll-up do: down for <wheel-down>, on to the next page at a page’s end (each turn a page, the page fitted whole), up for <wheel-up>.\\
\bottomrule
\end{xltabular}

\begin{xltabular}{\textwidth}{@{}l l l >{\raggedright\arraybackslash\hsize=1.000\hsize}X@{}}
\toprule
\textbf{Modal} & \textbf{Non-modal} & \textbf{Command} & \textbf{What it does}\\
\midrule
\endhead
none & \texttt{p}, \texttt{C-c p} & \texttt{pdf-\codebreak{}view-\codebreak{}fit-\codebreak{}page} & Draw the PDF’s pages as large as the window shows them whole.\\
none & \texttt{w}, \texttt{C-c w} & \texttt{pdf-\codebreak{}view-\codebreak{}fit-\codebreak{}width} & Draw the PDF’s pages as wide as the window, to scroll down.\\
none & \texttt{+}, \texttt{=}, \texttt{C-c +} & \texttt{pdf-\codebreak{}view-\codebreak{}zoom-in} & Draw the PDF’s pages larger, by \textbf{PDF-VIEW-ZOOM-STEP}; scroll with j, k, H and L.\\
none & \texttt{-}, \texttt{C-c -} & \texttt{pdf-\codebreak{}view-\codebreak{}zoom-\codebreak{}out} & Draw the PDF’s pages smaller, by \textbf{PDF-VIEW-ZOOM-STEP}.\\
none & \texttt{0}, \texttt{C-c 0} & \texttt{pdf-\codebreak{}view-\codebreak{}zoom-\codebreak{}reset} & Draw the PDF’s pages at the size of the fit chosen again, page or width (p, w), no longer zoomed in or out.\\
none & \texttt{d}, \texttt{C-c d} & \texttt{pdf-\codebreak{}view-\codebreak{}toggle-\codebreak{}dark} & Show the PDF’s pages dark, white as the theme’s background and black as its text, or as they are again; for every PDF while the editor runs (\textbf{PDF-VIEW-DARK-PAGES}).\\
\bottomrule
\end{xltabular}

\begin{xltabular}{\textwidth}{@{}l l l >{\raggedright\arraybackslash\hsize=1.000\hsize}X@{}}
\toprule
\textbf{Modal} & \textbf{Non-modal} & \textbf{Command} & \textbf{What it does}\\
\midrule
\endhead
none & \texttt{M-x pdf-\codebreak{}view-\codebreak{}search} & \texttt{pdf-\codebreak{}view-\codebreak{}search} & Find text in the PDF and turn to its page: from the page shown on, asked for; again with nothing typed, the next place of the same text.\\
none & \texttt{T}, \texttt{C-c T} & \texttt{pdf-\codebreak{}view-\codebreak{}toggle-\codebreak{}text} & Switch the PDF’s view between its pages and its text, to read, search and copy; the text opens at the page shown, and the pages at the page point is on.\\
none & \texttt{i}, \texttt{C-c i} & \texttt{pdf-\codebreak{}view-\codebreak{}facts} & Show what the PDF says of itself, as help: title, subject, author, the programs that made it (creator and producer), when it was made and changed, its version, encryption and what it allows, JavaScript, forms, pages and their size, as its program reads them in the sandbox; the file’s own name and size first.\\
none & \texttt{a}, \texttt{C-c a} & \texttt{pdf-\codebreak{}view-\codebreak{}attachments} & List the files embedded in the PDF, in a buffer beside it, as its program reads them in the sandbox; there, s saves one.\\
none & \texttt{s}, \texttt{RET} & \texttt{pdf-\codebreak{}attachment-\codebreak{}save} & Save the embedded file on the current line where you say, a new file, or one there after asking: taken out of the PDF in the sandbox, as its program reads it, and written as it is, never opened.\\
\bottomrule
\end{xltabular}

\begin{xltabular}{\textwidth}{@{}l l >{\raggedright\arraybackslash\hsize=1.000\hsize}X@{}}
\toprule
\textbf{Setting} & \textbf{Standard} & \textbf{What it changes}\\
\midrule
\endhead
\texttt{*pdf-\codebreak{}view-\codebreak{}dark-\codebreak{}pages*} & \texttt{nil} & True to show a PDF’s pages dark: white as the theme’s background, black as its text, the colours between in proportion.\\
\bottomrule
\end{xltabular}

\subsection{Office documents}
\label{sec:462}
Word documents, Excel workbooks and PowerPoint presentations open with
\texttt{C-x C-f} too, and so do the OpenDocument files that LibreOffice writes.
Tuatara shows them in one of two ways.

If LibreOffice is installed, Tuatara has it draw the pages of the document,
shut in the sandbox, and shows them as it shows a PDF. Every key of the PDF
section works: \texttt{l} and \texttt{h} turn pages, \texttt{w} fits the width, \texttt{C-s} searches.
You see the document as it would print, with its fonts, tables and pictures.

\texttt{T} switches to the document read as text. Tuatara reads the file itself and
writes what it finds as an outline: the headings of a Word document, the sheets
and cells of a workbook, the slides of a presentation. That is the view for
copying a paragraph out, or for searching a workbook for a number. It is
read-only. \texttt{T} there goes back to the pages.

Without LibreOffice, the document opens as text straight away, and the echo
area says why. You can ask for that always by setting \texttt{*office-\codebreak{}view-\codebreak{}pages*}
to \texttt{nil}, which also means LibreOffice is never run. The older kinds of Office
file, from before 2007 (\texttt{.doc}, \texttt{.xls}, \texttt{.ppt}), can only be shown as pages.
Without LibreOffice they open as raw octets.

Either way, the document is never changed, and nothing in it is run. Macros
stay asleep, and links to the internet are not followed. To edit the document,
\texttt{O} opens it in your desktop’s program. If that program offers to run the
document’s macros, say no unless you trust where the file came from.

Two settings keep a document read as text within bounds. A small Office file
can stand for a great deal of text. Reading stops at \texttt{*office-\codebreak{}note-\codebreak{}size-\codebreak{}limit*}
characters and the text says that the rest is left out. A single cell or
formula longer than \texttt{*office-\codebreak{}value-\codebreak{}length-\codebreak{}limit*} characters is cut short and
ends in \texttt{...}.

\begin{xltabular}{\textwidth}{@{}l >{\raggedright\arraybackslash\hsize=0.536\hsize}X l >{\raggedright\arraybackslash\hsize=1.464\hsize}X@{}}
\toprule
\textbf{Modal} & \textbf{Non-modal} & \textbf{Command} & \textbf{What it does}\\
\midrule
\endhead
none & \texttt{M-x office-\codebreak{}view-\codebreak{}pages} & \texttt{office-\codebreak{}view-\codebreak{}pages} & Show the pages of the Office document read as text here, in its place.\\
\bottomrule
\end{xltabular}

\begin{xltabular}{\textwidth}{@{}l l >{\raggedright\arraybackslash\hsize=1.000\hsize}X@{}}
\toprule
\textbf{Setting} & \textbf{Standard} & \textbf{What it changes}\\
\midrule
\endhead
\texttt{*office-\codebreak{}view-\codebreak{}pages*} & \texttt{t} & True to show a Word document, an Excel workbook, a PowerPoint presentation (and the older and OpenDocument kinds of them) opened as its pages, drawn by LibreOffice shut in by the sandbox, when it is installed; M-x pdf-view-toggle-text switches to it read as text.\\
\texttt{*office-\codebreak{}note-\codebreak{}size-\codebreak{}limit*} & \texttt{16777216} & The most characters a note read from a Word document, an Excel workbook, a PowerPoint presentation or an OpenDocument file may have: past them the rest of its text is left out, and the note says so, so that a file whose few octets stand for text copied many times cannot fill the memory.\\
\texttt{*office-\codebreak{}value-\codebreak{}length-\codebreak{}limit*} & \texttt{2000} & The most characters of one value an office file’s note shows (a cell of a sheet or a table, a formula): a longer one is cut after a word, with \ldots{}\\
\bottomrule
\end{xltabular}

\subsection{What a file says of itself}
\label{sec:463}
Many files carry facts about themselves that you do not see when you look at
them. A photo records the camera, the time and often the place where it was
taken. A Word document records who wrote it and who saved it last. A PDF
records the program that made it. These facts are called metadata.

\texttt{M-x show-\codebreak{}file-\codebreak{}metadata} shows them in a buffer of their own, which \texttt{q}
closes. It works on the file of the buffer you are in, and on the picture or
PDF you are looking at. In a folder listing the key is \texttt{g i}, and it reports
on the file the cursor is on. With \texttt{C-u} before it, it asks for a file.

The report has two parts. The first is what your computer knows about the
file: where it is, what kind of file its contents say it is, its size, who
owns it, when it was last changed, and its SHA-256 sum, a long number that is
the same for two files only when they are the same octet for octet. The second
part is what the file keeps inside itself. What is there depends on the kind
of file:

\begin{itemize}
\item a photo: the camera and its settings, the times, the place if it was
  stored, and whether a small preview picture is kept inside;
\item a PDF: the title, the author, the programs that made it, the dates, and
  how many times it was saved;
\item a Word, Excel or PowerPoint file: the author, the last person to save it,
  and the dates;
\item a sound or video file: its tags, such as the title and artist, and its
  times;
\item an encrypted file: which keys it was encrypted for.
\end{itemize}

There are everyday reasons to look. Before you put a photo of your house on
the internet, check whether it carries the place it was taken. Before you send
a document to a customer, check whose name is in it as the author. When two
dates disagree, the file’s own record of when it was made is often the one to
trust over the date in a folder listing, which changes each time the file is
copied.

The file is read in the background by Tuatara alone. Nothing in it is run, and
the time your computer keeps for when the file was last read is left as it
was.

For an encrypted file, the report can also say whose the keys are, by looking
them up in your own keyring. That runs the gpg program, so Tuatara asks first.
The setting \texttt{*metadata-\codebreak{}look-\codebreak{}up-\codebreak{}key-\codebreak{}owners*} says whether it asks every time
(\texttt{:ask}, the standard), always looks them up (\texttt{t}) or never does (\texttt{nil}). The
file itself is never handed to gpg.

\begin{try}
\textbf{Try it, with the holiday photo and the members}
\begin{enumerate}
\item Open \texttt{\textasciitilde{}/members.\codebreak{}csv}. Press \texttt{M-x}, type \texttt{show-\codebreak{}file-\codebreak{}metadata} and press
   \texttt{RET}. A buffer opens with the file’s path, kind, size, owner and times.
   A plain text file keeps nothing more of itself. Press \texttt{q}.
\item Press \texttt{C-x C-f}, type \texttt{\textasciitilde{}/Pictures/} and press \texttt{RET} to list the folder
   with your photo in it. Move to the photo with \texttt{j}.
\item Press \texttt{g i}. The report now has a second part, with what the camera
   wrote into the photo. Look for the time and for a position.
\item Press \texttt{q} to close the report, and \texttt{q} again to close the folder.
\end{enumerate}
\end{try}

\begin{xltabular}{\textwidth}{@{}l l l >{\raggedright\arraybackslash\hsize=1.000\hsize}X@{}}
\toprule
\textbf{Modal} & \textbf{Non-modal} & \textbf{Command} & \textbf{What it does}\\
\midrule
\endhead
none & \texttt{g i} & \texttt{show-\codebreak{}file-\codebreak{}metadata} & Show what the file shown is and what it says of itself, in a buffer of its own, which q closes: the file the buffer visits, the entry on the line in a list of a directory, the PDF or picture in its view; with C-u, or where there is none, asked for.\\
\bottomrule
\end{xltabular}

\begin{xltabular}{\textwidth}{@{}l l >{\raggedright\arraybackslash\hsize=1.000\hsize}X@{}}
\toprule
\textbf{Setting} & \textbf{Standard} & \textbf{What it changes}\\
\midrule
\endhead
\texttt{*metadata-\codebreak{}look-\codebreak{}up-\codebreak{}key-\codebreak{}owners*} & \texttt{:ask} & Whether M-x show-file-metadata looks up, in your own GnuPG keyring, whose the keys an encrypted file is for are, running gpg --list-keys (the file itself is never given to gpg): true looks them up, :ask asks each time, false (nil) never does, the keys then shown by their IDs alone.\\
\bottomrule
\end{xltabular}

\subsection{When a file cannot be drawn}
\label{sec:464}
Things go wrong with files. A download stops halfway and leaves half a
picture. A file has the name of one kind and the contents of another. A
terminal cannot draw pictures. In each case Tuatara tells you what happened in
the echo area and shows you what it can.

\begin{xltabular}{\textwidth}{@{}>{\raggedright\arraybackslash\hsize=1.170\hsize}X >{\raggedright\arraybackslash\hsize=0.830\hsize}X@{}}
\toprule
\textbf{What you opened} & \textbf{What you get}\\
\midrule
\endhead
A picture that is damaged or too large & Its octets in the hex view.\\
A picture, where pictures cannot be shown & Its octets in the hex view.\\
A PDF whose pages cannot be drawn & Its text.\\
A PDF, where pictures cannot be shown & Its text.\\
An Office document, without LibreOffice & Its text, or for the kinds from before 2007 its octets.\\
Any other file that is not text & Its octets in the hex view.\\
\bottomrule
\end{xltabular}

When a file falls back to the hex view, it is opened safely, as chapter 12
describes: read-only, with nothing in it acted on. The echo area gives the
reason and ends with \texttt{its octets, opened safely}.

The hex view is the right place to land, because it never fails. Every file
can be shown as its octets, and the first few often tell you what the file
really is. A file that starts with \texttt{\%PDF} is a PDF whatever its name says. One
that starts with \texttt{PK} is a zip archive, or a Word document, which is a zip
archive inside. One whose first row is all zeros is probably a download that
never finished. From there \texttt{M-x show-\codebreak{}file-\codebreak{}metadata} tells you what kind of
file the contents say it is, and \texttt{M-x hex-\codebreak{}search-\codebreak{}forward} finds what you are
looking for inside it.

\section{The terminal}
\label{sec:485}
A shell inside the editor, with its output kept as text you can search and copy.

Sooner or later you want to run a command while you are writing: check what
is in a folder, rename a batch of files, run a script, look at what a program
is doing. You can switch to another window on your desktop and back. Tuatara
offers another way. \texttt{M-x terminal} runs your shell in a buffer, so the shell
sits next to your text and what it prints can be read, searched and copied
like any other text.

This chapter shows how to open a terminal and where it appears, how typing
there differs from typing in a file, how to look back through what a program
printed, how to paste, how to use links and pictures, and what to do with a
program that wants a key Tuatara would keep for itself. The last part is about
the other direction: Tuatara running inside a terminal, which is how many
people start it, and what that changes for you.

A word on names. A \emph{shell} is the program that reads the commands you type
(\texttt{ls}, \texttt{cd}, \texttt{grep}) and runs them. It is bash unless you use another. A
\emph{terminal} in this chapter is a buffer that shows a shell, or any other
program, running on a screen of its own. The first part of the chapter is about
that. The last part is about the terminal your own desktop gives you, in which
Tuatara itself can run.

The state keys from chapter 1 apply as always, with one twist. In a terminal
you are typing into a program, so most keys do not mean what they mean in a
file. The sections below say which do.

\subsection{Opening a terminal}
\label{sec:469}
\texttt{M-x terminal} opens one. It starts your shell in the folder of the buffer you
came from, in a window across the bottom of the screen, about a third of the
screen high. If you were looking at \texttt{\textasciitilde{}/letters/\codebreak{}council.\codebreak{}txt}, the shell starts
in \texttt{\textasciitilde{}/letters/}. Type \texttt{ls} and you see your letters.

Each \texttt{M-x terminal} opens another terminal, with a name of its own:
\texttt{*terminal*}, \texttt{*terminal*<2>}, and so on. When the program ends, for example
when you type \texttt{exit}, the buffer and its window go with it. Killing the
buffer with \texttt{C-x k} ends the program too, without asking.

A terminal opened from a file you opened safely (chapter 12), or from a folder
you listed as safe or as evidence, starts in your home folder instead. The
shell’s prompt and its helpers then read nothing from that folder, which could
hold a hostile \texttt{.envrc} or \texttt{.git}. The echo area says so when it happens.

Figure 17.1 shows the screen with a letter above and a shell below.

\begin{figure}[H]
\caption{Figure 17.1. A letter above, a terminal across the bottom.}\label{org466}
\begin{Verbatim}[breaklines,breakautoindent=false,breakindent=1em,breakafter=/-.\,;:=,commandchars=\\\{\}]
 Dear Aroha,
 Thank you for the books. I will bring
 them on Saturday.

 TUATARA[N]  council.txt                 L3  C1  Text
 ~/letters$ ls
 council.txt  thanks.txt
 ~/letters$ wc -w council.txt
 41 council.txt
 ~/letters$ _
 TUATARA[I] P *terminal*  L5 C13  Terminal  ~/letters
\end{Verbatim}
\end{figure}

The lower window is the terminal. Its mode line says \texttt{Terminal}, and
\texttt{TUATARA[I]} says you are in INSERT state, typing into the program. The text
after \texttt{Terminal} is the title the shell gave its window, if it gave one (many
set \texttt{user@\codebreak{}host:\textasciitilde{}/\codebreak{}dir}).

\subsubsection{Where it appears}
\label{sec:467}
The bottom is the standard place. \texttt{*terminal-\codebreak{}place*} says where terminals
show: \texttt{:bottom}, \texttt{:top}, \texttt{:left}, \texttt{:right} or \texttt{:full}. The first four give a
window along that whole side of the screen. \texttt{:full} uses the window you are in,
and when you quit that terminal, the window shows what it showed before.

\texttt{*terminal-\codebreak{}size*} is how high the window is at the top or bottom, as a
fraction of the screen (\texttt{0.3}) or a number of rows (\texttt{15}). \texttt{*terminal-\codebreak{}width*}
is how wide it is at the left or right, as a fraction (\texttt{0.5}, half the screen)
or a number of columns (\texttt{80}).

Five commands open a terminal in a place of your choosing, whatever
\texttt{*terminal-\codebreak{}place*} says: \texttt{terminal-\codebreak{}bottom}, \texttt{terminal-\codebreak{}top}, \texttt{terminal-\codebreak{}left},
\texttt{terminal-\codebreak{}right} and \texttt{terminal-\codebreak{}full}. Use them when you want a terminal beside
a long document today, and your usual place is the bottom.

By default the shell is your own: \texttt{\$SHELL}, or \texttt{/bin/sh} if that is not set.
\texttt{*terminal-\codebreak{}program*} names another program with its arguments, for example
\texttt{("bash" "-l")} or \texttt{("fish")}.

\subsubsection{Showing and hiding}
\label{sec:468}
You often want the terminal for a few seconds. \texttt{M-x terminal-\codebreak{}toggle} is made
for that. It shows the terminal you opened last and types into it. Run it again
from inside the terminal and the window closes. The program goes on while it
is hidden, and when you show it again, the screen is as you left it. If there
is no terminal yet, it opens one.

\texttt{terminal-\codebreak{}toggle} has no key to begin with. If you use a terminal often, give
it one in your init file (chapter 20).

\begin{xltabular}{\textwidth}{@{}l l l >{\raggedright\arraybackslash\hsize=1.000\hsize}X@{}}
\toprule
\textbf{Modal} & \textbf{Non-modal} & \textbf{Command} & \textbf{What it does}\\
\midrule
\endhead
none & \texttt{M-x terminal} & \texttt{terminal} & Open a terminal in a buffer of its own, where \textbf{TERMINAL-PLACE} says, as OPEN-TERMINAL does.\\
none & \texttt{M-x terminal-\codebreak{}bottom} & \texttt{terminal-\codebreak{}bottom} & Open a new terminal in a window across the whole bottom, whatever \textbf{TERMINAL-PLACE} says, \textbf{TERMINAL-SIZE} high, as M-x terminal opens one.\\
none & \texttt{M-x terminal-\codebreak{}top} & \texttt{terminal-\codebreak{}top} & Open a new terminal in a window across the whole top, whatever \textbf{TERMINAL-PLACE} says, \textbf{TERMINAL-SIZE} high, as M-x terminal opens one.\\
none & \texttt{M-x terminal-\codebreak{}left} & \texttt{terminal-\codebreak{}left} & Open a new terminal in a window down the whole left side, whatever \textbf{TERMINAL-PLACE} says, \textbf{TERMINAL-WIDTH} wide, as M-x terminal opens one.\\
none & \texttt{M-x terminal-\codebreak{}right} & \texttt{terminal-\codebreak{}right} & Open a new terminal in a window down the whole right side, whatever \textbf{TERMINAL-PLACE} says, \textbf{TERMINAL-WIDTH} wide, as M-x terminal opens one.\\
none & \texttt{M-x terminal-\codebreak{}full} & \texttt{terminal-\codebreak{}full} & Open a new terminal in the selected window, whatever \textbf{TERMINAL-PLACE} says, as M-x terminal opens one; quitting it brings back what the window showed.\\
none & \texttt{M-x terminal-\codebreak{}toggle} & \texttt{terminal-\codebreak{}toggle} & Show the terminal and type into it, or hide it again: the terminal opened last, where \textbf{TERMINAL-PLACE} says (a new one when there is none); selected when it shows already; hidden when you are in it, its window closed (or, at :full, the buffer before it back).\\
\bottomrule
\end{xltabular}

\begin{xltabular}{\textwidth}{@{}l l >{\raggedright\arraybackslash\hsize=1.000\hsize}X@{}}
\toprule
\textbf{Setting} & \textbf{Standard} & \textbf{What it changes}\\
\midrule
\endhead
\texttt{*terminal-\codebreak{}place*} & \texttt{:bottom} & Where a terminal shows (M-x terminal, M-x terminal-toggle): :bottom or :top, a window across that whole side of the screen \textbf{TERMINAL-SIZE} high; :left or :right, one down that whole side \textbf{TERMINAL-WIDTH} wide; or :full, the selected window.\\
\texttt{*terminal-\codebreak{}size*} & \texttt{0.3} & How high a terminal’s window is at the bottom or top of the screen (\textbf{TERMINAL-PLACE}, M-x terminal-bottom, M-x terminal-top): a fraction of the screen’s height (0.3), or a number of rows (15).\\
\texttt{*terminal-\codebreak{}width*} & \texttt{0.5} & How wide a terminal’s window is at the left or right of the screen (\textbf{TERMINAL-PLACE}, M-x terminal-left, M-x terminal-right): a fraction of the screen’s width (0.5, half of it), or a number of columns (80).\\
\texttt{*terminal-\codebreak{}program*} & \texttt{nil} & The program M-x terminal runs, a list of it and its arguments, (“bash” ”-l”); NIL for your shell (\$SHELL, else /bin/sh).\\
\bottomrule
\end{xltabular}

\begin{try}
\textbf{Try it, running a command while you write}
\begin{enumerate}
\item Press \texttt{C-x C-f}, type \texttt{\textasciitilde{}/letter.\codebreak{}txt} and press \texttt{RET}. Press \texttt{i} and type
   \texttt{Dear Aroha,}. Press \texttt{C-g} and save with \texttt{C-x C-s}.
\item Run \texttt{M-x terminal}. A window opens across the bottom with a prompt.
\item Type \texttt{ls} and press \texttt{RET}. Your files are listed, \texttt{letter.\codebreak{}txt} among them.
\item Press \texttt{C-x o}. You are back in the letter. The shell is still there.
\item Run \texttt{M-x terminal-\codebreak{}toggle}. The terminal is selected and you are typing
   into it again.
\item Run \texttt{M-x terminal-\codebreak{}toggle} once more. The window closes, and the letter has
   the whole screen.
\item Run \texttt{M-x terminal-\codebreak{}toggle} again. The shell comes back as it was, with your
   \texttt{ls} still on its screen.
\item Type \texttt{exit} and press \texttt{RET}. The terminal and its window are gone.
\end{enumerate}
\end{try}

\subsection{Typing into a terminal}
\label{sec:471}
While you type into a terminal, every key goes to the program. That includes
Escape and \texttt{C-c}, which a shell uses to give up on a command, and the arrow
keys, which walk back through the commands you have typed. The text you type
does not go into the buffer. It goes to the program, and the program prints the
result.

A handful of keys stay the editor’s, so that you can always leave. The
setting \texttt{*terminal-\codebreak{}editor-\codebreak{}keys*} lists them. As standard they are:

\begin{longtable}{@{}l l l@{}}
\toprule
\textbf{Modal} & \textbf{Non-modal} & \textbf{What it does in a terminal}\\
\midrule
\endhead
none & \texttt{C-x} & The same as everywhere: windows, buffers, files and the rest\\
none & \texttt{M-x} & Run a command by name\\
none & \texttt{C-g} & Stop typing and look at the history (next section)\\
none & \texttt{C-\textbackslash{}} & The same as \texttt{C-g}, and also types into the terminal again\\
none & \texttt{C-q} & Send the next key to the program, whatever it is\\
none & \texttt{C-s} & Search the terminal’s history\\
\bottomrule
\end{longtable}

Everything else goes to the program. If your program needs one of these keys
itself, you have two ways out. \texttt{C-q} sends one key, once. Taking the key out of
\texttt{*terminal-\codebreak{}editor-\codebreak{}keys*} frees it for good. The next sections show both.

Each key is written as you would write it for \texttt{bind-\codebreak{}key} in chapter 20:
\texttt{"C-x"}, \texttt{"M-x"}, and so on. Your init file can change the list.

Programs that fill the screen, such as \texttt{less}, \texttt{htop} and \texttt{nvim}, work as they
do in any terminal. The mouse works too, when the program asks for it. Hold
\texttt{Shift} while you click if you would rather the editor take the mouse.

\begin{xltabular}{\textwidth}{@{}l >{\raggedright\arraybackslash\hsize=0.681\hsize}X >{\raggedright\arraybackslash\hsize=1.319\hsize}X@{}}
\toprule
\textbf{Setting} & \textbf{Standard} & \textbf{What it changes}\\
\midrule
\endhead
\texttt{*terminal-\codebreak{}editor-\codebreak{}keys*} & \texttt{("C-x" "M-x" "C-g" "C-\codebreak{}\textbackslash{}\textbackslash{}" "C-q" "C-\codebreak{}s")} & Keys that stay the editor’s while you type into a terminal (M-x terminal).\\
\bottomrule
\end{xltabular}

\subsubsection{Sending a key the editor keeps}
\label{sec:470}
You are running \texttt{nano} in a terminal, and it uses \texttt{C-x} to exit. Pressing
\texttt{C-x} would start an editor key sequence instead. Press \texttt{C-q} first. The echo
area says \texttt{Key for the program:}. The next key you press goes to \texttt{nano},
whatever it is. \texttt{C-q C-x} gives \texttt{nano} its \texttt{C-x}. \texttt{C-q C-g} gives the program
its \texttt{C-g}, which bash takes as “give up what you were doing” without leaving
what you typed.

If a program needs \texttt{C-g} all the time, as another editor does, take it out of
the list in your init file:

\begin{Verbatim}[breaklines,breakautoindent=false,breakindent=1em,breakafter=/-.\,;:=,commandchars=\\\{\}]
\color[HTML]{A85F00}(setf\normalcolor{} \color[HTML]{A85F00}*terminal-editor-keys*
      \color[HTML]{A85F00}(remove\normalcolor{} \color[HTML]{0A7AB3}"C-g"\normalcolor{} \color[HTML]{A85F00}*terminal-editor-keys*\normalcolor{} \color[HTML]{946C00}:\normalcolor{}test \color[HTML]{222222}#'\color[HTML]{A85F00}string=))
\end{Verbatim}

\texttt{C-g} then goes to the program, and \texttt{C-\textbackslash{}} alone stops typing.

\begin{xltabular}{\textwidth}{@{}l l l >{\raggedright\arraybackslash\hsize=1.000\hsize}X@{}}
\toprule
\textbf{Modal} & \textbf{Non-modal} & \textbf{Command} & \textbf{What it does}\\
\midrule
\endhead
none & \texttt{C-q} & \texttt{terminal-\codebreak{}send-\codebreak{}key} & Send the next key typed to the program in this terminal, whatever it is: one of those that stay the editor’s (C-g, C-x, M-x, C-q), say.\\
\bottomrule
\end{xltabular}

\subsection{The history view}
\label{sec:474}
A program prints more than fits on a screen. A long \texttt{ls} scrolls off the top.
A build prints hundreds of lines. In an ordinary terminal you scroll back with
the mouse. In Tuatara, press \texttt{C-g} or \texttt{C-\textbackslash{}}. The terminal turns into its
\emph{history view}: everything the program printed, including what went up off its
screen, shown as plain text with its colours. The program keeps running behind
it. The mode line says \texttt{Terminal history}, and the echo area says which key
types into it again.

In the history view you have the keys of any text. You can move with \texttt{j}, \texttt{k},
\texttt{h} and \texttt{l}, search with \texttt{C-s} and \texttt{C-r}, and select with \texttt{V} (chapter 3).
You press \texttt{y} to copy (chapter 4). Nothing you press goes to the program, so
there is no danger of typing a command by accident.

Figure 17.2 shows what you can do, with a build log as the example.

\begin{figure}[H]
\caption{Figure 17.2. Typing, and the history view.}\label{org472}
\begin{Verbatim}[breaklines,breakautoindent=false,breakindent=1em,breakafter=/-.\,;:=,commandchars=\\\{\}]
   typing                              history view
   ------                              ------------
   keys go to the program              keys move, search, select, copy
   mode line: Terminal                 mode line: Terminal history

        C-g or C-\char92{}  ------------------>
        <------------------  i, a or C-\char92{}
        <------------------  RET  (copies the selection first)
\end{Verbatim}
\end{figure}

To go back to typing, press \texttt{i} or \texttt{a}, or \texttt{C-\textbackslash{}} again. \texttt{RET} does the same,
and copies first whatever you selected, so that what you found is ready to
paste.

The history keeps 100,000 lines of what went off the screen as standard.
\texttt{*terminal-\codebreak{}scrollback*} changes the number. A terminal that has not printed
anything for ten seconds has its older history packed away in the background,
a little at a time, so that a long history costs a fraction of the memory.
Reading it unpacks it again without you noticing. \texttt{*terminal-\codebreak{}compress-\codebreak{}after*}
sets the ten seconds, and \texttt{nil} turns the packing off.

If you scroll the mouse wheel up while a program that does not read the wheel
is on the screen, Tuatara shows the history view for you, so that scrolling
back works as you expect. A pager such as \texttt{less} gets the wheel as the up and
down keys instead.

\subsubsection{Searching the history}
\label{sec:473}
\texttt{C-s} while you type does not go to the program. It opens the history view and
starts a search, back from where the program’s cursor is. Look for the line
you want, press \texttt{RET} to leave the search, and carry on in the history view.
In the history view, \texttt{C-s} and \texttt{C-r} search forward and back as in any text
(chapter 3).

This is the quickest way to find an error message in a long build.

\begin{xltabular}{\textwidth}{@{}l >{\raggedright\arraybackslash\hsize=0.622\hsize}X l >{\raggedright\arraybackslash\hsize=1.378\hsize}X@{}}
\toprule
\textbf{Modal} & \textbf{Non-modal} & \textbf{Command} & \textbf{What it does}\\
\midrule
\endhead
none & \texttt{M-x terminal-\codebreak{}stop-\codebreak{}typing} & \texttt{terminal-\codebreak{}stop-\codebreak{}typing} & Stop typing into the terminal: its history view, its whole history (what went up off its screen too) standing still as text with its colours, to move through, search, select and copy while its program goes on; with modal editing on in SUPPORT state (NORMAL’s keys that read and select, and the leader), else with the keys of any text.\\
none & \texttt{a}, \texttt{i}, \texttt{C-\textbackslash{}} & \texttt{terminal-\codebreak{}start-\codebreak{}typing} & Type into the terminal again after its history view (M-x terminal-stop-typing): everything typed goes to its program, the screen shown as it is now.\\
none & \texttt{RET} & \texttt{terminal-\codebreak{}copy-\codebreak{}and-\codebreak{}type} & In a terminal’s history view, copy what is selected (as M-w) and type into the terminal again; with nothing selected, type into it again.\\
none & \texttt{M-x terminal-\codebreak{}search} & \texttt{terminal-\codebreak{}search} & Search the terminal’s whole history, from where its program’s cursor is back (C-r as anywhere): typing there, its history view is shown first and the search starts once it is in; in the history view, search on from point (C-s).\\
\bottomrule
\end{xltabular}

\begin{xltabular}{\textwidth}{@{}l l >{\raggedright\arraybackslash\hsize=1.000\hsize}X@{}}
\toprule
\textbf{Setting} & \textbf{Standard} & \textbf{What it changes}\\
\midrule
\endhead
\texttt{*terminal-\codebreak{}scrollback*} & \texttt{100000} & How many lines a terminal (M-x terminal) keeps that have gone up off its screen.\\
\texttt{*terminal-\codebreak{}compress-\codebreak{}after*} & \texttt{10} & Seconds a terminal’s program has printed nothing after which its scrollback is compressed, in the background, a little at a time, so that a terminal with a long history holds a fraction of the memory; reading the history (its history view) undoes it as it goes, unseen.\\
\bottomrule
\end{xltabular}

\begin{try}
\textbf{Try it, looking back through output}
\begin{enumerate}
\item Run \texttt{M-x terminal}. Type \texttt{ls -l /usr/\codebreak{}bin} and press \texttt{RET}. Most of the
   list has gone off the top of the window.
\item Press \texttt{C-g}. The mode line says \texttt{Terminal history}.
\item Press \texttt{k} a few times to move up. You are reading the whole list as text,
   and the lines that went off the top are there.
\item Press \texttt{C-r}, type \texttt{bash} and press \texttt{RET}. The cursor moves to the nearest
   line above with \texttt{bash} in it.
\item Press \texttt{V} to select the line, then \texttt{RET}. The line is copied and you are
   typing into the shell again.
\item Type \texttt{echo} and a space, then press \texttt{C-V} (Control, Shift and \texttt{v}). The line
   is pasted after \texttt{echo}. It is not run. Press \texttt{C-u} to clear the line.
\end{enumerate}
\end{try}

\subsection{Pasting}
\label{sec:475}
Pasting into a terminal is different from pasting into a file. What you paste
goes to a program that may run it. A shell, given three lines, may run all
three at once, and a line that holds a hidden command could do harm.

Tuatara pastes carefully. The text is marked as pasted, which a shell such as
bash understands. It shows the lines at the prompt and runs none of them until
you press \texttt{RET}. Control characters in the text become spaces, so that pasted
text cannot send commands to the terminal itself. A program that does not mark
pastes would run each line as it arrives, so Tuatara asks first when the text
has more than one line.

While you type into the terminal, \texttt{C-V} and \texttt{S-<insert>} paste the text you
copied last, in Tuatara or in another program. \texttt{C-V} is Control, Shift and \texttt{v}
together, which is the paste key in most terminals. \texttt{*terminal-\codebreak{}paste-\codebreak{}keys*}
lists the keys, and you can change it. In the history view, \texttt{p} and \texttt{C-y} paste
and go back to typing. The middle mouse button pastes whatever you selected
last in any program, the way it does in a desktop terminal.

\begin{xltabular}{\textwidth}{@{}l l l >{\raggedright\arraybackslash\hsize=1.000\hsize}X@{}}
\toprule
\textbf{Modal} & \textbf{Non-modal} & \textbf{Command} & \textbf{What it does}\\
\midrule
\endhead
none & \texttt{p}, \texttt{C-y} & \texttt{terminal-\codebreak{}paste} & Paste into the terminal’s program the text copied last, in Tuatara or another program (as C-y takes it), and type into the terminal: marked as pasted for a program that asks, so that a shell runs none of its lines until you press RET, and its control characters made spaces, so that pasted text cannot work the terminal.\\
none & \texttt{<mouse-2>} & \texttt{terminal-\codebreak{}paste-\codebreak{}primary} & Paste into the terminal’s program the text selected last in any program (the primary selection), as the middle button does, and type into it.\\
\bottomrule
\end{xltabular}

\begin{xltabular}{\textwidth}{@{}l l >{\raggedright\arraybackslash\hsize=1.000\hsize}X@{}}
\toprule
\textbf{Setting} & \textbf{Standard} & \textbf{What it changes}\\
\midrule
\endhead
\texttt{*terminal-\codebreak{}paste-\codebreak{}keys*} & \texttt{("C-V" "S-<insert>")} & Keys that paste into a terminal while you type there (M-x terminal-paste): the text copied last, in Tuatara or another program, marked as pasted for a program that asks (a shell then runs none of its lines until you press RET).\\
\bottomrule
\end{xltabular}

\begin{try}
\textbf{Try it, pasting a command from a note}
\begin{enumerate}
\item Open \texttt{\textasciitilde{}/notes.\codebreak{}txt}, press \texttt{i} and type \texttt{ls -l \textasciitilde{}}. Press \texttt{C-g}.
\item Put the cursor on the line, press \texttt{V} and \texttt{y} to copy it. It says \texttt{Copied}.
\item Run \texttt{M-x terminal}. Press \texttt{C-V}. The command appears at the prompt. It has
   not run.
\item Press \texttt{RET}. The shell lists your home folder.
\end{enumerate}
\end{try}

\subsection{Links}
\label{sec:476}
Programs print web addresses and file names. A terminal buffer shows them
underlined, as \texttt{:link}. Press \texttt{C-c C-o} with the cursor on one, or hold
\texttt{Control} and click it with the mouse. Tuatara asks before doing anything, and
the question names where the link leads, because a link can show one thing and
go to another.

A file on your computer opens in Tuatara, in another window. A web or mail
address opens in your desktop’s own program. Programs can also write links of
their own (\texttt{ls --hyperlink}), and Tuatara follows those the same way.

The history view works the same: put the cursor on a link and press \texttt{C-c C-o}.

\begin{xltabular}{\textwidth}{@{}l >{\raggedright\arraybackslash\hsize=0.531\hsize}X l >{\raggedright\arraybackslash\hsize=1.469\hsize}X@{}}
\toprule
\textbf{Modal} & \textbf{Non-modal} & \textbf{Command} & \textbf{What it does}\\
\midrule
\endhead
none & \texttt{C-c C-o}, \texttt{C-<down-\codebreak{}mouse-1>} & \texttt{terminal-\codebreak{}open-\codebreak{}link} & Open the link under the mouse (C-click) or at point in a terminal: a link a program wrote (ls --hyperlink, OSC 8) or a web or file address written there.\\
\bottomrule
\end{xltabular}

\begin{longtable}{@{}l l@{}}
\toprule
\textbf{Face} & \textbf{What it paints}\\
\midrule
\endhead
\texttt{:link} & Links and URLs.\\
\bottomrule
\end{longtable}

\subsection{Pictures and colours from programs}
\label{sec:477}
A program that draws pictures with the kitty graphics protocol, such as
\texttt{timg} or \texttt{kitten icat}, can draw them in a terminal buffer, if Tuatara itself
can show pictures where it runs (chapter 16). In a window of your desktop it
can. In a terminal it can when that terminal can, such as kitty or Ghostty.
Where it cannot, such a program draws with text, as it would anywhere else.

Pictures take memory. \texttt{*terminal-\codebreak{}picture-\codebreak{}memory*} is the most a terminal keeps,
64 MiB as standard. Set it to \texttt{0} and programs always draw with text. A file
meant to slow a terminal down, by drawing a picture thousands of times, has
its pictures dropped when you print it with \texttt{cat}. The text stays.

Programs choose colours and styles, and Tuatara shows them as faces. Bold,
italic, underlined, faint and reversed text have a face each, listed below. The
sixteen basic colours a program may choose follow your theme, so the output of
\texttt{ls --color} looks at home next to your text. Colours beyond those sixteen
are shown as the program asked, rounded to the nearest of 4,096.

\begin{xltabular}{\textwidth}{@{}l l >{\raggedright\arraybackslash\hsize=1.000\hsize}X@{}}
\toprule
\textbf{Setting} & \textbf{Standard} & \textbf{What it changes}\\
\midrule
\endhead
\texttt{*terminal-\codebreak{}picture-\codebreak{}memory*} & \texttt{67108864} & The most octets of pictures a terminal keeps for the programs in it that draw them; 0 takes none, programs then drawing with text.\\
\bottomrule
\end{xltabular}

\begin{xltabular}{\textwidth}{@{}l >{\raggedright\arraybackslash\hsize=1.000\hsize}X@{}}
\toprule
\textbf{Face} & \textbf{What it paints}\\
\midrule
\endhead
\texttt{:terminal-\codebreak{}bold} & Text a program in a terminal set in bold.\\
\texttt{:terminal-\codebreak{}dim} & Text a program in a terminal set faint.\\
\texttt{:terminal-\codebreak{}italic} & Text a program in a terminal set in italics.\\
\texttt{:terminal-\codebreak{}reverse} & Text a program in a terminal shown in reverse, its colours swapped.\\
\texttt{:terminal-\codebreak{}underline} & Text a program in a terminal underlined.\\
\bottomrule
\end{xltabular}

\subsection{The title}
\label{sec:478}
Many shells set a title for their window, usually your name, the machine and
the folder. Tuatara shows it in the terminal’s mode line, after \texttt{Terminal}.
It is cut at 40 columns. \texttt{*terminal-\codebreak{}title-\codebreak{}limit*} changes that, and \texttt{nil} hides
the title. When you have several terminals open, the title tells you which is
in which folder.

\begin{xltabular}{\textwidth}{@{}l l >{\raggedright\arraybackslash\hsize=1.000\hsize}X@{}}
\toprule
\textbf{Setting} & \textbf{Standard} & \textbf{What it changes}\\
\midrule
\endhead
\texttt{*terminal-\codebreak{}title-\codebreak{}limit*} & \texttt{40} & The most columns of a terminal’s title (as the shell gives it, user@host:\textasciitilde{}/dir in many) the mode line shows; NIL shows none.\\
\bottomrule
\end{xltabular}

\subsection{Escape, nvim and the shell}
\label{sec:479}
Escape is the awkward key. A shell uses it for editing the line, and a
full-screen editor such as \texttt{nvim} uses it all the time to leave insert mode.
Tuatara, in its own state keys, may use Escape to go back to NORMAL state.

As standard, \texttt{*escape-\codebreak{}quits*} is off (chapter 8), and in a terminal Escape goes
to the program. In \texttt{nvim}, Escape does what you expect. To stop typing and
look at the history, press \texttt{C-g} or \texttt{C-\textbackslash{}}.

If you turned \texttt{*escape-\codebreak{}quits*} on, Escape stops typing in a terminal, as \texttt{C-g}
does, and the program never sees it. Two things then help. \texttt{C-q} followed by
Escape sends one Escape to the program. And \texttt{*terminal-\codebreak{}send-\codebreak{}escape*}, set to
\texttt{t}, gives Escape to a full-screen program (one on the alternate screen, such
as \texttt{nvim}, \texttt{less} or \texttt{htop}) while a shell at its prompt still lets Escape stop
typing. That is usually what you want with both settings on: Escape works
inside \texttt{nvim}, and at a shell it takes you to the history. \texttt{C-g} stops typing
in either case.

\begin{Verbatim}[breaklines,breakautoindent=false,breakindent=1em,breakafter=/-.\,;:=,commandchars=\\\{\}]
\color[HTML]{A85F00}(setf\normalcolor{} \color[HTML]{A85F00}*escape-quits*\normalcolor{} \color[HTML]{A85F00}t
      \color[HTML]{A85F00}*terminal-send-escape*\normalcolor{} \color[HTML]{A85F00}t)
\end{Verbatim}

\texttt{M-x terminal-\codebreak{}send-\codebreak{}escape} sends Escape to the program whatever it is, at a
shell prompt too. It has no key. If you want one that works while you type,
add it to the list of editor keys and bind it for terminals:

\begin{Verbatim}[breaklines,breakautoindent=false,breakindent=1em,breakafter=/-.\,;:=,commandchars=\\\{\}]
\color[HTML]{A85F00}(push\normalcolor{} \color[HTML]{0A7AB3}"C-;"\normalcolor{} \color[HTML]{A85F00}*terminal-editor-keys*)
\color[HTML]{A85F00}(bind-key\normalcolor{} \color[HTML]{0A7AB3}"C-;"\normalcolor{} \color[HTML]{222222}'\color[HTML]{A85F00}terminal-send-escape\normalcolor{} \color[HTML]{946C00}:\normalcolor{}mode \color[HTML]{222222}'\color[HTML]{A85F00}terminal)
\end{Verbatim}

\begin{xltabular}{\textwidth}{@{}l >{\raggedright\arraybackslash\hsize=0.596\hsize}X l >{\raggedright\arraybackslash\hsize=1.404\hsize}X@{}}
\toprule
\textbf{Modal} & \textbf{Non-modal} & \textbf{Command} & \textbf{What it does}\\
\midrule
\endhead
none & \texttt{M-x terminal-\codebreak{}send-\codebreak{}escape} & \texttt{terminal-\codebreak{}send-\codebreak{}escape} & Send Escape to the program in this terminal, at a shell’s prompt too, whatever \textbf{ESCAPE-QUITS} and \textbf{TERMINAL-SEND-ESCAPE} say.\\
\bottomrule
\end{xltabular}

\begin{xltabular}{\textwidth}{@{}l l >{\raggedright\arraybackslash\hsize=1.000\hsize}X@{}}
\toprule
\textbf{Setting} & \textbf{Standard} & \textbf{What it changes}\\
\midrule
\endhead
\texttt{*terminal-\codebreak{}send-\codebreak{}escape*} & \texttt{nil} & True to give Escape to a full-screen program in a terminal (one on the alternate screen: nvim, less, htop) even while \textbf{ESCAPE-QUITS} is on, Escape then still stopping typing at a shell’s prompt.\\
\bottomrule
\end{xltabular}

\begin{try}
\textbf{Try it, with a program like top}
\begin{enumerate}
\item Run \texttt{M-x terminal}. Type \texttt{top} and press \texttt{RET}. The screen fills with the
   list of programs, updating as you watch.
\item Press \texttt{C-g}. Typing stops. The mode line says \texttt{Terminal history}, and the
   screen text stands still to read, while \texttt{top} goes on behind it.
\item Press \texttt{i}. You are typing into \texttt{top} again and the screen is current.
\item Press \texttt{C-q}, then \texttt{q}. The echo area asks \texttt{Key for the program:} and the \texttt{q}
   goes to \texttt{top}, which ends. You are back at the shell.
\end{enumerate}
\end{try}

\subsection{The terminal library}
\label{sec:480}
The terminal buffer reads what programs print with a library from Ghostty.
\texttt{M-x terminal-\codebreak{}version} says which version your Tuatara runs on. It is useful
when you report a problem. If the library is missing, \texttt{M-x terminal} says so
instead of opening a terminal.

\begin{xltabular}{\textwidth}{@{}l l l >{\raggedright\arraybackslash\hsize=1.000\hsize}X@{}}
\toprule
\textbf{Modal} & \textbf{Non-modal} & \textbf{Command} & \textbf{What it does}\\
\midrule
\endhead
none & \texttt{M-x terminal-\codebreak{}version} & \texttt{terminal-\codebreak{}version} & Say which version of the terminal library (libghostty-vt, from Ghostty) M-x terminal runs on.\\
\bottomrule
\end{xltabular}

\subsection{Running Tuatara inside a terminal}
\label{sec:484}
Many people start Tuatara from the shell of their desktop terminal with
\texttt{tuatara}, and edit in that window. It works the same as the window version.
A few things depend on the terminal you use.

\subsubsection{Keys}
\label{sec:481}
Old terminals cannot tell some key presses apart. Escape and the start of an Alt
key both send an escape character, so a program has to wait a moment to see
whether another character follows. Tab and \texttt{C-i} are the same byte, and so are
\texttt{RET} and \texttt{C-m}.

Newer terminals, among them kitty, foot, WezTerm and Ghostty, can send keys
whole, using the kitty keyboard protocol. With it, Tuatara knows that Escape is
Escape, that \texttt{M-x} is \texttt{M-x}, and nothing waits 25 milliseconds. The setting
\texttt{*kitty-\codebreak{}keyboard*} decides whether Tuatara uses the protocol. The standard,
\texttt{:auto}, asks the terminal when Tuatara starts and uses it if the answer is yes.
Set it to \texttt{t} to use it always, or to \texttt{nil} never. You only need the setting if
your terminal says it speaks the protocol but does not work with it, or the
other way round.

If your terminal does not speak it, everything works, with one small delay
after Escape while Tuatara waits to see what follows.

\begin{xltabular}{\textwidth}{@{}l l >{\raggedright\arraybackslash\hsize=1.000\hsize}X@{}}
\toprule
\textbf{Setting} & \textbf{Standard} & \textbf{What it changes}\\
\midrule
\endhead
\texttt{*kitty-\codebreak{}keyboard*} & \texttt{:auto} & Whether the terminal frontend reads keys by the kitty keyboard protocol.\\
\bottomrule
\end{xltabular}

\subsubsection{Colours and pictures}
\label{sec:482}
\texttt{*tui-\codebreak{}colours*} says how many colours Tuatara uses. The standard, \texttt{:auto},
goes by what the terminal says. If the \texttt{COLORTERM} variable says \texttt{truecolor}
or \texttt{24bit}, your theme’s colours show exactly. Otherwise Tuatara uses as many as
\texttt{TERM} allows, and none where \texttt{NO\_\codebreak{}COLOR} is set. Setting it to \texttt{:24-\codebreak{}bit}, \texttt{256},
\texttt{16} or \texttt{8} overrides that, which helps when a terminal is wrong about itself,
for instance over a remote connection that does not pass \texttt{COLORTERM} on.

Pictures (chapter 16) need a terminal that speaks the kitty graphics protocol.
Under tmux or screen they stay off. The setting \texttt{*tui-\codebreak{}pictures*} is described
in chapter 16.

\subsubsection{Terminal inside the terminal}
\label{sec:483}
You can run \texttt{M-x terminal} while Tuatara runs in a terminal. A shell in a
terminal buffer in Tuatara in a terminal of your desktop is quite normal, and
works.

\begin{longtable}{@{}l l l@{}}
\toprule
\textbf{Setting} & \textbf{Standard} & \textbf{What it changes}\\
\midrule
\endhead
\texttt{*tui-\codebreak{}colours*} & \texttt{:auto} & How many colours Tuatara uses in a terminal.\\
\bottomrule
\end{longtable}

\section{The daemon and other machines}
\label{sec:513}
One editor that keeps running, and the screens that show it.

Until now each time you started \texttt{tuatara} you got an editor of its own. Its
buffers lived as long as it did. When you closed the terminal they were gone,
unless the session setting of chapter 7 brought them back. This chapter shows
another way. You start one editor with no screen of its own, and then look at
it through as many screens as you like. That editor is called a daemon. It
keeps running when you close a screen, so your buffers, your undo history and
your searches are still there when you come back.

The same idea reaches further. The daemon can run on another computer, a
server, and you can look at it from your own desk. The files, the programs and
the work all stay on the server. Your computer draws the screen and sends your
keys. This is not the same as opening a single file over ssh, which chapter 9
covers. There the editor runs here and only the file is far away. Here the
whole editor is far away.

You will need a terminal for the commands at the start. Everything inside the
editor is pressed in NORMAL state unless it starts with \texttt{C-x} or \texttt{C-c}. \texttt{i}
starts INSERT state, and \texttt{C-g} or \texttt{ESC} returns to NORMAL state.

\subsection{One editor in the background}
\label{sec:492}
A daemon is useful when you work on the same things all day and open and close
terminals as you go. You start it once, in the morning. Each terminal you open
after that shows it. Close the terminal at lunch and the editor does not
notice. Open another in the afternoon and everything is as you left it.

\subsubsection{The commands you type in a terminal}
\label{sec:486}
These are not editor commands. You type them in your shell, the program that
reads what you type in a terminal.

\begin{itemize}
\item \texttt{tuatara --daemon} starts the editor in the background and gives you your
  prompt back once it is ready. It reads your init file and loads your
  packages first, which can take a moment, and then says that the daemon runs.
\item \texttt{tuatara --client} shows the daemon in this terminal. With file names after
  it, those files are opened in the daemon. A client reads no init file of its
  own, because the daemon has already read yours.
\item \texttt{tuatara --client --start} does the same, but starts the daemon first if
  none is running. Many people use only this one.
\item \texttt{tuatara --client --gui} shows the daemon in a window of its own instead of
  in the terminal.
\item \texttt{tuatara --stop} stops the daemon. It asks about any change you have not
  saved, as described under “Stopping the daemon” below.
\end{itemize}

Words that mean the same as for plain \texttt{tuatara} work after \texttt{--client} too: a
line number such as \texttt{+40} before a file name, \texttt{--safe} and \texttt{--hex}. Chapter 12
explains the safe visit.

A \emph{client} is one screen showing the daemon, with a frame of its own and
buffers of its own. When you close a client, its buffers stay in the daemon.
Another client can take them up. A client started while the daemon is still
starting waits for it and says so.

\subsubsection{Names for several daemons}
\label{sec:487}
You can run more than one daemon. Give each a name, made of letters, digits,
\texttt{-} and \texttt{\_}, at most forty characters long. A daemon with no name is called
\texttt{default}.

\begin{Verbatim}[breaklines,breakautoindent=false,breakindent=1em,breakafter=/-.\,;:=,commandchars=\\\{\}]
tuatara --daemon=work
tuatara --client=work
tuatara --stop=work
\end{Verbatim}

One daemon for your work and another for your home projects keeps their
buffers apart. \texttt{--start} goes with a name as well: \texttt{tuatara --client=\codebreak{}work
--start}.

\subsubsection{Closing a client}
\label{sec:488}
\texttt{C-x C-c} in a client closes that client. Its screen goes, your terminal is
yours again, and the daemon carries on. Buffers with changes you have not
saved stay as they are, ready for the next client. To end the daemon itself,
use \texttt{daemon-\codebreak{}quit} or \texttt{tuatara --stop}.

A client also closes when its last buffer is closed. If you press \texttt{C-x k} on
the only buffer a client has, the client ends as if you had pressed \texttt{C-x C-c}.
The scratch buffer that every new client starts with does not count while you
have typed nothing in it.

\subsubsection{More than one window of a client}
\label{sec:489}
\texttt{C-x 5 2} opens another frame of the same client, with the same buffers. In a
terminal it takes turns with the first in that terminal. In a window of the
desktop it is a second window. \texttt{C-x 5 0} closes the frame you are in. The
client ends when its last frame closes.

\subsubsection{The D in the mode line}
\label{sec:490}
While your editor runs as a daemon, its mode line, and that of each client,
shows a \texttt{D} after the other signs. The sign \texttt{P} before it is for work in the
background that chapter 11 describes. The \texttt{D} is lit while there is a daemon
and faint otherwise, and it never moves, so nothing after it jumps about. The
setting \texttt{*daemon-\codebreak{}indicator*} is the letter, and \texttt{nil} takes it away. The face
\texttt{:mode-\codebreak{}line-\codebreak{}sign} is what it is drawn in.

\subsubsection{Commands, settings and faces for the daemon}
\label{sec:491}
\begin{xltabular}{\textwidth}{@{}l l l >{\raggedright\arraybackslash\hsize=1.000\hsize}X@{}}
\toprule
\textbf{Modal} & \textbf{Non-modal} & \textbf{Command} & \textbf{What it does}\\
\midrule
\endhead
none & \texttt{M-x daemon-\codebreak{}quit} & \texttt{daemon-\codebreak{}quit} & Stop the editor running in the background (tuatara --daemon), from any of its clients, asked first about every change not saved in it, the clients that ended included, folder by folder: s saves those of a folder, l leaves them unsaved, C-g keeps the daemon running; the saves begun are waited for (C-g there keeps it running too).\\
none & \texttt{C-x 5 2} & \texttt{make-\codebreak{}frame} & Open another frame showing the current buffer and work in it: a window of its own on the desktop, or in a terminal a whole screen of its own shown in place of the one before, the frames taking turns on the terminal, the mode line naming the one shown (F1, F2).\\
none & \texttt{C-x 5 0} & \texttt{delete-\codebreak{}frame} & Close the frame worked in, with its windows; its buffers stay, shown in the other frames as before.\\
\bottomrule
\end{xltabular}

\begin{xltabular}{\textwidth}{@{}l l >{\raggedright\arraybackslash\hsize=1.000\hsize}X@{}}
\toprule
\textbf{Setting} & \textbf{Standard} & \textbf{What it changes}\\
\midrule
\endhead
\texttt{*daemon-\codebreak{}indicator*} & \texttt{"D"} & What the mode line shows last among its signs before [-], always in the same place: in :mode-line-sign while this editor runs as a daemon (tuatara --daemon), in each of its clients, else faint, in :mode-line-sign-idle.\\
\bottomrule
\end{xltabular}

\begin{xltabular}{\textwidth}{@{}l >{\raggedright\arraybackslash\hsize=1.000\hsize}X@{}}
\toprule
\textbf{Face} & \textbf{What it paints}\\
\midrule
\endhead
\texttt{:mode-\codebreak{}line-\codebreak{}sign} & The mode line’s signs before [-] while their work goes on: P while a job of the buffer runs or waits its turn (\textbf{processing-indicator}), U while a directory of \textbf{indexed-directories} is being read (\textbf{indexing-indicator}), S while one is searched (\textbf{searching-indicator}), D while the editor runs as a daemon (\textbf{daemon-indicator}); in the theme’s warning colour, to stand out.\\
\bottomrule
\end{xltabular}

\subsection{Clients in the buffer list}
\label{sec:496}
Each client has its own buffers. \texttt{C-x b} offers yours, and the buffers that
belong to nobody, such as \texttt{*Messages*}. \texttt{C-x C-b} shows the list of chapter 7
with a section for each client in place of the sections for projects, so you
can see what each screen has open.

Figure 18.1 shows the list in a client that was started at two o’clock, with
one client that has already ended.

\begin{figure}[H]
\caption{Figure 18.1. The buffer list of a daemon with two clients.}\label{org493}
\begin{Verbatim}[breaklines,breakautoindent=false,breakindent=1em,breakafter=/-.\,;:=,commandchars=\\\{\}]
 {\tuataraFallbackTWO{}▾} This client: client 2 since 14:02, ~/notes (persistent)
   >*  todo.org         1,044  Org     ~/notes/todo.org
       *scratch*            0  Lisp
 {\tuataraFallbackTWO{}▸} client 3 since 14:40, ~/letters
 {\tuataraFallbackTWO{}▸} Ended 15:47 - client 1 since 14:02, ~/notes
 {\tuataraFallbackTWO{}▾} Every client's (1)
       *Messages*       1,210  Text
\end{Verbatim}
\end{figure}

Your own client comes first, marked \texttt{This client}, with the folder it was
started in. Other running clients follow, folded. The clients that have ended
come after them in a faint colour, the last to end first. Their buffers are
still in the daemon, and the list keeps them until you take them up or kill
them. The last section holds the buffers every client is offered.

On the heading of a client that has ended, \texttt{RET} takes its buffers up into
your client. \texttt{d} there marks every buffer of that folder for deletion, and \texttt{x}
kills them, asking first about unsaved changes. Starting \texttt{tuatara --client}
with no file names does the same as \texttt{RET} for the client that ended last. The
face \texttt{:list-\codebreak{}dim} is the faint colour of the ended clients.

\subsubsection{The client that stays}
\label{sec:494}
When the last buffer of a client is closed, the client closes. That would
leave a daemon with no screen to show you. So one client may stay open, left
on its scratch buffer. The setting \texttt{*client-\codebreak{}kept-\codebreak{}without-\codebreak{}buffers*} says which:

\begin{longtable}{@{}l l@{}}
\toprule
\textbf{Value} & \textbf{The client that stays}\\
\midrule
\endhead
\texttt{:first} & the one opened first of those still running (standard)\\
\texttt{:last} & only the last one left\\
\texttt{:all} & every client\\
\texttt{nil} & none; every client closes with its last buffer\\
\bottomrule
\end{longtable}

The heading of the client that stays says \texttt{(persistent)}, as in Figure 18.1.
To choose another, put the cursor in that client’s section of the buffer list
and press \texttt{K}. That client stays until it ends. With \texttt{:all} there is nothing to
choose.

The buffers of encrypted files are an exception. They close with their client,
and what was copied from them goes too, so that their text does not outlive
the person who opened them.

\subsubsection{Commands and settings for clients}
\label{sec:495}
\begin{xltabular}{\textwidth}{@{}l l l >{\raggedright\arraybackslash\hsize=1.000\hsize}X@{}}
\toprule
\textbf{Modal} & \textbf{Non-modal} & \textbf{Command} & \textbf{What it does}\\
\midrule
\endhead
none & \texttt{K}, \texttt{C-c K} & \texttt{buffers-\codebreak{}keep-\codebreak{}client} & Make the client whose folder the cursor is in, in the buffer list, the one that stays when its last buffer is closed, in place of the one \textbf{CLIENT-KEPT-WITHOUT-BUFFERS} names, until it ends.\\
\bottomrule
\end{xltabular}

\begin{xltabular}{\textwidth}{@{}l l >{\raggedright\arraybackslash\hsize=1.000\hsize}X@{}}
\toprule
\textbf{Setting} & \textbf{Standard} & \textbf{What it changes}\\
\midrule
\endhead
\texttt{*client-\codebreak{}kept-\codebreak{}without-\codebreak{}buffers*} & \texttt{:first} & Which clients of a daemon stay when their last buffer is closed: :FIRST, the one opened first of those running; :LAST, only the last one left; :ALL, every one; NIL, none.\\
\bottomrule
\end{xltabular}

\begin{xltabular}{\textwidth}{@{}l >{\raggedright\arraybackslash\hsize=1.000\hsize}X@{}}
\toprule
\textbf{Face} & \textbf{What it paints}\\
\midrule
\endhead
\texttt{:list-\codebreak{}dim} & What is of less interest in a list, set back: help buffers and the headings of clients that ended in the buffer list (C-x C-b), packages not loaded, bookmarks whose file is gone, files ignored in rimu and in a git status, a value missing in the health report (M-x check-health), a line saying more are not shown.\\
\bottomrule
\end{xltabular}

\begin{try}
\textbf{Try it, with one editor kept running all day}
\begin{enumerate}
\item In a terminal, type \texttt{tuatara --client --start \textasciitilde{}/notes.\codebreak{}txt} and press
   \texttt{RET}. The daemon starts and the file opens. The mode line shows \texttt{D}.
\item Press \texttt{i}, type \texttt{buy stamps}, then \texttt{C-g} and \texttt{C-x C-s} to save. Press \texttt{i}
   and type \texttt{call Hemi} without saving. Press \texttt{C-g}.
\item Press \texttt{C-x C-c}. The terminal is yours again. The editor is still running.
\item In another terminal, type \texttt{tuatara --client} and press \texttt{RET}. Press \texttt{C-x
   C-b}. The first client is listed as ended. Put the cursor on its heading
   and press \texttt{RET}. \texttt{notes.\codebreak{}txt} is here with \texttt{call Hemi} still in it.
\item Press \texttt{C-x C-c} again when you are done. The daemon stays for tomorrow.
\end{enumerate}
\end{try}

\subsection{Stopping the daemon}
\label{sec:497}
\texttt{M-x daemon-\codebreak{}quit} in any client, or \texttt{tuatara --stop} in a terminal, stops the
daemon. It asks about every change you have not saved, in every buffer of
every client, those of ended clients included, one folder at a time. For each
folder you press one of three keys:

\begin{itemize}
\item \texttt{s} saves the changes of that folder.
\item \texttt{l} leaves them unsaved. They are lost.
\item \texttt{C-g} stops the question and keeps the daemon running.
\end{itemize}

The daemon ends only when the saves you chose have been written. If you press
\texttt{C-g} while it waits, it carries on running and says that saves are still
under way. Every other client is told that the daemon ends, gives its terminal
back and exits.

\texttt{tuatara --stop} typed where there is no terminal, in a script, cannot ask. It
stops the daemon only when no change would be lost. It writes unsaved changes
as auto-saved copies, if the auto-save setting of chapter 7 allows it. When
some change has no copy, the daemon goes on, and the message says how to stop
it from a terminal. When the system itself asks the daemon to end, as at
logout, there is also no one to ask. It writes its auto-saved copies first and
waits up to thirty seconds for saves already begun.

A daemon told to stop always ends. If it has not ended in 60 seconds, or in 90
when there was no terminal to ask, the daemon is ended all the same, so that a
new daemon of that name can start. Before it goes, it writes in its log
(\texttt{\textasciitilde{}/.\codebreak{}cache/\codebreak{}tuatara/\codebreak{}daemon-\codebreak{}NAME.\codebreak{}log}) which threads were running and the calls
its editor was waiting in. The log names functions only, never what they were
given. Read it when a daemon seems to have stopped late.

\subsection{Sessions for each client}
\label{sec:498}
The session setting \texttt{*keep-\codebreak{}session*} of chapter 7 works for a daemon too. The
daemon keeps one session file for each daemon name, written client by client,
so the windows of one client are never mixed with another’s. If the daemon is
started again, each client’s folder is there as one that ended, and \texttt{tuatara
--client} with no file names takes up the one that ended last. Files opened
safely are not opened again by themselves. They wait, closed, in their folder,
and open safely when a client takes the folder up.

\begin{try}
\textbf{Try it, with the daemon stopped and started again}
\begin{enumerate}
\item Add \texttt{(setf *keep-\codebreak{}session* t)} to your init file and restart the daemon
   with \texttt{tuatara --stop} followed by \texttt{tuatara --daemon}.
\item Type \texttt{tuatara --client \textasciitilde{}/notes.\codebreak{}txt}, then press \texttt{C-x C-c}. Type \texttt{tuatara
   --stop} and answer \texttt{s} if asked.
\item Start the daemon again and type \texttt{tuatara --client}. The folder of the
   client that ended last is taken up, and \texttt{notes.\codebreak{}txt} is open.
\end{enumerate}
\end{try}

\subsection{A daemon on another machine}
\label{sec:504}
Suppose you have a small server called \texttt{pinetree}, in a cupboard in Taupo,
and you keep your notes and a few projects on it. With a daemon there you
open your editor from any computer and find it as you left it, with the files
next to the programs that build them. Language servers, git and terminal
buffers all run on pinetree. Your computer draws the screen.

\begin{figure}[H]
\caption{Figure 18.2. A client here, a daemon there, ssh between them.}\label{org499}
\begin{Verbatim}[breaklines,breakautoindent=false,breakindent=1em,breakafter=/-.\,;:=,commandchars=\\\{\}]
 your computer                              pinetree
 +-------------------+                      +---------------------+
 | tuatara --client= |    ssh (OpenSSH)     | tuatara daemon      |
 | draws the screen  |<====================>| files, git, builds, |
 | sends your keys   |    keys, screen      | language servers    |
 +-------------------+                      +---------------------+
   clipboard, agents and dropped files cross only for a trusted server
\end{Verbatim}
\end{figure}

\subsubsection{Starting, using and stopping}
\label{sec:500}
Name the server before the daemon’s name, with a colon between them:

\begin{Verbatim}[breaklines,breakautoindent=false,breakindent=1em,breakafter=/-.\,;:=,commandchars=\\\{\}]
tuatara --client=pinetree:
tuatara --client=aroha@pinetree.example:work
tuatara --client=pinetree: --start todo.org
tuatara --stop=pinetree:
\end{Verbatim}

The first shows the default daemon on pinetree. The second shows the daemon
called \texttt{work} there, as the user \texttt{aroha}. The third starts a daemon on
pinetree first if none runs, and opens \texttt{todo.\codebreak{}org} there. The last stops it.

A name with a dot or an \texttt{@} in it is taken as a server. A plain name such as
\texttt{--client=\codebreak{}work} is a daemon on this computer. A server named only by an alias
from your \texttt{\textasciitilde{}/.\codebreak{}ssh/\codebreak{}config} file needs a colon after it, as in \texttt{pinetree:}.

The files you name are the server’s. A relative name such as \texttt{notes.\codebreak{}org} is
\texttt{notes.\codebreak{}org} in your home folder on pinetree, and \texttt{/etc/\codebreak{}hosts} is the file of
that name there.

\subsubsection{What has to be in place}
\label{sec:501}
The connection is made by the ssh program of your system, as if you had typed
\texttt{ssh pinetree}, with your keys and your \texttt{\textasciitilde{}/.\codebreak{}ssh/\codebreak{}config}. If ssh needs a
passphrase it asks in your terminal first. You need to have connected to the
server once before, and checked its host key, because a host whose key is not
known is refused. You also need \texttt{setpriv}, a small program from the util-linux
package, which ties the connection to the client so that it ends when the
client does, even if the client is killed.

On the server you need the same build of Tuatara. If it is older, or not
there, the client tells you and exits. If Tuatara is not on the server’s
\texttt{PATH}, name it in \texttt{M-x server-\codebreak{}list} with \texttt{T}, as shown below.

The server has to run OpenSSH. A daemon takes a client from another machine
only when OpenSSH’s own server carries the connection. It refuses a
connection through any other ssh server. The daemon must run under an ordinary
account, directly on the server, not as root and not in a container.

If the server removes your runtime folder when you log out, the daemon’s
socket goes with it. \texttt{--start} says so. On a server run by systemd, the
command \texttt{loginctl enable-\codebreak{}linger} keeps the folder.

\subsubsection{Keys limited with command=}
\label{sec:502}
An entry in the server’s \texttt{authorized\_\codebreak{}keys} file may be limited with a
\texttt{command=} setting. That does not stop the key from forwarding a connection to
the daemon’s socket. If you give such a key to someone, add \texttt{restrict} or
\texttt{no-\codebreak{}port-\codebreak{}forwarding} to the same entry.

\subsubsection{What a server is given}
\label{sec:503}
A server you have not trusted gets the keys you type, the size of your screen
and the names of the files you open. Nothing else of your computer reaches it.
Pictures are drawn on the server and sent to your window as small pictures of
the size shown, so a page of a large scan does not cross the network whole.
Whatever your \texttt{\textasciitilde{}/.\codebreak{}ssh/\codebreak{}config} says, ssh is told not to forward X11, Kerberos
tickets, or an agent, and not to run local or remote commands.

A client closes with \texttt{C-x C-c} and leaves the daemon on pinetree running. A
client that is killed, or whose terminal is closed, takes its ssh connection
with it.

\subsection{The trusted servers}
\label{sec:512}
Some things are useful across the connection: the clipboard, your keys for
signing, files dropped on the window. Each is off until you trust a server
with it. \texttt{M-x server-\codebreak{}list} shows the trusted servers in a buffer called
\texttt{*Servers*}, a line for each. A server not on the list is given nothing.

\texttt{M-x server-\codebreak{}trust} adds a machine, written as \texttt{--client} takes it. Note that
\texttt{pinetree.\codebreak{}example} and \texttt{aroha@\codebreak{}pinetree.\codebreak{}example} are two entries. A new entry
has nothing turned on. Each thing is then turned on by a key on its line,
after Tuatara asks and says what the server will be given. Turning a thing
off asks nothing.

\subsubsection{The columns}
\label{sec:505}
\begin{Verbatim}[breaklines,breakautoindent=false,breakindent=1em,breakafter=/-.\,;:=,commandchars=\\\{\}]
 M Server            Agent Keys used for  GPG Clipboard Dropped files to Trusted
   aroha@pinetree    yes   github.com     no  yes       Dropped          2026-11-04
   napier.example    no                   no  no        no               2026-11-06
\end{Verbatim}

\begin{xltabular}{\textwidth}{@{}l >{\raggedright\arraybackslash\hsize=1.000\hsize}X@{}}
\toprule
\textbf{Column} & \textbf{What it says}\\
\midrule
\endhead
M & A mark (\texttt{>} or \texttt{D}) put on the line.\\
Server & The machine, as \texttt{--client} names it.\\
Agent & Whether your ssh agent is forwarded there.\\
Keys used for & The machines the agent’s keys may be used for from there; \texttt{any} when none is named.\\
GPG & Whether gpg-agent is forwarded there.\\
Clipboard & Whether the clipboard is shared with it.\\
Dropped files to & The folder there that dropped files are sent to, or \texttt{no}.\\
Trusted & The day you trusted the machine.\\
Tuatara there & The program to run there, if not the one on its \texttt{PATH}.\\
\bottomrule
\end{xltabular}

\subsubsection{The keys of the list}
\label{sec:506}
Each key acts on the servers marked with \texttt{m}, or on the line the cursor is on
when none is marked.

\begin{xltabular}{\textwidth}{@{}l l >{\raggedright\arraybackslash\hsize=1.000\hsize}X@{}}
\toprule
\textbf{Modal} & \textbf{Non-modal} & \textbf{What it does}\\
\midrule
\endhead
none & \texttt{+} & trust another machine, with nothing turned on\\
none & \texttt{A} & forward the ssh agent, or stop\\
none & \texttt{H} & name the machines the agent’s keys may be used for from there\\
none & \texttt{G} & forward gpg-agent, or stop\\
none & \texttt{C} & share the clipboard, or stop\\
none & \texttt{F} & send dropped files to a folder there, or stop\\
none & \texttt{T} & name the tuatara to run there\\
none & \texttt{D} & forget the server, after asking\\
\bottomrule
\end{xltabular}

The list is kept in a file named \texttt{ssh-\codebreak{}servers} in Tuatara’s state folder. A
line it cannot read gives its server nothing.

\subsubsection{The ssh agent}
\label{sec:507}
An ssh agent holds your private keys and signs for you, so a program that
needs to log in somewhere does not have to ask for your passphrase. With \texttt{A} on
for a server, the programs there, git for example, can sign with your keys
while a client shows that daemon. Tuatara does not hand over the agent you
keep your keys in. The client starts an agent of its own for this connection
and puts in it only the keys your ssh configuration names for that server. Each
key is added so that every use of it brings up a question on your screen, and
nothing is signed until you say yes. If there is no program to ask the
question, the client says so and does not connect.

\texttt{H} limits where the keys may be used from the server. With \texttt{github.\codebreak{}com}
named, the server can use them for github.com and nothing else. Write the
names as they appear in your \texttt{known\_\codebreak{}hosts} file. This needs OpenSSH 8.9 or
later on both machines. If the limit cannot be applied, the client refuses to
connect. To forward the agent with no limit, name no machine.

The agent reaches only the programs of that client’s frames. Another client
of the same daemon sees none of it, and it ends with the connection.

\subsubsection{The gpg agent}
\label{sec:508}
GnuPG is a program for signing and encrypting, and gpg-agent holds its keys.
With \texttt{G} on, gpg on the server signs and decrypts with your keys, which stay on
your computer. The passphrase is asked here. The forwarded socket cannot export
or change your keys, so nothing on the server can take them. While gpg-agent
here keeps the passphrase, the server can use that key without asking again.

Two things need setting on the server. gpg there must not start an agent of
its own, so put the line \texttt{no-\codebreak{}autostart} in \texttt{gpg.\codebreak{}conf} in GnuPG’s folder there
(\texttt{gpgconf --list-\codebreak{}dirs homedir} says where it is, usually \texttt{\textasciitilde{}/.\codebreak{}gnupg}). And the
ssh server should remove a leftover socket before forwarding a new one, which
means adding \texttt{StreamLocalBindU\codebreak{}nlink yes} to \texttt{/etc/\codebreak{}ssh/\codebreak{}sshd\_\codebreak{}config} as root.
When the client cannot tell whether that is set, it says what to add.

To encrypt a file on the server to someone else, gpg there needs their public
key. In the list of your keys (\texttt{M-x gpg-\codebreak{}list-\codebreak{}keys}), mark the keys with \texttt{m} and
run \texttt{M-x server-\codebreak{}copy-\codebreak{}keys}. It asks for the trusted server and sends the public
keys over ssh in the background. Nothing secret is sent.

\subsubsection{The clipboard and dropped files}
\label{sec:509}
With \texttt{C} on, the clipboard goes both ways while the client’s window has the
keyboard. What you copy in the daemon reaches the clipboard on your computer,
and the daemon may read what you copied here. Text that comes from the server
loses its control characters, so that pasting it into a terminal cannot carry
a hidden escape sequence. Without \texttt{C}, nothing of the clipboard crosses.

With \texttt{F} on, a file you drop onto the window showing the server’s daemon is
sent into a folder there, \texttt{Dropped} in your home folder unless you give another
name. A name beginning with \texttt{/} is taken from the root. The file is written
for you alone under a name nothing else has, and nothing in it is run. Only
files are sent, at most 100 MB each.

\subsubsection{A language server for a server}
\label{sec:510}
A language server runs where the daemon runs. If pinetree has none installed
and you cannot install one, \texttt{M-x server-\codebreak{}bring-\codebreak{}language-\codebreak{}server} copies one from
your computer, in the background, to a folder of the editor’s own on the
server. Nothing is fetched from the internet. It must be a program for the
server’s kind of machine that needs nothing but the usual system libraries.
\texttt{M-x server-\codebreak{}remove-\codebreak{}language-\codebreak{}server} takes it away again.

\subsubsection{Commands for the trusted servers}
\label{sec:511}
\begin{xltabular}{\textwidth}{@{}l >{\raggedright\arraybackslash\hsize=0.737\hsize}X l >{\raggedright\arraybackslash\hsize=1.263\hsize}X@{}}
\toprule
\textbf{Modal} & \textbf{Non-modal} & \textbf{Command} & \textbf{What it does}\\
\midrule
\endhead
none & \texttt{M-x server-\codebreak{}list} & \texttt{server-\codebreak{}list} & List the trusted servers, the machines whose daemons a client here (tuatara --client=HOST) may give more than the keys typed and the size of the screen, each with what it is given: the ssh agent, with the machines its keys may be used for from there; gpg-agent; the clipboard; files dropped on the window, sent to a folder there.\\
none & \texttt{+}, \texttt{C-c +} & \texttt{server-\codebreak{}trust} & Trust a machine asked for in the minibuffer, as tuatara --client= names it (build.example, me@build.example, or an alias of \textasciitilde{}/.ssh/config): it goes on the list of trusted servers (M-x server-list) with nothing given yet, each thing then turned on there by itself.\\
none & \texttt{A}, \texttt{C-c A} & \texttt{server-\codebreak{}toggle-\codebreak{}agent} & Forward your ssh agent to the server on the line, or those marked, or no longer: on, the programs there can sign with your keys while a client shows its daemon, each use asked about here, and only for the machines named by H where OpenSSH can limit them.\\
none & \texttt{H}, \texttt{C-c H} & \texttt{server-\codebreak{}agent-\codebreak{}destinations} & Name, in the minibuffer, the machines the keys of your ssh agent may be used for from the server on the line, or those marked, commas between them (github.com, me@git.example); none for any.\\
none & \texttt{G}, \texttt{C-c G} & \texttt{server-\codebreak{}toggle-\codebreak{}gpg} & Forward gpg-agent to the server on the line, or those marked, or no longer: on, gpg there signs and decrypts with your keys, which stay here, and cannot export or change them.\\
none & \texttt{C}, \texttt{C-c C} & \texttt{server-\codebreak{}toggle-\codebreak{}clipboard} & Share the clipboard with the server on the line, or those marked, or no longer: on, what is copied there reaches the clipboard here, as text with no control characters, and the daemon there may read what is copied here.\\
none & \texttt{F}, \texttt{C-c F} & \texttt{server-\codebreak{}toggle-\codebreak{}files} & Send the files dropped on a window showing the daemon of the server on the line, or those marked, to a folder there, asked for (from the home folder there unless it begins with /), or no longer: on, a file dropped is copied there, its size checked first, and nothing of it run.\\
none & \texttt{T}, \texttt{C-c T} & \texttt{server-\codebreak{}tuatara-\codebreak{}program} & Name, in the minibuffer, the tuatara to run on the server on the line, or those marked, when it is not the one on the PATH of the commands ssh runs there: /opt/tuatara/bin/tuatara, or a name found on that PATH; empty for the one on that PATH.\\
none & \texttt{M-x server-\codebreak{}bring-\codebreak{}language-\codebreak{}server} & \texttt{server-\codebreak{}bring-\codebreak{}language-\codebreak{}server} & Copy a language server installed on this machine to the server on the line of the list of trusted servers, or those marked, chosen in the minibuffer among those the editor knows that are here: nothing is fetched from anywhere.\\
none & \texttt{M-x server-\codebreak{}remove-\codebreak{}language-\codebreak{}server} & \texttt{server-\codebreak{}remove-\codebreak{}language-\codebreak{}server} & Take away a language server M-x server-bring-language-server brought to the server on the line of the list of trusted servers, or those marked, chosen in the minibuffer among those it lists there.\\
none & \texttt{M-x server-\codebreak{}copy-\codebreak{}keys} & \texttt{server-\codebreak{}copy-\codebreak{}keys} & Copy the public keys marked in the list of GnuPG keys (M-x gpg-list-keys), or the one at the cursor, into gpg’s keyring on a trusted server asked for, so that M-x gpg-encrypt-file there can encrypt to them: exported here and imported there over ssh, in the background.\\
\bottomrule
\end{xltabular}

\begin{try}
\textbf{Try it, with a daemon on pinetree}
\begin{enumerate}
\item On pinetree, over ssh, type \texttt{tuatara --daemon}. Then leave the ssh
   session. Check that you can type \texttt{ssh pinetree} without a surprise.
\item On your own computer, type \texttt{tuatara --client=\codebreak{}pinetree: --start todo.\codebreak{}org}.
   The file opens, in the home folder on pinetree. The mode line shows \texttt{D}.
\item Press \texttt{C-x C-b}. The client is listed with the folder it started in.
\item Press \texttt{M-x}, type \texttt{server-\codebreak{}trust}, press \texttt{RET}, type \texttt{pinetree} and press
   \texttt{RET}. Answer the question. Press \texttt{M-x} and type \texttt{server-\codebreak{}list}.
\item Put the cursor on the \texttt{pinetree} line and press \texttt{C}. Answer \texttt{y}. The
   Clipboard column now says \texttt{yes}. Copy a line in the daemon with \texttt{M-w} and
   paste it into a program of your own computer.
\item Press \texttt{C-x C-c}. The daemon on pinetree is still running. Type \texttt{tuatara
   --stop=\codebreak{}pinetree:} to end it.
\end{enumerate}
\end{try}

\begin{try}
\textbf{Try it, with a key that may be used for one place only}
\begin{enumerate}
\item Open \texttt{M-x server-\codebreak{}list}, put the cursor on a trusted server and press \texttt{H}.
\item Type \texttt{github.\codebreak{}com} and press \texttt{RET}. The Keys used for column shows it.
\item Press \texttt{A} and answer \texttt{y}. The next client of that server forwards your
   agent, each use confirmed on your screen, for github.com alone.
\end{enumerate}
\end{try}

\section{Packages}
\label{sec:544}
Adding to Tuatara what not everyone needs, and keeping it up to date.

Tuatara does a lot by itself, but some things are wanted by only a few
people. One person reads Outlook mailboxes, another keeps a calendar in step
with a phone, another wants a different set of colours. Putting all of that in
the editor would make it larger and slower for everyone. So these things are
kept apart, as packages, and you add the ones you want. This chapter explains
what a package is, how you name one in your init file, how Tuatara fetches and
checks it, and how you update it, switch it off or get rid of it. The end of
the chapter describes the packages the project publishes.

The init file is the file where you tell Tuatara how you like things. Chapter
20 covers it. Here you only add a line or two to it, which the steps show.
Every key in this chapter is pressed in NORMAL state unless it starts with
\texttt{C-x} or \texttt{C-c}, which work in both states.

\subsection{What a package is}
\label{sec:514}
A package is a folder of Lisp files that adds something to Tuatara: a
command, a way of opening a kind of file, a set of colours. Lisp is the
language Tuatara is written in. You do not need to read it to use a package,
but you can, and Tuatara offers to show it to you before it runs anything.

A package lives in one of two places. Most live in a git repository on the
internet, which is a folder of files with a history of every change made to
it. Tuatara fetches the version you ask for into a folder of its own. The
other kind is a folder on your own disk, which suits something you wrote
yourself. Tuatara calls the place a package comes from its source.

A package does three things for you that a loose file of code would not.
Tuatara compiles it, which turns the Lisp into a form the computer runs fast,
and it does this in the background while you go on typing. It loads the
package at every start, in a few thousandths of a second, from the compiled
form. And it checks who signed the version it fetched, so that a changed
package does not slip in unnoticed.

\subsection{Naming a package in your init file}
\label{sec:517}
You name a package with a line in your init file. The line says what the
package is called and where it comes from:

\begin{Verbatim}[breaklines,breakautoindent=false,breakindent=1em,breakafter=/-.\,;:=,commandchars=\\\{\}]
\color[HTML]{A85F00}(load-package\normalcolor{} \color[HTML]{0A7AB3}"themes"
              \color[HTML]{946C00}:\normalcolor{}source \color[HTML]{0A7AB3}"https://github.com/Tuatara-Editor/tuatara-themes"\color[HTML]{A85F00})
\end{Verbatim}

\texttt{load-\codebreak{}package} is the whole interface. The name comes first, in quotes, and
the rest are keywords, each followed by its value. Only the keywords you need
are written. Tuatara reads the init file when it starts, fetches the package
the first time, and asks you about it (see “The first load” below).

If you use a source for several packages, give it a short name once with
\texttt{package-\codebreak{}source}, and use that name after \texttt{:source}:

\begin{Verbatim}[breaklines,breakautoindent=false,breakindent=1em,breakafter=/-.\,;:=,commandchars=\\\{\}]
\color[HTML]{A85F00}(package-source\normalcolor{} \color[HTML]{0A7AB3}"mine"\normalcolor{} \color[HTML]{0A7AB3}"\color[HTML]{222222}~/Projects/\color[HTML]{0A7AB3}my-packages/"\color[HTML]{A85F00})
\color[HTML]{A85F00}(load-package\normalcolor{} \color[HTML]{0A7AB3}"hello"\normalcolor{} \color[HTML]{946C00}:\normalcolor{}source \color[HTML]{0A7AB3}"mine"\color[HTML]{A85F00})
\end{Verbatim}

An address that begins with \texttt{https://} or \texttt{ssh}, or a folder on your disk
(\texttt{\textasciitilde{}/} means your home folder), is accepted. An address that begins with
\texttt{http://} or \texttt{git://} is refused, because anyone on the way could read it and
change what is sent. Section “Sources that are not encrypted” says how to
allow one on a network you trust.

\subsubsection{The keywords}
\label{sec:515}
\begin{longtable}{@{}l l@{}}
\caption{Table 19.1. What can follow the name in a \texttt{load-\codebreak{}package} line.}\\
\toprule
\textbf{Keyword} & \textbf{What it says}\\
\midrule
\endfirsthead
\caption[]{Table 19.1. What can follow the name in a \texttt{load-\codebreak{}package} line. (continued)}\\
\toprule
\textbf{Keyword} & \textbf{What it says}\\
\midrule
\endhead
\texttt{:source} & Where the package comes from: the name given with\\
 & \texttt{package-\codebreak{}source}, an address, or a folder.\\
\texttt{:version} & Which version to fetch: a tag, a branch or a commit.\\
 & Left out, or \texttt{"release"}, it is the newest tag. With\\
 & \texttt{"latest"} it is the newest commit, and it is also the\\
 & one taken when the repository has no tag yet. A whole\\
 & 40-character commit keeps the package at that commit.\\
\texttt{:folder}, \texttt{:file} & The package’s folder inside the source, and its first\\
 & file, when they are not the package’s name. The folder\\
 & is \texttt{NAME/}, else the top of the source. The file is\\
 & \texttt{NAME.\codebreak{}lisp}, else \texttt{load.\codebreak{}lisp}.\\
\texttt{:signers} & The fingerprints of the keys that must have signed it.\\
\texttt{:plain t} & Allows an \texttt{http://} or \texttt{git://} address.\\
\texttt{:disabled t} & Leaves the whole line out, as if it were not there.\\
\texttt{:if FORM} & Leaves the line out when \texttt{FORM} gives \texttt{nil}. It is how\\
 & one init file serves two machines.\\
\texttt{:defer t} & Loads the package once Tuatara has started, not while\\
 & the init file is read.\\
\texttt{:after "other"} & Loads it only after the package or group \texttt{other} has.\\
\texttt{:commands} & Loads it when one of these commands is first run.\\
\texttt{:modes} & Loads it when a buffer first takes one of these modes.\\
\texttt{:init}, \texttt{:config} & Forms run before the package loads, and after it. Its\\
 & settings and keys go after \texttt{:config}.\\
\bottomrule
\end{longtable}

A form is a piece of Lisp in brackets, such as \texttt{(setf *garamond-\codebreak{}mode-\codebreak{}line*
nil)}. \texttt{:config} takes any number of them, one after another. This is the
place for a package’s settings, since they exist only once the package has
loaded:

\begin{Verbatim}[breaklines,breakautoindent=false,breakindent=1em,breakafter=/-.\,;:=,commandchars=\\\{\}]
\color[HTML]{A85F00}(load-package\normalcolor{} \color[HTML]{0A7AB3}"garamond"
              \color[HTML]{946C00}:\normalcolor{}source \color[HTML]{0A7AB3}"https://github.com/Tuatara-Editor/tuatara-garamond"
              \color[HTML]{946C00}:\normalcolor{}config
              \color[HTML]{A85F00}(setf\normalcolor{} \color[HTML]{A85F00}*garamond-mode-line*\normalcolor{} \color[HTML]{0A7AB3}nil\color[HTML]{A85F00}))
\end{Verbatim}

\texttt{:defer}, \texttt{:after}, \texttt{:commands} and \texttt{:modes} are for starting quickly. A
package you use once a week does not have to be ready at the first frame. With
\texttt{:commands} or \texttt{:modes} it waits until you ask for it. Even then you do not
have to remember to load it: running one of its commands loads it first.

\texttt{load-\codebreak{}package} with no \texttt{:source} loads nothing. It groups settings and keys
under a name, the way the starter init file does for \texttt{"keys"} and \texttt{"looks"}.
Tuatara also corrects a misspelt name: a name a letter or two away from a
package it knows is an error that names the one it thinks you meant.

\subsubsection{The first load}
\label{sec:516}
The first time Tuatara meets a package it fetches it, in the background, and
asks before it compiles and loads it:

\begin{Verbatim}[breaklines,breakautoindent=false,breakindent=1em,breakafter=/-.\,;:=,commandchars=\\\{\}]
Compile and load the package themes from https://github.com/...? (y, n, or v to view its source)
\end{Verbatim}

You answer with one key, with no \texttt{RET}. \texttt{y} goes on. \texttt{n} leaves the package
unloaded, and Tuatara asks again at the next start. \texttt{v} lists the package’s
folder so that you can read its source first, opened safely so that nothing in
it runs. \texttt{C-g} means “not now” and also asks again next time. Any other key
puts \texttt{Answer y, n or v.} before the question and asks again. Keys you typed
before the question appeared are thrown away, so a key meant for your text can
never answer it.

After \texttt{y} the compiling happens in a second copy of Tuatara that has no window
and runs apart from your editor. You can keep working. \texttt{C-x j} lists the job
and \texttt{C-x J} stops it. When it is done, the package loads in one step, between
two of your commands, so no command ever runs with a package half loaded.

If you trust your sources and want no question, set
\texttt{*package-\codebreak{}ask-\codebreak{}before-\codebreak{}loading*} to \texttt{nil}.

\subsection{The package manager}
\label{sec:520}
\texttt{M-x package-\codebreak{}manager} shows the packages your init file loads from a source,
one line each, in a buffer called \texttt{*Packages*}. It works as the buffer list
of chapter 7 does: you move over the lines, mark some, and act on the marked
ones.

\begin{figure}[H]
\caption{Figure 19.1. The package manager, with its README beside it.}\label{org518}
\begin{Verbatim}[breaklines,breakautoindent=false,breakindent=1em,breakafter=/-.\,;:=,commandchars=\\\{\}]
 M Name        Version  Commit   State     Signed by      Source
 - ----        -------  ------   -----     ---------      ------
   themes      v0.3     4c1e9a0  loaded    Aroha Ngata    https://github.com/T…
 > kitty       v0.2     b07d332  loaded    Aroha Ngata    https://github.com/T…
   hello       folder            loaded    unsigned       ~/Projects/my-pack…
   outlook     v1.1     91ac7e5  unloaded  Aroha Ngata    https://github.com/T…

   4 packages, 3 loaded
\end{Verbatim}
\end{figure}

The columns are the mark, the name, the version (\texttt{folder} for a folder of
your disk), the first letters of the commit, what became of the package this
session, who signs its source’s versions, and the source. The state is one of
\texttt{loaded}, \texttt{fetching}, \texttt{failed}, \texttt{refused}, \texttt{declined}, \texttt{kept}, \texttt{unloaded},
\texttt{removed} and \texttt{not fetched}. A package that is not loaded is drawn in the
faint colour, as in the figure. \texttt{declined} means you said no at the question.
\texttt{refused} means its code uses a name this Tuatara does not have, because the
package is newer than your Tuatara or a newer Tuatara renamed something. A
refused package is never run, and the reason shows beside the list where the
README would be. The README of the package under the cursor shows beside the
list, or below it when the window is narrow. \texttt{P} turns it off and on.

The keys are those of the lists you already know, with a few of its own.

\begin{longtable}{@{}l l l@{}}
\caption{Table 19.2. Keys in the package manager.}\\
\toprule
\textbf{Modal} & \textbf{Non-modal} & \textbf{What it does}\\
\midrule
\endfirsthead
\caption[]{Table 19.2. Keys in the package manager. (continued)}\\
\toprule
\textbf{Modal} & \textbf{Non-modal} & \textbf{What it does}\\
\midrule
\endhead
none & \texttt{RET}, \texttt{l} & Lists the package’s folder, to read its source.\\
none & \texttt{o} & The same, in another window.\\
none & \texttt{n}, \texttt{p} & Go to the next package, and the one before.\\
none & \texttt{m}, \texttt{u}, \texttt{U}, \texttt{t} & Mark, unmark, unmark all, swap the marks.\\
none & \texttt{f} & Fetches the marked packages again (see “Updating”).\\
none & \texttt{L} & Loads the marked ones.\\
none & \texttt{K} & Unloads the marked ones from the running editor.\\
none & \texttt{R} & Removes them: unloads them and keeps them from loading.\\
none & \texttt{d}, then \texttt{x} & Marks for deletion, then deletes what is kept of them.\\
none & \texttt{D} & Deletes what is kept of the marked ones.\\
none & \texttt{h}, \texttt{q} & Closes the list.\\
\bottomrule
\end{longtable}

With nothing marked, each of these acts on the package on the cursor’s line.

\texttt{M-x package-\codebreak{}scan} reads your init file again for \texttt{load-\codebreak{}package} lines that
have a \texttt{:source}, loads each package that is not loaded, as at a start, and
then shows the package manager. It is what you run after adding a line to the
init file, when you do not want to restart Tuatara.

\begin{try}
\textbf{Try it, adding the themes package}
\begin{enumerate}
\item Press \texttt{C-x C-f}, type \texttt{\textasciitilde{}/.\codebreak{}config/\codebreak{}tuatara/\codebreak{}init.\codebreak{}lisp} and press \texttt{RET}. This
   is your init file.
\item Press \texttt{g} then \texttt{e} to go to the end. Press \texttt{o} to open a new line and
   type:
   \texttt{(load-\codebreak{}package "themes" :source "https://\codebreak{}github.\codebreak{}com/\codebreak{}Tuatara-\codebreak{}Editor/\codebreak{}tuatara-\codebreak{}themes")}
\item Press \texttt{C-g} and then \texttt{C-x C-s} to save.
\item Run \texttt{M-x package-\codebreak{}scan}. \texttt{*Packages*} opens with a line for \texttt{themes}. In a
   moment the question appears. Press \texttt{v} to look at the source, then \texttt{q} to
   come back, and answer \texttt{y}. If the line says \texttt{failed} and the echo row says
   the package is refused because of a signer, your keyring lacks the key
   that signed it. Import it as described under “Signatures and trust”, run
   \texttt{M-x package-\codebreak{}update} and try again.
\item The echo row says the package is compiled and loaded. The line in the list
   now says \texttt{loaded}.
\item Run \texttt{M-x load-\codebreak{}theme}. The list now includes fourteen more themes. Choose
   \texttt{nord} and press \texttt{RET}. The colours change at once.
\end{enumerate}
\end{try}

\subsubsection{Commands of the package manager}
\label{sec:519}
\begin{xltabular}{\textwidth}{@{}l l l >{\raggedright\arraybackslash\hsize=1.000\hsize}X@{}}
\toprule
\textbf{Modal} & \textbf{Non-modal} & \textbf{Command} & \textbf{What it does}\\
\midrule
\endhead
none & \texttt{M-x package-\codebreak{}manager} & \texttt{package-\codebreak{}manager} & Manage the packages the init file loads from a source (load-package with :source) in a list, as the buffer list (C-x C-b) manages buffers.\\
none & \texttt{M-x package-\codebreak{}scan} & \texttt{package-\codebreak{}scan} & Look in the init file for the packages it loads from a source (each load-package with :source), and load each that is not loaded or on its way, as at the start: fetched in the background, asked about before it is compiled, loaded.\\
\bottomrule
\end{xltabular}

\subsection{Updating}
\label{sec:522}
A package fetched from a repository keeps the version you first took. When its
author publishes a newer one, \texttt{M-x package-\codebreak{}update} looks again. It asks each
repository for the version the \texttt{load-\codebreak{}package} line wants (the newest tag, the
newest commit, or the tag, branch or commit you named), fetches what is new in
the background and, when it is a different commit from the one you have, asks
the question of the first load again. Packages that are already up to date
are said to be so. Folders of your own disk are left alone, since you change
those yourself.

\texttt{M-x package-\codebreak{}update-\codebreak{}marked}, or \texttt{f} in the package manager, does the same for
only the packages you marked, or for the one on the line.

A package can say, in a comment at the top of its first file, the oldest
Tuatara it works with: \texttt{;;;; Needs: Tuatara 0.91.0}. Tuatara reads that line
without running anything. If the newest version needs a newer Tuatara than
yours, Tuatara looks back through the earlier versions for the newest that
fits, says which one it took, and loads that. If none fits, the version you
have stays. A package that says nothing loads in any Tuatara.

Tuatara keeps the version it took and the one before it, each with its
compiled files. When a new version fails to load, because it cannot be
compiled or uses a name your Tuatara lacks, the one before it is loaded in its
place, and a message tells you so once.

\begin{try}
\textbf{Try it, updating everything}
\begin{enumerate}
\item Run \texttt{M-x package-\codebreak{}update}. The echo row says it is fetching.
\item Press \texttt{C-x j} to see the job. When it finishes, packages with nothing new
   are said to be up to date.
\item If a package has a new commit, the question of the first load appears for
   it. Press \texttt{v} to read what changed, \texttt{q} to come back, then \texttt{y}.
\item Open \texttt{M-x package-\codebreak{}manager}. The Version and Commit columns show what you
   have now.
\item To update only some: press \texttt{m} on two packages, then \texttt{f}. Only those two
   are fetched again.
\end{enumerate}
\end{try}

\subsubsection{Commands for updating}
\label{sec:521}
\begin{xltabular}{\textwidth}{@{}l l l >{\raggedright\arraybackslash\hsize=1.000\hsize}X@{}}
\toprule
\textbf{Modal} & \textbf{Non-modal} & \textbf{Command} & \textbf{What it does}\\
\midrule
\endhead
none & \texttt{M-x package-\codebreak{}update} & \texttt{package-\codebreak{}update} & Fetch again each package load-package fetched (not the folders of your disk), in the background.\\
none & \texttt{f}, \texttt{C-c f} & \texttt{package-\codebreak{}update-\codebreak{}marked} & Fetch again the marked packages (>), or with none marked the one on the current line, as M-x package-update does them all.\\
\bottomrule
\end{xltabular}

\subsection{Signatures and trust}
\label{sec:526}
Code you load runs with all your rights. That is true of any editor’s
extensions. So Tuatara asks who signed each version, as one way of noticing
when something has changed hands.

The author signs a tag or a commit with a key, using GPG or SSH. Tuatara asks
git who signed the one it fetched. Git does the checking, with \texttt{gpg} and your
keyring for a GPG signature and with \texttt{ssh-\codebreak{}keygen} for an SSH one, and nothing
in the repository is run to do it. The question at the first load says who
signed the version, and your \texttt{y} makes that signer the one Tuatara remembers
for the source. The memory is kept in a file called \texttt{package-\codebreak{}trust} beside the
packages, which you can read. This is called trust on first use: you decide
once, and Tuatara holds the source to that decision.

From then on:

\begin{itemize}
\item A version signed by the same key is taken.
\item A version signed by another key, or unsigned when the versions before it
  were signed, or whose signature does not check out, is refused. The version
  you have keeps running.
\item A source whose first version was unsigned goes on having unsigned versions
  taken, since nobody has to sign. The first signed version makes its signer
  the one remembered.
\end{itemize}

A refusal says what changed:

\begin{Verbatim}[breaklines,breakautoindent=false,breakindent=1em,breakafter=/-.\,;:=,commandchars=\\\{\}]
The package tally from https://example.org/tally.git is refused: unsigned,
and the versions taken before were signed by Ann Author (key SHA256:...).
The version loaded now is kept; M-x package-trust-again trusts its source
anew, as the first time
\end{Verbatim}

Sometimes the change is honest: the author has a new key, or moved the
project. Then you decide, as at the first time. \texttt{M-x package-\codebreak{}trust-\codebreak{}again}
asks which package, forgets who signed its source, and fetches it again. The
question names the new signer, and your \texttt{y} makes it the one remembered. Think
before you press it. If you do not know why the key changed, ask the author.

A version signed with a GPG key your keyring lacks is refused, and Tuatara
remembers the key’s fingerprint. Import the key with \texttt{M-x gpg-\codebreak{}import-\codebreak{}keys} and
run \texttt{M-x package-\codebreak{}update}. Versions signed by that key, or by the key it
belongs to, are then taken.

\begin{xltabular}{\textwidth}{@{}l >{\raggedright\arraybackslash\hsize=0.542\hsize}X l >{\raggedright\arraybackslash\hsize=1.458\hsize}X@{}}
\toprule
\textbf{Modal} & \textbf{Non-modal} & \textbf{Command} & \textbf{What it does}\\
\midrule
\endhead
none & \texttt{M-x package-\codebreak{}trust-\codebreak{}again} & \texttt{package-\codebreak{}trust-\codebreak{}again} & Trust the source of a package asked for anew, as the first time: who signed the versions taken from it is forgotten, and the package is fetched again (M-x package-update); the version found is asked about with who signed it (\textbf{PACKAGE-ASK-BEFORE-LOADING}), and once it is let load, that signer is the one remembered.\\
none & \texttt{M-x gpg-\codebreak{}import-\codebreak{}keys} & \texttt{gpg-\codebreak{}import-\codebreak{}keys} & Import the keys in a file asked for (armoured or not) into your keyring; say what was imported.\\
\bottomrule
\end{xltabular}

\subsubsection{Naming the signers yourself}
\label{sec:523}
If you know whose keys you want, say so, and only versions signed by them are
taken, from the very first:

\begin{Verbatim}[breaklines,breakautoindent=false,breakindent=1em,breakafter=/-.\,;:=,commandchars=\\\{\}]
\color[HTML]{A85F00}(package-source\normalcolor{} \color[HTML]{0A7AB3}"mine"\normalcolor{} \color[HTML]{0A7AB3}"https://example.org/my-packages.git"
                \color[HTML]{946C00}:\normalcolor{}signers \color[HTML]{222222}'\color[HTML]{A85F00}(\color[HTML]{0A7AB3}"SHA256:hdVZ5mzqXcO1EQf8\discretionary{}{}{}N4u9MJvZb7TfsHkL\discretionary{}{}{}qR2wPaXyC0s"\color[HTML]{A85F00}))
\end{Verbatim}

An SSH key’s fingerprint is what \texttt{ssh-\codebreak{}keygen -lf KEY.\codebreak{}pub} prints. A GPG key’s
is the 40 hexadecimal digits that \texttt{gpg --fingerprint} shows. A key’s short ID
is refused, and so is a signing subkey’s fingerprint: give the key’s own.
\texttt{load-\codebreak{}package} takes \texttt{:signers} as well. Names in the init file always decide
for their source, and \texttt{package-\codebreak{}trust-\codebreak{}again} does not change that.

\subsubsection{Sources that are not encrypted}
\label{sec:524}
A server inside your organisation may speak only \texttt{http://}. You can allow it
for one package at a time:

\begin{Verbatim}[breaklines,breakautoindent=false,breakindent=1em,breakafter=/-.\,;:=,commandchars=\\\{\}]
\color[HTML]{A85F00}(load-package\normalcolor{} \color[HTML]{0A7AB3}"tally"\normalcolor{} \color[HTML]{946C00}:\normalcolor{}source \color[HTML]{0A7AB3}"http://git.example.org/tally.git"\normalcolor{} \color[HTML]{946C00}:\normalcolor{}plain \color[HTML]{A85F00}t)
\end{Verbatim}

Whoever controls the network between you and that server can hand you an older
version of the package than the newest, one with a bug that has since been
fixed. A signature does not stop that, because the old version is genuinely
signed. Use \texttt{:plain t} only on a network you trust. The \texttt{Signed by} column of
the package manager says \texttt{unencrypted} for these sources.

\subsubsection{Settings}
\label{sec:525}
\begin{xltabular}{\textwidth}{@{}l l >{\raggedright\arraybackslash\hsize=1.000\hsize}X@{}}
\toprule
\textbf{Setting} & \textbf{Standard} & \textbf{What it changes}\\
\midrule
\endhead
\texttt{*package-\codebreak{}check-\codebreak{}signatures*} & \texttt{t} & True to check who signed each version of a package, and take only those that M-x package-trust-again and the signers of a source (:signers) allow.\\
\bottomrule
\end{xltabular}

\texttt{*package-\codebreak{}check-\codebreak{}signatures*} is \texttt{t}, and you should leave it so. Setting it to
\texttt{nil} makes every version of every source be taken, signed by anyone or by no
one, and nothing is remembered. That is only for a machine whose sources are
trusted some other way. \texttt{M-x check-\codebreak{}health} says when signatures are not
checked.

\subsection{Unloading and removing}
\label{sec:530}
A package can be taken out of Tuatara in three steps, each going further than
the one before.

\texttt{M-x package-\codebreak{}unload} (\texttt{K} in the package manager) takes the package out of
the running editor: its commands, keys, modes, settings and hooks, and any
work it has in the background, which is stopped. A buffer in one of its modes
is put in the mode its file calls for, so nothing you typed is lost. The
package’s code stays on the disk and its \texttt{load-\codebreak{}package} line stays in your
init file, so it loads again at the next start. Use it for a package that is
misbehaving, or to try Tuatara without it for the afternoon.

\texttt{M-x package-\codebreak{}remove} (\texttt{R}) unloads it and keeps it from loading at the next
start as well, without changing your init file. Tuatara remembers in its state
folder that you removed the package, passes the \texttt{load-\codebreak{}package} line over at
each start and says so once in \texttt{*Messages*}. Its code stays on the disk, so
bringing it back is quick.

\texttt{M-x package-\codebreak{}delete} (\texttt{D}, or \texttt{d} then \texttt{x}) unloads and removes it and then
deletes what Tuatara keeps of it: the versions fetched and their compiled
files. For a package from a folder of your own disk only the compiled files go.
Your folder is never touched. It asks first, with the size of what goes:

\begin{Verbatim}[breaklines,breakautoindent=false,breakindent=1em,breakafter=/-.\,;:=,commandchars=\\\{\}]
Unload, remove and delete the package hello, 1.2 MB? (y or n, then RET)
\end{Verbatim}

\texttt{M-x package-\codebreak{}load} (\texttt{L}) goes the other way. It loads a package again by its
\texttt{load-\codebreak{}package} line, as at the start. That works for one you unloaded, one you
removed (it is then no longer removed), and one you deleted (it is fetched
again). In the package manager each of these four acts on the packages you
marked, or on the one on the line. Elsewhere the commands ask which package.

To stop a package for good, also take its line out of your init file, or put
\texttt{;;} in front of it, or add \texttt{:disabled t} to the line. Unloading alone does
not do that. A deleted package that is still named in the init file comes back
at the next start.

\begin{try}
\textbf{Try it, removing a package you no longer use}
\begin{enumerate}
\item Run \texttt{M-x package-\codebreak{}manager} and press \texttt{j} (non-modal: \texttt{C-n}) until the cursor
   is on \texttt{kitty}.
\item Press \texttt{R}. The package is unloaded and its state reads \texttt{removed}.
\item Restart Tuatara and open the package manager again. \texttt{kitty} is still there
   but is not loaded and was not asked about.
\item Press \texttt{L} on its line. It loads again.
\item To free the disk space as well, press \texttt{D}. The question names the size.
   Answer \texttt{y} and then \texttt{RET}. Remove the \texttt{load-\codebreak{}package} line from your init
   file when you are sure you are done with it.
\end{enumerate}
\end{try}

\subsubsection{Commands for unloading and removing}
\label{sec:527}
\begin{xltabular}{\textwidth}{@{}l l l >{\raggedright\arraybackslash\hsize=1.000\hsize}X@{}}
\toprule
\textbf{Modal} & \textbf{Non-modal} & \textbf{Command} & \textbf{What it does}\\
\midrule
\endhead
none & \texttt{K}, \texttt{C-c K} & \texttt{package-\codebreak{}unload} & Take a package out of the running editor: its commands and keys, with their labels, its modes, the places it gave Tuatara’s commands, its ways of showing text, its functions on Tuatara’s lists of hooks and of work to do now and then, its settings and its work in the background, which is stopped.\\
none & \texttt{R}, \texttt{C-c R} & \texttt{package-\codebreak{}remove} & Unload a package (M-x package-unload) and take it out of the packages loaded: it is not loaded from now on, nor at the next start, though the init file names it; the init file is left as it is, and Tuatara says in \textbf{Messages} once a session that it passes the package over.\\
none & \texttt{M-x package-\codebreak{}delete} & \texttt{package-\codebreak{}delete} & Unload a package (M-x package-unload), remove it from the packages loaded (M-x package-remove), and delete what Tuatara keeps of it: the versions of it fetched and their compiled files, or for a folder of your disk its compiled files alone, the folder of yours left as it is.\\
none & \texttt{L}, \texttt{C-c L} & \texttt{package-\codebreak{}load} & Load a package the init file loads from a source that is not loaded now, by its load-package there, as at the start: one unloaded (M-x package-unload), removed (M-x package-remove), whose loading failed, or one deleted, fetched again.\\
\bottomrule
\end{xltabular}

\subsubsection{Clearing out sources no longer used}
\label{sec:528}
Tuatara keeps what it fetches until you say otherwise. When you rename a
source or stop using one, the old copies stay in the packages folder.
\texttt{M-x package-\codebreak{}remove-\codebreak{}unused} finds them: every source that no \texttt{package-\codebreak{}source}
names and no \texttt{load-\codebreak{}package} fetches from, and every package of a source you
still use that no \texttt{load-\codebreak{}package} names. A \texttt{load-\codebreak{}package} turned off with
\texttt{:disabled t}, or left out by \texttt{:if}, still counts as used.

It lists what would go in \texttt{*Help*}, each with its size, and asks once.
Answer \texttt{n} and nothing is touched. Answer \texttt{y} and the folders are deleted in
the background. Each is looked at again just before it goes, so a source you
name in the meantime stays. Only Tuatara’s own packages folder is searched. A
folder of your own disk named as a source is never touched. If your init file
stopped at an error when Tuatara started, the command does nothing until the
file reads through, since everything after the error would look unused.

\begin{xltabular}{\textwidth}{@{}l >{\raggedright\arraybackslash\hsize=0.565\hsize}X l >{\raggedright\arraybackslash\hsize=1.435\hsize}X@{}}
\toprule
\textbf{Modal} & \textbf{Non-modal} & \textbf{Command} & \textbf{What it does}\\
\midrule
\endhead
none & \texttt{M-x package-\codebreak{}remove-\codebreak{}unused} & \texttt{package-\codebreak{}remove-\codebreak{}unused} & Delete the package sources Tuatara keeps that nothing in this session names any more: a source no package-source names and no load-package fetches from, or a package of a source no load-package names; one left out by :disabled or :if counts as named.\\
\bottomrule
\end{xltabular}

\subsubsection{Where packages are kept}
\label{sec:529}
Tuatara keeps what it fetches in \texttt{\textasciitilde{}/.\codebreak{}local/\codebreak{}share/\codebreak{}tuatara/\codebreak{}packages/}, or under
\texttt{\$XDG\_\codebreak{}DATA\_\codebreak{}HOME} when you have set that. \texttt{*packages-\codebreak{}directory*} puts them
somewhere else. It must be a folder you made yourself, and while it is missing
nothing is fetched or loaded from it. \texttt{*package-\codebreak{}size-\codebreak{}limit*} is the most a
package may take on the disk while it is fetched, its history and files
together. A fetch that grows past it is stopped and its files deleted. The
standard is 512 MiB. \texttt{*package-\codebreak{}sources*} holds the sources named with
\texttt{package-\codebreak{}source}. Set them with that function and not by hand.

\begin{xltabular}{\textwidth}{@{}l l >{\raggedright\arraybackslash\hsize=1.000\hsize}X@{}}
\toprule
\textbf{Setting} & \textbf{Standard} & \textbf{What it changes}\\
\midrule
\endhead
\texttt{*package-\codebreak{}sources*} & \texttt{nil} & The places packages come from, each (NAME .\\
\texttt{*packages-\codebreak{}directory*} & \texttt{nil} & The folder Tuatara keeps the packages it fetches in, in place of \textasciitilde{}/.local/share/tuatara/packages/ (\$XDG\_DATA\_HOME/tuatara/packages/); NIL for that standard place.\\
\texttt{*package-\codebreak{}size-\codebreak{}limit*} & \texttt{536870912} & The most octets a package may take on disk while it is fetched, what git keeps of its history and its files together: a fetch that grows past it is ended and its files deleted.\\
\texttt{*package-\codebreak{}ask-\codebreak{}before-\codebreak{}loading*} & \texttt{t} & True to be asked before a package is first compiled and loaded: one fetched the first time, and again when an update brings another commit; one of a folder of yours the first time only, what you change in it later being yours.\\
\bottomrule
\end{xltabular}

\subsection{What a package needs: check-health}
\label{sec:532}
Some packages need a program that is not part of Tuatara. csvdt needs csvdt,
outlook needs readpst, kitty needs kitty. Tuatara does not install these. It
looks for them, and \texttt{M-x check-\codebreak{}health} reports what it found. The report has a
section for Tuatara and a section for each package that is loaded, in a buffer
called \texttt{*Health*}. A check that passed is drawn in one colour, one that failed
in another, and a value that is missing in the faint one. When a package does
not work, this is the first place to look: it says that a program is not found
and, for some, where to get it.

\begin{figure}[H]
\caption{Figure 19.2. Part of a health report.}\label{org531}
\begin{Verbatim}[breaklines,breakautoindent=false,breakindent=1em,breakafter=/-.\,;:=,commandchars=\\\{\}]
Package csvdt
  csvdt           found     /home/aroha/.local/bin/csvdt
Package outlook
  readpst         not found     install libpst
  weasyprint      found     /usr/bin/weasyprint
\end{Verbatim}
\end{figure}

The check also tells you when signatures are not checked. It runs in the
background, so the report comes when it is ready.

\begin{xltabular}{\textwidth}{@{}l l l >{\raggedright\arraybackslash\hsize=1.000\hsize}X@{}}
\toprule
\textbf{Modal} & \textbf{Non-modal} & \textbf{Command} & \textbf{What it does}\\
\midrule
\endhead
none & \texttt{M-x check-\codebreak{}health} & \texttt{check-\codebreak{}health} & Show what each language and major mode has here.\\
\bottomrule
\end{xltabular}

\begin{xltabular}{\textwidth}{@{}l >{\raggedright\arraybackslash\hsize=1.000\hsize}X@{}}
\toprule
\textbf{Face} & \textbf{What it paints}\\
\midrule
\endhead
\texttt{:list-\codebreak{}good} & A state that is as it should be, in a list: a check that passed in the health report (M-x check-health).\\
\texttt{:list-\codebreak{}bad} & A state that is not as it should be, in a list: a check that failed in the health report (M-x check-health).\\
\texttt{:list-\codebreak{}dim} & What is of less interest in a list, set back: help buffers and the headings of clients that ended in the buffer list (C-x C-b), packages not loaded, bookmarks whose file is gone, files ignored in rimu and in a git status, a value missing in the health report (M-x check-health), a line saying more are not shown.\\
\bottomrule
\end{xltabular}

\subsection{The packages the project publishes}
\label{sec:542}
The project’s packages are not part of the editor. Each is a repository of its
own, \texttt{https://\codebreak{}github.\codebreak{}com/\codebreak{}Tuatara-\codebreak{}Editor/\codebreak{}tuatara-\codebreak{}NAME}, for the name of the
package. You load one with its address as the source, no \texttt{package-\codebreak{}source}
line needed. Each package’s README, in its repository, gives the steps and
the settings, and its settings show in \texttt{C-h v} once it is loaded. Every package
puts its commands under a prefix of its own, \texttt{C-c} and a letter.

\subsubsection{palimpsest}
\label{sec:533}
A PDF can be saved over and over, and each save can leave the earlier text in
the file. palimpsest takes a PDF apart into every save it holds, from the
original to the latest edit, and looks for the signs of editing that examiners
of documents know. It reads the file’s own bytes and runs nothing in it, and it
never changes the file. \texttt{M-x palimpsest} examines the PDF you are looking at,
or asks for one. The report is made in the background and lists findings with
the strongest first. In the report, \texttt{C-c p} is the prefix for its other
commands, such as saving a revision as a PDF of its own.

\subsubsection{csvdt}
\label{sec:534}
csvdt is a separate program that converts the times and durations in CSV
records. The csvdt package runs it on a CSV buffer, a selection or the lines
you banked, and shows the result in \texttt{*csvdt output*}. The buffer the records
came from is never changed. In a CSV or TSV buffer its commands are under
\texttt{C-c v}: \texttt{C-c v r} runs what you chose, and \texttt{C-c v ?} says what a run would
cover without running it. You install csvdt yourself.

\subsubsection{csv-columns}
\label{sec:535}
csv-columns rearranges the columns of every record the way awk’s \texttt{print}
does. \texttt{C-c r a} asks for the columns, counted from 0, so that \texttt{\$0,\$2,\$1} swaps
the second and third. It needs nothing installed. A large file is done in the
background, and each run is one step to undo.

\subsubsection{facematch}
\label{sec:536}
facematch ranks candidate photos by how probable it is that they show the
person in one or more master photos. You mark photos in the folder view, then
\texttt{C-c f f} runs the comparison, in the background and shut in Tuatara’s
sandbox. The comparing engine is a separate program that you download from its
release page, check against the list of sums beside it, and name in
\texttt{*facematch-\codebreak{}program*}. Tuatara never downloads it for you.

\subsubsection{garamond}
\label{sec:537}
Some writers put two spaces after a full stop, others one. garamond keeps the
spacing a document asks for, which the document says in its own \texttt{Local
Variables:} block. The mode line shows \texttt{2S} or \texttt{1S}, and typing one space after
a sentence becomes two where the document wants two. \texttt{M-x garamond-\codebreak{}mode}
turns it on by hand. It needs nothing installed.

\subsubsection{kitty}
\label{sec:538}
kitty opens windows of the kitty terminal, which Tuatara starts itself, so
that you get its fast drawing, its ligatures and its pictures. \texttt{M-x kitty}
opens one in the current folder. \texttt{M-x kitty-\codebreak{}input} types the selection or
the current line into it, and \texttt{M-x kitty-\codebreak{}text} brings its text, scrollback
included, into a buffer you can search. You install kitty yourself.

\subsubsection{outlook}
\label{sec:539}
outlook opens Outlook items (\texttt{.msg}), saved mail (\texttt{.eml}) and mailboxes
(\texttt{.pst}, \texttt{.ost}) as messages, read only. It was made for looking at mail as
evidence: nothing in a message runs, and attachments open only when you ask,
from private copies. In a message, \texttt{P} shows it as it looks, \texttt{O} opens an
attachment in the desktop’s own program after asking, and \texttt{C-c o} is the
prefix for the same commands. In a mailbox, \texttt{C} copies messages out. It needs
readpst and WeasyPrint.

\subsubsection{caldav}
\label{sec:540}
caldav shares Tuatara’s agenda with GNOME Calendar, Thunderbird and phones,
through a CalDAV server (Radicale, or EteSync through its local bridge). The
calendars are kept in one Org file, and only what is under a calendar’s
heading is synced. The password is never in your init file: it is in the
desktop’s keyring. In Org files and the agenda its commands are under
\texttt{C-c s}, with \texttt{C-c s s} syncing now. Its README sets a server up step by step.

\subsubsection{themes}
\label{sec:541}
themes adds fourteen colour themes, each a dark or light version of a palette
you may know from other editors: catppuccin-mocha, catppuccin-latte, dracula,
nord, one-dark, solarized-dark, solarized-light, tokyo-night,
tokyo-night-light, rose-pine, rose-pine-dawn, everforest-dark,
everforest-light and kanagawa-wave. It needs nothing installed.

\texttt{M-x load-\codebreak{}theme} lists every theme, with a description beside each, and
applies the one you choose. To have one at every start, put it after the
package line:

\begin{Verbatim}[breaklines,breakautoindent=false,breakindent=1em,breakafter=/-.\,;:=,commandchars=\\\{\}]
\color[HTML]{A85F00}(load-package\normalcolor{} \color[HTML]{0A7AB3}"themes"
              \color[HTML]{946C00}:\normalcolor{}source \color[HTML]{0A7AB3}"https://github.com/Tuatara-Editor/tuatara-themes"\color[HTML]{A85F00})
\color[HTML]{A85F00}(load-theme\normalcolor{} \color[HTML]{222222}'\color[HTML]{A85F00}catppuccin-mocha)
\end{Verbatim}

\texttt{M-x toggle-\codebreak{}light-\codebreak{}dark-\codebreak{}theme} switches between a light and a dark theme. The
standard pair is Tuatara’s own two, and \texttt{*light-\codebreak{}theme*} and \texttt{*dark-\codebreak{}theme*}
name others. If the window follows the desktop’s light or dark style, these
two are the ones it switches between. Any \texttt{set-\codebreak{}face} line in your init file
stays on top of a theme.

\begin{Verbatim}[breaklines,breakautoindent=false,breakindent=1em,breakafter=/-.\,;:=,commandchars=\\\{\}]
\color[HTML]{A85F00}(setf\normalcolor{} \color[HTML]{A85F00}*dark-theme*\normalcolor{} \color[HTML]{222222}'\color[HTML]{A85F00}tokyo-night
      \color[HTML]{A85F00}*light-theme*\normalcolor{} \color[HTML]{222222}'\color[HTML]{A85F00}tokyo-night-light)
\end{Verbatim}

\begin{xltabular}{\textwidth}{@{}l >{\raggedright\arraybackslash\hsize=0.545\hsize}X l >{\raggedright\arraybackslash\hsize=1.455\hsize}X@{}}
\toprule
\textbf{Modal} & \textbf{Non-modal} & \textbf{Command} & \textbf{What it does}\\
\midrule
\endhead
none & \texttt{M-x load-\codebreak{}theme} & \texttt{load-\codebreak{}theme} & Choose a colour theme, with each theme’s description beside it.\\
none & \texttt{M-x toggle-\codebreak{}light-\codebreak{}dark-\codebreak{}theme} & \texttt{toggle-\codebreak{}light-\codebreak{}dark-\codebreak{}theme} & Switch between the light and the dark theme: the two ao variants, unless the init file sets \textbf{light-theme} and \textbf{dark-theme} to others.\\
\bottomrule
\end{xltabular}

\begin{xltabular}{\textwidth}{@{}l l >{\raggedright\arraybackslash\hsize=1.000\hsize}X@{}}
\toprule
\textbf{Setting} & \textbf{Standard} & \textbf{What it changes}\\
\midrule
\endhead
\texttt{*light-\codebreak{}theme*} & \texttt{tuatara.\codebreak{}defaults::\codebreak{}ao-\codebreak{}light} & The light theme toggle-light-dark-theme switches to.\\
\texttt{*dark-\codebreak{}theme*} & \texttt{tuatara.\codebreak{}defaults::\codebreak{}ao-\codebreak{}dark} & The dark theme toggle-light-dark-theme switches to.\\
\bottomrule
\end{xltabular}

\subsection{Writing a package of your own}
\label{sec:543}
A package is a folder of Lisp files and can be small: one command in one file
is enough. How to write one, which parts of Tuatara it may use, and how to
publish it, is the subject of the developer handbook. To try a folder of your
own on your disk, name it with \texttt{package-\codebreak{}source} and load it as the lines near
the start of this chapter show. Tuatara compiles it and asks about it like any
other.

\section{Making it your own}
\label{sec:574}
Changing Tuatara to suit you, and keeping the changes.

Every chapter so far has said “as standard” about something: the font, the
colours, the keys, the way a journal entry is named. This chapter is about
changing those. You do it with a small file of your own, called the init
file, which Tuatara reads each time it starts. Nothing you put in it is ever
overwritten by an update.

The chapter shows what the file is and where Tuatara keeps it, how to change
a setting, how to give a key to a command, how to move or turn off the modal
states, how to pick a theme and a font, and how to write snippets of your
own. Two things that nothing else in the manual covers, the calendar and the
journal, have a short section each. The last section holds the odds and ends
that have no better home.

You do not need to know Lisp. Every change in this chapter is one line that
you copy, with a word or two changed.

\subsection{The init file}
\label{sec:545}
The init file is a plain text file called \texttt{init.\codebreak{}lisp} in a folder of its own,
\texttt{\textasciitilde{}/.\codebreak{}config/\codebreak{}tuatara/}. The first time Tuatara runs and finds no file there, it
writes one for you. The message \texttt{Wrote a starter init file} says so. After
that Tuatara never writes it again. It is yours.

The starter file lists every setting there is. Each one sits in a comment (a
line that starts with \texttt{;;} is a comment, and Tuatara skips it), with a few
words saying what the setting does and, after \texttt{Not set:}, what it is when you
leave it alone. To use a setting you take away the \texttt{;;} and the space after
it. To go back to the standard, you put them back.

\begin{Verbatim}[breaklines,breakautoindent=false,breakindent=1em,breakafter=/-.\,;:=,commandchars=\\\{\}]
  ;; Two keys typed quickly one after the other in INSERT state leave it for
  ;; NORMAL state, and neither stays in the text. Not set: nil.
  ;; (setf *insert-escape-keys* "jk")
\end{Verbatim}

The block above is as the starter file has it. With the last line’s \texttt{;;}
taken away, it reads \texttt{(setf *insert-\codebreak{}escape-\codebreak{}keys* "jk")} and the setting is on.
Everything is in groups, each begun by a banner of equals signs. Search for
three of them with \texttt{C-s} to jump from one group to the next.

Tuatara reads the file once, from the top, each time it starts. A change you
make is therefore there the next time you start. You do not have to start
again to try it. Put the cursor just after the closing bracket of the line
and press \texttt{C-x C-e}. That runs the line in the editor you are using, and the
echo area shows what it returned. The line stays in the file, so the change
is still there tomorrow.

If a line in the file is wrong, say a bracket is missing or a name is
misspelt, Tuatara tells you on the echo row, with the line, and goes on to
read the rest. It always starts. You can open the file and mend it.

To open the file, press \texttt{C-x C-f}, type \texttt{\textasciitilde{}/.\codebreak{}config/\codebreak{}tuatara/\codebreak{}init.\codebreak{}lisp} and
press \texttt{RET}.

\begin{try}
\textbf{Try it, with the init file}
\begin{enumerate}
\item Press \texttt{C-x C-f}, type \texttt{\textasciitilde{}/.\codebreak{}config/\codebreak{}tuatara/\codebreak{}init.\codebreak{}lisp} and press \texttt{RET}.
\item Press \texttt{g g} (non-modal: \texttt{M-<}), then \texttt{C-s}, type \texttt{fill-\codebreak{}column} and press
   \texttt{RET}. The search stops on the line \texttt{;; (setf *fill-\codebreak{}column* 72)}, which has a
   comment above it about the column that \texttt{M-q} (modal: \texttt{g Q}) fills to.
\item Press \texttt{g h} (non-modal: \texttt{C-a}) to go to the start of the line, then \texttt{l} until
   you are on the first \texttt{;}. Press \texttt{d} three times to take away both semicolons
   and the space.
\item Press \texttt{g l} (non-modal: \texttt{C-e}) to go to the end of the line, then \texttt{C-x C-e}.
   The echo area shows \texttt{72}.
\item Press \texttt{C-h v}, type \texttt{*fill-\codebreak{}column*} and press \texttt{RET}. The help says the
   value is 72.
\item Press \texttt{C-x C-s} to save. Starting Tuatara again keeps the 72.
\end{enumerate}
\end{try}

\subsection{Where Tuatara keeps its files}
\label{sec:549}
Tuatara keeps what it makes in three places, each with a different job. The
division follows the standard that most programs on Linux follow, so a tool
that backs up or cleans your home folder already knows about them.

\begin{itemize}
\item The configuration folder, \texttt{\textasciitilde{}/.\codebreak{}config/\codebreak{}tuatara/}, holds what you wrote or
  chose. This is the init file, your snippets, your note and export
  templates, and your logos. If you lose it you lose your choices.
\item The state folder, \texttt{\textasciitilde{}/.\codebreak{}local/\codebreak{}state/\codebreak{}tuatara/}, holds what Tuatara
  remembers about how you work: the files you opened lately, where you were
  in each, your bookmarks, the session, and the copies of unsaved changes.
  If you lose it nothing you wrote is gone, but Tuatara forgets its habits.
\item The cache folder, \texttt{\textasciitilde{}/.\codebreak{}cache/\codebreak{}tuatara/}, holds things Tuatara can make
  again: previews, pages drawn from documents, copies of files opened
  safely, compiled packages and the log of what happened. You can delete it
  at any time.
\end{itemize}

Fetched packages are kept in a fourth place, \texttt{\textasciitilde{}/.\codebreak{}local/\codebreak{}share/\codebreak{}tuatara/},
described in chapter 19.

Figure 20.1 shows the folders and what each holds.

\begin{figure}[H]
\caption{Figure 20.1. The folders Tuatara uses, and what is in each.}\label{org546}
\begin{Verbatim}[breaklines,breakautoindent=false,breakindent=1em,breakafter=/-.\,;:=,commandchars=\\\{\}]
 ~/.config/tuatara/            what you wrote
   init.lisp                     your settings and keys
   snippets/                     your snippets, a folder for each mode
   note-templates/ ...           templates and logos for notes and exports
   logos/

 ~/.local/state/tuatara/       what Tuatara remembers
   recent-files   places   bookmarks   session
   auto-save/                    copies of changes not yet saved

 ~/.cache/tuatara/             what can be made again
   debug.log   stderr.log        what Tuatara said
   previews, pages, compiled packages ...

 ~/.local/share/tuatara/       what was fetched
   packages/                     packages from other people
\end{Verbatim}
\end{figure}

If your computer sets \texttt{XDG\_\codebreak{}CONFIG\_\codebreak{}HOME}, \texttt{XDG\_\codebreak{}STATE\_\codebreak{}HOME}, \texttt{XDG\_\codebreak{}CACHE\_\codebreak{}HOME} or
\texttt{XDG\_\codebreak{}DATA\_\codebreak{}HOME}, Tuatara uses a folder called \texttt{tuatara} inside each of them
instead of the ones in the figure. You can also set the folders yourself,
which is useful when you keep them on an encrypted drive. The three settings
below do it. Make the folder first, with \texttt{mkdir -m 700}, so that only you can
read it. While the folder is not there, because the drive is not mounted,
Tuatara keeps those things in memory and says so.

\begin{Verbatim}[breaklines,breakautoindent=false,breakindent=1em,breakafter=/-.\,;:=,commandchars=\\\{\}]
\color[HTML]{A85F00}(setf\normalcolor{} \color[HTML]{A85F00}*state-directory*\normalcolor{} \color[HTML]{0A7AB3}"/run/media/aroha/Data/tuatara/state/"\color[HTML]{A85F00})
\end{Verbatim}

\subsubsection{What to copy to another computer}
\label{sec:547}
To have your Tuatara on another computer, copy the configuration folder,
\texttt{\textasciitilde{}/.\codebreak{}config/\codebreak{}tuatara/}, and nothing else. It holds the init file and everything
you wrote. Packages are fetched again on the new computer when the init file
asks for them.

Copy the state folder as well if you want the new computer to know your
recent files and bookmarks. Do not copy the cache. It is rebuilt as you
work, and parts of it are only right for the computer that made them.

Paths in the init file, such as a folder for your journal, must exist on
the new computer too.

\subsubsection{Settings for the folders}
\label{sec:548}
\begin{xltabular}{\textwidth}{@{}l l >{\raggedright\arraybackslash\hsize=1.000\hsize}X@{}}
\toprule
\textbf{Setting} & \textbf{Standard} & \textbf{What it changes}\\
\midrule
\endhead
\texttt{*state-\codebreak{}directory*} & \texttt{nil} & The folder Tuatara keeps what it remembers between sessions in, in place of \textasciitilde{}/.local/state/tuatara/ (\$XDG\_STATE\_HOME/tuatara/): the files visited lately, places in them, the session, auto-saved text, the packages fetched.\\
\texttt{*cache-\codebreak{}directory*} & \texttt{nil} & The folder Tuatara keeps its cache in, in place of \textasciitilde{}/.cache/tuatara/ (\$XDG\_CACHE\_HOME/tuatara/): files made again when they are gone, as compiled packages, previews and pages drawn, copies of files opened outside, the debug log and stderr.log.\\
\texttt{*packages-\codebreak{}directory*} & \texttt{nil} & The folder Tuatara keeps the packages it fetches in, in place of \textasciitilde{}/.local/share/tuatara/packages/ (\$XDG\_DATA\_HOME/tuatara/packages/); NIL for that standard place.\\
\bottomrule
\end{xltabular}

\begin{try}
\textbf{Try it, with the folders}
\begin{enumerate}
\item In a terminal, type \texttt{ls \textasciitilde{}/.\codebreak{}config/\codebreak{}tuatara/ \textasciitilde{}/.\codebreak{}local/\codebreak{}state/\codebreak{}tuatara/}. The
   first list holds \texttt{init.\codebreak{}lisp}. The second holds files such as \texttt{bookmarks}
   and \texttt{recent-\codebreak{}files}, and the folder \texttt{auto-\codebreak{}save}.
\item In Tuatara, press \texttt{C-h v}, type \texttt{*state-\codebreak{}directory*} and press \texttt{RET}. The
   value is \texttt{NIL}, which means the standard folder of the figure.
\end{enumerate}
\end{try}

\subsection{Changing a setting}
\label{sec:550}
A setting is a name that holds a value. Its name starts and ends with a star,
\texttt{*fill-\codebreak{}column*}, so that it is easy to tell from a command. You change one
with \texttt{setf}, which means “set to”:

\begin{Verbatim}[breaklines,breakautoindent=false,breakindent=1em,breakafter=/-.\,;:=,commandchars=\\\{\}]
\color[HTML]{A85F00}(setf\normalcolor{} \color[HTML]{A85F00}*fill-column*\normalcolor{} \color[HTML]{0A7AB3}72\color[HTML]{A85F00})
\end{Verbatim}

The line reads: set \texttt{*fill-\codebreak{}column*} to 72. The value can be a number, a word
in quotes, \texttt{t} for yes, \texttt{nil} for no, or a name that starts with a colon, such
as \texttt{:bar}. The help of each setting says which it takes.

\texttt{C-h v} asks for a setting’s name and shows its value, what it means, and
where it was defined. It is the quickest way to find out what you can put
after the name. \texttt{M-:} lets you type a line without opening the init file,
which is useful to try a value out before you save it.

Some settings hold a list of things, for example the modes that do something.
To add to one instead of replacing it, put your item in front of what it
holds:

\begin{Verbatim}[breaklines,breakautoindent=false,breakindent=1em,breakafter=/-.\,;:=,commandchars=\\\{\}]
\color[HTML]{A85F00}(push\normalcolor{} \color[HTML]{0A7AB3}"C-;"\normalcolor{} \color[HTML]{A85F00}*terminal-editor-keys*)
\end{Verbatim}

A few settings are better said with a function than with \texttt{setf}, because
they take a mode and a value. The starter file has these written out, for
example \texttt{(set-\codebreak{}tab-\codebreak{}width 'c-\codebreak{}mode 8)}. Keep to the form the starter file shows.

Chapter 21 lists every setting and what it is as standard.

\subsection{Keys of your own}
\label{sec:551}
\texttt{bind-\codebreak{}key} gives a key to a command. It is the main command of the init file.
You give it the keys and the command’s name, which is the name that \texttt{C-h f}
and \texttt{M-x} show. Keys are written as the help writes them: \texttt{C-} is Control,
\texttt{M-} is Meta (Alt on most keyboards) and \texttt{S-} is Shift. Keys that you press one
after the other are separated by a space.

\begin{Verbatim}[breaklines,breakautoindent=false,breakindent=1em,breakafter=/-.\,;:=,commandchars=\\\{\}]
\color[HTML]{A85F00}(bind-key\normalcolor{} \color[HTML]{0A7AB3}"C-c t"\normalcolor{} \color[HTML]{222222}'\color[HTML]{A85F00}toggle-light-dark-theme)
\end{Verbatim}

After that line, \texttt{C-c t} switches between the light and the dark theme in
every buffer and in both states. The mark before the command’s name is a
single quote, and it belongs there.

Where a key goes depends on the map you name with \texttt{:map}. Table 20.1 lists the
forms you will use. In all of them you can put any command in place of the one
shown.

\begin{xltabular}{\textwidth}{@{}>{\raggedright\arraybackslash\hsize=1.067\hsize}X >{\raggedright\arraybackslash\hsize=0.933\hsize}X@{}}
\caption{Table 20.1. The ways to bind a key.}\\
\toprule
\textbf{Form} & \textbf{What it does}\\
\midrule
\endfirsthead
\caption[]{Table 20.1. The ways to bind a key. (continued)}\\
\toprule
\textbf{Form} & \textbf{What it does}\\
\midrule
\endhead
\texttt{(bind-\codebreak{}key "C-c t" 'command)} & The key works in every buffer and both states\\
\texttt{(bind-\codebreak{}key "o" 'command :map :leader)} & \texttt{SPC o} in NORMAL state, and \texttt{C-x SPC o} everywhere\\
\texttt{(bind-\codebreak{}key "e" 'command :map :normal)} & \texttt{e} in NORMAL state only\\
\texttt{(bind-\codebreak{}key "C-c e" 'command :map 'csv-\codebreak{}mode)} & The key in CSV buffers only\\
\texttt{(bind-\codebreak{}key "m l" 'command :map :leader :mode '(org-\codebreak{}mode markdown-\codebreak{}mode))} & \texttt{SPC m l} in Org and Markdown only\\
\texttt{(bind-\codebreak{}key "t n" '("numbers" . command) :map :leader)} & \texttt{SPC t n}, shown as “numbers” in the key hints\\
\texttt{(bind-\codebreak{}key "t" nil :map :leader :label "toggles")} & The label of a prefix: \texttt{SPC} then shows \texttt{t +toggles}\\
\texttt{(unbind-\codebreak{}key "C-c t")} & Takes a key away\\
\bottomrule
\end{xltabular}

Three more things about the table. First, the maps. \texttt{:global} is the standard,
every buffer in both states. \texttt{:normal}, \texttt{:insert} and \texttt{:support} are the keys of
the NORMAL state, the INSERT state and the buffers that Tuatara makes itself,
such as lists and help. \texttt{:leader} is the leader, described in the next
section. Other maps are named after one kind of question: \texttt{:pattern}, for
instance, is every question that asks for a pattern to search for. A map can
also be the name of a major mode, as the CSV line shows. Then the key is
only there in buffers of that mode and the modes made from it.

Second, \texttt{:mode} is a quicker way to say the same. It takes one mode or a
list. A key bound for a mode never changes what the same key does in other
modes, so you can give \texttt{SPC m l} one meaning in Org and another in code
without either getting in the other’s way.

Third, labels. The key hints (chapter 8) show the name of a command beside
each key. A label is a name of your own to show instead. The two ways to give
one are in the table: put the label and the command together in a pair, or
bind a prefix to \texttt{nil} and give it a label.

Which keys are free? \texttt{C-h k} followed by a key says what the key does now. A key
with nothing bound says so. \texttt{C-c} followed by a letter is a good place to start,
and \texttt{C-c t} is free as standard. The letters \texttt{h j k l} only ever move the cursor
in NORMAL state. Do not give them another job.

\begin{try}
\textbf{Try it, with a key of your own}
\begin{enumerate}
\item Open the init file with \texttt{C-x C-f} and go to its end with \texttt{g e} (non-modal:
   \texttt{M->}).
\item Press \texttt{o} to open a new line and type this line.
   \texttt{(bind-\codebreak{}key "C-c t" 'toggle-\codebreak{}light-\codebreak{}dark-\codebreak{}theme)}
\item Press \texttt{C-g} to go back to NORMAL state, then \texttt{C-x C-e}. The echo area
   shows the name of the command.
\item Press \texttt{C-c t}. The colours turn light. Press it again and they are dark
   again.
\item Press \texttt{C-x C-s} to save. To take the key away again, put \texttt{(unbind-\codebreak{}key
   "C-c t")} under the line, or put \texttt{;;} in front of it.
\end{enumerate}
\end{try}

\subsection{A leader key of your own}
\label{sec:552}
The leader is a map of keys that belong to you. As Tuatara comes, nothing is
under it, so there is nothing to learn that you did not put there. In NORMAL
state, in lists and in the calendar, \texttt{SPC} opens it. \texttt{C-x SPC} opens it in
every buffer and state, INSERT state and a session with modal editing off
too, so a leader key works wherever you are.

The starter file lists, commented out, a set of leader keys that suit many
people: \texttt{SPC f} to find a file, \texttt{SPC b} to go to the other buffer, \texttt{SPC w} for
a family of window keys, and so on. Take away the \texttt{;;} in front of the ones you
want, or write your own with \texttt{:map :leader}.

\begin{Verbatim}[breaklines,breakautoindent=false,breakindent=1em,breakafter=/-.\,;:=,commandchars=\\\{\}]
\color[HTML]{A85F00}(bind-key\normalcolor{} \color[HTML]{0A7AB3}"o"\normalcolor{} \color[HTML]{222222}'\color[HTML]{A85F00}find-file\normalcolor{} \color[HTML]{946C00}:\normalcolor{}map \color[HTML]{946C00}:\normalcolor{}leader\color[HTML]{A85F00})
\color[HTML]{A85F00}(bind-key\normalcolor{} \color[HTML]{0A7AB3}"t"\normalcolor{} \color[HTML]{0A7AB3}nil\normalcolor{} \color[HTML]{946C00}:\normalcolor{}map \color[HTML]{946C00}:\normalcolor{}leader \color[HTML]{946C00}:\normalcolor{}label \color[HTML]{0A7AB3}"toggles"\color[HTML]{A85F00})
\color[HTML]{A85F00}(bind-key\normalcolor{} \color[HTML]{0A7AB3}"t l"\normalcolor{} \color[HTML]{222222}'\color[HTML]{A85F00}soft-wrap-mode\normalcolor{} \color[HTML]{946C00}:\normalcolor{}map \color[HTML]{946C00}:\normalcolor{}leader\color[HTML]{A85F00})
\color[HTML]{A85F00}(bind-key\normalcolor{} \color[HTML]{0A7AB3}"t h"\normalcolor{} \color[HTML]{222222}'\color[HTML]{A85F00}hl-line-mode\normalcolor{} \color[HTML]{946C00}:\normalcolor{}map \color[HTML]{946C00}:\normalcolor{}leader\color[HTML]{A85F00})
\end{Verbatim}

The four lines make \texttt{SPC o} open a file. They also make \texttt{SPC t} a prefix with
the label “toggles”, with \texttt{SPC t l} to wrap long lines and \texttt{SPC t h} to
highlight the line the cursor is on. After you press \texttt{SPC}, the key hints
show what is under it, so you do not have to remember.

If \texttt{SPC} is a key you would like for something else, move the leader.

\begin{Verbatim}[breaklines,breakautoindent=false,breakindent=1em,breakafter=/-.\,;:=,commandchars=\\\{\}]
\color[HTML]{A85F00}(setf\normalcolor{} \color[HTML]{A85F00}*leader-key*\normalcolor{} \color[HTML]{0A7AB3}","\color[HTML]{A85F00})
\end{Verbatim}

Now the leader is on \texttt{,} in NORMAL state, and \texttt{SPC} is free again. \texttt{C-x SPC}
does not move. Chapter 8 lists the leader’s settings with the settings for
modal editing.

\subsection{Modal editing, and turning it off}
\label{sec:553}
Tuatara comes with modal editing on. A file opens in NORMAL state, where
letters are commands. \texttt{i} starts INSERT state, where letters are typed, and
\texttt{C-g} goes back. Chapter 1 introduced it.

Some people would rather type at once and use the keys with Control and Meta
for commands, as other editors do. If that is you, turn modal editing off:

\begin{Verbatim}[breaklines,breakautoindent=false,breakindent=1em,breakafter=/-.\,;:=,commandchars=\\\{\}]
\color[HTML]{A85F00}(setf\normalcolor{} \color[HTML]{A85F00}*modal-editing*\normalcolor{} \color[HTML]{0A7AB3}nil\color[HTML]{A85F00})
\end{Verbatim}

Every buffer then takes the keys of its mode and the global keys. Typing puts
letters in the text. The states are gone from the mode line. Every command
is still there, to run with \texttt{M-x} or to bind to a key. The leader stays,
under \texttt{C-x SPC}. The commands that selecting, cursors, spelling and replacing
have in NORMAL state are under \texttt{C-z}, on the same letters: \texttt{C-z m f} selects
the function around the cursor, \texttt{C-z t} adds a cursor below.

\begin{try}
\textbf{Try it, with modal editing off}
\begin{enumerate}
\item In the init file, press \texttt{g e} (non-modal: \texttt{M->}) and \texttt{o}, and type
   \texttt{(setf *modal-\codebreak{}editing* nil)}. Press \texttt{C-g}, then \texttt{C-x C-e}. The echo area
   shows \texttt{NIL}.
\item Open another file with \texttt{C-x C-f}. The mode line no longer shows a state
   and the letters you type go into the text.
\item To go back, change the line to \texttt{(setf *modal-\codebreak{}editing* t)} and press
   \texttt{C-x C-e} again. You can also take the line out of the file.
\end{enumerate}
\end{try}

You can keep modal editing and make leaving INSERT state easier. \texttt{C-g} is
the key that does it as standard, and so does \texttt{ESC}. If you would rather type
two letters, set a pair that hardly ever follows one another in your
writing.

\begin{Verbatim}[breaklines,breakautoindent=false,breakindent=1em,breakafter=/-.\,;:=,commandchars=\\\{\}]
\color[HTML]{A85F00}(setf\normalcolor{} \color[HTML]{A85F00}*insert-escape-keys*\normalcolor{} \color[HTML]{0A7AB3}"jk"\color[HTML]{A85F00})
\end{Verbatim}

Typing \texttt{j} puts a \texttt{j} in the text at once, like any key. If \texttt{k} comes within a
fifth of a second, the \texttt{j} is taken out again and you are in NORMAL state.
Slower, or with another key between, both letters stay. It works while you
type in an ordinary buffer. It does not work in the minibuffer, in a terminal
or with several cursors.

Chapter 8 describes \texttt{*escape-\codebreak{}quits*}, which makes \texttt{ESC} do all that \texttt{C-g} does.

\begin{xltabular}{\textwidth}{@{}l l >{\raggedright\arraybackslash\hsize=1.000\hsize}X@{}}
\toprule
\textbf{Setting} & \textbf{Standard} & \textbf{What it changes}\\
\midrule
\endhead
\texttt{*insert-\codebreak{}escape-\codebreak{}keys*} & \texttt{nil} & NIL, or a string of two characters, “jk” say: typed one right after the other in INSERT state, they leave it for NORMAL state as M-x tuatara-normal does, and neither stays in the text.\\
\bottomrule
\end{xltabular}

\subsection{Themes and faces}
\label{sec:558}
A theme is a set of colours for everything Tuatara draws. A face is the look
of one kind of thing: how a comment is drawn, or the line you are on, or the
mode line. A theme gives every face its look.

\subsubsection{Choosing a theme}
\label{sec:554}
\texttt{M-x load-\codebreak{}theme} lists the themes, each with a line saying what it is like,
and puts the one you choose in force. There are four.

\begin{itemize}
\item \texttt{ao-\codebreak{}dark} is the standard: a deep navy ground with an amber accent.
\item \texttt{ao-\codebreak{}light} is the same palette worked for a white ground.
\item \texttt{terminal} uses the sixteen colours of your terminal, so Tuatara follows
  whatever colours the terminal uses.
\item \texttt{plain} uses no colours at all, only bold, underline, reverse video and
  dim. It is for terminals that cannot show colour.
\end{itemize}

The theme you choose with \texttt{M-x load-\codebreak{}theme} lasts until you leave. To keep it,
put it in the init file:

\begin{Verbatim}[breaklines,breakautoindent=false,breakindent=1em,breakafter=/-.\,;:=,commandchars=\\\{\}]
\color[HTML]{A85F00}(load-theme\normalcolor{} \color[HTML]{222222}'\color[HTML]{A85F00}ao-light)
\end{Verbatim}

In a window of its own, Tuatara follows the desktop. If the desktop turns
dark in the evening, Tuatara switches to its dark theme, and back again in
the morning. \texttt{*follow-\codebreak{}desktop-\codebreak{}style*} set to \texttt{nil} stops that and keeps one
theme. \texttt{M-x toggle-\codebreak{}light-\codebreak{}dark-\codebreak{}theme} switches between the dark and the light
theme by hand. \texttt{*light-\codebreak{}theme*} and \texttt{*dark-\codebreak{}theme*} say which two themes it,
and the desktop, switch between.

\begin{Verbatim}[breaklines,breakautoindent=false,breakindent=1em,breakafter=/-.\,;:=,commandchars=\\\{\}]
\color[HTML]{A85F00}(setf\normalcolor{} \color[HTML]{A85F00}*light-theme*\normalcolor{} \color[HTML]{222222}'\color[HTML]{A85F00}plain)
\color[HTML]{A85F00}(setf\normalcolor{} \color[HTML]{A85F00}*dark-theme*\normalcolor{} \color[HTML]{222222}'\color[HTML]{A85F00}terminal)
\end{Verbatim}

\begin{try}
\textbf{Try it, with a theme}
\begin{enumerate}
\item Press \texttt{M-x}, type \texttt{load-\codebreak{}theme} and press \texttt{RET}. Choose \texttt{ao-\codebreak{}light} with
   the arrow keys and \texttt{RET}. The colours turn light.
\item Press \texttt{M-x} and run \texttt{toggle-\codebreak{}light-\codebreak{}dark-\codebreak{}theme}. They turn dark again.
\item To keep the light theme, open the init file and add
   \texttt{(load-\codebreak{}theme 'ao-\codebreak{}light)} on a line of its own. Press \texttt{C-g}, then
   \texttt{C-x C-e} and \texttt{C-x C-s}.
\end{enumerate}
\end{try}

\subsubsection{Changing one face}
\label{sec:555}
\texttt{M-x list-\codebreak{}faces} shows every face, each with a sample drawn as the theme in
force draws it, and a line saying what it colours. \texttt{M-x describe-\codebreak{}face} shows
all that is known of one face, including what the theme and your init file
set for it. Both are covered in chapter 8.

\texttt{set-\codebreak{}face} changes a face. It takes the face’s name and then pairs of an
attribute and a value:

\begin{Verbatim}[breaklines,breakautoindent=false,breakindent=1em,breakafter=/-.\,;:=,commandchars=\\\{\}]
\color[HTML]{A85F00}(set-face\normalcolor{} \color[HTML]{946C00}:\normalcolor{}comment \color[HTML]{946C00}:\normalcolor{}fg \color[HTML]{946C00}:\normalcolor{}cyan \color[HTML]{946C00}:\normalcolor{}italic \color[HTML]{A85F00}t)
\color[HTML]{A85F00}(set-face\normalcolor{} \color[HTML]{946C00}:\normalcolor{}hl-line \color[HTML]{946C00}:\normalcolor{}bg \color[HTML]{0A7AB3}"#1a2a3a"\color[HTML]{A85F00})
\end{Verbatim}

The first line draws comments in cyan italics. The second gives the line you are
on a dark blue ground. A colour is a name such as \texttt{:red}, a number from 0 to
255, or a six-digit code in quotes. The attributes are \texttt{:fg} and \texttt{:bg} for the
text and the ground, \texttt{:bold}, \texttt{:italic}, \texttt{:dim} and \texttt{:reverse}, which each take
\texttt{t} or \texttt{nil}, and \texttt{:underline}. In a window of its own a face can also have a
\texttt{:font} and a \texttt{:height}.

What you set stays on top of any theme. Choose another theme and your comment
colour is still yours.

\subsubsection{A theme of your own}
\label{sec:556}
When a theme is nearly right, define one of your own that changes only what
you do not like. Name it, say in a few words what it is, say which theme to
start from, and give the colours that differ:

\begin{Verbatim}[breaklines,breakautoindent=false,breakindent=1em,breakafter=/-.\,;:=,commandchars=\\\{\}]
\color[HTML]{A85F00}(define-theme\normalcolor{} \color[HTML]{A85F00}my-ao
  \color[HTML]{0A7AB3}"Ao on black."
  \color[HTML]{946C00}:\normalcolor{}faces-from \color[HTML]{A85F00}ao-dark
  \color[HTML]{946C00}:\normalcolor{}palette \color[HTML]{A85F00}((ground\normalcolor{} \color[HTML]{0A7AB3}"#000000"\color[HTML]{A85F00})))
\color[HTML]{A85F00}(load-theme\normalcolor{} \color[HTML]{222222}'\color[HTML]{A85F00}my-ao)
\end{Verbatim}

The colours of a theme go by role: \texttt{ground} is the background, \texttt{text} the
text, \texttt{comment} the comments, \texttt{blue} and \texttt{green} and so on. The palette says
what each role is. In the example only \texttt{ground} changes, from navy to black,
and everything drawn on the ground follows. After \texttt{define-\codebreak{}theme}, \texttt{M-x
load-\codebreak{}theme} lists \texttt{my-ao} with the others.

\subsubsection{The faces}
\label{sec:557}
The faces in the table below are the ones the other chapters do not
describe. \texttt{:default} is the text and the ground, the face under all others.
The next four colour the lines of a log or a message when they say error,
warning, information or a hint.

\begin{xltabular}{\textwidth}{@{}l >{\raggedright\arraybackslash\hsize=1.000\hsize}X@{}}
\toprule
\textbf{Face} & \textbf{What it paints}\\
\midrule
\endhead
\texttt{:default} & Text with no other face, and the background: plain, so the terminal’s own colours show until a theme sets them.\\
\texttt{:error} & Errors marked in text, such as the error lines of a log.\\
\texttt{:warning} & Warnings marked in text, such as the warning lines of a log.\\
\texttt{:info} & Information marked in text, such as the info lines of a log.\\
\texttt{:hint} & Hints and debugging marked in text, such as the debug lines of a log.\\
\bottomrule
\end{xltabular}

Syntax colouring is the colour a language’s words get as you write code in
it: keywords, strings, numbers, the names of functions and so on. The faces
below name the kinds of word a language has. Chapter 13 explains how a file
is coloured. A theme gives each of these faces a colour, and you can
change any with \texttt{set-\codebreak{}face}. Their names are the ones tree-sitter, the
parser that finds the words, gives them.

\begin{xltabular}{\textwidth}{@{}l >{\raggedright\arraybackslash\hsize=1.000\hsize}X@{}}
\toprule
\textbf{Face} & \textbf{What it paints}\\
\midrule
\endhead
\texttt{:attribute} & Attributes, annotations and decorators.\\
\texttt{:constant.\codebreak{}character.\codebreak{}escape} & Escape sequences in strings, captured as tree-sitter names them.\\
\texttt{:constant.\codebreak{}numeric} & Numbers, captured as tree-sitter names them.\\
\texttt{:constructor} & Constructors: enum variants and structures built by name.\\
\texttt{:escape} & Escape sequences inside strings.\\
\texttt{:function.\codebreak{}macro} & Macros, where defined or called.\\
\texttt{:function.\codebreak{}method} & Methods, where defined or called.\\
\texttt{:function.\codebreak{}special} & Special functions, such as preprocessor macros.\\
\texttt{:keyword.\codebreak{}directive} & Preprocessor directives and pragmas, captured as tree-sitter names them.\\
\texttt{:keyword.\codebreak{}operator} & Keywords that are operators, such as and, or and as.\\
\texttt{:keyword.\codebreak{}storage.\codebreak{}modifier} & Storage modifiers such as mut, static and const.\\
\texttt{:label} & Labels.\\
\texttt{:namespace} & Namespaces, modules and packages.\\
\texttt{:property} & Fields, properties and keys.\\
\texttt{:punctuation.\codebreak{}special} & Punctuation with a meaning of its own, such as string interpolation.\\
\texttt{:string.\codebreak{}regexp} & Regular expressions written as strings.\\
\texttt{:variable.\codebreak{}other} & Other variables: fields and members.\\
\bottomrule
\end{xltabular}

Markdown, Org and the other languages for writing text have faces for the
parts of a page: headings, emphasis, links, quotations, code and lists.

\begin{xltabular}{\textwidth}{@{}l >{\raggedright\arraybackslash\hsize=1.000\hsize}X@{}}
\toprule
\textbf{Face} & \textbf{What it paints}\\
\midrule
\endhead
\texttt{:heading} & Headings in prose and markup.\\
\texttt{:emphasis} & Emphasised text in markup.\\
\texttt{:strong} & Strongly emphasised text in markup.\\
\texttt{:tag} & Tags in markup languages.\\
\texttt{:markup.\codebreak{}bold} & Bold text in markup.\\
\texttt{:markup.\codebreak{}italic} & Italic text in markup.\\
\texttt{:markup.\codebreak{}strikethrough} & Text struck through in markup; plain by default.\\
\texttt{:markup.\codebreak{}heading} & Headings in markup, and the marks that make them (the \# of Markdown).\\
\texttt{:markup.\codebreak{}heading.1} & Headings of the first level in markup.\\
\texttt{:markup.\codebreak{}heading.2} & Headings of the second level in markup.\\
\texttt{:markup.\codebreak{}heading.3} & Headings of the third level in markup.\\
\texttt{:markup.\codebreak{}heading.4} & Headings of the fourth level in markup.\\
\texttt{:markup.\codebreak{}heading.5} & Headings of the fifth level in markup.\\
\texttt{:markup.\codebreak{}heading.6} & Headings of the sixth level in markup.\\
\texttt{:markup.\codebreak{}link} & Links in markup.\\
\texttt{:markup.\codebreak{}link.\codebreak{}label} & The labels links in markup refer to, such as [label].\\
\texttt{:markup.\codebreak{}link.\codebreak{}text} & The text links in markup show.\\
\texttt{:markup.\codebreak{}link.\codebreak{}url} & The addresses links in markup lead to.\\
\texttt{:markup.\codebreak{}list} & The marks that begin the items of lists in markup.\\
\texttt{:markup.\codebreak{}quote} & Quotations in markup.\\
\texttt{:markup.\codebreak{}raw} & Code and verbatim text in markup.\\
\texttt{:markup.\codebreak{}raw.\codebreak{}block} & Blocks of code and verbatim text in markup.\\
\bottomrule
\end{xltabular}

\subsection{Fonts}
\label{sec:560}
In a window of its own (\texttt{tuatara --gui}), Tuatara draws the text in a font you
choose. In a terminal the font is the terminal’s own, and these commands
say so and change nothing.

To make the text larger or smaller for the moment, press \texttt{C-x C-+}. That makes
it a step larger, about a fifth. Each press of \texttt{+} after it goes on. \texttt{-} makes
it smaller, and \texttt{0} returns to the size it started with. Holding Ctrl and
turning the mouse wheel does the same. The size is for this session. When
you want a size for good, change the font in the init file.

\texttt{M-x font-\codebreak{}list} shows every font installed on the computer, the ones whose
letters are all the same width first, each with a sample written in it. Move
to a line and press \texttt{RET} to see the window in that font now. Press \texttt{w} to
copy the line that chooses it for good. Then paste it into your init file
with \texttt{C-y}.

\begin{Verbatim}[breaklines,breakautoindent=false,breakindent=1em,breakafter=/-.\,;:=,commandchars=\\\{\}]
\color[HTML]{A85F00}(setf\normalcolor{} \color[HTML]{A85F00}*gui-font*\normalcolor{} \color[HTML]{0A7AB3}"Adwaita Mono 14"\color[HTML]{A85F00})
\end{Verbatim}

A font is a family, an optional style and a size in points. The text window
works best with a font whose letters are all one width, since the text is
laid in columns. \texttt{M-x set-\codebreak{}frame-\codebreak{}font} is another way to the same change by
name: type a family and a size after it. \texttt{<f11>} makes the window fill the
screen, and \texttt{M-<f10>} makes it as large as the screen allows with its
title bar. Press either again to go back.

\begin{try}
\textbf{Try it, with a bigger font}
\begin{enumerate}
\item In a window (\texttt{tuatara --gui}), press \texttt{C-x C-+}. The text grows by a step.
   Press \texttt{+} twice more, then \texttt{0} to go back to the standard size.
\item Press \texttt{M-x}, type \texttt{font-\codebreak{}list} and press \texttt{RET}. Move to a line of a font
   you like with \texttt{j} and \texttt{k} and press \texttt{RET}. The window changes at once.
\item Press \texttt{w}. The echo area says the line is copied.
\item Open the init file, press \texttt{g e} (non-modal: \texttt{M->}), \texttt{o}, and paste with
   \texttt{C-y}. Change the size in the quotes to 14. Press \texttt{C-g}, \texttt{C-x C-e}, then
   \texttt{C-x C-s}.
\end{enumerate}
\end{try}

\subsubsection{Commands and faces for fonts and the window}
\label{sec:559}
\begin{xltabular}{\textwidth}{@{}l >{\raggedright\arraybackslash\hsize=0.636\hsize}X l >{\raggedright\arraybackslash\hsize=1.364\hsize}X@{}}
\toprule
\textbf{Modal} & \textbf{Non-modal} & \textbf{Command} & \textbf{What it does}\\
\midrule
\endhead
none & \texttt{M-x font-\codebreak{}list} & \texttt{font-\codebreak{}list} & List every font installed in \textbf{Fonts}, the fixed-width ones first: each family with its variants, the line of the init file that chooses each, and a sample written in it.\\
none & \texttt{RET} & \texttt{font-\codebreak{}list-\codebreak{}use} & Show the window in the font on the current line of the list of fonts, for this session; w copies the line that chooses it for good.\\
none & \texttt{w}, \texttt{C-c w} & \texttt{font-\codebreak{}list-\codebreak{}copy} & Copy the init-file line that chooses the font on the current line, to put in the init file (C-x C-f \textasciitilde{}/.config/tuatara/init.lisp).\\
none & \texttt{M-x set-\codebreak{}frame-\codebreak{}font} & \texttt{set-\codebreak{}frame-\codebreak{}font} & Choose the window’s font among the fixed-width fonts installed: a family, with any style and size typed after it (DejaVu Sans Mono Bold 12); without a size, the size in use stays, and so does any scaling by C-x C-+.\\
none & \texttt{C-x C-+}, \texttt{C-x C--}, \texttt{C-x C-0} & \texttt{text-\codebreak{}scale-\codebreak{}adjust} & Make the window’s font larger or smaller, by the last key of the keys that ran it: - a step smaller, 0 the size it started with, any other (+, =) a step larger, 1.2 times; with a prefix argument, that many steps.\\
none & \texttt{M-x text-\codebreak{}scale-\codebreak{}increase} & \texttt{text-\codebreak{}scale-\codebreak{}increase} & Make the window’s font a step larger (1.2 times), or the prefix argument’s worth of steps.\\
none & \texttt{M-x text-\codebreak{}scale-\codebreak{}decrease} & \texttt{text-\codebreak{}scale-\codebreak{}decrease} & Make the window’s font a step smaller, or the prefix argument’s worth of steps.\\
none & \texttt{C-<wheel-\codebreak{}up>}, \texttt{C-<wheel-\codebreak{}down>} & \texttt{mouse-\codebreak{}wheel-\codebreak{}text-\codebreak{}scale} & Make the window’s font a step larger as the mouse wheel turns up with Ctrl held (C-<wheel-up>), a step smaller as it turns down.\\
none & \texttt{<f11>} & \texttt{toggle-\codebreak{}frame-\codebreak{}fullscreen} & Make the window fill the screen, without its title bar, or go back to how it was.\\
none & \texttt{M-<f10>} & \texttt{toggle-\codebreak{}frame-\codebreak{}maximized} & Make the window as large as the screen allows, or go back to its size before.\\
\bottomrule
\end{xltabular}

\begin{xltabular}{\textwidth}{@{}l >{\raggedright\arraybackslash\hsize=1.000\hsize}X@{}}
\toprule
\textbf{Face} & \textbf{What it paints}\\
\midrule
\endhead
\texttt{:font-\codebreak{}list-\codebreak{}line} & The line for the init file that chooses a font, in the list of fonts (M-x font-list).\\
\texttt{:list-\codebreak{}name} & The name of a thing set apart in a list, above what it holds: a font’s family over its variants in M-x font-list.\\
\bottomrule
\end{xltabular}

\subsection{Your name}
\label{sec:561}
Some of what Tuatara writes carries your name: the author line of an
exported document, a note template and a snippet. As standard it uses the
full name your account has on the computer, or else the login name. To use
another, set it:

\begin{Verbatim}[breaklines,breakautoindent=false,breakindent=1em,breakafter=/-.\,;:=,commandchars=\\\{\}]
\color[HTML]{A85F00}(setf\normalcolor{} \color[HTML]{A85F00}*user-full-name*\normalcolor{} \color[HTML]{0A7AB3}"Aroha Smith"\color[HTML]{A85F00})
\end{Verbatim}

Its table is in chapter 13, with the snippets that use it.

\subsection{Snippets of your own}
\label{sec:562}
A snippet is a piece of text with places to fill in. You type its short key
and press \texttt{TAB}, and the text goes in with the cursor on the first place.
Chapter 13 tells all about the places. This section is about making a snippet
that no one else has, for something you write often.

A snippet is a file. It lives in \texttt{\textasciitilde{}/.\codebreak{}config/\codebreak{}tuatara/\codebreak{}snippets/}, in a folder
named for the mode that it is for. \texttt{fundamental-\codebreak{}mode} is for every buffer.
\texttt{org-\codebreak{}mode} is for Org notes. The file starts with three lines: the name, the
key you type, and a line with two dashes. Everything after the dashes is the
text that goes in. In the text, \texttt{\$1}, \texttt{\$2} and so on are places that \texttt{TAB} moves
through, \texttt{\$0} is where the cursor ends, and \texttt{\$\{1:\codebreak{}text\}} is a place that begins
with text you can type over. \texttt{\$AUTHOR} is your name, and \texttt{\$CURRENT\_\codebreak{}DATE},
\texttt{\$CURRENT\_\codebreak{}MONTH} and \texttt{\$CURRENT\_\codebreak{}YEAR} are today’s date.

\texttt{C-c y} (\texttt{insert-\codebreak{}snippet}) chooses one from a list by name. \texttt{M-x copy-\codebreak{}snippet}
copies one of Tuatara’s own into your folder, so that you can change it. Yours
then takes its place.

\begin{try}
\textbf{Try it, with a snippet for a letter’s greeting}
\begin{enumerate}
\item In a terminal, type these two commands.
\begin{Verbatim}[breaklines,breakautoindent=false,breakindent=1em,breakafter=/-.\,;:=,commandchars=\\\{\}]
\color[HTML]{7D2AE8}mkdir\normalcolor{} \color[HTML]{946C00}-p\normalcolor{} \color[HTML]{946C00}~/.config/tuatara/snippets/fundamental-mode
\color[HTML]{7D2AE8}printf\normalcolor{} \color[HTML]{0A7AB3}'# name: Greeting\char92{}n# key: dear\char92{}n# --\char92{}nKia ora $\char123{}1:Aroha\char125{},\char92{}n\char92{}n$0\char92{}n\char92{}nKind regards,\char92{}n$AUTHOR\char92{}n'\normalcolor{} \color[HTML]{222222}>\normalcolor{} \color[HTML]{0A7AB3}~/.config/tuatara/snippets/fundamental-mode/greeting
\end{Verbatim}
\item In Tuatara, open a new file with \texttt{C-x C-f} and the name \texttt{\textasciitilde{}/letter.\codebreak{}txt}.
   Press \texttt{i}, type \texttt{dear} and press \texttt{TAB}. The greeting goes in, with the name
   ready to type over.
\item Type \texttt{Hemi}, press \texttt{TAB}, and the cursor moves to the empty line for your
   letter. Your name is at the bottom.
\item Press \texttt{C-g}, then \texttt{u} until the text is gone. Delete the snippet file if
   you do not want it.
\end{enumerate}
\end{try}

\subsection{How lines are shown}
\label{sec:563}
Tuatara can show the same text in several ways without changing a character
of it: the line the cursor is on painted across the window, long lines
wrapped onto several rows, a line down the window at the fill column, the
column the cursor is in painted down the window, the text in a column in the
middle of the window, and the characters that draw nothing shown as their
codes. Each is a mode you switch on or off. They all work the same way.

A command with \texttt{-mode} at the end switches the thing for the buffer you are
in. That buffer keeps its choice. Most have a key under \texttt{C-x T}, for “toggle”.
The same command with \texttt{global-} in front switches it for every buffer that has
not chosen for itself, while Tuatara runs. A setting in the init file says
what every buffer has from the start. A list of modes, such as
\texttt{*soft-\codebreak{}wrap-\codebreak{}modes*}, says it for one kind of buffer, all Org notes or all code,
whatever the setting for every buffer says.

A buffer’s own choice counts first, then the list of modes, then the setting
for every buffer. Each switch says what it did on the echo row. With a count
before it, \texttt{C-u 1}, a command switches on whatever it was, and \texttt{C-u 0}
switches it off.

The commands and what they show, with the keys they have, are in the table
below. Some of them are shown in other chapters. Here they are together.

\begin{xltabular}{\textwidth}{@{}l >{\raggedright\arraybackslash\hsize=0.971\hsize}X l >{\raggedright\arraybackslash\hsize=1.029\hsize}X@{}}
\toprule
\textbf{Modal} & \textbf{Non-modal} & \textbf{Command} & \textbf{What it does}\\
\midrule
\endhead
none & \texttt{C-x T h} & \texttt{hl-\codebreak{}line-\codebreak{}mode} & Highlight the line point is on in the current buffer, or stop: every row of it, out to the window’s edge, in the :hl-line face, under the syntax colours and the selection, in the selected window.\\
none & \texttt{C-x T c} & \texttt{cursor-\codebreak{}column-\codebreak{}mode} & Highlight the column point is in, down the window, in the current buffer, or stop: in the :cursor-column face (the :hl-line face’s colour unless set apart), over the backgrounds of the text (a source block’s, the selection’s) and under the line point is on, past the end of short lines too, in the selected window.\\
none & \texttt{C-x T l} & \texttt{soft-\codebreak{}wrap-\codebreak{}mode} & Wrap the long lines of the current buffer at a word’s edge, or stop: a line longer than the window is shown on as many rows as it needs, at the fill column (C-x f), or the window’s edge when that is nearer; the text is not changed.\\
none & \texttt{C-x T v} & \texttt{visual-\codebreak{}line-\codebreak{}mode} & Mark the current buffer’s fill column (C-x f) with a line down the window, or stop: the last column a filled line may reach, the 80th for a fill column of 80, painted in the :visual-line face over the text’s backgrounds (a source block’s, the selection’s) and under the line point is on.\\
none & \texttt{C-x T o} & \texttt{centred-\codebreak{}text-\codebreak{}mode} & Show the text of the current buffer in a column of \textbf{CENTRED-TEXT-WIDTH} columns in the middle of the window, blank margins either side, or as it was: for writing and reading with nothing else about.\\
none & \texttt{C-x T i} & \texttt{invisible-\codebreak{}characters-\codebreak{}mode} & Show the characters that draw nothing in the current buffer as their codes, or stop: the byte-order mark as <FEFF>, a zero-width space as <200B>, the soft hyphen, the joiners, the word joiner, the marks that set the direction text is drawn in and the tags the same way, in the :invisible-character face, each taking the columns its code takes.\\
none & \texttt{M-x text-\codebreak{}direction-\codebreak{}mode} & \texttt{text-\codebreak{}direction-\codebreak{}mode} & Show the lines of the current buffer that hold right-to-left text (Hebrew, Arabic and the others) in the order they are read, or each line’s characters in the order they are stored.\\
none & \texttt{M-x global-\codebreak{}hl-\codebreak{}line-\codebreak{}mode} & \texttt{global-\codebreak{}hl-\codebreak{}line-\codebreak{}mode} & Highlight the line point is on, in every buffer, or stop, while the editor runs (\textbf{GLOBAL-HL-LINE-MODE} in the init file keeps it); a buffer that chose with C-x T h, and the modes \textbf{HL-LINE-MODES} names, keep theirs; with a count, on when it is positive, off when it is zero or less.\\
none & \texttt{M-x global-\codebreak{}soft-\codebreak{}wrap-\codebreak{}mode} & \texttt{global-\codebreak{}soft-\codebreak{}wrap-\codebreak{}mode} & Wrap long lines at a word’s edge, at the fill column, in every buffer, or stop, while the editor runs (\textbf{GLOBAL-SOFT-WRAP-MODE} in the init file keeps it); a buffer that chose with C-x T l, and the modes \textbf{SOFT-WRAP-MODES} names, keep theirs; with a count, on when it is positive, off when it is zero or less.\\
none & \texttt{M-x global-\codebreak{}visual-\codebreak{}line-\codebreak{}mode} & \texttt{global-\codebreak{}visual-\codebreak{}line-\codebreak{}mode} & Mark the fill column with a line down the window, in every buffer, or stop, while the editor runs (\textbf{GLOBAL-VISUAL-LINE-MODE} in the init file keeps it); a buffer that chose with C-x T v, and the modes \textbf{VISUAL-LINE-MODES} names, keep theirs; with a count, on when it is positive, off when it is zero or less.\\
none & \texttt{M-x global-\codebreak{}centred-\codebreak{}text-\codebreak{}mode} & \texttt{global-\codebreak{}centred-\codebreak{}text-\codebreak{}mode} & Show the text in a column in the middle of the window, in every buffer, or stop, while the editor runs (\textbf{GLOBAL-CENTRED-TEXT-MODE} in the init file keeps it); a buffer that chose with C-x T o, and the modes \textbf{CENTRED-TEXT-MODES} names, keep theirs; with a count, on when it is positive, off when it is zero or less.\\
none & \texttt{M-x global-\codebreak{}cursor-\codebreak{}column-\codebreak{}mode} & \texttt{global-\codebreak{}cursor-\codebreak{}column-\codebreak{}mode} & Highlight the column point is in, down the window, in every buffer, or stop, while the editor runs (\textbf{GLOBAL-CURSOR-COLUMN-MODE} in the init file keeps it); a buffer that chose with C-x T c, and the modes \textbf{CURSOR-COLUMN-MODES} names, keep theirs; with a count, on when it is positive, off when it is zero or less.\\
none & \texttt{M-x global-\codebreak{}invisible-\codebreak{}characters-\codebreak{}mode} & \texttt{global-\codebreak{}invisible-\codebreak{}characters-\codebreak{}mode} & Show the characters that draw nothing as their codes, <FEFF> for the byte-order mark, in every buffer, or stop, while the editor runs (\textbf{GLOBAL-INVISIBLE-CHARACTERS-MODE} in the init file keeps it); a buffer that chose with C-x T i, and the modes \textbf{INVISIBLE-CHARACTERS-MODES} names, keep theirs; with a count, on when it is positive, off when it is zero or less.\\
none & \texttt{M-x global-\codebreak{}text-\codebreak{}direction-\codebreak{}mode} & \texttt{global-\codebreak{}text-\codebreak{}direction-\codebreak{}mode} & Show right-to-left text in the order it is read, in every buffer, or stop, while the editor runs (\textbf{GLOBAL-TEXT-DIRECTION-MODE} in the init file keeps it); a buffer that chose with M-x text-direction-mode, and the modes \textbf{TEXT-DIRECTION-MODES} names, keep theirs; with a count, on when it is positive, off when it is zero or less.\\
\bottomrule
\end{xltabular}

Text that is written from right to left, such as Hebrew and Arabic, is
shown in the order it is read. \texttt{text-\codebreak{}direction-\codebreak{}mode} is for the rare text
where you want each line’s letters in the order they are stored instead.

The settings below say, in the init file, what every buffer has, which modes
differ, how wide a centred column is and whether a wrapped line wraps at the
fill column.

\begin{xltabular}{\textwidth}{@{}l l >{\raggedright\arraybackslash\hsize=1.000\hsize}X@{}}
\toprule
\textbf{Setting} & \textbf{Standard} & \textbf{What it changes}\\
\midrule
\endhead
\texttt{*global-\codebreak{}hl-\codebreak{}line-\codebreak{}mode*} & \texttt{t} & True to highlight the line point is on, in every buffer (T unless set); a buffer’s own M-x hl-line-mode, and the list for its major mode in \textbf{HL-LINE-MODES}, say otherwise for it.\\
\texttt{*hl-\codebreak{}line-\codebreak{}modes*} & \texttt{nil} & Major modes whose buffers highlight the line point is on, or not: a list of (MODE .\\
\texttt{*global-\codebreak{}soft-\codebreak{}wrap-\codebreak{}mode*} & \texttt{nil} & True to wrap long lines at a word’s edge, at the fill column, in every buffer (NIL unless set); a buffer’s own M-x soft-wrap-mode, and the list for its major mode in \textbf{SOFT-WRAP-MODES}, say otherwise for it.\\
\texttt{*soft-\codebreak{}wrap-\codebreak{}modes*} & \texttt{nil} & Major modes whose buffers wrap long lines at a word’s edge, at the fill column, or not: a list of (MODE .\\
\texttt{*soft-\codebreak{}wrap-\codebreak{}at-\codebreak{}fill-\codebreak{}column*} & \texttt{t} & True to wrap long lines at the buffer’s fill column (M-x set-fill-column, \textbf{FILL-COLUMN}), or at the window’s edge when that is nearer; NIL to wrap at the window’s edge always.\\
\texttt{*global-\codebreak{}visual-\codebreak{}line-\codebreak{}mode*} & \texttt{nil} & True to mark the fill column with a line down the window, in every buffer (NIL unless set); a buffer’s own M-x visual-line-mode, and the list for its major mode in \textbf{VISUAL-LINE-MODES}, say otherwise for it.\\
\texttt{*visual-\codebreak{}line-\codebreak{}modes*} & \texttt{nil} & Major modes whose buffers mark the fill column with a line down the window, or not: a list of (MODE .\\
\texttt{*global-\codebreak{}centred-\codebreak{}text-\codebreak{}mode*} & \texttt{nil} & True to show the text in a column in the middle of the window, in every buffer (NIL unless set); a buffer’s own M-x centred-text-mode, and the list for its major mode in \textbf{CENTRED-TEXT-MODES}, say otherwise for it.\\
\texttt{*centred-\codebreak{}text-\codebreak{}modes*} & \texttt{nil} & Major modes whose buffers show the text in a column in the middle of the window, or not: a list of (MODE .\\
\texttt{*centred-\codebreak{}text-\codebreak{}width*} & \texttt{80} & How many columns of text a buffer shows when its text is centred (M-x centred-text-mode): the column’s width, the rest of the window blank margins either side.\\
\texttt{*global-\codebreak{}cursor-\codebreak{}column-\codebreak{}mode*} & \texttt{nil} & True to highlight the column point is in, down the window, in every buffer (NIL unless set); a buffer’s own M-x cursor-column-mode, and the list for its major mode in \textbf{CURSOR-COLUMN-MODES}, say otherwise for it.\\
\texttt{*cursor-\codebreak{}column-\codebreak{}modes*} & \texttt{nil} & Major modes whose buffers highlight the column point is in, down the window, or not: a list of (MODE .\\
\texttt{*global-\codebreak{}invisible-\codebreak{}characters-\codebreak{}mode*} & \texttt{nil} & True to show the characters that draw nothing as their codes, <FEFF> for the byte-order mark, in every buffer (NIL unless set); a buffer’s own M-x invisible-characters-mode, and the list for its major mode in \textbf{INVISIBLE-CHARACTERS-MODES}, say otherwise for it.\\
\texttt{*invisible-\codebreak{}characters-\codebreak{}modes*} & \texttt{nil} & Major modes whose buffers show the characters that draw nothing as their codes, <FEFF> for the byte-order mark, or not: a list of (MODE .\\
\texttt{*global-\codebreak{}text-\codebreak{}direction-\codebreak{}mode*} & \texttt{:auto} & How lines holding right-to-left text are shown in buffers that have not chosen, and whose major mode \textbf{TEXT-DIRECTION-MODES} says nothing of.\\
\texttt{*text-\codebreak{}direction-\codebreak{}modes*} & \texttt{nil} & Major modes whose buffers show right-to-left text in the order it is read, or not: a list of (MODE .\\
\bottomrule
\end{xltabular}

The faces paint the line, the column and the characters. The first three
are the same colour unless you set them apart.

\begin{xltabular}{\textwidth}{@{}l >{\raggedright\arraybackslash\hsize=1.000\hsize}X@{}}
\toprule
\textbf{Face} & \textbf{What it paints}\\
\midrule
\endhead
\texttt{:hl-\codebreak{}line} & The line point is on, in the selected window, out to its edge (C-x T h, M-x global-hl-line-mode).\\
\texttt{:cursor-\codebreak{}column} & The column point is in, down the selected window (C-x T c, M-x global-cursor-column-mode): the same as :hl-line unless set apart, (set-face :cursor-column :bg \ldots{}) in the init file.\\
\texttt{:soft-\codebreak{}wrap-\codebreak{}row} & The row point is on, where long lines wrap (C-x T l), painted over the line’s :hl-line: the same as it unless set apart, (set-face :soft-wrap-row :bg \ldots{}) in the init file.\\
\texttt{:visual-\codebreak{}line} & The line down the window at the fill column (C-x T v): its background, over a source block’s or the selection’s, under the line point is on (:hl-line).\\
\texttt{:invisible-\codebreak{}character} & A character that draws nothing, shown as its code, <FEFF> for the byte-order mark, while its buffer shows such characters (C-x T i), and a character that sets the direction text is drawn in, shown as ? otherwise: as escape sequences are, unless set apart.\\
\texttt{:current-\codebreak{}line-\codebreak{}number} & The current line’s number in the gutter.\\
\texttt{:raw-\codebreak{}byte} & Octets that are not UTF-8, each shown as \texttt{\textbackslash{}xNN}.\\
\texttt{:right-\codebreak{}to-\codebreak{}left} & Text written from right to left, shown in the order it is read: nothing of its own unless set apart; the frontends find the runs whose letters they join by it.\\
\texttt{:stored-\codebreak{}order} & Right-to-left letters shown in the order they are stored, in a stretch of a line too long, or too far past an embedding or override, for the order they are read in to be worked out.\\
\bottomrule
\end{xltabular}

\subsection{The calendar}
\label{sec:565}
\texttt{C-x n a} shows the calendar in a window below the one you are in: the month
of today and the month before and after it, with the cursor on today.
\texttt{C-x n A} shows it in the window you are in instead, in place of its buffer.
\texttt{q} closes it and brings your buffer back.

The arrow keys, or \texttt{h} \texttt{l} for a day back or on and \texttt{k} \texttt{j} for a week back or
on, move through the days. \texttt{H} and \texttt{L} show the month before and after, and \texttt{J}
and \texttt{K} move the cursor a month. \texttt{.} goes to today, and \texttt{g l} goes back to where
you were before you moved. \texttt{g d} asks for a date and goes to it. You can
answer with \texttt{2026-12-24}, with \texttt{tomorrow}, \texttt{fri} or \texttt{+3d}, or as any date is
asked for in Tuatara. \texttt{o} shows the months around one you name, such as \texttt{jan 27}
for January 2027.

Days that have something on them are marked. An agenda entry (chapter 14) marks
its day, and a journal entry (the next section) underlines it. \texttt{RET} on a day
opens what it is marked for. When a day has both, Tuatara asks which, and each
answer is one key. \texttt{n} opens the journal entry for the day at the cursor, or
makes it. \texttt{m} marks the days that have agenda entries and \texttt{u} takes the marks
off. \texttt{p d} says which day of the year the cursor is on, and \texttt{p c} says the week
and the ISO date.

Three settings change how the calendar looks. By default it shows three months
and weeks begin on Monday. \texttt{*calendar-\codebreak{}months*} set to 12 shows a whole year.
\texttt{*calendar-\codebreak{}week-\codebreak{}start-\codebreak{}day*} set to 0 begins weeks on Sunday.
\texttt{*calendar-\codebreak{}months-\codebreak{}per-\codebreak{}row*} says how many months stand side by side. It is as
many as fit when it is \texttt{nil}.

\begin{Verbatim}[breaklines,breakautoindent=false,breakindent=1em,breakafter=/-.\,;:=,commandchars=\\\{\}]
\color[HTML]{A85F00}(setf\normalcolor{} \color[HTML]{A85F00}*calendar-week-start-day*\normalcolor{} \color[HTML]{0A7AB3}0\color[HTML]{A85F00})
\color[HTML]{A85F00}(setf\normalcolor{} \color[HTML]{A85F00}*calendar-months*\normalcolor{} \color[HTML]{0A7AB3}12\color[HTML]{A85F00})
\color[HTML]{A85F00}(setf\normalcolor{} \color[HTML]{A85F00}*calendar-months-per-row*\normalcolor{} \color[HTML]{0A7AB3}3\color[HTML]{A85F00})
\end{Verbatim}

\subsubsection{Commands, settings and faces for the calendar}
\label{sec:564}
\begin{xltabular}{\textwidth}{@{}l >{\raggedright\arraybackslash\hsize=0.848\hsize}X l >{\raggedright\arraybackslash\hsize=1.152\hsize}X@{}}
\toprule
\textbf{Modal} & \textbf{Non-modal} & \textbf{Command} & \textbf{What it does}\\
\midrule
\endhead
none & \texttt{C-f}, \texttt{<right>} & \texttt{calendar-\codebreak{}forward-\codebreak{}day} & Move point a day on in the calendar; with a count, that many.\\
none & \texttt{C-b}, \texttt{<left>} & \texttt{calendar-\codebreak{}backward-\codebreak{}day} & Move point a day back in the calendar; with a count, that many.\\
none & \texttt{C-n}, \texttt{<down>} & \texttt{calendar-\codebreak{}forward-\codebreak{}week} & Move point a week on in the calendar; with a count, that many.\\
none & \texttt{C-p}, \texttt{<up>} & \texttt{calendar-\codebreak{}backward-\codebreak{}week} & Move point a week back in the calendar; with a count, that many.\\
none & \texttt{M-\}} & \texttt{calendar-\codebreak{}forward-\codebreak{}month} & Move point a month on in the calendar, to the same day of the month (or its last); with a count, that many.\\
none & \texttt{M-\{} & \texttt{calendar-\codebreak{}backward-\codebreak{}month} & Move point a month back in the calendar, to the same day of the month (or its last); with a count, that many.\\
none & \texttt{C-x ]} & \texttt{calendar-\codebreak{}forward-\codebreak{}year} & Move point a year on in the calendar; with a count, that many.\\
none & \texttt{C-x [} & \texttt{calendar-\codebreak{}backward-\codebreak{}year} & Move point a year back in the calendar; with a count, that many.\\
none & \texttt{.}, \texttt{C-c .} & \texttt{calendar-\codebreak{}goto-\codebreak{}today} & Move point to today, the months around it shown.\\
none & \texttt{g d} & \texttt{calendar-\codebreak{}goto-\codebreak{}date} & Move point to a date asked for: 2026-12-24, 2026-12, or as any date asked for is read (tomorrow, fri, +3d, dec 24).\\
none & \texttt{M-x calendar-\codebreak{}jump-\codebreak{}back} & \texttt{calendar-\codebreak{}jump-\codebreak{}back} & Go back to the day point was on before it last moved; again, and back to the day after once more.\\
none & \texttt{o}, \texttt{C-c o} & \texttt{calendar-\codebreak{}other-\codebreak{}month} & Show the months around a month asked for, point on its first day: jan ’27 (or jan 27, january 2027) for January 2027, jan for this year’s, 2027-03, or a date (2027-03-14, fri, +3m) for its month.\\
none & \texttt{RET} & \texttt{calendar-\codebreak{}open-\codebreak{}day} & Open what the day at point is marked for: its agenda, its journal entry, whichever it has (\textbf{CALENDAR-DAY-OPENERS}); with more than one, asked which, each by the key that opens it alone.\\
none & \texttt{n}, \texttt{C-c n} & \texttt{calendar-\codebreak{}journal-\codebreak{}entry} & Open the journal’s entry for the calendar’s day at point, or make it: the calendar leaves its window, the entry shown where you were.\\
none & \texttt{m}, \texttt{C-c m} & \texttt{calendar-\codebreak{}mark-\codebreak{}days} & Mark the days with agenda entries in the calendar (as standard, with \textbf{calendar-mark-days}).\\
none & \texttt{u}, \texttt{C-c u} & \texttt{calendar-\codebreak{}unmark} & Take the marks off the calendar’s days, until m puts them back.\\
none & \texttt{C-a} & \texttt{calendar-\codebreak{}beginning-\codebreak{}of-\codebreak{}week} & Move point to the start of the week it is in (\textbf{calendar-week-start-day}); on it, a week back.\\
none & \texttt{C-e} & \texttt{calendar-\codebreak{}end-\codebreak{}of-\codebreak{}week} & Move point to the end of the week it is in; on it, a week on.\\
none & \texttt{M-a} & \texttt{calendar-\codebreak{}beginning-\codebreak{}of-\codebreak{}month} & Move point to the first day of its month.\\
none & \texttt{M-e} & \texttt{calendar-\codebreak{}end-\codebreak{}of-\codebreak{}month} & Move point to the last day of its month.\\
none & \texttt{M-<} & \texttt{calendar-\codebreak{}beginning-\codebreak{}of-\codebreak{}year} & Move point to the first day of its year.\\
none & \texttt{M->} & \texttt{calendar-\codebreak{}end-\codebreak{}of-\codebreak{}year} & Move point to the last day of its year.\\
none & \texttt{>}, \texttt{C-c >}, \texttt{C-x >} & \texttt{calendar-\codebreak{}scroll-\codebreak{}left} & Show the next month, the months moving left; with a count, that many.\\
none & \texttt{<}, \texttt{C-c <}, \texttt{C-x <} & \texttt{calendar-\codebreak{}scroll-\codebreak{}right} & Show the month before, the months moving right; with a count, that many.\\
none & \texttt{C-v}, \texttt{<next>} & \texttt{calendar-\codebreak{}scroll-\codebreak{}left-\codebreak{}three-\codebreak{}months} & Show the months after those shown, as many as are shown (\textbf{calendar-months}, three as standard); with a count, that many times as many.\\
none & \texttt{M-v}, \texttt{<prior>} & \texttt{calendar-\codebreak{}scroll-\codebreak{}right-\codebreak{}three-\codebreak{}months} & Show the months before those shown, as many as are shown (\textbf{calendar-months}, three as standard); with a count, that many times as many.\\
none & \texttt{p d} & \texttt{calendar-\codebreak{}print-\codebreak{}day-\codebreak{}of-\codebreak{}year} & Say which day of its year the day at point is.\\
none & \texttt{p c} & \texttt{calendar-\codebreak{}iso-\codebreak{}print-\codebreak{}date} & Say the ISO date of the day at point: its day of the week, its week and its year.\\
none & \texttt{C-x n A} & \texttt{calendar-\codebreak{}here} & Show the calendar in the window you are in, in place of its buffer, point on today; q brings that buffer back.\\
\bottomrule
\end{xltabular}

\begin{xltabular}{\textwidth}{@{}l l >{\raggedright\arraybackslash\hsize=1.000\hsize}X@{}}
\toprule
\textbf{Setting} & \textbf{Standard} & \textbf{What it changes}\\
\midrule
\endhead
\texttt{*calendar-\codebreak{}months*} & \texttt{3} & How many months the calendar shows: three as standard, the month point is in among them, as many before it as after (one more after for an even number); twelve for a year; NIL for as many as fit side by side in its window (\textbf{CALENDAR-MONTHS-PER-ROW}, or the window’s width).\\
\texttt{*calendar-\codebreak{}months-\codebreak{}per-\codebreak{}row*} & \texttt{nil} & How many months the calendar shows side by side, the rest of \textbf{CALENDAR-MONTHS} in rows of them below, a blank line between: NIL, as standard, for as many as its window’s width holds (three in 80 columns, four in 105), a number for that many.\\
\texttt{*calendar-\codebreak{}week-\codebreak{}start-\codebreak{}day*} & \texttt{1} & The day a week begins on in the calendar: 0 Sunday, 1 Monday (as standard), up to 6 Saturday.\\
\texttt{*calendar-\codebreak{}mark-\codebreak{}days*} & \texttt{t} & True for the calendar to mark the days \textbf{CALENDAR-MARK-FUNCTIONS} name (the days with agenda entries) as soon as it shows them; NIL to wait for M-x calendar-mark-days.\\
\bottomrule
\end{xltabular}

\begin{xltabular}{\textwidth}{@{}l >{\raggedright\arraybackslash\hsize=1.000\hsize}X@{}}
\toprule
\textbf{Face} & \textbf{What it paints}\\
\midrule
\endhead
\texttt{:calendar-\codebreak{}journal} & Days in the calendar with a journal entry, underlined and bold, in the text’s colour, or the colour of agenda entries on a day that has both.\\
\texttt{:calendar-\codebreak{}marked} & Days the calendar marks: those with agenda entries.\\
\texttt{:calendar-\codebreak{}month-\codebreak{}header} & The months’ titles in the calendar (October 2026).\\
\texttt{:calendar-\codebreak{}selected} & The day the cursor is on in the calendar, both its digits on the cursor’s colour as the cursor’s character is (set as the cursor’s colour changes).\\
\texttt{:calendar-\codebreak{}today} & Today in the calendar, on a bright ground so that it stands out.\\
\texttt{:calendar-\codebreak{}weekday-\codebreak{}header} & The names of the weekdays over the calendar’s columns (Mo Tu).\\
\texttt{:calendar-\codebreak{}weekend-\codebreak{}header} & The names of Saturday and Sunday over the calendar’s columns.\\
\bottomrule
\end{xltabular}

\subsection{The journal}
\label{sec:566}
A journal is a note for each day: what you did, what you thought, what to
remember. \texttt{C-x n j} opens today’s entry, or makes it when there is none. A new
entry is a file with a formatted name (chapter 9), in a folder called
\texttt{journal} inside the folder of your notes, with the day as its title, such as
\texttt{Friday 20 March 2026}, and the word \texttt{journal} among its keywords. It is saved
at once, so there is a file to find later. \texttt{M-x journal-\codebreak{}new-\codebreak{}entry} always
makes a new one, even when there is one for today. With \texttt{C-u} before it, it
asks for a date.

When there is more than one entry for the day, Tuatara asks which you want.
In the calendar, an entry marks its day, and \texttt{n} opens or makes the entry for
the day at the cursor. You can look back through the month and open any
day.

The journal is a set of ordinary notes, so you can search them, link to
them and export them like any others. Everything about how they are named
and written is a setting.

The folder is \texttt{*journal-\codebreak{}directory*}. The word that marks an entry, so that
entries can be told from your other notes, is \texttt{*journal-\codebreak{}keyword*}. An entry
lasts a day as standard. \texttt{*journal-\codebreak{}interval*} makes it a week, a month or a
year. The kind of file is \texttt{*journal-\codebreak{}file-\codebreak{}type*}, one of \texttt{:org}, \texttt{:markdown} or
\texttt{:text}. \texttt{*journal-\codebreak{}title-\codebreak{}format*} says what the title looks like. It may add
the time, or ask each time. \texttt{*journal-\codebreak{}file-\codebreak{}name*}, when set, gives entries plain
names such as \texttt{2026-03-20.\codebreak{}org} in place of formatted ones.
\texttt{*journal-\codebreak{}org-\codebreak{}header*} is the text an Org entry starts with.

\begin{Verbatim}[breaklines,breakautoindent=false,breakindent=1em,breakafter=/-.\,;:=,commandchars=\\\{\}]
\color[HTML]{A85F00}(setf\normalcolor{} \color[HTML]{A85F00}*journal-directory*\normalcolor{} \color[HTML]{0A7AB3}"\color[HTML]{222222}~/journal/\color[HTML]{0A7AB3}"\color[HTML]{A85F00})
\color[HTML]{A85F00}(setf\normalcolor{} \color[HTML]{A85F00}*journal-interval*\normalcolor{} \color[HTML]{946C00}:\normalcolor{}weekly\color[HTML]{A85F00})
\color[HTML]{A85F00}(setf\normalcolor{} \color[HTML]{A85F00}*journal-file-type*\normalcolor{} \color[HTML]{946C00}:\normalcolor{}markdown\color[HTML]{A85F00})
\end{Verbatim}

The three lines put the entries in \texttt{\textasciitilde{}/journal/}, one for each week, as
Markdown files.

\begin{xltabular}{\textwidth}{@{}l >{\raggedright\arraybackslash\hsize=0.684\hsize}X l >{\raggedright\arraybackslash\hsize=1.316\hsize}X@{}}
\toprule
\textbf{Modal} & \textbf{Non-modal} & \textbf{Command} & \textbf{What it does}\\
\midrule
\endhead
none & \texttt{C-x n j} & \texttt{journal-\codebreak{}new-\codebreak{}or-\codebreak{}existing-\codebreak{}entry} & Open today’s journal entry, or make it when there is none: in \textbf{journal-directory}, named and titled as \textbf{journal-file-name} and \textbf{journal-title-format} say, saved at once; with several for today, the one chosen.\\
none & \texttt{M-x journal-\codebreak{}new-\codebreak{}entry} & \texttt{journal-\codebreak{}new-\codebreak{}entry} & Make a new journal entry for today, whatever entries there are; with C-u, for a date asked for.\\
\bottomrule
\end{xltabular}

\begin{xltabular}{\textwidth}{@{}l >{\raggedright\arraybackslash\hsize=0.538\hsize}X >{\raggedright\arraybackslash\hsize=1.462\hsize}X@{}}
\toprule
\textbf{Setting} & \textbf{Standard} & \textbf{What it changes}\\
\midrule
\endhead
\texttt{*journal-\codebreak{}directory*} & \texttt{nil} & The directory of your journal’s entries: NIL, as standard, for journal/ in \textbf{FORMATTED-FILE-DIRECTORY} (such as \textasciitilde{}/notes/journal/), or a directory of its own, such as “\textasciitilde{}/journal/“.\\
\texttt{*journal-\codebreak{}keyword*} & \texttt{"journal"} & The keyword in the formatted name of each journal entry, by which entries are told from your other files: journal, as standard.\\
\texttt{*journal-\codebreak{}title-\codebreak{}format*} & \texttt{:day-\codebreak{}date-\codebreak{}month-\codebreak{}year} & The title of a new journal entry: :day (Friday), :day-date-month-year (Friday 20 March 2026, as standard), :day-date-month-year-24h (Friday 20 March 2026 13:39), :day-date-month-year-12h (Friday 20 March 2026 01:39 PM), a pattern in the \% specifiers of FORMAT-TIME (“Week \%V, \%A”), ”” for none, or NIL to be asked each time.\\
\texttt{*journal-\codebreak{}interval*} & \texttt{:daily} & How long a journal entry lasts: :daily, :weekly (the ISO week, from Monday), :monthly or :yearly.\\
\texttt{*journal-\codebreak{}file-\codebreak{}type*} & \texttt{nil} & The kind of file a new journal entry is: :org, :text or :markdown; NIL for \textbf{FORMATTED-FILE-TYPE}.\\
\texttt{*journal-\codebreak{}file-\codebreak{}name*} & \texttt{nil} & How journal entries are named.\\
\texttt{*journal-\codebreak{}org-\codebreak{}header*} & \texttt{"\#+title:      \{title\} \#+date:       \{da...} & The header a new journal entry in Org begins with, its name formatted: title, date, keywords and identifier lines, \{title\}, \{date\}, \{keywords\} and \{identifier\} standing for the entry’s; NIL for the Org header of your other files with formatted names (\textbf{FORMATTED-FILE-TYPES}).\\
\bottomrule
\end{xltabular}

\subsection{Odds and ends}
\label{sec:573}
The rest of the names that belong to this chapter are small. They are
grouped by what they do.

\subsubsection{Changing case, and what a character is}
\label{sec:567}
Three commands change the case of a word or a selection. They act on the
selection when there is one, and on the next word when there is not.
\texttt{capitalize-\codebreak{}region} is the same for a selection alone.
\texttt{delete-\codebreak{}trailing-\codebreak{}whitespace} takes the spaces and tabs from the ends of lines.
It works on the selection if there is one and on the whole buffer if not. \texttt{C-x
=} says which character the cursor is on, with its code in three bases.

\begin{xltabular}{\textwidth}{@{}l >{\raggedright\arraybackslash\hsize=0.829\hsize}X l >{\raggedright\arraybackslash\hsize=1.171\hsize}X@{}}
\toprule
\textbf{Modal} & \textbf{Non-modal} & \textbf{Command} & \textbf{What it does}\\
\midrule
\endhead
none & \texttt{M-x upcase-\codebreak{}dwim} & \texttt{upcase-\codebreak{}dwim} & Put the region in upper case, or with none the next word, as upcase-word does.\\
none & \texttt{M-x downcase-\codebreak{}dwim} & \texttt{downcase-\codebreak{}dwim} & Put the region in lower case, or with none the next word, as downcase-word does.\\
none & \texttt{M-x capitalize-\codebreak{}dwim} & \texttt{capitalize-\codebreak{}dwim} & Capitalise the region, or with none the next word, as capitalize-word does.\\
none & \texttt{M-x capitalize-\codebreak{}region} & \texttt{capitalize-\codebreak{}region} & Capitalise each word of the region.\\
none & \texttt{M-x delete-\codebreak{}trailing-\codebreak{}whitespace} & \texttt{delete-\codebreak{}trailing-\codebreak{}whitespace} & Delete the spaces and tabs at the end of every line, of the selection when there is one, else of the whole buffer (of the part shown of a narrowed one), and then the blank lines at its end too (\textbf{DELETE-TRAILING-LINES}).\\
none & \texttt{C-x =} & \texttt{what-\codebreak{}cursor-\codebreak{}position} & Say in the echo area which character is at point, with its code in decimal, octal and hexadecimal, and where point is: its place counted in characters from 1, of how many there are and as a share of them, its octet and its column.\\
none & \texttt{C-x n e} & \texttt{new-\codebreak{}empty-\codebreak{}buffer} & Show a new empty buffer, untitled (untitled<2> beside one already there), for text that has no file yet: C-x C-s asks where to save it.\\
\bottomrule
\end{xltabular}

\subsubsection{Buttons and the mouse}
\label{sec:568}
A button is a piece of text that does something when you press \texttt{RET} on it:
a link in a help page, a name in a list. \texttt{TAB} goes to the next one in a
buffer and \texttt{S-TAB} goes back. Moving past the last goes round to the first.
\texttt{push-\codebreak{}button} does what the button at the cursor does. It is on \texttt{RET} and \texttt{l},
and on \texttt{C-c l}. Links in Markdown are followed with a command of their own.

The mouse works as in other programs. A click puts the cursor, a drag
selects, two clicks select a word and three a line. Letting go of the
button makes what you selected the desktop’s primary selection, which the
middle button pastes in Tuatara and in other programs.

\begin{xltabular}{\textwidth}{@{}l l l >{\raggedright\arraybackslash\hsize=1.000\hsize}X@{}}
\toprule
\textbf{Modal} & \textbf{Non-modal} & \textbf{Command} & \textbf{What it does}\\
\midrule
\endhead
none & \texttt{TAB} & \texttt{forward-\codebreak{}button} & Move point to the next button, from the last round to the first.\\
none & \texttt{S-TAB} & \texttt{backward-\codebreak{}button} & Move point to the button before, from the first round to the last.\\
none & \texttt{l}, \texttt{RET}, \texttt{C-c l} & \texttt{push-\codebreak{}button} & Do what the button at point does: open what it names.\\
none & \texttt{<down-\codebreak{}mouse-1>} & \texttt{push-\codebreak{}button-\codebreak{}at-\codebreak{}click} & Put point where the first mouse button went down, as <down-mouse-1> does, and do what a button there does.\\
none & \texttt{C-c C-o} & \texttt{markdown-\codebreak{}follow-\codebreak{}link-\codebreak{}at-\codebreak{}point} & Follow the Markdown link point is on (C-c C-o, and RET in NORMAL state): to the heading a table of contents’s link (\#an-anchor) leads to, a web or mail address in the desktop’s program after asking, or a file, visited.\\
none & \texttt{<drag-\codebreak{}mouse-1>} & \texttt{mouse-\codebreak{}drag-\codebreak{}region} & Select from where the first mouse button went down to where the mouse is now, point following it, in the window it went down in; held above or below that window’s text, it scrolls a line at a time.\\
none & \texttt{<up-\codebreak{}mouse-1>} & \texttt{mouse-\codebreak{}release} & The first mouse button let go: a drag ends there, what it selected staying selected and becoming the desktop’s primary selection, which the middle button pastes in other programs.\\
none & \texttt{<double-\codebreak{}mouse-1>} & \texttt{mouse-\codebreak{}select-\codebreak{}word} & Select the word the first mouse button went down on twice, point at its end: the run of letters, digits and \_ there, or of spaces, or of other signs; on a bracket, the bracketed group.\\
none & \texttt{<triple-\codebreak{}mouse-1>} & \texttt{mouse-\codebreak{}select-\codebreak{}line} & Select the line the first mouse button went down on three times, its line break too.\\
\bottomrule
\end{xltabular}

\subsubsection{Modes for other kinds of text}
\label{sec:569}
Tuatara picks a mode for a file from its name and what is in it. When it
picks wrong, or you want another, run the mode by name: \texttt{M-x markdown-\codebreak{}mode}
switches the buffer to Markdown. The commands below are the modes of
written languages, data and system files that no other chapter lists, and a
few for note-taking formats such as todo.txt and Ledger. A mode that is on
as soon as you open a file has a name ending in \texttt{-mode} like these.

\begin{xltabular}{\textwidth}{@{}l >{\raggedright\arraybackslash\hsize=0.818\hsize}X l >{\raggedright\arraybackslash\hsize=1.182\hsize}X@{}}
\toprule
\textbf{Modal} & \textbf{Non-modal} & \textbf{Command} & \textbf{What it does}\\
\midrule
\endhead
none & \texttt{M-x beancount-\codebreak{}mode} & \texttt{beancount-\codebreak{}mode} & Switch the current buffer to beancount-mode.\\
none & \texttt{M-x bibtex-\codebreak{}mode} & \texttt{bibtex-\codebreak{}mode} & Switch the current buffer to bibtex-mode.\\
none & \texttt{M-x djot-\codebreak{}mode} & \texttt{djot-\codebreak{}mode} & Switch the current buffer to djot-mode.\\
none & \texttt{M-x eml-\codebreak{}mode} & \texttt{eml-\codebreak{}mode} & Switch the current buffer to eml-mode.\\
none & \texttt{M-x etc-\codebreak{}passwd-\codebreak{}generic-\codebreak{}mode} & \texttt{etc-\codebreak{}passwd-\codebreak{}generic-\codebreak{}mode} & Switch the current buffer to etc-passwd-generic-mode.\\
none & \texttt{M-x gemini-\codebreak{}mode} & \texttt{gemini-\codebreak{}mode} & Switch the current buffer to gemini-mode.\\
none & \texttt{M-x ghostty-\codebreak{}mode} & \texttt{ghostty-\codebreak{}mode} & Switch the current buffer to ghostty-mode.\\
none & \texttt{M-x hosts-\codebreak{}generic-\codebreak{}mode} & \texttt{hosts-\codebreak{}generic-\codebreak{}mode} & Switch the current buffer to hosts-generic-mode.\\
none & \texttt{M-x jj-\codebreak{}description-\codebreak{}mode} & \texttt{jj-\codebreak{}description-\codebreak{}mode} & Switch the current buffer to jj-description-mode.\\
none & \texttt{M-x klog-\codebreak{}mode} & \texttt{klog-\codebreak{}mode} & Switch the current buffer to klog-mode.\\
none & \texttt{M-x latex-\codebreak{}mode} & \texttt{latex-\codebreak{}mode} & Switch the current buffer to latex-mode.\\
none & \texttt{M-x ldif-\codebreak{}mode} & \texttt{ldif-\codebreak{}mode} & Switch the current buffer to ldif-mode.\\
none & \texttt{M-x ledger-\codebreak{}mode} & \texttt{ledger-\codebreak{}mode} & Switch the current buffer to ledger-mode.\\
none & \texttt{M-x markdoc-\codebreak{}mode} & \texttt{markdoc-\codebreak{}mode} & Switch the current buffer to markdoc-mode.\\
none & \texttt{M-x markdown-\codebreak{}mode} & \texttt{markdown-\codebreak{}mode} & Switch the current buffer to markdown-mode.\\
none & \texttt{M-x mediawiki-\codebreak{}mode} & \texttt{mediawiki-\codebreak{}mode} & Switch the current buffer to mediawiki-mode.\\
none & \texttt{M-x pem-\codebreak{}mode} & \texttt{pem-\codebreak{}mode} & Switch the current buffer to pem-mode.\\
none & \texttt{M-x po-\codebreak{}mode} & \texttt{po-\codebreak{}mode} & Switch the current buffer to po-mode.\\
none & \texttt{M-x pod-\codebreak{}mode} & \texttt{pod-\codebreak{}mode} & Switch the current buffer to pod-mode.\\
none & \texttt{M-x poly-\codebreak{}markdown+r-\codebreak{}mode} & \texttt{poly-\codebreak{}markdown+r-\codebreak{}mode} & Switch the current buffer to poly-markdown+r-mode.\\
none & \texttt{M-x quarto-\codebreak{}mode} & \texttt{quarto-\codebreak{}mode} & Switch the current buffer to quarto-mode.\\
none & \texttt{M-x robots-\codebreak{}txt-\codebreak{}mode} & \texttt{robots-\codebreak{}txt-\codebreak{}mode} & Switch the current buffer to robots-txt-mode.\\
none & \texttt{M-x rst-\codebreak{}mode} & \texttt{rst-\codebreak{}mode} & Switch the current buffer to rst-mode.\\
none & \texttt{M-x strace-\codebreak{}mode} & \texttt{strace-\codebreak{}mode} & Switch the current buffer to strace-mode.\\
none & \texttt{M-x task-\codebreak{}mode} & \texttt{task-\codebreak{}mode} & Switch the current buffer to task-mode.\\
none & \texttt{M-x todotxt-\codebreak{}mode} & \texttt{todotxt-\codebreak{}mode} & Switch the current buffer to todotxt-mode.\\
none & \texttt{M-x xit-\codebreak{}mode} & \texttt{xit-\codebreak{}mode} & Switch the current buffer to xit-mode.\\
\bottomrule
\end{xltabular}

Grammars are what colour and parse a language. Tuatara turns on those of the
common languages by itself. For another, a file asks whether to turn its
grammar on, and \texttt{M-x turn-\codebreak{}on-\codebreak{}grammar} does it later. \texttt{M-x turn-\codebreak{}off-\codebreak{}grammar}
stops a language being coloured or parsed, in later sessions too.

\begin{xltabular}{\textwidth}{@{}l l l >{\raggedright\arraybackslash\hsize=1.000\hsize}X@{}}
\toprule
\textbf{Modal} & \textbf{Non-modal} & \textbf{Command} & \textbf{What it does}\\
\midrule
\endhead
none & \texttt{M-x turn-\codebreak{}on-\codebreak{}grammar} & \texttt{turn-\codebreak{}on-\codebreak{}grammar} & Turn on the grammar of a language, asked for (the current buffer’s first when its grammar is off): its files are coloured and parsed from now on, and text in it inside other files, in later sessions too.\\
none & \texttt{M-x turn-\codebreak{}off-\codebreak{}grammar} & \texttt{turn-\codebreak{}off-\codebreak{}grammar} & Turn off the grammar of a language, asked for (the current buffer’s first): its files are no longer coloured or parsed, nor text in it inside other files, in later sessions too, and no file asks to turn it on again; M-x turn-on-grammar turns it on again.\\
\bottomrule
\end{xltabular}

\subsubsection{Lists and the cursor’s popups}
\label{sec:570}
When Tuatara lists completions beside the cursor with \texttt{C-M-i}, the faces
below say how the list is drawn. In a buffer you can also have the entries
of a directory listed as you type a path, after the slash.

\begin{xltabular}{\textwidth}{@{}l l >{\raggedright\arraybackslash\hsize=1.000\hsize}X@{}}
\toprule
\textbf{Setting} & \textbf{Standard} & \textbf{What it changes}\\
\midrule
\endhead
\texttt{*complete-\codebreak{}paths*} & \texttt{t} & True to list the entries of a directory beside the cursor when a slash is typed in a buffer after its name: a path is a word with a slash in it (\emph{etc}, \textasciitilde{}/notes/, ./src/, or a file: link’s), one not starting with / or \textasciitilde{} taken from the directory the buffer stands in.\\
\texttt{*complete-\codebreak{}paths-\codebreak{}modes*} & \texttt{nil} & Major modes whose buffers list a path’s directory as it is typed, or not, whatever \textbf{COMPLETE-PATHS} says: a list of (MODE .\\
\texttt{*relative-\codebreak{}age-\codebreak{}limit*} & \texttt{1209600} & Seconds within which the last change of a file among completion candidates is shown as an age (5 min ago, 3 days ago); older ones show the date, and with 0 every one does.\\
\texttt{*preview-\codebreak{}buffers*} & \texttt{nil} & True to show, while you choose a buffer with M-x switch-to-buffer (and M-x switch-to-buffer-other-window, and every other question for a buffer), the buffer selected in the window you are in, what it showed coming back when the question ends.\\
\bottomrule
\end{xltabular}

\begin{xltabular}{\textwidth}{@{}l >{\raggedright\arraybackslash\hsize=1.000\hsize}X@{}}
\toprule
\textbf{Face} & \textbf{What it paints}\\
\midrule
\endhead
\texttt{:popup} & The list of completions beside the cursor (C-M-i), and the documentation beside it: their text and ground.\\
\texttt{:popup-\codebreak{}annotation} & What is said of a completion in the list beside the cursor: what it is.\\
\texttt{:popup-\codebreak{}bar} & The scroll bar of the list of completions beside the cursor.\\
\texttt{:popup-\codebreak{}current} & The chosen completion in the list beside the cursor.\\
\texttt{:popup-\codebreak{}heading} & The first line of the documentation beside the list of completions: the name and its parameters.\\
\texttt{:completion-\codebreak{}mode-\codebreak{}directory} & The d of a directory’s permissions beside it among completion candidates.\\
\texttt{:completion-\codebreak{}mode-\codebreak{}link} & The l of a symbolic link’s permissions beside it among completion candidates.\\
\texttt{:input-\codebreak{}method-\codebreak{}pending} & Keys typed under an input method that wait for the next key to settle them (an a under swedish-postfix, waiting to see whether an a, an e or another key comes), underlined.\\
\bottomrule
\end{xltabular}

Reports and lists, such as the health report of \texttt{M-x check-\codebreak{}health}, have faces
for the name of a field and for what cannot be changed:

\begin{xltabular}{\textwidth}{@{}l >{\raggedright\arraybackslash\hsize=1.000\hsize}X@{}}
\toprule
\textbf{Face} & \textbf{What it paints}\\
\midrule
\endhead
\texttt{:list-\codebreak{}label} & The name of a field before its value, in a report or a message’s details (Subject:, Verdict).\\
\texttt{:list-\codebreak{}read-\codebreak{}only} & Something that cannot be changed, in a list: the names of read-only buffers in the buffer list (C-x C-b).\\
\bottomrule
\end{xltabular}

\subsubsection{Limits and safety}
\label{sec:571}
Four settings are limits that keep a mistake from filling your computer. Each
has a sensible standard, and a reason to change it is rare.

\begin{itemize}
\item \texttt{*insert-\codebreak{}file-\codebreak{}limit*} is how much of another file \texttt{M-x
  insert-\codebreak{}formatted-\codebreak{}note-\codebreak{}template} reads. A larger file is refused.
\item \texttt{*program-\codebreak{}output-\codebreak{}limit*} is how much of what a program writes Tuatara
  takes. A program that writes more is stopped.
\item \texttt{*wipe-\codebreak{}passes*} is how many times a file that you wipe is overwritten with
  random octets before it is deleted. Once is enough for a drive that writes
  in place.
\item \texttt{*markdown-\codebreak{}toc-\codebreak{}on-\codebreak{}save*} makes the table of contents of a Markdown file
  again each time you save, under the heading tagged \texttt{:TOC:}.
\end{itemize}

\begin{xltabular}{\textwidth}{@{}l l >{\raggedright\arraybackslash\hsize=1.000\hsize}X@{}}
\toprule
\textbf{Setting} & \textbf{Standard} & \textbf{What it changes}\\
\midrule
\endhead
\texttt{*insert-\codebreak{}file-\codebreak{}limit*} & \texttt{16777216} & The most octets M-x insert-formatted-note-template reads from a file that is no note template; a larger file is refused, and said so.\\
\texttt{*program-\codebreak{}output-\codebreak{}limit*} & \texttt{67108864} & The most octets PROGRAM-OUTPUT takes of what a program writes (64 MB): one that writes more, a clipboard of a huge copy or a program gone wrong, is stopped and its output not taken, so that it cannot fill memory.\\
\texttt{*wipe-\codebreak{}passes*} & \texttt{1} & How many times a file wiped is overwritten with random octets before it is deleted: once is enough for a disk that writes in place; more do not help one that does not (a copy-on-write file system, a solid-state drive).\\
\texttt{*markdown-\codebreak{}toc-\codebreak{}on-\codebreak{}save*} & \texttt{t} & True to make the table of contents of a Markdown file again each time it is saved, under its first heading tagged :TOC: (\#\# Contents <!-- :TOC: -->); NIL to leave it to M-x toc-org-insert-toc.\\
\bottomrule
\end{xltabular}

\subsubsection{About, the tutor and Lisp}
\label{sec:572}
\texttt{C-h C-a} shows \texttt{*About Tuatara*}, the page that introduces the program. The
tutor, a buffer of lessons, teaches the modal keys by editing it. They have
faces of their own, so that a theme can colour them.

\begin{xltabular}{\textwidth}{@{}l >{\raggedright\arraybackslash\hsize=1.000\hsize}X@{}}
\toprule
\textbf{Face} & \textbf{What it paints}\\
\midrule
\endhead
\texttt{:about-\codebreak{}heading} & The headings of \textbf{About Tuatara}.\\
\texttt{:about-\codebreak{}name} & Tuatara’s name in large letters in \textbf{About Tuatara}.\\
\texttt{:tutor-\codebreak{}arrow} & The arrow before each line to practise on in the tutor.\\
\texttt{:tutor-\codebreak{}key} & Keys in the tutor.\\
\texttt{:tutor-\codebreak{}state} & The names of the modal states in the tutor: NORMAL, INSERT, SUPPORT, TUATARA[N].\\
\bottomrule
\end{xltabular}

\texttt{C-c M-m} shows what a macro call in Lisp code expands to, which is useful when
you write your own functions in the init file. \texttt{pattern-\codebreak{}finish} takes what you
typed in a question that asks for a pattern, and says why when it cannot.

{\small
\begin{xltabular}{\textwidth}{@{}l l l >{\raggedright\arraybackslash\hsize=1.000\hsize}X@{}}
\toprule
\textbf{Modal} & \textbf{Non-modal} & \textbf{Command} & \textbf{What it does}\\
\midrule
\endhead
none & \texttt{C-c M-m} & \texttt{pp-\codebreak{}macroexpand-\codebreak{}all-\codebreak{}last-\codebreak{}expression} & Show what the macro call at point expands to, every macro in it expanded too (SB-CLTL2:MACROEXPAND-ALL), pretty-printed in \textbf{Pp Macroexpand Output}; with a prefix argument (C-u), put it in after the call instead.\\
none & \texttt{M-x pattern-\codebreak{}finish} & \texttt{pattern-\codebreak{}finish} & In a question that takes a pattern: take what is typed, or when it cannot be used, say why in the prompt and wait for it to be mended (M-x minibuffer-cancel gives up).\\
\bottomrule
\end{xltabular}
}

\section{Every setting}
\label{sec:575}
The name, standard value and purpose of every setting you can change.

The earlier chapters describe the settings of each subject where you meet it.
This chapter lists them all in one table, in the order of the alphabet, so that
you can look one up by name. The table is made from the program each time
this manual is made, so it never lists a setting that Tuatara does not have.

A setting is a name that starts and ends with a star, such as
\texttt{*auto-\codebreak{}save*}. You change one in your init file, as in
\texttt{(setf *auto-\codebreak{}save* nil)}. The table has three columns. The first is the name
of the setting. The second is its standard value, the one Tuatara has until
you change it, and a long value is cut short with three dots. The third says
what the setting changes.

Only the settings that Tuatara promises to keep are listed. A setting that
is not here may change from one version to the next, and should not be set in
an init file you keep for years.

To find a setting from within the editor, press \texttt{C-h v} and type a few
letters of its name. Tuatara shows its value and its full description. To
find the settings that mention a word, press \texttt{C-h a} and type the word.

\begin{xltabular}{\textwidth}{@{}l >{\raggedright\arraybackslash\hsize=0.727\hsize}X >{\raggedright\arraybackslash\hsize=1.273\hsize}X@{}}
\toprule
\textbf{Setting} & \textbf{Standard} & \textbf{What it changes}\\
\midrule
\endhead
\texttt{*archive-\codebreak{}free-\codebreak{}reserve*} & \texttt{1073741824} & The octets left free on the drive an archive is unpacked onto (1 GB): unpacking stops before the drive has less, so that an archive holding far more than it takes up cannot fill the drive.\\
\texttt{*archive-\codebreak{}listing-\codebreak{}limit*} & \texttt{67108864} & The most octets of an archive’s list of contents shown (64 MB); a longer one is not shown, and that is said.\\
\texttt{*archive-\codebreak{}seconds*} & \texttt{86400} & The seconds tar, 7-Zip or bsdtar may take to make, list or unpack one archive before it is ended: a day, for the largest; M-x stop-job ends it sooner.\\
\texttt{*auto-\codebreak{}save*} & \texttt{t} & True for buffers visiting files to have their changes not saved written now and then to a copy of their own in \textbf{AUTO-SAVE-DIRECTORY}, never to the file itself; M-x recover-file takes such a copy back.\\
\texttt{*auto-\codebreak{}save-\codebreak{}directory*} & \texttt{nil} & The folder auto-saved copies are kept in, made for you alone when it is not there: NIL for auto-save/ in Tuatara’s state folder (\$XDG\_STATE\_HOME/tuatara/, or \textasciitilde{}/.local/state/tuatara/), or a folder’s name, “\textasciitilde{}/backups/tuatara/“.\\
\texttt{*auto-\codebreak{}save-\codebreak{}idle*} & \texttt{5} & Seconds without a change after which a buffer’s changes not saved are auto-saved (\textbf{AUTO-SAVE}), 5 as standard.\\
\texttt{*auto-\codebreak{}save-\codebreak{}interval*} & \texttt{60} & The most seconds changes not saved go without an auto-save while they go on without a pause (\textbf{AUTO-SAVE}), 60 as standard.\\
\texttt{*auto-\codebreak{}save-\codebreak{}size-\codebreak{}limit*} & \texttt{268435456} & The largest buffer auto-saved, in octets, 256 MB as standard: a larger one is not copied (each copy would be written whole), so a crash loses its changes not saved.\\
\texttt{*auto-\codebreak{}scratch-\codebreak{}buffers*} & \texttt{("*lisp-\codebreak{}scratch*" "*org-\codebreak{}scratch*")} & The scratch buffers made without being asked for: at startup, and when the last buffer is killed.\\
\texttt{*blink-\codebreak{}matching-\codebreak{}paren*} & \texttt{t} & True to say, when a closing bracket is typed, what its partner is when it is out of sight (“Matches (defun foo ()”), or that it has none.\\
\texttt{*bookmark-\codebreak{}set-\codebreak{}asks-\codebreak{}categories*} & \texttt{t} & True for M-x bookmark-set to ask, after the name, which categories the bookmark is in (those used already offered, new ones typed); NIL never to ask.\\
\texttt{*bookmarks-\codebreak{}kept*} & \texttt{1000} & How many bookmarks you may have before setting another says you have more.\\
\texttt{*cache-\codebreak{}directory*} & \texttt{nil} & The folder Tuatara keeps its cache in, in place of \textasciitilde{}/.cache/tuatara/ (\$XDG\_CACHE\_HOME/tuatara/): files made again when they are gone, as compiled packages, previews and pages drawn, copies of files opened outside, the debug log and stderr.log.\\
\texttt{*calendar-\codebreak{}mark-\codebreak{}days*} & \texttt{t} & True for the calendar to mark the days \textbf{CALENDAR-MARK-FUNCTIONS} name (the days with agenda entries) as soon as it shows them; NIL to wait for M-x calendar-mark-days.\\
\texttt{*calendar-\codebreak{}months*} & \texttt{3} & How many months the calendar shows: three as standard, the month point is in among them, as many before it as after (one more after for an even number); twelve for a year; NIL for as many as fit side by side in its window (\textbf{CALENDAR-MONTHS-PER-ROW}, or the window’s width).\\
\texttt{*calendar-\codebreak{}months-\codebreak{}per-\codebreak{}row*} & \texttt{nil} & How many months the calendar shows side by side, the rest of \textbf{CALENDAR-MONTHS} in rows of them below, a blank line between: NIL, as standard, for as many as its window’s width holds (three in 80 columns, four in 105), a number for that many.\\
\texttt{*calendar-\codebreak{}week-\codebreak{}start-\codebreak{}day*} & \texttt{1} & The day a week begins on in the calendar: 0 Sunday, 1 Monday (as standard), up to 6 Saturday.\\
\texttt{*capture-\codebreak{}template-\codebreak{}directories*} & \texttt{nil} & The folders your capture templates are in, a template a file, as M-x org-capture offers them with Tuatara’s own (t, n, j, m); NIL, as standard, for capture-templates/ in your configuration folder (\textasciitilde{}/.config/tuatara/capture-templates/, XDG\_CONFIG\_HOME’s when it is set).\\
\texttt{*centred-\codebreak{}text-\codebreak{}modes*} & \texttt{nil} & Major modes whose buffers show the text in a column in the middle of the window, or not: a list of (MODE .\\
\texttt{*centred-\codebreak{}text-\codebreak{}width*} & \texttt{80} & How many columns of text a buffer shows when its text is centred (M-x centred-text-mode): the column’s width, the rest of the window blank margins either side.\\
\texttt{*client-\codebreak{}kept-\codebreak{}without-\codebreak{}buffers*} & \texttt{:first} & Which clients of a daemon stay when their last buffer is closed: :FIRST, the one opened first of those running; :LAST, only the last one left; :ALL, every one; NIL, none.\\
\texttt{*clipboard-\codebreak{}size-\codebreak{}limit*} & \texttt{8388608} & The most characters a kill may have and still go on the system clipboard, 8 million as standard; a longer one stays on the kill ring only.\\
\texttt{*code-\codebreak{}files-\codebreak{}make-\codebreak{}folders*} & \texttt{nil} & NIL for tangle-code-files to make the folders a file goes in only for blocks with :mkdirp yes, so that an Org file writes the same everywhere; true to make them for every block, as :mkdirp yes does (a block’s :mkdirp no still says not to).\\
\texttt{*compile-\codebreak{}command*} & \texttt{"make -k "} & The command M-x compile offers the first time, before you have typed one; after that it offers the one typed last.\\
\texttt{*compile-\codebreak{}output-\codebreak{}limit*} & \texttt{268435456} & The most octets of a command’s output M-x compile keeps in \textbf{compile}: 256 MB as standard.\\
\texttt{*complete-\codebreak{}paths*} & \texttt{t} & True to list the entries of a directory beside the cursor when a slash is typed in a buffer after its name: a path is a word with a slash in it (\emph{etc}, \textasciitilde{}/notes/, ./src/, or a file: link’s), one not starting with / or \textasciitilde{} taken from the directory the buffer stands in.\\
\texttt{*complete-\codebreak{}paths-\codebreak{}modes*} & \texttt{nil} & Major modes whose buffers list a path’s directory as it is typed, or not, whatever \textbf{COMPLETE-PATHS} says: a list of (MODE .\\
\texttt{*csv-\codebreak{}header*} & \texttt{t} & True when the first line of the files CSV mode shows is their header, naming the columns: always shown by search filters, never deleted by M-x list-keep-records, its names offered for a column.\\
\texttt{*csv-\codebreak{}separator*} & \texttt{:detect} & What separates the fields of the files CSV mode shows.\\
\texttt{*cursor-\codebreak{}colour*} & \texttt{nil} & The cursor’s colour, as faces take colours: a name such as :red, a number 0-255, or “\#rrggbb”; NIL (as standard) for the colour of the text, the character under a block then in the colour of the ground.\\
\texttt{*cursor-\codebreak{}column-\codebreak{}modes*} & \texttt{nil} & Major modes whose buffers highlight the column point is in, down the window, or not: a list of (MODE .\\
\texttt{*cursor-\codebreak{}shape*} & \texttt{:block} & The cursor’s shape: :block, :bar (a line before the character) or :underline, in the window and the terminal alike, unless the modal state has its own (\textbf{TUATARA-INSERT-CURSOR-SHAPE} and the rest).\\
\texttt{*cursors-\codebreak{}limit*} & \texttt{10000} & The most cursors M-x add-cursor-below and M-x add-cursor-above make, or NIL for as many as the memory has room for: every key but typing runs once at every cursor, so a great many make every key slow; what is typed is written as it is typed only at the cursors in sight, the others when the typing ends.\\
\texttt{*daemon-\codebreak{}indicator*} & \texttt{"D"} & What the mode line shows last among its signs before [-], always in the same place: in :mode-line-sign while this editor runs as a daemon (tuatara --daemon), in each of its clients, else faint, in :mode-line-sign-idle.\\
\texttt{*dark-\codebreak{}theme*} & \texttt{tuatara.\codebreak{}defaults::\codebreak{}ao-\codebreak{}dark} & The dark theme toggle-light-dark-theme switches to.\\
\texttt{*debug-\codebreak{}log-\codebreak{}limit*} & \texttt{10000} & The lines \textbf{Debug} keeps, the oldest going first, 10000 as standard; NIL keeps them all, and the buffer then grows for as long as debug mode is on.\\
\texttt{*debug-\codebreak{}mode*} & \texttt{nil} & True to start the editor with debug mode on, as M-x debug-mode switches it: what the editor’s parts do is reported in \textbf{Debug}.\\
\texttt{*debug-\codebreak{}on-\codebreak{}error*} & \texttt{nil} & True to show the calls under way when a command or an evaluation meets an error, in \textbf{Backtrace}; M-x toggle-debug-on-error switches it.\\
\texttt{*default-\codebreak{}directory*} & \texttt{"\textasciitilde{}/"} & The directory of a buffer with no file (\textbf{scratch}, the lists): a file name typed there is taken from it, and M-x search-in-files searches it.\\
\texttt{*default-\codebreak{}input-\codebreak{}method*} & \texttt{nil} & The input method M-x toggle-input-method turns on in a buffer that has had none on yet, such as “swedish-postfix”; NIL for it to ask.\\
\bottomrule
\end{xltabular}

{\small
\begin{xltabular}{\textwidth}{@{}l >{\raggedright\arraybackslash\hsize=0.810\hsize}X >{\raggedright\arraybackslash\hsize=1.190\hsize}X@{}}
\toprule
\textbf{Setting} & \textbf{Standard} & \textbf{What it changes}\\
\midrule
\endhead
\texttt{*differences-\codebreak{}limit*} & \texttt{67108864} & The most octets of differences shown (64 MB); when there are more, none are shown, and that is said.\\
\texttt{*electric-\codebreak{}pair-\codebreak{}mode*} & \texttt{nil} & True to close delimiters as they are typed: an opening one writes its closing half after point, a closing half typed where it stands steps over it, and M-x delete-backward-char between an empty pair takes both.\\
\texttt{*electric-\codebreak{}pair-\codebreak{}safe-\codebreak{}exclusions*} & \texttt{(\#\textbackslash{}< \#\textbackslash{}' \#\textbackslash{}` \#\textbackslash{}} \#\textbackslash{}* \#\textbackslash{}\textasciitilde{} \#\textbackslash{}\(\vert{}\) \#\textbackslash{}\textbackslash{} \#\textbackslash{}/ \#\textbackslash{}:\ldots{}= & The delimiters :safe leaves out: ordinary text far more often than a delimiter, such as the apostrophe in don’t, a colon, less-than.\\
\texttt{*electric-\codebreak{}pair-\codebreak{}typing*} & \texttt{:safe} & Which delimiters close themselves as they are typed: :safe, the pairs of \textbf{PAIR-DELIMITERS} less those \textbf{ELECTRIC-PAIR-SAFE-EXCLUSIONS} names; :all, every one of them; or a list of the opening characters, taken as it is (an empty list pairs nothing).\\
\texttt{*escape-\codebreak{}quits*} & \texttt{nil} & True to have Escape do what C-g does everywhere.\\
\texttt{*evidence-\codebreak{}folders*} & \texttt{nil} & The folders whose files are evidence: each opens safely however it is opened (read-only, no mode, hooks or version control run on it), and nothing is written in such a folder or below it, by saving, saving as, exporting, copying, moving, renaming, deleting or making folders, nor any of the editor’s own files (backups, auto-saves, locks).\\
\texttt{*export-\codebreak{}trusted-\codebreak{}folders*} & \texttt{nil} & The folders, besides a note’s own folder and the folders in it and your note templates’ folder (NOTE-TEMPLATES-DIRECTORY), that exporting the note (to PDF or to a Word document) may take files from: the setup files its \#+SETUPFILE lines name and the pictures its \href{...}{\ldots{}} links show.\\
\texttt{*extra-\codebreak{}sentence-\codebreak{}abbreviations*} & \texttt{nil} & Abbreviations of your own whose full stop ends no sentence, added to Tuatara’s (\textbf{SENTENCE-ABBREVIATIONS}), each written as typed: ’(“Ph.D.” “et al.” “Inc.”).\\
\texttt{*file-\codebreak{}local-\codebreak{}variables*} & \texttt{t} & True for a file’s own settings (-\textbf{- mode: org -}- on its first line, a Local Variables block near its end) to be taken as it is opened, namely its major mode, tab width, tabs or spaces, fill column and spaces between sentences, and those a package takes; nothing else.\\
\texttt{*fill-\codebreak{}ask-\codebreak{}column*} & \texttt{t} & True to ask for the column each time M-x fill-paragraph or M-x fill-region fills (M-x fill-paragraph and M-x fill-region too), the buffer’s fill column offered: M-x minibuffer-finish fills to it, and another number typed fills to that column and makes it the buffer’s from then on, while the editor runs.\\
\texttt{*fill-\codebreak{}column*} & \texttt{80} & The column filled lines end by, unless a buffer has its own: one given to a fill command as its count, or set with M-x set-fill-column, kept while the editor runs.\\
\texttt{*fill-\codebreak{}paragraph-\codebreak{}limit*} & \texttt{1048576} & The most octets a paragraph may span to be filled, by itself (M-x fill-paragraph, M-x fill-paragraph) or in a selection (M-x fill-region); a larger one is not filled, and says so, so that filling never holds the editor up for long.\\
\texttt{*folder-\codebreak{}tree-\codebreak{}depth-\codebreak{}limit*} & \texttt{200} & The most folders, one inside another, that copying, moving or deleting a folder goes down through.\\
\texttt{*follow-\codebreak{}desktop-\codebreak{}style*} & \texttt{t} & True to have the window’s theme follow the desktop’s light or dark style: \textbf{dark-theme} while it is dark, \textbf{light-theme} while it is light, switching when the desktop does.\\
\texttt{*format-\codebreak{}in-\codebreak{}sandbox*} & \texttt{t} & True for formatters to run shut in the sandbox: the project’s folder and the formatter’s own seen, read only, nothing written, no network.\\
\texttt{*format-\codebreak{}on-\codebreak{}save*} & \texttt{nil} & The languages whose buffers are formatted (M-x format-buffer) before each save, the save going on once the text is laid out: a list of language names, (rust python), or T for every language with a formatter installed.\\
\texttt{*format-\codebreak{}seconds*} & \texttt{30} & The seconds a formatter may take before it is stopped and nothing is changed, 30 as standard; a language server laying out code that has not answered by then is no longer waited for.\\
\texttt{*formatted-\codebreak{}file-\codebreak{}directory*} & \texttt{"\textasciitilde{}/"} & The directory of your files with formatted names: open-file-formatted lists those in it and its subdirectories, new-file-formatted makes new ones there, and the keywords in their names are offered for every new file with a formatted name.\\
\texttt{*formatted-\codebreak{}file-\codebreak{}extension*} & \texttt{""} & The end of a new file’s name that find-file-formatted offers when it asks for one (New file EXTENSION): nothing typed there takes it.\\
\texttt{*formatted-\codebreak{}file-\codebreak{}keywords*} & \texttt{nil} & Keywords offered for a new file with a formatted name, besides the keywords in the names of the files in \textbf{FORMATTED-FILE-DIRECTORY}: none as standard.\\
\texttt{*formatted-\codebreak{}file-\codebreak{}prompts*} & \texttt{(:title :subdirectory :keywords :file-\codebreak{}ty...} & What new-file-formatted asks for a new file, in this order: :title, :subdirectory (of \textbf{FORMATTED-FILE-DIRECTORY}), :keywords, :file-type (one of \textbf{FORMATTED-FILE-TYPES}; else \textbf{FORMATTED-FILE-TYPE} is used) and :template (blank, or one of your note templates’ files, asked only when there are any).\\
\texttt{*formatted-\codebreak{}file-\codebreak{}sort-\codebreak{}keywords*} & \texttt{t} & True to put the keywords of a formatted name in alphabetical order.\\
\texttt{*formatted-\codebreak{}file-\codebreak{}type*} & \texttt{:org} & The kind of a new file when \textbf{FORMATTED-FILE-PROMPTS} does not ask, and the one offered first when it does: a name in \textbf{FORMATTED-FILE-TYPES}.\\
\texttt{*formatted-\codebreak{}file-\codebreak{}types*} & \texttt{((:\codebreak{}org ".org" :org "\#+title: \{title\} \#+d...} & The kinds of files with formatted names, each (NAME EXTENSION STYLE HEADER): the end of the file’s name, how its header writes the date and keywords, and the header a new file begins with.\\
\texttt{*git-\codebreak{}commit-\codebreak{}fill-\codebreak{}column*} & \texttt{72} & The column a git message’s lines are filled to: 72, as git recommends for the body; the summary on the first line is best kept to 50.\\
\texttt{*git-\codebreak{}diff-\codebreak{}lines*} & \texttt{1000} & How many lines of a file’s changes M-x git-toggle-diff shows under it at first, and how many more each M-x git-toggle-diff on the line after them shows.\\
\texttt{*git-\codebreak{}log-\codebreak{}length*} & \texttt{200} & How many commits M-x git-log lists.\\
\texttt{*git-\codebreak{}menus*} & \texttt{:slim} & How much the git menus (branch, worktree \ldots{}) show: :SLIM the actions \textbf{GIT-SLIM-ACTIONS} names, the ones most used, :FULL every one.\\
\texttt{*git-\codebreak{}recent-\codebreak{}commits*} & \texttt{10} & How many of the latest commits M-x git shows.\\
\texttt{*git-\codebreak{}section-\codebreak{}lines*} & \texttt{1000} & How many files each section of a status (M-x git) lists at first, and how many more each M-x git-toggle-diff on the line after them lists.\\
\texttt{*git-\codebreak{}slim-\codebreak{}actions*} & \texttt{(tuatara.\codebreak{}git::\codebreak{}git-\codebreak{}commit tuatara.\codebreak{}git::\codebreak{}gi...} & The actions the slim git menus show (\textbf{GIT-MENUS}), by command: push one to show it too, or take one away.\\
\texttt{*git-\codebreak{}watch-\codebreak{}seconds*} & \texttt{2} & The least time, in seconds, between two readings of a status on the screen when git changes its repository, from a terminal too (GIT-WATCH-TICK), or NIL for a status never to follow it: then M-x revert-buffer reads it again.\\
\texttt{*global-\codebreak{}centred-\codebreak{}text-\codebreak{}mode*} & \texttt{nil} & True to show the text in a column in the middle of the window, in every buffer (NIL unless set); a buffer’s own M-x centred-text-mode, and the list for its major mode in \textbf{CENTRED-TEXT-MODES}, say otherwise for it.\\
\texttt{*global-\codebreak{}cursor-\codebreak{}column-\codebreak{}mode*} & \texttt{nil} & True to highlight the column point is in, down the window, in every buffer (NIL unless set); a buffer’s own M-x cursor-column-mode, and the list for its major mode in \textbf{CURSOR-COLUMN-MODES}, say otherwise for it.\\
\texttt{*global-\codebreak{}hl-\codebreak{}line-\codebreak{}mode*} & \texttt{t} & True to highlight the line point is on, in every buffer (T unless set); a buffer’s own M-x hl-line-mode, and the list for its major mode in \textbf{HL-LINE-MODES}, say otherwise for it.\\
\texttt{*global-\codebreak{}invisible-\codebreak{}characters-\codebreak{}mode*} & \texttt{nil} & True to show the characters that draw nothing as their codes, <FEFF> for the byte-order mark, in every buffer (NIL unless set); a buffer’s own M-x invisible-characters-mode, and the list for its major mode in \textbf{INVISIBLE-CHARACTERS-MODES}, say otherwise for it.\\
\texttt{*global-\codebreak{}soft-\codebreak{}wrap-\codebreak{}mode*} & \texttt{nil} & True to wrap long lines at a word’s edge, at the fill column, in every buffer (NIL unless set); a buffer’s own M-x soft-wrap-mode, and the list for its major mode in \textbf{SOFT-WRAP-MODES}, say otherwise for it.\\
\texttt{*global-\codebreak{}spellcheck-\codebreak{}mode*} & \texttt{nil} & True to mark the words spelt wrong, in every buffer (NIL unless set); a buffer’s own M-x spellcheck-mode, and the list for its major mode in \textbf{SPELLCHECK-MODES}, say otherwise for it.\\
\texttt{*global-\codebreak{}text-\codebreak{}direction-\codebreak{}mode*} & \texttt{:auto} & How lines holding right-to-left text are shown in buffers that have not chosen, and whose major mode \textbf{TEXT-DIRECTION-MODES} says nothing of.\\
\texttt{*global-\codebreak{}visual-\codebreak{}line-\codebreak{}mode*} & \texttt{nil} & True to mark the fill column with a line down the window, in every buffer (NIL unless set); a buffer’s own M-x visual-line-mode, and the list for its major mode in \textbf{VISUAL-LINE-MODES}, say otherwise for it.\\
\bottomrule
\end{xltabular}
}

\begin{xltabular}{\textwidth}{@{}l >{\raggedright\arraybackslash\hsize=0.810\hsize}X >{\raggedright\arraybackslash\hsize=1.190\hsize}X@{}}
\toprule
\textbf{Setting} & \textbf{Standard} & \textbf{What it changes}\\
\midrule
\endhead
\texttt{*gpg-\codebreak{}loopback-\codebreak{}warning*} & \texttt{t} & When true, the minibuffer asking for a passphrase for gpg’s loopback shows above it, while it asks, a window saying why it asks, how the passphrase goes and what that risks (the garbage collector’s copies, memory not locked, no keyboard grab, a look-alike).\\
\texttt{*gpg-\codebreak{}passphrase-\codebreak{}from*} & \texttt{:auto} & Who asks for a passphrase gpg needs.\\
\texttt{*gpg-\codebreak{}program*} & \texttt{"gpg"} & The GnuPG program that decrypts and encrypts .gpg files, by its full path, from the root (/opt/gnupg/bin/gpg), or by name, then looked for only in the system’s own folders of programs (/usr/local/bin, /usr/bin, /bin), never in a folder of yours your PATH may name, nor in the folder the editor is in.\\
\texttt{*grammars-\codebreak{}on*} & \texttt{(:common)} & The languages whose grammars read text without asking: a file in such a language is coloured and parsed as it opens, and so is the text in it inside another file (a code block in a note or a Markdown file).\\
\texttt{*gui-\codebreak{}font*} & \texttt{nil} & The font of Tuatara’s window, as Pango reads it, “Adwaita Mono 12” for one; NIL for the desktop’s monospace font.\\
\texttt{*gui-\codebreak{}padding*} & \texttt{8} & The room, in pixels, between the edges of Tuatara’s window and its text, in the ground’s colour.\\
\texttt{*gui-\codebreak{}renderer*} & \texttt{:gl} & How Tuatara’s window is drawn: :gl, by the graphics card through OpenGL, as standard; :vulkan, by the graphics card through Vulkan; :cairo, by the processor alone; NIL, as GTK chooses.\\
\texttt{*gui-\codebreak{}title*} & \texttt{"\textasciitilde{}a - Tuatara"} & The title of the window tuatara --gui opens, where the desktop names it: a format control given the name of the buffer the selected window shows.\\
\texttt{*help-\codebreak{}window-\codebreak{}select*} & \texttt{nil} & True to move to the window a help buffer is shown in (M-x describe-function, M-x describe-key and the rest); M-x quit-window there then closes it and goes back; NIL (the standard) leaves you where you were.\\
\texttt{*hex-\codebreak{}view-\codebreak{}binary-\codebreak{}files*} & \texttt{t} & True to show a file that looks binary in the hex view (M-x hex-mode) when it is opened: one with a null (NUL) octet in its first 64 KB, as executables, archives and disk images have and text has not.\\
\texttt{*hl-\codebreak{}line-\codebreak{}modes*} & \texttt{nil} & Major modes whose buffers highlight the line point is on, or not: a list of (MODE .\\
\texttt{*idle-\codebreak{}collection-\codebreak{}seconds*} & \texttt{0.3} & Seconds you leave the editor alone, no key typed, after which the garbage made since the last collection is collected, when it is at least half of what makes one (\textbf{NURSERY-OCTETS}): the pause comes while you think, not between two keys.\\
\texttt{*indent-\codebreak{}new-\codebreak{}lines*} & \texttt{t} & True to indent a line as it is begun: after RET, and on a line opened with NORMAL state’s M-x tuatara-open-below or M-x tuatara-open-above, the new line goes to its column as the language’s mode indents it (Lisp, Rust and shell scripts among them).\\
\texttt{*indent-\codebreak{}tabs-\codebreak{}mode*} & \texttt{nil} & True to put tab characters in where M-x indent-for-tab-command and indentation go on to a tab stop, NIL to put spaces in; for every buffer whose major mode says nothing of its own (set-indent-tabs-mode).\\
\texttt{*indexed-\codebreak{}directories*} & \texttt{nil} & Directories Tuatara keeps read in memory, so that M-x open-file-formatted, M-x search-in-formatted and M-x search-in-files answer at once there, and in the directories under them: the names of their files and the text of those a search reads (at most \textbf{SEARCH-FILES-LIMIT} octets each), read in the background when the editor starts and again whenever something in them changes on disk.\\
\texttt{*indexes-\codebreak{}kept-\codebreak{}on-\codebreak{}disk*} & \texttt{t} & True to keep what is read of \textbf{INDEXED-DIRECTORIES} on disk between sessions, in the cache directory (\$XDG\_CACHE\_HOME/tuatara/indexes/, readable by you alone): at the next start the directories are only listed again, and only the files changed meanwhile read, rather than all.\\
\texttt{*indexing-\codebreak{}indicator*} & \texttt{"U"} & What the mode line shows while directories are kept read (\textbf{INDEXED-DIRECTORIES}), always in the same place, before [-], after the sign of work for the buffer (\textbf{PROCESSING-INDICATOR}): in :mode-line-sign while one is being read in the background (at least \textbf{SIGN-LIT-SECONDS}), else faint, in :mode-line-sign-idle, so that nothing after it moves as reading starts and ends.\\
\texttt{*input-\codebreak{}method-\codebreak{}keys*} & \texttt{nil} & Input methods of your own under the leader’s i, each (KEY LABEL METHOD): KEY a key, LABEL what it is called, METHOD an input method’s name \textbackslash{}(M-x set-input-method lists them), such as ’((“s” “swedish” “swedish-postfix”)).\\
\texttt{*insert-\codebreak{}escape-\codebreak{}keys*} & \texttt{nil} & NIL, or a string of two characters, “jk” say: typed one right after the other in INSERT state, they leave it for NORMAL state as M-x tuatara-normal does, and neither stays in the text.\\
\texttt{*insert-\codebreak{}file-\codebreak{}limit*} & \texttt{16777216} & The most octets M-x insert-formatted-note-template reads from a file that is no note template; a larger file is refused, and said so.\\
\texttt{*invisible-\codebreak{}characters-\codebreak{}modes*} & \texttt{nil} & Major modes whose buffers show the characters that draw nothing as their codes, <FEFF> for the byte-order mark, or not: a list of (MODE .\\
\texttt{*invisible-\codebreak{}characters-\codebreak{}when-\codebreak{}safe*} & \texttt{t} & True to show the characters that draw nothing as their codes (<FEFF>, <200B>) in every file opened safely (M-x find-file-safely) that has not chosen otherwise with M-x invisible-characters-mode; NIL to leave those buffers to \textbf{INVISIBLE-CHARACTERS-MODES} and \textbf{GLOBAL-INVISIBLE-CHARACTERS-MODE} as any other.\\
\texttt{*journal-\codebreak{}directory*} & \texttt{nil} & The directory of your journal’s entries: NIL, as standard, for journal/ in \textbf{FORMATTED-FILE-DIRECTORY} (such as \textasciitilde{}/notes/journal/), or a directory of its own, such as “\textasciitilde{}/journal/“.\\
\texttt{*journal-\codebreak{}file-\codebreak{}name*} & \texttt{nil} & How journal entries are named.\\
\texttt{*journal-\codebreak{}file-\codebreak{}type*} & \texttt{nil} & The kind of file a new journal entry is: :org, :text or :markdown; NIL for \textbf{FORMATTED-FILE-TYPE}.\\
\texttt{*journal-\codebreak{}interval*} & \texttt{:daily} & How long a journal entry lasts: :daily, :weekly (the ISO week, from Monday), :monthly or :yearly.\\
\texttt{*journal-\codebreak{}keyword*} & \texttt{"journal"} & The keyword in the formatted name of each journal entry, by which entries are told from your other files: journal, as standard.\\
\texttt{*journal-\codebreak{}org-\codebreak{}header*} & \texttt{"\#+title:      \{title\} \#+date:       \{da...} & The header a new journal entry in Org begins with, its name formatted: title, date, keywords and identifier lines, \{title\}, \{date\}, \{keywords\} and \{identifier\} standing for the entry’s; NIL for the Org header of your other files with formatted names (\textbf{FORMATTED-FILE-TYPES}).\\
\texttt{*journal-\codebreak{}title-\codebreak{}format*} & \texttt{:day-\codebreak{}date-\codebreak{}month-\codebreak{}year} & The title of a new journal entry: :day (Friday), :day-date-month-year (Friday 20 March 2026, as standard), :day-date-month-year-24h (Friday 20 March 2026 13:39), :day-date-month-year-12h (Friday 20 March 2026 01:39 PM), a pattern in the \% specifiers of FORMAT-TIME (“Week \%V, \%A”), ”” for none, or NIL to be asked each time.\\
\texttt{*keep-\codebreak{}session*} & \texttt{nil} & True for the editor to keep its session: the files and directories open and how the windows show them, written as it stops (and now and then while it runs), opened again as it was when it starts without files.\\
\texttt{*keep-\codebreak{}session-\codebreak{}interval*} & \texttt{30} & Seconds between the writings of the session while the editor runs, when it has changed (\textbf{KEEP-SESSION}), 30 as standard: what a crash can lose.\\
\texttt{*key-\codebreak{}hint-\codebreak{}rows*} & \texttt{nil} & Most rows the key hints take after a prefix key (C-x, SPC), the echo row under them included, or NIL for as many as the frame has room for, as a menu takes; where they do not all fit, the last row ends in \ldots{}\\
\texttt{*kill-\codebreak{}ring-\codebreak{}max*} & \texttt{120} & The most kills the kill ring keeps, 120 as standard: older ones are forgotten.\\
\texttt{*kitty-\codebreak{}keyboard*} & \texttt{:auto} & Whether the terminal frontend reads keys by the kitty keyboard protocol.\\
\texttt{*language-\codebreak{}formatters*} & \texttt{nil} & The formatters of languages you choose yourself, in place of the language’s usual one: each (LANGUAGE PROGRAM ARGUMENT\ldots{}), the code given on its standard input and taken back laid out from its standard output; (LANGUAGE) alone for none, the language server or Tuatara’s own indentation laying it out instead.\\
\texttt{*language-\codebreak{}servers*} & \texttt{nil} & The language servers of languages you choose yourself, in place of the language’s usual ones: each (LANGUAGE SERVER\ldots{}), LANGUAGE a symbol (python, rust), each SERVER a known server’s name (“basedpyright”) or a list of a program and its arguments ((“my-server” ”--stdio”)); (LANGUAGE) alone for none.\\
\texttt{*last-\codebreak{}used-\codebreak{}first*} & \texttt{t} & True for the questions that choose a file or a folder (M-x find-file, M-x rimu and the rest) to start in the folder of the path chosen last in any of them, and to list first, in each folder, what was chosen there last (or something in it), TAB then going the way last gone; the others follow in order of name.\\
\texttt{*latex-\codebreak{}engines*} & \texttt{("xelatex" "lualatex")} & The programs that make a PDF of \LaTeX{}, tried in this order: when one fails, the next is tried.\\
\texttt{*latex-\codebreak{}main-\codebreak{}font*} & \texttt{nil} & The font the enhanced export sets a document’s text in, by its family’s name as fontconfig knows it (“Noto Serif”); NIL for the first installed of Libertinus Serif, Noto Serif, Source Serif 4, Liberation Serif, DejaVu Serif.\\
\texttt{*latex-\codebreak{}mono-\codebreak{}font*} & \texttt{nil} & The fixed-width font the enhanced export sets code in, as \textbf{LATEX-MAIN-FONT}; NIL for your window’s font (\textbf{gui-font}) when it is installed, else the first of JetBrains Mono, Cascadia Code, Noto Sans Mono, Source Code Pro, Liberation Mono, DejaVu Sans Mono.\\
\bottomrule
\end{xltabular}

\begin{xltabular}{\textwidth}{@{}l >{\raggedright\arraybackslash\hsize=0.889\hsize}X >{\raggedright\arraybackslash\hsize=1.111\hsize}X@{}}
\toprule
\textbf{Setting} & \textbf{Standard} & \textbf{What it changes}\\
\midrule
\endhead
\texttt{*latex-\codebreak{}redaction-\codebreak{}commands*} & \texttt{("R" "D" "redact")} & The \LaTeX{} commands that hide their argument in your \LaTeX{} files, \texttt{\textbackslash{}\textbackslash{}R\{John} Otago\} for one, by their names without the backslash: what M-x export-check-redactions and M-x export-latex-compile-pdf look for in a \LaTeX{} file, as hidden text to be found nowhere else.\\
\texttt{*latex-\codebreak{}sans-\codebreak{}font*} & \texttt{nil} & The sans-serif font of the enhanced export, as \textbf{LATEX-MAIN-FONT}; NIL for the first installed of Libertinus Sans, Noto Sans, Source Sans 3, Liberation Sans, DejaVu Sans.\\
\texttt{*leader-\codebreak{}key*} & \texttt{"SPC"} & The key that opens the leader’s keys (bind-key’s :map :leader) in NORMAL and support states, in lists and in the calendar: SPC, or another written as bind-key writes keys (“,”, “C-c SPC”).\\
\texttt{*light-\codebreak{}theme*} & \texttt{tuatara.\codebreak{}defaults::\codebreak{}ao-\codebreak{}light} & The light theme toggle-light-dark-theme switches to.\\
\texttt{*line-\codebreak{}numbers-\codebreak{}default*} & \texttt{nil} & Line numbers for buffers that have not chosen, and whose major mode \textbf{LINE-NUMBERS-MODES} says nothing of: NIL for none, :absolute, or :relative (the current line shows its own number, the others their distance).\\
\texttt{*line-\codebreak{}numbers-\codebreak{}modes*} & \texttt{nil} & Major modes whose buffers show line numbers in the gutter, or not: a list of (MODE .\\
\texttt{*line-\codebreak{}record-\codebreak{}modes*} & \texttt{(tuatara.\codebreak{}defaults::\codebreak{}log-\codebreak{}mode tuatara.\codebreak{}defa...} & The major modes whose files hold a record a line, as a log does: search filters (M-x search-filter-regex) narrow them, and the modes made from them, by hiding the lines they leave out, as they narrow a CSV.\\
\texttt{*list-\codebreak{}header-\codebreak{}lock*} & \texttt{nil} & True to keep the header of a file of records at the top of the window as the rest scrolls under it, wherever there is one: a CSV’s first line, a log’s once M-x list-toggle-header makes it a header; as M-x list-toggle-header-lock does for one buffer.\\
\texttt{*logos-\codebreak{}directory*} & \texttt{nil} & The folder of the pictures your templates ask a note for, the logo its \#+logo: names (LOGOS-DIRECTORY); NIL, as standard, for logos/ in the configuration directory.\\
\texttt{*lsp-\codebreak{}code-\codebreak{}lenses*} & \texttt{nil} & True for what language servers say of lines (3 references, Run test) to be shown faint at each line’s end.\\
\texttt{*lsp-\codebreak{}diagnostics-\codebreak{}display*} & \texttt{:box} & How what a language server found where the cursor rests is said: :box, in a box under the cursor; :echo, in the echo area; NIL, not at all (the list of M-x lsp-diagnostics still has it).\\
\texttt{*lsp-\codebreak{}diagnostics-\codebreak{}style*} & \texttt{:wave} & How the text a language server finds something wrong in is marked: :wave, errors and warnings underlined wavy, information and hints dotted; :dotted, all dotted; :line, all with a straight line; :background, the text on a ground of its kind’s colour (the theme’s error, warning and information grounds), not underlined.\\
\texttt{*lsp-\codebreak{}highlight-\codebreak{}symbol*} & \texttt{t} & True for the name the cursor rests on, and its other uses the language server finds in the file, to be shown in :symbol-highlight.\\
\texttt{*lsp-\codebreak{}hover-\codebreak{}echo*} & \texttt{t} & True for the echo area to show what the thing the cursor rests on is (its type, its signature), as the language server knows it, when the cursor is in no call whose parameters are shown.\\
\texttt{*lsp-\codebreak{}idle-\codebreak{}delay*} & \texttt{0.5} & The seconds the cursor rests before what is known of where it stands is shown: what a server found wrong there, the parameters of the call, what the thing is, 0.5 as standard.\\
\texttt{*lsp-\codebreak{}inlay-\codebreak{}hints*} & \texttt{t} & True for language servers’ inlay hints to be shown within the code: the types of bindings, the names of the parameters arguments go to.\\
\texttt{*lsp-\codebreak{}languages*} & \texttt{t} & The languages whose language servers are used: T, every language with a server installed; a list of languages, those alone, ’(rust go python); or :NOT and languages, every one but those, ’(:not markdown latex).\\
\texttt{*lsp-\codebreak{}outline*} & \texttt{t} & True for M-x buffer-outline to list a file’s names as its language server knows them (functions, types, their fields and methods, nested), where it has one that does; else as the mode finds its headings.\\
\texttt{*lsp-\codebreak{}semantic-\codebreak{}faces*} & \texttt{(("namespace" . :namespace) ("type" . :t...} & The face each kind of name a language server tells apart is coloured in, as (KIND .\\
\texttt{*lsp-\codebreak{}semantic-\codebreak{}tokens*} & \texttt{nil} & True for code to be coloured as its language server sees it: each name in the face \textbf{LSP-SEMANTIC-FACES} gives its kind, over the mode’s colours.\\
\texttt{*lsp-\codebreak{}signature-\codebreak{}help*} & \texttt{t} & True for the echo area to show the parameters of the call the cursor is in, as the language server knows them, once the cursor rests there.\\
\texttt{*lsp-\codebreak{}trusted-\codebreak{}projects*} & \texttt{nil} & Folders whose projects you trust to run their own programs: their language servers start by themselves as their files are opened, git runs the programs they name for themselves (hooks, filters), and M-x compile runs their build commands without asking.\\
\texttt{*lsp-\codebreak{}wait*} & \texttt{2} & The seconds a command waits for a language server’s answer it needs at once (a completion, a definition), a key stopping the wait sooner, 2 as standard.\\
\texttt{*lsp-\codebreak{}watch-\codebreak{}files*} & \texttt{t} & True for language servers that ask to be told of their project’s files changed on disk, made or deleted, by programs other than the editor, to be told.\\
\texttt{*lsp-\codebreak{}watch-\codebreak{}skipped*} & \texttt{("target" "node\_\codebreak{}modules" "\_\_\codebreak{}pycache\_\_" "...} & The folders of a project whose files are never watched for its language servers, by name (hidden folders, whose names begin with a dot, never are): what a build makes.\\
\texttt{*mark-\codebreak{}ring-\codebreak{}max*} & \texttt{16} & How many marks each buffer’s mark ring keeps, and the global mark ring across buffers; past that the oldest goes.\\
\texttt{*markdown-\codebreak{}toc-\codebreak{}on-\codebreak{}save*} & \texttt{t} & True to make the table of contents of a Markdown file again each time it is saved, under its first heading tagged :TOC: (\#\# Contents <!-- :TOC: -->); NIL to leave it to M-x toc-org-insert-toc.\\
\texttt{*message-\codebreak{}log-\codebreak{}max*} & \texttt{1000} & Lines the \textbf{Messages} buffer keeps of the messages shown, the oldest going first, 1000 as standard; NIL keeps none.\\
\texttt{*metadata-\codebreak{}look-\codebreak{}up-\codebreak{}key-\codebreak{}owners*} & \texttt{:ask} & Whether M-x show-file-metadata looks up, in your own GnuPG keyring, whose the keys an encrypted file is for are, running gpg --list-keys (the file itself is never given to gpg): true looks them up, :ask asks each time, false (nil) never does, the keys then shown by their IDs alone.\\
\texttt{*modal-\codebreak{}editing*} & \texttt{t} & True for Tuatara’s modal editing: NORMAL and INSERT states in editing buffers, NORMAL’s reading keys in support buffers, the leader on SPC.\\
\texttt{*mode-\codebreak{}line-\codebreak{}name-\codebreak{}limit*} & \texttt{nil} & The most columns a buffer’s name takes in the mode line, or NIL for as many as it has: a longer name keeps its start and its end, … standing for its middle.\\
\texttt{*mode-\codebreak{}line-\codebreak{}tuatara*} & \texttt{t} & True to write TUATARA before the modal state in the mode line, TUATARA[N]; NIL for the state alone, [N].\\
\texttt{*notify-\codebreak{}program*} & \texttt{"notify-\codebreak{}send"} & The program that shows a notification on the desktop: a name looked for on the PATH, or a path.\\
\texttt{*office-\codebreak{}note-\codebreak{}size-\codebreak{}limit*} & \texttt{16777216} & The most characters a note read from a Word document, an Excel workbook, a PowerPoint presentation or an OpenDocument file may have: past them the rest of its text is left out, and the note says so, so that a file whose few octets stand for text copied many times cannot fill the memory.\\
\texttt{*office-\codebreak{}value-\codebreak{}length-\codebreak{}limit*} & \texttt{2000} & The most characters of one value an office file’s note shows (a cell of a sheet or a table, a formula): a longer one is cut after a word, with \ldots{}\\
\texttt{*office-\codebreak{}view-\codebreak{}pages*} & \texttt{t} & True to show a Word document, an Excel workbook, a PowerPoint presentation (and the older and OpenDocument kinds of them) opened as its pages, drawn by LibreOffice shut in by the sandbox, when it is installed; M-x pdf-view-toggle-text switches to it read as text.\\
\texttt{*open-\codebreak{}binary-\codebreak{}files-\codebreak{}safely*} & \texttt{t} & True to open safely, as M-x find-file-safely does, a file that looks binary: one with a null (NUL) octet in its first 64 KB, as executables, archives and disk images have and text has not.\\
\texttt{*open-\codebreak{}dropped-\codebreak{}files-\codebreak{}safely*} & \texttt{nil} & True to open files dropped on the window safely, as M-x universal-argument M-x find-file does: read-only, no mode, nothing written beside them.\\
\texttt{*open-\codebreak{}externally-\codebreak{}ask*} & \texttt{t} & True for M-x open-externally to ask before opening a file in the desktop’s own program, which reads it outside the sandbox; NIL to open it at once.\\
\texttt{*open-\codebreak{}externally-\codebreak{}kinds*} & \texttt{(:png :jpeg :gif :webp :bmp :tiff :raw :...} & The kinds of file M-x open-externally opens in the desktop’s own program, as their first octets tell them (pictures, PDFs, sound and video, and Office documents, as standard): keywords as TUATARA.IO:OCTETS-KIND gives them, :png, :jpeg, :pdf, :docx and the like.\\
\bottomrule
\end{xltabular}

{\small
\begin{xltabular}{\textwidth}{@{}l >{\raggedright\arraybackslash\hsize=0.935\hsize}X >{\raggedright\arraybackslash\hsize=1.065\hsize}X@{}}
\toprule
\textbf{Setting} & \textbf{Standard} & \textbf{What it changes}\\
\midrule
\endhead
\texttt{*open-\codebreak{}externally-\codebreak{}program*} & \texttt{"xdg-\codebreak{}open"} & The program that opens a file in the desktop’s own program for its kind, given the file’s name: xdg-open, looked for on the PATH.\\
\texttt{*open-\codebreak{}externally-\codebreak{}programs*} & \texttt{nil} & Programs to open files and addresses in instead of the one the desktop would choose (xdg-open), by Multipurpose Internet Mail Extensions (MIME) type as xdg-mime names them: a list of (TYPE .\\
\texttt{*org-\codebreak{}agenda-\codebreak{}autosave-\codebreak{}files*} & \texttt{nil} & The agenda files saved as soon as a change made from the agenda (a keyword, a date, tags, a note, the clock) changes them: NIL for none (M-x org-agenda-save-all-buffers saves them), T for every agenda file, or a list of files, folders and patterns as \textbf{ORG-AGENDA-FILES} takes them: ’(“\textasciitilde{}/notes/agenda.org”).\\
\texttt{*org-\codebreak{}agenda-\codebreak{}category-\codebreak{}width*} & \texttt{12} & How wide the agenda’s first column is, its category (the file’s name, or its \#+CATEGORY:) and a colon, in columns: a longer one is cut short with …, so that every line’s time and text stand in one column.\\
\texttt{*org-\codebreak{}agenda-\codebreak{}current-\codebreak{}time-\codebreak{}string*} & \texttt{"← now ────────────────\codebreak{}────────────────\codebreak{}─...} & The line that shows the time now in today’s time grid; NIL for none.\\
\texttt{*org-\codebreak{}agenda-\codebreak{}custom-\codebreak{}commands*} & \texttt{(("u" "Untagged tasks" :tags-\codebreak{}todo "-\{.*\}...} & Views of your own that M-x org-agenda offers beside its own: each (KEY DESCRIPTION TYPE MATCH :header HEADER).\\
\texttt{*org-\codebreak{}agenda-\codebreak{}date-\codebreak{}in-\codebreak{}title*} & \texttt{nil} & True to keep the dates written in a heading or a checkbox’s line in its title on the agenda’s days, as Org shows them; NIL takes them out there, the day and the time column saying them already (the TODO list and matches keep them).\\
\texttt{*org-\codebreak{}agenda-\codebreak{}event-\codebreak{}kind-\codebreak{}text*} & \texttt{"Event:"} & What the agenda’s column of kinds (Scheduled:, Deadline:) says on the line of an event, a date in the text of an entry with no TODO keyword (:org-agenda-event-kind), Event:.\\
\texttt{*org-\codebreak{}agenda-\codebreak{}files*} & \texttt{("\textasciitilde{}/\codebreak{}org/\codebreak{}agenda.\codebreak{}org")} & The files the agenda reads: each a file, a folder (every .org file in it, those whose names begin with a dot left out) or a pattern of names in a folder, * for any characters and ? for one (“\textasciitilde{}/notes/*.org”), names beginning with a dot only for a pattern that does.\\
\texttt{*org-\codebreak{}agenda-\codebreak{}log-\codebreak{}mode*} & \texttt{nil} & True for the agenda to list what was closed and clocked on each day too.\\
\texttt{*org-\codebreak{}agenda-\codebreak{}refresh-\codebreak{}minutes*} & \texttt{10} & Minutes after which an agenda shown is drawn again by itself, so that the time now in its grid and the days left to deadlines stay true; NIL for never.\\
\texttt{*org-\codebreak{}agenda-\codebreak{}span*} & \texttt{:week} & How many days M-x org-agenda’s agenda shows: :day, :week, :month or :year, or a number of days.\\
\texttt{*org-\codebreak{}agenda-\codebreak{}start-\codebreak{}on-\codebreak{}weekday*} & \texttt{1} & The day a week’s agenda starts on, 1 for Monday to 6 for Saturday and 7 or 0 for Sunday; NIL to start on today.\\
\texttt{*org-\codebreak{}agenda-\codebreak{}task-\codebreak{}kind-\codebreak{}text*} & \texttt{"Task:"} & What the agenda’s column of kinds says on the line of a date in the text of a task, an entry with a TODO keyword, not on its SCHEDULED: or DEADLINE: line: Task:, told apart from an event’s (\textbf{ORG-AGENDA-EVENT-KIND-TEXT}, which says whether the column is there at all).\\
\texttt{*org-\codebreak{}agenda-\codebreak{}time-\codebreak{}grid*} & \texttt{((:\codebreak{}daily :today :require-\codebreak{}timed) (800 100...} & Today’s time grid in the agenda: (WHEN TIMES AFTER-TIME LINE).\\
\texttt{*org-\codebreak{}agenda-\codebreak{}untimed-\codebreak{}text*} & \texttt{"all day"} & What the agenda’s time column says on a line with no time of day, faint: “all day”, so that every line’s kind and text stand in one column; ”” leaves the column empty, lined up all the same; NIL leaves it out, the line’s text where the time would have been.\\
\texttt{*org-\codebreak{}bullets*} & \texttt{("{\tuataraFallbackTWO{}◉}" "○" "{\tuataraFallbackONE{}✸}" "{\tuataraFallbackONE{}✿}")} & The signs Org headlines show in place of their last star, by level, the first for the first level and round again past the last; NIL to show the stars.\\
\texttt{*org-\codebreak{}checkbox-\codebreak{}signs*} & \texttt{nil} & The signs list items’ checkboxes show, each (CHECKBOX .\\
\texttt{*org-\codebreak{}clock-\codebreak{}into-\codebreak{}drawer*} & \texttt{t} & True for clock lines to go into a :LOGBOOK: drawer of their headline; NIL for them to go right under it.\\
\texttt{*org-\codebreak{}complete-\codebreak{}checkboxes*} & \texttt{t} & True for [ typed right after a list item’s bullet and its space (- [) to make it a checkbox at once, - [ ] with point after the space, ready for the item’s text; NIL to type [ as anywhere else.\\
\texttt{*org-\codebreak{}cycle-\codebreak{}hide-\codebreak{}drawers*} & \texttt{t} & True for drawers (:LOGBOOK:, :PROPERTIES:) and a task’s notes (ORG-NOTES-RANGE) to stay folded when TAB shows a headline’s or an item’s children or all of it; TAB on a drawer’s line, or on the first of the notes, opens it.\\
\texttt{*org-\codebreak{}cycle-\codebreak{}include-\codebreak{}plain-\codebreak{}lists*} & \texttt{t} & True for TAB on a list item with sub-items or text under its first line to fold it, as on a headline: folded, then its sub-items, then all of it; NIL for TAB there to indent.\\
\texttt{*org-\codebreak{}deadline-\codebreak{}warning-\codebreak{}days*} & \texttt{14} & How many days before its deadline a task is listed on today; a deadline’s own -3d says otherwise.\\
\texttt{*org-\codebreak{}depth-\codebreak{}limit*} & \texttt{200} & The most elements and objects of Org text the parser reads one within another: what lies deeper is read as plain text, so that a file made of nothing but nesting cannot exhaust the parser.\\
\texttt{*org-\codebreak{}fold-\codebreak{}metadata*} & \texttt{t} & True for an Org file to open with what is kept of each task folded: its drawers (:PROPERTIES:, :LOGBOOK:) and its notes (- State “DONE” \ldots{}, - Note taken on \ldots{}, CLOCK: lines) each to one line, those of headings and list items alike; its SCHEDULED: and DEADLINE: shown.\\
\texttt{*org-\codebreak{}formula-\codebreak{}depth-\codebreak{}limit*} & \texttt{200} & The most a formula’s parts may be nested, within brackets, signs and powers: a formula nested deeper is refused.\\
\texttt{*org-\codebreak{}hide-\codebreak{}leading-\codebreak{}stars*} & \texttt{t} & True to hide a headline’s stars but the last, with \textbf{ORG-BULLETS} or in the indented view; NIL to show them faint (the :org-leading-stars face).\\
\texttt{*org-\codebreak{}link-\codebreak{}descriptive*} & \texttt{t} & True for Org links to be shown as their labels, \hyperref[org376]{label} as label and \hyperref[org376]{target} as target, and entities as their characters, \texttt{\textbackslash{}dollar\{\}} as \$; NIL to show them as written.\\
\texttt{*org-\codebreak{}list-\codebreak{}bullets*} & \texttt{nil} & The signs list items show in place of their bullets, each (BULLET .\\
\texttt{*org-\codebreak{}log-\codebreak{}checkboxes*} & \texttt{t} & True to log when a list item’s checkbox is ticked or cleared, in a :LOGBOOK: drawer under the item: - State “DONE” from “TODO” and the time, as a TODO headline’s change of state is logged; NIL to log nothing.\\
\texttt{*org-\codebreak{}log-\codebreak{}done*} & \texttt{:time} & What is noted when a task is done: :TIME (as standard) a CLOSED: line with the time, which the agenda’s log mode lists; NIL nothing.\\
\texttt{*org-\codebreak{}log-\codebreak{}into-\codebreak{}drawer*} & \texttt{t} & True for the notes logged of a task (a change of state, a repeat, a reschedule, a note taken) to go into the entry’s :LOGBOOK: drawer, where they fold with it (\textbf{ORG-FOLD-METADATA}); NIL for them to go right after its planning line and properties, outside any drawer.\\
\texttt{*org-\codebreak{}log-\codebreak{}redeadline*} & \texttt{nil} & What Org notes when a task’s deadline is changed or taken away: NIL nothing, :TIME the deadline it had and when.\\
\texttt{*org-\codebreak{}log-\codebreak{}repeat*} & \texttt{:time} & What Org notes when a task that repeats is done: :TIME the change of state and LAST\_REPEAT with the time; NIL nothing.\\
\texttt{*org-\codebreak{}log-\codebreak{}reschedule*} & \texttt{nil} & What Org notes when a task’s scheduled date is changed or taken away: NIL nothing, :TIME the date it had and when.\\
\texttt{*org-\codebreak{}notify-\codebreak{}before*} & \texttt{10} & How many minutes before its time an entry tagged :notify: is told, as well as when the time comes; 0 or NIL for only then.\\
\texttt{*org-\codebreak{}notify-\codebreak{}day-\codebreak{}time*} & \texttt{"08:00"} & When an entry tagged :notify: with a date and no time of day is told on that day, or as Tuatara starts when that is later; NIL never tells them.\\
\texttt{*org-\codebreak{}notify-\codebreak{}tag*} & \texttt{"notify"} & The tag of the agenda’s entries told on the desktop as their time comes (a heading’s own or one it takes from those above it); NIL tells none.\\
\texttt{*org-\codebreak{}read-\codebreak{}date-\codebreak{}in-\codebreak{}calendar*} & \texttt{t} & True for Org’s date question (M-x org-schedule, M-x org-deadline, M-x org-time-stamp, a capture template’s \%\textasciicircum{}T) to begin in the calendar shown below it (\textbf{ORG-READ-DATE-POPUP-CALENDAR}).\\
\texttt{*org-\codebreak{}read-\codebreak{}date-\codebreak{}popup-\codebreak{}calendar*} & \texttt{t} & True for Org’s date prompt (M-x org-schedule, M-x org-deadline) to show the calendar below while it asks: the day read from what is typed marked there, S-<left> and S-<right> a day either way, S-<up> and S-<down> a week, M-S-<left> and M-S-<right> a month, M-S-<up> and M-S-<down> a year, C-.\\
\bottomrule
\end{xltabular}
}

\begin{xltabular}{\textwidth}{@{}l >{\raggedright\arraybackslash\hsize=0.857\hsize}X >{\raggedright\arraybackslash\hsize=1.143\hsize}X@{}}
\toprule
\textbf{Setting} & \textbf{Standard} & \textbf{What it changes}\\
\midrule
\endhead
\texttt{*org-\codebreak{}refile-\codebreak{}targets*} & \texttt{((:\codebreak{}agenda-\codebreak{}files . 3) (:this-\codebreak{}file . 3))} & Where M-x org-refile offers to move an entry: a list of (WHERE .\\
\texttt{*org-\codebreak{}return-\codebreak{}cycles-\codebreak{}checkboxes*} & \texttt{t} & True for M-x tuatara-normal-return in NORMAL state anywhere on a list item’s line with a checkbox to cycle it, [ ] to [-] to [X] and round, the counts above brought up to date (a link point is on is followed first); NIL to leave it to M-x org-ctrl-c-ctrl-c.\\
\texttt{*org-\codebreak{}return-\codebreak{}cycles-\codebreak{}todo*} & \texttt{t} & True for M-x tuatara-normal-return in NORMAL state on a headline’s line with a TODO keyword to change it to the next, as M-x org-todo does, but never to none: TODO to WAIT to DONE and round, or round the document’s \#+TODO: keywords, those of \textbf{org-todo-keywords-not-cycled} passed over (a link point is on is followed first).\\
\texttt{*org-\codebreak{}return-\codebreak{}follows-\codebreak{}link*} & \texttt{t} & True for RET on an Org link to follow it; M-x org-open-at-point follows one always.\\
\texttt{*org-\codebreak{}src-\codebreak{}languages*} & \texttt{(("emacs-\codebreak{}lisp" . "elisp") ("C++" . "cpp"...} & The language each name a source block is given (\#+begin\_src emacs-lisp) means, where Org’s name is not the language’s own: each (NAME .\\
\texttt{*org-\codebreak{}standard-\codebreak{}todo-\codebreak{}keywords*} & \texttt{(("TODO" "WAIT") "DONE" "CANCELLED")} & The TODO keywords of an Org file with no \#+TODO: line of its own, as (NOT-DONE .\\
\texttt{*org-\codebreak{}startup-\codebreak{}indented*} & \texttt{t} & True for Org files to open in the indented view (M-x org-indent-mode), as they do: each line shown indented by the level of the headline it is under, a headline’s stars but the last hidden, the file itself as it is; NIL for them to open as they are written.\\
\texttt{*org-\codebreak{}tags-\codebreak{}column*} & \texttt{0} & Where M-x org-set-tags-command and M-x org-align-tags put a headline’s tags: 0, one space after the headline’s text; negative, so that they end at that column (-77 to end at the 77th); positive, so that they start there.\\
\texttt{*org-\codebreak{}toc-\codebreak{}max-\codebreak{}depth*} & \texttt{2} & How deep the headlines of a table of contents go when its tag says nothing (:TOC:), 2 for those of one and two stars.\\
\texttt{*org-\codebreak{}toc-\codebreak{}on-\codebreak{}save*} & \texttt{t} & True to make the table of contents of an Org file again each time it is saved, under its first headline tagged :TOC:; NIL to leave it to M-x toc-org-insert-toc.\\
\texttt{*org-\codebreak{}toc-\codebreak{}style*} & \texttt{"gh"} & How a table of contents links to its headlines when its tag says nothing, “gh” for GitHub’s anchors (\#a-headline, which GitHub shows a README’s headlines by), “org” for their titles.\\
\texttt{*org-\codebreak{}todo-\codebreak{}follows-\codebreak{}checkboxes*} & \texttt{t} & True for a TODO headline with a statistics cookie ([2/3], [66\%]) to become done by itself once all it counts are done, its TODO sub-headlines (or, with none, the checkboxes right under it), with CLOSED: and the time, and not done again when one is not yet done; up the tree, a headline above with a cookie following in turn.\\
\texttt{*org-\codebreak{}todo-\codebreak{}keyword-\codebreak{}faces*} & \texttt{(("WAIT" . :org-\codebreak{}wait) ("CANCELLED" . :or...} & The TODO keywords with a face of their own, each (KEYWORD .\\
\texttt{*org-\codebreak{}todo-\codebreak{}keywords-\codebreak{}not-\codebreak{}cycled*} & \texttt{("CANCELLED" "CANCELED")} & The TODO keywords passed over as a task’s keyword goes round to the next: by M-x org-todo, M-x org-shiftright and M-x org-shiftleft, by RET on a task’s headline and by the agenda’s M-x org-agenda-todo.\\
\texttt{*package-\codebreak{}ask-\codebreak{}before-\codebreak{}loading*} & \texttt{t} & True to be asked before a package is first compiled and loaded: one fetched the first time, and again when an update brings another commit; one of a folder of yours the first time only, what you change in it later being yours.\\
\texttt{*package-\codebreak{}check-\codebreak{}signatures*} & \texttt{t} & True to check who signed each version of a package, and take only those that M-x package-trust-again and the signers of a source (:signers) allow.\\
\texttt{*package-\codebreak{}size-\codebreak{}limit*} & \texttt{536870912} & The most octets a package may take on disk while it is fetched, what git keeps of its history and its files together: a fetch that grows past it is ended and its files deleted.\\
\texttt{*package-\codebreak{}sources*} & \texttt{nil} & The places packages come from, each (NAME .\\
\texttt{*packages-\codebreak{}directory*} & \texttt{nil} & The folder Tuatara keeps the packages it fetches in, in place of \textasciitilde{}/.local/share/tuatara/packages/ (\$XDG\_DATA\_HOME/tuatara/packages/); NIL for that standard place.\\
\texttt{*pair-\codebreak{}delimiters*} & \texttt{((\#\textbackslash{}\{ . \#\textbackslash{}\}) (\#\textbackslash{}[ . \#\textbackslash{}]) (\#\textbackslash{}( . \#\textbackslash{})) (\#\textbackslash{}...} & The pairs of delimiters, each (OPEN .\\
\texttt{*pdf-\codebreak{}view-\codebreak{}dark-\codebreak{}pages*} & \texttt{nil} & True to show a PDF’s pages dark: white as the theme’s background, black as its text, the colours between in proportion.\\
\texttt{*picture-\codebreak{}memory*} & \texttt{4294967296} & The most memory the program drawing a large picture is given (4 GB): a picture’s own size, read from its file, tells how much it needs (about 32 octets a pixel, a hundred-megapixel PNG some 3.9 GB), never less than \textbf{SANDBOX-MEMORY}; one needing more than this is not drawn, and says so.\\
\texttt{*preview-\codebreak{}buffers*} & \texttt{nil} & True to show, while you choose a buffer with M-x switch-to-buffer (and M-x switch-to-buffer-other-window, and every other question for a buffer), the buffer selected in the window you are in, what it showed coming back when the question ends.\\
\texttt{*preview-\codebreak{}file-\codebreak{}search*} & \texttt{nil} & True to show, while you search the files of a project or a directory with M-x search-in-files, the file of the line selected in the window you are in, at that line, opened safely, as \textbf{PREVIEW-FILES} shows files.\\
\texttt{*preview-\codebreak{}files*} & \texttt{nil} & True to show, while you choose a file with M-x find-file (and the other ways of finding a file) or M-x find-recent-file, the file selected in the window you are in, opened safely: read-only, no mode or hook run on it, its times untouched, a binary one in the hex view; what the window showed comes back when the question ends.\\
\texttt{*preview-\codebreak{}formatted-\codebreak{}files*} & \texttt{nil} & True to show, while you choose one of your files with formatted names with M-x open-file-formatted, the file selected in the window you are in, opened safely, as \textbf{PREVIEW-FILES} shows files.\\
\texttt{*preview-\codebreak{}formatted-\codebreak{}search*} & \texttt{nil} & True to show, while you search your files with formatted names with M-x search-in-formatted, the file of the line selected in the window you are in, at that line, opened safely, as \textbf{PREVIEW-FILES} shows files.\\
\texttt{*processing-\codebreak{}indicator*} & \texttt{"P"} & What the mode line shows first, before [-] and the signs of directories kept read (\textbf{INDEXING-INDICATOR}), always in the same place: in :mode-line-sign while a job runs or waits its turn in the buffer, or another program works for it (a language server indexing), else faint, in :mode-line-sign-idle, so that nothing after it moves as work starts and ends.\\
\texttt{*program-\codebreak{}output-\codebreak{}limit*} & \texttt{67108864} & The most octets PROGRAM-OUTPUT takes of what a program writes (64 MB): one that writes more, a clipboard of a huge copy or a program gone wrong, is stopped and its output not taken, so that it cannot fill memory.\\
\texttt{*project-\codebreak{}root-\codebreak{}markers*} & \texttt{(".\codebreak{}git" ".project")} & The names of the files or directories that mark the root of a project: M-x search-in-files searches all of the project a buffer’s file is in, from the nearest directory above the file holding one of them.\\
\texttt{*prompt-\codebreak{}above-\codebreak{}candidates*} & \texttt{t} & True to show the prompt of a completing read above its candidates, the list reading down from it; NIL to show it on the bottom row, under them.\\
\texttt{*recent-\codebreak{}files-\codebreak{}auto-\codebreak{}cleanup*} & \texttt{t} & True to take out of the list of recent files, when the editor starts, the files no longer there; NIL keeps them, shown as not found.\\
\texttt{*recent-\codebreak{}files-\codebreak{}limit*} & \texttt{25} & How many files the list of recent files keeps, 25 as standard.\\
\texttt{*rectangle-\codebreak{}cursors-\codebreak{}limit*} & \texttt{10000} & The most rows M-x tuatara-change on a rectangle (M-x rectangle-mark-mode) puts a cursor on, or NIL for as many as the memory has room for: the block is emptied and what is typed appears on every row, as at the cursors of M-x add-cursor-below.\\
\texttt{*relative-\codebreak{}age-\codebreak{}limit*} & \texttt{1209600} & Seconds within which the last change of a file among completion candidates is shown as an age (5 min ago, 3 days ago); older ones show the date, and with 0 every one does.\\
\texttt{*reload-\codebreak{}changed-\codebreak{}files*} & \texttt{nil} & True to have every buffer keep up with its file on disk by itself.\\
\texttt{*remember-\codebreak{}recent-\codebreak{}files*} & \texttt{t} & True to keep a list of the files you visit, the latest first, for find-recent-file (M-x find-recent-file, M-x find-recent-file) to offer; NIL adds nothing to it.\\
\texttt{*remote-\codebreak{}preview-\codebreak{}size-\codebreak{}limit*} & \texttt{20000000} & The most octets a file on another machine may hold to be shown in the preview of its list (M-x list-toggle-preview), 20 MB by default: all of it has to come from the server first.\\
\texttt{*remote-\codebreak{}seconds*} & \texttt{30} & Seconds a server reached for its files (/ssh:HOST:/path) has to answer each request, and ssh to connect: past them the connection is closed and the work in hand fails, saying so.\\
\texttt{*removed-\codebreak{}sentence-\codebreak{}abbreviations*} & \texttt{nil} & Abbreviations of Tuatara’s (\textbf{SENTENCE-ABBREVIATIONS}) that should end a sentence after all, written as they are there: ’(“St.” “No.”).\\
\bottomrule
\end{xltabular}

\begin{xltabular}{\textwidth}{@{}l >{\raggedright\arraybackslash\hsize=0.711\hsize}X >{\raggedright\arraybackslash\hsize=1.289\hsize}X@{}}
\toprule
\textbf{Setting} & \textbf{Standard} & \textbf{What it changes}\\
\midrule
\endhead
\texttt{*rimu-\codebreak{}delete-\codebreak{}to-\codebreak{}trash*} & \texttt{t} & True for rimu’s deletions (M-x list-delete-marked, M-x list-execute) to move files and directories to the trash the desktop’s file managers use, from which they can be put back; NIL to delete them for good, after asking twice.\\
\texttt{*rimu-\codebreak{}details*} & \texttt{nil} & True for rimu to show the details of each entry before its name: its inode number, as ls -i does, then as ls -l does its permissions, links, owner, group, size and when it last changed, and after the name where a symbolic link leads.\\
\texttt{*rimu-\codebreak{}directories-\codebreak{}first*} & \texttt{t} & True for rimu to list directories, and the symbolic links that lead to directories, before everything else, as ls --group-directories-first does.\\
\texttt{*rimu-\codebreak{}directory-\codebreak{}sizes*} & \texttt{nil} & True for rimu’s lists to show each entry’s size before its name, a directory’s with everything under it; NIL, as standard, to show them only on M-x rimu-toggle-sizes in a list.\\
\texttt{*rimu-\codebreak{}dotfiles*} & \texttt{nil} & True for rimu to show the entries whose names start with a dot, .\\
\texttt{*rimu-\codebreak{}follow-\codebreak{}changes*} & \texttt{t} & True for rimu’s lists to be read again by themselves when their directory changes on disk: a file made, deleted, renamed or finished being written by another program, many at once as one; NIL to read one again only with M-x revert-buffer.\\
\texttt{*rimu-\codebreak{}git-\codebreak{}info*} & \texttt{nil} & True for rimu’s lists to show, after each name, the entry’s last commit in git’s history: when it was made and the first line of its message.\\
\texttt{*rimu-\codebreak{}hidden-\codebreak{}columns*} & \texttt{nil} & The columns rimu’s lists never show, however they are asked for: among the details :INODE, :PERMISSIONS, :LINKS, :OWNER, :GROUP, :SIZE and :TIME, and the columns of their own :SIZES (directories’ sizes) and :GIT (the last commits).\\
\texttt{*rimu-\codebreak{}ignored-\codebreak{}extensions*} & \texttt{(".o" ".so" ".a" ".fasl" ".elc" ".pyc" "...} & Endings of names rimu shows set back, in the :list-dim face: compiled files, backups (\textasciitilde{}) and the files some editors save as you type (\#).\\
\texttt{*rimu-\codebreak{}permissions*} & \texttt{:symbolic} & How rimu’s details write each entry’s permissions: :SYMBOLIC, as ls -l does, drwxr-xr-x; :OCTAL, its kind’s letter and the number chmod takes, d755, -644, -4755 with the set-user, set-group or sticky bit.\\
\texttt{*rimu-\codebreak{}preview*} & \texttt{t} & True for rimu lists to show the entry on the current line in a window beside the list, once point has rested on it \textbf{RIMU-PREVIEW-DELAY} seconds: the start of a file, a binary one in the hex view; the list of a directory; what a named pipe, a socket or a device is, never opening it.\\
\texttt{*rimu-\codebreak{}preview-\codebreak{}delay*} & \texttt{0.1} & Seconds point rests on an entry of a rimu list before the preview shows it, so that moving over many lines quickly reads none of them, 0.1 as standard.\\
\texttt{*rimu-\codebreak{}preview-\codebreak{}size*} & \texttt{1048576} & The most of a file rimu’s preview shows, in octets from its start: 1 MB by default.\\
\texttt{*rimu-\codebreak{}reuse-\codebreak{}lists*} & \texttt{t} & True for a folder a bookmark goes to (M-x bookmark-jump), or a list goes to across a project’s edge, to be listed in a rimu list there is of its project, moved there, the one used last (for a folder in no project, one in none): a list of its own only when there is none.\\
\texttt{*rimu-\codebreak{}sort*} & \texttt{:name} & The order rimu lists entries in: :NAME, by name, octet by octet as ls orders names in the C locale (capitals before small letters); :TIME, those changed last first.\\
\texttt{*safe-\codebreak{}directories*} & \texttt{nil} & Directories whose files always open safely, as M-x find-file-safely opens a file: read-only, with no major mode, hooks or version control, and nothing written beside them; for example ’(“\textasciitilde{}/Evidence/”).\\
\texttt{*sandbox-\codebreak{}input-\codebreak{}limit*} & \texttt{1073741824} & The largest file a program decoding files is given (1 GB): it gets a copy in memory, which a larger file, or one claiming a huge size with nothing in it, would fill.\\
\texttt{*sandbox-\codebreak{}memory*} & \texttt{1073741824} & The octets of memory a program decoding a file may take (1 GB); NIL for no limit.\\
\texttt{*sandbox-\codebreak{}program*} & \texttt{"bwrap"} & The program that shuts in the programs decoding files for the editor, and \LaTeX{} making PDFs: bubblewrap, named bwrap, looked for on the PATH (or a full path to it).\\
\texttt{*sandbox-\codebreak{}seccomp*} & \texttt{t} & True to have the kernel refuse a program decoding a file the system calls no decoder needs (see the file’s header), on x86-64 and AArch64; NIL to leave them to it, the rest of the sandbox kept.\\
\texttt{*sandbox-\codebreak{}seconds*} & \texttt{10} & The seconds a program decoding a file may run before it is ended, and the processor time it may take, 10 as standard.\\
\texttt{*sandbox-\codebreak{}total-\codebreak{}memory*} & \texttt{2147483648} & The octets of memory all the programs decoding one file may take together (2 GB), where the system can hold them to it (systemd-run, a scope of the user’s own); NIL for no such limit, each program still held to \textbf{SANDBOX-MEMORY}.\\
\texttt{*save-\codebreak{}in-\codebreak{}place*} & \texttt{nil} & True to save a large file by writing only what changed into the file itself, where that is quicker than writing it whole.\\
\texttt{*save-\codebreak{}place-\codebreak{}limit*} & \texttt{400} & How many files’ places are kept, the latest first, 400 as standard.\\
\texttt{*save-\codebreak{}place-\codebreak{}mode*} & \texttt{t} & True to remember where point was in each file you visit and put it back there when you visit the file again, in this session or a later one (on as standard); NIL to open files at their top.\\
\texttt{*scroll-\codebreak{}margin*} & \texttt{0} & How many rows a window keeps between point and its top and bottom edges: moving point nearer an edge scrolls the window, as long as there is text beyond that edge.\\
\texttt{*scroll-\codebreak{}recenter*} & \texttt{nil} & True to put point’s line in the middle of the window whenever point leaves the window or comes within \textbf{SCROLL-MARGIN} rows of its edge, as M-x recenter does; NIL to scroll only as far as keeps point in view, that many rows from the edge.\\
\texttt{*search-\codebreak{}files-\codebreak{}limit*} & \texttt{134217728} & The most octets of text kept in memory for a search of many files (M-x search-in-files, M-x search-in-formatted) and for each directory of \textbf{INDEXED-DIRECTORIES} (and on disk with it, \textbf{INDEXES-KEPT-ON-DISK}): 128 MB as standard.\\
\texttt{*search-\codebreak{}files-\codebreak{}lines-\codebreak{}limit*} & \texttt{2000} & The most matching lines search-in-files and search-in-formatted list, 2000 as standard.\\
\texttt{*search-\codebreak{}files-\codebreak{}min-\codebreak{}input*} & \texttt{3} & How many characters search-in-files and search-in-formatted wait for before searching.\\
\texttt{*search-\codebreak{}results-\codebreak{}save-\codebreak{}files*} & \texttt{t} & True to save the files that lines found by M-x search-in-files are written back into (M-x search-results-finish), those that were not open too, closed again after; NIL to leave each changed and not saved in a buffer of its own, for you to look at and save (M-x save-buffer).\\
\texttt{*searching-\codebreak{}indicator*} & \texttt{"S"} & What the mode line shows after \textbf{INDEXING-INDICATOR} while directories are kept read, always in the same place: in :mode-line-sign while a search runs over a directory kept read (M-x search-in-files and the like; at least \textbf{SIGN-LIT-SECONDS}), else faint, in :mode-line-sign-idle.\\
\texttt{*sentence-\codebreak{}end-\codebreak{}double-\codebreak{}space*} & \texttt{nil} & True for two spaces between sentences, as a typist writes them: a sentence then ends at .\\
\texttt{*set-\codebreak{}mark-\codebreak{}command-\codebreak{}repeat-\codebreak{}pop*} & \texttt{nil} & True for M-x set-mark-command pressed again right after M-x universal-argument M-x set-mark-command (or M-x pop-global-mark) to go on back to the mark before; NIL, as standard, for it to set the mark.\\
\texttt{*show-\codebreak{}completion-\codebreak{}count*} & \texttt{t} & True to show before the prompt of a completing read the place of the selected candidate and how many there are, as they change with what is typed: 3/120 (*/0 when nothing matches and the text typed can be taken, !/0 when it cannot); before a search’s, the match the cursor is at and how many there are.\\
\texttt{*show-\codebreak{}paren-\codebreak{}mode*} & \texttt{t} & True to show the bracket that pairs with the one at point: the one before point when it closes, else the one after it when it opens.\\
\texttt{*snippet-\codebreak{}directories*} & \texttt{nil} & The folders your snippets are in, each with a folder for each major mode (org-mode/, prog-mode/, fundamental-mode/), a snippet a file in it; NIL, as standard, for snippets/ in your configuration folder (\textasciitilde{}/.config/tuatara/snippets/, XDG\_CONFIG\_HOME’s when it is set).\\
\texttt{*snippet-\codebreak{}tab*} & \texttt{t} & True for a snippet’s key before the cursor and TAB to put the snippet in.\\
\texttt{*soft-\codebreak{}wrap-\codebreak{}at-\codebreak{}fill-\codebreak{}column*} & \texttt{t} & True to wrap long lines at the buffer’s fill column (M-x set-fill-column, \textbf{FILL-COLUMN}), or at the window’s edge when that is nearer; NIL to wrap at the window’s edge always.\\
\texttt{*soft-\codebreak{}wrap-\codebreak{}indicator*} & \texttt{"{\tuataraFallbackTWO{}↪}"} & What the gutter shows beside a row that goes on from the one above when lines wrap, in the :wrap-indicator face, where the line number would be; NIL for nothing.\\
\bottomrule
\end{xltabular}

\begin{xltabular}{\textwidth}{@{}l >{\raggedright\arraybackslash\hsize=0.930\hsize}X >{\raggedright\arraybackslash\hsize=1.070\hsize}X@{}}
\toprule
\textbf{Setting} & \textbf{Standard} & \textbf{What it changes}\\
\midrule
\endhead
\texttt{*soft-\codebreak{}wrap-\codebreak{}modes*} & \texttt{nil} & Major modes whose buffers wrap long lines at a word’s edge, at the fill column, or not: a list of (MODE .\\
\texttt{*sort-\codebreak{}fold-\codebreak{}case*} & \texttt{nil} & True to ignore the case of ASCII letters when sorting: “b” then sorts before “C”.\\
\texttt{*spellcheck-\codebreak{}comment-\codebreak{}faces*} & \texttt{(:comment)} & In code, the faces of comments, whose words are checked (unless \textbf{SPELLCHECK-MODES} says :strings): those named after them with a dot too (:comment.documentation).\\
\texttt{*spellcheck-\codebreak{}excluded-\codebreak{}faces*} & \texttt{(:org-\codebreak{}code :org-\codebreak{}verbatim :org-\codebreak{}link :org-...} & In text (Org, Markdown, \LaTeX{} and the like), the faces whose words are not checked: code, links, dates, tags, keywords.\\
\texttt{*spellcheck-\codebreak{}languages*} & \texttt{nil} & The languages words are checked in, a list of names as enchant knows them, such as ’(“en\_NZ” “sv\_SE”): a word is right when one of them takes it.\\
\texttt{*spellcheck-\codebreak{}modes*} & \texttt{nil} & Major modes whose buffers mark the words spelt wrong, or not: a list of (MODE .\\
\texttt{*spellcheck-\codebreak{}string-\codebreak{}faces*} & \texttt{(:string)} & In code, the faces of strings, whose words are checked (unless \textbf{SPELLCHECK-MODES} says :comments): those named after them with a dot too.\\
\texttt{*split-\codebreak{}height-\codebreak{}threshold*} & \texttt{80} & The fewest rows a window must have for a buffer shown in another window to split it into one above the other; a lone window is split so whenever it has room.\\
\texttt{*split-\codebreak{}limit*} & \texttt{100000} & The most places M-x split-selection holds, or NIL for as many as the memory has room for, 100000 as standard.\\
\texttt{*split-\codebreak{}limit-\codebreak{}ask*} & \texttt{nil} & True to be asked, when the places would pass their limit, whether to hold them all the same, the memory having room, rather than those in sight only.\\
\texttt{*split-\codebreak{}width-\codebreak{}threshold*} & \texttt{160} & The fewest columns a window must have for a buffer shown in another window to split it into two side by side.\\
\texttt{*startup-\codebreak{}buffer*} & \texttt{"*lisp-\codebreak{}scratch*"} & The buffer shown when the editor starts without files: “\textbf{lisp-scratch}” or “\textbf{org-scratch}” (or any buffer the init file makes).\\
\texttt{*state-\codebreak{}directory*} & \texttt{nil} & The folder Tuatara keeps what it remembers between sessions in, in place of \textasciitilde{}/.local/state/tuatara/ (\$XDG\_STATE\_HOME/tuatara/): the files visited lately, places in them, the session, auto-saved text, the packages fetched.\\
\texttt{*syntax-\codebreak{}helper-\codebreak{}unconfined*} & \texttt{nil} & True to have the grammars run, where there is no sandbox, for the files of projects you trust, in a program of their own that is not shut in.\\
\texttt{*tab-\codebreak{}always-\codebreak{}indent*} & \texttt{:complete} & What M-x indent-for-tab-command does where the major mode indents lines (Lisp, Rust, C, shell scripts): :COMPLETE (the standard) indents the current line, and when that changes nothing completes the name at point, as M-x completion-at-point does in Lisp.\\
\texttt{*tab-\codebreak{}width*} & \texttt{4} & The columns from one tab stop to the next: a tab character shows its text on at the next stop, and M-x indent-for-tab-command goes on to it.\\
\texttt{*terminal-\codebreak{}compress-\codebreak{}after*} & \texttt{10} & Seconds a terminal’s program has printed nothing after which its scrollback is compressed, in the background, a little at a time, so that a terminal with a long history holds a fraction of the memory; reading the history (its history view) undoes it as it goes, unseen.\\
\texttt{*terminal-\codebreak{}editor-\codebreak{}keys*} & \texttt{("C-x" "M-x" "C-g" "C-\codebreak{}\textbackslash{}\textbackslash{}" "C-q" "C-\codebreak{}s")} & Keys that stay the editor’s while you type into a terminal (M-x terminal).\\
\texttt{*terminal-\codebreak{}paste-\codebreak{}keys*} & \texttt{("C-V" "S-<insert>")} & Keys that paste into a terminal while you type there (M-x terminal-paste): the text copied last, in Tuatara or another program, marked as pasted for a program that asks (a shell then runs none of its lines until you press RET).\\
\texttt{*terminal-\codebreak{}picture-\codebreak{}memory*} & \texttt{67108864} & The most octets of pictures a terminal keeps for the programs in it that draw them; 0 takes none, programs then drawing with text.\\
\texttt{*terminal-\codebreak{}place*} & \texttt{:bottom} & Where a terminal shows (M-x terminal, M-x terminal-toggle): :bottom or :top, a window across that whole side of the screen \textbf{TERMINAL-SIZE} high; :left or :right, one down that whole side \textbf{TERMINAL-WIDTH} wide; or :full, the selected window.\\
\texttt{*terminal-\codebreak{}program*} & \texttt{nil} & The program M-x terminal runs, a list of it and its arguments, (“bash” ”-l”); NIL for your shell (\$SHELL, else /bin/sh).\\
\texttt{*terminal-\codebreak{}scrollback*} & \texttt{100000} & How many lines a terminal (M-x terminal) keeps that have gone up off its screen.\\
\texttt{*terminal-\codebreak{}send-\codebreak{}escape*} & \texttt{nil} & True to give Escape to a full-screen program in a terminal (one on the alternate screen: nvim, less, htop) even while \textbf{ESCAPE-QUITS} is on, Escape then still stopping typing at a shell’s prompt.\\
\texttt{*terminal-\codebreak{}size*} & \texttt{0.3} & How high a terminal’s window is at the bottom or top of the screen (\textbf{TERMINAL-PLACE}, M-x terminal-bottom, M-x terminal-top): a fraction of the screen’s height (0.3), or a number of rows (15).\\
\texttt{*terminal-\codebreak{}title-\codebreak{}limit*} & \texttt{40} & The most columns of a terminal’s title (as the shell gives it, user@host:\textasciitilde{}/dir in many) the mode line shows; NIL shows none.\\
\texttt{*terminal-\codebreak{}width*} & \texttt{0.5} & How wide a terminal’s window is at the left or right of the screen (\textbf{TERMINAL-PLACE}, M-x terminal-left, M-x terminal-right): a fraction of the screen’s width (0.5, half of it), or a number of columns (80).\\
\texttt{*text-\codebreak{}direction-\codebreak{}modes*} & \texttt{nil} & Major modes whose buffers show right-to-left text in the order it is read, or not: a list of (MODE .\\
\texttt{*text-\codebreak{}direction-\codebreak{}when-\codebreak{}safe*} & \texttt{nil} & True to show right-to-left text in the order it is read in a file opened safely (M-x find-file-safely) too; NIL to show it there in the order it is stored, unless the buffer chose otherwise with M-x text-direction-mode.\\
\texttt{*tuatara-\codebreak{}insert-\codebreak{}cursor-\codebreak{}colour*} & \texttt{nil} & The cursor’s colour in INSERT state, as faces take colours; NIL (as standard) for \textbf{CURSOR-COLOUR}.\\
\texttt{*tuatara-\codebreak{}insert-\codebreak{}cursor-\codebreak{}shape*} & \texttt{nil} & The cursor’s shape in INSERT state, while typing: :block, :bar or :underline; NIL (as standard) for \textbf{CURSOR-SHAPE}.\\
\texttt{*tuatara-\codebreak{}normal-\codebreak{}cursor-\codebreak{}colour*} & \texttt{nil} & The cursor’s colour in NORMAL state, as faces take colours; NIL (as standard) for \textbf{CURSOR-COLOUR}.\\
\texttt{*tuatara-\codebreak{}normal-\codebreak{}cursor-\codebreak{}shape*} & \texttt{nil} & The cursor’s shape in NORMAL state: :block, :bar or :underline; NIL (as standard) for \textbf{CURSOR-SHAPE}.\\
\texttt{*tuatara-\codebreak{}normal-\codebreak{}tab-\codebreak{}modes*} & \texttt{(org-\codebreak{}mode)} & The major modes, by name, where M-x tuatara-normal-tab in NORMAL state does what TAB does in the mode, the modes derived from them too; in the others it does nothing, as the keys that would type do nothing there.\\
\texttt{*tuatara-\codebreak{}support-\codebreak{}cursor-\codebreak{}colour*} & \texttt{nil} & The cursor’s colour in support buffers, as faces take colours; NIL (as standard) for \textbf{CURSOR-COLOUR}.\\
\texttt{*tuatara-\codebreak{}support-\codebreak{}cursor-\codebreak{}shape*} & \texttt{nil} & The cursor’s shape in support buffers (lists, help, \textbf{Messages}): :block, :bar or :underline; NIL (as standard) for \textbf{CURSOR-SHAPE}.\\
\texttt{*tui-\codebreak{}colours*} & \texttt{:auto} & How many colours Tuatara uses in a terminal.\\
\texttt{*tui-\codebreak{}mouse*} & \texttt{t} & Whether the mouse works in Tuatara in a terminal as in its window: the first button puts point, drags select (past the text, scrolling), twice selects a word and three times a line, the wheel scrolls, and a terminal’s program in M-x terminal gets them.\\
\texttt{*tui-\codebreak{}padding*} & \texttt{0} & The room, in whole cells, between the edges of the terminal and Tuatara’s text.\\
\texttt{*tui-\codebreak{}pictures*} & \texttt{:auto} & Whether pictures (a preview of an image, a PDF’s pages) are shown in the terminal.\\
\bottomrule
\end{xltabular}

\begin{xltabular}{\textwidth}{@{}l l >{\raggedright\arraybackslash\hsize=1.000\hsize}X@{}}
\toprule
\textbf{Setting} & \textbf{Standard} & \textbf{What it changes}\\
\midrule
\endhead
\texttt{*typing-\codebreak{}pause*} & \texttt{1/2} & Seconds without a key typed before work over a text too large to go through at each key starts while a pattern is typed: the count of a search’s matches (M-x regex-search-forward) and a filter being typed (M-x search-filter-regex), its lines in the window shown at once.\\
\texttt{*undo-\codebreak{}limit*} & \texttt{nil} & The most memory, in octets, the undo history of one buffer may hold; NIL, the default, for no limit.\\
\texttt{*use-\codebreak{}system-\codebreak{}clipboard*} & \texttt{t} & True to put what is killed or copied on the system clipboard as well, where other programs paste it from, and to paste what another program copied there since; NIL keeps killed text to the editor.\\
\texttt{*user-\codebreak{}full-\codebreak{}name*} & \texttt{nil} & Your full name, as “Ann Taranaki”: the author of an exported document with no \#+author:, \%author\% in a note template and \$AUTHOR in a snippet.\\
\texttt{*visual-\codebreak{}line-\codebreak{}modes*} & \texttt{nil} & Major modes whose buffers mark the fill column with a line down the window, or not: a list of (MODE .\\
\texttt{*window-\codebreak{}colours-\codebreak{}from-\codebreak{}desktop*} & \texttt{nil} & True to paint text and ground the theme leaves without a colour (the terminal and plain themes in a window) in the colours of the desktop’s GTK theme, following its light or dark style as it changes; NIL, as standard, for black text on a white ground, as a window with no theme shows it.\\
\texttt{*window-\codebreak{}min-\codebreak{}height*} & \texttt{4} & The fewest rows, its mode line included, a window may be left with when it is split or made smaller.\\
\texttt{*window-\codebreak{}min-\codebreak{}width*} & \texttt{10} & The fewest columns a window may be left with when it is split or made narrower.\\
\texttt{*wipe-\codebreak{}passes*} & \texttt{1} & How many times a file wiped is overwritten with random octets before it is deleted: once is enough for a disk that writes in place; more do not help one that does not (a copy-on-write file system, a solid-state drive).\\
\texttt{*word-\codebreak{}completion-\codebreak{}keep-\codebreak{}case*} & \texttt{t} & True to find words that start with what is typed whatever the case, and to keep the case typed: with “foo” typed, a word Foobar is put in as foobar, and with “FOO” typed, it is put in as FOOBAR.\\
\bottomrule
\end{xltabular}

\section{Every command}
\label{sec:577}
The name, keys and purpose of every command Tuatara has.

The earlier chapters describe the commands of each subject where you meet
them. This chapter lists them all in one table, in the alphabetical order of
their names, so that you can look one up by name. The table is made from the
program each time this manual is made, so it never lists a command that Tuatara
does not have.

The table has four columns. The first two are keys, as chapter 1 explains: the
modal keys, which work only with modal editing on, and the non-modal keys, which
work with it off too. Each shows up to three keys, the shortest first, or
“none”. A command with no key at all says \texttt{M-x} and its name in the second
column, because you can only run it by its name. The third column is the name
of the command and the fourth says what it does, in one sentence. The keys of a
command can differ from one major mode to another, and the table shows keys from
every mode. A key written \texttt{C-x} means hold down Control and press x.

To run a command by its name, press \texttt{M-x}, type a few letters of the name and
press \texttt{RET}. To find a command from within the editor, press \texttt{C-h f} and type
a few letters of its name. Tuatara shows its keys and its full description.
To find the command that a key runs, press \texttt{C-h k} and then the key. To find
the commands that mention a word, press \texttt{C-h a} and type the word.

{\small
\begin{xltabular}{\textwidth}{@{}l >{\raggedright\arraybackslash\hsize=0.696\hsize}X l >{\raggedright\arraybackslash\hsize=1.304\hsize}X@{}}
\toprule
\textbf{Modal} & \textbf{Non-modal} & \textbf{Command} & \textbf{What it does}\\
\midrule
\endhead
none & \texttt{C-h C-a} & \texttt{about-\codebreak{}tuatara} & Show \textbf{About Tuatara}, the page that introduces Tuatara.\\
\texttt{T} & \texttt{C-z T} & \texttt{add-\codebreak{}cursor-\codebreak{}above} & Add a cursor on the line above, at the same column: the first press makes cursors of point and the line above, each one after adds one above the first, and a count adds that many (negative, below).\\
\texttt{t} & \texttt{C-z t} & \texttt{add-\codebreak{}cursor-\codebreak{}below} & Add a cursor on the line below, at the same column.\\
none & \texttt{M-x align-\codebreak{}cursors} & \texttt{align-\codebreak{}cursors} & Put every cursor at the real cursor’s column, or its line’s end where the line is shorter; with cursors selecting, wrap the selections in = as any delimiter does.\\
none & \texttt{C-h a} & \texttt{apropos-\codebreak{}command} & List the commands and variables whose names or documentation match the words typed, in \textbf{Apropos}.\\
none & \texttt{M-m} & \texttt{back-\codebreak{}to-\codebreak{}indentation} & Move point to the first character of its line that is not a space or a tab.\\
none & \texttt{S-TAB} & \texttt{backward-\codebreak{}button} & Move point to the button before, from the first round to the last.\\
\texttt{h} & \texttt{C-b}, \texttt{<left>} & \texttt{backward-\codebreak{}char} & Move point backward one character (or the prefix argument’s worth); in the hex view, one octet.\\
none & \texttt{C-M-\codebreak{}DEL} & \texttt{backward-\codebreak{}kill-\codebreak{}sexp} & Kill the balanced expression before point: to where backward-sexp goes.\\
none & \texttt{C-M-p} & \texttt{backward-\codebreak{}list} & Move backward over the bracketed group before point, to its opening bracket; the opposite of forward-list.\\
none & \texttt{M-\{} & \texttt{backward-\codebreak{}paragraph} & Move point to the start of the paragraph: the start of the blank line before it, or the start of the text.\\
none & \texttt{M-a} & \texttt{backward-\codebreak{}sentence} & Move point to the start of the sentence: after the end of the one before, or the start of the paragraph’s text.\\
\texttt{B} & \texttt{C-M-b} & \texttt{backward-\codebreak{}sexp} & Move backward over a balanced expression, the opposite of forward-sexp: a quote before the expression goes with it.\\
none & \texttt{C-M-u} & \texttt{backward-\codebreak{}up-\codebreak{}list} & Move back out of the bracketed group around point, to its opening bracket.\\
\texttt{b} & \texttt{M-b} & \texttt{backward-\codebreak{}word} & Move point to the start of the previous word (or the prefix argument’s worth, forward with a negative one).\\
none & \texttt{C-x +} & \texttt{balance-\codebreak{}windows} & Make the windows the same size: each split shares its room equally.\\
\texttt{m DEL}, \texttt{m <delete>} & \texttt{C-z m DEL}, \texttt{C-z m <delete>} & \texttt{bank-\codebreak{}clear} & Let go of every banked line of this buffer; the text is left as it is.\\
\texttt{m l} & \texttt{C-z m l} & \texttt{bank-\codebreak{}lines} & Bank the lines the selection touches, or the line at point: set them aside, highlighted, and let go of the selection, to go on and bank more elsewhere.\\
\texttt{m u} & \texttt{C-z m u} & \texttt{bank-\codebreak{}release-\codebreak{}line} & Take the line at point out of the bank, the other banked lines kept; on a line not banked, say so.\\
\texttt{m U} & \texttt{C-z m U} & \texttt{bank-\codebreak{}release-\codebreak{}section} & Take the run of banked lines point is on out of the bank, all the lines beside each other at once; the banks elsewhere are kept.\\
\texttt{g g} & \texttt{M-<} & \texttt{beginning-\codebreak{}of-\codebreak{}buffer} & Move point to the start of the buffer.\\
none & \texttt{C-M-a} & \texttt{beginning-\codebreak{}of-\codebreak{}defun} & Move back to the start of the definition point is in, or of the one before; with a count, that many definitions back, or forward to the start of the ones after when negative.\\
\texttt{g h} & \texttt{C-a}, \texttt{<home>} & \texttt{beginning-\codebreak{}of-\codebreak{}line} & Move point to the start of the current line; where long lines wrap at a word’s edge (C-x T l), of the screen row it is on; in the hex view, of the row of octets.\\
none & \texttt{C-x m d} & \texttt{bookmark-\codebreak{}delete} & Forget a bookmark, chosen by its name.\\
none & \texttt{C-x m b} & \texttt{bookmark-\codebreak{}jump} & Go to a bookmark, chosen by its name (each listed with its file): its file visited, point where it was set, found again when the text above it changed.\\
none & \texttt{C-x m l} & \texttt{bookmark-\codebreak{}manager} & Manage the bookmarks in a list, as the buffer list (C-x C-b) manages buffers, in sections: one for each category (C-x m m asks for them, c here sets them), its heading saying how many it holds, then one for the bookmarks in none; a bookmark in two categories is in both.\\
none & \texttt{r}, \texttt{C-c r} & \texttt{bookmark-\codebreak{}rename} & Give the bookmark on the current line another name, in every section it is in; a name another bookmark has already is taken from it after asking.\\
none & \texttt{C-x m m} & \texttt{bookmark-\codebreak{}set} & Keep where point is in the current file (or the folder a directory list shows and the entry the cursor is on) under a name you give, the file’s name offered, then in the categories you choose, those used already offered, new ones typed (\textbf{BOOKMARK-SET-ASKS-CATEGORIES}).\\
none & \texttt{c}, \texttt{C-c c} & \texttt{bookmark-\codebreak{}set-\codebreak{}categories} & Say which categories the bookmark on the current line is in, or the marked ones (>) are, the sections holding them unfolded first: those used already offered, new ones typed, several apart by commas; the bookmark’s own written in first, to change.\\
none & \texttt{M-g o} & \texttt{buffer-\codebreak{}outline} & Go to a heading of the current buffer, chosen from a list of them: Org’s headings, Lisp’s comments of three semicolons and top-level forms, and in a language parsed with tree-sitter its definitions and headings (a function, a type, a Markdown heading), each shown as its line after its number, in the order of the text.\\
none & \texttt{C-x C-b} & \texttt{buffers} & List every live buffer in \textbf{Buffers}, those used most recently first, in a section for each project (a folder holding .git or .project, \textbf{PROJECT-ROOT-MARKERS}: its files, the folders of it listed and the buffers standing in them), the project of the buffer used last first, then one for those in none.\\
none & \texttt{K}, \texttt{C-c K} & \texttt{buffers-\codebreak{}keep-\codebreak{}client} & Make the client whose folder the cursor is in, in the buffer list, the one that stays when its last buffer is closed, in place of the one \textbf{CLIENT-KEPT-WITHOUT-BUFFERS} names, until it ends.\\
none & \texttt{s} & \texttt{buffers-\codebreak{}save} & Save the marked buffers (>) that have a file, or with none marked the buffer on the current line.\\
none & \texttt{C-x n a} & \texttt{calendar} & Show the calendar: the month of today and the months either side, below the window you are in, point on today.\\
none & \texttt{C-b}, \texttt{<left>} & \texttt{calendar-\codebreak{}backward-\codebreak{}day} & Move point a day back in the calendar; with a count, that many.\\
none & \texttt{M-\{} & \texttt{calendar-\codebreak{}backward-\codebreak{}month} & Move point a month back in the calendar, to the same day of the month (or its last); with a count, that many.\\
none & \texttt{C-p}, \texttt{<up>} & \texttt{calendar-\codebreak{}backward-\codebreak{}week} & Move point a week back in the calendar; with a count, that many.\\
none & \texttt{C-x [} & \texttt{calendar-\codebreak{}backward-\codebreak{}year} & Move point a year back in the calendar; with a count, that many.\\
none & \texttt{M-a} & \texttt{calendar-\codebreak{}beginning-\codebreak{}of-\codebreak{}month} & Move point to the first day of its month.\\
none & \texttt{C-a} & \texttt{calendar-\codebreak{}beginning-\codebreak{}of-\codebreak{}week} & Move point to the start of the week it is in (\textbf{calendar-week-start-day}); on it, a week back.\\
\bottomrule
\end{xltabular}
}

{\small
\begin{xltabular}{\textwidth}{@{}l >{\raggedright\arraybackslash\hsize=0.803\hsize}X l >{\raggedright\arraybackslash\hsize=1.197\hsize}X@{}}
\toprule
\textbf{Modal} & \textbf{Non-modal} & \textbf{Command} & \textbf{What it does}\\
\midrule
\endhead
none & \texttt{M-<} & \texttt{calendar-\codebreak{}beginning-\codebreak{}of-\codebreak{}year} & Move point to the first day of its year.\\
none & \texttt{M-e} & \texttt{calendar-\codebreak{}end-\codebreak{}of-\codebreak{}month} & Move point to the last day of its month.\\
none & \texttt{C-e} & \texttt{calendar-\codebreak{}end-\codebreak{}of-\codebreak{}week} & Move point to the end of the week it is in; on it, a week on.\\
none & \texttt{M->} & \texttt{calendar-\codebreak{}end-\codebreak{}of-\codebreak{}year} & Move point to the last day of its year.\\
none & \texttt{C-f}, \texttt{<right>} & \texttt{calendar-\codebreak{}forward-\codebreak{}day} & Move point a day on in the calendar; with a count, that many.\\
none & \texttt{M-\}} & \texttt{calendar-\codebreak{}forward-\codebreak{}month} & Move point a month on in the calendar, to the same day of the month (or its last); with a count, that many.\\
none & \texttt{C-n}, \texttt{<down>} & \texttt{calendar-\codebreak{}forward-\codebreak{}week} & Move point a week on in the calendar; with a count, that many.\\
none & \texttt{C-x ]} & \texttt{calendar-\codebreak{}forward-\codebreak{}year} & Move point a year on in the calendar; with a count, that many.\\
none & \texttt{g d} & \texttt{calendar-\codebreak{}goto-\codebreak{}date} & Move point to a date asked for: 2026-12-24, 2026-12, or as any date asked for is read (tomorrow, fri, +3d, dec 24).\\
none & \texttt{.}, \texttt{C-c .} & \texttt{calendar-\codebreak{}goto-\codebreak{}today} & Move point to today, the months around it shown.\\
none & \texttt{C-x n A} & \texttt{calendar-\codebreak{}here} & Show the calendar in the window you are in, in place of its buffer, point on today; q brings that buffer back.\\
none & \texttt{p c} & \texttt{calendar-\codebreak{}iso-\codebreak{}print-\codebreak{}date} & Say the ISO date of the day at point: its day of the week, its week and its year.\\
none & \texttt{n}, \texttt{C-c n} & \texttt{calendar-\codebreak{}journal-\codebreak{}entry} & Open the journal’s entry for the calendar’s day at point, or make it: the calendar leaves its window, the entry shown where you were.\\
none & \texttt{M-x calendar-\codebreak{}jump-\codebreak{}back} & \texttt{calendar-\codebreak{}jump-\codebreak{}back} & Go back to the day point was on before it last moved; again, and back to the day after once more.\\
none & \texttt{m}, \texttt{C-c m} & \texttt{calendar-\codebreak{}mark-\codebreak{}days} & Mark the days with agenda entries in the calendar (as standard, with \textbf{calendar-mark-days}).\\
none & \texttt{RET} & \texttt{calendar-\codebreak{}open-\codebreak{}day} & Open what the day at point is marked for: its agenda, its journal entry, whichever it has (\textbf{CALENDAR-DAY-OPENERS}); with more than one, asked which, each by the key that opens it alone.\\
none & \texttt{o}, \texttt{C-c o} & \texttt{calendar-\codebreak{}other-\codebreak{}month} & Show the months around a month asked for, point on its first day: jan ’27 (or jan 27, january 2027) for January 2027, jan for this year’s, 2027-03, or a date (2027-03-14, fri, +3m) for its month.\\
none & \texttt{p d} & \texttt{calendar-\codebreak{}print-\codebreak{}day-\codebreak{}of-\codebreak{}year} & Say which day of its year the day at point is.\\
none & \texttt{>}, \texttt{C-c >}, \texttt{C-x >} & \texttt{calendar-\codebreak{}scroll-\codebreak{}left} & Show the next month, the months moving left; with a count, that many.\\
none & \texttt{C-v}, \texttt{<next>} & \texttt{calendar-\codebreak{}scroll-\codebreak{}left-\codebreak{}three-\codebreak{}months} & Show the months after those shown, as many as are shown (\textbf{calendar-months}, three as standard); with a count, that many times as many.\\
none & \texttt{<}, \texttt{C-c <}, \texttt{C-x <} & \texttt{calendar-\codebreak{}scroll-\codebreak{}right} & Show the month before, the months moving right; with a count, that many.\\
none & \texttt{M-v}, \texttt{<prior>} & \texttt{calendar-\codebreak{}scroll-\codebreak{}right-\codebreak{}three-\codebreak{}months} & Show the months before those shown, as many as are shown (\textbf{calendar-months}, three as standard); with a count, that many times as many.\\
none & \texttt{u}, \texttt{C-c u} & \texttt{calendar-\codebreak{}unmark} & Take the marks off the calendar’s days, until m puts them back.\\
none & \texttt{M-x capitalize-\codebreak{}dwim} & \texttt{capitalize-\codebreak{}dwim} & Capitalise the region, or with none the next word, as capitalize-word does.\\
none & \texttt{M-x capitalize-\codebreak{}region} & \texttt{capitalize-\codebreak{}region} & Capitalise each word of the region.\\
none & \texttt{M-c} & \texttt{capitalize-\codebreak{}word} & Capitalise the next word, or the count’s: its first letter in upper case and the rest in lower.\\
none & \texttt{C-x T o} & \texttt{centred-\codebreak{}text-\codebreak{}mode} & Show the text of the current buffer in a column of \textbf{CENTRED-TEXT-WIDTH} columns in the middle of the window, blank margins either side, or as it was: for writing and reading with nothing else about.\\
none & \texttt{M-x check-\codebreak{}health} & \texttt{check-\codebreak{}health} & Show what each language and major mode has here.\\
none & \texttt{C-x r c} & \texttt{clear-\codebreak{}rectangle} & Blank out the rectangle between the mark and point: its columns become spaces on each line, so that what comes after stays where it was; a line that ends inside the block is cut at its left column instead, leaving no spaces at its end, and a line shorter than the block is left alone.\\
none & \texttt{M-;} & \texttt{comment-\codebreak{}dwim} & Comment what suits where point is (M-;): with a region, its lines, or uncomment them when they all are (as M-x toggle-comment-lines-ret); on a line with a comment, go into it; on a blank line, begin a comment; after code, begin one at the line’s end, at \textbf{COMMENT-COLUMN}.\\
none & \texttt{M-x compile} & \texttt{compile} & Run a command asked for, a build or the tests (make -k, cargo build, pytest), with your shell from the root of the project the current buffer’s file is in (see \textbf{PROJECT-ROOT-MARKERS}), or from its directory when it is in none, and show what it prints in \textbf{compile} as it comes, the editor free meanwhile.\\
none & \texttt{C-c C-k} & \texttt{compile-\codebreak{}stop} & Stop the command running in this compile’s output, and every program it started: what it printed so far stays.\\
none & \texttt{C-M-/} & \texttt{complete-\codebreak{}word} & List the words in the buffers that start with the text before point beside the cursor, the nearest first, each with the buffer it is in when that is not this one; M-/ puts them in one after another.\\
none & \texttt{C-M-i}, \texttt{M-TAB} & \texttt{completion-\codebreak{}at-\codebreak{}point} & Complete the word at point from a list beside it.\\
none & \texttt{M-x conflict-\codebreak{}keep-\codebreak{}base} & \texttt{conflict-\codebreak{}keep-\codebreak{}base} & Put the conflict point is in right by keeping its base, the lines after its \(\vert{}\)\(\vert{}\)\(\vert{}\)\(\vert{}\)\(\vert{}\)\(\vert{}\)\(\vert{}\) line, as they were before either side changed them, the sides and the marker lines taken away; then go to the next conflict.\\
none & \texttt{M-x conflict-\codebreak{}keep-\codebreak{}both-\codebreak{}sides} & \texttt{conflict-\codebreak{}keep-\codebreak{}both-\codebreak{}sides} & Put the conflict point is in right by keeping both its sides, the first then the second, the base and the marker lines taken away; then go to the next conflict.\\
none & \texttt{M-x conflict-\codebreak{}keep-\codebreak{}first-\codebreak{}side} & \texttt{conflict-\codebreak{}keep-\codebreak{}first-\codebreak{}side} & Put the conflict point is in right by keeping its first side, the lines after its <<<<<<< line, and taking away the rest and its marker lines; then go to the next conflict.\\
none & \texttt{M-x conflict-\codebreak{}keep-\codebreak{}second-\codebreak{}side} & \texttt{conflict-\codebreak{}keep-\codebreak{}second-\codebreak{}side} & Put the conflict point is in right by keeping its second side, the lines after its \texttt{=====} line, and taking away the rest and its marker lines; then go to the next conflict.\\
none & \texttt{M-x conflict-\codebreak{}markers-\codebreak{}mode} & \texttt{conflict-\codebreak{}markers-\codebreak{}mode} & Show the conflicts a merge left in the current buffer, and give it the keys that put them right, or stop.\\
none & \texttt{M-x conflict-\codebreak{}next} & \texttt{conflict-\codebreak{}next} & Go to the start of the next conflict, its <<<<<<< line, and say which of how many it is.\\
\bottomrule
\end{xltabular}
}

\begin{xltabular}{\textwidth}{@{}l >{\raggedright\arraybackslash\hsize=0.744\hsize}X l >{\raggedright\arraybackslash\hsize=1.256\hsize}X@{}}
\toprule
\textbf{Modal} & \textbf{Non-modal} & \textbf{Command} & \textbf{What it does}\\
\midrule
\endhead
none & \texttt{M-x conflict-\codebreak{}previous} & \texttt{conflict-\codebreak{}previous} & Go to the start of the conflict before point, its <<<<<<< line, and say which of how many it is.\\
none & \texttt{M-x conflict-\codebreak{}show-\codebreak{}sides-\codebreak{}diff} & \texttt{conflict-\codebreak{}show-\codebreak{}sides-\codebreak{}diff} & Show how the two sides of the conflict point is in differ, in a buffer of their own in another window, which q closes: the lines only the first side has after -, those only the second has after +, three lines around each change, the sides named by their marker lines.\\
none & \texttt{M-x copy-\codebreak{}from-\codebreak{}evidence} & \texttt{copy-\codebreak{}from-\codebreak{}evidence} & Copy files out of an evidence folder (\textbf{EVIDENCE-FOLDERS}) to a folder asked for, which may not be in one: in a rimu list the marked files, or the one on the line; elsewhere the file at point or the buffer’s file, each of them in an evidence folder.\\
none & \texttt{C-x r M-w} & \texttt{copy-\codebreak{}rectangle-\codebreak{}as-\codebreak{}kill} & Copy the rectangle between the mark and point to the rectangle store that C-x r y pastes from, the text staying as it is.\\
\texttt{y} & none & \texttt{copy-\codebreak{}selection-\codebreak{}or-\codebreak{}char} & Copy the selection, or COUNT characters at point when nothing is selected.\\
none & \texttt{M-x copy-\codebreak{}snippet} & \texttt{copy-\codebreak{}snippet} & Copy one of Tuatara’s own snippets, chosen by its name, into your snippets folder (the first of \textbf{snippet-directories}, in the folder of its mode) and visit it: your copy, changed as you like, takes the place of Tuatara’s.\\
none & \texttt{M-=} & \texttt{count-\codebreak{}words} & Say how many lines, sentences, words and characters the selection has, or the whole buffer without one (the part shown of a narrowed one): “Region has 2 lines, 2 sentences, 9 words, and 48 characters”.\\
none & \texttt{M-x credits} & \texttt{credits} & Show where Tuatara’s ideas come from: the projects it owes, and what it took from each.\\
none & \texttt{C-c C-s} & \texttt{csv-\codebreak{}define-\codebreak{}separator} & Ask for the character that separates the fields of this buffer and show them separated by it.\\
none & \texttt{C-x T c} & \texttt{cursor-\codebreak{}column-\codebreak{}mode} & Highlight the column point is in, down the window, in the current buffer, or stop: in the :cursor-column face (the :hl-line face’s colour unless set apart), over the backgrounds of the text (a source block’s, the selection’s) and under the line point is on, past the end of short lines too, in the selected window.\\
none & \texttt{M-x cursors-\codebreak{}append} & \texttt{cursors-\codebreak{}append} & Type after the character under every cursor, or after its selection; a cursor at its line’s end types there.\\
none & \texttt{M-x cursors-\codebreak{}append-\codebreak{}line} & \texttt{cursors-\codebreak{}append-\codebreak{}line} & Type at the end of every cursor’s line.\\
none & \texttt{M-x cursors-\codebreak{}bank} & \texttt{cursors-\codebreak{}bank} & Bank every cursor’s line, or take them all out where every one is banked; the bank is there for y, d and f once the cursors end.\\
none & \texttt{M-x cursors-\codebreak{}capitalise} & \texttt{cursors-\codebreak{}capitalise} & Capitalise every cursor’s selection, or the count’s words from every cursor.\\
none & \texttt{M-x cursors-\codebreak{}change} & \texttt{cursors-\codebreak{}change} & Empty every cursor’s selection, or the count’s characters at every cursor, then type there; a selection goes on the kill ring as one kill.\\
none & \texttt{M-x cursors-\codebreak{}comment} & \texttt{cursors-\codebreak{}comment} & Comment or uncomment every cursor’s line.\\
none & \texttt{M-x cursors-\codebreak{}copy} & \texttt{cursors-\codebreak{}copy} & Copy every cursor’s selection, or the count’s characters at every cursor, as one kill, one line a cursor; the selections are let go.\\
none & \texttt{M-x cursors-\codebreak{}delete} & \texttt{cursors-\codebreak{}delete} & Delete every cursor’s selection, a selected line with its newline, or the count’s characters at every cursor; one change, the cursors kept.\\
none & \texttt{M-x cursors-\codebreak{}downcase} & \texttt{cursors-\codebreak{}downcase} & Put every cursor’s selection in lower case, or the count’s words from every cursor.\\
none & \texttt{M-x cursors-\codebreak{}indent} & \texttt{cursors-\codebreak{}indent} & Indent every cursor’s line as the mode indents it.\\
none & \texttt{M-x cursors-\codebreak{}insert} & \texttt{cursors-\codebreak{}insert} & Type before every cursor, or its selection: what is typed at one appears at all.\\
none & \texttt{M-x cursors-\codebreak{}insert-\codebreak{}line} & \texttt{cursors-\codebreak{}insert-\codebreak{}line} & Type at the start of every cursor’s line.\\
none & \texttt{M-x cursors-\codebreak{}kill-\codebreak{}line} & \texttt{cursors-\codebreak{}kill-\codebreak{}line} & Kill from every cursor to the end of its line, as one kill and one change; a cursor at its line’s end kills nothing.\\
none & \texttt{M-x cursors-\codebreak{}line-\codebreak{}end} & \texttt{cursors-\codebreak{}line-\codebreak{}end} & Put every cursor at the end of its line, keeping its selection.\\
none & \texttt{M-x cursors-\codebreak{}move-\codebreak{}past-\codebreak{}limit} & \texttt{cursors-\codebreak{}move-\codebreak{}past-\codebreak{}limit} & Move cursors past their limit, those out of sight with those in sight, by the NORMAL state motion of the keys pressed: a character right or left (with a count, that many), or to the start or the end of the line.\\
none & \texttt{M-x cursors-\codebreak{}open-\codebreak{}above} & \texttt{cursors-\codebreak{}open-\codebreak{}above} & Open a line above every cursor’s line and type on all of them.\\
none & \texttt{M-x cursors-\codebreak{}open-\codebreak{}below} & \texttt{cursors-\codebreak{}open-\codebreak{}below} & Open a line below every cursor’s line and type on all of them.\\
none & \texttt{M-x cursors-\codebreak{}redo} & \texttt{cursors-\codebreak{}redo} & Redo, keeping the cursors.\\
none & \texttt{M-x cursors-\codebreak{}repeat} & \texttt{cursors-\codebreak{}repeat} & Run the cursors’ last command again at every cursor, with the count it had: a motion, a selection, or deleting, pasting, killing and banking.\\
none & \texttt{M-x cursors-\codebreak{}replay} & \texttt{cursors-\codebreak{}replay} & Run the command of the keys just pressed at every cursor: the NORMAL state motion or selection they are on, or in a mark run, the run’s.\\
none & \texttt{M-x cursors-\codebreak{}run-\codebreak{}step-\codebreak{}back} & \texttt{cursors-\codebreak{}run-\codebreak{}step-\codebreak{}back} & Put every cursor’s selection back where the run’s last press found it.\\
none & \texttt{M-x cursors-\codebreak{}run-\codebreak{}step-\codebreak{}forward} & \texttt{cursors-\codebreak{}run-\codebreak{}step-\codebreak{}forward} & Put every cursor’s selection forward again, a step u took back.\\
none & \texttt{M-x cursors-\codebreak{}run-\codebreak{}toggle} & \texttt{cursors-\codebreak{}run-\codebreak{}toggle} & Start a mark run at every cursor, or end it and let go of the selections: each cursor keeps what it selected, or selects its word, and the run’s keys w W b B s S grow and h l g h g l * adjust every selection at once; u and U step them back and forward.\\
none & \texttt{M-x cursors-\codebreak{}select} & \texttt{cursors-\codebreak{}select} & Start a selection at every cursor, where each is: the motions then grow them.\\
none & \texttt{M-x cursors-\codebreak{}select-\codebreak{}lines} & \texttt{cursors-\codebreak{}select-\codebreak{}lines} & Select every cursor’s whole line, or let go where every cursor has.\\
none & \texttt{M-x cursors-\codebreak{}unbank} & \texttt{cursors-\codebreak{}unbank} & Take every cursor’s line out of the bank, leaving every other bank alone.\\
none & \texttt{M-x cursors-\codebreak{}undo} & \texttt{cursors-\codebreak{}undo} & Undo, keeping the cursors: a change made at them all is taken back at all.\\
none & \texttt{M-x cursors-\codebreak{}upcase} & \texttt{cursors-\codebreak{}upcase} & Put every cursor’s selection in upper case, or the count’s words from every cursor.\\
none & \texttt{M-x cursors-\codebreak{}yank} & \texttt{cursors-\codebreak{}yank} & Paste at every cursor, over its selection.\\
none & \texttt{M-SPC} & \texttt{cycle-\codebreak{}spacing} & Change the spaces around point in turn: once, one space (M-x just-one-space, a count as there); pressed again, none; a third time, the spaces as they were.\\
\bottomrule
\end{xltabular}

{\small
\begin{xltabular}{\textwidth}{@{}l >{\raggedright\arraybackslash\hsize=0.755\hsize}X l >{\raggedright\arraybackslash\hsize=1.245\hsize}X@{}}
\toprule
\textbf{Modal} & \textbf{Non-modal} & \textbf{Command} & \textbf{What it does}\\
\midrule
\endhead
none & \texttt{M-x daemon-\codebreak{}quit} & \texttt{daemon-\codebreak{}quit} & Stop the editor running in the background (tuatara --daemon), from any of its clients, asked first about every change not saved in it, the clients that ended included, folder by folder: s saves those of a folder, l leaves them unsaved, C-g keeps the daemon running; the saves begun are waited for (C-g there keeps it running too).\\
none & \texttt{C-x 8 -} & \texttt{deactivate-\codebreak{}input-\codebreak{}method} & Turn off this buffer’s input method, whichever it is.\\
none & \texttt{M-x debug-\codebreak{}mode} & \texttt{debug-\codebreak{}mode} & Switch debug mode on or off.\\
none & \texttt{DEL} & \texttt{delete-\codebreak{}backward-\codebreak{}char} & Delete the character before point (with a count, that many): of a letter with accents, the last accent alone, of an emoji with a skin tone, the tone; of a carriage return and a newline, both, which end one line.\\
none & \texttt{C-x C-o} & \texttt{delete-\codebreak{}blank-\codebreak{}lines} & Delete the blank lines about point: on a blank line, all the blank lines around it but one; on a blank line alone, that one; on a line with text, the blank lines after it.\\
none & \texttt{C-d}, \texttt{<delete>} & \texttt{delete-\codebreak{}char} & Delete the character after point (with a count, that many; a negative one, before point), a letter with its accents or an emoji with its skin tone whole (CLUSTER-FORWARD), a carriage return and a newline both (TUATARA.BUFFER:WHOLE-LINE-ENDS).\\
none & \texttt{M-x delete-\codebreak{}duplicate-\codebreak{}lines} & \texttt{delete-\codebreak{}duplicate-\codebreak{}lines} & Delete all but one of each set of identical lines the selection touches: the first of them is kept, with C-u the last; with C-u C-u only a line identical to the one before it goes, which is quicker for sorted text; with C-u C-u C-u blank lines are all kept.\\
none & \texttt{C-x 5 0} & \texttt{delete-\codebreak{}frame} & Close the frame worked in, with its windows; its buffers stay, shown in the other frames as before.\\
none & \texttt{M-\textbackslash{}} & \texttt{delete-\codebreak{}horizontal-\codebreak{}space} & Delete the spaces and tabs around point; with C-u, only those before it.\\
none & \texttt{C-x 5 1} & \texttt{delete-\codebreak{}other-\codebreak{}frames} & Close every frame but the one worked in; in a client of an editor running in the background, every other frame of that client.\\
\texttt{F} & \texttt{C-x 1} & \texttt{delete-\codebreak{}other-\codebreak{}windows} & Make the selected window fill the screen, closing the others.\\
none & \texttt{M-x delete-\codebreak{}pair} & \texttt{delete-\codebreak{}pair} & Delete the brackets or quote marks around the balanced expression after point: the character there, after any spaces and tabs, and the one that closes the expression it begins.\\
none & \texttt{C-x r d} & \texttt{delete-\codebreak{}rectangle} & Delete the rectangle between the mark and point, keeping it nowhere: the rectangle store keeps what it had.\\
\texttt{d} & none & \texttt{delete-\codebreak{}selection-\codebreak{}or-\codebreak{}char} & Delete the selection into the kill ring; without one, delete COUNT characters at point, which are not saved: a character under the cursor is a typo being fixed, not a cut.\\
none & \texttt{M-x delete-\codebreak{}trailing-\codebreak{}whitespace} & \texttt{delete-\codebreak{}trailing-\codebreak{}whitespace} & Delete the spaces and tabs at the end of every line, of the selection when there is one, else of the whole buffer (of the part shown of a narrowed one), and then the blank lines at its end too (\textbf{DELETE-TRAILING-LINES}).\\
none & \texttt{C-x 0} & \texttt{delete-\codebreak{}window} & Close the selected window: its room goes to the window beside it, and the window used before it is selected.\\
\texttt{?} & \texttt{C-h b} & \texttt{describe-\codebreak{}bindings} & List every key binding in effect, keymap by keymap, each followed by the keymaps it falls through to.\\
none & \texttt{M-x describe-\codebreak{}buffer-\codebreak{}places} & \texttt{describe-\codebreak{}buffer-\codebreak{}places} & Show where buffers are shown, the rules place-buffers makes, the one made last first: which buffers, the place, the size, and whether the window is kept for them and selected.\\
none & \texttt{M-x describe-\codebreak{}char} & \texttt{describe-\codebreak{}char} & Show in \textbf{Help} all that is known of the character at point: where it is, its code in decimal, octal and hexadecimal, its Unicode name and kind, its octets and the columns it takes, the face it shows in and the major mode.\\
none & \texttt{C-h x} & \texttt{describe-\codebreak{}command} & Show what a command does, which keys run it and where it is defined, the file a link that shows the definition.\\
none & \texttt{C-h C-c} & \texttt{describe-\codebreak{}copying} & Show the license Tuatara comes under.\\
none & \texttt{M-x describe-\codebreak{}evidence} & \texttt{describe-\codebreak{}evidence} & Say whether the file or folder here is evidence, and by which entry of \textbf{EVIDENCE-FOLDERS}: a rimu list’s entry on the line, a path written at point, or the buffer’s file; a link is evidence when what it leads to is.\\
none & \texttt{M-x describe-\codebreak{}export-\codebreak{}settings} & \texttt{describe-\codebreak{}export-\codebreak{}settings} & Show the settings an :EXPORT: drawer may hold in an Org note, what each sets and what it takes: how the enhanced export sets a heading, a table, a list, a block or a picture, in words rather than \LaTeX{}.\\
none & \texttt{M-x describe-\codebreak{}face} & \texttt{describe-\codebreak{}face} & Show all that is known of a face: a sample drawn in it, what it colours, its definition, what the theme in force and your init file set for it, the face it inherits from and each of its attributes as it is drawn now.\\
none & \texttt{C-h f} & \texttt{describe-\codebreak{}function} & Show what a function, macro or command does: its parameters, its documentation, the keys that run it when it is a command, and where it is defined, the file a link that shows the definition.\\
none & \texttt{C-h I} & \texttt{describe-\codebreak{}input-\codebreak{}method} & Show an input method’s rules, the keys typed in turn and what they give (C-h I): this buffer’s, else one asked for.\\
none & \texttt{C-h k} & \texttt{describe-\codebreak{}key} & Read a key sequence and show the command it runs and that command’s documentation, the question taken away once the keys are read.\\
none & \texttt{C-h C-w} & \texttt{describe-\codebreak{}no-\codebreak{}warranty} & Show the part of the license that says Tuatara comes with no warranty.\\
none & \texttt{C-h v} & \texttt{describe-\codebreak{}variable} & Show a variable’s value and documentation, and where it is defined, the file a link that shows the definition.\\
none & \texttt{M-x diff-\codebreak{}files} & \texttt{diff-\codebreak{}files} & Compare the file shown with another asked for, navigating to it as C-x C-f does, and show the differences in a buffer of their own, which q closes: the entry on the line in a list of a directory (there = compares the two marked); with C-u, or where there is none, both asked for.\\
none & \texttt{C-x 8 .} & \texttt{digraph-\codebreak{}input-\codebreak{}mode} & Type digraphs as you go: turn on the rfc1345 input method, so that typed text takes \& and a digraph’s keys for its character, \&e’ for e with an acute; \& and keys that make none stay as typed.\\
none & \texttt{M-x display-\codebreak{}line-\codebreak{}numbers-\codebreak{}mode} & \texttt{display-\codebreak{}line-\codebreak{}numbers-\codebreak{}mode} & Show absolute line numbers in the gutter of the current buffer, or none: each line its own number, from 1.\\
none & \texttt{M-x display-\codebreak{}relative-\codebreak{}line-\codebreak{}numbers-\codebreak{}mode} & \texttt{display-\codebreak{}relative-\codebreak{}line-\codebreak{}numbers-\codebreak{}mode} & Show relative line numbers in the gutter of the current buffer, or none: the current line shows its number, the others how many lines away they are, which is what counted motion like C-u 7 j expects; with a count, on when it is positive, off when it is zero or less.\\
none & \texttt{C-M-d} & \texttt{down-\codebreak{}list} & Move forward into the next bracketed group, just past its opening bracket, the words, strings and comments before it passed over.\\
none & \texttt{M-x downcase-\codebreak{}dwim} & \texttt{downcase-\codebreak{}dwim} & Put the region in lower case, or with none the next word, as downcase-word does.\\
none & \texttt{C-x C-l} & \texttt{downcase-\codebreak{}region} & Put the region in lower case.\\
none & \texttt{M-l} & \texttt{downcase-\codebreak{}word} & Put the next word, or the count’s, in lower case and move over it; with a negative count, the words before point, point staying.\\
none & \texttt{M-x drop-\codebreak{}last-\codebreak{}cursor} & \texttt{drop-\codebreak{}last-\codebreak{}cursor} & Take back the cursor added last, keeping the others; the real cursor is never taken back, and C-g ends the cursors.\\
none & \texttt{M-x electric-\codebreak{}pair-\codebreak{}mode} & \texttt{electric-\codebreak{}pair-\codebreak{}mode} & Turn closing delimiters as they are typed on or off.\\
\texttt{g e} & \texttt{M->} & \texttt{end-\codebreak{}of-\codebreak{}buffer} & Move point to the end of the buffer.\\
\bottomrule
\end{xltabular}
}

{\small
\begin{xltabular}{\textwidth}{@{}l >{\raggedright\arraybackslash\hsize=0.794\hsize}X l >{\raggedright\arraybackslash\hsize=1.206\hsize}X@{}}
\toprule
\textbf{Modal} & \textbf{Non-modal} & \textbf{Command} & \textbf{What it does}\\
\midrule
\endhead
none & \texttt{C-M-e} & \texttt{end-\codebreak{}of-\codebreak{}defun} & Move forward to the end of the definition point is in, or of the one after, and on to the start of the next line when only blanks or a comment follow it there; with a count, that many definitions on, or back to the ends of those before when negative.\\
\texttt{g l} & \texttt{C-e}, \texttt{<end>} & \texttt{end-\codebreak{}of-\codebreak{}line} & Move point to the end of the current line, before its newline; where long lines wrap at a word’s edge (C-x T l), of the screen row it is on; in the hex view, of the row of octets.\\
none & \texttt{M-x end-\codebreak{}safe-\codebreak{}visit} & \texttt{end-\codebreak{}safe-\codebreak{}visit} & Open the file of this buffer, opened safely, as usual from now on, where it is.\\
none & \texttt{C-x \textasciicircum{}} & \texttt{enlarge-\codebreak{}window} & Make the selected window a row taller, or the prefix argument’s rows; they come from the windows below it, then above.\\
none & \texttt{C-x \}} & \texttt{enlarge-\codebreak{}window-\codebreak{}horizontally} & Make the selected window a column wider, or the prefix argument’s columns; they come from the windows to its right, then its left.\\
none & \texttt{M-x eval-\codebreak{}buffer} & \texttt{eval-\codebreak{}buffer} & Evaluate every Lisp form in the current buffer one after another, as a file’s are, and show the last one’s value; not a buffer of more than \textbf{REWRITE-THRESHOLD} octets.\\
none & \texttt{M-:} & \texttt{eval-\codebreak{}expression} & Read a Lisp expression in the minibuffer and evaluate it in the running editor, in the current buffer’s package; show its value, or with a prefix argument (C-u) insert it at point.\\
none & \texttt{C-x C-e} & \texttt{eval-\codebreak{}last-\codebreak{}expression} & Evaluate the expression before point in the running editor and show its value in the echo area; with a prefix argument (C-u), insert the value after point.\\
none & \texttt{C-j} & \texttt{eval-\codebreak{}print-\codebreak{}last-\codebreak{}expression} & Evaluate the expression before point in the running editor and put what it printed and its value on the lines after it.\\
none & \texttt{M-x eval-\codebreak{}region} & \texttt{eval-\codebreak{}region} & Evaluate the Lisp forms in the selection one after another, as a file’s are (IN-PACKAGE carries on to the next), and show the last one’s value.\\
none & \texttt{C-M-x} & \texttt{eval-\codebreak{}top-\codebreak{}level-\codebreak{}form} & Evaluate the top-level form around point (the one starting with an open bracket at the start of a line) in the running editor and show its value.\\
none & \texttt{C-x C-x} & \texttt{exchange-\codebreak{}point-\codebreak{}and-\codebreak{}mark} & Swap point and mark.\\
none & \texttt{M-x} & \texttt{execute-\codebreak{}extended-\codebreak{}command} & Run a command chosen by name, with completion over every command: the candidates show each command’s keys and the first line of its documentation.\\
none & \texttt{M-o} & \texttt{expand-\codebreak{}selection} & Select the smallest piece of syntax around the selection, or around the character at point: the word, then the expression holding it, then the statement, and so on out; in a buffer with no syntax tree, as its language server sees the code.\\
none & \texttt{M-/} & \texttt{expand-\codebreak{}word} & Put in place of the start of the word before point the nearest word that starts with it, looked for in this buffer, then in the others (those shown first).\\
none & \texttt{M-x export-\codebreak{}check-\codebreak{}redactions} & \texttt{export-\codebreak{}check-\codebreak{}redactions} & Look for the texts this Org or \LaTeX{} file hides shown plainly elsewhere in it, and list each place in \textbf{Redactions}, where l shows it.\\
none & \texttt{M-x export-\codebreak{}latex-\codebreak{}compile-\codebreak{}pdf} & \texttt{export-\codebreak{}latex-\codebreak{}compile-\codebreak{}pdf} & Make the \LaTeX{} file of the current buffer into a PDF beside it, in the background, saving it first when it has changes, by xelatex, or by lualatex when xelatex fails (\textbf{latex-engines}), or first by the one its first lines name (“\% Intended \LaTeX{} compiler: lualatex”, “\%!TEX program = lualatex”), run as often as its cross-references need.\\
none & \texttt{M-x export-\codebreak{}latex-\codebreak{}compile-\codebreak{}pdf-\codebreak{}buffer} & \texttt{export-\codebreak{}latex-\codebreak{}compile-\codebreak{}pdf-\codebreak{}buffer} & Make the \LaTeX{} file of the current buffer into a PDF beside it, as M-x export-latex-compile-pdf does, and show the PDF: down the right side of the screen when it is wide enough (\textbf{split-width-threshold} columns), else alone on it.\\
none & \texttt{M-x export-\codebreak{}latex-\codebreak{}compile-\codebreak{}pdf-\codebreak{}xdg} & \texttt{export-\codebreak{}latex-\codebreak{}compile-\codebreak{}pdf-\codebreak{}xdg} & Make the \LaTeX{} file of the current buffer into a PDF beside it, as M-x export-latex-compile-pdf does, and open the PDF in the desktop’s program for PDFs (xdg-open), or the one \textbf{open-externally-programs} gives “application/pdf”, asking first unless \textbf{open-externally-ask} is NIL.\\
none & \texttt{M-x export-\codebreak{}latex-\codebreak{}template-\codebreak{}copy} & \texttt{export-\codebreak{}latex-\codebreak{}template-\codebreak{}copy} & Copy a built-in template of the enhanced export, asked for with completion (article, its own look, or org, the plain export’s document), into your own templates, \textasciitilde{}/.config/tuatara/latex-templates/NAME.tex, and visit it there to change: from then on \#+template: NAME takes your copy.\\
none & \texttt{M-x export-\codebreak{}template-\codebreak{}from-\codebreak{}docx} & \texttt{export-\codebreak{}template-\codebreak{}from-\codebreak{}docx} & Make a Word document given to you a template to write reports in.\\
none & \texttt{M-x export-\codebreak{}template-\codebreak{}from-\codebreak{}latex} & \texttt{export-\codebreak{}template-\codebreak{}from-\codebreak{}latex} & Make a \LaTeX{} document given to you a template to write reports in.\\
\texttt{g Q} & \texttt{M-q} & \texttt{fill-\codebreak{}paragraph} & Fill the paragraph at point: its words re-flowed so that no line passes the fill column (80, \textbf{FILL-COLUMN}, or the buffer’s own), the text before them on each line kept (an indentation, a comment’s ;;, a quotation’s >).\\
\texttt{g q} & \texttt{C-z g q} & \texttt{fill-\codebreak{}region} & Fill the lines the selection touches, each paragraph in them as fill-paragraph fills one, the lines before and after the selection left as they are: the column asked for first, or given by a count, and kept as the buffer’s.\\
none & \texttt{C-x C-f} & \texttt{find-\codebreak{}file} & Visit a file, completing its name against the folder being typed.\\
none & \texttt{C-x n f} & \texttt{find-\codebreak{}file-\codebreak{}formatted} & Visit a file, or make a new one with a formatted name.\\
none & \texttt{C-x 5 f} & \texttt{find-\codebreak{}file-\codebreak{}other-\codebreak{}frame} & Visit a file in another frame and work in it, asking for its name as C-x C-f does; a folder is listed there.\\
none & \texttt{C-x 4 f}, \texttt{C-x 4 C-f} & \texttt{find-\codebreak{}file-\codebreak{}other-\codebreak{}window} & Visit a file in another window and select it, completing its name as find-file does.\\
none & \texttt{M-x find-\codebreak{}file-\codebreak{}safely} & \texttt{find-\codebreak{}file-\codebreak{}safely} & Open a file safely, for looking at a file from a source you do not trust: read-only, with no major mode, no hooks and no version control, so that nothing in the file is acted on; its file system left exactly as it was (no copy beside it), and its access time too where the system allows.\\
none & \texttt{C-x C-r} & \texttt{find-\codebreak{}recent-\codebreak{}file} & Visit a file visited lately: the latest first, each with its permissions, size and when it last changed beside it.\\
none & \texttt{M-x flip-\codebreak{}window-\codebreak{}layout-\codebreak{}horizontally} & \texttt{flip-\codebreak{}window-\codebreak{}layout-\codebreak{}horizontally} & Mirror the layout of windows left to right.\\
none & \texttt{M-x flip-\codebreak{}window-\codebreak{}layout-\codebreak{}vertically} & \texttt{flip-\codebreak{}window-\codebreak{}layout-\codebreak{}vertically} & Mirror the layout of windows top to bottom.\\
none & \texttt{M-x flush-\codebreak{}lines} & \texttt{flush-\codebreak{}lines} & Delete the lines the selection touches in which a regular expression, asked for in the minibuffer, finds a match.\\
none & \texttt{M-x follow-\codebreak{}file} & \texttt{follow-\codebreak{}file} & Keep the current buffer up with its file on disk, or stop.\\
none & \texttt{M-x font-\codebreak{}list} & \texttt{font-\codebreak{}list} & List every font installed in \textbf{Fonts}, the fixed-width ones first: each family with its variants, the line of the init file that chooses each, and a sample written in it.\\
none & \texttt{w}, \texttt{C-c w} & \texttt{font-\codebreak{}list-\codebreak{}copy} & Copy the init-file line that chooses the font on the current line, to put in the init file (C-x C-f \textasciitilde{}/.config/tuatara/init.lisp).\\
none & \texttt{RET} & \texttt{font-\codebreak{}list-\codebreak{}use} & Show the window in the font on the current line of the list of fonts, for this session; w copies the line that chooses it for good.\\
none & \texttt{C-c f} & \texttt{format-\codebreak{}buffer} & Lay out the code of the current buffer as its language’s formatter does (\textbf{LANGUAGE-FORMATTERS}, else the language’s usual one: rustfmt, ruff, clang-format, shfmt\ldots{}), in the background, shut in the sandbox; the change is one that C-/ takes back whole.\\
none & \texttt{C-c F} & \texttt{format-\codebreak{}region} & Lay out the code selected as the language server does; with no server that lays out a part, each line selected indented as the mode indents it.\\
none & \texttt{TAB} & \texttt{forward-\codebreak{}button} & Move point to the next button, from the last round to the first.\\
\bottomrule
\end{xltabular}
}

\begin{xltabular}{\textwidth}{@{}l >{\raggedright\arraybackslash\hsize=0.619\hsize}X l >{\raggedright\arraybackslash\hsize=1.381\hsize}X@{}}
\toprule
\textbf{Modal} & \textbf{Non-modal} & \textbf{Command} & \textbf{What it does}\\
\midrule
\endhead
\texttt{l} & \texttt{C-f}, \texttt{<right>} & \texttt{forward-\codebreak{}char} & Move point forward one character (or the prefix argument’s worth); in the hex view, one octet.\\
none & \texttt{C-M-n} & \texttt{forward-\codebreak{}list} & Move forward over the next bracketed group: the words, strings and comments before it passed over, the group taken whole.\\
none & \texttt{M-\}} & \texttt{forward-\codebreak{}paragraph} & Move point to the end of the paragraph: the start of the blank line after it, or the end of the text.\\
none & \texttt{M-e} & \texttt{forward-\codebreak{}sentence} & Move point to the end of the sentence: after its .\\
\texttt{W} & \texttt{C-M-f} & \texttt{forward-\codebreak{}sexp} & Move forward over a balanced expression.\\
\texttt{w} & \texttt{M-f} & \texttt{forward-\codebreak{}word} & Move point to the end of the next word (or the prefix argument’s worth, back with a negative one): past what is not a word, then past the word.\\
none & \texttt{C-x g} & \texttt{git} & Show the state of the git repository the current file (or folder) is in: the branch, the files not tracked, changed and staged, the latest commits; C-h b there lists its keys.\\
none & \texttt{M-x git-am} & \texttt{git-am} & Take in patch files asked for, as sent by mail (git format-patch writes them), each made a commit with its author and message; one that does not fit stops it half way (w c goes on, w s skips it, w a stops it).\\
none & \texttt{M-x git-\codebreak{}am-\codebreak{}abort} & \texttt{git-\codebreak{}am-\codebreak{}abort} & Stop the taking in of patches, after asking: the branch and the files go back to where they were before it began.\\
none & \texttt{M-x git-\codebreak{}am-\codebreak{}continue} & \texttt{git-\codebreak{}am-\codebreak{}continue} & Go on with the taking in of patches stopped half way, once the files git could not join are put right and staged.\\
none & \texttt{w}, \texttt{C-c w} & \texttt{git-\codebreak{}am-\codebreak{}menu} & Open the menu of patches by mail: take patch files in as commits, or go on with, skip one of, or stop a taking in stopped half way.\\
none & \texttt{M-x git-\codebreak{}am-\codebreak{}skip} & \texttt{git-\codebreak{}am-\codebreak{}skip} & Skip the commit the taking in of patches stopped at and go on with the next.\\
none & \texttt{M-x git-\codebreak{}apply-\codebreak{}patch} & \texttt{git-\codebreak{}apply-\codebreak{}patch} & Apply a patch file asked for to the files, as git apply does; the patch menu’s switches stage it too (--index), only stage it (--cached), or join it three ways where it does not fit (those for writing patches are left out).\\
none & \texttt{M-x git-\codebreak{}bisect-\codebreak{}bad} & \texttt{git-\codebreak{}bisect-\codebreak{}bad} & Say the commit checked out has the fault: git tries one before it.\\
none & \texttt{M-x git-\codebreak{}bisect-\codebreak{}good} & \texttt{git-\codebreak{}bisect-\codebreak{}good} & Say the commit checked out has not the fault: git tries one after it.\\
none & \texttt{B}, \texttt{C-c B} & \texttt{git-\codebreak{}bisect-\codebreak{}menu} & Open the bisect menu: look for the commit that brought a fault in, saying of each commit tried whether it is good or bad, or letting a command say.\\
none & \texttt{M-x git-\codebreak{}bisect-\codebreak{}reset} & \texttt{git-\codebreak{}bisect-\codebreak{}reset} & Stop looking: HEAD goes back to where it was when the bisection began.\\
none & \texttt{M-x git-\codebreak{}bisect-\codebreak{}run} & \texttt{git-\codebreak{}bisect-\codebreak{}run} & Let a command asked for, run by the shell in the repository, say of each commit tried whether it has the fault (it fails) or not (it succeeds, 125 for one not to be tried), until the first bad commit is found.\\
none & \texttt{M-x git-\codebreak{}bisect-\codebreak{}skip} & \texttt{git-\codebreak{}bisect-\codebreak{}skip} & Say the commit checked out cannot be tried: git tries one beside it.\\
none & \texttt{M-x git-\codebreak{}bisect-\codebreak{}start} & \texttt{git-\codebreak{}bisect-\codebreak{}start} & Begin looking for the commit that brought a fault in: a revision that has it (HEAD when nothing is typed) and one that has not are asked for, and git checks out one between to try; B g, B b and B s then say it is good, bad, or not to be tried.\\
none & \texttt{M-x git-\codebreak{}blame} & \texttt{git-\codebreak{}blame} & List the current file (as saved), or the one on point’s line of a status, a line each with the commit that last changed it, its author and date; RET on a line shows its commit.\\
none & \texttt{l}, \texttt{RET}, \texttt{C-c l} & \texttt{git-\codebreak{}blame-\codebreak{}show} & Show the commit that last changed the line of the blame point is on.\\
none & \texttt{M-x git-\codebreak{}branch-\codebreak{}and-\codebreak{}checkout} & \texttt{git-\codebreak{}branch-\codebreak{}and-\codebreak{}checkout} & Make a new branch, asked for, starting at a revision asked for (the branch checked out when nothing is typed), and check it out.\\
none & \texttt{M-x git-\codebreak{}branch-\codebreak{}create} & \texttt{git-\codebreak{}branch-\codebreak{}create} & Make a new branch, asked for, starting at a revision asked for, without checking it out.\\
none & \texttt{M-x git-\codebreak{}branch-\codebreak{}delete} & \texttt{git-\codebreak{}branch-\codebreak{}delete} & Delete a branch, asked for, not the one checked out: one with commits no other branch has is deleted only after asking again, as they would be found only in git’s reflog then.\\
none & \texttt{b}, \texttt{C-c b} & \texttt{git-\codebreak{}branch-\codebreak{}menu} & Open the branch menu: check out a branch or a revision, make one (checked out or not, or spun off with the work not pushed), rename, move, delete one, or set its upstream.\\
none & \texttt{M-x git-\codebreak{}branch-\codebreak{}rename} & \texttt{git-\codebreak{}branch-\codebreak{}rename} & Rename a branch, asked for (the one checked out when nothing is typed), to a name asked for.\\
none & \texttt{M-x git-\codebreak{}branch-\codebreak{}reset} & \texttt{git-\codebreak{}branch-\codebreak{}reset} & Move a branch, asked for, to a revision asked for, after asking: the branch checked out keeps the changes not committed (git refuses when they would be lost); its commits after the revision stay in git’s reflog.\\
none & \texttt{M-x git-\codebreak{}branch-\codebreak{}set-\codebreak{}upstream} & \texttt{git-\codebreak{}branch-\codebreak{}set-\codebreak{}upstream} & Set the upstream of a branch, asked for (the one checked out when nothing is typed), to a branch of a remote asked for: the one pull brings from and push sends to, and the status counts ahead and behind of.\\
none & \texttt{M-x git-\codebreak{}branch-\codebreak{}spin-\codebreak{}off} & \texttt{git-\codebreak{}branch-\codebreak{}spin-\codebreak{}off} & Make a new branch, asked for, of the work on the branch checked out and check it out: the commits the branch has that its upstream has not go with the new branch, and the branch goes back to where it and its upstream parted; changes not committed come along.\\
none & \texttt{M-x git-\codebreak{}bundle-\codebreak{}create} & \texttt{git-\codebreak{}bundle-\codebreak{}create} & Write the whole repository, every branch and tag, as a single bundle file asked for: one file to carry it, from which it can be cloned.\\
none & \texttt{M-x git-\codebreak{}bundle-\codebreak{}verify} & \texttt{git-\codebreak{}bundle-\codebreak{}verify} & Check a bundle file asked for: whether this repository has what it needs to take it in.\\
none & \texttt{M-x git-\codebreak{}checkout} & \texttt{git-\codebreak{}checkout} & Check out a branch or a revision, asked for: a branch by its name, else a revision (a commit, a tag, origin/main), which leaves HEAD detached.\\
none & \texttt{M-x git-\codebreak{}cherry-\codebreak{}pick} & \texttt{git-\codebreak{}cherry-\codebreak{}pick} & Make a commit, asked for (the one on point’s line of a status or a log first, or another branch’s newest by the branch’s name), again on the branch checked out, with its message.\\
none & \texttt{M-x git-\codebreak{}cherry-\codebreak{}pick-\codebreak{}abort} & \texttt{git-\codebreak{}cherry-\codebreak{}pick-\codebreak{}abort} & Stop the cherry-pick, after asking: the branch and the files go back to where they were before it began.\\
none & \texttt{M-x git-\codebreak{}cherry-\codebreak{}pick-\codebreak{}continue} & \texttt{git-\codebreak{}cherry-\codebreak{}pick-\codebreak{}continue} & Go on with the cherry-pick stopped half way, once the files git could not join are put right and staged.\\
none & \texttt{A}, \texttt{C-c A} & \texttt{git-\codebreak{}cherry-\codebreak{}pick-\codebreak{}menu} & Open the cherry-pick menu: make a commit again on this branch, or bring its changes in uncommitted; go on with, skip or stop one stopped half way.\\
none & \texttt{M-x git-\codebreak{}cherry-\codebreak{}pick-\codebreak{}no-\codebreak{}commit} & \texttt{git-\codebreak{}cherry-\codebreak{}pick-\codebreak{}no-\codebreak{}commit} & Bring the changes of a commit, asked for, into the files and stage them, not committed.\\
none & \texttt{M-x git-\codebreak{}cherry-\codebreak{}pick-\codebreak{}skip} & \texttt{git-\codebreak{}cherry-\codebreak{}pick-\codebreak{}skip} & Skip the commit the cherry-pick stopped at and go on with the next.\\
none & \texttt{M-x git-\codebreak{}commit} & \texttt{git-\codebreak{}commit} & Write a commit of what is staged: its message in a buffer of its own, what is staged listed under it as comments; C-c C-c commits it, C-c C-k gives it up.\\
\bottomrule
\end{xltabular}

\begin{xltabular}{\textwidth}{@{}l >{\raggedright\arraybackslash\hsize=0.542\hsize}X l >{\raggedright\arraybackslash\hsize=1.458\hsize}X@{}}
\toprule
\textbf{Modal} & \textbf{Non-modal} & \textbf{Command} & \textbf{What it does}\\
\midrule
\endhead
none & \texttt{M-x git-\codebreak{}commit-\codebreak{}amend} & \texttt{git-\codebreak{}commit-\codebreak{}amend} & Amend the last commit: what is staged goes into it, and its message, written in a buffer to begin with, is written again; C-c C-c amends it.\\
none & \texttt{C-c C-k} & \texttt{git-\codebreak{}commit-\codebreak{}cancel} & Give up this commit’s message: nothing is committed.\\
none & \texttt{M-x git-\codebreak{}commit-\codebreak{}extend} & \texttt{git-\codebreak{}commit-\codebreak{}extend} & Extend the last commit with what is staged, its message as it is.\\
none & \texttt{C-c C-c} & \texttt{git-\codebreak{}commit-\codebreak{}finish} & Commit with the message of this buffer (lines beginning with \# left out), as git commit does with your settings (a signature among them) and the arguments the commit was begun with (--amend for an amended one), and close it.\\
none & \texttt{M-x git-\codebreak{}commit-\codebreak{}fixup} & \texttt{git-\codebreak{}commit-\codebreak{}fixup} & Make a commit of what is staged fixing up one asked for: a rebase with autosquash folds it into that one, its message left as it was.\\
none & \texttt{c}, \texttt{C-c c} & \texttt{git-\codebreak{}commit-\codebreak{}menu} & Open the commit menu: write a commit, amend, reword or extend the last one, or make one fixing up or squashed into another; its switches go to git commit.\\
none & \texttt{M-x git-\codebreak{}commit-\codebreak{}reword} & \texttt{git-\codebreak{}commit-\codebreak{}reword} & Reword the last commit: its message, written in a buffer to begin with, is written again, nothing staged going into it; C-c C-c rewords it.\\
none & \texttt{M-x git-\codebreak{}commit-\codebreak{}squash} & \texttt{git-\codebreak{}commit-\codebreak{}squash} & Make a commit of what is staged to be squashed into one asked for: a rebase with autosquash folds it in, its message added to that one’s.\\
none & \texttt{M-x git-\codebreak{}delete-\codebreak{}file} & \texttt{git-\codebreak{}delete-\codebreak{}file} & Delete the file you are in, or the one on point’s line of a status, from the files and from git, after asking; the deletion is staged.\\
none & \texttt{M-x git-\codebreak{}diff} & \texttt{git-\codebreak{}diff} & Show every change of the repository not yet committed, staged or not, in a buffer of its own; the diff menu’s switches shape it.\\
none & \texttt{M-x git-\codebreak{}diff-\codebreak{}commit} & \texttt{git-\codebreak{}diff-\codebreak{}commit} & Show a commit asked for (the one on the line first): what it says and its changes.\\
none & \texttt{M-x git-\codebreak{}diff-\codebreak{}file} & \texttt{git-\codebreak{}diff-\codebreak{}file} & Show the changes of the file you are in not yet committed, staged or not.\\
none & \texttt{d}, \texttt{C-c d} & \texttt{git-\codebreak{}diff-\codebreak{}menu} & Open the diff menu: show the changes not committed, what is staged, what is not, the files against a revision, between two revisions, or a commit; its switches shape them.\\
none & \texttt{M-x git-\codebreak{}diff-\codebreak{}range} & \texttt{git-\codebreak{}diff-\codebreak{}range} & Show what changed between two revisions asked for: from the first to the second.\\
none & \texttt{M-x git-\codebreak{}diff-\codebreak{}revision} & \texttt{git-\codebreak{}diff-\codebreak{}revision} & Show how the files differ from a revision asked for (a branch, a tag, a commit).\\
none & \texttt{M-x git-\codebreak{}diff-\codebreak{}staged} & \texttt{git-\codebreak{}diff-\codebreak{}staged} & Show what is staged: the changes the next commit makes.\\
none & \texttt{M-x git-\codebreak{}diff-\codebreak{}unstaged} & \texttt{git-\codebreak{}diff-\codebreak{}unstaged} & Show the changes not staged: those in the files beyond what is staged.\\
none & \texttt{M-x git-\codebreak{}fetch} & \texttt{git-\codebreak{}fetch} & Fetch the commits of the branch’s upstream remote (or origin) that this repository has not, the branches checked out left as they are.\\
none & \texttt{M-x git-\codebreak{}fetch-\codebreak{}all} & \texttt{git-\codebreak{}fetch-\codebreak{}all} & Fetch from every remote.\\
none & \texttt{f}, \texttt{C-c f} & \texttt{git-\codebreak{}fetch-\codebreak{}menu} & Open the fetch menu: fetch from the upstream remote, a remote asked for, or every remote.\\
none & \texttt{M-x git-\codebreak{}fetch-\codebreak{}remote} & \texttt{git-\codebreak{}fetch-\codebreak{}remote} & Fetch from a remote asked for.\\
none & \texttt{M-x git-\codebreak{}file-\codebreak{}menu} & \texttt{git-\codebreak{}file-\codebreak{}menu} & Open the menu of the file you are in (or the one on point’s line of a status): stage or unstage it or its hunk at point, show its changes, its log or who changed each line, go back through its history, put it back to a revision, rename or delete it.\\
none & \texttt{M-x git-\codebreak{}format-\codebreak{}patch} & \texttt{git-\codebreak{}format-\codebreak{}patch} & Write the commits after a revision asked for as patch files, a commit a file, in a folder asked for; the patch menu’s switches number them, add a cover letter, or sign them off (those for applying are left out).\\
none & \texttt{M-x git-\codebreak{}ignore-\codebreak{}folder} & \texttt{git-\codebreak{}ignore-\codebreak{}folder} & Have git ignore a pattern asked for in the .gitignore of the current file’s folder.\\
none & \texttt{i}, \texttt{C-c i} & \texttt{git-\codebreak{}ignore-\codebreak{}menu} & Open the ignore menu: have git ignore a file or a pattern, for everyone (.gitignore at the top or in the file’s folder) or here alone.\\
none & \texttt{M-x git-\codebreak{}ignore-\codebreak{}private} & \texttt{git-\codebreak{}ignore-\codebreak{}private} & Have git ignore a pattern asked for here alone, in git’s own exclude file, which is never committed.\\
none & \texttt{M-x git-\codebreak{}ignore-\codebreak{}shared} & \texttt{git-\codebreak{}ignore-\codebreak{}shared} & Have git ignore a pattern asked for (the file on point’s line, or the current one, offered), in the repository’s .gitignore, committed with it.\\
none & \texttt{M-x git-\codebreak{}log} & \texttt{git-\codebreak{}log} & List the commits of the branch checked out, the latest first, a line each (\textbf{GIT-LOG-LENGTH} of them); RET shows one.\\
none & \texttt{M-x git-\codebreak{}log-\codebreak{}all} & \texttt{git-\codebreak{}log-\codebreak{}all} & List the commits of every reference: branches, remotes’ branches, tags, stashes.\\
none & \texttt{M-x git-\codebreak{}log-\codebreak{}branches} & \texttt{git-\codebreak{}log-\codebreak{}branches} & List the commits of every branch, the repository’s own and its remotes’.\\
none & \texttt{M-x git-\codebreak{}log-\codebreak{}file} & \texttt{git-\codebreak{}log-\codebreak{}file} & List the commits that changed the file on point’s line of a status, or the current file, following it across renames.\\
none & \texttt{L}, \texttt{C-c L} & \texttt{git-\codebreak{}log-\codebreak{}menu} & Open the log menu: list the commits of the branch checked out, another, every branch or reference, or a file’s, or the reflog; its switches and options shape the list.\\
none & \texttt{M-x git-\codebreak{}log-\codebreak{}other} & \texttt{git-\codebreak{}log-\codebreak{}other} & List the commits of a branch or revision asked for.\\
none & \texttt{l}, \texttt{RET}, \texttt{C-c l} & \texttt{git-\codebreak{}log-\codebreak{}show} & Show the commit on the current line of the log.\\
none & \texttt{M-x git-\codebreak{}merge} & \texttt{git-\codebreak{}merge} & Merge a branch, asked for, into the one checked out, committing it with git’s own message; files changed on both sides stop it half way, said so in the status (c c commits once they are put right, m a stops it).\\
none & \texttt{M-x git-\codebreak{}merge-\codebreak{}abort} & \texttt{git-\codebreak{}merge-\codebreak{}abort} & Stop the merge, after asking: the branch and the files go back to where they were before it began.\\
none & \texttt{M-x git-\codebreak{}merge-\codebreak{}continue} & \texttt{git-\codebreak{}merge-\codebreak{}continue} & Go on with the merge stopped half way, once the files git could not join are put right and staged.\\
none & \texttt{m}, \texttt{C-c m} & \texttt{git-\codebreak{}merge-\codebreak{}menu} & Open the merge menu: merge a branch, committed or not, squash its changes into one, or stop a merge stopped half way.\\
none & \texttt{M-x git-\codebreak{}merge-\codebreak{}no-\codebreak{}commit} & \texttt{git-\codebreak{}merge-\codebreak{}no-\codebreak{}commit} & Merge a branch, asked for, without committing: the merge is left staged, to look at and commit (c c) or stop (m a).\\
none & \texttt{M-x git-\codebreak{}merge-\codebreak{}squash} & \texttt{git-\codebreak{}merge-\codebreak{}squash} & Bring the changes of a branch, asked for, into the files and stage them as one change, not a merge: c c commits them.\\
\bottomrule
\end{xltabular}

\begin{xltabular}{\textwidth}{@{}l >{\raggedright\arraybackslash\hsize=0.756\hsize}X l >{\raggedright\arraybackslash\hsize=1.244\hsize}X@{}}
\toprule
\textbf{Modal} & \textbf{Non-modal} & \textbf{Command} & \textbf{What it does}\\
\midrule
\endhead
none & \texttt{n}, \texttt{C-c n} & \texttt{git-\codebreak{}next-\codebreak{}section} & Go to the next file or hunk of the status, or the next title.\\
none & \texttt{M-x git-\codebreak{}notes-\codebreak{}edit} & \texttt{git-\codebreak{}notes-\codebreak{}edit} & Write a note on a commit asked for (the one on point’s line first, else HEAD), its note so far shown to begin with: a note says more of a commit without changing it.\\
none & \texttt{T}, \texttt{C-c T} & \texttt{git-\codebreak{}notes-\codebreak{}menu} & Open the notes menu: write a note on a commit, take one away, or prune those on commits gone.\\
none & \texttt{M-x git-\codebreak{}notes-\codebreak{}prune} & \texttt{git-\codebreak{}notes-\codebreak{}prune} & Take away the notes on commits the repository no longer has.\\
none & \texttt{M-x git-\codebreak{}notes-\codebreak{}remove} & \texttt{git-\codebreak{}notes-\codebreak{}remove} & Take away the note on a commit asked for.\\
none & \texttt{W}, \texttt{C-c W} & \texttt{git-\codebreak{}patch-\codebreak{}menu} & Open the patch menu: write commits as patch files, apply a patch file, save the changes not committed as one, or write or check a bundle.\\
none & \texttt{p}, \texttt{C-c p} & \texttt{git-\codebreak{}previous-\codebreak{}section} & Go to the previous file or hunk of the status, or the previous title.\\
none & \texttt{M-x git-\codebreak{}pull} & \texttt{git-\codebreak{}pull} & Bring in the upstream’s commits: only when the branch can simply move on to them (fast-forward), unless the pull menu says to rebase (-r) or merge (-m) when it cannot.\\
none & \texttt{M-x git-\codebreak{}pull-\codebreak{}from} & \texttt{git-\codebreak{}pull-\codebreak{}from} & Bring in the commits of a branch, asked for, of a remote asked for, as M-x git-pull does.\\
none & \texttt{F}, \texttt{C-c F} & \texttt{git-\codebreak{}pull-\codebreak{}menu} & Open the pull menu: bring in the upstream’s commits, or a branch’s of a remote asked for.\\
none & \texttt{M-x git-\codebreak{}push} & \texttt{git-\codebreak{}push} & Push the branch to its upstream, in the background; the push menu’s switches go with it (a forced push asked first).\\
none & \texttt{M-x git-\codebreak{}push-\codebreak{}elsewhere} & \texttt{git-\codebreak{}push-\codebreak{}elsewhere} & Push the branch checked out to a remote asked for, under its own name; -u in the push menu makes that its upstream.\\
none & \texttt{P}, \texttt{C-c P} & \texttt{git-\codebreak{}push-\codebreak{}menu} & Open the push menu: push the branch to its upstream or a remote asked for, every tag or one; its switches force, set the upstream, try only, skip the hooks.\\
none & \texttt{M-x git-\codebreak{}push-\codebreak{}tag} & \texttt{git-\codebreak{}push-\codebreak{}tag} & Push a tag, asked for, to a remote asked for.\\
none & \texttt{M-x git-\codebreak{}push-\codebreak{}tags} & \texttt{git-\codebreak{}push-\codebreak{}tags} & Push every tag to the upstream remote (or origin).\\
none & \texttt{M-x git-\codebreak{}rebase-\codebreak{}abort} & \texttt{git-\codebreak{}rebase-\codebreak{}abort} & Stop the rebase, after asking: the branch and the files go back to where they were before it began.\\
none & \texttt{M-x git-\codebreak{}rebase-\codebreak{}continue} & \texttt{git-\codebreak{}rebase-\codebreak{}continue} & Go on with the rebase stopped half way, once the files git could not join are put right and staged.\\
none & \texttt{M-x git-\codebreak{}rebase-\codebreak{}interactive} & \texttt{git-\codebreak{}rebase-\codebreak{}interactive} & Rebase the commits after a revision asked for (the commit on point’s line first) as a list says.\\
none & \texttt{R}, \texttt{C-c R} & \texttt{git-\codebreak{}rebase-\codebreak{}menu} & Open the rebase menu: rebase the branch on its upstream or a revision, or go on with, skip a commit of, or stop a rebase stopped half way.\\
none & \texttt{M-x git-\codebreak{}rebase-\codebreak{}onto} & \texttt{git-\codebreak{}rebase-\codebreak{}onto} & Rebase the branch checked out on a revision asked for: its commits since they parted made again after it.\\
none & \texttt{M-x git-\codebreak{}rebase-\codebreak{}skip} & \texttt{git-\codebreak{}rebase-\codebreak{}skip} & Skip the commit the rebase stopped at and go on with the next.\\
none & \texttt{C-c C-k} & \texttt{git-\codebreak{}rebase-\codebreak{}todo-\codebreak{}cancel} & Give up this rebase’s list: nothing is rebased.\\
none & \texttt{C-c C-c} & \texttt{git-\codebreak{}rebase-\codebreak{}todo-\codebreak{}finish} & Rebase as this list says: each reword line asks the commit’s new subject first (the rest of its message kept), then git rebases in the background; the list’s folder, the new messages in it, is deleted once git is done.\\
none & \texttt{M-x git-\codebreak{}rebase-\codebreak{}upstream} & \texttt{git-\codebreak{}rebase-\codebreak{}upstream} & Rebase the branch checked out on its upstream: its own commits made again after the upstream’s.\\
none & \texttt{M-x git-\codebreak{}reflog} & \texttt{git-\codebreak{}reflog} & List where HEAD has been, the latest first: each commit it was at and what moved it there (a commit, a checkout, a reset), as git’s reflog keeps it.\\
none & \texttt{l}, \texttt{RET}, \texttt{C-c l} & \texttt{git-\codebreak{}refs-\codebreak{}log} & List the log of the branch or tag on the line of the list of references; on a section’s heading, fold or unfold the section.\\
none & \texttt{M-x git-\codebreak{}remote-\codebreak{}add} & \texttt{git-\codebreak{}remote-\codebreak{}add} & Add a remote: a name asked for (origin when nothing is typed and there is none) and its address; its branches are fetched with f e.\\
none & \texttt{M}, \texttt{C-c M} & \texttt{git-\codebreak{}remote-\codebreak{}menu} & Open the remote menu: add a remote, rename one, give it a new address, remove it, or forget its branches gone.\\
none & \texttt{M-x git-\codebreak{}remote-\codebreak{}prune} & \texttt{git-\codebreak{}remote-\codebreak{}prune} & Forget the branches of a remote, asked for, that it no longer has.\\
none & \texttt{M-x git-\codebreak{}remote-\codebreak{}remove} & \texttt{git-\codebreak{}remote-\codebreak{}remove} & Remove a remote, asked for, after asking: its branches here go with it.\\
none & \texttt{M-x git-\codebreak{}remote-\codebreak{}rename} & \texttt{git-\codebreak{}remote-\codebreak{}rename} & Rename a remote, asked for, to a name asked for.\\
none & \texttt{M-x git-\codebreak{}remote-\codebreak{}set-\codebreak{}address} & \texttt{git-\codebreak{}remote-\codebreak{}set-\codebreak{}address} & Give a remote, asked for, a new address, its own shown to begin with.\\
none & \texttt{M-x git-\codebreak{}rename-\codebreak{}file} & \texttt{git-\codebreak{}rename-\codebreak{}file} & Rename the file you are in, or the one on point’s line of a status, to a path asked for within the repository (into a folder there, when it names one); the rename is staged, and the buffers visiting the file follow it.\\
none & \texttt{M-x git-\codebreak{}reset-\codebreak{}file} & \texttt{git-\codebreak{}reset-\codebreak{}file} & Put a file, the one on point’s line of a status or the current one, back to how a revision asked for has it, staged, after asking: its changes not committed are lost.\\
none & \texttt{M-x git-\codebreak{}reset-\codebreak{}hard} & \texttt{git-\codebreak{}reset-\codebreak{}hard} & Move the branch checked out to a revision asked for, what is staged and the files with it, after asking: changes not committed are lost.\\
none & \texttt{M-x git-\codebreak{}reset-\codebreak{}keep} & \texttt{git-\codebreak{}reset-\codebreak{}keep} & Move the branch checked out to a revision asked for, the files with it, keeping changes not committed; git refuses when they would be lost.\\
none & \texttt{X}, \texttt{C-c X} & \texttt{git-\codebreak{}reset-\codebreak{}menu} & Open the reset menu: move the branch to a revision, keeping what is staged, the files, or changes not committed, or nothing; or put a file back to a revision.\\
none & \texttt{M-x git-\codebreak{}reset-\codebreak{}mixed} & \texttt{git-\codebreak{}reset-\codebreak{}mixed} & Move the branch checked out to a revision asked for, what is staged put back to it too, the files as they are (their changes shown as not staged).\\
none & \texttt{M-x git-\codebreak{}reset-\codebreak{}soft} & \texttt{git-\codebreak{}reset-\codebreak{}soft} & Move the branch checked out to a revision asked for, what is staged and the files as they are (the commits’ changes shown as staged).\\
none & \texttt{x} & \texttt{git-\codebreak{}reverse} & Reverse the changes of the file on the line of the status (on a section’s title, of all its files), after asking, which says it cannot be undone: git kept no copy of them.\\
\bottomrule
\end{xltabular}

\begin{xltabular}{\textwidth}{@{}l >{\raggedright\arraybackslash\hsize=0.711\hsize}X l >{\raggedright\arraybackslash\hsize=1.289\hsize}X@{}}
\toprule
\textbf{Modal} & \textbf{Non-modal} & \textbf{Command} & \textbf{What it does}\\
\midrule
\endhead
none & \texttt{M-x git-\codebreak{}revert} & \texttt{git-\codebreak{}revert} & Make a commit undoing a commit asked for (the one on point’s line of a status or a log first), with git’s message saying so.\\
none & \texttt{M-x git-\codebreak{}revert-\codebreak{}abort} & \texttt{git-\codebreak{}revert-\codebreak{}abort} & Stop the revert, after asking: the branch and the files go back to where they were before it began.\\
none & \texttt{M-x git-\codebreak{}revert-\codebreak{}continue} & \texttt{git-\codebreak{}revert-\codebreak{}continue} & Go on with the revert stopped half way, once the files git could not join are put right and staged.\\
none & \texttt{\_}, \texttt{C-c \_} & \texttt{git-\codebreak{}revert-\codebreak{}menu} & Open the revert menu: make a commit undoing another, or undo it in the files uncommitted; go on with, skip or stop one stopped half way.\\
none & \texttt{M-x git-\codebreak{}revert-\codebreak{}no-\codebreak{}commit} & \texttt{git-\codebreak{}revert-\codebreak{}no-\codebreak{}commit} & Undo the changes of a commit, asked for, in the files and stage that, not committed.\\
none & \texttt{M-x git-\codebreak{}revert-\codebreak{}skip} & \texttt{git-\codebreak{}revert-\codebreak{}skip} & Skip the commit the revert stopped at and go on with the next.\\
none & \texttt{C-x G} & \texttt{git-\codebreak{}rewind} & Go back through git’s history a commit at a time, and come back.\\
none & \texttt{M-x git-\codebreak{}rewind-\codebreak{}copy-\codebreak{}hash} & \texttt{git-\codebreak{}rewind-\codebreak{}copy-\codebreak{}hash} & Copy the short hash of the commit that left the file or folder of this history as it is shown.\\
none & \texttt{M-x git-\codebreak{}rewind-\codebreak{}newer} & \texttt{git-\codebreak{}rewind-\codebreak{}newer} & Show the file or folder of this history as the commit after left it.\\
none & \texttt{M-x git-\codebreak{}rewind-\codebreak{}older} & \texttt{git-\codebreak{}rewind-\codebreak{}older} & Show the file or folder of this history as the commit before left it.\\
none & \texttt{M-x git-\codebreak{}rewind-\codebreak{}show-\codebreak{}commit} & \texttt{git-\codebreak{}rewind-\codebreak{}show-\codebreak{}commit} & Show the commit that left the file or folder of this history as it is shown: what it says and its changes.\\
none & \texttt{M-x git-\codebreak{}rewind-to} & \texttt{git-\codebreak{}rewind-to} & Show the file or folder of this history as a commit chosen from its log left it, the latest first, by its hash, date, author or subject.\\
none & \texttt{!}, \texttt{C-c !} & \texttt{git-\codebreak{}run-\codebreak{}command} & Run git with the words asked for (log --oneline -5, say; quotes group words as in a shell) in the repository at hand, no shell between, and show what it printed; for what the menus have no action for.\\
none & \texttt{M-x git-\codebreak{}save-\codebreak{}patch} & \texttt{git-\codebreak{}save-\codebreak{}patch} & Save every change not committed, staged or not, as a new patch file asked for, its octets as git writes them.\\
none & \texttt{y} & \texttt{git-\codebreak{}show-\codebreak{}refs} & List the repository’s branches (the one checked out marked *, each with its upstream), its remotes’ branches and its tags, a line each with the commit it is at and its subject, in sections folded as in every list in sections; RET on one lists its log.\\
none & \texttt{M-x git-\codebreak{}sparse-\codebreak{}add} & \texttt{git-\codebreak{}sparse-\codebreak{}add} & Check out folders asked for too, in a sparse checkout.\\
none & \texttt{M-x git-\codebreak{}sparse-\codebreak{}disable} & \texttt{git-\codebreak{}sparse-\codebreak{}disable} & Check the whole repository out again.\\
none & \texttt{>}, \texttt{C-c >} & \texttt{git-\codebreak{}sparse-\codebreak{}menu} & Open the sparse checkout menu: check out only some folders, more of them, everything again, or narrow it again.\\
none & \texttt{M-x git-\codebreak{}sparse-\codebreak{}reapply} & \texttt{git-\codebreak{}sparse-\codebreak{}reapply} & Keep out of the files again what the sparse checkout leaves out, after a merge or a rebase brought some in.\\
none & \texttt{M-x git-\codebreak{}sparse-\codebreak{}set} & \texttt{git-\codebreak{}sparse-\codebreak{}set} & Check out only folders asked for (and the files at the top), the rest of the repository kept out of the files.\\
none & \texttt{s} & \texttt{git-\codebreak{}stage} & Stage the file on the line of the status, or on a section’s title all of its files; on a hunk, the hunk; with lines selected in hunks, those lines.\\
none & \texttt{M-x git-\codebreak{}stage-\codebreak{}file} & \texttt{git-\codebreak{}stage-\codebreak{}file} & Stage the file you are in (as saved), or the one on point’s line of a status.\\
none & \texttt{M-x git-\codebreak{}stage-\codebreak{}hunk} & \texttt{git-\codebreak{}stage-\codebreak{}hunk} & Stage the hunk at point: in the status (C-x g) the one on point’s line, as s does; in a file’s own buffer, saved, the change git sees at point’s line, as the file is on disk.\\
none & \texttt{M-x git-\codebreak{}stash} & \texttt{git-\codebreak{}stash} & Put every change aside as a stash, staged or not, with a message asked for; the files go back to the last commit.\\
none & \texttt{M-x git-\codebreak{}stash-\codebreak{}apply} & \texttt{git-\codebreak{}stash-\codebreak{}apply} & Bring the changes of a stash, asked for, back into the files, the stash kept; git says when they meet changes of their own.\\
none & \texttt{M-x git-\codebreak{}stash-\codebreak{}branch} & \texttt{git-\codebreak{}stash-\codebreak{}branch} & Make a new branch, asked for, of a stash, asked for: checked out at the commit the stash was made on, its changes brought back, the stash dropped.\\
none & \texttt{M-x git-\codebreak{}stash-\codebreak{}drop} & \texttt{git-\codebreak{}stash-\codebreak{}drop} & Drop a stash, asked for, after asking: its changes are gone.\\
none & \texttt{M-x git-\codebreak{}stash-\codebreak{}index} & \texttt{git-\codebreak{}stash-\codebreak{}index} & Put only what is staged aside as a stash, with a message asked for; the changes not staged stay.\\
none & \texttt{M-x git-\codebreak{}stash-\codebreak{}keep-\codebreak{}index} & \texttt{git-\codebreak{}stash-\codebreak{}keep-\codebreak{}index} & Put every change aside as a stash, with a message asked for, keeping what is staged in place in the files and staged.\\
none & \texttt{z}, \texttt{C-c z} & \texttt{git-\codebreak{}stash-\codebreak{}menu} & Open the stash menu: put changes aside (all, only what is staged, or keeping it), apply, pop, drop or show a stash, or make a branch of one.\\
none & \texttt{M-x git-\codebreak{}stash-\codebreak{}pop} & \texttt{git-\codebreak{}stash-\codebreak{}pop} & Bring the changes of a stash, asked for, back into the files and drop it; one whose changes meet others is kept, said so.\\
none & \texttt{M-x git-\codebreak{}stash-\codebreak{}show} & \texttt{git-\codebreak{}stash-\codebreak{}show} & Show a stash, asked for: the files it changes and its changes.\\
none & \texttt{M-x git-\codebreak{}submodule-\codebreak{}add} & \texttt{git-\codebreak{}submodule-\codebreak{}add} & Add a repository, its address asked for, as a submodule at a path asked for (named after the address when nothing is typed), cloned there.\\
none & \texttt{o}, \texttt{C-c o} & \texttt{git-\codebreak{}submodule-\codebreak{}menu} & Open the submodule menu: add a submodule, check them out, take their addresses again, remove one or show its status.\\
none & \texttt{M-x git-\codebreak{}submodule-\codebreak{}remove} & \texttt{git-\codebreak{}submodule-\codebreak{}remove} & Remove a submodule, asked for, after asking: its folder and its entry in .gitmodules go, the removal staged.\\
none & \texttt{M-x git-\codebreak{}submodule-\codebreak{}status} & \texttt{git-\codebreak{}submodule-\codebreak{}status} & Show the status of a submodule, asked for (the one on point’s line first), as C-x g shows a repository’s.\\
none & \texttt{M-x git-\codebreak{}submodule-\codebreak{}sync} & \texttt{git-\codebreak{}submodule-\codebreak{}sync} & Take each submodule’s address again from .gitmodules, after it changed there.\\
none & \texttt{M-x git-\codebreak{}submodule-\codebreak{}update} & \texttt{git-\codebreak{}submodule-\codebreak{}update} & Check every submodule out at the commit this repository records (cloning those not checked out yet), and theirs in turn; with -r in the submodule menu at the newest commit of the branch each follows instead.\\
none & \texttt{M-x git-\codebreak{}subtree-\codebreak{}add} & \texttt{git-\codebreak{}subtree-\codebreak{}add} & Bring another repository’s branch into a new folder of this one as a subtree, its history joined in (one commit with -s in the subtree menu).\\
none & \texttt{O}, \texttt{C-c O} & \texttt{git-\codebreak{}subtree-\codebreak{}menu} & Open the subtree menu: bring another repository into a folder, pull its newer commits, push the folder’s, or split the folder’s history off.\\
\bottomrule
\end{xltabular}

{\small
\begin{xltabular}{\textwidth}{@{}l >{\raggedright\arraybackslash\hsize=0.766\hsize}X l >{\raggedright\arraybackslash\hsize=1.234\hsize}X@{}}
\toprule
\textbf{Modal} & \textbf{Non-modal} & \textbf{Command} & \textbf{What it does}\\
\midrule
\endhead
none & \texttt{M-x git-\codebreak{}subtree-\codebreak{}pull} & \texttt{git-\codebreak{}subtree-\codebreak{}pull} & Bring the newer commits of another repository’s branch into the subtree folder.\\
none & \texttt{M-x git-\codebreak{}subtree-\codebreak{}push} & \texttt{git-\codebreak{}subtree-\codebreak{}push} & Send the commits that changed the subtree folder to another repository’s branch.\\
none & \texttt{M-x git-\codebreak{}subtree-\codebreak{}split} & \texttt{git-\codebreak{}subtree-\codebreak{}split} & Make a history of the commits that changed a folder asked for, as if it were a repository of its own, and say the commit it ends at (a branch of it with =b in the subtree menu).\\
none & \texttt{M-x git-\codebreak{}tag-\codebreak{}create} & \texttt{git-\codebreak{}tag-\codebreak{}create} & Tag a revision asked for (the commit on point’s line first, else HEAD) with a name asked for; the tag menu’s -a makes it annotated with a message asked for, -s signed (with a message too), -f moves a tag of that name.\\
none & \texttt{M-x git-\codebreak{}tag-\codebreak{}delete} & \texttt{git-\codebreak{}tag-\codebreak{}delete} & Delete a tag, asked for, here (not on any remote), after asking.\\
none & \texttt{t} & \texttt{git-\codebreak{}tag-\codebreak{}menu} & Open the tag menu: tag a revision, or delete a tag; P t and P T in the push menu send tags.\\
none & \texttt{TAB} & \texttt{git-\codebreak{}toggle-\codebreak{}diff} & Show the changes of the file on the line of the status under it, in hunks, the first \textbf{GIT-DIFF-LINES} of them, or no longer.\\
none & \textasciitilde{}, C-c \textasciitilde{} & \texttt{git-\codebreak{}toggle-\codebreak{}full-\codebreak{}menus} & Make the git menus of this status show every action, or only the ones most used again (\textbf{GIT-MENUS}, \textbf{GIT-SLIM-ACTIONS}).\\
none & \texttt{u}, \texttt{C-c u} & \texttt{git-\codebreak{}unstage} & Unstage the file on the line of the status, or on a section’s title all of its files, a rename by both its names; on a hunk, the hunk; with lines selected in hunks, those lines.\\
none & \texttt{M-x git-\codebreak{}unstage-\codebreak{}file} & \texttt{git-\codebreak{}unstage-\codebreak{}file} & Unstage the file you are in, or the one on point’s line of a status, in a repository with no commit yet too.\\
none & \texttt{M-x git-\codebreak{}unstage-\codebreak{}hunk} & \texttt{git-\codebreak{}unstage-\codebreak{}hunk} & Unstage the hunk at point: in the status (C-x g) the one on point’s line, as u does; in a file’s own buffer, saved, the change staged at point’s line.\\
none & \texttt{l}, \texttt{RET}, \texttt{C-c l} & \texttt{git-\codebreak{}visit} & Visit the file on the line of the status, at the line of the hunk’s line point is on, a file in conflict at its first conflict, its conflicts shown (M-x conflict-markers-mode); or show the commit or the stash on it, or a submodule’s status; a file deleted is said to be.\\
none & \texttt{M-x git-\codebreak{}worktree-\codebreak{}branch} & \texttt{git-\codebreak{}worktree-\codebreak{}branch} & Make a new branch, asked for, starting at a revision asked for, and a worktree for it in a folder asked for.\\
none & \texttt{M-x git-\codebreak{}worktree-\codebreak{}checkout} & \texttt{git-\codebreak{}worktree-\codebreak{}checkout} & Make a worktree for a branch asked for, in a folder asked for: the branch checked out there, to work on beside this one.\\
none & \texttt{M-x git-\codebreak{}worktree-\codebreak{}delete} & \texttt{git-\codebreak{}worktree-\codebreak{}delete} & Delete a worktree, asked for, not the main one, after asking: its folder goes, its branch stays.\\
none & \texttt{Z}, \texttt{C-c Z} & \texttt{git-\codebreak{}worktree-\codebreak{}menu} & Open the worktree menu: make a worktree for a branch, or for a new one, show a worktree’s status, move or delete one.\\
none & \texttt{M-x git-\codebreak{}worktree-\codebreak{}move} & \texttt{git-\codebreak{}worktree-\codebreak{}move} & Move a worktree, asked for, not the main one, to a folder asked for.\\
none & \texttt{M-x git-\codebreak{}worktree-\codebreak{}status} & \texttt{git-\codebreak{}worktree-\codebreak{}status} & Show the status of a worktree, asked for, as C-x g shows a repository’s.\\
none & \texttt{M-x global-\codebreak{}centred-\codebreak{}text-\codebreak{}mode} & \texttt{global-\codebreak{}centred-\codebreak{}text-\codebreak{}mode} & Show the text in a column in the middle of the window, in every buffer, or stop, while the editor runs (\textbf{GLOBAL-CENTRED-TEXT-MODE} in the init file keeps it); a buffer that chose with C-x T o, and the modes \textbf{CENTRED-TEXT-MODES} names, keep theirs; with a count, on when it is positive, off when it is zero or less.\\
none & \texttt{M-x global-\codebreak{}cursor-\codebreak{}column-\codebreak{}mode} & \texttt{global-\codebreak{}cursor-\codebreak{}column-\codebreak{}mode} & Highlight the column point is in, down the window, in every buffer, or stop, while the editor runs (\textbf{GLOBAL-CURSOR-COLUMN-MODE} in the init file keeps it); a buffer that chose with C-x T c, and the modes \textbf{CURSOR-COLUMN-MODES} names, keep theirs; with a count, on when it is positive, off when it is zero or less.\\
none & \texttt{M-x global-\codebreak{}display-\codebreak{}line-\codebreak{}numbers-\codebreak{}mode} & \texttt{global-\codebreak{}display-\codebreak{}line-\codebreak{}numbers-\codebreak{}mode} & Show line numbers in the gutter, in every buffer, or stop, while the editor runs (\textbf{LINE-NUMBERS-DEFAULT} in the init file keeps it); a buffer that chose with M-x display-line-numbers-mode, and the modes \textbf{LINE-NUMBERS-MODES} names, keep theirs; with a count, on when it is positive, off when it is zero or less.\\
none & \texttt{M-x global-\codebreak{}hl-\codebreak{}line-\codebreak{}mode} & \texttt{global-\codebreak{}hl-\codebreak{}line-\codebreak{}mode} & Highlight the line point is on, in every buffer, or stop, while the editor runs (\textbf{GLOBAL-HL-LINE-MODE} in the init file keeps it); a buffer that chose with C-x T h, and the modes \textbf{HL-LINE-MODES} names, keep theirs; with a count, on when it is positive, off when it is zero or less.\\
none & \texttt{M-x global-\codebreak{}invisible-\codebreak{}characters-\codebreak{}mode} & \texttt{global-\codebreak{}invisible-\codebreak{}characters-\codebreak{}mode} & Show the characters that draw nothing as their codes, <FEFF> for the byte-order mark, in every buffer, or stop, while the editor runs (\textbf{GLOBAL-INVISIBLE-CHARACTERS-MODE} in the init file keeps it); a buffer that chose with C-x T i, and the modes \textbf{INVISIBLE-CHARACTERS-MODES} names, keep theirs; with a count, on when it is positive, off when it is zero or less.\\
none & \texttt{M-x global-\codebreak{}soft-\codebreak{}wrap-\codebreak{}mode} & \texttt{global-\codebreak{}soft-\codebreak{}wrap-\codebreak{}mode} & Wrap long lines at a word’s edge, at the fill column, in every buffer, or stop, while the editor runs (\textbf{GLOBAL-SOFT-WRAP-MODE} in the init file keeps it); a buffer that chose with C-x T l, and the modes \textbf{SOFT-WRAP-MODES} names, keep theirs; with a count, on when it is positive, off when it is zero or less.\\
none & \texttt{M-x global-\codebreak{}spellcheck-\codebreak{}mode} & \texttt{global-\codebreak{}spellcheck-\codebreak{}mode} & Mark the words spelt wrong, in every buffer, or stop, while the editor runs (\textbf{GLOBAL-SPELLCHECK-MODE} in the init file keeps it); a buffer that chose with M-x spellcheck-mode, and the modes \textbf{SPELLCHECK-MODES} names, keep theirs; with a count, on when it is positive, off when it is zero or less.\\
none & \texttt{M-x global-\codebreak{}text-\codebreak{}direction-\codebreak{}mode} & \texttt{global-\codebreak{}text-\codebreak{}direction-\codebreak{}mode} & Show right-to-left text in the order it is read, in every buffer, or stop, while the editor runs (\textbf{GLOBAL-TEXT-DIRECTION-MODE} in the init file keeps it); a buffer that chose with M-x text-direction-mode, and the modes \textbf{TEXT-DIRECTION-MODES} names, keep theirs; with a count, on when it is positive, off when it is zero or less.\\
none & \texttt{M-x global-\codebreak{}visual-\codebreak{}line-\codebreak{}mode} & \texttt{global-\codebreak{}visual-\codebreak{}line-\codebreak{}mode} & Mark the fill column with a line down the window, in every buffer, or stop, while the editor runs (\textbf{GLOBAL-VISUAL-LINE-MODE} in the init file keeps it); a buffer that chose with C-x T v, and the modes \textbf{VISUAL-LINE-MODES} names, keep theirs; with a count, on when it is positive, off when it is zero or less.\\
\texttt{:} & \texttt{M-g g}, \texttt{M-g M-g} & \texttt{goto-\codebreak{}line} & Move point to the start of a line number read from the minibuffer (1-based), the mark left where it was unless the region is active.\\
none & \texttt{M-x gpg-\codebreak{}choose-\codebreak{}encryption-\codebreak{}keys} & \texttt{gpg-\codebreak{}choose-\codebreak{}encryption-\codebreak{}keys} & Choose again whom the encrypted file shown is encrypted to, or a passphrase, in the minibuffer: from its next save it is encrypted so.\\
none & \texttt{M-x gpg-\codebreak{}copy-\codebreak{}to-\codebreak{}clipboard} & \texttt{gpg-\codebreak{}copy-\codebreak{}to-\codebreak{}clipboard} & Put the text copied last on the system clipboard, even text of an encrypted file, which is otherwise kept within Tuatara: other programs (clipboard managers among them) can then read it, and keep it.\\
none & \texttt{M-x gpg-\codebreak{}decrypt-\codebreak{}file} & \texttt{gpg-\codebreak{}decrypt-\codebreak{}file} & Decrypt the file shown (in a list of a directory, the one on the line) into a file asked for, the name without .gpg offered; after asking, since its text is then on disk where it can be read (made for you alone).\\
none & \texttt{M-x gpg-\codebreak{}decrypt-\codebreak{}region} & \texttt{gpg-\codebreak{}decrypt-\codebreak{}region} & Decrypt in place the armoured block (-----BEGIN PGP MESSAGE-----) in the text selected, or without a selection the one point is in, else the one nearest to point.\\
none & \texttt{M-x gpg-\codebreak{}decrypt-\codebreak{}with-\codebreak{}passphrase} & \texttt{gpg-\codebreak{}decrypt-\codebreak{}with-\codebreak{}passphrase} & Ask for the passphrase of the encrypted file this buffer waits for, shown as stars, and decrypt it with it.\\
none & \texttt{M-x gpg-\codebreak{}encrypt-\codebreak{}file} & \texttt{gpg-\codebreak{}encrypt-\codebreak{}file} & Encrypt the file shown (in a list of a directory, the one on the line; with C-u, or where there is none, asked for) into FILE.gpg, to the keys or with a passphrase chosen in the minibuffer, commas between several.\\
none & \texttt{M-x gpg-\codebreak{}encrypt-\codebreak{}region} & \texttt{gpg-\codebreak{}encrypt-\codebreak{}region} & Encrypt the text selected in place, as an armoured block (-----BEGIN PGP MESSAGE-----), to the keys or with a passphrase chosen in the minibuffer.\\
none & \texttt{M-x gpg-\codebreak{}import-\codebreak{}keys} & \texttt{gpg-\codebreak{}import-\codebreak{}keys} & Import the keys in a file asked for (armoured or not) into your keyring; say what was imported.\\
none & \texttt{M-x gpg-\codebreak{}import-\codebreak{}keys-\codebreak{}region} & \texttt{gpg-\codebreak{}import-\codebreak{}keys-\codebreak{}region} & Import the keys in the text selected (an armoured -----BEGIN PGP PUBLIC KEY BLOCK-----) into your keyring.\\
none & \texttt{M-x gpg-\codebreak{}keys-\codebreak{}details} & \texttt{gpg-\codebreak{}keys-\codebreak{}details} & Show the details of the key at the cursor: its fingerprint, its names, its subkeys and what they can do, when they expire.\\
none & \texttt{M-x gpg-\codebreak{}keys-\codebreak{}export} & \texttt{gpg-\codebreak{}keys-\codebreak{}export} & Export the public part of the keys marked (or the one at the cursor), armoured, to a file asked for.\\
none & \texttt{M-x gpg-\codebreak{}keys-\codebreak{}insert} & \texttt{gpg-\codebreak{}keys-\codebreak{}insert} & Put the public part of the keys marked (or the one at the cursor), armoured, at point in the buffer the list was opened from.\\
\bottomrule
\end{xltabular}
}

{\footnotesize
\begin{xltabular}{\textwidth}{@{}l >{\raggedright\arraybackslash\hsize=0.688\hsize}X l >{\raggedright\arraybackslash\hsize=1.312\hsize}X@{}}
\toprule
\textbf{Modal} & \textbf{Non-modal} & \textbf{Command} & \textbf{What it does}\\
\midrule
\endhead
none & \texttt{M-x gpg-\codebreak{}list-\codebreak{}keys} & \texttt{gpg-\codebreak{}list-\codebreak{}keys} & List your GnuPG keys and those you have of others: there m marks one, M-x gpg-keys-details shows its details, M-x gpg-import-keys imports keys from a file, M-x gpg-keys-export exports, M-x gpg-keys-insert inserts and D deletes those marked.\\
none & \texttt{M-x gpg-\codebreak{}lock} & \texttt{gpg-\codebreak{}lock} & Make gpg-agent forget every passphrase it keeps for a while (ten minutes, as it comes, its default-cache-ttl), of encrypted files and of your keys: the next to open one asks again.\\
none & \texttt{M-x gpg-\codebreak{}sign-\codebreak{}file} & \texttt{gpg-\codebreak{}sign-\codebreak{}file} & Sign the file shown (in a list of a directory, the one on the line) with your key: a detached signature (FILE.sig, the file left as it is), clear-signed text (FILE.asc) or a signed file (FILE.gpg), asked which.\\
none & \texttt{M-x gpg-\codebreak{}sign-\codebreak{}region} & \texttt{gpg-\codebreak{}sign-\codebreak{}region} & Clear-sign the text selected in place (-----BEGIN PGP SIGNED MESSAGE-----), with your key.\\
none & \texttt{M-x gpg-\codebreak{}verify-\codebreak{}file} & \texttt{gpg-\codebreak{}verify-\codebreak{}file} & Check the signature of the file shown (in a list of a directory, the one on the line): a detached signature (FILE.sig, FILE.asc) against its file, else the signature within the file; say who signed it, when, and how far that key is trusted.\\
none & \texttt{M-x gpg-\codebreak{}verify-\codebreak{}region} & \texttt{gpg-\codebreak{}verify-\codebreak{}region} & Check the signature of the clear-signed block (-----BEGIN PGP SIGNED MESSAGE-----) in the text selected, or without a selection the one point is in, else the one nearest to point: who signed it, when, and how far that key is trusted.\\
none & \texttt{M-x hex-\codebreak{}mode} & \texttt{hex-\codebreak{}mode} & Show the current buffer as a hex dump, or as text again: each row the offset, sixteen octets in hex and the text they make, a dot for an octet that is no printable ASCII character.\\
none & \texttt{M-x hex-\codebreak{}search-\codebreak{}backward} & \texttt{hex-\codebreak{}search-\codebreak{}backward} & Find the last place before the cursor that holds the octets asked for in hex, and show it in the hex view.\\
none & \texttt{M-x hex-\codebreak{}search-\codebreak{}forward} & \texttt{hex-\codebreak{}search-\codebreak{}forward} & Find the next place after the cursor that holds the octets asked for in hex (“de ad be ef”), and show it in the hex view.\\
none & \texttt{M-x hex-\codebreak{}set-\codebreak{}byte} & \texttt{hex-\codebreak{}set-\codebreak{}byte} & In the hex view, write one octet over the one at the cursor, asking for it in hex (41 is A), and move on to the next.\\
none & \texttt{M-x hex-\codebreak{}switch-\codebreak{}column} & \texttt{hex-\codebreak{}switch-\codebreak{}column} & In the hex view, go from the hex digits to the text beside them, or back: typing then writes characters there, or hex digits here, over the octets.\\
none & \texttt{C-x T h} & \texttt{hl-\codebreak{}line-\codebreak{}mode} & Highlight the line point is on in the current buffer, or stop: every row of it, out to the window’s edge, in the :hl-line face, under the syntax colours and the selection, in the selected window.\\
none & \texttt{M-x ielm} & \texttt{ielm} & Show \textbf{ielm}, where Lisp typed after the prompt is evaluated in the running editor: made with a prompt the first time, as it was after.\\
none & \texttt{C-c C-b} & \texttt{ielm-\codebreak{}change-\codebreak{}working-\codebreak{}buffer} & Choose the working buffer of \textbf{ielm}: the one the Lisp typed there sees as current, (insert “x”) putting text in it.\\
none & \texttt{C-c C-u} & \texttt{ielm-\codebreak{}kill-\codebreak{}input} & Take away the input after \textbf{ielm}’s last prompt.\\
none & \texttt{M-n} & \texttt{ielm-\codebreak{}next-\codebreak{}input} & Put the input after the one shown in \textbf{ielm}’s history in place of the input.\\
none & \texttt{M-p} & \texttt{ielm-\codebreak{}previous-\codebreak{}input} & Put the input before the one shown in \textbf{ielm}’s history in place of the input.\\
none & \texttt{RET} & \texttt{ielm-\codebreak{}return} & RET in \textbf{ielm}: after the last prompt, evaluate the input when it is whole expressions, else begin a new line of it; on an earlier input, evaluate that again.\\
none & \texttt{C-j} & \texttt{ielm-\codebreak{}send-\codebreak{}input} & Evaluate the input after \textbf{ielm}’s last prompt as it is, whole or not.\\
none & \texttt{TAB} & \texttt{ielm-\codebreak{}tab} & TAB in \textbf{ielm}: where only blanks come before point on its line of the input, indent the line as Lisp (never the prompt’s line); else complete the name at point from a list beside it, as C-M-i does.\\
none & \texttt{TAB} & \texttt{indent-\codebreak{}for-\codebreak{}tab-\codebreak{}command} & Indent the current line as its major mode indents it: point in the blanks the line begins with goes to where its text begins, point in its text stays with the text.\\
none & \texttt{C-M-\textbackslash{}} & \texttt{indent-\codebreak{}region} & Indent each line the selection touches as its major mode indents it; the lines inside a comment or a string begun on a line before keep their blanks, which are text.\\
\texttt{>} & \texttt{C-z >} & \texttt{indent-\codebreak{}selection-\codebreak{}or-\codebreak{}line} & Indent the lines the selection touches, or the current line, as the language’s mode indents them.\\
none & \texttt{M-x indent-\codebreak{}tabs-\codebreak{}mode} & \texttt{indent-\codebreak{}tabs-\codebreak{}mode} & Switch the current buffer between tab characters and spaces where TAB and indentation go on to a tab stop; with a count, tabs when it is positive, spaces when it is zero or less.\\
none & \texttt{C-x 8 RET} & \texttt{insert-\codebreak{}char} & Insert a character by its Unicode name or its code (C-x 8 RET).\\
none & \texttt{C-x 8 \&} & \texttt{insert-\codebreak{}digraph} & Type the character a digraph stands for, the count’s times: asks for its keys, e’ for e with an acute, -> for an arrow, << for a guillemet; M-x list-digraphs lists them all.\\
none & \texttt{M-x insert-\codebreak{}formatted-\codebreak{}backlinks-\codebreak{}block} & \texttt{insert-\codebreak{}formatted-\codebreak{}backlinks-\codebreak{}block} & Put in a dynamic block of links to the files with formatted names that link to this one (a block also written denote-backlinks).\\
none & \texttt{M-x insert-\codebreak{}formatted-\codebreak{}files-\codebreak{}as-\codebreak{}headings-\codebreak{}block} & \texttt{insert-\codebreak{}formatted-\codebreak{}files-\codebreak{}as-\codebreak{}headings-\codebreak{}block} & Put in a dynamic block holding the text of the Org files with formatted names whose names a regular expression matches, each under a headline of its title and tags, its own headlines a level down (a block also written denote-files-as-headings): the regular expression and the order of the files asked for.\\
none & \texttt{M-x insert-\codebreak{}formatted-\codebreak{}files-\codebreak{}block} & \texttt{insert-\codebreak{}formatted-\codebreak{}files-\codebreak{}block} & Put in a dynamic block holding the text of the files with formatted names whose names a regular expression matches (a block also written denote-files): the regular expression and the order of the files asked for.\\
none & \texttt{M-x insert-\codebreak{}formatted-\codebreak{}links-\codebreak{}block} & \texttt{insert-\codebreak{}formatted-\codebreak{}links-\codebreak{}block} & Put in a dynamic block of links to the files with formatted names whose names a regular expression matches (a block also written denote-links): asked for, the files it matches listed as it is typed; the links labelled with the files’ signatures and titles.\\
none & \texttt{M-x insert-\codebreak{}formatted-\codebreak{}missing-\codebreak{}links-\codebreak{}block} & \texttt{insert-\codebreak{}formatted-\codebreak{}missing-\codebreak{}links-\codebreak{}block} & Put in a dynamic block of links to the files with formatted names whose names a regular expression matches and that this file does not link to yet, (a block also written denote-missing-links): the regular expression asked for, the files it matches listed as it is typed.\\
none & \texttt{C-x i} & \texttt{insert-\codebreak{}formatted-\codebreak{}note-\codebreak{}template} & Read a note template, or any file, into the current buffer at point: one of your templates (blank, and each file in \textasciitilde{}/.config/tuatara/note-templates/), \%date\% and \%author\% filled in, or with “Another file\ldots{}” any file, as it is.\\
none & \texttt{M-x insert-\codebreak{}formatted-\codebreak{}sequence-\codebreak{}block} & \texttt{insert-\codebreak{}formatted-\codebreak{}sequence-\codebreak{}block} & Put in a dynamic block listing a file with a sequence for its signature and those under it (a block also written denote-sequence): the file asked for, and how many levels to list, its own the first.\\
none & \texttt{M-x insert-\codebreak{}keys-\codebreak{}counter} & \texttt{insert-\codebreak{}keys-\codebreak{}counter} & Insert the counter at point, written as M-x set-keys-counter-format set (the number alone as standard), and add 1 to it, or with a prefix argument N, add N.\\
none & \texttt{M-(} & \texttt{insert-\codebreak{}pair} & Put the pair of the delimiter of the key pressed around the region (M-( for parentheses), point on the first character inside; when the pair stands just outside the region already, take it off.\\
none & \texttt{M-x insert-\codebreak{}pasted-\codebreak{}text} & \texttt{insert-\codebreak{}pasted-\codebreak{}text} & Put in the text just pasted into the terminal Tuatara runs in, as text: none of it is taken for keys or commands, whatever it holds, in any editing state.\\
none & \texttt{M-x insert-\codebreak{}recorded-\codebreak{}keys-\codebreak{}definition} & \texttt{insert-\codebreak{}recorded-\codebreak{}keys-\codebreak{}definition} & Insert at point the Lisp text that makes the keys recorded last a command again and binds it to a key, both asked for, for an init file: a name (the one M-x name-recorded-keys gave, as standard) and a key (nothing typed for none).\\
none & \texttt{C-c y} & \texttt{insert-\codebreak{}snippet} & Put in a snippet chosen by its name among those of the current buffer’s mode (each listed with its key), at point; with the selection active, in its place, the selected text where it has \$TM\_SELECTED\_TEXT.\\
none & \texttt{M-x insert-\codebreak{}tab} & \texttt{insert-\codebreak{}tab} & Insert a tab character at point, whatever the tab settings say; C-z TAB (tab-to-tab-stop) goes on to the next tab stop as they say.\\
none & \texttt{C-x T i} & \texttt{invisible-\codebreak{}characters-\codebreak{}mode} & Show the characters that draw nothing in the current buffer as their codes, or stop: the byte-order mark as <FEFF>, a zero-width space as <200B>, the soft hyphen, the joiners, the word joiner, the marks that set the direction text is drawn in and the tags the same way, in the :invisible-character face, each taking the columns its code takes.\\
\bottomrule
\end{xltabular}
}

\begin{xltabular}{\textwidth}{@{}l >{\raggedright\arraybackslash\hsize=0.737\hsize}X l >{\raggedright\arraybackslash\hsize=1.263\hsize}X@{}}
\toprule
\textbf{Modal} & \textbf{Non-modal} & \textbf{Command} & \textbf{What it does}\\
\midrule
\endhead
none & \texttt{M-\textasciicircum{}} & \texttt{join-\codebreak{}line} & Join this line to the one before it: the line end between them goes, and the spaces about it become one space, or none where a line begins or ends, before a closing bracket or after an opening one.\\
\texttt{g j} & none & \texttt{join-\codebreak{}selected-\codebreak{}lines} & Pull the following line up onto this one, COUNT times; with a selection, join the lines it touches into one, in the background when it holds more than \textbf{REWRITE-THRESHOLD} octets.\\
none & \texttt{M-x journal-\codebreak{}new-\codebreak{}entry} & \texttt{journal-\codebreak{}new-\codebreak{}entry} & Make a new journal entry for today, whatever entries there are; with C-u, for a date asked for.\\
none & \texttt{C-x n j} & \texttt{journal-\codebreak{}new-\codebreak{}or-\codebreak{}existing-\codebreak{}entry} & Open today’s journal entry, or make it when there is none: in \textbf{journal-directory}, named and titled as \textbf{journal-file-name} and \textbf{journal-title-format} say, saved at once; with several for today, the one chosen.\\
\texttt{S} & \texttt{C-z S} & \texttt{jump-\codebreak{}back} & Rotate to the next stored position in the ring and jump there; press again to cycle.\\
none & \texttt{M-x just-\codebreak{}one-\codebreak{}space} & \texttt{just-\codebreak{}one-\codebreak{}space} & Leave one space in place of the spaces and tabs around point, point after it; with a count N, N spaces; a negative count takes the line ends among them too, leaving -N spaces.\\
none & \texttt{M-x keep-\codebreak{}lines} & \texttt{keep-\codebreak{}lines} & Delete the lines the selection touches in which a regular expression, asked for in the minibuffer, finds no match, keeping those it matches.\\
none & \texttt{M-x keyboard-\codebreak{}quit} & \texttt{keyboard-\codebreak{}quit} & Abandon the current operation and deactivate the region; what stays shown until then goes, such as the matches of the last search.\\
none & \texttt{C-x k} & \texttt{kill-\codebreak{}buffer} & Kill the current buffer, asking first when it has changes not saved.\\
none & \texttt{C-x 4 0} & \texttt{kill-\codebreak{}buffer-\codebreak{}and-\codebreak{}window} & Kill the selected window’s buffer, asking first when it has changes not saved, and close the window when there are others.\\
\texttt{D} & \texttt{C-k} & \texttt{kill-\codebreak{}line} & Kill the rest of the current line; at the end of a line, kill the newline.\\
none & \texttt{C-x p k} & \texttt{kill-\codebreak{}project-\codebreak{}buffers} & Kill every buffer of the project the current buffer is in (a folder above it holding one of \textbf{PROJECT-ROOT-MARKERS}, .git or .project), its files, its folders listed (C-x d), and the buffers standing in them; after asking, saying how many are modified, each modified one asked about as it goes.\\
none & \texttt{C-x r k} & \texttt{kill-\codebreak{}rectangle} & Kill the rectangle between the mark and point: the block of columns from one to the other on each line between, a short line’s part padded with spaces, goes to the rectangle store that C-x r y pastes from.\\
none & \texttt{C-w} & \texttt{kill-\codebreak{}region} & Kill the text between point and mark; a rectangle (C-x r SPC) is killed as a block, which C-y pastes back as one.\\
none & \texttt{M-w} & \texttt{kill-\codebreak{}ring-\codebreak{}save} & Copy the text between point and mark to the kill ring without deleting it; a rectangle (C-x r SPC) is copied as a block, which C-y pastes as one.\\
none & \texttt{M-x kill-\codebreak{}sexp} & \texttt{kill-\codebreak{}sexp} & Kill the balanced expression after point: to where forward-sexp goes.\\
none & \texttt{C-S-\codebreak{}DEL} & \texttt{kill-\codebreak{}whole-\codebreak{}line} & Kill the current line and its line end; with a count N, N lines from this one; with a negative count, -N lines back, with the line end before them rather than after; with 0, the line without its line end.\\
none & \texttt{M-x lisp-\codebreak{}scratch} & \texttt{lisp-\codebreak{}scratch} & Show \textbf{lisp-scratch}, the Lisp buffer that is not saved, where C-x C-e evaluates in the running editor and C-j puts what the expression before point printed and its value on the lines after it; when it was killed, a new one.\\
none & \texttt{M-x list-\codebreak{}auto-\codebreak{}saves} & \texttt{list-\codebreak{}auto-\codebreak{}saves} & List the auto-saved copies of changes not saved (\textbf{AUTO-SAVE-DIRECTORY}), the newest first: after a crash or a quit without saving, what can be taken back.\\
none & \texttt{D}, \texttt{C-c D} & \texttt{list-\codebreak{}delete-\codebreak{}marked} & Delete the marked entries of the list (>), or with none marked the one on the current line, after asking, as the list deletes, as M-x list-execute does those marked for deletion; in rimu, with C-u, for good rather than to the trash.\\
none & \texttt{C-x / d} & \texttt{list-\codebreak{}delete-\codebreak{}records} & Delete the records whose line a regular expression matches, keeping the rest: asked for as C-x / k asks, in the same files.\\
none & \texttt{M-x list-\codebreak{}digraphs} & \texttt{list-\codebreak{}digraphs} & Show every digraph: the keys that type a character a keyboard lacks, those English text wants most first, then all of them by script, with the code point and name of each, and how to type them.\\
none & \texttt{x} & \texttt{list-\codebreak{}execute} & Delete the entries of the list marked for deletion (D), after asking, as the list deletes: the buffer list kills its buffers, rimu puts files in the trash, the list of jobs drops those waiting and stops those running.\\
none & \texttt{M-x list-\codebreak{}faces} & \texttt{list-\codebreak{}faces} & List every face a theme can colour in \textbf{Faces}, in sections by the first word of their names: each face’s name, a sample drawn in it as the theme in force draws it, and the first sentence of what it colours.\\
none & \texttt{C-x j} & \texttt{list-\codebreak{}jobs} & List the jobs in the background in \textbf{Jobs}: in each buffer those running, with how far they have got, and those waiting their turn, in the order they will run, their text painted in the buffer meanwhile.\\
none & \texttt{C-x / k} & \texttt{list-\codebreak{}keep-\codebreak{}records} & Keep the records whose line a regular expression matches and delete the rest, in a file of records, a line each: a CSV or TSV, a log, a JSON Lines file, todo.txt and the others \textbf{LINE-RECORD-MODES} names, the column asked for first as C-x / f asks.\\
none & \texttt{m}, \texttt{C-c m} & \texttt{list-\codebreak{}mark-\codebreak{}entry} & Mark the entry on the current line of the list (>), for the list’s commands that act on those marked, and move down; with a count, as many entries, up the list for a negative one.\\
none & \texttt{d}, \texttt{C-c d} & \texttt{list-\codebreak{}mark-\codebreak{}for-\codebreak{}deletion} & Mark the entry on the current line of the list for deletion (D), for the list’s command that deletes them, and move down; with a count, as many entries, up the list for a negative one.\\
none & \texttt{n}, \texttt{C-c n} & \texttt{list-\codebreak{}next-\codebreak{}entry} & Go to the next entry of the list, the start of its first line, passing over the headings of its sections.\\
none & \texttt{p}, \texttt{C-c p} & \texttt{list-\codebreak{}previous-\codebreak{}entry} & Go to the entry of the list before the one on the current line, the start of its first line, passing over the headings of its sections.\\
none & \texttt{S-TAB} & \texttt{list-\codebreak{}toggle-\codebreak{}all-\codebreak{}folds} & Fold every section of the list, or unfold them all when every one is folded; a section not in the list now keeps its fold for when it comes back.\\
none & \texttt{TAB} & \texttt{list-\codebreak{}toggle-\codebreak{}fold} & Fold the section of the list the current line is in, its heading left showing, or unfold it when it is folded.\\
none & \texttt{C-x / t} & \texttt{list-\codebreak{}toggle-\codebreak{}header} & Take the first line of a file of records as its header, or as a record like the rest; the same again switches back; with a count, a header when it is positive, a record when it is zero or less.\\
none & \texttt{C-x / L} & \texttt{list-\codebreak{}toggle-\codebreak{}header-\codebreak{}lock} & Keep the first line of a file of records, its header, at the top of the window as the rest scrolls under it, coloured as it is below; the same again lets it scroll away; with a count, kept when it is positive, let go when it is zero or less.\\
none & \texttt{t} & \texttt{list-\codebreak{}toggle-\codebreak{}marks} & Mark the entries the list shows that are not marked and take the mark off those marked; those marked for deletion (D), or given a mark of another kind, stay so, as do the entries a filter or a folded section hides.\\
none & \texttt{P}, \texttt{C-c P} & \texttt{list-\codebreak{}toggle-\codebreak{}preview} & Turn the preview beside this list on or off: the entry on the current line shown in a window beside the list, on its right when there is room, else below it, following the cursor from line to line.\\
none & \texttt{U}, \texttt{C-c U} & \texttt{list-\codebreak{}unmark-\codebreak{}all} & Take every mark off the entries of the list, of any kind; in rimu, those of the files marked in other directories too.\\
none & \texttt{u}, \texttt{C-c u} & \texttt{list-\codebreak{}unmark-\codebreak{}entry} & Take the mark off the entry on the current line of the list, of any kind, and move down; with a count, off as many entries, up the list for a negative one.\\
none & \texttt{l}, \texttt{RET}, \texttt{C-c l} & \texttt{list-\codebreak{}visit} & Open the entry on the current line of the list: the buffer, the bookmark’s place, the file; on a section’s heading, fold or unfold the section.\\
none & \texttt{o}, \texttt{C-c o} & \texttt{list-\codebreak{}visit-\codebreak{}other} & Open the entry on the current line of the list in another window; on a section’s heading, fold or unfold the section.\\
\bottomrule
\end{xltabular}

\begin{xltabular}{\textwidth}{@{}l >{\raggedright\arraybackslash\hsize=0.650\hsize}X l >{\raggedright\arraybackslash\hsize=1.350\hsize}X@{}}
\toprule
\textbf{Modal} & \textbf{Non-modal} & \textbf{Command} & \textbf{What it does}\\
\midrule
\endhead
none & \texttt{C-M-r} & \texttt{literal-\codebreak{}search-\codebreak{}backward} & Search backward as you type for literal text; see literal-search-forward.\\
none & \texttt{C-M-s} & \texttt{literal-\codebreak{}search-\codebreak{}forward} & Search forward as you type for literal text: keys as in regex-search-forward, nothing typed special.\\
none & \texttt{M-x load-\codebreak{}theme} & \texttt{load-\codebreak{}theme} & Choose a colour theme, with each theme’s description beside it.\\
none & \texttt{C-c L} & \texttt{lsp} & Start the language servers of the current buffer’s language for its project, after asking whether you trust the project: a server runs the project’s code (build scripts, macros, plugins).\\
none & \texttt{C-c a} & \texttt{lsp-\codebreak{}code-\codebreak{}actions} & Offer the language server’s fixes and changes for the selection, or for point: those for what it found wrong there first; the one chosen is made.\\
none & \texttt{C-c k} & \texttt{lsp-\codebreak{}code-\codebreak{}lens} & Run what a language server says of the current line (a code lens: Run test, 3 references), chosen when there are several, as the server’s own command.\\
none & \texttt{C-c K} & \texttt{lsp-\codebreak{}code-\codebreak{}lenses-\codebreak{}mode} & Show what language servers say of lines (3 references, Run test) faint at each line’s end, or no longer; M-x lsp-code-lens runs one.\\
none & \texttt{C-c h} & \texttt{lsp-\codebreak{}describe-\codebreak{}thing-\codebreak{}at-\codebreak{}point} & Show what the language server knows of the thing at point (its type, its documentation) in \textbf{Help}.\\
none & \texttt{C-c l} & \texttt{lsp-\codebreak{}diagnostics} & List what the language servers found wrong in the open files, a line each, then in the files of their projects not open, in a window across the bottom; RET there goes to it, opening the file.\\
none & \texttt{M-x lsp-\codebreak{}events} & \texttt{lsp-\codebreak{}events} & List the last messages to and from the language servers the current buffer is open in (all those running when it is in none), the newest last, and what each said on its error output: to see what a server was asked and answered when it does not do what it should; the last 500 of each are kept.\\
none & \texttt{C-c I} & \texttt{lsp-\codebreak{}find-\codebreak{}implementations} & Go to the implementations of the thing at point (of a trait, an interface, an abstract method), as its language server finds them: there at once when it is one, else listed in \textbf{xref}.\\
none & \texttt{C-c r} & \texttt{lsp-\codebreak{}find-\codebreak{}references} & List where the thing at point is used, as its language server finds it, in \textbf{xref}: RET goes to one.\\
none & \texttt{C-c g} & \texttt{lsp-\codebreak{}find-\codebreak{}symbol} & Go to a function, type or other name of the current buffer’s project, as its language server finds them by what you type (the name at point offered): each is listed with what it is and where, and the one chosen is gone to.\\
none & \texttt{C-c T} & \texttt{lsp-\codebreak{}find-\codebreak{}type-\codebreak{}definition} & Go to the definition of the type of the thing at point (a variable’s type, a function’s result), as its language server finds it.\\
none & \texttt{C-c z} & \texttt{lsp-\codebreak{}fold} & Hide the lines of the block the cursor is in (a function’s body, a struct, a run of imports), as its language server knows them, under the block’s first line; on a line whose block is hidden, show it again.\\
none & \texttt{C-c Z} & \texttt{lsp-\codebreak{}fold-\codebreak{}all} & Hide every outermost block of the file, as its language server knows them, under its first line; with any hidden, show all again.\\
none & \texttt{M-x lsp-\codebreak{}forget-\codebreak{}project} & \texttt{lsp-\codebreak{}forget-\codebreak{}project} & Trust the current buffer’s project no more: its language servers end, and no longer start by themselves.\\
none & \texttt{C-c c} & \texttt{lsp-\codebreak{}incoming-\codebreak{}calls} & List the places that call the function at point, as its language server knows them, in \textbf{xref}: RET goes to one.\\
none & \texttt{C-c i} & \texttt{lsp-\codebreak{}inlay-\codebreak{}hints-\codebreak{}mode} & Show the language servers’ inlay hints within the code, or no longer: the types of bindings, the names of the parameters arguments go to.\\
none & \texttt{C-c n} & \texttt{lsp-\codebreak{}next-\codebreak{}diagnostic} & Go to the next place the language servers found something wrong, and say what.\\
none & \texttt{C-c o} & \texttt{lsp-\codebreak{}organise-\codebreak{}imports} & Have the language server put the file’s imports in order, as its own action for that does it (unused ones taken out, where it does that), once it answers, when the file is still as it was asked of.\\
none & \texttt{C-c C} & \texttt{lsp-\codebreak{}outgoing-\codebreak{}calls} & List the functions the function at point calls, as its language server knows them, in \textbf{xref}: RET goes to one.\\
none & \texttt{C-c p} & \texttt{lsp-\codebreak{}previous-\codebreak{}diagnostic} & Go to the place before point where the language servers found something wrong, and say what.\\
none & \texttt{C-c R} & \texttt{lsp-\codebreak{}rename} & Rename the thing at point everywhere the language server knows it is used, the new name asked for: each file’s changes one change undo takes back, nothing saved (C-x s saves them).\\
none & \texttt{M-x lsp-\codebreak{}restart} & \texttt{lsp-\codebreak{}restart} & End the language servers the current buffer is open in and start them again; in a project not trusted, none is started (M-x lsp asks first).\\
none & \texttt{C-c S} & \texttt{lsp-\codebreak{}semantic-\codebreak{}tokens-\codebreak{}mode} & Colour code as its language server sees it, or no longer: each name in the face of its kind (a parameter, a field, a macro), over the mode’s colours.\\
none & \texttt{M-x lsp-\codebreak{}shutdown} & \texttt{lsp-\codebreak{}shutdown} & End the language servers the current buffer is open in; M-x lsp starts them again.\\
none & \texttt{C-x 5 2} & \texttt{make-\codebreak{}frame} & Open another frame showing the current buffer and work in it: a window of its own on the desktop, or in a terminal a whole screen of its own shown in place of the one before, the frames taking turns on the terminal, the mode line naming the one shown (F1, F2).\\
none & \texttt{M-x man} & \texttt{man} & Show a manual page in a buffer of its own, read-only: ask for it by name (ls), with its section (ls(1)) or by section and name (1 ls), completing on the names of the pages the system has.\\
none & \texttt{h}, \texttt{C-c h} & \texttt{man-\codebreak{}back} & Show the page a link was followed from; close the window when this page was not reached by a link.\\
none & \texttt{n}, \texttt{C-c n} & \texttt{man-\codebreak{}next-\codebreak{}section} & Go to the start of the next section heading of the manual page.\\
none & \texttt{p}, \texttt{C-c p} & \texttt{man-\codebreak{}previous-\codebreak{}section} & Go to the start of the section heading before the line point is on.\\
none & \texttt{M-x mark-\codebreak{}defun} & \texttt{mark-\codebreak{}defun} & Select the definition point is in, or the one after: point at its start, with the comments right above it and the blank lines above those, and the mark after it.\\
none & \texttt{M-x mark-\codebreak{}end-\codebreak{}of-\codebreak{}sentence} & \texttt{mark-\codebreak{}end-\codebreak{}of-\codebreak{}sentence} & Set the mark at the end of the sentence, point staying, and activate the region; with a count, that many sentences.\\
none & \texttt{M-h} & \texttt{mark-\codebreak{}paragraph} & Select the paragraph at point: point at the blank line before it (or the start of the text), the mark at the blank line after it.\\
none & \texttt{M-x mark-\codebreak{}run-\codebreak{}adopt} & \texttt{mark-\codebreak{}run-\codebreak{}adopt} & Take the selection into a mark run: its keys then grow it from where it is.\\
none & \texttt{M-x mark-\codebreak{}run-\codebreak{}buffer-\codebreak{}end} & \texttt{mark-\codebreak{}run-\codebreak{}buffer-\codebreak{}end} & Stretch the run’s end on to the end of the text, or with no run go there.\\
none & \texttt{M-x mark-\codebreak{}run-\codebreak{}buffer-\codebreak{}start} & \texttt{mark-\codebreak{}run-\codebreak{}buffer-\codebreak{}start} & Stretch the run’s start back to the start of the text, or with no run go there.\\
none & \texttt{M-x mark-\codebreak{}run-\codebreak{}cancel} & \texttt{mark-\codebreak{}run-\codebreak{}cancel} & Drop the selection and end the mark run.\\
none & \texttt{M-x mark-\codebreak{}run-\codebreak{}down} & \texttt{mark-\codebreak{}run-\codebreak{}down} & Move point down a line, or the count’s, without ending the mark run.\\
\bottomrule
\end{xltabular}

{\small
\begin{xltabular}{\textwidth}{@{}l >{\raggedright\arraybackslash\hsize=0.690\hsize}X l >{\raggedright\arraybackslash\hsize=1.310\hsize}X@{}}
\toprule
\textbf{Modal} & \textbf{Non-modal} & \textbf{Command} & \textbf{What it does}\\
\midrule
\endhead
none & \texttt{M-x mark-\codebreak{}run-\codebreak{}exchange} & \texttt{mark-\codebreak{}run-\codebreak{}exchange} & Swap point and the mark, so that h j k l adjust the run’s other end.\\
none & \texttt{M-x mark-\codebreak{}run-\codebreak{}left} & \texttt{mark-\codebreak{}run-\codebreak{}left} & Move point back a character, or the count’s, without ending the mark run.\\
none & \texttt{M-x mark-\codebreak{}run-\codebreak{}line-\codebreak{}backward} & \texttt{mark-\codebreak{}run-\codebreak{}line-\codebreak{}backward} & Select the line point is on, or grow the run’s start up by whole lines: a run starting mid-line first takes the start of that line.\\
none & \texttt{M-x mark-\codebreak{}run-\codebreak{}line-\codebreak{}end} & \texttt{mark-\codebreak{}run-\codebreak{}line-\codebreak{}end} & Stretch the run’s end on to the end of its line, or with no run go there.\\
none & \texttt{M-x mark-\codebreak{}run-\codebreak{}line-\codebreak{}forward} & \texttt{mark-\codebreak{}run-\codebreak{}line-\codebreak{}forward} & Select the line point is on, its line end included, or grow the run’s end down by whole lines: a run ending mid-line first takes the rest of that line.\\
none & \texttt{M-x mark-\codebreak{}run-\codebreak{}line-\codebreak{}start} & \texttt{mark-\codebreak{}run-\codebreak{}line-\codebreak{}start} & Stretch the run’s start back to the start of its line, or with no run go there.\\
none & \texttt{M-x mark-\codebreak{}run-\codebreak{}right} & \texttt{mark-\codebreak{}run-\codebreak{}right} & Move point forward a character, or the count’s, without ending the mark run.\\
none & \texttt{M-x mark-\codebreak{}run-\codebreak{}step-\codebreak{}back} & \texttt{mark-\codebreak{}run-\codebreak{}step-\codebreak{}back} & Put the mark run back where its last press found it; U puts it forward again.\\
none & \texttt{M-x mark-\codebreak{}run-\codebreak{}step-\codebreak{}forward} & \texttt{mark-\codebreak{}run-\codebreak{}step-\codebreak{}forward} & Put the mark run forward where u stepped it back from.\\
\texttt{M} & \texttt{C-z M} & \texttt{mark-\codebreak{}run-\codebreak{}toggle} & Start a mark run: the m prefix held down.\\
none & \texttt{M-x mark-\codebreak{}run-up} & \texttt{mark-\codebreak{}run-up} & Move point up a line, or the count’s, without ending the mark run.\\
none & \texttt{C-M-@}, \texttt{C-M-\codebreak{}SPC} & \texttt{mark-\codebreak{}sexp} & Set the mark after the next balanced expression, point staying, and activate it; with a count, that many expressions, or before point when negative.\\
none & \texttt{C-x h} & \texttt{mark-\codebreak{}whole-\codebreak{}buffer} & Select the whole buffer: mark at the end, point at the start; where point was goes on the mark ring, with the mark before.\\
none & \texttt{M-@} & \texttt{mark-\codebreak{}word} & Set the mark at the end of the next word, point staying, and activate the region; with a count, that many words.\\
none & \texttt{C-c C-o} & \texttt{markdown-\codebreak{}follow-\codebreak{}link-\codebreak{}at-\codebreak{}point} & Follow the Markdown link point is on (C-c C-o, and RET in NORMAL state): to the heading a table of contents’s link (\#an-anchor) leads to, a web or mail address in the desktop’s program after asking, or a file, visited.\\
none & \texttt{M-x match-\codebreak{}denote-\codebreak{}files-\codebreak{}link} & \texttt{match-\codebreak{}denote-\codebreak{}files-\codebreak{}link} & Insert a link to your files with formatted names whose names match what is asked for, the files it matches listed as it is typed, as a denote: link, which other programs reading those links follow too, written LABEL; in Markdown [LABEL](denote:query-filenames:TEXT) and in text files <denote:query-filenames:TEXT>.\\
none & \texttt{C-x n l} & \texttt{match-\codebreak{}formatted-\codebreak{}files-\codebreak{}link} & Insert a link to your files with formatted names whose names match a regular expression, asked for, the files it matches listed as it is typed: LABEL.\\
none & \texttt{M-x menu-\codebreak{}describe} & \texttt{menu-\codebreak{}describe} & List in \textbf{Help} every entry of the open menu: its keys, what it says and the command or argument behind it, those not shown now too.\\
none & \texttt{M-x menu-\codebreak{}leave} & \texttt{menu-\codebreak{}leave} & Leave the open menu, asked first when anything in it was changed; it stays open unless the answer is yes.\\
none & \texttt{M-x menu-\codebreak{}set-\codebreak{}option} & \texttt{menu-\codebreak{}set-\codebreak{}option} & Give the option of the open menu whose key was typed a value, asked for, or take its value away when it has one, or ask it again with that value when the option says so (MENU-ENTRY-REPLACE), or say it is left out of the menu now; the menu stays open.\\
none & \texttt{M-x menu-\codebreak{}toggle-\codebreak{}switch} & \texttt{menu-\codebreak{}toggle-\codebreak{}switch} & Turn the switch of the open menu whose key was typed on, or off, or say it is left out of the menu now; the menu stays open.\\
none & \texttt{<drag-\codebreak{}mouse-1>} & \texttt{mouse-\codebreak{}drag-\codebreak{}region} & Select from where the first mouse button went down to where the mouse is now, point following it, in the window it went down in; held above or below that window’s text, it scrolls a line at a time.\\
none & \texttt{<up-\codebreak{}mouse-1>} & \texttt{mouse-\codebreak{}release} & The first mouse button let go: a drag ends there, what it selected staying selected and becoming the desktop’s primary selection, which the middle button pastes in other programs.\\
none & \texttt{<triple-\codebreak{}mouse-1>} & \texttt{mouse-\codebreak{}select-\codebreak{}line} & Select the line the first mouse button went down on three times, its line break too.\\
none & \texttt{<double-\codebreak{}mouse-1>} & \texttt{mouse-\codebreak{}select-\codebreak{}word} & Select the word the first mouse button went down on twice, point at its end: the run of letters, digits and \_ there, or of spaces, or of other signs; on a bracket, the bracketed group.\\
none & \texttt{<down-\codebreak{}mouse-1>} & \texttt{mouse-\codebreak{}set-\codebreak{}point} & Put point where the first mouse button went down, selecting the window there, and let go of the selection; dragging then selects from there.\\
none & \texttt{C-<wheel-\codebreak{}up>}, \texttt{C-<wheel-\codebreak{}down>} & \texttt{mouse-\codebreak{}wheel-\codebreak{}text-\codebreak{}scale} & Make the window’s font a step larger as the mouse wheel turns up with Ctrl held (C-<wheel-up>), a step smaller as it turns down.\\
none & \texttt{<mouse-2>} & \texttt{mouse-\codebreak{}yank-\codebreak{}primary} & Insert where the middle button was pressed the text selected last in any program (the desktop’s primary selection); the mark at its start.\\
none & \texttt{M-x name-\codebreak{}recorded-\codebreak{}keys} & \texttt{name-\codebreak{}recorded-\codebreak{}keys} & Give the keys recorded last a name, asked for, which then plays them as a command: M-x NAME, or a key bound to it, as C-x e does, a prefix argument as with it.\\
none & \texttt{C-x x n} & \texttt{narrow-\codebreak{}to-\codebreak{}region} & Show only the selected text, and keep every command to it: motion, search and replacement, editing and selecting stop at its edges, and the rest of the buffer is out of sight and out of reach until C-x x w.\\
none & \texttt{C-x x s} & \texttt{narrow-\codebreak{}to-\codebreak{}section} & Show only the section point is in, and keep every command to it, as C-x x n does for the selection: in Org, the subtree of the headline point is under, the headline with it; in code, the definition point is in (or the one after it), as C-M-a and C-M-e find it.\\
none & \texttt{C-x n e} & \texttt{new-\codebreak{}empty-\codebreak{}buffer} & Show a new empty buffer, untitled (untitled<2> beside one already there), for text that has no file yet: C-x C-s asks where to save it.\\
none & \texttt{C-x n n} & \texttt{new-\codebreak{}file-\codebreak{}formatted} & Make a new file with a formatted name in \textbf{FORMATTED-FILE-DIRECTORY}: asks what \textbf{FORMATTED-FILE-PROMPTS} says (as standard the title, the subdirectory, the keywords, the kind of file and the template), and names it as find-file-formatted does, with the header for its kind.\\
none & \texttt{C-j}, \texttt{RET} & \texttt{newline} & Insert a line end at point: a newline, or in a buffer whose lines end with a carriage return and a newline (a Windows file), both; with a count, that many, and none for 0.\\
\texttt{j} & \texttt{C-n}, \texttt{<down>} & \texttt{next-\codebreak{}line} & Move point down one line (or the prefix argument’s worth), keeping the column when possible; where long lines wrap at a word’s edge (C-x T l), one row.\\
none & \texttt{M-g n} & \texttt{next-\codebreak{}place} & Go to the next place of the list of places used last: the next line naming a file and line in a compile’s output (M-x compile), the next line a search of files found (M-s g), or the next place the language servers found something wrong in a buffer.\\
none & \texttt{C-x r N} & \texttt{number-\codebreak{}rectangle-\codebreak{}lines} & Number the lines of the rectangle between the mark and point: 1, 2, 3 and on, each with a space after it, put in at the block’s left column, a line shorter than that padded out to it; the numbers are right-aligned to the width of the last, and the text of the block moves right to make room.\\
none & \texttt{M-x number-\codebreak{}rectangle-\codebreak{}lines-\codebreak{}replace} & \texttt{number-\codebreak{}rectangle-\codebreak{}lines-\codebreak{}replace} & Replace the rectangle between the mark and point with numbers, 1, 2, 3 and on down its lines, nothing after them: to renumber a list, or to turn its bullets into numbers.\\
none & \texttt{M-x office-\codebreak{}view-\codebreak{}pages} & \texttt{office-\codebreak{}view-\codebreak{}pages} & Show the pages of the Office document read as text here, in its place.\\
none & \texttt{M-x open-\codebreak{}dropped-\codebreak{}files} & \texttt{open-\codebreak{}dropped-\codebreak{}files} & Visit the files dropped on the window, in the window they were dropped on, as C-x C-f does, the last one shown; safely, as C-u C-x C-f does, when \textbf{OPEN-DROPPED-FILES-SAFELY} is true.\\
\bottomrule
\end{xltabular}
}

\begin{xltabular}{\textwidth}{@{}l >{\raggedright\arraybackslash\hsize=0.700\hsize}X l >{\raggedright\arraybackslash\hsize=1.300\hsize}X@{}}
\toprule
\textbf{Modal} & \textbf{Non-modal} & \textbf{Command} & \textbf{What it does}\\
\midrule
\endhead
none & \texttt{O}, \texttt{C-c O} & \texttt{open-\codebreak{}externally} & Open the file shown, in the desktop’s own program for its kind (xdg-open), or the one \textbf{OPEN-EXTERNALLY-PROGRAMS} gives its Multipurpose Internet Mail Extensions (MIME) type, for what Tuatara does not do: printing, forms, playing.\\
none & \texttt{C-x n o} & \texttt{open-\codebreak{}file-\codebreak{}formatted} & Open one of your files with formatted names, those in \textbf{FORMATTED-FILE-DIRECTORY} and its subdirectories, choosing among them by any part of the name: listed under their subdirectory or their kind, the ones changed last first, each as its date, its title and its keywords, with its size and when it changed at the right.\\
none & \texttt{C-o} & \texttt{open-\codebreak{}line} & Insert a line end after point, point staying before it; with a count, that many line ends.\\
none & \texttt{C-x C-p} & \texttt{open-\codebreak{}path-\codebreak{}at-\codebreak{}point} & Open the file or folder whose path is at the cursor in the text, as C-x C-f opens it: a folder (\emph{etc}, \textasciitilde{}/notes/) in a list of it, a file \textbackslash{}(\textasciitilde{}/notes/todo.org, ./README.org, or a file: link’s) in its buffer.\\
none & \texttt{C-c C-z} & \texttt{org-\codebreak{}add-\codebreak{}note} & Add a note to the entry point is under (C-c C-z): the note is asked for, and logged where Org logs notes, at the top of its :LOGBOOK: drawer with \textbf{ORG-LOG-INTO-DRAWER}, else after its planning line and properties, in the form - Note taken on \textit{[2026-10-01 Thu 18:00]} \textbackslash{}\textbackslash{} and the note under it.\\
none & \texttt{a}, \texttt{C-c a}, \texttt{C-x a a} & \texttt{org-\codebreak{}agenda} & Show the agenda, asking which view: a the agenda of the week (or \textbf{ORG-AGENDA-SPAN}), t the TODO list, m the headlines a match selects (work+urgent-home, TODO=”NEXT”), or a view of your own (\textbf{ORG-AGENDA-CUSTOM-COMMANDS}, u the tasks with no tag), or c the entries whose dates look wrong (M-x org-agenda-check-dates), read from \textbf{ORG-AGENDA-FILES}.\\
none & \texttt{n}, \texttt{C-c n}, \texttt{C-c C-z} & \texttt{org-\codebreak{}agenda-\codebreak{}add-\codebreak{}note} & Add a note to the entry on the agenda’s line, as C-c C-z adds one.\\
none & \texttt{C-c C-x C} & \texttt{org-\codebreak{}agenda-\codebreak{}cancel-\codebreak{}task} & Cancel the task on the agenda’s line, as C-c C-x C does in its file: its keyword made CANCELLED and noted with CLOSED:.\\
none & \texttt{M-x org-\codebreak{}agenda-\codebreak{}check-\codebreak{}dates} & \texttt{org-\codebreak{}agenda-\codebreak{}check-\codebreak{}dates} & Show the entries of \textbf{ORG-AGENDA-FILES} whose dates look wrong, each with what is wrong (ORG-ENTRY-DATE-PROBLEM).\\
none & \texttt{C}, \texttt{C-c C} & \texttt{org-\codebreak{}agenda-\codebreak{}clock-\codebreak{}cancel} & Cancel the clock, as C-c C-x C-q does, from the agenda: the clock’s file saved at once when it is one of \textbf{ORG-AGENDA-AUTOSAVE-FILES}.\\
none & \texttt{I}, \texttt{C-c I} & \texttt{org-\codebreak{}agenda-\codebreak{}clock-in} & Clock in on the entry on the agenda’s line, as M-x org-clock-in does: a clock running is clocked out first.\\
none & \texttt{O}, \texttt{C-c O} & \texttt{org-\codebreak{}agenda-\codebreak{}clock-\codebreak{}out} & Clock out, as C-c C-x C-o does, from the agenda: the clock’s file saved at once when it is one of \textbf{ORG-AGENDA-AUTOSAVE-FILES}.\\
none & \texttt{D}, \texttt{C-c D} & \texttt{org-\codebreak{}agenda-\codebreak{}day-\codebreak{}view} & Show the agenda of one day: the day at point in the agenda, else today.\\
none & \texttt{d}, \texttt{C-c d}, \texttt{C-c C-d} & \texttt{org-\codebreak{}agenda-\codebreak{}deadline} & Give the entry on the agenda’s line a deadline, the date asked for as C-c C-d asks it; with C-u, take it away.\\
none & \texttt{h}, \texttt{C-c h} & \texttt{org-\codebreak{}agenda-\codebreak{}earlier} & Show the agenda’s previous span: the previous week in a week’s agenda, the previous day in a day’s.\\
none & \texttt{TAB} & \texttt{org-\codebreak{}agenda-\codebreak{}goto} & Show the entry on the agenda’s line in the other window and go there.\\
none & \texttt{c}, \texttt{C-c c} & \texttt{org-\codebreak{}agenda-\codebreak{}goto-\codebreak{}calendar} & Show the calendar on the day of the agenda line at point.\\
none & \texttt{g h} & \texttt{org-\codebreak{}agenda-\codebreak{}goto-\codebreak{}today} & Show the agenda’s span holding today, point on today’s date.\\
none & \texttt{l}, \texttt{C-c l} & \texttt{org-\codebreak{}agenda-\codebreak{}later} & Show the agenda’s next span: the next week in a week’s agenda, the next day in a day’s.\\
none & \texttt{M-x org-\codebreak{}agenda-\codebreak{}list} & \texttt{org-\codebreak{}agenda-\codebreak{}list} & Show the agenda of this week (or \textbf{ORG-AGENDA-SPAN}): each day with its scheduled tasks, deadlines and timestamps.\\
none & \texttt{L}, \texttt{C-c L} & \texttt{org-\codebreak{}agenda-\codebreak{}log-\codebreak{}mode} & List what was closed and clocked on each day too, or no longer; \textbf{ORG-AGENDA-LOG-MODE} says how the agenda starts.\\
none & \texttt{m}, \texttt{C-c m}, \texttt{C-x a m} & \texttt{org-\codebreak{}agenda-\codebreak{}month-\codebreak{}view} & Show the agenda of a month: the month of the day at point in the agenda, else this month.\\
none & \texttt{g l} & \texttt{org-\codebreak{}agenda-\codebreak{}previous-\codebreak{}view} & Go back to the view or span the agenda showed before this one, point on the line it was on, as cd - goes back to the directory before: pressed again, return.\\
none & \texttt{C-c C-w} & \texttt{org-\codebreak{}agenda-\codebreak{}refile} & Move the entry on the agenda’s line under another heading, as C-c C-w moves one.\\
none & \texttt{s} & \texttt{org-\codebreak{}agenda-\codebreak{}save-\codebreak{}all-\codebreak{}buffers} & Save the buffers of the agenda files (\textbf{org-agenda-files}, and those added to them) with changes not saved, those the agenda changed among them; other Org buffers are left as they are (M-x org-save-all-org-buffers saves them all).\\
none & \texttt{C-c C-s} & \texttt{org-\codebreak{}agenda-\codebreak{}schedule} & Schedule the entry on the agenda’s line, the date asked for as C-c C-s asks it; with C-u, take its scheduling away.\\
none & \texttt{t}, \texttt{C-c C-q} & \texttt{org-\codebreak{}agenda-\codebreak{}set-\codebreak{}tags} & Set the tags of the entry on the agenda’s line, asked for as C-c C-q asks them.\\
none & \texttt{RET} & \texttt{org-\codebreak{}agenda-\codebreak{}switch-to} & Open the entry on the agenda’s line in this window, point on its headline.\\
none & \texttt{T}, \texttt{C-c T}, \texttt{C-c C-t} & \texttt{org-\codebreak{}agenda-\codebreak{}todo} & Change the TODO keyword of the entry on the agenda’s line to its next, as C-c C-t does in its file: a task done is noted with CLOSED:, one that repeats moves on to its next date.\\
none & \texttt{u}, \texttt{C-/}, \texttt{C-\_} & \texttt{org-\codebreak{}agenda-\codebreak{}undo} & Undo the last change made from this agenda, in the file it was made in, and draw the agenda again; when the file was changed there since, nothing is undone and you are told to undo in the file first.\\
none & \texttt{w}, \texttt{C-c w}, \texttt{C-x a w} & \texttt{org-\codebreak{}agenda-\codebreak{}week-\codebreak{}view} & Show the agenda of a week: the week of the day at point in the agenda, else this week.\\
none & \texttt{Y}, \texttt{C-c Y}, \texttt{C-x a y} & \texttt{org-\codebreak{}agenda-\codebreak{}year-\codebreak{}view} & Show the agenda of a year: the year of the day at point in the agenda, else this year.\\
none & \texttt{M-x org-\codebreak{}align-\codebreak{}all-\codebreak{}tags} & \texttt{org-\codebreak{}align-\codebreak{}all-\codebreak{}tags} & Put the tags of every headline of the buffer where \textbf{ORG-TAGS-COLUMN} says: by default one space after each one’s text, in one step to undo.\\
none & \texttt{M-x org-\codebreak{}align-\codebreak{}tags} & \texttt{org-\codebreak{}align-\codebreak{}tags} & Put the tags of the headline point is under where \textbf{ORG-TAGS-COLUMN} says: by default one space after its text.\\
none & \texttt{M-x org-\codebreak{}calendar-\codebreak{}backward-\codebreak{}day} & \texttt{org-\codebreak{}calendar-\codebreak{}backward-\codebreak{}day} & In Org’s date prompt: the calendar’s day a day back, written at the prompt.\\
none & \texttt{M-x org-\codebreak{}calendar-\codebreak{}backward-\codebreak{}month} & \texttt{org-\codebreak{}calendar-\codebreak{}backward-\codebreak{}month} & In Org’s date prompt: the calendar’s day a month back, written at the prompt as the same day of the month, or the month’s last when it has fewer, as the calendar moves; -1m typed counts on past the month’s end instead.\\
none & \texttt{M-x org-\codebreak{}calendar-\codebreak{}backward-\codebreak{}week} & \texttt{org-\codebreak{}calendar-\codebreak{}backward-\codebreak{}week} & In Org’s date prompt: the calendar’s day a week back, written at the prompt.\\
none & \texttt{M-x org-\codebreak{}calendar-\codebreak{}backward-\codebreak{}year} & \texttt{org-\codebreak{}calendar-\codebreak{}backward-\codebreak{}year} & In Org’s date prompt: the calendar’s day a year back, written at the prompt, so 29 February a year back is 28 February, as the calendar moves; -1y typed is 1 March.\\
none & \texttt{M-x org-\codebreak{}calendar-\codebreak{}forward-\codebreak{}day} & \texttt{org-\codebreak{}calendar-\codebreak{}forward-\codebreak{}day} & In Org’s date prompt: the calendar’s day a day on, written at the prompt.\\
none & \texttt{M-x org-\codebreak{}calendar-\codebreak{}forward-\codebreak{}month} & \texttt{org-\codebreak{}calendar-\codebreak{}forward-\codebreak{}month} & In Org’s date prompt: the calendar’s day a month on, written at the prompt as the same day of the month, or the month’s last when it has fewer (31 January a month on is 28 February), as the calendar moves; +1m typed counts on past the month’s end instead (3 March).\\
\bottomrule
\end{xltabular}

{\small
\begin{xltabular}{\textwidth}{@{}l >{\raggedright\arraybackslash\hsize=0.694\hsize}X l >{\raggedright\arraybackslash\hsize=1.306\hsize}X@{}}
\toprule
\textbf{Modal} & \textbf{Non-modal} & \textbf{Command} & \textbf{What it does}\\
\midrule
\endhead
none & \texttt{M-x org-\codebreak{}calendar-\codebreak{}forward-\codebreak{}week} & \texttt{org-\codebreak{}calendar-\codebreak{}forward-\codebreak{}week} & In Org’s date prompt: the calendar’s day a week on, written at the prompt.\\
none & \texttt{M-x org-\codebreak{}calendar-\codebreak{}forward-\codebreak{}year} & \texttt{org-\codebreak{}calendar-\codebreak{}forward-\codebreak{}year} & In Org’s date prompt: the calendar’s day a year on, written at the prompt, so 29 February a year on is 28 February, as the calendar moves; +1y typed is 1 March.\\
none & \texttt{a}, \texttt{C-c a} & \texttt{org-\codebreak{}calendar-\codebreak{}goto-\codebreak{}agenda} & Show the agenda of the week holding the calendar’s day at point, point on that day.\\
none & \texttt{M-x org-\codebreak{}calendar-\codebreak{}goto-\codebreak{}today} & \texttt{org-\codebreak{}calendar-\codebreak{}goto-\codebreak{}today} & In Org’s date prompt: the calendar on today, written at the prompt.\\
none & \texttt{C-c C-x C} & \texttt{org-\codebreak{}cancel-\codebreak{}task} & Cancel the task of the headline point is in: its keyword made CANCELLED, a done keyword that M-x org-todo and the other keys that go round the keywords pass over, noted with CLOSED: and when as a task done is.\\
none & \texttt{C-x a c} & \texttt{org-\codebreak{}capture} & Take a note, a task or a journal entry down from wherever you are: press the key of one of the capture templates (Tuatara’s and yours, in \textbf{CAPTURE-TEMPLATE-DIRECTORIES}), answer what it asks, write in the \textbf{Capture} buffer and file it with C-c C-c (or throw it away with C-c C-k); you are back where you were.\\
none & \texttt{C-c C-c} & \texttt{org-\codebreak{}capture-\codebreak{}finalize} & File the text of this \textbf{Capture} buffer where its template says, and go back to where you were; with its template’s clock, clock it.\\
none & \texttt{C-c C-k} & \texttt{org-\codebreak{}capture-\codebreak{}kill} & Throw away what this \textbf{Capture} buffer holds, and go back to where you were.\\
none & \texttt{i}, \texttt{C-c i} & \texttt{org-\codebreak{}capture-\codebreak{}on-\codebreak{}day} & Capture, as C-x a c does, for the day at point in the calendar or on the agenda (the day of its line, or of the date above it), so that the agenda lists the entry that day: the template’s \%t and \%T are that day’s date, and a template with neither gets that date on the line after its heading.\\
none & \texttt{C-c C-w} & \texttt{org-\codebreak{}capture-\codebreak{}refile} & File the text of this \textbf{Capture} buffer under a heading chosen as C-c C-w chooses one, instead of where its template says.\\
none & \texttt{M-x org-\codebreak{}capture-\codebreak{}template-\codebreak{}copy} & \texttt{org-\codebreak{}capture-\codebreak{}template-\codebreak{}copy} & Copy one of Tuatara’s own capture templates, chosen by its name, into your capture templates folder (the first of \textbf{CAPTURE-TEMPLATE-DIRECTORIES}) and visit it: your copy, changed as you like (where it files, its text), takes the place of Tuatara’s on the same key.\\
none & \texttt{C-c C-x C-q} & \texttt{org-\codebreak{}clock-\codebreak{}cancel} & Stop the clock running and take its line away (C-c C-x C-q): its :LOGBOOK: drawer too, when nothing else is in it; saved when that is its buffer’s only change, as C-x a o does.\\
none & \texttt{C-x a j}, \texttt{C-c C-x C-j} & \texttt{org-\codebreak{}clock-\codebreak{}goto} & Go to the task the clock runs on (C-c C-x C-j, from any buffer C-x a j): its file opened again when its buffer was killed; with none running, one left open in the agenda files, as after a restart, runs again (ORG-CLOCK-RUNNING).\\
none & \texttt{C-c C-x C-i} & \texttt{org-\codebreak{}clock-in} & Clock in on the task point is under (C-c C-x C-i): CLOCK: and the time go into its :LOGBOOK: drawer, made when it has none, after its planning line and properties; a clock running is clocked out first.\\
none & \texttt{C-x a o}, \texttt{C-c C-x C-o} & \texttt{org-\codebreak{}clock-\codebreak{}out} & Clock out (C-c C-x C-o, from any buffer C-x a o): the clock’s line finished with the time it ended and how long it ran; with none running, of the one left open in the agenda files (ORG-CLOCK-RUNNING).\\
none & \texttt{M-x org-\codebreak{}clock-\codebreak{}report} & \texttt{org-\codebreak{}clock-\codebreak{}report} & Make the clock table point is in again, or put a new one in: its scope the subtree point is in, or the file before the first headline, its depth two levels.\\
none & \texttt{C-c C-c} & \texttt{org-\codebreak{}ctrl-\codebreak{}c-\codebreak{}ctrl-c} & Do what Org’s C-c C-c does where point is: on a dynamic block’s \#+BEGIN: line, make the block again (M-x org-dblock-update).\\
none & \texttt{C-c -} & \texttt{org-\codebreak{}ctrl-\codebreak{}c-\codebreak{}minus} & Do what Org’s C-c - does where point is: in a table, a rule line below the row (org-table-insert-hline); in a list, the next kind of bullet (org-cycle-list-bullet).\\
none & \texttt{TAB} & \texttt{org-\codebreak{}cycle} & On a headline, fold its subtree to its headline alone, then show its children’s headlines, then all of it, and fold it again, a step each time.\\
none & \texttt{M-x org-\codebreak{}cycle-\codebreak{}list-\codebreak{}bullet} & \texttt{org-\codebreak{}cycle-\codebreak{}list-\codebreak{}bullet} & Give the list point is in its next kind of bullet (C-c -): - to + to * to 1.\\
none & \texttt{M-x org-\codebreak{}date-\codebreak{}choose-\codebreak{}mark} & \texttt{org-\codebreak{}date-\codebreak{}choose-\codebreak{}mark} & In Org’s date question, answered in the calendar: mark the calendar’s day, beside those marked where the question takes several days, else in place of the one marked.\\
none & \texttt{M-x org-\codebreak{}date-\codebreak{}choose-\codebreak{}unmark} & \texttt{org-\codebreak{}date-\codebreak{}choose-\codebreak{}unmark} & In Org’s date question, answered in the calendar: take the mark off the calendar’s day.\\
none & \texttt{M-x org-\codebreak{}date-\codebreak{}choose-\codebreak{}write} & \texttt{org-\codebreak{}date-\codebreak{}choose-\codebreak{}write} & In Org’s date question, answered in the calendar: write the days chosen at the prompt (the calendar’s day when none is marked), to add a time or a repeat to; RET there takes it.\\
none & \texttt{M-x org-\codebreak{}date-\codebreak{}repeat-\codebreak{}menu} & \texttt{org-\codebreak{}date-\codebreak{}repeat-\codebreak{}menu} & At the date prompt of M-x org-schedule, M-x org-deadline and M-x org-time-stamp: leave it for the menu of a repeat (REPEAT-MENU), started from what was typed, or the calendar’s day.\\
none & \texttt{C-c C-x C-u} & \texttt{org-\codebreak{}dblock-\codebreak{}update} & Make the dynamic block point is in again (C-c C-x C-u): the text between \#+BEGIN: and \#+END: written anew by the writer \#+BEGIN: names, in the background, the editor free meanwhile.\\
none & \texttt{C-c C-d} & \texttt{org-\codebreak{}deadline} & Give the task point is under a deadline (C-c C-d): the date asked for as M-x org-schedule asks it, DEADLINE: on its planning line, a repeater and warning it had kept.\\
none & \texttt{M-x org-\codebreak{}demote-\codebreak{}subtree} & \texttt{org-\codebreak{}demote-\codebreak{}subtree} & Demote the headline point is in a level, and every headline under it.\\
none & \texttt{M-x org-\codebreak{}do-\codebreak{}demote} & \texttt{org-\codebreak{}do-\codebreak{}demote} & Demote the headline point is in a level, its subtree’s other headlines as they are.\\
none & \texttt{M-x org-\codebreak{}do-\codebreak{}promote} & \texttt{org-\codebreak{}do-\codebreak{}promote} & Promote the headline point is in a level, its subtree’s other headlines as they are.\\
none & \texttt{C-c C-x x} & \texttt{org-\codebreak{}dynamic-\codebreak{}block-\codebreak{}insert-\codebreak{}dblock} & Put a new dynamic block in (C-c C-x x): its kind asked for among those that can be (a clock table, links to files with formatted names\ldots{}), then what that kind asks.\\
none & \texttt{M-x org-\codebreak{}export-\codebreak{}docx} & \texttt{org-\codebreak{}export-\codebreak{}docx} & Write the Org note of this buffer as a Word document (.docx) beside it, or where \#+EXPORT\_FILE\_NAME says (its .docx): read as the enhanced PDF export reads it, its :EXPORT: drawers’ settings and all, what redact and protect hide a black bar, never its text, in the file; in the background.\\
none & \texttt{M-x org-\codebreak{}export-\codebreak{}docx-\codebreak{}xdg} & \texttt{org-\codebreak{}export-\codebreak{}docx-\codebreak{}xdg} & Write the Org note of this buffer as a Word document, as M-x org-export-docx does, and open it in the desktop’s program for Word documents.\\
none & \texttt{M-x org-\codebreak{}export-\codebreak{}html} & \texttt{org-\codebreak{}export-\codebreak{}html} & Write the Org note of this buffer as a web page (.html) beside it, or where \#+EXPORT\_FILE\_NAME says (its .html), one file with its stylesheet inside, in the Ao theme’s colours, light or dark as the reader’s system prefers: read as the enhanced PDF export reads it, what redact and protect hide a black bar, never its text.\\
none & \texttt{M-x org-\codebreak{}export-\codebreak{}latex} & \texttt{org-\codebreak{}export-\codebreak{}latex} & Write the Org file of the current buffer as \LaTeX{} beside it (or where its \#+export\_file\_name says), the plain export with no template or layout of Tuatara’s own, in the background, without making a PDF: the \LaTeX{} file M-x org-export-latex-pdf makes its PDF from.\\
none & \texttt{M-x org-\codebreak{}export-\codebreak{}latex-\codebreak{}enhanced} & \texttt{org-\codebreak{}export-\codebreak{}latex-\codebreak{}enhanced} & Write the Org file of the current buffer as \LaTeX{} as M-x org-export-latex does, with Tuatara’s own improvements (its \#+template, tables that fit, redaction): the \LaTeX{} file M-x org-export-latex-pdf-enhanced makes its PDF from.\\
none & \texttt{M-x org-\codebreak{}export-\codebreak{}latex-\codebreak{}pdf} & \texttt{org-\codebreak{}export-\codebreak{}latex-\codebreak{}pdf} & Make the Org file of the current buffer into a PDF beside it, in the background: the plain export, the note set as standard \LaTeX{} sets it, with no template or layout of Tuatara’s own.\\
none & \texttt{M-x org-\codebreak{}export-\codebreak{}latex-\codebreak{}pdf-\codebreak{}buffer} & \texttt{org-\codebreak{}export-\codebreak{}latex-\codebreak{}pdf-\codebreak{}buffer} & Make the Org file of the current buffer into a PDF beside it, as M-x org-export-latex-pdf does, and show the PDF: down the right side of the screen when it is wide enough (\textbf{split-width-threshold} columns), else alone on it.\\
none & \texttt{M-x org-\codebreak{}export-\codebreak{}latex-\codebreak{}pdf-\codebreak{}enhanced} & \texttt{org-\codebreak{}export-\codebreak{}latex-\codebreak{}pdf-\codebreak{}enhanced} & Make the Org file of the current buffer into a PDF beside it, as M-x org-export-latex-pdf does, with Tuatara’s own improvements: its \#+template, tables that fit and redaction (M-x org-export-latex-enhanced).\\
none & \texttt{M-x org-\codebreak{}export-\codebreak{}latex-\codebreak{}pdf-\codebreak{}enhanced-\codebreak{}buffer} & \texttt{org-\codebreak{}export-\codebreak{}latex-\codebreak{}pdf-\codebreak{}enhanced-\codebreak{}buffer} & Make the Org file of the current buffer into a PDF as M-x org-export-latex-pdf-enhanced does, and show it as M-x org-export-latex-pdf-buffer shows its PDF.\\
none & \texttt{M-x org-\codebreak{}export-\codebreak{}latex-\codebreak{}pdf-\codebreak{}enhanced-\codebreak{}xdg} & \texttt{org-\codebreak{}export-\codebreak{}latex-\codebreak{}pdf-\codebreak{}enhanced-\codebreak{}xdg} & Make the Org file of the current buffer into a PDF as M-x org-export-latex-pdf-enhanced does, and open it in the desktop’s program as M-x org-export-latex-pdf-xdg opens its PDF.\\
\bottomrule
\end{xltabular}
}

\begin{xltabular}{\textwidth}{@{}l >{\raggedright\arraybackslash\hsize=0.791\hsize}X l >{\raggedright\arraybackslash\hsize=1.209\hsize}X@{}}
\toprule
\textbf{Modal} & \textbf{Non-modal} & \textbf{Command} & \textbf{What it does}\\
\midrule
\endhead
none & \texttt{M-x org-\codebreak{}export-\codebreak{}latex-\codebreak{}pdf-\codebreak{}xdg} & \texttt{org-\codebreak{}export-\codebreak{}latex-\codebreak{}pdf-\codebreak{}xdg} & Make the Org file of the current buffer into a PDF beside it, as M-x org-export-latex-pdf does, and open the PDF in the desktop’s program for PDFs (xdg-open), or the one \textbf{open-externally-programs} gives “application/pdf”, asking first unless \textbf{open-externally-ask} is NIL.\\
none & \texttt{S-TAB} & \texttt{org-\codebreak{}global-\codebreak{}cycle} & Cycle the whole buffer: its overview (the headlines of the first level), then its contents (every headline), then everything, and the overview again.\\
none & \texttt{M-x org-\codebreak{}indent-\codebreak{}mode} & \texttt{org-\codebreak{}indent-\codebreak{}mode} & Show the current Org buffer indented, or as it is again: each line indented by the level of the headline it is under, a headline’s stars but the last hidden.\\
none & \texttt{M-x org-\codebreak{}insert-\codebreak{}export-\codebreak{}drawer} & \texttt{org-\codebreak{}insert-\codebreak{}export-\codebreak{}drawer} & Write an :EXPORT: drawer for the table, list, block or picture point is on, or else for the heading it is under, holding the settings that fit it, each auto to change (M-x describe-export-settings says what each takes).\\
none & \texttt{C-c C-l} & \texttt{org-\codebreak{}insert-\codebreak{}link} & Insert an Org link (C-c C-l): it asks for the link, offering those org-store-link (C-x a l) kept and the kinds of link (file asks for a file, denote for one of your notes), then for its label, one offered selected (the stored link’s, the note’s title) to take or type over; the text selected becomes the label instead.\\
none & \texttt{M-S-\codebreak{}RET} & \texttt{org-\codebreak{}insert-\codebreak{}todo-\codebreak{}heading} & A new TODO headline of this level after this subtree.\\
none & \texttt{C-c M} & \texttt{org-\codebreak{}macro-\codebreak{}unwrap} & In an Org note, take away the macro call point is in, its text left as it was: \texorpdfstring{\redactedbox{nnnnnnnnnn}}{[redacted]} becomes John Otago.\\
none & \texttt{C-c m} & \texttt{org-\codebreak{}macro-\codebreak{}wrap} & In an Org note, wrap the selection, or the word at point, in a macro chosen from a list: R, \texorpdfstring{\redactedbox{nnnnnnnnnn}}{[redacted]}, hidden in the enhanced export’s PDF as a black box; P, hidden there only with \#+options: protect:t; or one the note declares with a \#+macro: line, what it stands for shown beside it.\\
none & \texttt{M-RET} & \texttt{org-\codebreak{}meta-\codebreak{}return} & A new item or headline after what point is in and all under it: in a list, an item of the same kind after the one point is in and all under it (its bullet, the next number, [ ] when it has a checkbox).\\
none & \texttt{M-j}, \texttt{M-<down>} & \texttt{org-\codebreak{}metadown} & In an Org table, move the row down; else move the subtree down (M-j).\\
none & \texttt{M-h}, \texttt{M-<left>} & \texttt{org-\codebreak{}metaleft} & In an Org table, move the column left; else promote the headline (M-h).\\
none & \texttt{M-l}, \texttt{M-<right>} & \texttt{org-\codebreak{}metaright} & In an Org table, move the column right; else demote the headline (M-l).\\
none & \texttt{M-k}, \texttt{M-<up>} & \texttt{org-\codebreak{}metaup} & In an Org table, move the row up; else move the subtree up (M-k).\\
none & \texttt{M-x org-\codebreak{}move-\codebreak{}subtree-\codebreak{}down} & \texttt{org-\codebreak{}move-\codebreak{}subtree-\codebreak{}down} & Move the subtree point is in down, past its next sibling.\\
none & \texttt{M-x org-\codebreak{}move-\codebreak{}subtree-up} & \texttt{org-\codebreak{}move-\codebreak{}subtree-up} & Move the subtree point is in up, before its sibling above.\\
none & \texttt{C-c C-o} & \texttt{org-\codebreak{}open-\codebreak{}at-\codebreak{}point} & Follow the link point is on (C-c C-o): to a headline, a custom ID or a target here, a file, an ID, a note by its identifier, or a web or mail address in the desktop’s program, after asking.\\
none & \texttt{M-x org-\codebreak{}promote-\codebreak{}subtree} & \texttt{org-\codebreak{}promote-\codebreak{}subtree} & Promote the headline point is in a level, and every headline under it.\\
none & \texttt{C-c C-w} & \texttt{org-\codebreak{}refile} & Move the heading point is under, with all under it, to another heading: chosen by its path (agenda.org/Commitments and Jobs) among the headings \textbf{ORG-REFILE-TARGETS} offers, put at the end of what is under it, a level below; a clock running in it goes on there.\\
none & \texttt{RET} & \texttt{org-\codebreak{}return} & RET in Org: on a link, follow it (\textbf{org-return-follows-link}); in a table, go to the row below, aligned; else a new line, as RET does in text.\\
none & \texttt{M-x org-\codebreak{}save-\codebreak{}all-\codebreak{}org-\codebreak{}buffers} & \texttt{org-\codebreak{}save-\codebreak{}all-\codebreak{}org-\codebreak{}buffers} & Save every Org buffer with changes not saved: those the agenda changed among them, and any other.\\
none & \texttt{C-c C-s} & \texttt{org-\codebreak{}schedule} & Schedule the task point is under (C-c C-s): the date is asked for (fri, +3d, 10-15, sep 30, 2026-12-24 14:00), what it reads shown as it is typed, and SCHEDULED: goes on its planning line, a repeater and warning it had kept.\\
none & \texttt{M-x org-\codebreak{}scratch} & \texttt{org-\codebreak{}scratch} & Show \textbf{org-scratch}, the Org buffer for notes that are not saved; when it was killed, a new one.\\
none & \texttt{C-c C-q} & \texttt{org-\codebreak{}set-\codebreak{}tags-\codebreak{}command} & Set the tags of the headline point is under (C-c C-q): asked for with its tags offered, work urgent or :work:urgent:, nothing for none; they go at the end of its line, where \textbf{ORG-TAGS-COLUMN} says, by default one space after its text.\\
none & \texttt{S-<left>} & \texttt{org-\codebreak{}shiftleft} & On a headline, its TODO keyword before, none before the first.\\
none & \texttt{M-S-<down>} & \texttt{org-\codebreak{}shiftmetadown} & In an Org table, put an empty row above this one (M-S-<down>).\\
none & \texttt{M-H}, \texttt{M-S-<left>} & \texttt{org-\codebreak{}shiftmetaleft} & In an Org table, delete the column; else promote the subtree (M-H).\\
none & \texttt{M-L}, \texttt{M-S-<right>} & \texttt{org-\codebreak{}shiftmetaright} & In an Org table, put an empty column before this one; else demote the subtree (M-L).\\
none & \texttt{M-S-<up>} & \texttt{org-\codebreak{}shiftmetaup} & In an Org table, delete the row point is in (M-S-<up>).\\
none & \texttt{S-<right>} & \texttt{org-\codebreak{}shiftright} & On a headline, its next TODO keyword, as M-x org-todo does.\\
none & \texttt{C-x a l} & \texttt{org-\codebreak{}store-\codebreak{}link} & Keep a link to where point is, for org-insert-link (M-x org-insert-link) to offer: under an Org headline, to it (by its CUSTOM\_ID, else its title); in one of your notes, to the note (denote:ID); in another file, to the line point is on.\\
none & \texttt{M-x org-\codebreak{}table-\codebreak{}align} & \texttt{org-\codebreak{}table-\codebreak{}align} & Align the Org table point is in: each column as wide as its widest field, numbers to the right, rules as wide as the table.\\
none & \texttt{C-c =} & \texttt{org-\codebreak{}table-\codebreak{}eval-\codebreak{}formula} & Ask for the formula of the column point is in, the one it has offered, and work it out (C-c =); with C-u, the formula of the field point is in.\\
none & \texttt{M-x org-\codebreak{}table-\codebreak{}insert-\codebreak{}hline} & \texttt{org-\codebreak{}table-\codebreak{}insert-\codebreak{}hline} & Put a rule below the row point is in, in an Org table (C-c -).\\
none & \texttt{C-c *} & \texttt{org-\codebreak{}table-\codebreak{}recalculate} & Work out the formulas of the Org table point is in, its \#+TBLFM: line (C-c *): column formulas (\$3=\$1+\$2) below the first rule (in every row when it has none), field formulas (@4\$2=vsum(@2..@3)) and range formulas (@2\$2..@3\$3=\$1*10), worked out with the table arithmetic.\\
none & \texttt{M-x org-\codebreak{}tags-\codebreak{}view} & \texttt{org-\codebreak{}tags-\codebreak{}view} & Show the headlines of \textbf{ORG-AGENDA-FILES} an Org match selects, read in the minibuffer: work, +work-home, work\(\vert{}\)proj, TODO=”NEXT”, LEVEL>1, -\{.*\} for those with no tag.\\
none & \texttt{C-c .} & \texttt{org-\codebreak{}time-\codebreak{}stamp} & Put in a timestamp at point (C-c .), the date asked for as M-x org-schedule asks it: \textit{<2026-10-13 Tue>}, or with a time \textit{<2026-10-13 Tue 14:00>}.\\
none & \texttt{C-c !} & \texttt{org-\codebreak{}time-\codebreak{}stamp-\codebreak{}inactive} & Put in an inactive timestamp at point (C-c !): as M-x org-time-stamp, written \textit{[2026-10-13 Tue]}, which the agenda leaves out.\\
none & \texttt{C-c C-t} & \texttt{org-\codebreak{}todo} & Change the TODO keyword of the headline point is in to the next of the document’s (\#+TODO:, else TODO, WAIT and DONE), none after the last, those of \textbf{org-todo-keywords-not-cycled} passed over (CANCELLED, which M-x org-cancel-task sets); a task done is noted with CLOSED: and when.\\
none & \texttt{M-x org-\codebreak{}todo-\codebreak{}list} & \texttt{org-\codebreak{}todo-\codebreak{}list} & Show the TODO list: every task not done in \textbf{ORG-AGENDA-FILES}.\\
none & \texttt{C-c C-x L} & \texttt{org-\codebreak{}toggle-\codebreak{}link-\codebreak{}display} & Show this Org buffer’s links as written, \hyperref[org376]{label}, and its entities as written, \texttt{\textbackslash{}dollar\{\}}, or the links as their labels and the entities as their characters again (C-c C-x L).\\
\bottomrule
\end{xltabular}

\begin{xltabular}{\textwidth}{@{}l >{\raggedright\arraybackslash\hsize=0.722\hsize}X l >{\raggedright\arraybackslash\hsize=1.278\hsize}X@{}}
\toprule
\textbf{Modal} & \textbf{Non-modal} & \textbf{Command} & \textbf{What it does}\\
\midrule
\endhead
none & \texttt{M-x org-\codebreak{}update-\codebreak{}all-\codebreak{}dblocks} & \texttt{org-\codebreak{}update-\codebreak{}all-\codebreak{}dblocks} & Make every dynamic block of the current Org buffer again: each written from the text as it is now, in the background, the editor free meanwhile, and put in together, in one step to undo.\\
none & \texttt{C-x 5 o} & \texttt{other-\codebreak{}frame} & Go to the next frame and bring its window before the others, or in a terminal show it in place of the frame shown; with a count N, the Nth frame on, back when N is negative.\\
none & \texttt{C-x o} & \texttt{other-\codebreak{}window} & Select the next window, from the top left, the first again after the last.\\
none & \texttt{M-x package-\codebreak{}delete} & \texttt{package-\codebreak{}delete} & Unload a package (M-x package-unload), remove it from the packages loaded (M-x package-remove), and delete what Tuatara keeps of it: the versions of it fetched and their compiled files, or for a folder of your disk its compiled files alone, the folder of yours left as it is.\\
none & \texttt{L}, \texttt{C-c L} & \texttt{package-\codebreak{}load} & Load a package the init file loads from a source that is not loaded now, by its load-package there, as at the start: one unloaded (M-x package-unload), removed (M-x package-remove), whose loading failed, or one deleted, fetched again.\\
none & \texttt{M-x package-\codebreak{}manager} & \texttt{package-\codebreak{}manager} & Manage the packages the init file loads from a source (load-package with :source) in a list, as the buffer list (C-x C-b) manages buffers.\\
none & \texttt{R}, \texttt{C-c R} & \texttt{package-\codebreak{}remove} & Unload a package (M-x package-unload) and take it out of the packages loaded: it is not loaded from now on, nor at the next start, though the init file names it; the init file is left as it is, and Tuatara says in \textbf{Messages} once a session that it passes the package over.\\
none & \texttt{M-x package-\codebreak{}remove-\codebreak{}unused} & \texttt{package-\codebreak{}remove-\codebreak{}unused} & Delete the package sources Tuatara keeps that nothing in this session names any more: a source no package-source names and no load-package fetches from, or a package of a source no load-package names; one left out by :disabled or :if counts as named.\\
none & \texttt{M-x package-\codebreak{}scan} & \texttt{package-\codebreak{}scan} & Look in the init file for the packages it loads from a source (each load-package with :source), and load each that is not loaded or on its way, as at the start: fetched in the background, asked about before it is compiled, loaded.\\
none & \texttt{M-x package-\codebreak{}trust-\codebreak{}again} & \texttt{package-\codebreak{}trust-\codebreak{}again} & Trust the source of a package asked for anew, as the first time: who signed the versions taken from it is forgotten, and the package is fetched again (M-x package-update); the version found is asked about with who signed it (\textbf{PACKAGE-ASK-BEFORE-LOADING}), and once it is let load, that signer is the one remembered.\\
none & \texttt{K}, \texttt{C-c K} & \texttt{package-\codebreak{}unload} & Take a package out of the running editor: its commands and keys, with their labels, its modes, the places it gave Tuatara’s commands, its ways of showing text, its functions on Tuatara’s lists of hooks and of work to do now and then, its settings and its work in the background, which is stopped.\\
none & \texttt{M-x package-\codebreak{}update} & \texttt{package-\codebreak{}update} & Fetch again each package load-package fetched (not the folders of your disk), in the background.\\
none & \texttt{f}, \texttt{C-c f} & \texttt{package-\codebreak{}update-\codebreak{}marked} & Fetch again the marked packages (>), or with none marked the one on the current line, as M-x package-update does them all.\\
none & \texttt{M-x pattern-\codebreak{}finish} & \texttt{pattern-\codebreak{}finish} & In a question that takes a pattern: take what is typed, or when it cannot be used, say why in the prompt and wait for it to be mended (M-x minibuffer-cancel gives up).\\
none & \texttt{s}, \texttt{RET} & \texttt{pdf-\codebreak{}attachment-\codebreak{}save} & Save the embedded file on the current line where you say, a new file, or one there after asking: taken out of the PDF in the sandbox, as its program reads it, and written as it is, never opened.\\
none & \texttt{a}, \texttt{C-c a} & \texttt{pdf-\codebreak{}view-\codebreak{}attachments} & List the files embedded in the PDF, in a buffer beside it, as its program reads them in the sandbox; there, s saves one.\\
none & \texttt{i}, \texttt{C-c i} & \texttt{pdf-\codebreak{}view-\codebreak{}facts} & Show what the PDF says of itself, as help: title, subject, author, the programs that made it (creator and producer), when it was made and changed, its version, encryption and what it allows, JavaScript, forms, pages and their size, as its program reads them in the sandbox; the file’s own name and size first.\\
none & \texttt{g h} & \texttt{pdf-\codebreak{}view-\codebreak{}first-\codebreak{}page} & Turn to the first page of the PDF.\\
none & \texttt{p}, \texttt{C-c p} & \texttt{pdf-\codebreak{}view-\codebreak{}fit-\codebreak{}page} & Draw the PDF’s pages as large as the window shows them whole.\\
none & \texttt{w}, \texttt{C-c w} & \texttt{pdf-\codebreak{}view-\codebreak{}fit-\codebreak{}width} & Draw the PDF’s pages as wide as the window, to scroll down.\\
none & \texttt{g p} & \texttt{pdf-\codebreak{}view-\codebreak{}goto-\codebreak{}page} & Turn to a page of the PDF by its number, asked for; in its text, go to where it begins.\\
none & \texttt{g l} & \texttt{pdf-\codebreak{}view-\codebreak{}last-\codebreak{}page} & Turn to the last page of the PDF.\\
none & \texttt{l}, \texttt{C-c l} & \texttt{pdf-\codebreak{}view-\codebreak{}next-\codebreak{}page} & Turn to the next page of the PDF.\\
none & \texttt{<next>} & \texttt{pdf-\codebreak{}view-\codebreak{}next-\codebreak{}page-\codebreak{}anywhere} & Turn to the next page of the PDF; in its text, go to where it begins.\\
none & \texttt{g e} & \texttt{pdf-\codebreak{}view-\codebreak{}page-\codebreak{}end} & Show the bottom of the page.\\
none & \texttt{g g} & \texttt{pdf-\codebreak{}view-\codebreak{}page-\codebreak{}top} & Show the top of the page.\\
none & \texttt{h}, \texttt{C-c h} & \texttt{pdf-\codebreak{}view-\codebreak{}previous-\codebreak{}page} & Turn to the page before in the PDF.\\
none & \texttt{<prior>} & \texttt{pdf-\codebreak{}view-\codebreak{}previous-\codebreak{}page-\codebreak{}anywhere} & Turn to the page before in the PDF; in its text, go to where it begins.\\
none & \texttt{<wheel-\codebreak{}up>}, \texttt{<wheel-\codebreak{}down>} & \texttt{pdf-\codebreak{}view-\codebreak{}scroll-\codebreak{}by-\codebreak{}wheel} & Scroll the page as the mouse wheel turned, as M-x pdf-view-scroll-down and M-x pdf-view-scroll-up do: down for <wheel-down>, on to the next page at a page’s end (each turn a page, the page fitted whole), up for <wheel-up>.\\
none & \texttt{j}, \texttt{C-c j} & \texttt{pdf-\codebreak{}view-\codebreak{}scroll-\codebreak{}down} & Scroll the page down, and at its end on to the top of the next.\\
none & \texttt{H}, \texttt{C-c H} & \texttt{pdf-\codebreak{}view-\codebreak{}scroll-\codebreak{}left} & Scroll a page wider than the window to the left.\\
none & \texttt{L}, \texttt{C-c L} & \texttt{pdf-\codebreak{}view-\codebreak{}scroll-\codebreak{}right} & Scroll a page wider than the window to the right.\\
none & \texttt{k}, \texttt{C-c k} & \texttt{pdf-\codebreak{}view-\codebreak{}scroll-up} & Scroll the page up, and at its top on to the end of the page before.\\
none & \texttt{M-x pdf-\codebreak{}view-\codebreak{}search} & \texttt{pdf-\codebreak{}view-\codebreak{}search} & Find text in the PDF and turn to its page: from the page shown on, asked for; again with nothing typed, the next place of the same text.\\
none & \texttt{d}, \texttt{C-c d} & \texttt{pdf-\codebreak{}view-\codebreak{}toggle-\codebreak{}dark} & Show the PDF’s pages dark, white as the theme’s background and black as its text, or as they are again; for every PDF while the editor runs (\textbf{PDF-VIEW-DARK-PAGES}).\\
none & \texttt{T}, \texttt{C-c T} & \texttt{pdf-\codebreak{}view-\codebreak{}toggle-\codebreak{}text} & Switch the PDF’s view between its pages and its text, to read, search and copy; the text opens at the page shown, and the pages at the page point is on.\\
none & \texttt{+}, \texttt{=}, \texttt{C-c +} & \texttt{pdf-\codebreak{}view-\codebreak{}zoom-in} & Draw the PDF’s pages larger, by \textbf{PDF-VIEW-ZOOM-STEP}; scroll with j, k, H and L.\\
none & \texttt{-}, \texttt{C-c -} & \texttt{pdf-\codebreak{}view-\codebreak{}zoom-\codebreak{}out} & Draw the PDF’s pages smaller, by \textbf{PDF-VIEW-ZOOM-STEP}.\\
none & \texttt{0}, \texttt{C-c 0} & \texttt{pdf-\codebreak{}view-\codebreak{}zoom-\codebreak{}reset} & Draw the PDF’s pages at the size of the fit chosen again, page or width (p, w), no longer zoomed in or out.\\
none & \texttt{C-x e} & \texttt{play-\codebreak{}recorded-\codebreak{}keys} & Play the keys recorded last (C-x ( \ldots{}\\
\bottomrule
\end{xltabular}

{\small
\begin{xltabular}{\textwidth}{@{}l >{\raggedright\arraybackslash\hsize=0.661\hsize}X l >{\raggedright\arraybackslash\hsize=1.339\hsize}X@{}}
\toprule
\textbf{Modal} & \textbf{Non-modal} & \textbf{Command} & \textbf{What it does}\\
\midrule
\endhead
none & \texttt{C-x C-SPC} & \texttt{pop-\codebreak{}global-\codebreak{}mark} & Go to the newest mark of the global mark ring, in its buffer, that mark going to the ring’s far end.\\
none & \texttt{M-x pop-\codebreak{}to-\codebreak{}mark-\codebreak{}command} & \texttt{pop-\codebreak{}to-\codebreak{}mark-\codebreak{}command} & Go to the mark, and take the mark before it off the mark ring as the mark.\\
none & \texttt{C-c M-m} & \texttt{pp-\codebreak{}macroexpand-\codebreak{}all-\codebreak{}last-\codebreak{}expression} & Show what the macro call at point expands to, every macro in it expanded too (SB-CLTL2:MACROEXPAND-ALL), pretty-printed in \textbf{Pp Macroexpand Output}; with a prefix argument (C-u), put it in after the call instead.\\
none & \texttt{C-c C-m}, \texttt{C-c RET} & \texttt{pp-\codebreak{}macroexpand-\codebreak{}last-\codebreak{}expression} & Show what the macro call at point expands to, once (MACROEXPAND-1), pretty-printed in \textbf{Pp Macroexpand Output}: the call beginning at point when an opening bracket is there, else the one before point.\\
\texttt{k} & \texttt{C-p}, \texttt{<up>} & \texttt{previous-\codebreak{}line} & Move point up one line (or the prefix argument’s worth), keeping the column when possible; where long lines wrap at a word’s edge (C-x T l), one row.\\
none & \texttt{M-g p} & \texttt{previous-\codebreak{}place} & Go to the place before the one gone to last in the list of places used last: a compile’s output (M-x compile), the lines a search of files found (M-s g), or the places the language servers found something wrong in a buffer, as M-g n goes to the next.\\
none & \texttt{l}, \texttt{RET}, \texttt{C-c l} & \texttt{push-\codebreak{}button} & Do what the button at point does: open what it names.\\
none & \texttt{<down-\codebreak{}mouse-1>} & \texttt{push-\codebreak{}button-\codebreak{}at-\codebreak{}click} & Put point where the first mouse button went down, as <down-mouse-1> does, and do what a button there does.\\
none & \texttt{M-x push-\codebreak{}mark-\codebreak{}command} & \texttt{push-\codebreak{}mark-\codebreak{}command} & Set the mark at point and activate the region, the mark before going on the mark ring; when the mark is at point already, only activate it.\\
none & \texttt{M-\%} & \texttt{query-\codebreak{}replace} & Replace literal text, asking at each match: y replaces, n skips, ! replaces all the rest, .\\
\texttt{r q} & \texttt{C-z r q} & \texttt{query-\codebreak{}replace-\codebreak{}regexp} & Replace a regular expression as C-z r r does, asking at each match: y replaces, n leaves it and goes to the next, ! replaces all the rest, .\\
none & \texttt{M-x query-\codebreak{}replace-\codebreak{}regexp-\codebreak{}thorough} & \texttt{query-\codebreak{}replace-\codebreak{}regexp-\codebreak{}thorough} & Replace a regular expression asking at each match, as query-replace-regexp does, finding every match of up to a megabyte whole however many lines it spans, with the cost and the limit of regex-search-forward-thorough: a match of more than a megabyte is still cut short, or missed when it needs text further on.\\
none & \texttt{h}, \texttt{q}, \texttt{C-c h} & \texttt{quit-\codebreak{}window} & Put the selected window back as it was before its buffer was shown: close it when it was made to show the buffer (a help window, \textbf{Messages}), else show the buffer it showed before.\\
none & \texttt{C-q} & \texttt{quoted-\codebreak{}insert} & Read the next key and put in the character it stands for, whatever that key does otherwise: C-q TAB a tab character where TAB puts spaces in, C-q C-l a form feed, C-q C-m a carriage return.\\
none & \texttt{C-x C-q} & \texttt{read-\codebreak{}only-\codebreak{}mode} & Make the current buffer read-only, or editable again if it is; with a count, read-only when it is positive, editable when it is zero or less.\\
\texttt{z z} & \texttt{C-l} & \texttt{recenter} & Scroll so that the current line is in the middle of the window.\\
none & \texttt{M-x recover-\codebreak{}file} & \texttt{recover-\codebreak{}file} & Take into the current buffer the changes auto-saved for its file (\textbf{AUTO-SAVE}) after asking, as one edit: C-/ takes it back, a save writes it to the file.\\
\texttt{m v} & \texttt{C-x r SPC} & \texttt{rectangle-\codebreak{}mark-\codebreak{}mode} & Start a rectangle: the block of columns between point and the mark on each line from one to the other; any motion sizes it, and its corner can stand past the end of a short line.\\
\texttt{U} & \texttt{C-?}, \texttt{M-\_}, \texttt{C-M-\_} & \texttt{redo} & Redo the last undone change: the cursor goes after what it makes again, but for a change made in the background (a large replacement, a sort), after which it stays where it is.\\
none & \texttt{C-r} & \texttt{regex-\codebreak{}search-\codebreak{}backward} & Search backward as you type for a regular expression; keys, replacing and the limit on matches of more than 4 KB as in regex-search-forward.\\
none & \texttt{M-x regex-\codebreak{}search-\codebreak{}backward-\codebreak{}thorough} & \texttt{regex-\codebreak{}search-\codebreak{}backward-\codebreak{}thorough} & Search backward as regex-search-backward does, finding every match of up to a megabyte whole however many lines it spans, with the cost and the limit of regex-search-forward-thorough: a match of more than a megabyte is still cut short, or missed when it needs text further on.\\
none & \texttt{C-s} & \texttt{regex-\codebreak{}search-\codebreak{}forward} & Search forward as you type for a regular expression (PCRE syntax); a pattern without upper-case letters ignores case, unless it names a case: :upper:, :lower:, \texttt{\textbackslash{}p\{Lu\}}, \texttt{\textbackslash{}p\{Ll\}} and \texttt{\textbackslash{}p\{Lt\}} heed it, as a capital letter does.\\
none & \texttt{M-x regex-\codebreak{}search-\codebreak{}forward-\codebreak{}thorough} & \texttt{regex-\codebreak{}search-\codebreak{}forward-\codebreak{}thorough} & Search forward as regex-search-forward does, finding every match of up to a megabyte (about a million characters of ASCII) whole however many lines it spans, where regex-search-forward may miss one of more than 4 KB that runs across lines and needs text further on (a block of lines matched by \texttt{BEGIN[\textbackslash{}s\textbackslash{}S]*?END}, say).\\
none & \texttt{C-x r +} & \texttt{register-\codebreak{}add} & Add to what a register keeps: a number grows by the prefix argument (1 without one, below 0 to take away), text has the selected text put after it.\\
none & \texttt{C-x r s} & \texttt{register-\codebreak{}copy} & Keep the selection under a character you type, as long as the editor runs: the text between point and mark, or the block of columns when a rectangle is selected (C-x r SPC).\\
none & \texttt{C-x r i} & \texttt{register-\codebreak{}insert} & Put what a register keeps into the text at point, asked for by its character: text as it was, with point after it and the mark before; a block (C-x r s of a rectangle) as a block of columns, as C-x r y pastes one, with a prefix argument each row that many times across; a number as digits.\\
none & \texttt{C-x r j} & \texttt{register-\codebreak{}jump} & Go to the place a register keeps (C-x r p), asked for by its character: in its buffer, shown, point where it was set, moved with the text around it.\\
none & \texttt{C-x r l} & \texttt{register-\codebreak{}list} & List the registers in \textbf{Registers}: the character that names each, what it keeps (a place, text, a rectangle, a number) and a look at that.\\
none & \texttt{C-x r n} & \texttt{register-\codebreak{}number} & Keep a whole number under a character you type: the prefix argument, else one asked for.\\
none & \texttt{C-x r p} & \texttt{register-\codebreak{}point} & Keep where point is under a character you type, as long as the editor runs: C-x r j goes back there, the text above it changed or not, and even after the buffer is closed when it was a file.\\
none & \texttt{M-x register-\codebreak{}prepend} & \texttt{register-\codebreak{}prepend} & Put the selected text before the text a register keeps, asked for by its character, as C-x r + puts it after.\\
none & \texttt{C-x n r} & \texttt{rename-\codebreak{}formatted} & Give a file a formatted name: its identifier kept, or made from when it was last changed; a title and keywords asked for, the file’s own offered; its extension kept.\\
none & \texttt{C-x n k} & \texttt{rename-\codebreak{}formatted-\codebreak{}keywords} & Give files formatted names with other keywords, their titles kept, asked for once for them all: the current buffer’s file, or in a list of a directory the marked files or the one on the line.\\
none & \texttt{M-x rename-\codebreak{}visited-\codebreak{}file} & \texttt{rename-\codebreak{}visited-\codebreak{}file} & Give the file the current buffer visits a new name, asked for, its own offered to change: a name alone keeps it in its folder, a path (\textasciitilde{}/notes/a.org, ../a.org) moves it there on the same drive.\\
\texttt{.} & \texttt{C-z .} & \texttt{repeat} & Run the last command again, other than repeat itself, with the prefix argument it had (one given to repeat is used instead) and as the key that ran it (C-x C-- shrinks again, it does not grow); repeat again, and it runs that command once more, as often as you like.\\
none & \texttt{M-x repeat-\codebreak{}menu} & \texttt{repeat-\codebreak{}menu} & Open the menu of a repeat: its start, time, how it repeats and when it ends, what they come to shown under it; RET writes it.\\
none & \texttt{M-x repeat-\codebreak{}menu-\codebreak{}write} & \texttt{repeat-\codebreak{}menu-\codebreak{}write} & Write the repeat of the open repeat menu where the date prompt it came from would have written its date; say why when it cannot be written.\\
\texttt{r r} & \texttt{C-z r r} & \texttt{replace-\codebreak{}regexp} & Replace a regular expression, as one edit: pattern/replacement, typed on one line or in two parts, as RET after the pattern types the / before the replacement.\\
none & \texttt{M-x replace-\codebreak{}regexp-\codebreak{}thorough} & \texttt{replace-\codebreak{}regexp-\codebreak{}thorough} & Replace a regular expression as replace-regexp does (the first match on each line, every match after a leading \%), finding every match of up to a megabyte whole however many lines it spans, with the cost and the limit of regex-search-forward-thorough: a match of more than a megabyte is still cut short, or missed when it needs text further on.\\
none & \texttt{M-x reverse-\codebreak{}region} & \texttt{reverse-\codebreak{}region} & Reverse the order of the lines the selection touches.\\
\bottomrule
\end{xltabular}
}

\begin{xltabular}{\textwidth}{@{}l >{\raggedright\arraybackslash\hsize=0.744\hsize}X l >{\raggedright\arraybackslash\hsize=1.256\hsize}X@{}}
\toprule
\textbf{Modal} & \textbf{Non-modal} & \textbf{Command} & \textbf{What it does}\\
\midrule
\endhead
none & \texttt{r}, \texttt{g r}, \texttt{C-c r} & \texttt{revert-\codebreak{}buffer} & Read the current buffer’s file again, replacing its text: one step that undo takes back, unless the text could no longer be trusted (then the history goes too).\\
none & \texttt{C-x d} & \texttt{rimu} & List a folder with rimu, asked for with completion from the folder of the path chosen last in a question for a path, else the current buffer’s: its name, then its entries a line each, folders first.\\
none & \texttt{c}, \texttt{C-c c} & \texttt{rimu-\codebreak{}archive-\codebreak{}create} & Pack the marked entries, from every directory they are marked in, or the one on the line, into one tar archive asked for, navigating to it as C-x C-f does: the directory of the other list beside this one offered first.\\
none & \texttt{Z}, \texttt{C-c Z} & \texttt{rimu-\codebreak{}archive-\codebreak{}extract} & Extract the archives marked, from every directory they are marked in, or the one on the line, into a directory asked for, navigating to it as C-x C-f does, made when it is not there: for one archive, a directory of its name beside it offered first.\\
none & \texttt{C-c e} & \texttt{rimu-\codebreak{}archive-\codebreak{}list} & Show what the archive shown holds, in a buffer of its own, which q closes: the entry on the line in a list of a directory (C-c e there); with C-u, or where there is none, asked for.\\
none & \texttt{C}, \texttt{C-c C} & \texttt{rimu-\codebreak{}copy} & Copy the marked entries, or the one on the line, to a place asked for: the directory of the other list beside this one offered first, or this one; one entry may be given a new name.\\
none & \texttt{: d} & \texttt{rimu-\codebreak{}decrypt-\codebreak{}files} & Decrypt the files marked (or the one on the line), each into its name without .gpg, for you alone, after asking once: their text is then on disk.\\
none & \texttt{=}, \texttt{C-c =} & \texttt{rimu-\codebreak{}diff} & Compare the two entries marked, from every directory they are marked in, the one shown first taken as the old, and show the differences in a buffer of their own, which q closes, as M-x diff-files does: unified, the old’s lines taken out after -, the other’s put in after +; two directories with all in them.\\
none & \texttt{C-o} & \texttt{rimu-\codebreak{}display} & Show the entry on the current line in another window, staying in the list: a directory in a list of its own, a file as l opens it.\\
none & \texttt{M-x rimu-\codebreak{}duplicates} & \texttt{rimu-\codebreak{}duplicates} & List in a rimu list the files that are the same as another, in groups: in the directories marked in this rimu list, or else under a directory asked for.\\
none & \texttt{C-x C-q} & \texttt{rimu-\codebreak{}edit-\codebreak{}names} & Edit the names of this list’s entries as text, a line each, in its window, with every key that edits text: C-c C-c renames those whose line changed, C-c C-k leaves them.\\
none & \texttt{: e} & \texttt{rimu-\codebreak{}encrypt-\codebreak{}files} & Encrypt the files marked (or the one on the line), each into FILE.gpg, to the keys or with a passphrase chosen in the minibuffer, asked once; and, when you say so, asked once too, wipe each file once it is encrypted (M-x wipe-file): with a passphrase, only after it, asked again, opens what was made.\\
none & \texttt{M-x rimu-\codebreak{}find} & \texttt{rimu-\codebreak{}find} & List in a rimu list the files and directories under a directory asked for whose names match a pattern: * stands for any characters, ? for one (\textbf{.pdf, report-202?.}); upper and lower case alike unless the pattern has an upper-case letter.\\
none & \texttt{M-x rimu-\codebreak{}find-\codebreak{}text} & \texttt{rimu-\codebreak{}find-\codebreak{}text} & List in a rimu list the files under a directory asked for whose text matches what you type, as M-s g matches a line: each word a regular expression or its own characters, all on one line.\\
none & \texttt{g f} & \texttt{rimu-\codebreak{}goto-\codebreak{}file} & Go to an entry of this list, its name asked for with completion; among those the search filters leave, while they narrow the list.\\
none & \texttt{L}, \texttt{C-c L} & \texttt{rimu-\codebreak{}hardlink} & Give the marked files, or the one on the line, a second name in a place asked for, as C asks: one file, two names, on the same drive only, and no directory.\\
none & \texttt{g h} & \texttt{rimu-\codebreak{}home} & List your home directory in this list’s place.\\
none & \texttt{C-x C-j} & \texttt{rimu-\codebreak{}jump} & List the directory of the current buffer’s file with rimu, point on the file; a buffer with no file lists the directory it stands in.\\
none & \texttt{g l} & \texttt{rimu-\codebreak{}last-\codebreak{}directory} & List again the directory this list listed before, point on the entry it was on; again, back to this one, as cd - goes back and forth.\\
none & \texttt{+}, \texttt{C-c +} & \texttt{rimu-\codebreak{}make-\codebreak{}directory} & Make a directory, asked for, in this list’s directory unless given elsewhere, with those above it that are not there; the list shows it.\\
none & \texttt{A}, \texttt{C-c A} & \texttt{rimu-\codebreak{}make-\codebreak{}formatted-\codebreak{}file} & Make a new file with a formatted name in this list’s directory, asking what new-file-formatted asks but the subdirectory (the title, the keywords, the kind and the template), with the header of its kind, saved at once; the list shows it, point on it.\\
none & \texttt{C-c f m} & \texttt{rimu-\codebreak{}mark-\codebreak{}master} & Mark the file on the line of this rimu list as a master (M), a file that the files marked (m) are compared with by the commands that compare files, a photo of the person looked for among photos, and go to the next line; u takes the mark away.\\
none & \texttt{M}, \texttt{C-c M} & \texttt{rimu-\codebreak{}move} & Move the marked entries, or the one on the line, into a directory asked for, navigating to it as C-x C-f does, only directories offered: the directory of the other list beside this one offered first.\\
none & \texttt{C-c C-k} & \texttt{rimu-\codebreak{}names-\codebreak{}cancel} & Leave the names as they are, and go back to the list.\\
none & \texttt{C-c C-c} & \texttt{rimu-\codebreak{}names-\codebreak{}finish} & Rename the entries whose line changed, and go back to the list: all or, when a line does not do, none, and why is said.\\
none & \texttt{>}, \texttt{C-c >} & \texttt{rimu-\codebreak{}next-\codebreak{}directory} & Go to the next directory of this list.\\
none & \texttt{l}, \texttt{RET}, \texttt{C-c l} & \texttt{rimu-\codebreak{}open} & Open the entry on the current line: a directory is listed in its place in this list; a file is visited as C-x C-f visits it, safely where that would be, and with a prefix argument (C-u) safely always.\\
none & \texttt{C-x 5 d} & \texttt{rimu-\codebreak{}other-\codebreak{}frame} & List a folder with rimu in another frame and work in it, the folder asked for as C-x d asks.\\
none & \texttt{<}, \texttt{C-c <} & \texttt{rimu-\codebreak{}previous-\codebreak{}directory} & Go to the previous directory of this list.\\
none & \texttt{o}, \texttt{C-c o} & \texttt{rimu-\codebreak{}quick-\codebreak{}edit} & Edit the file on the current line where the preview shows it, the quick edit: the preview becomes the file itself, opened as C-x C-f opens it (safely where that would be, and with C-u safely always), and the cursor goes there.\\
none & \texttt{M-x rimu-\codebreak{}quick-\codebreak{}edit-\codebreak{}finish} & \texttt{rimu-\codebreak{}quick-\codebreak{}edit-\codebreak{}finish} & Go back from the quick edit of this file to the preview beside the rimu list (as the cursor going back to the list does), asking first whether to save the changes made: y saves them; n lets them go (a buffer open before the quick edit keeps them); M-x keyboard-quit goes on with the quick edit.\\
none & \texttt{R}, \texttt{C-c R} & \texttt{rimu-\codebreak{}rename} & Rename the entry on the line, or the one marked, in its place, the new name asked for, the old offered: a name only, no directory; a name there already is refused.\\
none & \texttt{: s} & \texttt{rimu-\codebreak{}sign-\codebreak{}files} & Sign the files marked (or the one on the line) with your key, the kind of signature asked once.\\
none & \texttt{Y}, \texttt{C-c Y} & \texttt{rimu-\codebreak{}symlink} & Make symbolic links to the marked entries, or the one on the line, in a place asked for, as C asks: the directory of the other list beside this one offered first; one entry’s link may be given a new name.\\
none & \texttt{K}, \texttt{C-c K} & \texttt{rimu-\codebreak{}take-\codebreak{}out-\codebreak{}lines} & Take the lines of the entries marked in this directory out of this list, or with none marked here the line point is on, their marks taken away: the files stay as they are, only the list shows fewer, to look at what is left.\\
none & \texttt{I}, \texttt{C-c I} & \texttt{rimu-\codebreak{}toggle-\codebreak{}details} & Show or hide the details before each name in this list: its inode number, as ls -i shows it, then its permissions, links, owner, group, size and when it last changed, as ls -l shows them, and where a symbolic link leads.\\
none & \texttt{H}, \texttt{C-c H} & \texttt{rimu-\codebreak{}toggle-\codebreak{}dotfiles} & Show or hide the entries of this list whose names start with a dot, .\\
none & \texttt{G}, \texttt{C-c G} & \texttt{rimu-\codebreak{}toggle-\codebreak{}git} & Show or hide, after each name in this list, the entry’s last commit in git’s history: when it was made and the first line of its message.\\
none & \texttt{T}, \texttt{C-c T} & \texttt{rimu-\codebreak{}toggle-\codebreak{}picture} & Switch the preview between the picture of a file and its text: the first page of a PDF or its text, an image or its octets in the hex view.\\
none & \texttt{S}, \texttt{C-c S} & \texttt{rimu-\codebreak{}toggle-\codebreak{}sizes} & Show or hide, before each name in this list, how much it takes: a file its size, a directory the size of everything in it and under it, hidden files too, folders of another drive mounted in it left out, worked out in the background while the list shows it.\\
\bottomrule
\end{xltabular}

{\small
\begin{xltabular}{\textwidth}{@{}l >{\raggedright\arraybackslash\hsize=0.661\hsize}X l >{\raggedright\arraybackslash\hsize=1.339\hsize}X@{}}
\toprule
\textbf{Modal} & \textbf{Non-modal} & \textbf{Command} & \textbf{What it does}\\
\midrule
\endhead
none & \texttt{s} & \texttt{rimu-\codebreak{}toggle-\codebreak{}sort} & Order this list by name or by time, those changed last first; the mode line says which.\\
none & \texttt{a}, \texttt{C-c a} & \texttt{rimu-\codebreak{}touch} & Touch a name asked for in this list’s directory, as touch does: a name not there is made a new, empty file, the list on it for l to open it; a file or directory there has its times set to now (a link’s own, not those of where it leads).\\
none & \texttt{g t} & \texttt{rimu-\codebreak{}trash} & List what is in the trash, as the buffer list (C-x C-b) lists buffers: a section for your own trash and one for each drive’s, the last to go there first, each thing’s name, when it went and where it was.\\
none & \texttt{X}, \texttt{C-c X} & \texttt{rimu-\codebreak{}trash-\codebreak{}empty} & Empty the trash: everything the list of the trash shows deleted for good, after asking twice, in the background.\\
none & \texttt{R}, \texttt{C-c R} & \texttt{rimu-\codebreak{}trash-\codebreak{}restore} & Put the marked things (>), or with none marked the one on the current line, back where they were, after asking, where they go said; another name asked for when one’s is taken, the directories above it made again when they are gone.\\
none & \texttt{h}, \texttt{C-c h} & \texttt{rimu-up} & List the directory holding this one, in this list, point on the one left; in a list of files found, list the directory they were found in.\\
none & \texttt{: v} & \texttt{rimu-\codebreak{}verify-\codebreak{}files} & Check the signatures of the files marked (or the one on the line), each said in turn and kept in \textbf{Messages}.\\
none & \texttt{g w} & \texttt{rimu-\codebreak{}wipe-\codebreak{}files} & Wipe the files marked (or the one on the line) as M-x wipe-file does, asked once: overwritten, then deleted.\\
none & \texttt{M-x rotate-\codebreak{}colouring} & \texttt{rotate-\codebreak{}colouring} & Colour the current buffer’s language in the next style: rich, then sparse, then none, then rich again.\\
none & \texttt{C-x w ,} & \texttt{rotate-\codebreak{}window-\codebreak{}layout-\codebreak{}anticlockwise} & Turn the layout of windows a quarter turn anticlockwise: windows one above the other come to stand side by side, the top one on the left, and windows side by side come to stand one above the other, the right one on top.\\
none & \texttt{C-x w .} & \texttt{rotate-\codebreak{}window-\codebreak{}layout-\codebreak{}clockwise} & Turn the layout of windows a quarter turn clockwise: windows one above the other come to stand side by side, the top one on the right, and windows side by side come to stand one above the other, the left one on top.\\
none & \texttt{C-x w r} & \texttt{rotate-\codebreak{}windows} & Move each window’s buffer to the next window, from the top left, and the last window’s to the first; you go with your buffer.\\
none & \texttt{C-x w R} & \texttt{rotate-\codebreak{}windows-\codebreak{}back} & Move each window’s buffer to the window before it, from the top left, and the first window’s to the last; you go with your buffer.\\
none & \texttt{C-x C-s} & \texttt{save-\codebreak{}buffer} & Save the current buffer to its file, asking for one if it has none, as C-x C-f asks.\\
none & \texttt{C-x C-c} & \texttt{save-\codebreak{}buffers-\codebreak{}kill-\codebreak{}editor} & Exit the editor, asking first if a buffer visiting a file, or one made by C-x n e, has unsaved changes.\\
none & \texttt{C-x s} & \texttt{save-\codebreak{}some-\codebreak{}buffers} & Save every buffer visiting a file with changes not saved, after asking once, the files named: the files a rename across a project changed, say.\\
none & \texttt{<wheel-\codebreak{}up>}, \texttt{<wheel-\codebreak{}down>} & \texttt{scroll-\codebreak{}by-\codebreak{}wheel} & Scroll the window under the mouse, else the selected one, as the mouse wheel turned: \textbf{WHEEL-SCROLL-LINES} lines on for <wheel-down>, back for <wheel-up>; that window’s point stays in view.\\
none & \texttt{M-v}, \texttt{<prior>} & \texttt{scroll-\codebreak{}down-\codebreak{}command} & Scroll the window backward by nearly a screenful: by the rows it shows, where long lines wrap.\\
none & \texttt{C-M-v} & \texttt{scroll-\codebreak{}other-\codebreak{}window} & Scroll the next window forward by nearly a screenful, leaving the selected one as it is.\\
none & \texttt{C-M-V} & \texttt{scroll-\codebreak{}other-\codebreak{}window-\codebreak{}down} & Scroll the next window backward by nearly a screenful, leaving the selected one as it is.\\
none & \texttt{C-v}, \texttt{<next>} & \texttt{scroll-\codebreak{}up-\codebreak{}command} & Scroll the window forward by nearly a screenful: by the rows it shows, where long lines wrap.\\
none & \texttt{M-x search-\codebreak{}abort} & \texttt{search-\codebreak{}abort} & Cancel the search and return to where it began; a search filter being typed goes, the list shown as the filters kept leave it.\\
none & \texttt{M-x search-\codebreak{}exit} & \texttt{search-\codebreak{}exit} & End the search, leaving point at the match and the mark where the search began; when the string holds a replacement, carry it out; for a search filter, keep it on the list.\\
none & \texttt{M-x search-\codebreak{}files-\codebreak{}list} & \texttt{search-\codebreak{}files-\codebreak{}list} & While searching files (M-s g), list the lines found so far in a buffer of their own, one each as path:line: text, to go through with RET or to change at once, with C-x C-q, and write back into their files.\\
none & \texttt{C-x / h} & \texttt{search-\codebreak{}filter-\codebreak{}hide-\codebreak{}records} & Hide the records whose line a regular expression matches, showing the others, deleting nothing: asked for as C-x / f asks, kept with the other filters.\\
none & \texttt{C-x / l} & \texttt{search-\codebreak{}filter-\codebreak{}list} & List the search filters narrowing the list here (C-x / s) in a buffer of their own, each with how many entries it leaves; there d takes away the filter on the line and D all of them.\\
none & \texttt{C-x / s} & \texttt{search-\codebreak{}filter-\codebreak{}regex} & Narrow the list here to the entries a regular expression matches, as it is typed: those whose line matches, or what more the list holds of them (a mailbox the headers and contents of its messages).\\
none & \texttt{d}, \texttt{C-c d} & \texttt{search-\codebreak{}filter-\codebreak{}remove} & Take away the search filter on the current line of the list of filters (C-x / l): the list it narrowed widens, the filters after it narrowing again what those before it leave, in the background where that takes time.\\
none & \texttt{D}, \texttt{C-c D}, \texttt{C-x / r} & \texttt{search-\codebreak{}filter-\codebreak{}remove-\codebreak{}all} & Take away every search filter narrowing the list here, or the list the list of filters here is for: the whole list is shown again.\\
none & \texttt{C-x / f} & \texttt{search-\codebreak{}filter-\codebreak{}show-\codebreak{}records} & Show only the records whose line a regular expression matches, hiding the others, as C-x / s does, deleting nothing.\\
none & \texttt{M-s g} & \texttt{search-\codebreak{}in-\codebreak{}files} & Find text in the files of the project the current buffer’s file is in (see \textbf{PROJECT-ROOT-MARKERS}), or under its directory when it is in none, or under the directory it lists (C-x d), or \textbf{DEFAULT-DIRECTORY} for another buffer with no file; with a prefix argument, in a directory asked for.\\
none & \texttt{C-x n g} & \texttt{search-\codebreak{}in-\codebreak{}formatted} & Find text in \textbf{FORMATTED-FILE-DIRECTORY}, where your files with formatted names are: every file under it, as search-in-files finds text in a directory; the lines with the words typed are listed as you type, the files read in the background meanwhile, M-r switching to literal text and back, and M-x minibuffer-finish visits one.\\
none & \texttt{M-x search-\codebreak{}insert-\codebreak{}symbol} & \texttt{search-\codebreak{}insert-\codebreak{}symbol} & During a search: add the symbol at point where the search began to the search string, in a regular expression as a whole symbol, never found inside a longer one.\\
none & \texttt{M-x search-\codebreak{}next-\codebreak{}match} & \texttt{search-\codebreak{}next-\codebreak{}match} & Go to the next match of the search just ended with RET, without its prompt: C-s right after the search does, again and again, the last match followed by the first.\\
none & \texttt{M-x search-\codebreak{}previous-\codebreak{}match} & \texttt{search-\codebreak{}previous-\codebreak{}match} & Go to the previous match of the search just ended with RET, without its prompt: C-r right after the search does, again and again, the first match followed by the last.\\
none & \texttt{M-x search-\codebreak{}repeat-\codebreak{}backward} & \texttt{search-\codebreak{}repeat-\codebreak{}backward} & During a search: the next match backward.\\
none & \texttt{M-x search-\codebreak{}repeat-\codebreak{}forward} & \texttt{search-\codebreak{}repeat-\codebreak{}forward} & During a search: the next match forward (with an empty string, the last search string).\\
none & \texttt{C-c C-k} & \texttt{search-\codebreak{}results-\codebreak{}cancel} & Throw away the changes made to the lines found, and make them read-only again; the files are not touched.\\
none & \texttt{C-x C-q} & \texttt{search-\codebreak{}results-\codebreak{}edit} & Make the text after the path and line number of the lines found changeable, with the keys that change text: the path and line number stay as they are, and a line is not split or joined.\\
none & \texttt{C-c C-c} & \texttt{search-\codebreak{}results-\codebreak{}finish} & Write the lines changed here into their files.\\
\bottomrule
\end{xltabular}
}

\begin{xltabular}{\textwidth}{@{}l >{\raggedright\arraybackslash\hsize=0.571\hsize}X l >{\raggedright\arraybackslash\hsize=1.429\hsize}X@{}}
\toprule
\textbf{Modal} & \textbf{Non-modal} & \textbf{Command} & \textbf{What it does}\\
\midrule
\endhead
none & \texttt{C-x T s} & \texttt{search-\codebreak{}toggle-\codebreak{}highlight} & Light the matches of the last search that ended at a match again, in its buffer, or take them away when they are lit, as M-x keyboard-quit does.\\
none & \texttt{M-x search-\codebreak{}toggle-\codebreak{}regex} & \texttt{search-\codebreak{}toggle-\codebreak{}regex} & During a search, or a search of files (M-s g): switch between a regular expression and literal text.\\
\texttt{m C} & \texttt{C-z m C} & \texttt{select-\codebreak{}comment} & Select the comment around point with the comment lines right before and after it.\\
\texttt{m c} & \texttt{C-z m c} & \texttt{select-\codebreak{}comment-\codebreak{}inside} & Select the comment around point.\\
\texttt{m X} & \texttt{C-z m X} & \texttt{select-\codebreak{}element} & Select all of the element around point, its tags too.\\
\texttt{m x} & \texttt{C-z m x} & \texttt{select-\codebreak{}element-\codebreak{}inside} & Select what is between the tags of the element around point: in HTML, XML, JSX, Svelte and Vue.\\
\texttt{m E} & \texttt{C-z m E} & \texttt{select-\codebreak{}entry} & Select all of the entry around point, its key and value or the list’s item.\\
\texttt{m e} & \texttt{C-z m e} & \texttt{select-\codebreak{}entry-\codebreak{}inside} & Select the key or the value at point of the entry around it: of a key and its value in JSON, YAML or TOML, a field of a struct, an item of a list, as the language’s rules know them.\\
\texttt{m F} & \texttt{C-z m F} & \texttt{select-\codebreak{}function} & Select all of the function around point.\\
\texttt{m f} & \texttt{C-z m f} & \texttt{select-\codebreak{}function-\codebreak{}inside} & Select the body of the function around point: the block of a Rust fn or closure, the body of a shell function, the forms of a Lisp defun, defmacro, defgeneric, defmethod or lambda after its lambda list.\\
\texttt{m i} & \texttt{C-z m i} & \texttt{select-\codebreak{}inside-\codebreak{}pair} & Select what is inside a pair of delimiters, the delimiters left out: the pair of the delimiter at point, or the one whose half is typed next ((, ), [, “, ‘, =, * and the rest of \textbf{pair-delimiters}).\\
\texttt{V} & \texttt{C-z V} & \texttt{select-\codebreak{}lines} & Start a line selection anchored on the current line, or cancel a live one.\\
\texttt{J} & \texttt{C-z J} & \texttt{select-\codebreak{}next-\codebreak{}line} & Move down, extending the line selection if one is active.\\
none & \texttt{M-n} & \texttt{select-\codebreak{}next-\codebreak{}sibling} & Select the next piece of syntax beside the selection, or beside the one at point: the next argument, statement or item, and after the last one, the next beside the piece around them.\\
\texttt{m a} & \texttt{C-z m a} & \texttt{select-\codebreak{}pair} & Select a pair of delimiters and what it holds, as select-inside-pair finds it.\\
\texttt{m p} & \texttt{C-z m p} & \texttt{select-\codebreak{}paragraph} & Select the paragraph at point with one blank line, the one before it (the one after, for the first paragraph), so that deleting it leaves one blank line between its neighbours: from a blank line between two, the one below.\\
\texttt{m P} & \texttt{C-z m P} & \texttt{select-\codebreak{}paragraph-\codebreak{}backward} & Select the paragraph at point, or grow a paragraph selection backward: from a blank line between two, the one above.\\
\texttt{m R} & \texttt{C-z m R} & \texttt{select-\codebreak{}parameter} & Select the parameter or argument around point with the comma after it.\\
\texttt{m r} & \texttt{C-z m r} & \texttt{select-\codebreak{}parameter-\codebreak{}inside} & Select the parameter or argument around point: of a Rust function’s parameters, a call’s arguments or a generic’s parameters; of a shell command’s arguments; of a Lisp lambda list or a call’s arguments.\\
\texttt{K} & \texttt{C-z K} & \texttt{select-\codebreak{}previous-\codebreak{}line} & Move up, extending the line selection if one is active.\\
none & \texttt{M-p} & \texttt{select-\codebreak{}previous-\codebreak{}sibling} & Select the piece of syntax before the selection, or before the one at point: the argument, statement or item before it, and before the first one, the one before the piece around them.\\
\texttt{m s} & \texttt{C-z m s} & \texttt{select-\codebreak{}sentence} & Select the sentence at point, as far as its full stop; from the gap between two, the one ahead.\\
\texttt{m S} & \texttt{C-z m S} & \texttt{select-\codebreak{}sentence-\codebreak{}backward} & Select the sentence at point, or grow a sentence selection backward: from the gap between two, the one behind.\\
\texttt{m A} & \texttt{C-z m A} & \texttt{select-\codebreak{}sexp} & Select the bracketed expression around point, its brackets included: point at the start, the mark at the end.\\
\texttt{m I} & \texttt{C-z m I} & \texttt{select-\codebreak{}sexp-\codebreak{}inside} & Select what is inside the bracketed expression around point, its brackets left out: point at the start, the mark at the end.\\
\texttt{m W} & \texttt{C-z m W} & \texttt{select-\codebreak{}symbol} & Select the symbol at point (a name with its \_ and -, Lisp’s own characters in Lisp), punctuation at either end left out: from the gap between two, the one ahead, brackets and quotes crossed.\\
\texttt{m B} & \texttt{C-z m B} & \texttt{select-\codebreak{}symbol-\codebreak{}backward} & Select the symbol at point, or grow a symbol selection backward: from the gap between two, the one behind.\\
\texttt{m Y} & \texttt{C-z m Y} & \texttt{select-\codebreak{}test} & Select all of the test around point.\\
\texttt{m y} & \texttt{C-z m y} & \texttt{select-\codebreak{}test-\codebreak{}inside} & Select the body of the test around point: a Rust function marked \#[test], a Python function whose name begins with test, a Go TestX or BenchmarkX, and the tests of Zig, Elixir, Dart, Julia and others, as the language’s rules know them.\\
\texttt{m T} & \texttt{C-z m T} & \texttt{select-\codebreak{}type} & Select all of the type around point.\\
\texttt{m t} & \texttt{C-z m t} & \texttt{select-\codebreak{}type-\codebreak{}inside} & Select the body of the type around point: the fields, variants or items of a Rust struct, enum, union, trait or impl; in Lisp, what follows the name in a defclass, defstruct, define-condition or deftype.\\
\texttt{\%} & none & \texttt{select-\codebreak{}whole-\codebreak{}buffer} & Select the whole buffer: point at the start, the mark at the end.\\
none & \texttt{C-M-k} & \texttt{select-\codebreak{}window-\codebreak{}above} & Select the window above the selected one: the one over point’s column.\\
none & \texttt{C-M-j} & \texttt{select-\codebreak{}window-\codebreak{}below} & Select the window below the selected one: the one under point’s column.\\
none & \texttt{C-M-h} & \texttt{select-\codebreak{}window-\codebreak{}left} & Select the window to the left of the selected one: the one beside point’s line.\\
none & \texttt{C-M-l} & \texttt{select-\codebreak{}window-\codebreak{}right} & Select the window to the right of the selected one: the one beside point’s line.\\
\texttt{m w} & \texttt{C-z m w} & \texttt{select-\codebreak{}word} & Select the word at point; from the gap between two words, the one ahead (the last one at the end of the text).\\
\texttt{m b} & \texttt{C-z m b} & \texttt{select-\codebreak{}word-\codebreak{}backward} & Select the word at point, or grow a word selection backward, as select-word grows it forward: from the gap between two words, the one behind.\\
none & \texttt{M-x self-\codebreak{}insert-\codebreak{}command} & \texttt{self-\codebreak{}insert-\codebreak{}command} & Insert the character that invoked this command, prefix-argument times.\\
none & \texttt{H}, \texttt{C-c H} & \texttt{server-\codebreak{}agent-\codebreak{}destinations} & Name, in the minibuffer, the machines the keys of your ssh agent may be used for from the server on the line, or those marked, commas between them (github.com, me@git.example); none for any.\\
\bottomrule
\end{xltabular}

{\small
\begin{xltabular}{\textwidth}{@{}l >{\raggedright\arraybackslash\hsize=0.822\hsize}X l >{\raggedright\arraybackslash\hsize=1.178\hsize}X@{}}
\toprule
\textbf{Modal} & \textbf{Non-modal} & \textbf{Command} & \textbf{What it does}\\
\midrule
\endhead
none & \texttt{M-x server-\codebreak{}bring-\codebreak{}language-\codebreak{}server} & \texttt{server-\codebreak{}bring-\codebreak{}language-\codebreak{}server} & Copy a language server installed on this machine to the server on the line of the list of trusted servers, or those marked, chosen in the minibuffer among those the editor knows that are here: nothing is fetched from anywhere.\\
none & \texttt{M-x server-\codebreak{}copy-\codebreak{}keys} & \texttt{server-\codebreak{}copy-\codebreak{}keys} & Copy the public keys marked in the list of GnuPG keys (M-x gpg-list-keys), or the one at the cursor, into gpg’s keyring on a trusted server asked for, so that M-x gpg-encrypt-file there can encrypt to them: exported here and imported there over ssh, in the background.\\
none & \texttt{M-x server-\codebreak{}list} & \texttt{server-\codebreak{}list} & List the trusted servers, the machines whose daemons a client here (tuatara --client=HOST) may give more than the keys typed and the size of the screen, each with what it is given: the ssh agent, with the machines its keys may be used for from there; gpg-agent; the clipboard; files dropped on the window, sent to a folder there.\\
none & \texttt{M-x server-\codebreak{}remove-\codebreak{}language-\codebreak{}server} & \texttt{server-\codebreak{}remove-\codebreak{}language-\codebreak{}server} & Take away a language server M-x server-bring-language-server brought to the server on the line of the list of trusted servers, or those marked, chosen in the minibuffer among those it lists there.\\
none & \texttt{A}, \texttt{C-c A} & \texttt{server-\codebreak{}toggle-\codebreak{}agent} & Forward your ssh agent to the server on the line, or those marked, or no longer: on, the programs there can sign with your keys while a client shows its daemon, each use asked about here, and only for the machines named by H where OpenSSH can limit them.\\
none & \texttt{C}, \texttt{C-c C} & \texttt{server-\codebreak{}toggle-\codebreak{}clipboard} & Share the clipboard with the server on the line, or those marked, or no longer: on, what is copied there reaches the clipboard here, as text with no control characters, and the daemon there may read what is copied here.\\
none & \texttt{F}, \texttt{C-c F} & \texttt{server-\codebreak{}toggle-\codebreak{}files} & Send the files dropped on a window showing the daemon of the server on the line, or those marked, to a folder there, asked for (from the home folder there unless it begins with /), or no longer: on, a file dropped is copied there, its size checked first, and nothing of it run.\\
none & \texttt{G}, \texttt{C-c G} & \texttt{server-\codebreak{}toggle-\codebreak{}gpg} & Forward gpg-agent to the server on the line, or those marked, or no longer: on, gpg there signs and decrypts with your keys, which stay here, and cannot export or change them.\\
none & \texttt{+}, \texttt{C-c +} & \texttt{server-\codebreak{}trust} & Trust a machine asked for in the minibuffer, as tuatara --client= names it (build.example, me@build.example, or an alias of \textasciitilde{}/.ssh/config): it goes on the list of trusted servers (M-x server-list) with nothing given yet, each thing then turned on there by itself.\\
none & \texttt{T}, \texttt{C-c T} & \texttt{server-\codebreak{}tuatara-\codebreak{}program} & Name, in the minibuffer, the tuatara to run on the server on the line, or those marked, when it is not the one on the PATH of the commands ssh runs there: /opt/tuatara/bin/tuatara, or a name found on that PATH; empty for the one on that PATH.\\
none & \texttt{C-x f} & \texttt{set-\codebreak{}fill-\codebreak{}column} & Set the column the current buffer’s lines are filled to by M-q and C-z g q, wrapped at by C-x T l and marked by C-x T v, while the editor runs: a count given first (C-u 72 C-x f), or else the column asked for, the buffer’s offered, which M-x minibuffer-finish alone keeps.\\
none & \texttt{M-x set-\codebreak{}frame-\codebreak{}font} & \texttt{set-\codebreak{}frame-\codebreak{}font} & Choose the window’s font among the fixed-width fonts installed: a family, with any style and size typed after it (DejaVu Sans Mono Bold 12); without a size, the size in use stays, and so does any scaling by C-x C-+.\\
none & \texttt{C-x RET C-\textbackslash{}} & \texttt{set-\codebreak{}input-\codebreak{}method} & Choose an input method and turn it on in this buffer (C-x RET C-\textbackslash{}).\\
none & \texttt{M-x set-\codebreak{}keys-\codebreak{}counter} & \texttt{set-\codebreak{}keys-\codebreak{}counter} & Set the number M-x insert-keys-counter inserts next, asked for, or the prefix argument when one is given.\\
none & \texttt{M-x set-\codebreak{}keys-\codebreak{}counter-\codebreak{}format} & \texttt{set-\codebreak{}keys-\codebreak{}counter-\codebreak{}format} & Set how M-x insert-keys-counter writes the counter: \%d for the number, \%03d for it in three digits with zeros before (001), \%3d with spaces before, \%-3d with spaces after, \%\% a per cent sign, the rest as it is (Item \%d.\\
none & \texttt{C-SPC} & \texttt{set-\codebreak{}mark-\codebreak{}command} & Set the mark at point and activate the region (C-SPC).\\
none & \texttt{M-x sha256-\codebreak{}sum} & \texttt{sha256-\codebreak{}sum} & Show the SHA-256 sum of the file the buffer visits as it is on disk, or with a selection of the octets selected, or in a buffer with no file of its whole text; keep it in \textbf{Messages} and put it on the kill ring.\\
none & \texttt{g i} & \texttt{show-\codebreak{}file-\codebreak{}metadata} & Show what the file shown is and what it says of itself, in a buffer of its own, which q closes: the file the buffer visits, the entry on the line in a list of a directory, the PDF or picture in its view; with C-u, or where there is none, asked for.\\
none & \texttt{M-x show-\codebreak{}file-\codebreak{}name} & \texttt{show-\codebreak{}file-\codebreak{}name} & Say the full name of the file the current buffer shows, its directory and name, in the echo area, where \textbf{Messages} keeps it too: in a rimu list, that of the entry on the current line; in a buffer with no file, the directory it stands in.\\
none & \texttt{M-x show-\codebreak{}paren-\codebreak{}mode} & \texttt{show-\codebreak{}paren-\codebreak{}mode} & Switch on or off showing the bracket that pairs with the one at point, in every buffer: when on, the bracket before point when it closes, else the one after it when it opens, is shown with its partner in :paren-match, or in :paren-mismatch when the partner is of another kind or there is none.\\
none & \texttt{M-i} & \texttt{shrink-\codebreak{}selection} & Go back to the selection before the last M-o, or else select the first piece of syntax inside the selection; without a selection, the smallest piece of syntax at point.\\
none & \texttt{M-x shrink-\codebreak{}window} & \texttt{shrink-\codebreak{}window} & Make the selected window a row shorter, or the prefix argument’s rows; they go to the window below it, or above.\\
none & \texttt{C-x \{} & \texttt{shrink-\codebreak{}window-\codebreak{}horizontally} & Make the selected window a column narrower, or the prefix argument’s columns; they go to the window to its right, or its left.\\
none & \texttt{C-x -} & \texttt{shrink-\codebreak{}window-\codebreak{}if-\codebreak{}larger-\codebreak{}than-\codebreak{}buffer} & Make the selected window as short as its buffer’s lines need, when it is taller.\\
none & \texttt{M-x snippet-\codebreak{}next-\codebreak{}field} & \texttt{snippet-\codebreak{}next-\codebreak{}field} & Go to the next field of the snippet being filled in; past the last, end it where it puts the cursor.\\
none & \texttt{M-x snippet-\codebreak{}previous-\codebreak{}field} & \texttt{snippet-\codebreak{}previous-\codebreak{}field} & Go back to the field of the snippet before the one the cursor is in.\\
none & \texttt{C-x T l} & \texttt{soft-\codebreak{}wrap-\codebreak{}mode} & Wrap the long lines of the current buffer at a word’s edge, or stop: a line longer than the window is shown on as many rows as it needs, at the fill column (C-x f), or the window’s edge when that is nearer; the text is not changed.\\
none & \texttt{M-x sort-\codebreak{}columns} & \texttt{sort-\codebreak{}columns} & Sort the lines from the one where the selection starts to the one where it ends by the text between the columns of its two ends, as text; with a prefix argument, in descending order.\\
none & \texttt{M-x sort-\codebreak{}columns-\codebreak{}reverse} & \texttt{sort-\codebreak{}columns-\codebreak{}reverse} & Sort the lines from the one where the selection starts to the one where it ends by the text between the columns of its two ends, in descending order: as M-x sort-columns does with a prefix argument.\\
none & \texttt{M-x sort-\codebreak{}fields} & \texttt{sort-\codebreak{}fields} & Sort the lines the selection touches by one field of each, as text: fields are separated by spaces and tabs, and the prefix argument names the field (the first by default; a negative one counts from the end of the line).\\
none & \texttt{M-x sort-\codebreak{}fields-\codebreak{}reverse} & \texttt{sort-\codebreak{}fields-\codebreak{}reverse} & Sort the lines the selection touches by one field of each, as text, in descending order: as M-x sort-fields, the prefix argument naming the field (a negative one counting from the end of the line).\\
none & \texttt{M-x sort-\codebreak{}lines} & \texttt{sort-\codebreak{}lines} & Sort the lines the selection touches by their text: by the codes of their characters, the case of ASCII letters ignored when \textbf{SORT-FOLD-CASE} is true.\\
none & \texttt{M-x sort-\codebreak{}lines-\codebreak{}reverse} & \texttt{sort-\codebreak{}lines-\codebreak{}reverse} & Sort the lines the selection touches by their text in descending order, as M-x sort-lines does with a prefix argument: lines that compare equal come in reverse order too.\\
none & \texttt{M-x sort-\codebreak{}numeric-\codebreak{}fields} & \texttt{sort-\codebreak{}numeric-\codebreak{}fields} & Sort the lines the selection touches by the number at the start of one field of each (the prefix argument’s, as for sort-fields): 0x and hexadecimal digits, or a number in \textbf{SORT-NUMERIC-BASE}, decimal by default and then with a fraction and exponent; a field without a number counts as 0.\\
none & \texttt{M-x sort-\codebreak{}numeric-\codebreak{}fields-\codebreak{}reverse} & \texttt{sort-\codebreak{}numeric-\codebreak{}fields-\codebreak{}reverse} & Sort the lines the selection touches by the number at the start of one field of each, largest first: as M-x sort-numeric-fields, the prefix argument naming the field.\\
none & \texttt{M-x sort-\codebreak{}pages} & \texttt{sort-\codebreak{}pages} & Sort the pages the selection touches by their text, pages ending with a form feed at the start of a line.\\
none & \texttt{M-x sort-\codebreak{}pages-\codebreak{}reverse} & \texttt{sort-\codebreak{}pages-\codebreak{}reverse} & Sort the pages the selection touches by their text in descending order: as M-x sort-pages does with a prefix argument.\\
none & \texttt{M-x sort-\codebreak{}paragraphs} & \texttt{sort-\codebreak{}paragraphs} & Sort the paragraphs the selection touches by their text: paragraphs are separated by blank lines, which stay where they were.\\
none & \texttt{M-x sort-\codebreak{}paragraphs-\codebreak{}reverse} & \texttt{sort-\codebreak{}paragraphs-\codebreak{}reverse} & Sort the paragraphs the selection touches by their text in descending order: as M-x sort-paragraphs does with a prefix argument.\\
none & \texttt{M-x sort-\codebreak{}regexp-\codebreak{}fields} & \texttt{sort-\codebreak{}regexp-\codebreak{}fields} & Sort the matches of a regular expression within the selection: asks for the pattern of the records, then for their key, either \textbackslash{}\& for the whole record, \textbackslash{}1 to \textbackslash{}9 for a group of the record’s pattern, or a pattern whose first match in the record is its key (an empty key when there is none).\\
\bottomrule
\end{xltabular}
}

\begin{xltabular}{\textwidth}{@{}l >{\raggedright\arraybackslash\hsize=0.895\hsize}X l >{\raggedright\arraybackslash\hsize=1.105\hsize}X@{}}
\toprule
\textbf{Modal} & \textbf{Non-modal} & \textbf{Command} & \textbf{What it does}\\
\midrule
\endhead
none & \texttt{M-x sort-\codebreak{}regexp-\codebreak{}fields-\codebreak{}reverse} & \texttt{sort-\codebreak{}regexp-\codebreak{}fields-\codebreak{}reverse} & Sort the matches of a regular expression within the selection by their keys in descending order: as M-x sort-regexp-fields does with a prefix argument, asking for the records’ pattern and their key.\\
\texttt{z g} & \texttt{C-z z g} & \texttt{spellcheck-\codebreak{}add-\codebreak{}word} & Add the word at point to a language’s personal word list, kept by enchant in \textasciitilde{}/.config/enchant/, the list other programs using enchant share: asked which language when words are checked in more than one.\\
\texttt{z f} & \texttt{C-z z f} & \texttt{spellcheck-\codebreak{}add-\codebreak{}word-\codebreak{}to-\codebreak{}file} & Keep the word at point as right in this file only: added to the file’s own word list, a LocalWords line at its end, written as a comment (\# in Org, ;; in Lisp) that other spelling checkers read too.\\
\texttt{z =} & \texttt{C-z z =} & \texttt{spellcheck-\codebreak{}correct} & Correct the word spelt wrong at point, or else the one nearest point in the window: choose one of enchant’s suggestions, or type the word as it should be, or keep the word, while the editor runs, in the file’s own word list (as C-z z f does), or for good in a language’s word list (as C-z z g does).\\
none & \texttt{M-x spellcheck-\codebreak{}languages} & \texttt{spellcheck-\codebreak{}languages} & Set the languages the current buffer’s words are checked in, by their names as enchant knows them, apart by spaces (en\_NZ sv\_SE): a word is right when one of them takes it.\\
none & \texttt{M-x spellcheck-\codebreak{}mode} & \texttt{spellcheck-\codebreak{}mode} & Mark the words spelt wrong in the current buffer, or stop: as they are shown, underlined in the :spellcheck face; in code only in comments and strings, unless \textbf{SPELLCHECK-MODES} says otherwise for its mode (:comments, :strings, or :all for the code too).\\
\texttt{z n} & \texttt{C-z z n} & \texttt{spellcheck-\codebreak{}next} & Go to the next word spelt wrong after point (the count’s worth), as M-x spellcheck-mode marks them, though they are not shown yet.\\
\texttt{z p} & \texttt{C-z z p} & \texttt{spellcheck-\codebreak{}previous} & Go to the word spelt wrong before point (the count’s worth), as M-x spellcheck-mode marks them.\\
none & \texttt{M-x split-\codebreak{}append} & \texttt{split-\codebreak{}append} & Type after every place of the split, keeping what is there.\\
none & \texttt{M-x split-\codebreak{}change} & \texttt{split-\codebreak{}change} & Empty every place of the split, then type at all of them.\\
none & \texttt{M-x split-\codebreak{}delete} & \texttt{split-\codebreak{}delete} & Delete every place of the split, one change for undo, and end the split.\\
none & \texttt{M-x split-\codebreak{}insert} & \texttt{split-\codebreak{}insert} & Type before every place of the split, keeping what is there: what is typed at one appears at all.\\
none & \texttt{M-x split-\codebreak{}quit} & \texttt{split-\codebreak{}quit} & End the split, keeping what it changed.\\
\texttt{f} & \texttt{C-z f} & \texttt{split-\codebreak{}selection} & Hold every match of a regexp in the selection, then wait for a verb: i types before every match and a after, c empties them first, d deletes them, w or any delimiter wraps them (again, to take the pair off), and C-g lets go.\\
none & \texttt{M-x split-\codebreak{}selection-\codebreak{}thorough} & \texttt{split-\codebreak{}selection-\codebreak{}thorough} & Hold every match of a regexp in the selection, as f does, finding every match of up to a megabyte (about a million characters of ASCII) whole however many lines it spans, where f may miss one of more than 4 KB that runs across lines and needs text further on; the writing out of sight the same.\\
none & \texttt{C-x 2} & \texttt{split-\codebreak{}window-\codebreak{}below} & Split the selected window into two, one above the other, both showing its buffer from the same line; the upper one stays selected.\\
none & \texttt{C-x 3} & \texttt{split-\codebreak{}window-\codebreak{}right} & Split the selected window into two side by side, both showing its buffer from the same line; the left one stays selected.\\
none & \texttt{M-x split-\codebreak{}wrap} & \texttt{split-\codebreak{}wrap} & Wrap every place of the split in the pair of the delimiter asked for, or take the pair off where it stands outside every place already.\\
none & \texttt{M-x split-\codebreak{}wrap-\codebreak{}key} & \texttt{split-\codebreak{}wrap-\codebreak{}key} & Wrap every place of the split in the pair of the delimiter pressed, or take it off; see split-wrap.\\
none & \texttt{C-x (} & \texttt{start-\codebreak{}recording-\codebreak{}keys} & Begin recording the keys you type, the answers typed to questions too, until C-x ); C-x e then plays them again, as if typed once more: C-x ( C-s date RET C-k C-x ) records a search for date and a kill there, and each play searches for date again and kills from where it found it.\\
\texttt{v} & none & \texttt{start-\codebreak{}selection} & Start a selection at point: any motion extends it.\\
none & \texttt{C-x J} & \texttt{stop-\codebreak{}job} & Stop the work running in the background in the current buffer, after asking.\\
none & \texttt{C-x )} & \texttt{stop-\codebreak{}recording-\codebreak{}keys} & End the recording C-x ( began: the keys typed since, but those of this command, become the keys C-x e plays.\\
none & \texttt{C-x r t} & \texttt{string-\codebreak{}rectangle} & Replace the rectangle between the mark and point with a text read in the minibuffer, on each of its lines; with no width, the text goes in at its column on each line.\\
none & \texttt{C-x w k} & \texttt{swap-\codebreak{}window-\codebreak{}above} & Swap the selected window’s buffer with the window above; you go with your buffer, so it moves up.\\
none & \texttt{C-x w j} & \texttt{swap-\codebreak{}window-\codebreak{}below} & Swap the selected window’s buffer with the window below; you go with your buffer, so it moves down.\\
none & \texttt{C-x w h} & \texttt{swap-\codebreak{}window-\codebreak{}left} & Swap the selected window’s buffer with the window to the left; you go with your buffer, so it moves left.\\
none & \texttt{C-x w l} & \texttt{swap-\codebreak{}window-\codebreak{}right} & Swap the selected window’s buffer with the window to the right; you go with your buffer, so it moves right.\\
none & \texttt{C-x x b} & \texttt{switch-\codebreak{}other-\codebreak{}buffer} & Switch to the buffer used before this one (one no other window shows, when there is one), so that doing it again comes back.\\
none & \texttt{C-x b} & \texttt{switch-\codebreak{}to-\codebreak{}buffer} & Show another buffer, completing on the live buffers, the ones used lately first (M-x minibuffer-finish alone goes back to the one used last); a new name creates one.\\
none & \texttt{C-x 5 b} & \texttt{switch-\codebreak{}to-\codebreak{}buffer-\codebreak{}other-\codebreak{}frame} & Show a buffer in another frame and work in it, completing on the live buffers as C-x b does, the ones used lately first; a new name makes one.\\
none & \texttt{C-x 4 b} & \texttt{switch-\codebreak{}to-\codebreak{}buffer-\codebreak{}other-\codebreak{}window} & Show a buffer in another window and select it, completing on the live buffers as C-x b does, the ones used lately first; a new name makes one.\\
none & \texttt{C-z TAB} & \texttt{tab-\codebreak{}to-\codebreak{}tab-\codebreak{}stop} & Go on to the next tab stop, putting spaces in up to it, or a tab character where the buffer wants tabs (\textbf{indent-tabs-mode}, set-indent-tabs-mode, M-x indent-tabs-mode); with a count, that many stops.\\
none & \texttt{C-c C-v t}, \texttt{C-c C-v C-t} & \texttt{tangle-\codebreak{}code-\codebreak{}files} & Write the code files of this Org file’s source blocks (C-c C-v t), each block with :tangle FILE (on its \#+begin\_src line, in a subtree’s property :header-args: or :header-args:LANGUAGE:, or in \#+PROPERTY: header-args for the file) goes to FILE.\\
none & \texttt{M-x terminal} & \texttt{terminal} & Open a terminal in a buffer of its own, where \textbf{TERMINAL-PLACE} says, as OPEN-TERMINAL does.\\
none & \texttt{M-x terminal-\codebreak{}bottom} & \texttt{terminal-\codebreak{}bottom} & Open a new terminal in a window across the whole bottom, whatever \textbf{TERMINAL-PLACE} says, \textbf{TERMINAL-SIZE} high, as M-x terminal opens one.\\
none & \texttt{RET} & \texttt{terminal-\codebreak{}copy-\codebreak{}and-\codebreak{}type} & In a terminal’s history view, copy what is selected (as M-w) and type into the terminal again; with nothing selected, type into it again.\\
none & \texttt{M-x terminal-\codebreak{}full} & \texttt{terminal-\codebreak{}full} & Open a new terminal in the selected window, whatever \textbf{TERMINAL-PLACE} says, as M-x terminal opens one; quitting it brings back what the window showed.\\
none & \texttt{M-x terminal-\codebreak{}left} & \texttt{terminal-\codebreak{}left} & Open a new terminal in a window down the whole left side, whatever \textbf{TERMINAL-PLACE} says, \textbf{TERMINAL-WIDTH} wide, as M-x terminal opens one.\\
none & \texttt{C-c C-o}, \texttt{C-<down-\codebreak{}mouse-1>} & \texttt{terminal-\codebreak{}open-\codebreak{}link} & Open the link under the mouse (C-click) or at point in a terminal: a link a program wrote (ls --hyperlink, OSC 8) or a web or file address written there.\\
\bottomrule
\end{xltabular}

\begin{xltabular}{\textwidth}{@{}l >{\raggedright\arraybackslash\hsize=0.762\hsize}X l >{\raggedright\arraybackslash\hsize=1.238\hsize}X@{}}
\toprule
\textbf{Modal} & \textbf{Non-modal} & \textbf{Command} & \textbf{What it does}\\
\midrule
\endhead
none & \texttt{p}, \texttt{C-y} & \texttt{terminal-\codebreak{}paste} & Paste into the terminal’s program the text copied last, in Tuatara or another program (as C-y takes it), and type into the terminal: marked as pasted for a program that asks, so that a shell runs none of its lines until you press RET, and its control characters made spaces, so that pasted text cannot work the terminal.\\
none & \texttt{<mouse-2>} & \texttt{terminal-\codebreak{}paste-\codebreak{}primary} & Paste into the terminal’s program the text selected last in any program (the primary selection), as the middle button does, and type into it.\\
none & \texttt{M-x terminal-\codebreak{}right} & \texttt{terminal-\codebreak{}right} & Open a new terminal in a window down the whole right side, whatever \textbf{TERMINAL-PLACE} says, \textbf{TERMINAL-WIDTH} wide, as M-x terminal opens one.\\
none & \texttt{M-x terminal-\codebreak{}search} & \texttt{terminal-\codebreak{}search} & Search the terminal’s whole history, from where its program’s cursor is back (C-r as anywhere): typing there, its history view is shown first and the search starts once it is in; in the history view, search on from point (C-s).\\
none & \texttt{M-x terminal-\codebreak{}send-\codebreak{}escape} & \texttt{terminal-\codebreak{}send-\codebreak{}escape} & Send Escape to the program in this terminal, at a shell’s prompt too, whatever \textbf{ESCAPE-QUITS} and \textbf{TERMINAL-SEND-ESCAPE} say.\\
none & \texttt{C-q} & \texttt{terminal-\codebreak{}send-\codebreak{}key} & Send the next key typed to the program in this terminal, whatever it is: one of those that stay the editor’s (C-g, C-x, M-x, C-q), say.\\
none & \texttt{a}, \texttt{i}, \texttt{C-\textbackslash{}} & \texttt{terminal-\codebreak{}start-\codebreak{}typing} & Type into the terminal again after its history view (M-x terminal-stop-typing): everything typed goes to its program, the screen shown as it is now.\\
none & \texttt{M-x terminal-\codebreak{}stop-\codebreak{}typing} & \texttt{terminal-\codebreak{}stop-\codebreak{}typing} & Stop typing into the terminal: its history view, its whole history (what went up off its screen too) standing still as text with its colours, to move through, search, select and copy while its program goes on; with modal editing on in SUPPORT state (NORMAL’s keys that read and select, and the leader), else with the keys of any text.\\
none & \texttt{M-x terminal-\codebreak{}toggle} & \texttt{terminal-\codebreak{}toggle} & Show the terminal and type into it, or hide it again: the terminal opened last, where \textbf{TERMINAL-PLACE} says (a new one when there is none); selected when it shows already; hidden when you are in it, its window closed (or, at :full, the buffer before it back).\\
none & \texttt{M-x terminal-\codebreak{}top} & \texttt{terminal-\codebreak{}top} & Open a new terminal in a window across the whole top, whatever \textbf{TERMINAL-PLACE} says, \textbf{TERMINAL-SIZE} high, as M-x terminal opens one.\\
none & \texttt{M-x terminal-\codebreak{}version} & \texttt{terminal-\codebreak{}version} & Say which version of the terminal library (libghostty-vt, from Ghostty) M-x terminal runs on.\\
none & \texttt{M-x text-\codebreak{}direction-\codebreak{}mode} & \texttt{text-\codebreak{}direction-\codebreak{}mode} & Show the lines of the current buffer that hold right-to-left text (Hebrew, Arabic and the others) in the order they are read, or each line’s characters in the order they are stored.\\
none & \texttt{C-x C-+}, \texttt{C-x C--}, \texttt{C-x C-0} & \texttt{text-\codebreak{}scale-\codebreak{}adjust} & Make the window’s font larger or smaller, by the last key of the keys that ran it: - a step smaller, 0 the size it started with, any other (+, =) a step larger, 1.2 times; with a prefix argument, that many steps.\\
none & \texttt{M-x text-\codebreak{}scale-\codebreak{}decrease} & \texttt{text-\codebreak{}scale-\codebreak{}decrease} & Make the window’s font a step smaller, or the prefix argument’s worth of steps.\\
none & \texttt{M-x text-\codebreak{}scale-\codebreak{}increase} & \texttt{text-\codebreak{}scale-\codebreak{}increase} & Make the window’s font a step larger (1.2 times), or the prefix argument’s worth of steps.\\
none & \texttt{M-x toc-\codebreak{}org-\codebreak{}insert-\codebreak{}toc} & \texttt{toc-\codebreak{}org-\codebreak{}insert-\codebreak{}toc} & Make the table of contents of the current Org file again, under its first headline tagged :TOC: (as saving does, with \textbf{org-toc-on-save}), a list of links to its headlines, two levels deep unless the tag says how deep (:TOC\_3:), to GitHub’s anchors unless it says :TOC\_2\_org: for their titles; :TOC:QUOTE: puts it in a quote block.\\
\texttt{C} & \texttt{C-z C} & \texttt{toggle-\codebreak{}comment-\codebreak{}lines} & Comment the lines the selection touches, or the current line, or uncomment them when they are all commented: with the language’s marker (;; in Lisp, \emph{/ in Rust, \# in shell scripts and Org) at their smallest indentation, or each line between the two a comment has where it has an end (<!-- --> in HTML, /* *} in CSS).\\
none & \texttt{M-x toggle-\codebreak{}comment-\codebreak{}lines-\codebreak{}ret} & \texttt{toggle-\codebreak{}comment-\codebreak{}lines-\codebreak{}ret} & Comment the lines the region touches, or the line point is on, then going to the next line, or uncomment them when they are all commented; it has no key of its own.\\
none & \texttt{M-x toggle-\codebreak{}debug-\codebreak{}on-\codebreak{}error} & \texttt{toggle-\codebreak{}debug-\codebreak{}on-\codebreak{}error} & Switch debug-on-error on or off: while it is on, an error in a command, or in Lisp evaluated in the editor (C-x C-e, C-M-x, M-:, C-j), shows the calls under way when it was signalled in \textbf{Backtrace}, innermost first, besides the message.\\
none & \texttt{<f11>} & \texttt{toggle-\codebreak{}frame-\codebreak{}fullscreen} & Make the window fill the screen, without its title bar, or go back to how it was.\\
none & \texttt{M-<f10>} & \texttt{toggle-\codebreak{}frame-\codebreak{}maximized} & Make the window as large as the screen allows, or go back to its size before.\\
none & \texttt{C-\textbackslash{}} & \texttt{toggle-\codebreak{}input-\codebreak{}method} & Turn this buffer’s input method off, or on again (C-\textbackslash{}).\\
none & \texttt{M-x toggle-\codebreak{}light-\codebreak{}dark-\codebreak{}theme} & \texttt{toggle-\codebreak{}light-\codebreak{}dark-\codebreak{}theme} & Switch between the light and the dark theme: the two ao variants, unless the init file sets \textbf{light-theme} and \textbf{dark-theme} to others.\\
none & \texttt{M-x toggle-\codebreak{}line-\codebreak{}numbers} & \texttt{toggle-\codebreak{}line-\codebreak{}numbers} & Cycle the gutter of the current buffer through no line numbers, absolute, relative.\\
none & \texttt{C-x T t} & \texttt{toggle-\codebreak{}mode-\codebreak{}line-\codebreak{}tuatara} & Write TUATARA before the modal state in the mode line, TUATARA[N], or the state alone, [N]; \textbf{mode-line-tuatara} says which as Tuatara starts.\\
none & \texttt{C-t} & \texttt{transpose-\codebreak{}chars} & Interchange the characters about point, moving on past them; at the end of a line, the two before point.\\
none & \texttt{C-x C-t} & \texttt{transpose-\codebreak{}lines} & Interchange this line and the one before it, point after them; with a count N, carry the line before point past N lines (back when N is negative); with 0, interchange the lines of point and the mark.\\
none & \texttt{C-M-t} & \texttt{transpose-\codebreak{}sexps} & Interchange the balanced expressions about point, point after them; with a count N, carry the expression before point past N expressions (back when N is negative); with 0, interchange the expressions at point and at the mark.\\
none & \texttt{C-x w t} & \texttt{transpose-\codebreak{}window-\codebreak{}layout} & Make windows one above the other stand side by side, and windows side by side one above the other, each keeping its place in the order; again turns them back.\\
none & \texttt{M-t} & \texttt{transpose-\codebreak{}words} & Interchange the words about point, point after them; with a count N, carry the word before or around point past N words (back when N is negative); with 0, interchange the words at point and at the mark.\\
none & \texttt{M-x trust-\codebreak{}project} & \texttt{trust-\codebreak{}project} & Ask again whether you trust the project of the current buffer (the git repository its file or folder is in, else its project, or the repository a git buffer shows), whatever was answered before.\\
\texttt{c} & none & \texttt{tuatara-\codebreak{}change} & Delete the selection (into the kill ring) or the character at point (not saved) and enter INSERT state.\\
\texttt{i} & none & \texttt{tuatara-\codebreak{}insert} & Insert at point, and enter INSERT state.\\
\texttt{a} & none & \texttt{tuatara-\codebreak{}insert-\codebreak{}after} & Insert after the character at point, and enter INSERT state.\\
\texttt{A} & none & \texttt{tuatara-\codebreak{}insert-\codebreak{}line-\codebreak{}end} & Enter INSERT state at the end of the line.\\
\texttt{I} & none & \texttt{tuatara-\codebreak{}insert-\codebreak{}line-\codebreak{}start} & Enter INSERT state at the start of the line.\\
\texttt{C-g} & none & \texttt{tuatara-\codebreak{}normal} & Return to NORMAL state; the selection is dropped, a list of completions beside the cursor closes, and a line begun and left blank loses the blanks it was indented with.\\
\texttt{RET} & none & \texttt{tuatara-\codebreak{}normal-\codebreak{}return} & RET in NORMAL state: what the buffer’s mode does with it (\textbf{enter-dwim-functions}), as following the Org link point is on; else, and always in the hex view, nothing, so that it changes no text.\\
\texttt{TAB} & none & \texttt{tuatara-\codebreak{}normal-\codebreak{}tab} & TAB in NORMAL state: in a buffer whose major mode is one of \textbf{tuatara-normal-tab-modes} (Org), or derives from one, what TAB does in that mode, as in INSERT state; in the hex view, from its hex digits to its text or back (hex-switch-column); elsewhere nothing, so that it changes no text.\\
\texttt{O} & none & \texttt{tuatara-\codebreak{}open-\codebreak{}above} & Open a line above the current one and enter INSERT state on it, indented as the language’s mode indents it (see \textbf{INDENT-NEW-LINES}).\\
\bottomrule
\end{xltabular}

\begin{xltabular}{\textwidth}{@{}l >{\raggedright\arraybackslash\hsize=0.619\hsize}X l >{\raggedright\arraybackslash\hsize=1.381\hsize}X@{}}
\toprule
\textbf{Modal} & \textbf{Non-modal} & \textbf{Command} & \textbf{What it does}\\
\midrule
\endhead
\texttt{o} & none & \texttt{tuatara-\codebreak{}open-\codebreak{}below} & Open a line below the current one and enter INSERT state on it, indented as the language’s mode indents it.\\
none & \texttt{C-g} & \texttt{tuatara-\codebreak{}quit} & C-g: in INSERT state return to NORMAL; in NORMAL drop the selection.\\
none & \texttt{M-x tuatara-\codebreak{}tutor} & \texttt{tuatara-\codebreak{}tutor} & Open the tutor: an ordinary buffer holding its own lessons, to learn the modal keys by editing it as it says, with every key the one you have.\\
none & \texttt{M-x tuatara-\codebreak{}version} & \texttt{tuatara-\codebreak{}version} & Say which Tuatara this is: its version, the commit it was built from and when, the Steel Bank Common Lisp (SBCL) it runs on, and where it shows (a window, the terminal).\\
none & \texttt{M-x turn-\codebreak{}off-\codebreak{}grammar} & \texttt{turn-\codebreak{}off-\codebreak{}grammar} & Turn off the grammar of a language, asked for (the current buffer’s first): its files are no longer coloured or parsed, nor text in it inside other files, in later sessions too, and no file asks to turn it on again; M-x turn-on-grammar turns it on again.\\
none & \texttt{M-x turn-\codebreak{}on-\codebreak{}grammar} & \texttt{turn-\codebreak{}on-\codebreak{}grammar} & Turn on the grammar of a language, asked for (the current buffer’s first when its grammar is off): its files are coloured and parsed from now on, and text in it inside other files, in later sessions too.\\
\texttt{u} & \texttt{C-/}, \texttt{C-\_}, \texttt{C-x u} & \texttt{undo} & Undo the last change.\\
none & \texttt{C-u} & \texttt{universal-\codebreak{}argument} & Begin a numeric prefix argument: C-u alone means 4, C-u C-u 16, digits after it give their number, a leading minus negates.\\
none & \texttt{M-x up-\codebreak{}list} & \texttt{up-\codebreak{}list} & Move forward out of the bracketed group around point, past its closing bracket; out of a string first to its end.\\
none & \texttt{M-x upcase-\codebreak{}dwim} & \texttt{upcase-\codebreak{}dwim} & Put the region in upper case, or with none the next word, as upcase-word does.\\
none & \texttt{C-x C-u} & \texttt{upcase-\codebreak{}region} & Put the region in upper case.\\
none & \texttt{M-u} & \texttt{upcase-\codebreak{}word} & Put the next word, or the count’s, in upper case and move over it; with a negative count, the words before point, point staying.\\
none & \texttt{M-x vc-\codebreak{}gutter-\codebreak{}mode} & \texttt{vc-\codebreak{}gutter-\codebreak{}mode} & Switch the version-control gutter on or off for all buffers; with a count, on when it is positive, off when it is zero or less.\\
none & \texttt{M-x vc-\codebreak{}gutter-\codebreak{}refresh} & \texttt{vc-\codebreak{}gutter-\codebreak{}refresh} & Recompute the version-control marks of the current buffer from git.\\
none & \texttt{C-h E} & \texttt{view-\codebreak{}debug-\codebreak{}buffer} & Show \textbf{Debug}, the reports of debug mode (M-x debug-mode), where \textbf{Messages} is shown, its end in sight.\\
none & \texttt{C-h e} & \texttt{view-\codebreak{}echo-\codebreak{}area-\codebreak{}messages} & Show the \textbf{Messages} buffer: the messages shown in the echo area so far, repeats counted, the latest at the end.\\
none & \texttt{C-h N} & \texttt{view-\codebreak{}notifications} & Show \textbf{Notifications}, what was told so far this session, the agenda’s entries tagged :notify: as their time came and what packages told, a line each with the date and time it was told and who told it, the latest at the end.\\
none & \texttt{C-x T v} & \texttt{visual-\codebreak{}line-\codebreak{}mode} & Mark the current buffer’s fill column (C-x f) with a line down the window, or stop: the last column a filled line may reach, the 80th for a fill column of 80, painted in the :visual-line face over the text’s backgrounds (a source block’s, the selection’s) and under the line point is on.\\
none & \texttt{C-x =} & \texttt{what-\codebreak{}cursor-\codebreak{}position} & Say in the echo area which character is at point, with its code in decimal, octal and hexadecimal, and where point is: its place counted in characters from 1, of how many there are and as a share of them, its octet and its column.\\
none & \texttt{C-x x w} & \texttt{widen} & Show all of the buffer again after C-x x n or C-x x s, every command reaching all of it, point staying where it is.\\
none & \texttt{C-x w x} & \texttt{window-\codebreak{}swap-\codebreak{}states} & Swap what the selected window shows with the next window (as C-x o would select it); you go with your buffer.\\
none & \texttt{M-x wipe-\codebreak{}file} & \texttt{wipe-\codebreak{}file} & Wipe the file shown (in a list of a directory, the one on the line; with C-u, or where there is none, asked for) after asking: its contents overwritten with random octets, put on disk, then the file renamed and deleted, in the background, so that it is not found again where the disk writes in place.\\
\texttt{"}, \texttt{\$}, \texttt{'} & \texttt{C-z "}, \texttt{C-z \$}, \texttt{C-z '} & \texttt{wrap-\codebreak{}selection-\codebreak{}with-\codebreak{}key} & Wrap the selection in the delimiter pressed and its partner, or take them off when they stand just outside it: ” on a word quotes it, and ” on the word of a quoted word unquotes it.\\
none & \texttt{C-x C-w} & \texttt{write-\codebreak{}file} & Write the current buffer to another file, asked for as C-x C-f asks (C-x C-w), from the folder of the path chosen last in a question for a path, else the buffer’s own folder.\\
none & \texttt{M-.}, \texttt{RET}, \texttt{C-c d} & \texttt{xref-\codebreak{}find-\codebreak{}definitions} & Go to the definition of the name at point.\\
none & \texttt{M-,}, \texttt{C-c D} & \texttt{xref-\codebreak{}go-\codebreak{}back} & Go back to where M-.\\
none & \texttt{C-M-,} & \texttt{xref-\codebreak{}go-\codebreak{}forward} & Go forward again to where M-, came back from.\\
none & \texttt{C-y} & \texttt{yank} & Insert the most recently killed text at point, or what another program copied to the system clipboard since, its lines ending as the buffer’s do.\\
\texttt{p} & none & \texttt{yank-\codebreak{}over-\codebreak{}selection} & Paste the last kill COUNT times, replacing the selection if there is one, its lines ending as the buffer’s do; what another program copied to the system clipboard since counts as the last kill.\\
\texttt{P} & \texttt{C-x r y} & \texttt{yank-\codebreak{}rectangle} & Paste the rectangle store (what y, d or c last took from a rectangle) as a block of columns: from point, one row a line, or over a selected rectangle of as many rows in its place.\\
none & \texttt{M-x zap-\codebreak{}to-\codebreak{}char} & \texttt{zap-\codebreak{}to-\codebreak{}char} & Kill up to and including the next character typed after it; with a count N, to the Nth, back when N is negative.\\
none & \texttt{M-z} & \texttt{zap-\codebreak{}up-\codebreak{}to-\codebreak{}char} & Kill up to, but not including, the next character typed after it; with a count N, to the Nth, back when N is negative.\\
\bottomrule
\end{xltabular}

\subsection{Commands that switch to a major mode}
\label{sec:576}
A major mode sets how Tuatara treats one kind of text, as chapter 13
describes. Each major mode has a command named after it that switches the
current buffer to that mode. Most have no key, because Tuatara picks the mode
from the file’s name and contents when it opens the file. The commands are
listed here, apart from the rest, so that the table above can be read through.

\begin{xltabular}{\textwidth}{@{}l >{\raggedright\arraybackslash\hsize=0.875\hsize}X l >{\raggedright\arraybackslash\hsize=1.125\hsize}X@{}}
\toprule
\textbf{Modal} & \textbf{Non-modal} & \textbf{Command} & \textbf{What it does}\\
\midrule
\endhead
none & \texttt{M-x ada-\codebreak{}mode} & \texttt{ada-\codebreak{}mode} & Switch the current buffer to ada-mode.\\
none & \texttt{M-x adl-\codebreak{}mode} & \texttt{adl-\codebreak{}mode} & Switch the current buffer to adl-mode.\\
none & \texttt{M-x agda2-\codebreak{}mode} & \texttt{agda2-\codebreak{}mode} & Switch the current buffer to agda2-mode.\\
none & \texttt{M-x alloy-\codebreak{}mode} & \texttt{alloy-\codebreak{}mode} & Switch the current buffer to alloy-mode.\\
none & \texttt{M-x amber-\codebreak{}mode} & \texttt{amber-\codebreak{}mode} & Switch the current buffer to amber-mode.\\
none & \texttt{M-x asm-\codebreak{}mode} & \texttt{asm-\codebreak{}mode} & Switch the current buffer to asm-mode.\\
none & \texttt{M-x astro-\codebreak{}mode} & \texttt{astro-\codebreak{}mode} & Switch the current buffer to astro-mode.\\
none & \texttt{M-x awk-\codebreak{}mode} & \texttt{awk-\codebreak{}mode} & Switch the current buffer to awk-mode.\\
none & \texttt{M-x bash-\codebreak{}mode} & \texttt{bash-\codebreak{}mode} & Switch the current buffer to bash-mode.\\
none & \texttt{M-x basic-\codebreak{}mode} & \texttt{basic-\codebreak{}mode} & Switch the current buffer to basic-mode.\\
none & \texttt{M-x bass-\codebreak{}mode} & \texttt{bass-\codebreak{}mode} & Switch the current buffer to bass-mode.\\
none & \texttt{M-x bazel-\codebreak{}starlark-\codebreak{}mode} & \texttt{bazel-\codebreak{}starlark-\codebreak{}mode} & Switch the current buffer to bazel-starlark-mode.\\
none & \texttt{M-x beancount-\codebreak{}mode} & \texttt{beancount-\codebreak{}mode} & Switch the current buffer to beancount-mode.\\
none & \texttt{M-x bibtex-\codebreak{}mode} & \texttt{bibtex-\codebreak{}mode} & Switch the current buffer to bibtex-mode.\\
none & \texttt{M-x bicep-\codebreak{}mode} & \texttt{bicep-\codebreak{}mode} & Switch the current buffer to bicep-mode.\\
none & \texttt{M-x bitbake-\codebreak{}mode} & \texttt{bitbake-\codebreak{}mode} & Switch the current buffer to bitbake-mode.\\
none & \texttt{M-x blade-\codebreak{}mode} & \texttt{blade-\codebreak{}mode} & Switch the current buffer to blade-mode.\\
none & \texttt{M-x blueprint-\codebreak{}mode} & \texttt{blueprint-\codebreak{}mode} & Switch the current buffer to blueprint-mode.\\
none & \texttt{M-x c++-\codebreak{}mode} & \texttt{c++-\codebreak{}mode} & Switch the current buffer to c++-mode.\\
none & \texttt{M-x c-mode} & \texttt{c-mode} & Switch the current buffer to c-mode.\\
none & \texttt{M-x c3-\codebreak{}mode} & \texttt{c3-\codebreak{}mode} & Switch the current buffer to c3-mode.\\
none & \texttt{M-x caddyfile-\codebreak{}mode} & \texttt{caddyfile-\codebreak{}mode} & Switch the current buffer to caddyfile-mode.\\
none & \texttt{M-x cairo-\codebreak{}mode} & \texttt{cairo-\codebreak{}mode} & Switch the current buffer to cairo-mode.\\
none & \texttt{M-x capnp-\codebreak{}mode} & \texttt{capnp-\codebreak{}mode} & Switch the current buffer to capnp-mode.\\
none & \texttt{M-x cel-\codebreak{}mode} & \texttt{cel-\codebreak{}mode} & Switch the current buffer to cel-mode.\\
none & \texttt{M-x chuck-\codebreak{}mode} & \texttt{chuck-\codebreak{}mode} & Switch the current buffer to chuck-mode.\\
none & \texttt{M-x circom-\codebreak{}mode} & \texttt{circom-\codebreak{}mode} & Switch the current buffer to circom-mode.\\
none & \texttt{M-x clarity-\codebreak{}mode} & \texttt{clarity-\codebreak{}mode} & Switch the current buffer to clarity-mode.\\
none & \texttt{M-x clojure-\codebreak{}mode} & \texttt{clojure-\codebreak{}mode} & Switch the current buffer to clojure-mode.\\
none & \texttt{M-x cmake-\codebreak{}mode} & \texttt{cmake-\codebreak{}mode} & Switch the current buffer to cmake-mode.\\
none & \texttt{M-x concerto-\codebreak{}mode} & \texttt{concerto-\codebreak{}mode} & Switch the current buffer to concerto-mode.\\
none & \texttt{M-x conf-\codebreak{}javaprop-\codebreak{}mode} & \texttt{conf-\codebreak{}javaprop-\codebreak{}mode} & Switch the current buffer to conf-javaprop-mode.\\
none & \texttt{M-x containerfile-\codebreak{}mode} & \texttt{containerfile-\codebreak{}mode} & Switch the current buffer to containerfile-mode.\\
none & \texttt{M-x cpon-\codebreak{}mode} & \texttt{cpon-\codebreak{}mode} & Switch the current buffer to cpon-mode.\\
none & \texttt{M-x crystal-\codebreak{}mode} & \texttt{crystal-\codebreak{}mode} & Switch the current buffer to crystal-mode.\\
none & \texttt{M-x csharp-\codebreak{}mode} & \texttt{csharp-\codebreak{}mode} & Switch the current buffer to csharp-mode.\\
none & \texttt{M-x css-\codebreak{}mode} & \texttt{css-\codebreak{}mode} & Switch the current buffer to css-mode.\\
none & \texttt{M-x csv-\codebreak{}mode} & \texttt{csv-\codebreak{}mode} & Switch the current buffer to csv-mode.\\
none & \texttt{M-x cue-\codebreak{}mode} & \texttt{cue-\codebreak{}mode} & Switch the current buffer to cue-mode.\\
none & \texttt{M-x cylc-\codebreak{}mode} & \texttt{cylc-\codebreak{}mode} & Switch the current buffer to cylc-mode.\\
\bottomrule
\end{xltabular}

\begin{xltabular}{\textwidth}{@{}l >{\raggedright\arraybackslash\hsize=0.818\hsize}X l >{\raggedright\arraybackslash\hsize=1.182\hsize}X@{}}
\toprule
\textbf{Modal} & \textbf{Non-modal} & \textbf{Command} & \textbf{What it does}\\
\midrule
\endhead
none & \texttt{M-x cython-\codebreak{}mode} & \texttt{cython-\codebreak{}mode} & Switch the current buffer to cython-mode.\\
none & \texttt{M-x d-mode} & \texttt{d-mode} & Switch the current buffer to d-mode.\\
none & \texttt{M-x dart-\codebreak{}mode} & \texttt{dart-\codebreak{}mode} & Switch the current buffer to dart-mode.\\
none & \texttt{M-x dbml-\codebreak{}mode} & \texttt{dbml-\codebreak{}mode} & Switch the current buffer to dbml-mode.\\
none & \texttt{M-x debian-\codebreak{}control-\codebreak{}mode} & \texttt{debian-\codebreak{}control-\codebreak{}mode} & Switch the current buffer to debian-control-mode.\\
none & \texttt{M-x dhall-\codebreak{}mode} & \texttt{dhall-\codebreak{}mode} & Switch the current buffer to dhall-mode.\\
none & \texttt{M-x diff-\codebreak{}mode} & \texttt{diff-\codebreak{}mode} & Switch the current buffer to diff-mode.\\
none & \texttt{M-x djot-\codebreak{}mode} & \texttt{djot-\codebreak{}mode} & Switch the current buffer to djot-mode.\\
none & \texttt{M-x docker-\codebreak{}compose-\codebreak{}mode} & \texttt{docker-\codebreak{}compose-\codebreak{}mode} & Switch the current buffer to docker-compose-mode.\\
none & \texttt{M-x dotenv-\codebreak{}mode} & \texttt{dotenv-\codebreak{}mode} & Switch the current buffer to dotenv-mode.\\
none & \texttt{M-x doxyfile-\codebreak{}mode} & \texttt{doxyfile-\codebreak{}mode} & Switch the current buffer to doxyfile-mode.\\
none & \texttt{M-x dtd-\codebreak{}mode} & \texttt{dtd-\codebreak{}mode} & Switch the current buffer to dtd-mode.\\
none & \texttt{M-x dts-\codebreak{}mode} & \texttt{dts-\codebreak{}mode} & Switch the current buffer to dts-mode.\\
none & \texttt{M-x dune-\codebreak{}mode} & \texttt{dune-\codebreak{}mode} & Switch the current buffer to dune-mode.\\
none & \texttt{M-x dunstrc-\codebreak{}mode} & \texttt{dunstrc-\codebreak{}mode} & Switch the current buffer to dunstrc-mode.\\
none & \texttt{M-x earthfile-\codebreak{}mode} & \texttt{earthfile-\codebreak{}mode} & Switch the current buffer to earthfile-mode.\\
none & \texttt{M-x ebnf-\codebreak{}mode} & \texttt{ebnf-\codebreak{}mode} & Switch the current buffer to ebnf-mode.\\
none & \texttt{M-x eex-\codebreak{}mode} & \texttt{eex-\codebreak{}mode} & Switch the current buffer to eex-mode.\\
none & \texttt{M-x eiffel-\codebreak{}mode} & \texttt{eiffel-\codebreak{}mode} & Switch the current buffer to eiffel-mode.\\
none & \texttt{M-x ejs-\codebreak{}mode} & \texttt{ejs-\codebreak{}mode} & Switch the current buffer to ejs-mode.\\
none & \texttt{M-x elixir-\codebreak{}mode} & \texttt{elixir-\codebreak{}mode} & Switch the current buffer to elixir-mode.\\
none & \texttt{M-x elm-\codebreak{}mode} & \texttt{elm-\codebreak{}mode} & Switch the current buffer to elm-mode.\\
none & \texttt{M-x elvish-\codebreak{}mode} & \texttt{elvish-\codebreak{}mode} & Switch the current buffer to elvish-mode.\\
none & \texttt{M-x emacs-\codebreak{}lisp-\codebreak{}mode} & \texttt{emacs-\codebreak{}lisp-\codebreak{}mode} & Switch the current buffer to emacs-lisp-mode.\\
none & \texttt{M-x eml-\codebreak{}mode} & \texttt{eml-\codebreak{}mode} & Switch the current buffer to eml-mode.\\
none & \texttt{M-x erb-\codebreak{}mode} & \texttt{erb-\codebreak{}mode} & Switch the current buffer to erb-mode.\\
none & \texttt{M-x erlang-\codebreak{}mode} & \texttt{erlang-\codebreak{}mode} & Switch the current buffer to erlang-mode.\\
none & \texttt{M-x esdl-\codebreak{}mode} & \texttt{esdl-\codebreak{}mode} & Switch the current buffer to esdl-mode.\\
none & \texttt{M-x etc-\codebreak{}passwd-\codebreak{}generic-\codebreak{}mode} & \texttt{etc-\codebreak{}passwd-\codebreak{}generic-\codebreak{}mode} & Switch the current buffer to etc-passwd-generic-mode.\\
none & \texttt{M-x feature-\codebreak{}mode} & \texttt{feature-\codebreak{}mode} & Switch the current buffer to feature-mode.\\
none & \texttt{M-x fennel-\codebreak{}mode} & \texttt{fennel-\codebreak{}mode} & Switch the current buffer to fennel-mode.\\
none & \texttt{M-x fga-\codebreak{}mode} & \texttt{fga-\codebreak{}mode} & Switch the current buffer to fga-mode.\\
none & \texttt{M-x fidl-\codebreak{}mode} & \texttt{fidl-\codebreak{}mode} & Switch the current buffer to fidl-mode.\\
none & \texttt{M-x fish-\codebreak{}mode} & \texttt{fish-\codebreak{}mode} & Switch the current buffer to fish-mode.\\
none & \texttt{M-x flatbuffers-\codebreak{}mode} & \texttt{flatbuffers-\codebreak{}mode} & Switch the current buffer to flatbuffers-mode.\\
none & \texttt{M-x forth-\codebreak{}mode} & \texttt{forth-\codebreak{}mode} & Switch the current buffer to forth-mode.\\
none & \texttt{M-x fortran-\codebreak{}mode} & \texttt{fortran-\codebreak{}mode} & Switch the current buffer to fortran-mode.\\
none & \texttt{M-x freebasic-\codebreak{}mode} & \texttt{freebasic-\codebreak{}mode} & Switch the current buffer to freebasic-mode.\\
none & \texttt{M-x fsharp-\codebreak{}mode} & \texttt{fsharp-\codebreak{}mode} & Switch the current buffer to fsharp-mode.\\
none & \texttt{M-x fundamental-\codebreak{}mode} & \texttt{fundamental-\codebreak{}mode} & Switch the current buffer to fundamental-mode.\\
\bottomrule
\end{xltabular}

\begin{xltabular}{\textwidth}{@{}l >{\raggedright\arraybackslash\hsize=0.875\hsize}X l >{\raggedright\arraybackslash\hsize=1.125\hsize}X@{}}
\toprule
\textbf{Modal} & \textbf{Non-modal} & \textbf{Command} & \textbf{What it does}\\
\midrule
\endhead
none & \texttt{M-x gdscript-\codebreak{}mode} & \texttt{gdscript-\codebreak{}mode} & Switch the current buffer to gdscript-mode.\\
none & \texttt{M-x gemini-\codebreak{}mode} & \texttt{gemini-\codebreak{}mode} & Switch the current buffer to gemini-mode.\\
none & \texttt{M-x ghostty-\codebreak{}mode} & \texttt{ghostty-\codebreak{}mode} & Switch the current buffer to ghostty-mode.\\
none & \texttt{M-x git-\codebreak{}commit-\codebreak{}mode} & \texttt{git-\codebreak{}commit-\codebreak{}mode} & Switch the current buffer to git-commit-mode.\\
none & \texttt{M-x git-\codebreak{}rebase-\codebreak{}mode} & \texttt{git-\codebreak{}rebase-\codebreak{}mode} & Switch the current buffer to git-rebase-mode.\\
none & \texttt{M-x gitattributes-\codebreak{}mode} & \texttt{gitattributes-\codebreak{}mode} & Switch the current buffer to gitattributes-mode.\\
none & \texttt{M-x gitconfig-\codebreak{}mode} & \texttt{gitconfig-\codebreak{}mode} & Switch the current buffer to gitconfig-mode.\\
none & \texttt{M-x github-\codebreak{}action-\codebreak{}mode} & \texttt{github-\codebreak{}action-\codebreak{}mode} & Switch the current buffer to github-action-mode.\\
none & \texttt{M-x gitignore-\codebreak{}mode} & \texttt{gitignore-\codebreak{}mode} & Switch the current buffer to gitignore-mode.\\
none & \texttt{M-x gitlab-\codebreak{}ci-\codebreak{}mode} & \texttt{gitlab-\codebreak{}ci-\codebreak{}mode} & Switch the current buffer to gitlab-ci-mode.\\
none & \texttt{M-x gleam-\codebreak{}mode} & \texttt{gleam-\codebreak{}mode} & Switch the current buffer to gleam-mode.\\
none & \texttt{M-x glimmer-\codebreak{}js-\codebreak{}mode} & \texttt{glimmer-\codebreak{}js-\codebreak{}mode} & Switch the current buffer to glimmer-js-mode.\\
none & \texttt{M-x glimmer-\codebreak{}mode} & \texttt{glimmer-\codebreak{}mode} & Switch the current buffer to glimmer-mode.\\
none & \texttt{M-x glimmer-\codebreak{}ts-\codebreak{}mode} & \texttt{glimmer-\codebreak{}ts-\codebreak{}mode} & Switch the current buffer to glimmer-ts-mode.\\
none & \texttt{M-x glsl-\codebreak{}mode} & \texttt{glsl-\codebreak{}mode} & Switch the current buffer to glsl-mode.\\
none & \texttt{M-x gn-\codebreak{}mode} & \texttt{gn-\codebreak{}mode} & Switch the current buffer to gn-mode.\\
none & \texttt{M-x gnuplot-\codebreak{}mode} & \texttt{gnuplot-\codebreak{}mode} & Switch the current buffer to gnuplot-mode.\\
none & \texttt{M-x go-\codebreak{}mod-\codebreak{}mode} & \texttt{go-\codebreak{}mod-\codebreak{}mode} & Switch the current buffer to go-mod-mode.\\
none & \texttt{M-x go-\codebreak{}mode} & \texttt{go-\codebreak{}mode} & Switch the current buffer to go-mode.\\
none & \texttt{M-x go-\codebreak{}template-\codebreak{}mode} & \texttt{go-\codebreak{}template-\codebreak{}mode} & Switch the current buffer to go-template-mode.\\
none & \texttt{M-x go-\codebreak{}work-\codebreak{}mode} & \texttt{go-\codebreak{}work-\codebreak{}mode} & Switch the current buffer to go-work-mode.\\
none & \texttt{M-x godot-\codebreak{}resource-\codebreak{}mode} & \texttt{godot-\codebreak{}resource-\codebreak{}mode} & Switch the current buffer to godot-resource-mode.\\
none & \texttt{M-x gpr-\codebreak{}mode} & \texttt{gpr-\codebreak{}mode} & Switch the current buffer to gpr-mode.\\
none & \texttt{M-x graphql-\codebreak{}mode} & \texttt{graphql-\codebreak{}mode} & Switch the current buffer to graphql-mode.\\
none & \texttt{M-x graphviz-\codebreak{}dot-\codebreak{}mode} & \texttt{graphviz-\codebreak{}dot-\codebreak{}mode} & Switch the current buffer to graphviz-dot-mode.\\
none & \texttt{M-x gren-\codebreak{}mode} & \texttt{gren-\codebreak{}mode} & Switch the current buffer to gren-mode.\\
none & \texttt{M-x groovy-\codebreak{}mode} & \texttt{groovy-\codebreak{}mode} & Switch the current buffer to groovy-mode.\\
none & \texttt{M-x hare-\codebreak{}mode} & \texttt{hare-\codebreak{}mode} & Switch the current buffer to hare-mode.\\
none & \texttt{M-x haskell-\codebreak{}mode} & \texttt{haskell-\codebreak{}mode} & Switch the current buffer to haskell-mode.\\
none & \texttt{M-x haxe-\codebreak{}mode} & \texttt{haxe-\codebreak{}mode} & Switch the current buffer to haxe-mode.\\
none & \texttt{M-x hcl-\codebreak{}mode} & \texttt{hcl-\codebreak{}mode} & Switch the current buffer to hcl-mode.\\
none & \texttt{M-x hdl-\codebreak{}mode} & \texttt{hdl-\codebreak{}mode} & Switch the current buffer to hdl-mode.\\
none & \texttt{M-x heex-\codebreak{}mode} & \texttt{heex-\codebreak{}mode} & Switch the current buffer to heex-mode.\\
none & \texttt{M-x hocon-\codebreak{}mode} & \texttt{hocon-\codebreak{}mode} & Switch the current buffer to hocon-mode.\\
none & \texttt{M-x hoon-\codebreak{}mode} & \texttt{hoon-\codebreak{}mode} & Switch the current buffer to hoon-mode.\\
none & \texttt{M-x hosts-\codebreak{}generic-\codebreak{}mode} & \texttt{hosts-\codebreak{}generic-\codebreak{}mode} & Switch the current buffer to hosts-generic-mode.\\
none & \texttt{M-x html-\codebreak{}mode} & \texttt{html-\codebreak{}mode} & Switch the current buffer to html-mode.\\
none & \texttt{M-x hurl-\codebreak{}mode} & \texttt{hurl-\codebreak{}mode} & Switch the current buffer to hurl-mode.\\
none & \texttt{M-x hy-\codebreak{}mode} & \texttt{hy-\codebreak{}mode} & Switch the current buffer to hy-mode.\\
none & \texttt{M-x hyprlang-\codebreak{}mode} & \texttt{hyprlang-\codebreak{}mode} & Switch the current buffer to hyprlang-mode.\\
\bottomrule
\end{xltabular}

\begin{xltabular}{\textwidth}{@{}l >{\raggedright\arraybackslash\hsize=0.783\hsize}X l >{\raggedright\arraybackslash\hsize=1.217\hsize}X@{}}
\toprule
\textbf{Modal} & \textbf{Non-modal} & \textbf{Command} & \textbf{What it does}\\
\midrule
\endhead
none & \texttt{M-x ini-\codebreak{}mode} & \texttt{ini-\codebreak{}mode} & Switch the current buffer to ini-mode.\\
none & \texttt{M-x ink-\codebreak{}mode} & \texttt{ink-\codebreak{}mode} & Switch the current buffer to ink-mode.\\
none & \texttt{M-x inko-\codebreak{}mode} & \texttt{inko-\codebreak{}mode} & Switch the current buffer to inko-mode.\\
none & \texttt{M-x janet-\codebreak{}mode} & \texttt{janet-\codebreak{}mode} & Switch the current buffer to janet-mode.\\
none & \texttt{M-x java-\codebreak{}mode} & \texttt{java-\codebreak{}mode} & Switch the current buffer to java-mode.\\
none & \texttt{M-x jinja2-\codebreak{}mode} & \texttt{jinja2-\codebreak{}mode} & Switch the current buffer to jinja2-mode.\\
none & \texttt{M-x jj-\codebreak{}description-\codebreak{}mode} & \texttt{jj-\codebreak{}description-\codebreak{}mode} & Switch the current buffer to jj-description-mode.\\
none & \texttt{M-x jq-\codebreak{}mode} & \texttt{jq-\codebreak{}mode} & Switch the current buffer to jq-mode.\\
none & \texttt{M-x js-\codebreak{}jsx-\codebreak{}mode} & \texttt{js-\codebreak{}jsx-\codebreak{}mode} & Switch the current buffer to js-jsx-mode.\\
none & \texttt{M-x js-\codebreak{}mode} & \texttt{js-\codebreak{}mode} & Switch the current buffer to js-mode.\\
none & \texttt{M-x json-\codebreak{}mode} & \texttt{json-\codebreak{}mode} & Switch the current buffer to json-mode.\\
none & \texttt{M-x json5-\codebreak{}mode} & \texttt{json5-\codebreak{}mode} & Switch the current buffer to json5-mode.\\
none & \texttt{M-x jsonc-\codebreak{}mode} & \texttt{jsonc-\codebreak{}mode} & Switch the current buffer to jsonc-mode.\\
none & \texttt{M-x jsonl-\codebreak{}mode} & \texttt{jsonl-\codebreak{}mode} & Switch the current buffer to jsonl-mode.\\
none & \texttt{M-x jsonnet-\codebreak{}mode} & \texttt{jsonnet-\codebreak{}mode} & Switch the current buffer to jsonnet-mode.\\
none & \texttt{M-x julia-\codebreak{}mode} & \texttt{julia-\codebreak{}mode} & Switch the current buffer to julia-mode.\\
none & \texttt{M-x just-\codebreak{}mode} & \texttt{just-\codebreak{}mode} & Switch the current buffer to just-mode.\\
none & \texttt{M-x kcl-\codebreak{}mode} & \texttt{kcl-\codebreak{}mode} & Switch the current buffer to kcl-mode.\\
none & \texttt{M-x kconfig-\codebreak{}mode} & \texttt{kconfig-\codebreak{}mode} & Switch the current buffer to kconfig-mode.\\
none & \texttt{M-x kdl-\codebreak{}mode} & \texttt{kdl-\codebreak{}mode} & Switch the current buffer to kdl-mode.\\
none & \texttt{M-x klog-\codebreak{}mode} & \texttt{klog-\codebreak{}mode} & Switch the current buffer to klog-mode.\\
none & \texttt{M-x koka-\codebreak{}mode} & \texttt{koka-\codebreak{}mode} & Switch the current buffer to koka-mode.\\
none & \texttt{M-x kotlin-\codebreak{}mode} & \texttt{kotlin-\codebreak{}mode} & Switch the current buffer to kotlin-mode.\\
none & \texttt{M-x koto-\codebreak{}mode} & \texttt{koto-\codebreak{}mode} & Switch the current buffer to koto-mode.\\
none & \texttt{M-x latex-\codebreak{}mode} & \texttt{latex-\codebreak{}mode} & Switch the current buffer to latex-mode.\\
none & \texttt{M-x ld-\codebreak{}script-\codebreak{}mode} & \texttt{ld-\codebreak{}script-\codebreak{}mode} & Switch the current buffer to ld-script-mode.\\
none & \texttt{M-x ldif-\codebreak{}mode} & \texttt{ldif-\codebreak{}mode} & Switch the current buffer to ldif-mode.\\
none & \texttt{M-x lean4-\codebreak{}mode} & \texttt{lean4-\codebreak{}mode} & Switch the current buffer to lean4-mode.\\
none & \texttt{M-x ledger-\codebreak{}mode} & \texttt{ledger-\codebreak{}mode} & Switch the current buffer to ledger-mode.\\
none & \texttt{M-x less-\codebreak{}css-\codebreak{}mode} & \texttt{less-\codebreak{}css-\codebreak{}mode} & Switch the current buffer to less-css-mode.\\
none & \texttt{M-x lisp-\codebreak{}interaction-\codebreak{}mode} & \texttt{lisp-\codebreak{}interaction-\codebreak{}mode} & Switch the current buffer to lisp-interaction-mode.\\
none & \texttt{M-x lisp-\codebreak{}mode} & \texttt{lisp-\codebreak{}mode} & Switch the current buffer to lisp-mode.\\
none & \texttt{M-x literate-\codebreak{}haskell-\codebreak{}mode} & \texttt{literate-\codebreak{}haskell-\codebreak{}mode} & Switch the current buffer to literate-haskell-mode.\\
none & \texttt{M-x llvm-\codebreak{}mode} & \texttt{llvm-\codebreak{}mode} & Switch the current buffer to llvm-mode.\\
none & \texttt{M-x log-\codebreak{}mode} & \texttt{log-\codebreak{}mode} & Switch the current buffer to log-mode.\\
none & \texttt{M-x lpf-\codebreak{}mode} & \texttt{lpf-\codebreak{}mode} & Switch the current buffer to lpf-mode.\\
none & \texttt{M-x lua-\codebreak{}mode} & \texttt{lua-\codebreak{}mode} & Switch the current buffer to lua-mode.\\
none & \texttt{M-x luau-\codebreak{}mode} & \texttt{luau-\codebreak{}mode} & Switch the current buffer to luau-mode.\\
none & \texttt{M-x makefile-\codebreak{}mode} & \texttt{makefile-\codebreak{}mode} & Switch the current buffer to makefile-mode.\\
none & \texttt{M-x markdoc-\codebreak{}mode} & \texttt{markdoc-\codebreak{}mode} & Switch the current buffer to markdoc-mode.\\
\bottomrule
\end{xltabular}

\begin{xltabular}{\textwidth}{@{}l >{\raggedright\arraybackslash\hsize=0.870\hsize}X l >{\raggedright\arraybackslash\hsize=1.130\hsize}X@{}}
\toprule
\textbf{Modal} & \textbf{Non-modal} & \textbf{Command} & \textbf{What it does}\\
\midrule
\endhead
none & \texttt{M-x markdown-\codebreak{}mode} & \texttt{markdown-\codebreak{}mode} & Switch the current buffer to markdown-mode.\\
none & \texttt{M-x matlab-\codebreak{}mode} & \texttt{matlab-\codebreak{}mode} & Switch the current buffer to matlab-mode.\\
none & \texttt{M-x mediawiki-\codebreak{}mode} & \texttt{mediawiki-\codebreak{}mode} & Switch the current buffer to mediawiki-mode.\\
none & \texttt{M-x mermaid-\codebreak{}mode} & \texttt{mermaid-\codebreak{}mode} & Switch the current buffer to mermaid-mode.\\
none & \texttt{M-x meson-\codebreak{}mode} & \texttt{meson-\codebreak{}mode} & Switch the current buffer to meson-mode.\\
none & \texttt{M-x mojo-\codebreak{}mode} & \texttt{mojo-\codebreak{}mode} & Switch the current buffer to mojo-mode.\\
none & \texttt{M-x move-\codebreak{}mode} & \texttt{move-\codebreak{}mode} & Switch the current buffer to move-mode.\\
none & \texttt{M-x nasm-\codebreak{}mode} & \texttt{nasm-\codebreak{}mode} & Switch the current buffer to nasm-mode.\\
none & \texttt{M-x nearley-\codebreak{}mode} & \texttt{nearley-\codebreak{}mode} & Switch the current buffer to nearley-mode.\\
none & \texttt{M-x nestedtext-\codebreak{}mode} & \texttt{nestedtext-\codebreak{}mode} & Switch the current buffer to nestedtext-mode.\\
none & \texttt{M-x nginx-\codebreak{}mode} & \texttt{nginx-\codebreak{}mode} & Switch the current buffer to nginx-mode.\\
none & \texttt{M-x nickel-\codebreak{}mode} & \texttt{nickel-\codebreak{}mode} & Switch the current buffer to nickel-mode.\\
none & \texttt{M-x nim-\codebreak{}mode} & \texttt{nim-\codebreak{}mode} & Switch the current buffer to nim-mode.\\
none & \texttt{M-x nix-\codebreak{}mode} & \texttt{nix-\codebreak{}mode} & Switch the current buffer to nix-mode.\\
none & \texttt{M-x nushell-\codebreak{}mode} & \texttt{nushell-\codebreak{}mode} & Switch the current buffer to nushell-mode.\\
none & \texttt{M-x ocaml-\codebreak{}interface-\codebreak{}mode} & \texttt{ocaml-\codebreak{}interface-\codebreak{}mode} & Switch the current buffer to ocaml-interface-mode.\\
none & \texttt{M-x ocaml-\codebreak{}mode} & \texttt{ocaml-\codebreak{}mode} & Switch the current buffer to ocaml-mode.\\
none & \texttt{M-x odin-\codebreak{}mode} & \texttt{odin-\codebreak{}mode} & Switch the current buffer to odin-mode.\\
none & \texttt{M-x ohm-\codebreak{}mode} & \texttt{ohm-\codebreak{}mode} & Switch the current buffer to ohm-mode.\\
none & \texttt{M-x org-\codebreak{}mode} & \texttt{org-\codebreak{}mode} & Switch the current buffer to org-mode.\\
none & \texttt{M-x pascal-\codebreak{}mode} & \texttt{pascal-\codebreak{}mode} & Switch the current buffer to pascal-mode.\\
none & \texttt{M-x pem-\codebreak{}mode} & \texttt{pem-\codebreak{}mode} & Switch the current buffer to pem-mode.\\
none & \texttt{M-x penrose-\codebreak{}mode} & \texttt{penrose-\codebreak{}mode} & Switch the current buffer to penrose-mode.\\
none & \texttt{M-x perl-\codebreak{}mode} & \texttt{perl-\codebreak{}mode} & Switch the current buffer to perl-mode.\\
none & \texttt{M-x pest-\codebreak{}mode} & \texttt{pest-\codebreak{}mode} & Switch the current buffer to pest-mode.\\
none & \texttt{M-x php-\codebreak{}mode} & \texttt{php-\codebreak{}mode} & Switch the current buffer to php-mode.\\
none & \texttt{M-x picat-\codebreak{}mode} & \texttt{picat-\codebreak{}mode} & Switch the current buffer to picat-mode.\\
none & \texttt{M-x pip-\codebreak{}requirements-\codebreak{}mode} & \texttt{pip-\codebreak{}requirements-\codebreak{}mode} & Switch the current buffer to pip-requirements-mode.\\
none & \texttt{M-x pkl-\codebreak{}mode} & \texttt{pkl-\codebreak{}mode} & Switch the current buffer to pkl-mode.\\
none & \texttt{M-x po-\codebreak{}mode} & \texttt{po-\codebreak{}mode} & Switch the current buffer to po-mode.\\
none & \texttt{M-x pod-\codebreak{}mode} & \texttt{pod-\codebreak{}mode} & Switch the current buffer to pod-mode.\\
none & \texttt{M-x poly-\codebreak{}markdown+r-\codebreak{}mode} & \texttt{poly-\codebreak{}markdown+r-\codebreak{}mode} & Switch the current buffer to poly-markdown+r-mode.\\
none & \texttt{M-x ponylang-\codebreak{}mode} & \texttt{ponylang-\codebreak{}mode} & Switch the current buffer to ponylang-mode.\\
none & \texttt{M-x powershell-\codebreak{}mode} & \texttt{powershell-\codebreak{}mode} & Switch the current buffer to powershell-mode.\\
none & \texttt{M-x prisma-\codebreak{}mode} & \texttt{prisma-\codebreak{}mode} & Switch the current buffer to prisma-mode.\\
none & \texttt{M-x prog-\codebreak{}mode} & \texttt{prog-\codebreak{}mode} & Switch the current buffer to prog-mode.\\
none & \texttt{M-x prolog-\codebreak{}mode} & \texttt{prolog-\codebreak{}mode} & Switch the current buffer to prolog-mode.\\
none & \texttt{M-x protobuf-\codebreak{}mode} & \texttt{protobuf-\codebreak{}mode} & Switch the current buffer to protobuf-mode.\\
none & \texttt{M-x proverif-\codebreak{}mode} & \texttt{proverif-\codebreak{}mode} & Switch the current buffer to proverif-mode.\\
none & \texttt{M-x prql-\codebreak{}mode} & \texttt{prql-\codebreak{}mode} & Switch the current buffer to prql-mode.\\
\bottomrule
\end{xltabular}

\begin{xltabular}{\textwidth}{@{}l >{\raggedright\arraybackslash\hsize=0.800\hsize}X l >{\raggedright\arraybackslash\hsize=1.200\hsize}X@{}}
\toprule
\textbf{Modal} & \textbf{Non-modal} & \textbf{Command} & \textbf{What it does}\\
\midrule
\endhead
none & \texttt{M-x ptx-\codebreak{}mode} & \texttt{ptx-\codebreak{}mode} & Switch the current buffer to ptx-mode.\\
none & \texttt{M-x pug-\codebreak{}mode} & \texttt{pug-\codebreak{}mode} & Switch the current buffer to pug-mode.\\
none & \texttt{M-x purescript-\codebreak{}mode} & \texttt{purescript-\codebreak{}mode} & Switch the current buffer to purescript-mode.\\
none & \texttt{M-x python-\codebreak{}mode} & \texttt{python-\codebreak{}mode} & Switch the current buffer to python-mode.\\
none & \texttt{M-x ql-\codebreak{}mode} & \texttt{ql-\codebreak{}mode} & Switch the current buffer to ql-mode.\\
none & \texttt{M-x qml-\codebreak{}mode} & \texttt{qml-\codebreak{}mode} & Switch the current buffer to qml-mode.\\
none & \texttt{M-x quarto-\codebreak{}mode} & \texttt{quarto-\codebreak{}mode} & Switch the current buffer to quarto-mode.\\
none & \texttt{M-x quint-\codebreak{}mode} & \texttt{quint-\codebreak{}mode} & Switch the current buffer to quint-mode.\\
none & \texttt{M-x r-mode} & \texttt{r-mode} & Switch the current buffer to r-mode.\\
none & \texttt{M-x racket-\codebreak{}mode} & \texttt{racket-\codebreak{}mode} & Switch the current buffer to racket-mode.\\
none & \texttt{M-x rego-\codebreak{}mode} & \texttt{rego-\codebreak{}mode} & Switch the current buffer to rego-mode.\\
none & \texttt{M-x rescript-\codebreak{}mode} & \texttt{rescript-\codebreak{}mode} & Switch the current buffer to rescript-mode.\\
none & \texttt{M-x ripple-\codebreak{}mode} & \texttt{ripple-\codebreak{}mode} & Switch the current buffer to ripple-mode.\\
none & \texttt{M-x robot-\codebreak{}mode} & \texttt{robot-\codebreak{}mode} & Switch the current buffer to robot-mode.\\
none & \texttt{M-x robots-\codebreak{}txt-\codebreak{}mode} & \texttt{robots-\codebreak{}txt-\codebreak{}mode} & Switch the current buffer to robots-txt-mode.\\
none & \texttt{M-x ron-\codebreak{}mode} & \texttt{ron-\codebreak{}mode} & Switch the current buffer to ron-mode.\\
none & \texttt{M-x rpm-\codebreak{}spec-\codebreak{}mode} & \texttt{rpm-\codebreak{}spec-\codebreak{}mode} & Switch the current buffer to rpm-spec-mode.\\
none & \texttt{M-x rst-\codebreak{}mode} & \texttt{rst-\codebreak{}mode} & Switch the current buffer to rst-mode.\\
none & \texttt{M-x ruby-\codebreak{}mode} & \texttt{ruby-\codebreak{}mode} & Switch the current buffer to ruby-mode.\\
none & \texttt{M-x rust-\codebreak{}mode} & \texttt{rust-\codebreak{}mode} & Switch the current buffer to rust-mode.\\
none & \texttt{M-x sage-\codebreak{}mode} & \texttt{sage-\codebreak{}mode} & Switch the current buffer to sage-mode.\\
none & \texttt{M-x scad-\codebreak{}mode} & \texttt{scad-\codebreak{}mode} & Switch the current buffer to scad-mode.\\
none & \texttt{M-x scala-\codebreak{}mode} & \texttt{scala-\codebreak{}mode} & Switch the current buffer to scala-mode.\\
none & \texttt{M-x scfg-\codebreak{}mode} & \texttt{scfg-\codebreak{}mode} & Switch the current buffer to scfg-mode.\\
none & \texttt{M-x scheme-\codebreak{}mode} & \texttt{scheme-\codebreak{}mode} & Switch the current buffer to scheme-mode.\\
none & \texttt{M-x sclang-\codebreak{}mode} & \texttt{sclang-\codebreak{}mode} & Switch the current buffer to sclang-mode.\\
none & \texttt{M-x scss-\codebreak{}mode} & \texttt{scss-\codebreak{}mode} & Switch the current buffer to scss-mode.\\
none & \texttt{M-x shellcheckrc-\codebreak{}mode} & \texttt{shellcheckrc-\codebreak{}mode} & Switch the current buffer to shellcheckrc-mode.\\
none & \texttt{M-x slang-\codebreak{}mode} & \texttt{slang-\codebreak{}mode} & Switch the current buffer to slang-mode.\\
none & \texttt{M-x slint-\codebreak{}mode} & \texttt{slint-\codebreak{}mode} & Switch the current buffer to slint-mode.\\
none & \texttt{M-x slisp-\codebreak{}mode} & \texttt{slisp-\codebreak{}mode} & Switch the current buffer to slisp-mode.\\
none & \texttt{M-x smali-\codebreak{}mode} & \texttt{smali-\codebreak{}mode} & Switch the current buffer to smali-mode.\\
none & \texttt{M-x smithy-\codebreak{}mode} & \texttt{smithy-\codebreak{}mode} & Switch the current buffer to smithy-mode.\\
none & \texttt{M-x sml-\codebreak{}mode} & \texttt{sml-\codebreak{}mode} & Switch the current buffer to sml-mode.\\
none & \texttt{M-x snakemake-\codebreak{}mode} & \texttt{snakemake-\codebreak{}mode} & Switch the current buffer to snakemake-mode.\\
none & \texttt{M-x solidity-\codebreak{}mode} & \texttt{solidity-\codebreak{}mode} & Switch the current buffer to solidity-mode.\\
none & \texttt{M-x sourcepawn-\codebreak{}mode} & \texttt{sourcepawn-\codebreak{}mode} & Switch the current buffer to sourcepawn-mode.\\
none & \texttt{M-x spade-\codebreak{}mode} & \texttt{spade-\codebreak{}mode} & Switch the current buffer to spade-mode.\\
none & \texttt{M-x spicedb-\codebreak{}mode} & \texttt{spicedb-\codebreak{}mode} & Switch the current buffer to spicedb-mode.\\
none & \texttt{M-x sql-\codebreak{}mode} & \texttt{sql-\codebreak{}mode} & Switch the current buffer to sql-mode.\\
\bottomrule
\end{xltabular}

\begin{xltabular}{\textwidth}{@{}l >{\raggedright\arraybackslash\hsize=0.857\hsize}X l >{\raggedright\arraybackslash\hsize=1.143\hsize}X@{}}
\toprule
\textbf{Modal} & \textbf{Non-modal} & \textbf{Command} & \textbf{What it does}\\
\midrule
\endhead
none & \texttt{M-x ssh-\codebreak{}config-\codebreak{}mode} & \texttt{ssh-\codebreak{}config-\codebreak{}mode} & Switch the current buffer to ssh-config-mode.\\
none & \texttt{M-x strace-\codebreak{}mode} & \texttt{strace-\codebreak{}mode} & Switch the current buffer to strace-mode.\\
none & \texttt{M-x strictdoc-\codebreak{}mode} & \texttt{strictdoc-\codebreak{}mode} & Switch the current buffer to strictdoc-mode.\\
none & \texttt{M-x styx-\codebreak{}mode} & \texttt{styx-\codebreak{}mode} & Switch the current buffer to styx-mode.\\
none & \texttt{M-x svelte-\codebreak{}mode} & \texttt{svelte-\codebreak{}mode} & Switch the current buffer to svelte-mode.\\
none & \texttt{M-x sway-\codebreak{}mode} & \texttt{sway-\codebreak{}mode} & Switch the current buffer to sway-mode.\\
none & \texttt{M-x swift-\codebreak{}mode} & \texttt{swift-\codebreak{}mode} & Switch the current buffer to swift-mode.\\
none & \texttt{M-x systemd-\codebreak{}mode} & \texttt{systemd-\codebreak{}mode} & Switch the current buffer to systemd-mode.\\
none & \texttt{M-x systemverilog-\codebreak{}mode} & \texttt{systemverilog-\codebreak{}mode} & Switch the current buffer to systemverilog-mode.\\
none & \texttt{M-x t32-\codebreak{}mode} & \texttt{t32-\codebreak{}mode} & Switch the current buffer to t32-mode.\\
none & \texttt{M-x tablegen-\codebreak{}mode} & \texttt{tablegen-\codebreak{}mode} & Switch the current buffer to tablegen-mode.\\
none & \texttt{M-x tact-\codebreak{}mode} & \texttt{tact-\codebreak{}mode} & Switch the current buffer to tact-mode.\\
none & \texttt{M-x task-\codebreak{}mode} & \texttt{task-\codebreak{}mode} & Switch the current buffer to task-mode.\\
none & \texttt{M-x tcl-\codebreak{}mode} & \texttt{tcl-\codebreak{}mode} & Switch the current buffer to tcl-mode.\\
none & \texttt{M-x teal-\codebreak{}mode} & \texttt{teal-\codebreak{}mode} & Switch the current buffer to teal-mode.\\
none & \texttt{M-x templ-\codebreak{}mode} & \texttt{templ-\codebreak{}mode} & Switch the current buffer to templ-mode.\\
none & \texttt{M-x tera-\codebreak{}mode} & \texttt{tera-\codebreak{}mode} & Switch the current buffer to tera-mode.\\
none & \texttt{M-x text-\codebreak{}mode} & \texttt{text-\codebreak{}mode} & Switch the current buffer to text-mode.\\
none & \texttt{M-x textproto-\codebreak{}mode} & \texttt{textproto-\codebreak{}mode} & Switch the current buffer to textproto-mode.\\
none & \texttt{M-x thrift-\codebreak{}mode} & \texttt{thrift-\codebreak{}mode} & Switch the current buffer to thrift-mode.\\
none & \texttt{M-x tla-\codebreak{}mode} & \texttt{tla-\codebreak{}mode} & Switch the current buffer to tla-mode.\\
none & \texttt{M-x todotxt-\codebreak{}mode} & \texttt{todotxt-\codebreak{}mode} & Switch the current buffer to todotxt-mode.\\
none & \texttt{M-x tolk-\codebreak{}mode} & \texttt{tolk-\codebreak{}mode} & Switch the current buffer to tolk-mode.\\
none & \texttt{M-x toml-\codebreak{}mode} & \texttt{toml-\codebreak{}mode} & Switch the current buffer to toml-mode.\\
none & \texttt{M-x tql-\codebreak{}mode} & \texttt{tql-\codebreak{}mode} & Switch the current buffer to tql-mode.\\
none & \texttt{M-x tsv-\codebreak{}mode} & \texttt{tsv-\codebreak{}mode} & Switch the current buffer to tsv-mode.\\
none & \texttt{M-x tsx-\codebreak{}mode} & \texttt{tsx-\codebreak{}mode} & Switch the current buffer to tsx-mode.\\
none & \texttt{M-x twig-\codebreak{}mode} & \texttt{twig-\codebreak{}mode} & Switch the current buffer to twig-mode.\\
none & \texttt{M-x typescript-\codebreak{}mode} & \texttt{typescript-\codebreak{}mode} & Switch the current buffer to typescript-mode.\\
none & \texttt{M-x typespec-\codebreak{}mode} & \texttt{typespec-\codebreak{}mode} & Switch the current buffer to typespec-mode.\\
none & \texttt{M-x typst-\codebreak{}mode} & \texttt{typst-\codebreak{}mode} & Switch the current buffer to typst-mode.\\
none & \texttt{M-x ungrammar-\codebreak{}mode} & \texttt{ungrammar-\codebreak{}mode} & Switch the current buffer to ungrammar-mode.\\
none & \texttt{M-x unison-\codebreak{}mode} & \texttt{unison-\codebreak{}mode} & Switch the current buffer to unison-mode.\\
none & \texttt{M-x uxntal-\codebreak{}mode} & \texttt{uxntal-\codebreak{}mode} & Switch the current buffer to uxntal-mode.\\
none & \texttt{M-x v-mode} & \texttt{v-mode} & Switch the current buffer to v-mode.\\
none & \texttt{M-x vala-\codebreak{}mode} & \texttt{vala-\codebreak{}mode} & Switch the current buffer to vala-mode.\\
none & \texttt{M-x vento-\codebreak{}mode} & \texttt{vento-\codebreak{}mode} & Switch the current buffer to vento-mode.\\
none & \texttt{M-x verilog-\codebreak{}mode} & \texttt{verilog-\codebreak{}mode} & Switch the current buffer to verilog-mode.\\
none & \texttt{M-x vhdl-\codebreak{}mode} & \texttt{vhdl-\codebreak{}mode} & Switch the current buffer to vhdl-mode.\\
none & \texttt{M-x vhs-\codebreak{}mode} & \texttt{vhs-\codebreak{}mode} & Switch the current buffer to vhs-mode.\\
\bottomrule
\end{xltabular}

\begin{xltabular}{\textwidth}{@{}l >{\raggedright\arraybackslash\hsize=0.694\hsize}X l >{\raggedright\arraybackslash\hsize=1.306\hsize}X@{}}
\toprule
\textbf{Modal} & \textbf{Non-modal} & \textbf{Command} & \textbf{What it does}\\
\midrule
\endhead
none & \texttt{M-x vimrc-\codebreak{}mode} & \texttt{vimrc-\codebreak{}mode} & Switch the current buffer to vimrc-mode.\\
none & \texttt{M-x vue-\codebreak{}mode} & \texttt{vue-\codebreak{}mode} & Switch the current buffer to vue-mode.\\
none & \texttt{M-x wast-\codebreak{}mode} & \texttt{wast-\codebreak{}mode} & Switch the current buffer to wast-mode.\\
none & \texttt{M-x wat-\codebreak{}mode} & \texttt{wat-\codebreak{}mode} & Switch the current buffer to wat-mode.\\
none & \texttt{M-x werk-\codebreak{}mode} & \texttt{werk-\codebreak{}mode} & Switch the current buffer to werk-mode.\\
none & \texttt{M-x wesl-\codebreak{}mode} & \texttt{wesl-\codebreak{}mode} & Switch the current buffer to wesl-mode.\\
none & \texttt{M-x wgsl-\codebreak{}mode} & \texttt{wgsl-\codebreak{}mode} & Switch the current buffer to wgsl-mode.\\
none & \texttt{M-x wit-\codebreak{}mode} & \texttt{wit-\codebreak{}mode} & Switch the current buffer to wit-mode.\\
none & \texttt{M-x woodpecker-\codebreak{}ci-\codebreak{}mode} & \texttt{woodpecker-\codebreak{}ci-\codebreak{}mode} & Switch the current buffer to woodpecker-ci-mode.\\
none & \texttt{M-x wren-\codebreak{}mode} & \texttt{wren-\codebreak{}mode} & Switch the current buffer to wren-mode.\\
none & \texttt{M-x xit-\codebreak{}mode} & \texttt{xit-\codebreak{}mode} & Switch the current buffer to xit-mode.\\
none & \texttt{M-x xml-\codebreak{}mode} & \texttt{xml-\codebreak{}mode} & Switch the current buffer to xml-mode.\\
none & \texttt{M-x xtc-\codebreak{}mode} & \texttt{xtc-\codebreak{}mode} & Switch the current buffer to xtc-mode.\\
none & \texttt{M-x yaml-\codebreak{}mode} & \texttt{yaml-\codebreak{}mode} & Switch the current buffer to yaml-mode.\\
none & \texttt{M-x yara-\codebreak{}mode} & \texttt{yara-\codebreak{}mode} & Switch the current buffer to yara-mode.\\
none & \texttt{M-x yuck-\codebreak{}mode} & \texttt{yuck-\codebreak{}mode} & Switch the current buffer to yuck-mode.\\
none & \texttt{M-x zig-\codebreak{}mode} & \texttt{zig-\codebreak{}mode} & Switch the current buffer to zig-mode.\\
\bottomrule
\end{xltabular}


\end{document}
